\documentclass[a4paper]{article}

\usepackage{etex}
\usepackage[utf8]{inputenc}
\usepackage[T1]{fontenc}	 
\usepackage{lmodern} 

\usepackage[frenchb]{babel}
%
\usepackage{amsmath,amsfonts}
%\usepackage{amssymb}
%\usepackage{pdfpages}  
\usepackage{textcomp}
\usepackage{makeidx}
\usepackage{pst-all}

\usepackage{graphicx}
\usepackage{hyperref}
\usepackage{pst-fr3d}
\usepackage{infix-RPN}
\usepackage{pst-infixplot}

\usepackage{pst-eucl}
\usepackage{pst-slpe}
\usepackage{wrapfig}
\usepackage{pst-solides3d}
\usepackage{pst-3dplot}
\usepackage{pst-fun}
\usepackage{pst-func}
\usepackage{pst-bezier}
\usepackage{pst-bspline}
\usepackage{animate}
\usepackage{makeidx}
\usepackage{pst-tools}
\usepackage{pst-poly}
%\usepackage{pst-gantt}

%\usepackage{showidx}
%\usepackage{pst-abspos}
%====================================================================

\makeindex

\newcommand{\AC}[1]{\{#1\}}

\newcommand{\BS}[1]{$\backslash$#1}

%\newcommand{\BSS}[1]{\textbf{\cyan $\backslash${#1}}\index{Commandes!\textbackslash{}#1}}
\newcommand{\BSS}[1]{\textbf{\textbackslash{#1}}\index{\textbf{1) Commandes}!#1 @\textbackslash{}#1}}

\newcommand{\BSI}[2]{ \index{\textbf{4) Par modules }!\textbf{#2}!#1@\textbackslash{}#1 (M)}}

%\newcommand{\BSI}[2]{\cite{#2} \index{\textbf{4) Par modules }!\textbf{#2}!#1@\textbackslash{}#1 (M)}}

\newcommand{\RDI}[2]{ \index{\textbf{4) Par modules }!\textbf{#2}!#1@#1 (P)}}

%\newcommand{\RDI}[2]{\cite{#2} \index{\textbf{4) Par modules }!\textbf{#2}!#1@#1 (P)}}

\newcommand{\BDI}[2]{\cite{#2} \index{\textbf{4) Par modules }!\textbf{#2}!#1@#1 (V) }}

\newcommand{\RDD}[1]{{\red #1}\index{\textbf{2) Paramètres et options}!#1}}

\newcommand{\RDDD}[2]{{\red #1}\index{\textbf{2) Paramètres et options}!#2!#1}}

\newcommand{\BDD}[1]{{\blue #1}\index{\textbf{3) Variables PsTricks}!#1}}

\newcommand{\DFR}{\psset{unit=.25cm,fillstyle=solid,linewidth=0pt}\begin{pspicture*}(3,1.5) \psframe[fillcolor=blue](1,1.5)\psframe[fillcolor=white](1,0)(2,1.5)\psframe[fillcolor=red](2,0)(3,1.5)\end{pspicture*}
\psset{unit=1cm,fillstyle=solid,linewidth=1pt}}

\newcommand{\DGB}{\psset{unit=.25cm,fillstyle=solid,linewidth=0pt}\begin{pspicture*}(3,1.5) \psframe[fillcolor=blue](3,1.5) \psline[linecolor=white,linewidth=.1cm](3,1.5)\psline[linecolor=white,linewidth=.1cm](0,1.5)(3,0) \psline[linecolor=white,linewidth=.1cm](0,.75)(3,.75)\psline[linecolor=white,linewidth=.1cm](1.5,0)(1.5,1.5)
\psline[linecolor=red,linewidth=.05cm](3,1.5) \psline[linecolor=red,linewidth=.05cm](0,1.5)(3,0)
\psline[linecolor=red,linewidth=.05cm](0,.75)(3,.75) \psline[linecolor=red,linewidth=.05cm](1.5,0)(1.5,1.5) \end{pspicture*}

\psset{unit=1cm,fillstyle=solid,linewidth=1pt}}

\newcommand{\maboite}[1]{\begin{center} \psframebox[fillcolor=yellow,fillstyle=solid]{ \parbox{.5\linewidth} {\centering #1}} \end{center}}

\newcommand{\mydessin}[1]{%
 \begin{pspicture}(-5,-3)(5,5)
\psset{unit=.5cm}
\psset{solidmemory}
\psset{viewpoint=50 #1  20 rtp2xyz, Decran=50}
\psSolid[object=cylindrecreux,h=10,r=2,fillcolor=white,mode=4,name=A1,action=none,incolor=green!50,SphericalCoor,viewpoint=50 30 20](0,0,-3)%
%
\psSolid[object=conecreux,h=15,r=2,RotY=-60,fillcolor=white,incolor=red!50,action=none,mode=5,name=B1,SphericalCoor,viewpoint=50 30 20](4,0,0)%
\psSolid[object=fusion,action=draw**,base=A1 B1,SphericalCoor,](0,0,0)%
\composeSolid%
 \end{pspicture}}
 
\newcommand{\essais}[1]{}

\newcommand{\MonMotif}{%
\begin{pspicture}(1,.5)
 \psframe[dimen=middle,fillcolor=yellow,fillstyle=solid,linecolor=blue](1,.5)
 \pscircle[dimen=middle,fillcolor=magenta,fillstyle=solid,linecolor=magenta](.5,.25){.1}
\end{pspicture}}






\newcommand{\xx}{\unskip\kern.05em:} 

% Version française

\newcommand{\TFRGB}[2]{#1} % #1 en francais #2 in english
\newcommand{\dft}{Par défaut}
\newcommand{\SSCT}[2]{\section{#1}}
\newcommand{\SbSSCT}[2]{\subsection{#1}}
\newcommand{\SbSbSSCT}[2]{\subsubsection{#1}}
\newcommand{\SSCTTC}[4]{\section[#1]{#2}}
\newcommand{\SbSSCTTC}[4]{\subsection[#1]{#2}}
\newcommand{\SbSbSSCTTC}[4]{\subsubsection[#1]{#2}}
 % version française
%%English version
\newcommand{\TFRGB}[2]{#2} % #1 en francais #2 in english

\newcommand{\dft}{By default}
\newcommand{\SSCT}[2]{ \section{#2}}
\newcommand{\SbSSCT}[2]{ \subsection{#2}}
\newcommand{\SbSbSSCT}[2]{ \subsubsection{#2}}
\newcommand{\SSCTTC}[4]{ \section[#2]{#2}}
\newcommand{\SbSSCTTC}[4]{ \subsection[#2]{#2}}
\newcommand{\SbSbSSCTTC}[4]{ \subsubsection[#3]{#4}} % english version

 %======================================

 \begin{document}
  \newrgbcolor{Vert}{0 .8 0}
 
  
\author{{\Huge Jean Pierre Casteleyn } \\ {\Huge IUT GTE } \\ {\Huge Dunkerque, France }}

\DeclareFixedFont{\RM}{T1}{ptm}{b}{n}{2cm}

\DeclareFixedFont{\RMM}{T1}{ptm}{b}{n}{1cm}


\TFRGB{
\title{ \pscharpath[fillstyle=gradient,gradbegin=red,gradend=cyan,shadow=true,doubleline=true]{\RM Visuel PSTricks} \\ \vspace{1cm} \pscharpath[fillstyle=gradient,gradbegin=red,gradend=cyan,shadow=true,doubleline=true]{\RMM Version 2.30}\\ \vspace{1cm}%
\includegraphics[width=6cm]{logoiut.eps} %
}}
{\title{ \pscharpath[fillstyle=gradient,gradbegin=red,gradend=cyan,shadow=true,doubleline=true]{\RM Visual PSTricks} \\ \vspace{1cm} \pscharpath[fillstyle=gradient,gradbegin=red,gradend=cyan,shadow=true,doubleline=true]{\RMM Version 2.30}\\ \vspace{1cm}%
\includegraphics[width=6cm]{logoiut.eps} %
}}



\date{{\LARGE mis à jour le \today}}

 

 
 \maketitle



\newpage


\TFRGB{
\textbf{Objectifs }: 

\begin{itemize}
\item Avoir une image par  commande ou par paramètre.
\item Avoir un texte réduit au strict minimum.
\item être le plus complet possible.
\end{itemize} 
}
{\textbf{Objectives }: 

\begin{itemize}
\item an image per command or parameter.
\item the minimum text possible.
\item the most compleate possible.
\end{itemize}}
\bigskip

\psset{llx=-.5cm,lly=-.5cm,urx=.2cm,ury=.2cm,xAxisLabel={}, yAxisLabel={}}
\psset{tickcolor=black!10}

\TFRGB{\textbf{Légende} :}{\textbf{Légend}}

\psset{llx=-.2cm,lly=-.2cm,urx=.2cm,ury=.2cm,xAxisLabel={}, yAxisLabel={}}
 
\begin{tabular}{|c|l|} \hline
\begin{psgraph}[axesstyle=none,xticksize=0 2,yticksize=0 2,subticks=0,shift=*,labelsep = 2pt,labelFontSize=\scriptscriptstyle](0,0)(2,2){1cm}{1cm}
\dotnode[dotstyle=*,linecolor=blue,dotscale=2](1,1){A} \nput{45}{A}{A} 
 \end{psgraph}
&  
\TFRGB{un n\oe ud de base}{Basic node}
\\ \hline  
\begin{psgraph}[axesstyle=none,xticksize=0 2,yticksize=0 2,subticks=0,shift=*,labelsep = 2pt,labelFontSize=\scriptscriptstyle](0,0)(2,2){1cm}{1cm}
\dotnode[dotstyle=*,linecolor=red,dotscale=2](1,1){A} \nput{45}{A}{A} 
 \end{psgraph}
&  
\TFRGB{un n\oe ud calculé}{Calculated node}
\\ \hline
\begin{psgraph}[axesstyle=none,xticksize=0 2,yticksize=0 2,subticks=0,shift=*,labelsep = 2pt,labelFontSize=\scriptscriptstyle](0,0)(2,2){1cm}{1cm}
\dotnode[dotstyle=*,dotscale=2,linecolor=blue](1,1){A}
% Asterisk 
 \end{psgraph}
&  
\TFRGB{un point}{a point}
\\ \hline  
\begin{psgraph}[axesstyle=none,xticksize=0 2,yticksize=0 2,subticks=0,shift=*,labelsep = 2pt,labelFontSize=\scriptscriptstyle](0,0)(2,2){1cm}{1cm}
\psline[linewidth=2pt,linecolor=blue](2,2)
 \end{psgraph}
&  
\TFRGB{un élément de base}[Base element] 
\\ \hline 
\begin{psgraph}[axesstyle=none,xticksize=0 2,yticksize=0 2,subticks=0,shift=*,labelsep = 2pt,labelFontSize=\scriptscriptstyle](0,0)(2,2){1cm}{1cm}
\psline[linewidth=2pt,linecolor=red](2,2)
 \end{psgraph}
&  
\TFRGB{un élément supplémentaire}{Additional element} 
\\ \hline 
\begin{psgraph}[axesstyle=none,xticksize=0 2,yticksize=0 2,subticks=0,shift=*,labelsep = 2pt,labelFontSize=\scriptscriptstyle](0,0)(2,2){1cm}{1cm}
\psline[linewidth=2pt,linecolor=green](2,2)
\psline[linewidth=2pt,linecolor=cyan](1,0)(2,1)
 \end{psgraph}
&  
\TFRGB{autres éléments supplémentaires}{Other additional element} 
\\ \hline  
\begin{psgraph}[axesstyle=none,xticksize=0 2,yticksize=0 2,subticks=0,shift=*,labelsep = 2pt,labelFontSize=\scriptscriptstyle](0,0)(2,2){1cm}{1cm}
\psline[linestyle=dashed](2,2)
\psline[linestyle=dotted](1,0)(2,1)
 \end{psgraph}
&  
\TFRGB{construction pour expliquer une commande, une option ou un paramètre}{to  highlight the command, the option oe a parameter }
\\ \hline  
\begin{psgraph}[axesstyle=none,xticksize=0 2,yticksize=0 2,subticks=0,shift=*,labelsep = 2pt,labelFontSize=\scriptscriptstyle](0,0)(2,2){1cm}{1cm}
\psframe[fillstyle=solid,fillcolor=yellow](.5,.5)(1.5,1.5)

 \end{psgraph}
&  
\TFRGB{couleur de remplissage par défaut  : blanc}{Filling color (\dft : white)}
\\ \hline 
\end{tabular}  

\psset{llx=-.5cm,lly=-.5cm,urx=.2cm,ury=.2cm,xAxisLabel={}, yAxisLabel={}} 



\bigskip

\TFRGB{
\textbf{Vous pouvez me contacter à}
 \href{mailto:jpcdk@yahoo.fr}{mon e-mail personnel} pour

\begin{itemize}
\item me signaler les erreurs que vous avez constatés 
\item me faire part de vos commentaires, suggestions \dots
\end{itemize}}
{
\textbf{You can contact me at }
 \href{mailto:jpcdk@yahoo.fr}{my personal email} to

\begin{itemize}
\item let me know the mistakes found
\item give me your commentaries, your suggestions \dots
\end{itemize}}
 


\bigskip

\textbf{\TFRGB{Merci à }{Thanks to}:}

Alain Bécue ,
Denis Bitouzé,
Jean Côme charpentier,
Martin Giese,
Denis Girou,
Alexander Grahn,
Christophe Jorssen,
Dr. Uwe Kern,
Manuel Luque,
Dominique Rodriguez,
Michael Sharpe,
Tobias Nähring,
Herbert Voß,
Timothy Van Zandt.

\newpage
\TFRGB{
\textbf{Quoi de neuf dans cette mise à jour}

\begin{itemize}
\item Ajout des modules :
\begin{itemize}
\item pst-bezier : voir page \pageref{bezier}
\item  pst-fun : voir page \pageref{fun}
\item    pst-func : voir page \pageref{func}
\end{itemize}
\item différentes mises à jour du module pstricks-add : 
\begin{itemize}
\item pscspline  : voir page \pageref{pscspline}
\item  psellipseAB  : voir page \pageref{psellipseAB}
\item  Notion de chemin PSTricks : voir page \pageref{chem}
\item Commenter un élément : voir page \pageref{com}
\item Homothétie : voir page \pageref{hom}
\item Commande psrotate : voir page \pageref{rot}
\item Annuler des objets : voir page \pageref{ann}
\item Des dés : voir page \pageref{des}
\item paramètres d’un graphe en barres : voir page \pageref{bar}
\item Options VarStep et VarStepEpsilon  : voir page \pageref{var}
\item Macro psVectorfield : voir page \pageref{vec}
\end{itemize}
\end{itemize}
}{}
 %ù==========================================================
 
\addtolength{\hoffset}{-1.5cm}

\setcounter{tocdepth}{1}
%\setcounter{tocdepth}{3}
 \tableofcontents
 
 
 \newpage

\SSCT{Les figures de base}{basic figures}
%\section{Les figures de base}

\input{pst1}

\newpage



\SSCT{Les paramètres disponibles}{Parameters avaliable}

\SbSSCT{Epaisseur du trait}{Linewidth}
\psset{linecolor=blue}


\begin{tabular}{|c|c|} \hline  
\begin{psgraph}[axesstyle=none,xticksize=0 1,yticksize=0 4,subticks=0,shift=*](0,0)(4,1){4cm}{1cm} 
\psline[linewidth=10mm](2,0)(2,1) 
 \end{psgraph}
 &
\parbox[c]{10cm}{  
\BS{psline}[\RDD{linewidth}=10mm](2,0)(2,1)  \\
}
\\ \hline 
\begin{psgraph}[axesstyle=none,xticksize=0 1,yticksize=0 4,subticks=0,shift=*](0,0)(4,1){4cm}{1cm}  
\psline[linewidth=1cm](2,0)(2,1) 
 \end{psgraph}
 &
\parbox[c]{10cm}{  
\BS{psline}[\RDD{linewidth}=1cm](2,0)(2,1) \\
}
\\ \hline 
\begin{psgraph}[axesstyle=none,xticksize=0 1,yticksize=0 4,subticks=0,shift=*](0,0)(4,1){4cm}{1cm} 
\psline[linewidth=1in](2,0)(2,1) 
 \end{psgraph}
 &
\parbox[c]{10cm}{  
\BS{psline}[\RDD{linewidth}=1in](2,0)(2,1) \\
}
\\ \hline 
\begin{psgraph}[axesstyle=none,xticksize=0 1,yticksize=0 4,subticks=0,shift=*](0,0)(4,1){4cm}{1cm} 
\psline[linewidth=10pt](2,0)(2,1) 
 \end{psgraph}
 &
\parbox[c]{10cm}{  
\BS{psline}[\RDD{linewidth}=10pt](2,0)(2,1) \\
}
\\ \hline
\multicolumn{2}{|c|}{ {\blue \dft{} : linewidth = 0.8pt} } \\ \hline
\end{tabular}

\bigskip



\begin{tabular}{|c|c|} \hline 
\multicolumn{2}{|c|}{{\emph Dimensions en fonction de la taille de la police }} \\  \hline
\begin{psgraph}[axesstyle=none,xticksize=0 1,yticksize=0 4,subticks=0,shift=*](0,0)(4,1){4cm}{1cm} 
\psline[linewidth=1em](2,0)(2,1) 
\rput[c](2,0.5){\green  m}
 \end{psgraph}
 &
\parbox[c]{10cm}{  
\BS{psline}[\RDD{linewidth}=1em](2,0)(2,1) \\
}
\\ \hline 
\begin{psgraph}[axesstyle=none,xticksize=0 1,yticksize=0 4,subticks=0,shift=*](0,0)(4,1){4cm}{1cm} 
{\Huge \psline[linewidth=1em](2,0)(2,1) 
\rput[c](2,0.5){\green  m}}
 \end{psgraph}
 &
\parbox[c]{10cm}{  
\AC{\BSS{Huge} \BS{psline}[\RDD{linewidth}=1em](2,0)(2,1)  }\\
}
\\ \hline
\begin{psgraph}[axesstyle=none,xticksize=0 1,yticksize=0 4,subticks=0,shift=*](0,0)(4,1){4cm}{1cm} 
\psline[linewidth=1ex](0,0.5)(4,0.5) 
\rput[c](2,0.5){\green  x}
 \end{psgraph}
 &
\parbox[c]{10cm}{  
\BS{psline}[\RDD{linewidth}=1ex](0,0.5)(4,0.5)  \\
}
\\ \hline
\begin{psgraph}[axesstyle=none,xticksize=0 1,yticksize=0 4,subticks=0,shift=*](0,0)(4,1){4cm}{1cm} 
{\Huge \psline[linewidth=1ex](0,0.5)(4,0.5) 
\rput[c](2,0.5){\green x}}
 \end{psgraph}
 &
\parbox[c]{10cm}{  
\AC{\BSS{Huge} \BS{psline}[\RDD{linewidth}=1ex](0,0.5)(4,0.5)   }\\
}
\\ \hline  
\end{tabular}


%-------------------------------------------------
\SbSSCT{Couleur de ligne}{Line color}


\begin{tabular}{|c|c|} \hline  
\begin{psgraph}[axesstyle=none,xticksize=0 1,yticksize=0 4,subticks=0,shift=*](0,-.5)(4,.5){4cm}{1cm}  
 \psline[linewidth=0.5cm,linecolor=green ](4,0)
 \end{psgraph}
 &
\parbox[c]{10cm}{  
\BS{psline} [linewidth=0.5cm,\RDD{linecolor}=green ] (4,0)\\
}
\\ \hline 
\end{tabular} 

\psset{unit=.8cm}
\smallskip

\begin{tabular}{|l|l|l|l|p{.8cm}|p{.8cm}|l|l|l|l|l|} \hline

\multicolumn{10}{|c|}{  \TFRGB{{\emph Couleurs disponibles}}{\emph Colors available} }  \\
\hline  black & darkgray & gray  & lightgray  &white  &  red & green & blue & cyan &  magenta   \\ \hline
\rule[-.3cm]{0pt}{.6cm}
\psline[linewidth=10pt, linecolor=black ](1,0)&
\psline[linewidth=10pt, linecolor= darkgray](1,0) &
\psline[linewidth=10pt, linecolor=gray ](1,0) &
\psline[linewidth=10pt, linecolor=lightgray ](1,0)  &
\psline[linewidth=10pt, linecolor=white ](1,0)  &
\psline[linewidth=10pt, linecolor= red ](1,0)  &
\psline[linewidth=10pt, linecolor= green ](1,0) & 
\psline[linewidth=10pt, linecolor=blue ](1,0) & 
\psline[linewidth=10pt, linecolor= cyan ](1,0) &
\psline[linewidth=10pt, linecolor=magenta ](1,0) \\  \hline  
brown & lime & olive & orange & pink  & purple & teal &violet & yellow &  \\ \hline  
\rule[-.3cm]{0pt}{.6cm}
\psline[linewidth=10pt, linecolor=brown](1,0)&
\psline[linewidth=10pt, linecolor= lime](1,0) &
\psline[linewidth=10pt, linecolor=olive](1,0) &
\psline[linewidth=10pt, linecolor=orange ](1,0)  &
\psline[linewidth=10pt, linecolor=pink ](1,0)  &
\psline[linewidth=10pt, linecolor= purple ](1,0)  &
\psline[linewidth=10pt, linecolor=teal ](1,0) &
\psline[linewidth=10pt, linecolor=violet ](1,0) &
\psline[linewidth=10pt, linecolor= yellow ](1,0) &
\psline[linewidth=10pt, linecolor=white ](1,0) \\  \hline  
\multicolumn{10}{|c|}{ {\blue \dft : linecolor = black }} \\   \hline  
\end{tabular} 
%___________________________________________________________________

\SbSSCT{Styles de ligne}{Line style}

\begin{tabular}{|c|c|} \hline  
\begin{psgraph}[axesstyle=none,xticksize=0 1,yticksize=0 4,subticks=0,shift=*](0,-.5)(4,.5){4cm}{1cm} 
\psline[linewidth=0.5cm,linestyle= dotted ](4,0)
 \end{psgraph}
 &
\parbox[c]{10cm}{  
\BS{psline}[linewidth=0.5cm,\RDD{linestyle}= \RDD{dotted} ](4,0) \\
}
\\ \hline 
\begin{psgraph}[axesstyle=none,xticksize=0 1,yticksize=0 4,subticks=0,shift=*](0,-.5)(4,.5){4cm}{1cm} 
\psline[linewidth=0.5cm,linestyle= dashed  ](4,0)
 \end{psgraph}
 &
\parbox[c]{10cm}{  
\BS{psline}[linewidth=0.5cm,\RDD{linestyle}= \RDD{dashed} ](4,0) \\
}
\\ \hline
\begin{psgraph}[axesstyle=none,xticksize=0 1,yticksize=0 4,subticks=0,shift=*](0,-.5)(4,.5){4cm}{1cm} 
\psline[linewidth=0.5cm,linestyle= none  ](4,0)
 \end{psgraph}
 &
\parbox[c]{10cm}{  
\BS{psline}[linewidth=0.5cm,\RDD{linestyle}= \RDD{none}](4,0) \\
}
\\ \hline
 \multicolumn{2}{|c|}{ {\blue \dft{} : linestyle = solid } }
\\ \hline   
\end{tabular}

\bigskip

\begin{tabular}{|c|c|} \hline 
\begin{psgraph}[axesstyle=none,xticksize=0 1,yticksize=0 4,subticks=0,shift=*](0,-.5)(4,.5){4cm}{1cm} 
\psline[linewidth=0.5cm,linestyle=dotted,dotsep =0.5cm](4,0)
 \end{psgraph}
 &
\parbox[c]{10cm}{  
\BS{psline}[linewidth=0.5cm,linestyle= dotted ,\RDD{dotsep} ={\red 1cm}](4,0) \\
}
\\ \hline
\begin{psgraph}[axesstyle=none,xticksize=0 1,yticksize=0 4,subticks=0,shift=*](0,-.5)(4,.5){4cm}{1cm} 
\psline[linewidth=0.5cm,linestyle=dashed,dash=1cm ](4,0)
 \end{psgraph}
 &
\parbox[c]{10cm}{  
\BS{psline}[linewidth=0.5cm,linestyle= dashed ,\RDD{dash}={\red 1cm} ](4,0) \\
}
\\ \hline 
\begin{psgraph}[axesstyle=none,xticksize=0 1,yticksize=0 4,subticks=0,shift=*](0,-.5)(4,.5){4cm}{1cm} 
\psline[linewidth=0.5cm,linestyle=dashed,dash=1cm 0.5cm](4,0)
 \end{psgraph}
 &
\parbox[c]{10cm}{  
\BS{psline}[linewidth=0.5cm,linestyle= dashed ,\\
\RDD{dash}={\red1cm 0.5cm}](4,0) \\
}
\\ \hline
 \multicolumn{2}{|c|}{ {\blue \dft{} : dotsep = 3pt  dashsep= 5pt 3pt} }
\\ \hline  
\end{tabular} 

\bigskip


\textbf{\TFRGB{Nouvelle option}{New option}} : dashcolor \cite{pst-news10} 

 
\smallskip

\begin{tabular}{|c|c|} \hline  
\begin{psgraph}[axesstyle=none,xticksize=0 1,yticksize=0 4,subticks=0,shift=*](0,-.5)(4,.5){4cm}{1cm} 
\psline[linewidth=.5cm,linestyle=dashed,dashcolor=red](0,0)(4,0)
 \end{psgraph}
 &
\parbox[c]{10cm}{  
\BS{psline}[linewidth=.5cm,linestyle=dashed,\RDD{dashcolor}=red](0,0)(4,0) \\
}
\\ \hline 
\begin{psgraph}[axesstyle=none,xticksize=0 1,yticksize=0 4,subticks=0,shift=*](0,-.5)(4,.5){4cm}{1cm} 
\psline[linewidth=0.5cm,linestyle=dashed,linecolor=black,dashcolor=black!40,dash=5mm 5mm](0,0)(4,0)
 \end{psgraph}
 &
\parbox[c]{10cm}{  
\BS{psline}[linewidth=0.5cm,linestyle=dashed, \\ linecolor=black,\RDD{dashcolor}=black!40,dash=5mm 5mm](0,0)(4,0) \\
}
\\ \hline 
\end{tabular}  
 
 
 
%%---------------------------------------------------------------------------------------
\SbSSCT{ Lignes doubles}{Double line}

\begin{tabular}{|c|c|} \hline  
\begin{psgraph}[axesstyle=none,xticksize=0 1,yticksize=0 4,subticks=0,shift=*](0,-.5)(4,.5){4cm}{1cm}  
 \psline[doubleline=true](4,0)
 \end{psgraph}
 &
\parbox[c]{10cm}{  
\BS{psline}[\RDD{doubleline}=true](4,0)\\
}
\\ \hline 
\end{tabular} 

\smallskip

\begin{tabular}{|c|c|} \hline  
\begin{psgraph}[axesstyle=none,xticksize=0 1,yticksize=0 4,subticks=0,shift=*](0,-.5)(4,.5){4cm}{1cm}  
 \psline[linewidth=0.25cm,doubleline=true,doublesep=.3cm](4,0)
 \end{psgraph}
 &
\parbox[c]{10cm}{  
\BS{psline}[linewidth=0.25cm,doubleline=true,\RDD{doublesep}=.3cm](4,0)\\
}
\\ \hline 
\multicolumn{2}{|c|}{ {\blue \dft{} :  doublesep = 1.25\BS{pslinewidth}} } \\ \hline 
\end{tabular} 

\smallskip

\begin{tabular}{|c|c|} \hline 

\begin{psgraph}[axesstyle=none,xticksize=0 1,yticksize=0 4,subticks=0,shift=*](0,-.5)(4,.5){4cm}{1cm}  
 \psline[linewidth=0.25cm,doubleline=true,,doublecolor=red](4,0)
 \end{psgraph}
 &
\parbox[c]{10cm}{  
\BS{psline}[linewidth=0.25cm,doubleline=true,\RDD{doublecolor}=red](4,0)\\
}
\\ \hline
\multicolumn{2}{|c|}{ {\blue \dft{} : doublecolor = white } } \\ \hline 
\end{tabular} 



%%============================================

\SbSSCT{Bordure de ligne}{Border}


\begin{tabular}{|c|c|} \hline  
\begin{psgraph}[axesstyle=none,xticksize=0 1,yticksize=0 4,subticks=0,shift=*](0,-.5)(4,.5){4cm}{1cm}  
 \psline[linewidth=0.5cm,border=0.25cm,bordercolor=red](3,0)
 \end{psgraph}
 &
\parbox[c]{10cm}{  
\BS{psline}[linewidth=0.5cm,border=0.25cm,bordercolor=red](3,0) \\
}
\\ \hline 
\end{tabular} 

\smallskip

\begin{tabular}{|c|c|} \hline  
\begin{psgraph}[axesstyle=none,xticksize=- 1,yticksize=-1 ,subticks=0,shift=*](0,0)(-1,-1)(1,1){2cm}{2cm}  
 \psline[linewidth=3pt,linecolor=red](-1,-1)(1,1)
 \psline[linewidth=3pt,linecolor=blue,border=0.25cm](1,-1)(-1,1)
 \end{psgraph}
 &
\parbox[c]{10cm}{  
\BS{psline}[linewidth=3pt,linecolor=red](-1,-1)(1,1)\\
\BS{psline}[linewidth=3pt,linecolor=blue,border==0.25cm](1,-1)(-1,1)\\
}
\\ \hline 
\end{tabular} 


%%-------------------------------------------------------------------------------
\SbSSCT{Ombrage de ligne}{Shadow}

\begin{tabular}{|c|c|} \hline  
\begin{psgraph}[axesstyle=none,xticksize=0 1,yticksize=0 4,subticks=0,shift=*](0,-.5)(4,.5){4cm}{1cm}  
 \psline[linewidth=.5cm,shadow=true](3,0)
 \end{psgraph}
 &
\parbox[c]{10cm}{  
\BS{psline}[linecolor=red,\RDD{shadow}=true](3,0)\\
}
\\ \hline 
\multicolumn{2}{|c|}{ {\blue \dft{} : shadow = false } } \\ \hline 
\end{tabular} 

\smallskip

\begin{tabular}{|c|c|} \hline  
\begin{psgraph}[axesstyle=none,xticksize=0 1,yticksize=0 4,subticks=0,shift=*](0,-.5)(4,.5){4cm}{1cm}  
 \psline[linewidth=.5cm,shadow=true,shadowsize=.5cm](3,0)
 \end{psgraph}
 &
\parbox[c]{10cm}{  
\BS{psline}[linewidth=.5cm,shadow=true,\RDD{shadowsize}=.5cm](3,0)\\
}
\\ \hline
\multicolumn{2}{|c|}{ {\blue \dft{} : shadowsize = 3pt } } \\ \hline  
\end{tabular} 

\smallskip

\begin{tabular}{|c|c|} \hline  
\begin{psgraph}[axesstyle=none,xticksize=0 1,yticksize=0 4,subticks=0,shift=*](0,-.5)(4,.5){4cm}{1cm}  
 \psline[linewidth=.5cm,shadow=true,shadowcolor=green](3,0)
 \end{psgraph}
 &
\parbox[c]{10cm}{  
\BS{psline}[linewidth=.5cm,shadow=true,\RDD{shadowcolor}=green](3,0)\\
}
\\ \hline
\multicolumn{2}{|c|}{ {\blue \dft{} : shadowcolor = darkgray } } \\ \hline  
\end{tabular} 

\smallskip

\begin{tabular}{|c|c|} \hline  
\begin{psgraph}[axesstyle=none,xticksize=0 1,yticksize=0 4,subticks=0,shift=*](0,-.5)(4,.5){4cm}{1cm}  
 \psline[linewidth=.5cm,shadow=true,shadowsize=1cm,shadowangle=15](3,0)
 \end{psgraph}
 &
\parbox[c]{10cm}{  
\BS{psline}[linewidth=.5cm,shadow=true,shadowsize=1cm,\\
\RDD{shadowangle}=15](3,0)\\
}
\\ \hline
\multicolumn{2}{|c|}{ {\blue \dft{} : ,shadowangle = - 45 } } \\ \hline 
\end{tabular} 



%%-----------------------------------------------------------
\SbSSCT{Paramètres spécifiques}{Specific parameters }

\begin{tabular}{|c|c|} \hline 
\multicolumn{2}{|c|}{ \textbf{\BS{}psdiamond} } \\ \hline  
 \begin{psgraph}[axesstyle=none,xticksize=0 2cm,yticksize=0 4cm,subticks=0,shift=*](0,0)(4,2){4cm}{2cm} 
  \psdiamond[linestyle=dotted](2,1)(1,1)
 \psdiamond[gangle=30](2,1)(1,1)
  \end{psgraph}
 &
\parbox[c]{10cm}{  
\BS{psdiamond}[linestyle=dotted](2,1)(1,1)\\
\BS{psdiamond}[\RDD{gangle}=30](2,1)(1,1)\\
}
\\ \hline
\end{tabular}



\bigskip
\begin{tabular}{|c|c|} \hline 
\multicolumn{2}{|c|}{ \textbf{\BS{}pssarc} } \\ \hline  
\begin{psgraph}[axesstyle=none,xticksize=0 2cm,yticksize=0 2cm,subticks=0,shift=*](0,0)(2,2){2cm}{2cm}
\psarc[linecolor=green,linewidth=0.5cm,linejoin=0](0,0){2}{0}{90}  
\psarc[arcsep=1cm](0,0){2}{0}{90} 
\end{psgraph}
 &
\parbox[c]{10cm}{  
\BS{psarc}[linecolor=green,linewidth=0.5cm](0,0)\AC{2}\AC{0}\AC{90}\\
\BS{psarc}[\RDD{arcsep}=1cm](0,0)\AC{2}\AC{0}\AC{90}\\
}
\\ \hline
\begin{psgraph}[axesstyle=none,xticksize=0 2cm,yticksize=0 2cm,subticks=0,shift=*](0,0)(2,2){2cm}{2cm}
\psarc[linecolor=green,linewidth=0.5cm,linejoin=0](0,0){2}{0}{90}  
\psarc[arcsepA=1cm](0,0){2}{0}{90} 
\end{psgraph}
 &
\parbox[c]{10cm}{  
\BS{psarc}[linecolor=green,linewidth=0.5cm](0,0)\AC{2}\AC{0}\AC{90}\\
\BS{psarc}[\RDD{arcsepA}=1cm](0,0)\AC{2}\AC{0}\AC{90}\\
}
\\ \hline
\begin{psgraph}[axesstyle=none,xticksize=0 2cm,yticksize=0 2cm,subticks=0,shift=*](0,0)(2,2){2cm}{2cm}
\psarc[linecolor=green,linewidth=0.5cm,linejoin=0](0,0){2}{0}{90}  
\psarc[arcsepB=1cm](0,0){2}{0}{90} 
\end{psgraph}
 &
\parbox[c]{10cm}{  
\BS{psarc}[linecolor=green,linewidth=0.5cm](0,0)\AC{2}\AC{0}\AC{90}\\
\BS{psarc}[\RDD{arcsepB}=1cm](0,0)\AC{2}\AC{0}\AC{90}\\
 }
\\ \hline
%\multicolumn{2}{|c|}{ {\blue \dft{} : ,shadowangle = - 45 } } \\ \hline 
\end{tabular}

\smallskip

%%--------------------------------------------------------------
\SbSSCT{Remplissage de surface}{Filling }

\psset{fillcolor=yellow}

\begin{tabular}{|c|c|} \hline  
\begin{psgraph}[axesstyle=none,xticksize=0 1.5,yticksize=0 1.5,subticks=0,shift=*](0,0)(1.5,1.5){1.5cm}{1.5cm} 
 \psframe[fillstyle=solid](1,1)
 \end{psgraph}
 &
\parbox[c]{10cm}{  
\BS{psframe}[\RDD{fillstyle}=solid](1,1)\\
}
\\ \hline 
\end{tabular} 

\smallskip

\begin{tabular}{|c|c|c|c|c|c|c|} \hline
 \multicolumn{7}{|c|}{  Types de remplissages disponibles } \\ \hline
 \RDDD{none}{fillstyle}& \RDDD{solid}{fillstyle}& \RDDD{vlines}{fillstyle}  &  \RDDD{hlines}{fillstyle}  & \RDDD{crosshatch}{fillstyle}  & \RDDD{penrose}{fillstyle}   & dots\\  \hline
 & 
\begin{pspicture}[shift=*](-.2,-.2)(1.2,1.2)
 \psframe[fillstyle=solid](1,1)
\end{pspicture}
 & 
 \begin{pspicture}[shift=*](-.2,-.2)(1.2,1.2)
 \psframe[fillstyle=vlines ](1,1)
 \end{pspicture}
 & 
 \begin{pspicture}[shift=*](-.2,-.2)(1.2,1.2)
\psframe[fillstyle= hlines](1,1)
 \end{pspicture}

 & 
 \begin{pspicture}[shift=*](-.2,-.2)(1.2,1.2)
 \psframe[fillstyle=crosshatch](1,1)
  \end{pspicture}

&
 \begin{pspicture}[shift=*](-.2,-.2)(1.2,1.2)
\psframe[fillstyle= penrose,psscale=.5](1,1)
  \end{pspicture}

 &
  \begin{pspicture}[shift=*](-.2,-.2)(1.2,1.2)
 \psframe[fillstyle= dots,psscale=.5](1,1)
   \end{pspicture}
\\ \hline
& &  \RDD{vlines*} & \RDD{hlines*} & \RDD{crosshatch*} & \RDD{penrose*} &
\\ \hline
&
&
 \begin{pspicture}[shift=*](-.2,-.2)(1.2,1.2)
\psframe[fillstyle=vlines*](1,1)
 \end{pspicture}
 & 
  \begin{pspicture}[shift=*](-.2,-.2)(1.2,1.2)
 \psframe[fillstyle=hlines*](1,1)
  \end{pspicture}
 & 
 \begin{pspicture}[shift=*](-.2,-.2)(1.2,1.2)
\psframe[fillstyle= crosshatch*](1,1) 
  \end{pspicture}
 &
 \begin{pspicture}[shift=*](-.2,-.2)(1.2,1.2)
\psframe[fillstyle= penrose*,psscale=.5](1,1)
  \end{pspicture}
 &
  
 \\ \hline
\end{tabular}


\bigskip
\textbf{\TFRGB{Options disponibles}{Options available } }:
 
\begin{tabular}{|c|c|} \hline  
\begin{psgraph}[axesstyle=none,xticksize=0 1,yticksize=0 1,subticks=0,shift=*](0,0)(1,1){1cm}{1cm} 
\psframe[fillstyle=hlines*,fillcolor=green](1,1)
 \end{psgraph}
 &
\parbox[c]{10cm}{  
\BS{psframe}[fillstyle=hlines*,\RDD{fillcolor}=green](1,1)\\
}
\\ \hline 
\multicolumn{2}{|c|}{ {\blue  fillcolor = white  } } \\ \hline 
\begin{psgraph}[axesstyle=none,xticksize=0 1,yticksize=0 1,subticks=0,shift=*](0,0)(1,1){1cm}{1cm} 
\psframe[fillstyle=hlines*,hatchwidth=3pt](1,1)
\end{psgraph}
 &
\parbox[c]{10cm}{  
\BS{psframe}[fillstyle=hlines*,\RDD{hatchwidth}=3pt](1,1)\\
}
\\ \hline 
\multicolumn{2}{|c|}{ {\blue  hatchwidth = 0.8pt } } \\ \hline
\begin{psgraph}[axesstyle=none,xticksize=0 1,yticksize=0 1,subticks=0,shift=*](0,0)(1,1){1cm}{1cm} 
\psframe[fillstyle=hlines*,hatchsep=10pt](1,1)
\end{psgraph}
 &
\parbox[c]{10cm}{  
\BS{psframe}[fillstyle=hlines*,\RDD{hatchsep}=10pt](1,1)\\
}
\\ \hline
\multicolumn{2}{|c|}{ {\blue  hatchsep = 4pt } } \\ \hline
\begin{psgraph}[axesstyle=none,xticksize=0 1,yticksize=0 1,subticks=0,shift=*](0,0)(1,1){1cm}{1cm} 
\psframe[fillstyle=hlines*,hatchcolor=red](1,1)
\end{psgraph}
 &
\parbox[c]{10cm}{  
\BS{psframe}[fillstyle=hlines*,\RDD{hatchcolor}=red](1,1)\\
}
\\ \hline
\multicolumn{2}{|c|}{ {\blue hatchcolor = black } } \\ \hline
\begin{psgraph}[axesstyle=none,xticksize=0 1,yticksize=0 1,subticks=0,shift=*](0,0)(1,1){1cm}{1cm} 
\psframe[fillstyle=hlines*,hatchangle=25](1,1)
\end{psgraph}
 &
\parbox[c]{10cm}{  
\BS{psframe}[fillstyle=hlines*,\RDD{hatchangle}=25](1,1)\\
}
\\ \hline  
\multicolumn{2}{|c|}{ {\blue  hatchangle =  45 } } \\ \hline 

\end{tabular}  
 

\bigskip
\textbf{\TFRGB{Nouvelle option}{New option} : hatchwidthinc  hatchsepinc} \cite{pst-news05} 



\begin{tabular}{|c|} \hline 

\BS{}psframe[fillstyle=vlines,\RDD{hatchwidthinc}=2pt](14,1)
\\ \hline  
\begin{pspicture}(\linewidth,1.3)
\psframe[fillstyle=vlines,hatchwidthinc=2pt](14,1)
\end{pspicture}
\\ \hline 
\end{tabular} 

\smallskip
\begin{tabular}{|c|} \hline  
\BS{}psframe[fillstyle=vlines,\RDD{hatchsepinc}=2pt](14,1)
\\ \hline  

 \begin{pspicture}(\linewidth,1.3)
\psframe[fillstyle=vlines,hatchsepinc=2pt](14,1)
\end{pspicture}
\\ \hline 
\end{tabular} 

%%----------------------------------------------------------------------------------------
\SbSSCT{Coins arrondis}{Line arc}

\begin{tabular}{|c|c|} \hline  
\begin{psgraph}[axesstyle=none,xticksize=-1 1,yticksize=0 2,subticks=0,shift=*](0,0)(0,-1)(2,1){2cm}{2cm} 
\psline[linearc=0.5cm](0,1)(1,-1)(2,1)
\psline[linestyle=dotted,linewidth=1pt](0,1)(1,-1)(2,1)
 \end{psgraph}
 &
\parbox[c]{10cm}{  
\BS{psline}[\RDD{linearc}=0.5cm](0,1)(1,-1)(2,1)\\
\BS{psline}[linestyle=dotted](0,1)(1,-1)(2,1)\\
% \BS{end}\AC{psgraph} \\
 }\\ \hline
  \multicolumn{2}{|c|}{  {\blue  \dft{}: linearc  = 0 } } \\ \hline 
\end{tabular}


%%----------------------------------------------------

\SbSSCT{Types de points}{Dot style}


\begin{center}
\begin{tabular}{|c | p{3cm} || c | p{3cm} |}\hline
 \multicolumn{3}{|c|}{ \BS{}psdots [\RDD{dotstyle}=pentagon*](.5,0)(1.5,0)(2.5,0)  } 
& \psdots[dotstyle=pentagon*](.5,0)(1.5,0)(2.5,0)  \\ \hline 

* 	& \psdots[dotstyle=*](.5,0)(1.5,0)(2.5,0)  	&  o 	& \psdots[dotstyle=o](.5,0)(1.5,0)(2.5,0)  \\
x 	& \psdots[dotstyle=x](.5,0)(1.5,0)(2.5,0)	& + 	& \psdots[dotstyle=+](.5,0)(1.5,0)(2.5,0) \\
Bo 	&\psdots[dotstyle=Bo](.5,0)(1.5,0)(2.5,0)	&  B+	& \psdots[dotstyle=B+](.5,0)(1.5,0)(2.5,0) \\
asterisk & \psdots[dotstyle=asterisk](.5,0)(1.5,0)(2.5,0) &  Basterisk 	& \psdots[dotstyle=Basterisk](.5,0)(1.5,0)(2.5,0) 	\\
Asterisk & \psdots[dotstyle=Asterisk](.5,0)(1.5,0)(2.5,0) & BoldAsterisk & \psdots[dotstyle=BoldAsterisk](.5,0)(1.5,0)(2.5,0) \\
SolidAsterisk 	& \psdots[dotstyle=SolidAsterisk](.5,0)(1.5,0)(2.5,0) 	& oplus & \psdots[dotstyle=oplus](.5,0)(1.5,0)(2.5,0) \\
BoldOplus & \psdots[dotstyle=BoldOplus](.5,0)(1.5,0)(2.5,0) &  SolidOplus & \psdots[dotstyle=SolidOplus](.5,0)(1.5,0)(2.5,0) \\
otimes 	& \psdots[dotstyle=otimes](.5,0)(1.5,0)(2.5,0) 	& | 	& \psdots[dotstyle=|](.5,0)(1.5,0)(2.5,0) \\
square 	& \psdots[dotstyle=square](.5,0)(1.5,0)(2.5,0)	&  Bsquare 	& \psdots[dotstyle=Bsquare](.5,0)(1.5,0)(2.5,0)	  \\ 
square* 	& \psdots[dotstyle=square*](.5,0)(1.5,0)(2.5,0)	&  diamond 	& \psdots[dotstyle=diamond](.5,0)(1.5,0)(2.5,0)	 \\ 
diamond* 	& \psdots[dotstyle=diamond*](.5,0)(1.5,0)(2.5,0) &  triangle & \psdots[dotstyle=triangle ](.5,0)(1.5,0)(2.5,0)	 \\
Btriangle 	&   \psdots[dotstyle=Btriangle](.5,0)(1.5,0)(2.5,0)	&  triangle* &  \psdots[dotstyle=triangle*](.5,0)(1.5,0)(2.5,0)	 \\
pentagon 	& \psdots[dotstyle=pentagon](.5,0)(1.5,0)(2.5,0)	&  Bpentagon &  \psdots[dotstyle=Bpentagon](.5,0)(1.5,0)(2.5,0)	\\
pentagon* 	& \psdots[dotstyle=pentagon*](.5,0)(1.5,0)(2.5,0)	&  Hexagon 	& \psdots[dotstyle=Hexagon](.5,0)(1.5,0)(2.5,0)	\\
BoldHexagon & \psdots[dotstyle= BoldHexagon](.5,0)(1.5,0)(2.5,0) & SolidHexagon & \psdots[dotstyle=SolidHexagon](.5,0)(1.5,0)(2.5,0)\\

 Octogon & \psdots[dotstyle= Octogon](.5,0)(1.5,0)(2.5,0)&  BoldOctogon 	& \psdots[dotstyle= BoldOctogon](.5,0)(1.5,0)(2.5,0)\\
SolidOctogon 	& \psdots[dotstyle= SolidOctogon](.5,0)(1.5,0)(2.5,0)	& & \\

\hline 
 \multicolumn{4}{|c|}{ {\blue \dft{} : dotstyle = * } } \\ \hline 
\end{tabular} 
\end{center}
\footnotetext[1]{linecolor=blue,fillcolor=yellow}

%%----------------------------------------------------
\SbSSCT{Paramètres des points}{Parameters of the points }

\begin{tabular}{|c|c|} \hline  
\begin{psgraph}[axesstyle=none,xticksize=0 2,yticksize=0 3,subticks=0,shift=*](0,-1)(3,1){3cm}{2cm}
\psdots[dotstyle=diamond*, dotsize= 1cm](0.5,0)(1.5,0)(2.5,0)
 \end{psgraph}
 &
\parbox[c]{10cm}{  
\BS{psdots}[dotstyle=diamond*, \RDD{dotsize}= 1cm] (0.5,0)(1.5,0)(2.5,0)\\
}
\\ \hline
\begin{psgraph}[axesstyle=none,xticksize=0 2,yticksize=0 3,subticks=0,shift=*](0,-1)(3,1){3cm}{2cm}
\psdots[dotstyle=diamond*, dotsize= .5cm 10] (0.5,0)(1.5,0)(2.5,0)
\psdots[linecolor=red,dotstyle=diamond*, dotsize= 0.5cm ] (0.5,0)(1.5,0)(2.5,0)
 \end{psgraph}
 &
\parbox[c]{10cm}{  
\BS{psdots}[dotstyle=diamond*, \RDD{dotsize}= 0.5cm 10] \\ (0.5,0)(1.5,0)(2.5,0)\\
\BS{psdots}[linecolor=red,dotstyle=diamond*, \RDD{dotsize}= 0.5cm] \\ 
(0.5,0)(1.5,0)(2.5,0)\\
}
\\ \hline
\multicolumn{2}{|c|}{ {\blue  dotsize= 0.5pt 2.5 } } \\ \hline  
\begin{psgraph}[axesstyle=none,xticksize=0 1,yticksize=0 3,subticks=0,shift=*](0,-.5)(3,.5){3cm}{1cm}
\psdots[dotstyle=diamond*, dotscale= 5](0.5,0)(1.5,0)(2.5,0)

 \end{psgraph}
 &
\parbox[c]{10cm}{  
\BS{psdots}[dotstyle=diamond*, \RDD{dotscale}= 5] \\ (0.5,0)(1.5,0)(2.5,0)\\
}
\\ \hline
\begin{psgraph}[axesstyle=none,xticksize=0 1,yticksize=0 3,subticks=0,shift=*](0,-.5)(3,.5){3cm}{1cm}
\psdots[dotstyle=diamond*, dotscale= 5 2](0.5,0)(1.5,0)(2.5,0)
 \end{psgraph}
 &
\parbox[c]{10cm}{  
\BS{psdots}[dotstyle=diamond*, \RDD{dotscale}= 5 2] (0.5,0)(1.5,0)(2.5,0)\\
 }
\\ \hline
\multicolumn{2}{|c|}{ {\blue  dotscale= 1 } } \\ \hline 
\begin{psgraph}[axesstyle=none,xticksize=0 1,yticksize=0 3,subticks=0,shift=*](0,-.5)(3,.5){3cm}{1cm}
\psdots[dotstyle=diamond*,dotscale= 5, dotangle= 30](0.5,0)(1.5,0)(2.5,0)
 \end{psgraph}
 &
\parbox[c]{10cm}{  
\BS{psdots}[dotstyle=diamond*,dotscale= 5, \RDD{dotangle}= 30] \\
(0.5,0)(1.5,0)(2.5,0)\\
}
\\ \hline
\multicolumn{2}{|c|}{ {\blue  dotangle= 0 } } \\ \hline 
\end{tabular} 


\newpage


\SSCT{Les extrémités}{Arrowheads and such}
 
\SbSSCT{Les types d'extrémités disponibles}{Types of extremities available}

\psset{arrowscale=2,unit=1cm,linewidth=2pt}

\begin{center}
\begin{tabular}{|c|c|c|c|}
\hline  
\multicolumn{4}{|c|}{ \TFRGB{Extrémités à l'échelle 2}{Ends on scale 2}} \\ \hline

{\large \AC{-} }		& 
\begin{pspicture}[shift=*](0,-.5)(4.5,.5)
	\psline{-}(0.5,0)(4,0) 
\end{pspicture} 	&  &
\\
{\large \AC{<->}}		&
\begin{pspicture}[shift=*](0,-.5)(4.5,.5)
\psline{<->}(0.5,0)(4,0) 
\end{pspicture}  
&
{\large  \AC{>-<}}		&
\begin{pspicture}[shift=*](-1,-.5)(5.5,.5)
\psline{>-<}(0.5,0)(4,0) 
\end{pspicture} 
\\
{\large \AC{<<->>}}	&
\begin{pspicture}[shift=*](0,-.5)(4.5,.5)
\psline{<<->>}(0.5,0)(4,0) 
\end{pspicture} 
 	&
{\large \AC{>>-<<}}	&
\begin{pspicture}[shift=*](0,-.5)(4.5,.5)
\psline{>>-<<}(0.5,0)(4,0) 
\end{pspicture} 
	  \\
{\large \AC{|-| }}		&
\begin{pspicture}[shift=*](0,-.5)(4.5,.5)
\psline{|-|}(0.5,0)(4,0) 
\end{pspicture} 
	&
{\large \AC{|*-|*}}	&
\begin{pspicture}[shift=*](0,-.5)(4.5,.5)
\psline{|*-|*}(0.5,0)(4,0) 
\end{pspicture} 
\\

{\large \AC{[-]}}		&
\begin{pspicture}[shift=*](0,-.5)(4.5,.5)
\psline{[-]}(0.5,0)(4,0) 
\end{pspicture} 
	&
{\large \AC{]-[}}		& 
\begin{pspicture}[shift=*](0,-.5)(4.5,.5)
\psline{]-[}(0.5,0)(4,0) 
\end{pspicture} 
	 \\	

{\large \AC{(-)}}		&
\begin{pspicture}[shift=*](0,-.5)(4.5,.5)
\psline{(-)}(0.5,0)(4,0) 
\end{pspicture} 
	&
{\large \AC{)-(}}		&
\begin{pspicture}[shift=*](0,-.5)(4.5,.5)
\psline{)-(}(0.5,0)(4,0) 
\end{pspicture} 
 \\	
{\large \AC{o-o}}		&
\begin{pspicture}[shift=*](0,-.5)(4.5,.5)
\psline{o-o}(0.5,0)(4,0) 
\end{pspicture} 
	&
{\large \AC{*-*}}		&
\begin{pspicture}[shift=*](0,-.5)(4.5,.5)
\psline{*-*}(0.5,0)(4,0) 
\end{pspicture} 
	\\
{\large \AC{oo-oo}}		&
\begin{pspicture}[shift=*](0,-.5)(4.5,.5)
\psline{oo-oo}(0.5,0)(4,0) 
\end{pspicture} 
	 &
{\large \AC{**-**}}		&
\begin{pspicture}[shift=*](0,-.5)(4.5,.5)
\psline{**-**}(0.5,0)(4,0) 
\end{pspicture}
\\

{\large \AC{|<->|}}		&
\begin{pspicture}[shift=*](0,-.5)(4.5,.5)
\psline{|<->|}(0.5,0)(4,0) 
\end{pspicture}
 &
{\large \AC{|>-<|}}		& 
\begin{pspicture}[shift=*](0,-.5)(4.5,.5)
\psline{|>-<|}(0.5,0)(4,0) 
\end{pspicture}
	\\

{\large \AC{|<->|}}		&
\begin{pspicture}[shift=*](0,-.5)(4.5,.5)
\psline{|<->|}(0.5,0)(4,0) 
\end{pspicture}
 &
{\large \AC{|>-<|}}		& 
\begin{pspicture}[shift=*](0,-.5)(4.5,.5)
\psline{|>-<|}(0.5,0)(4,0) 
\end{pspicture}
	\\

{\large \AC{h-h}}		&
\begin{pspicture}[shift=*](0,-.5)(4.5,.5)
\psline{h-h}(0.5,0)(4,0) 
\end{pspicture}
 &
{\large \AC{H-H}}		&
\begin{pspicture}[shift=*](0,-.5)(4.5,.5)
\psline{H-H}(0.5,0)(4,0) 
\end{pspicture} 
	\\

{\large \AC{v-v}}		&
\begin{pspicture}[shift=*](0,-.5)(4.5,.5)
\psline{v-v}(0.5,0)(4,0) 
\end{pspicture} 
	 &
{\large \AC{V-V}}		&
\begin{pspicture}[shift=*](0,-.5)(4.5,.5)
\psline{V-V}(0.5,0)(4,0) 
\end{pspicture} 
	\\

{\large \AC{f-f}}		&
\begin{pspicture}[shift=*](0,-.5)(4.5,.5)
\psline{f-f}(0.5,0)(4,0) 
\end{pspicture} 
	&
{\large \AC{F-F}}		&
\begin{pspicture}[shift=*](0,-.5)(4.5,.5)
\psline{F-F}(0.5,0)(4,0) 
\end{pspicture} 
	\\

{\large \AC{t-t}}		& 
\begin{pspicture}[shift=*](0,-.5)(4.5,.5)
\psline{t-t}(0.5,0)(4,0) 
\end{pspicture}
 &
{\large \AC{T-T}}		& 
\begin{pspicture}[shift=*](0,-.5)(4.5,.5)
\psline{T-T}(0.5,0)(4,0) 
\end{pspicture}
	\\

{\large \AC{<D-D>}}		&
\begin{pspicture}[shift=*](0,-.5)(4.5,.5)
\psline[arrowscale=5]{<D-D>}(0.5,0)(4,0) 
\end{pspicture}
&	
{\large \AC{D>-<D}}		& 
\begin{pspicture}[shift=*](0,-.5)(4.5,.5)
\psline[arrowscale=5]{D>-<D}(0.5,0)(4,0) 
\end{pspicture}	\\

\hline
\multicolumn{4}{|c|}{\TFRGB{Largeur de ligne}{linewidth} : 0,3cm} \\ \hline
{\large \AC{-}}		& 
\begin{pspicture}[shift=*](0,-.5)(4.5,.5)
\psline[arrowscale=1,linewidth=.3cm]{-}(0.5,0)(4,0)
\psline[linecolor=red,linewidth=1pt](0.5,0)(4,0)
\end{pspicture}
			&
{\large  \AC{c-c}}		&
\begin{pspicture}[shift=*](0,-.5)(4.5,.5)
\psline[arrowscale=1,linewidth=.3cm]{c-c}(0.5,0)(4,0)
\psline[linecolor=red,linewidth=1pt](0.5,0)(4,0)
\end{pspicture}
	\\
{\large  \AC{C-C}}		& 
\begin{pspicture}[shift=*](0,-.5)(4.5,.5)
\psline[arrowscale=1,linewidth=.3cm]{C-C}(0.5,0)(4,0)
\psline[linecolor=red,linewidth=1pt](0.5,0)(4,0)
\end{pspicture}
	&
{\large \AC{cc-cc}}		&
\begin{pspicture}[shift=*](0,-.5)(4.5,.5)
\psline[arrowscale=1,linewidth=.3cm]{cc-cc}(0.5,0)(4,0)
\psline[linecolor=red,linewidth=1pt](0.5,0)(4,0)
\end{pspicture}
	\\
\hline
\end{tabular}
\end{center}

%---------------------------------------------------------------------------
\subsection[Linejoin linecap]{Linejoin linecap \cite{pst-news08} }
\begin{tabular}{|c|c|c|}
\hline 
\multicolumn{3}{|c|}{\BS{}psline[\RDD{linecap}=0,linewidth=10pt](2,0.5)(2,2.5) } \\
\hline 
 \begin{pspicture}(4,3)
 \psline[linecap=0,linewidth=10pt](2,.5)(2,2.5)
  \psline[linecolor=red,linewidth=1pt](2,.5)(2,2.5) 
  \end{pspicture} 
 &  \begin{pspicture}(4,3)
 \psline[linecap=1,linewidth=10pt] (2,.5)(2,2.5)
 \psline[linecolor=red,linewidth=1pt](2,.5)(2,2.5)  
  \end{pspicture} & 
   \begin{pspicture}(4,3)
 \psline[linecap=2,linewidth=10pt](2,.5)(2,2.5)
 \psline[linecolor=red,linewidth=1pt](2,.5)(2,2.5) 
    \end{pspicture} \\ 
\hline linecap=0 & linecap=1 & linecap=2 \\ 
\hline \hline
\multicolumn{3}{|c|}{\BS{}pstriangle[\RDD{linejoin}=0,linewidth=10pt] (2,0.5)(2,2) } \\
\hline 
 \begin{pspicture}(4,3)
 \pstriangle[linejoin=0,linewidth=10pt] (2,.5)(2,2)
 \pstriangle[linecolor=red,linewidth=1pt] (2,.5)(2,2)
  \end{pspicture} 
 &  \begin{pspicture}(4,3)
 \pstriangle[linejoin=1,linewidth=10pt] (2,.5)(2,2)
  \pstriangle[linecolor=red,linewidth=1pt] (2,.5)(2,2) 
  \end{pspicture} & 
   \begin{pspicture}(4,3)
 \pstriangle[linejoin=2,linewidth=10pt] (2,.5)(2,2)
  \pstriangle[linecolor=red,linewidth=1pt] (2,.5)(2,2)  \end{pspicture} \\ 
\hline linejoin=0 & linejoin=1 & linejoin=2 \\ 
\hline 
\end{tabular} 

%------------------------------------------------------------------------------------------
\SbSSCT{Flèches multiples}{Multiple arrows}


\begin{center}
\begin{tabular}{|p{6cm}|p{6cm}|}
\hline  
\multicolumn{2}{|c|}{ \textbf{nArrows}} \\ \hline

\rule[-.5cm]{0pt}{1cm}	\psline[nArrows=5]{>>->>}(0.5,0)(5,0)  		&
\rule[-.5cm]{0pt}{1cm} 	\psline[nArrows=5]{<<-<<}(0.5,0)(5,0) 	\\ \hline
\BS{psline}[\RDD{nArrows}=5]\AC{>>->>}(0.5,0)(5,0) \RDI{nArrows}{pstricks-add} & 
\BS{psline}[\RDD{nArrows}=5]\AC{<<-<<}(0.5,0)(5,0)  \\ \hline
	\end{tabular}
\end{center}


\begin{center}
\begin{tabular}{|p{6cm}|p{6cm}|}
\hline  
\multicolumn{2}{|c|}{ \textbf{nArrowsA ArrowsB}} \\ \hline

\rule[-.5cm]{0pt}{1cm}	\psline[nArrowsA=5]{>>->>}(0.5,0)(5,0)  		&
\rule[-.5cm]{0pt}{1cm} 	\psline[nArrowsB=5]{>>->>}(0.5,0)(5,0) 	\\ \hline
\BS{}psline[\RDD{nArrowsA}=5]\AC{>>->>}(0.5,0)(5,0) \RDI{nArrowsA}{pstricks-add} & 
\BS{}psline[\RDD{nArrowsB}=5]\AC{>>->>}(0.5,0)(5,0) \RDI{nArrowsB}{pstricks-add} \\ \hline
\rule[-.5cm]{0pt}{1cm}	\psline[nArrowsA=5]{<<-<<}(0.5,0)(5,0)  		&
\rule[-.5cm]{0pt}{1cm} 	\psline[nArrowsB=5]{<<-<<}(0.5,0)(5,0) 	\\ \hline
\BS{}psline[{\red nArrowsA=5}]\AC{<<-<<}(0.5,0)(5,0) & 
\BS{}psline[{\red nArrowsB=5}]\AC{<<-<<}(0.5,0)(5,0) \\ \hline

\end{tabular}
\end{center}

\SbSSCT{Paramètres des extrémités}{Parameters available}
\label{parametres}

\begin{center}
\begin{tabular}{|p{7cm}|p{7cm}|} \hline 
\rule[-.5cm]{0pt}{1cm}	\psline[arrowsize=3pt 3]{->}(0.5,0)(5,0)  		
&
\rule[-.5cm]{0pt}{1cm} 	\psline[arrowlength=5]{->}(0.5,0)(5,0) 	\\ \hline
\BS{psline}[\RDD{Arrowsize}=3pt 3]\AC{->} & \BS{psline}[\RDD{arrowlength}= 5]\AC{->}\\ \hline
{\blue \dft{}  : arrowsize= 1.5pt 2} & {\blue \dft{}  : arrowlength= 1.4 }\\ \hline
\end{tabular}
\end{center}
 
\begin{center}
\begin{tabular}{|p{7cm}|p{7cm}|} \hline 
\rule[-.5cm]{0pt}{1cm}	\psline[arrowinset=0,arrowscale=4]{->}(0.5,0)(5,0)  		&
\rule[-.5cm]{0pt}{1cm} 	\psline[arrowinset=.8,arrowscale=4]{->}(0.5,0)(5,0) 	\\ \hline
\BS{}psline[\RDD{arrowinset}=0]\AC{->} &
 \BS{}psline[{\red arrowinset=.8}]\AC{->}\\ \hline
%par défaut :  1.5pt 2 & par défaut : 1.4 \\ \hline
\multicolumn{2}{|c|}{\blue \dft{}  : arrowinset=.4  (40\% )} \\ \hline
\end{tabular}
\end{center}
 
\begin{center}
\begin{tabular}{|p{7cm}|p{7cm}|} \hline 
\rule[-.5cm]{0pt}{1cm}	\psline[tbarsize=4pt 2]{|<-|}(0.5,0)(5,0)  
&
\rule[-.5cm]{0pt}{1cm} 	\psline[tbarsize=4pt 2,linewidth=4pt]{|<-|}(0.5,0)(5,0) 	
\\ \hline
linewidth=2pt & linewidth=4pt \\ \hline
\multicolumn{2}{|c|}{\BS{}psline[\RDD{tbarsize}=4pt 2]\AC{|<-|} \hspace{1cm} {\blue \dft{}  : tbarsize=2pt 5 } } \\ \hline
\end{tabular}
\end{center}
 
\begin{center}
\begin{tabular}{|p{7cm}|p{7cm}|} \hline 
\rule[-.5cm]{0pt}{1cm}	\psline[bracketlength=.5]{-]}(0.5,0)(5,0)  		&
\rule[-.5cm]{0pt}{1cm} 	\psline[rbracketlength=.5]{-)}(0.5,0)(5,0) 	\\ \hline
\BS{}psline[\RDD{bracketlength}=.5]\AC{-]} & 
\BS{}psline[\RDD{rbracketlength}=.5]\AC{-)}\\ \hline
{\blue \dft{}  : bracketlength= 0.15} & {\blue \dft{}  : rbracketlength=0.15}\\ \hline
\end{tabular}
\end{center}
 
\begin{center}
\begin{tabular}{|p{7cm}|p{7cm}|} \hline 
\rule[-.5cm]{0pt}{1cm}	\psline[arrowscale=5]{->}(0.5,0)(5,0)  		&
\rule[-.5cm]{0pt}{1cm} 	\psline[arrowscale= 5 10]{->}(0.5,0)(5,0) 	\\ \hline
\BS{}psline[\RDD{arrowscale}=5]\AC{->} & 
\BS{}psline[{\red arrowscale= 5 10}]\AC{->}\\ \hline
\multicolumn{2}{|c|}{\blue  \dft{}  : arrowscale=1 } \\ \hline
\end{tabular}
\end{center}

\begin{center}
\begin{tabular}{|p{7cm}|p{7cm}|} \hline 
\rule[-.5cm]{0pt}{1cm}	\psline[hooklength=10mm]{-H}(0.5,0)(5,0)  		&
\rule[-.5cm]{0pt}{1cm} 	\psline[hookwidth=3mm]{-H}(0.5,-.3)(5,-.3) 	\\ \hline
\BS{}psline[\RDD{hooklength}=10mm ]\AC{-H} & 
\BS{}psline[\RDD{hookwidth}=3mm]\AC{-H}\\ \hline
{\blue \dft{}  : hooklength=3mm } & {\blue \dft{}  : hookwidth=1mm }	\\ \hline
\end{tabular}
\end{center}

 
\begin{center}
\begin{tabular}{|p{7cm}|p{7cm}|} \hline 
\rule[-.5cm]{0pt}{1cm}	\psline[arrowLW=1pt]{o-*}(0.5,0)(5,0)  		&
\rule[-.5cm]{0pt}{1cm} 	\psline[arrowLW=1mm]{*-o}(0.5,0)(5,0) 	\\ \hline
\BS{}psline[\RDD{arrowLW}=1pt]\AC{o-*} & 
\BS{}psline[{\red arrowLW=1mm}]\AC{*-o}\\ \hline
%\multicolumn{2}{|c|}{ par défaut : arrowscale=1 } \\ \hline
\end{tabular}
\end{center}

\begin{center}
\begin{tabular}{|p{7cm}|p{7cm}|} \hline 
\rule[-.5cm]{0pt}{1cm}	\psline[veearrowlength=.5cm]{v-V}(0.5,0)(2,0)  		&
\rule[-.5cm]{0pt}{1cm} 	\psline[veearrowangle=60]{v-V}(0.5,0)(2,0) 	\\ \hline
\BS{psline}[\RDD{veearrowlength}=.5cm ]\AC{v-V} \RDI{veearrowlength}{pstricks-add} & 
\BS{psline}[\RDD{veearrowangle}=60]\AC{v-V} \RDI{veearrowangle}{pstricks-add}\\ \hline
{\blue \dft{}  : veearrowlength = 3mm} & {\blue \dft{}  : veearrowangle = 30 }	\\ \hline
\end{tabular}
\end{center}

\begin{center}
\begin{tabular}{|p{7cm}|p{7cm}|} \hline 
\rule[-.5cm]{0pt}{1cm}	\psline[veearrowlinewidth =.5mm]{v-V}(0.5,0)(2,0)  		&
\rule[-.5cm]{0pt}{1cm} 	\psline[filledveearrowlength = 5mm]{f-F}(0.5,0)(2,0) 	\\ \hline
\BS{}psline[\RDD{veearrowlinewidth} =.5mm ]\AC{v-V} \RDI{veearrowlinewidth}{pstricks-add} & 
\BS{}psline[\RDD{filledveearrowlength} = 5mm]\AC{f-F} \RDI{filledveearrowlength}{pstricks-add}\\ \hline
{\blue  \dft{}  : veearrowlinewidth = 0.35mm} & {\blue \dft{}  : filledveearrowlength = 3mm }	\\ \hline
\end{tabular}
\end{center}

\begin{center}
\begin{tabular}{|p{7cm}|p{7cm}|} \hline 
\rule[-.5cm]{0pt}{1cm}	\psline[filledveearrowangle = 45]{f-F}(0.5,0)(2,0)  		&
\rule[-.5cm]{0pt}{1cm} 	\psline[filledveearrowlinewidth =1mm]{f-F}(0.5,0)(2,0) 	\\ \hline
\BS{psline}[\RDD{filledveearrowangle} = 45 ]\AC{f-F} \RDI{filledveearrowangle}{pstricks-add} & 
\BS{psline}[\RDD{filledveearrowlinewidth} =1mm]\AC{f-F} \RDI{filledveearrowlinewidth}{pstricks-add} \\ \hline
{\blue \dft{}  : filledveearrowangle = 15} & {\blue \dft{}  : filledveearrowlinewidth =0.35mm }	\\ \hline
\end{tabular}
\end{center}

\begin{center}
\begin{tabular}{|p{7cm}|p{7cm}|} \hline 
\rule[-.5cm]{0pt}{1cm}	\psline[tickarrowlength=2.5mm]{t-T}(0.5,0)(2,0)  		&
\rule[-.5cm]{0pt}{1cm} 	\psline[tickarrowlinewidth=1mm]{t-T}(0.5,0)(2,0) 	\\ \hline
\BS{psline}[\RDD{tickarrowlength}=2.5mm ]\AC{t-T} \RDI{tickarrowlength}{pstricks-add} & 
\BS{psline}[\RDD{tickarrowlinewidth}=1mm]\AC{t-T} \RDI{tickarrowlinewidth}{pstricks-add} \\ \hline
{\blue \dft{}  : tickarrowlength= 1.5mm } &{\blue  \dft{}  : tickarrowlinewidth=0.35mm }	\\ \hline
\end{tabular}
\end{center}

\begin{center}
\begin{tabular}{|p{7cm}|p{7cm}|} \hline 
\rule[-.5cm]{0pt}{1cm}	\psline[arrowlinestyle=dotted]{t-T}(0.5,0)(2,0)  		&
\rule[-.5cm]{0pt}{1cm} 	\psline[arrowlinestyle=dashed]{v-V}(0.5,0)(2,0) 	\\ \hline
\BS{psline}[\RDD{arrowlinestyle}=dotted ]\AC{t-T} & 
\BS{psline}[\RDD{arrowlinestyle}=dashed]\AC{v-V}\\ \hline
\multicolumn{2}{|c|}{arrowlinestyle= solid } \\ \hline
\end{tabular}
\end{center}


\begin{center}
\begin{tabular}{|p{7cm}|p{7cm}|} \hline 
\rule[-.5cm]{0pt}{1cm}	\psline[ArrowFill=false,arrowinset=0]{<->}(0.5,0)(2,0)  		&
\rule[-.5cm]{0pt}{1cm} 	\psline[ArrowFill=false]{>-<}(0.5,0)(2,0) 	\\ \hline
\BS{}psline[\RDD{ArrowFill}=false,arrowinset=0 ]\AC{>-<} \RDI{ArrowFill}{pstricks-add}  & 
\BS{}psline[{\red ArrowFill=false}]\AC{>-<}\\ \hline
%\multicolumn{2}{|c|}{arrowlinestyle= solid } \\ \hline
\end{tabular}
\end{center}


\begin{center}
\begin{tabular}{|p{7cm}|p{7cm}|} \hline 
\rule[-.5cm]{0pt}{1cm}	\psline[arrowsize=3pt 3]{->}(0.5,0)(5,0)  		&
\rule[-.5cm]{0pt}{1cm} 	\psline[arrowlength=5]{->}(0.5,0)(5,0) 	\\ \hline
\BS{}psline[\RDD{Arrowsize}=3]\AC{->} & \BS{}psline[\RDD{arrowlength}= 5]\AC{->}\\ \hline
{\blue \dft{}  : arrowsize= 1.5pt 2 } & {\blue \dft{}  : arrowlength= 1.4 } \\ \hline
\end{tabular}
\end{center}
 
\begin{center}
\begin{tabular}{|p{7cm}|p{7cm}|} \hline 
\rule[-.5cm]{0pt}{1cm}	\psline[arrowinset=0,arrowscale=4]{->}(0.5,0)(5,0)  		&
\rule[-.5cm]{0pt}{1cm} 	\psline[arrowinset=.8,arrowscale=4]{->}(0.5,0)(5,0) 	\\ \hline
\BS{}psline[\RDD{arrowinset}{\red =0}]\AC{->} &
 \BS{}psline[{\red arrowinset=.8}]\AC{->}\\ \hline
%par défaut :  1.5pt 2 & par défaut : 1.4 \\ \hline
\multicolumn{2}{|c|}{\blue \dft{}  : arrowscale=.4  (40\% )} \\ \hline
\end{tabular}
\end{center}
 
\begin{center}
\begin{tabular}{|p{7cm}|p{7cm}|} \hline 
\rule[-.5cm]{0pt}{1cm}	\psline[tbarsize=4pt 2]{|<-|}(0.5,0)(5,0)  		&
\rule[-.5cm]{0pt}{1cm} 	\psline[tbarsize=4pt 2,linewidth=4pt]{|<-|}(0.5,0)(5,0) 	\\ \hline
%\BS{}psline[{\red tbarsize=4pt 2}](0.5,0)(2,0) & \BS{}psline[{\red arrowinset=.8}](0.5,0)(2,0)\\ \hline
linewidth=2pt & linewidth=4pt \\ \hline
\multicolumn{2}{|c|}{\BS{}psline[\RDD{tbarsize}=4pt 2]\AC{|<-|} \hspace{1cm} {\blue \dft{}  : tbarsize=2pt 5 } } \\ \hline
\end{tabular}
\end{center}
 
\begin{center}
\begin{tabular}{|p{7cm}|p{7cm}|} \hline 
\rule[-.5cm]{0pt}{1cm}	\psline[bracketlength=.5]{-]}(0.5,0)(5,0)  		&
\rule[-.5cm]{0pt}{1cm} 	\psline[rbracketlength=.5]{-)}(0.5,0)(5,0) 	\\ \hline
\BS{psline}[\RDD{bracketlength}=.5]\AC{-]} & 
\BS{psline}[\RDD{rbracketlength}=.5]\AC{-)}\\ \hline
{\blue \dft{}  0.15} & {\blue \dft{}  0.15 } \\ \hline
%\multicolumn{2}{|c|}{\BS{}psline[{\red tbarsize=4pt 2}](0.5,0)(2,0) par défaut : tbarsize=2pt 5  } \\ \hline
\end{tabular}
\end{center}
 
\begin{center}
\begin{tabular}{|p{7cm}|p{7cm}|} \hline 
\rule[-.5cm]{0pt}{1cm}	\psline[arrowscale=5]{->}(0.5,0)(5,0)  		&
\rule[-.5cm]{0pt}{1cm} 	\psline[arrowscale= 5 10]{->}(0.5,0)(5,0) 	\\ \hline
\BS{psline}[\RDD{arrowscale}=5]\AC{-]} & 
\BS{psline}[{\red arrowscale= 5 10}]\AC{-)}\\ \hline
\multicolumn{2}{|c|}{\blue  \dft{}  : arrowscale=1 } \\ \hline
\end{tabular}
\end{center}

\begin{center}
\begin{tabular}{|p{7cm}|p{7cm}|} \hline 
\rule[-.5cm]{0pt}{1cm}	\psline[hooklength=10mm]{-H}(0.5,0)(5,0)  		&
\rule[-.5cm]{0pt}{1cm} 	\psline[hookwidth=3mm]{-H}(0.5,-.3)(5,-.3) 	\\ \hline
\BS{psline}[\RDD{hooklength}=10mm ]\AC{-H} \RDI{hooklength}{pstricks-add} & 
\BS{psline}[\RDD{hookwidth}=3mm]\AC{-H} \RDI{hookwidth}{pstricks-add}  \\ \hline
{\blue \dft{}  : hooklength=3mm } & {\blue \dft{}  : hookwidth=1mm }	\\ \hline
\end{tabular}
\end{center}

 
\begin{center}
\begin{tabular}{|p{7cm}|p{7cm}|} \hline 
\rule[-.5cm]{0pt}{1cm}	\psline[arrowLW=1pt]{o-*}(0.5,0)(5,0)  		&
\rule[-.5cm]{0pt}{1cm} 	\psline[arrowLW=1mm]{*-o}(0.5,0)(5,0) 	\\ \hline
\BS{psline}[\RDD{arrowLW}=1pt]\AC{o-*} & 
\BS{psline}[{\red arrowLW=1mm}]\AC{*-o}\\ \hline
%\multicolumn{2}{|c|}{ par défaut : arrowscale=1 } \\ \hline
\end{tabular}
\end{center}

\begin{center}
\begin{tabular}{|p{7cm}|p{7cm}|} \hline 
\rule[-.5cm]{0pt}{1cm}	\psline[veearrowlength=.5cm]{v-V}(0.5,0)(2,0)  		&
\rule[-.5cm]{0pt}{1cm} 	\psline[veearrowangle=60]{v-V}(0.5,0)(2,0) 	\\ \hline
\BS{}psline[\RDD{veearrowlength}=.5cm ]\AC{v-V} & 
\BS{}psline[\RDD{veearrowangle}=60]\AC{v-V}\\ \hline
{\blue \dft{}  : veearrowlength = 3mm} & {\blue \dft{}  : veearrowangle = 30}	\\ \hline
\end{tabular}
\end{center}

\begin{center}
\begin{tabular}{|p{7cm}|p{7cm}|} \hline 
\rule[-.5cm]{0pt}{1cm}	\psline[veearrowlinewidth =.5mm]{v-V}(0.5,0)(2,0)  		&
\rule[-.5cm]{0pt}{1cm} 	\psline[filledveearrowlength = 5mm]{f-F}(0.5,0)(2,0) 	\\ \hline
\BS{}psline[\RDD{veearrowlinewidth}=.5mm]\AC{v-V} & 
\BS{}psline[\RDD{filledveearrowlength} = 5mm]\AC{f-F}\\ \hline
{\blue \dft{}  : veearrowlinewidth = 0.35mm} & {\blue  \dft{}  : filledveearrowlength = 3mm }	\\ \hline
\end{tabular}
\end{center}

\begin{center}
\begin{tabular}{|p{7cm}|p{7cm}|} \hline 
\rule[-.5cm]{0pt}{1cm}	\psline[filledveearrowangle = 45]{f-F}(0.5,0)(2,0)  		&
\rule[-.5cm]{0pt}{1cm} 	\psline[filledveearrowlinewidth =1mm]{f-F}(0.5,0)(2,0) 	\\ \hline
\BS{}psline[\RDD{filledveearrowangle} = 45 ]\AC{f-F} & 
\BS{}psline[\RDD{filledveearrowlinewidth} =1mm]\AC{f-F}\\ \hline
{\blue \dft{}  : filledveearrowangle = 15 } & {\blue  \dft{}  : filledveearrowlinewidth =0.35mm }	\\ \hline
\end{tabular}
\end{center}

\begin{center}
\begin{tabular}{|p{7cm}|p{7cm}|} \hline 
\rule[-.5cm]{0pt}{1cm}	\psline[tickarrowlength=2.5mm]{t-T}(0.5,0)(2,0)  		&
\rule[-.5cm]{0pt}{1cm} 	\psline[tickarrowlinewidth=1mm]{t-T}(0.5,0)(2,0) 	\\ \hline
\BS{}psline[\RDD{tickarrowlength}=2.5mm ]\AC{t-T} & 
\BS{}psline[\RDD{tickarrowlinewidth}=1mm]\AC{t-T}\\ \hline
{\blue \dft{}  : tickarrowlength= 1.5mm } & {\blue \dft{}  : tickarrowlinewidth=0.35mm }	\\ \hline
\end{tabular}
\end{center}

\begin{center}
\begin{tabular}{|p{7cm}|p{7cm}|} \hline 
\rule[-.5cm]{0pt}{1cm}	\psline[arrowlinestyle=dotted]{t-T}(0.5,0)(2,0)  		&
\rule[-.5cm]{0pt}{1cm} 	\psline[arrowlinestyle=dashed]{v-V}(0.5,0)(2,0) 	\\ \hline
\BS{}psline[\RDD{arrowlinestyle}=dotted ]\AC{t-T} & 
\BS{}psline[\RDD{arrowlinestyle}=dashed]\AC{v-V}\\ \hline
\multicolumn{2}{|c|}{arrowlinestyle= solid } \\ \hline
\end{tabular}
\end{center}

\begin{center}
\begin{tabular}{|p{7cm}|p{7cm}|} \hline 
\rule[-.5cm]{0pt}{1cm}	\psline[ArrowFill=false,arrowinset=0]{<->}(0.5,0)(2,0)  		&
\rule[-.5cm]{0pt}{1cm} 	\psline[ArrowFill=false]{>-<}(0.5,0)(2,0) 	\\ \hline
\BS{psline}[\RDD{ArrowFill}=false,arrowinset=0 ]\AC{>-<} & 
\BS{psline}[{\red ArrowFill=false}]\AC{>-<}\\ \hline

\end{tabular}
\end{center}

\psset{arrowscale=1}
%\vspace{2cm}

\newpage

\label{poly}
\section[Des polygones avec pst poly]{Des polygones avec pst poly  \cite{pst-poly}}
\SSCT{Des polygones avec pst poly}{Plolygons with pstpoly}

\psset{PstPicture=false}

Utilisation du module pst-poly (consultez le fichier pst-poly-doc.pdf )

\begin{tabular}{|l|} \hline  
\begin{psgraph}[axesstyle=none,xticksize=-1 1,yticksize=-1 1,subticks=0](0,0)(-1,-1)(1,1){4cm}{4cm} 
  \PstPolygon
  \pscircle[linestyle=dotted]{1}
\end{psgraph}
\\ \hline  
  \BSS{PstPolygon}  \BSI{PstPolygon}{pst-poly} \\
  \BS{pscircle}[linestyle=dotted]\AC{1}
\\ \hline 
\end{tabular}

\bigskip


\begin{tabular}{|c|c|} \hline  
\rule[-1.5cm]{0pt}{1cm}texte avant \PstPolygon[PstPicture=false] texte après
&  
texte avant \PstPolygon[PstPicture=true] texte après
\\ \hline  
texte avant \BS{PstPolygon}[\RDD{PstPicture=false}] texte après 
&  
texte avant \BS{PstPolygon}[\RDD{PstPicture=true}] texte après
\\ \hline 
\multicolumn{2}{|c|}{\dft :  PstPicture=true}  \RDI{PstPicture}{pst-poly} 
\\ \hline 
\end{tabular} 

\bigskip

\begin{tabular}{|c|c|} \hline  
\begin{psgraph}[axesstyle=none,xticksize=-2 2,yticksize=-2 2,subticks=0](0,0)(-2,-2)(2,2){4cm}{4cm} 
   \PstPolygon[unit=2]
   \pscircle[linestyle=dotted]{2}
\end{psgraph}
&  
\begin{psgraph}[axesstyle=none,xticksize=-2 2,yticksize=-2 2,subticks=0](0,0)(-2,-2)(2,2){4cm}{4cm} 
   \PstPolygon[unit=.5]
   \pscircle[linestyle=dotted]{.5}
\end{psgraph}
\\ \hline  
\BS{PstPolygon}[\RDD{unit}=2]   \RDI{unit}{pst-poly} 
&
\BS{PstPolygon}[\RDD{unit}=.5]
\\
\BS{pscircle}[linestyle=dotted]\AC{2}
&
\BS{pscircle}[linestyle=dotted]\AC{.5}
\\ \hline 
\end{tabular} 

\subsection{Options}


\begin{tabular}{|c|c|c|} \hline  
\begin{psgraph}[axesstyle=none,xticksize=-1 1,yticksize=-1 1,subticks=0](0,0)(-1,-1)(1,1){4cm}{4cm} 
  \PstPolygon[linestyle=dashed,PstPicture=false]
\PstPolygon[PolyRotation=18,PstPicture=false]
\end{psgraph}
&  
\begin{psgraph}[axesstyle=none,xticksize=-1 1,yticksize=-1 1,subticks=0](0,0)(-1,-1)(1,1){4cm}{4cm}
  \PstPolygon[linestyle=dashed,PstPicture=false]
\PstPolygon[PolyRotation=36,PstPicture=false]
\end{psgraph}
&  
\begin{psgraph}[axesstyle=none,xticksize=-1 1,yticksize=-1 1,subticks=0](0,0)(-1,-1)(1,1){4cm}{4cm}
  \PstPolygon[linestyle=dashed,PstPicture=false]
\PstPolygon[PolyRotation=45,PstPicture=false]
\end{psgraph}
\\ \hline 
\RDD{PolyRotation}=18  \RDI{PolyRotation}{pst-poly} & 
\RDD{PolyRotation}=36 & 
\RDD{PolyRotation}=45 \\ 
\hline 
\end{tabular}


\bigskip


\begin{tabular}{|c|c|c|} \hline  
\begin{psgraph}[axesstyle=none,xticksize=-1 1,yticksize=-1 1,subticks=0](0,0)(-1,-1)(1,1){4cm}{4cm}
  \PstPolygon[linestyle=dashed]
\PstPolygon[PolyNbSides=3]
\end{psgraph}
&  
\begin{psgraph}[axesstyle=none,xticksize=-1 1,yticksize=-1 1,subticks=0](0,0)(-1,-1)(1,1){4cm}{4cm}
  \PstPolygon[linestyle=dashed]
\PstPolygon[PolyNbSides=7]
\end{psgraph}
&  
\begin{psgraph}[axesstyle=none,xticksize=-1 1,yticksize=-1 1,subticks=0](0,0)(-1,-1)(1,1){4cm}{4cm}
  \PstPolygon[linestyle=dashed]
\PstPolygon[PolyNbSides=50]
\end{psgraph}
\\ \hline 
\RDD{PolyNbSides}=3 & \RDD{PolyNbSides}=7 & \RDD{PolyNbSides}=50  \RDI{PolyNbSides}{pst-poly} \\ 
\hline 
\end{tabular}

\bigskip  

\begin{tabular}{|c|c|c|} \hline  
\begin{psgraph}[axesstyle=none,xticksize=-1 1,yticksize=-1 1,subticks=0](0,0)(-1,-1)(1,1){4cm}{4cm} 
  \PstPolygon[linestyle=dashed,]
\PstPolygon[PolyOffset=2]
\end{psgraph}
&  
\begin{psgraph}[axesstyle=none,xticksize=-1 1,yticksize=-1 1,subticks=0](0,0)(-1,-1)(1,1){4cm}{4cm}
  \PstPolygon[linestyle=dashed,PolyNbSides=7]
\PstPolygon[PolyOffset=2,PolyNbSides=7]
\end{psgraph}
&  
\begin{psgraph}[axesstyle=none,xticksize=-1 1,yticksize=-1 1,subticks=0](0,0)(-1,-1)(1,1){4cm}{4cm}
  \PstPolygon[linestyle=dashed,PolyNbSides=7]
\PstPolygon[PolyOffset=4,PolyNbSides=7]
\end{psgraph}
\\ \hline 
\RDD{PolyOffset}=2 & \RDD{PolyOffset}=2 & \RDD{PolyOffset}=4  \RDI{PolyOffset}{pst-poly} \\ 
  & PolyNbSides=7 & PolyNbSides=7 \\
\hline 
\end{tabular}

\bigskip

\begin{tabular}{|c|c|c|} \hline 
\multicolumn{3}{|c|}{\BS{PstPolygon}[\RDD{PolyIntermediatePoint}=0.5] \RDI{PolyIntermediatePoint}{pst-poly} }
\\ \hline 
\begin{psgraph}[axesstyle=none,xticksize=-2 2,yticksize=-2 2,subticks=0](0,0)(-2,-2)(2,2){4cm}{4cm} 
\PstPolygon[linestyle=dotted,]
   \PstPolygon[PolyIntermediatePoint=0.5]
   \pscircle[linestyle=dotted,linecolor=red]{0.5}
\end{psgraph}
&  
\begin{psgraph}[axesstyle=none,xticksize=-2 2,yticksize=-2 2,subticks=0](0,0)(-2,-2)(2,2){4cm}{4cm}
\PstPolygon[linestyle=dotted,] 
   \PstPolygon[PolyIntermediatePoint=1.5]
   \pscircle[linestyle=dotted,linecolor=red]{1.5}
\end{psgraph}
&  
\begin{psgraph}[axesstyle=none,xticksize=-2 2,yticksize=-2 2,subticks=0](0,0)(-2,-2)(2,2){4cm}{4cm}
\PstPolygon[linestyle=dotted,] 
   \PstPolygon[PolyIntermediatePoint=-0.5]
  \pscircle[linestyle=dotted,linecolor=red]{0.5}
\end{psgraph}
\\ \hline  
[\RDD{PolyIntermediatePoint}=0.5]
&
[\RDD{PolyIntermediatePoint}=1.5]
&
[\RDD{PolyIntermediatePoint}=-0.5] 
\\
\BS{pscircle}\AC{0.5}
&
\BS{pscircle}\AC{1.5}
&
\BS{pscircle}\AC{.5}
\\ \hline 
\end{tabular} 


\bigskip

\begin{tabular}{|c|c|c|} \hline 
\multicolumn{3}{|c|}{\BS{PstPolygon}[PolyIntermediatePoint=0.5,\RDD{PolyCurves}=true]  \RDI{PolyCurves}{pst-poly}  }
\\ \hline 
\begin{psgraph}[axesstyle=none,xticksize=-2 2,yticksize=-2 2,subticks=0](0,0)(-2,-2)(2,2){4cm}{4cm} 
\PstPolygon[linestyle=dotted,]
   \PstPolygon[PolyIntermediatePoint=0.5,PolyCurves=true]
   \pscircle[linestyle=dotted,linecolor=red]{0.5}
\end{psgraph}
&  
\begin{psgraph}[axesstyle=none,xticksize=-2 2,yticksize=-2 2,subticks=0](0,0)(-2,-2)(2,2){4cm}{4cm}
	\PstPolygon[linestyle=dotted,] 
   \PstPolygon[PolyIntermediatePoint=1.5,PolyCurves=true]
   \pscircle[linestyle=dotted,linecolor=red]{1.5}
\end{psgraph}
&  
\begin{psgraph}[axesstyle=none,xticksize=-2 2,yticksize=-2 2,subticks=0](0,0)(-2,-2)(2,2){4cm}{4cm}
\PstPolygon[linestyle=dotted,] 
   \PstPolygon[PolyIntermediatePoint=-0.5,PolyCurves=true]
  \pscircle[linestyle=dotted,linecolor=red]{0.5}
\end{psgraph}
\\ \hline  
[PolyIntermediatePoint=0.5]
&
[PolyIntermediatePoint=1.5]
&
[PolyIntermediatePoint=-0.5]
\\
\BS{pscircle}\AC{0.5}
&
\BS{pscircle}\AC{.5}
&
\BS{pscircle}\AC{.5}
\\ \hline 
\end{tabular} 


\bigskip



\begin{tabular}{|c|c|c|} \hline 
\multicolumn{3}{|c|}{\BS{PstPolygon}[PolyOffset=2,PolyNbSides=40,\RDD{PolyEpicycloid}=true]  \RDI{PolyEpicycloid}{pst-poly}  }
\\ \hline 
\begin{psgraph}[axesstyle=none,xticksize=-1.2 1.2,yticksize=-1.2 1.2,subticks=0](0,0)(-1.2,-1.2)(1.2,1.2){4cm}{4cm}
\PstPolygon[PolyNbSides=40,linestyle=dotted,] 
\PstPolygon[PolyOffset=2,PolyNbSides=40,PolyEpicycloid=true]
\end{psgraph}
&  
\begin{psgraph}[axesstyle=none,xticksize=-1.2 1.2,yticksize=-1.2 1.2,subticks=0](0,0)(-1.2,-1.2)(1.2,1.2){4cm}{4cm}
\PstPolygon[PolyNbSides=7,linestyle=dotted,] 
\PstPolygon[PolyOffset=3,PolyNbSides=40,PolyEpicycloid=true]
\end{psgraph}
&  
\begin{psgraph}[axesstyle=none,xticksize=-1.2 1.2,yticksize=-1.2 1.2,subticks=0](0,0)(-1.2,-1.2)(1.2,1.2){4cm}{4cm}
\PstPolygon[PolyNbSides=40,linestyle=dotted,] 
\PstPolygon[PolyOffset=4,PolyNbSides=40,PolyEpicycloid=true]
\end{psgraph}
\\ \hline  
PolyOffset=2
&
PolyOffset=3
&
PolyOffset=4
\\ \hline 
\begin{psgraph}[axesstyle=none,xticksize=-1.2 1.2,yticksize=-1.2 1.2,subticks=0](0,0)(-1.2,-1.2)(1.2,1.2){4cm}{4cm}
\PstPolygon[PolyNbSides=40,linestyle=dotted,] 
\PstPolygon[PolyOffset=5,PolyNbSides=40,PolyEpicycloid=true]
\end{psgraph}
&  
\begin{psgraph}[axesstyle=none,xticksize=-1.2 1.2,yticksize=-1.2 1.2,subticks=0](0,0)(-1.2,-1.2)(1.2,1.2){4cm}{4cm}
\PstPolygon[PolyNbSides=7,linestyle=dotted,] 
\PstPolygon[PolyOffset=40,PolyNbSides=40,PolyEpicycloid=true]
\end{psgraph}
&  
\begin{psgraph}[axesstyle=none,xticksize=-1.2 1.2,yticksize=-1.2 1.2,subticks=0](0,0)(-1.2,-1.2)(1.2,1.2){4cm}{4cm}
\PstPolygon[PolyNbSides=40,linestyle=dotted,] 
\PstPolygon[PolyOffset=41,PolyNbSides=40,PolyEpicycloid=true]
\end{psgraph}
\\ \hline  
PolyOffset=5
&
PolyOffset=40
&
PolyOffset=41
\\ \hline 

\end{tabular}

\bigskip

\begin{tabular}{|c|c|} \hline \multicolumn{2}{|c|}{\BS{PstPolygon}[unit=1,\RDD{PolyName}=A,PolyNbSides=6] \RDI{PolyName}{pst-poly} } \\
\multicolumn{2}{|c|}{\BS{PstPolygon}[unit=2,\RDD{PolyName}=B,PolyNbSides=6]}
\\ \hline  
\begin{psgraph}[axesstyle=none,xticksize=-2 2,yticksize=-2 2,subticks=0](0,0)(-2,-2)(2,2){4cm}{4cm} 
\PstPolygon[unit=1,PolyName=A,PolyNbSides=6]
\PstPolygon[unit=2,PolyName=B,PolyNbSides=12]
\psline[linecolor=red,linewidth=2pt](A1)(B1)
\end{psgraph}
&  
\begin{psgraph}[axesstyle=none,xticksize=-2 2,yticksize=-2 2,subticks=0](0,0)(-2,-2)(2,2){4cm}{4cm} 
\PstPolygon[unit=1,PolyName=A,PolyNbSides=6]
\PstPolygon[unit=2,PolyName=B,PolyNbSides=12]
\psline[linecolor=red,linewidth=2pt](A2)(B5)
\end{psgraph}
\\ \hline  
\BS{psline}(A1)(B1)
&
\BS{psline}(A2)(B5)  
\\ \hline 
\end{tabular} 

%-------------------------------------------

%\subsection{Polygones prédéfinis}
\SbSSCT{Polygones prédéfinis}{Predefined polygon}

\begin{tabular}{|c|c|c|c|} \hline 
\begin{psgraph}[axesstyle=none,xticksize=-1.2 1.2,yticksize=-1.2 1.2,subticks=0](0,0)(-1.2,-1.2)(1.2,1.2){2cm}{2cm} 
\PstTriangle 
\end{psgraph}
& 
\begin{psgraph}[axesstyle=none,xticksize=-1.2 1.2,yticksize=-1.2 1.2,subticks=0](0,0)(-1.2,-1.2)(1.2,1.2){2cm}{2cm}   
\PstSquare
\end{psgraph}
&  
\begin{psgraph}[axesstyle=none,xticksize=-1.2 1.2,yticksize=-1.2 1.2,subticks=0](0,0)(-1.2,-1.2)(1.2,1.2){2cm}{2cm} 
\PstPentagon
\end{psgraph}
& 
\begin{psgraph}[axesstyle=none,xticksize=-1.2 1.2,yticksize=-1.2 1.2,subticks=0](0,0)(-1.2,-1.2)(1.2,1.2){2cm}{2cm}   
\PstHexagon
\end{psgraph}
\\ \hline  
\BSS{PstTriangle}  \BSI{PstTriangle}{pst-poly} 
&  
\BSS{PstSquare}  \BSI{PstSquare}{pst-poly} 
&  
\BSS{PstPentagon}  \BSI{PstPentagon}{pst-poly} 
&  
\BSS{PstHexagon}  \BSI{PstHexagon}{pst-poly} 
\\ \hline 
\begin{psgraph}[axesstyle=none,xticksize=-1.2 1.2,yticksize=-1.2 1.2,subticks=0](0,0)(-1.2,-1.2)(1.2,1.2){2cm}{2cm}  
\PstHeptagon
\end{psgraph}
& 
\begin{psgraph}[axesstyle=none,xticksize=-1.2 1.2,yticksize=-1.2 1.2,subticks=0](0,0)(-1.2,-1.2)(1.2,1.2){2cm}{2cm} 
\PstOctogon
\end{psgraph}
& 
\begin{psgraph}[axesstyle=none,xticksize=-1.2 1.2,yticksize=-1.2 1.2,subticks=0](0,0)(-1.2,-1.2)(1.2,1.2){2cm}{2cm}   
\PstNonagon
\end{psgraph}
&
\begin{psgraph}[axesstyle=none,xticksize=-1.2 1.2,yticksize=-1.2 1.2,subticks=0](0,0)(-1.2,-1.2)(1.2,1.2){2cm}{2cm}   
\PstDecagon
\end{psgraph}
\\ \hline  
\BSS{PstHeptagon}  \BSI{PstHeptagon}{pst-poly} & 
\BSS{PstOctogon}  \BSI{PstOctogon}{pst-poly} & 
\BSS{PstNonagon}  \BSI{PstNonagon}{pst-poly} & 
\BSS{PstDecagon}  \BSI{PstDecagon}{pst-poly} \\ 
\hline 
\begin{psgraph}[axesstyle=none,xticksize=-1.2 1.2,yticksize=-1.2 1.2,subticks=0](0,0)(-1.2,-1.2)(1.2,1.2){2cm}{2cm}   
\PstDodecagon
\end{psgraph} 
&
\begin{psgraph}[axesstyle=none,xticksize=-1.2 1.2,yticksize=-1.2 1.2,subticks=0](0,0)(-1.2,-1.2)(1.2,1.2){2cm}{2cm}   
\PstStarFiveLines
\end{psgraph}
& 
\begin{psgraph}[axesstyle=none,xticksize=-1.2 1.2,yticksize=-1.2 1.2,subticks=0](0,0)(-1.2,-1.2)(1.2,1.2){2cm}{2cm}   
\PstStarFive
\end{psgraph}
&  

\\ \hline  
\BSS{PstDodecagon}  \BSI{PstDodecagon}{pst-poly} & \BSS{PstStarFiveLines} \BSI{PstStarFiveLines}{pst-poly} & \BSS{PstStarFive}  \BSI{PstStarFive}{pst-poly} &  \\ 
\hline 
\end{tabular} 

\bigskip

\begin{tabular}{|c|c|c|c|} \hline 
\begin{psgraph}[axesstyle=none,xticksize=-1.2 1.2,yticksize=-1.2 1.2,subticks=0](0,0)(-1.2,-1.2)(1.2,1.2){2cm}{2cm}  
\PstTriangle* 
\end{psgraph}
& 
\begin{psgraph}[axesstyle=none,xticksize=-1.2 1.2,yticksize=-1.2 1.2,subticks=0](0,0)(-1.2,-1.2)(1.2,1.2){2cm}{2cm}   
\PstSquare*
\end{psgraph}
&  
\begin{psgraph}[axesstyle=none,xticksize=-1.2 1.2,yticksize=-1.2 1.2,subticks=0](0,0)(-1.2,-1.2)(1.2,1.2){2cm}{2cm} 
\PstPentagon*
\end{psgraph}
& 
\begin{psgraph}[axesstyle=none,xticksize=-1.2 1.2,yticksize=-1.2 1.2,subticks=0](0,0)(-1.2,-1.2)(1.2,1.2){2cm}{2cm}  
\PstHexagon*
\end{psgraph}
\\ \hline  
\BSS{PstTriangle*}  \BSI{PstTriangle*}{pst-poly}
&  
\BSS{PstSquare*}  \BSI{PstSquare*}{pst-poly}
&  
\BSS{PstPentagon*}  \BSI{PstPentagon*}{pst-poly}
&  
\BSS{PstHexagon*}  \BSI{PstHexagon*}{pst-poly}
\\ \hline 
\begin{psgraph}[axesstyle=none,xticksize=-1.2 1.2,yticksize=-1.2 1.2,subticks=0](0,0)(-1.2,-1.2)(1.2,1.2){2cm}{2cm}  
\PstHeptagon*
\end{psgraph}
& 
\begin{psgraph}[axesstyle=none,xticksize=-1.2 1.2,yticksize=-1.2 1.2,subticks=0](0,0)(-1.2,-1.2)(1.2,1.2){2cm}{2cm}  
\PstOctogon*
\end{psgraph}
& 
\begin{psgraph}[axesstyle=none,xticksize=-1.2 1.2,yticksize=-1.2 1.2,subticks=0](0,0)(-1.2,-1.2)(1.2,1.2){2cm}{2cm} 
\PstNonagon*
\end{psgraph}
&
\begin{psgraph}[axesstyle=none,xticksize=-1.2 1.2,yticksize=-1.2 1.2,subticks=0](0,0)(-1.2,-1.2)(1.2,1.2){2cm}{2cm}   
\PstDecagon*
\end{psgraph}
\\ \hline  
\BSS{PstHeptagon*}  \BSI{PstHeptagon*}{pst-poly} & \BSS{PstOctogon*}  \BSI{PstOctogon*}{pst-poly} & \BSS{PstNonagon*}  \BSI{PstNonagon*}{pst-poly} & \BSS{PstDecagon*}  \BSI{PstDecagon*}{pst-poly} \\ 
\hline 
\begin{psgraph}[axesstyle=none,xticksize=-1.2 1.2,yticksize=-1.2 1.2,subticks=0](0,0)(-1.2,-1.2)(1.2,1.2){2cm}{2cm}  
\PstDodecagon*
\end{psgraph} 
&
\begin{psgraph}[axesstyle=none,xticksize=-1.2 1.2,yticksize=-1.2 1.2,subticks=0](0,0)(-1.2,-1.2)(1.2,1.2){2cm}{2cm}   
\PstStarFiveLines*
\end{psgraph}
& 
\begin{psgraph}[axesstyle=none,xticksize=-1.2 1.2,yticksize=-1.2 1.2,subticks=0](0,0)(-1.2,-1.2)(1.2,1.2){2cm}{2cm} 
\PstStarFive*
\end{psgraph}
&  

\\ \hline  
\BSS{PstDodecagon*}  \BSI{PstDodecagon*}{pst-poly} & \BSS{PstStarFiveLines*}  \BSI{PstStarFiveLines*}{pst-poly} & \BSS{PstStarFive*}  \BSI{PstStarFive*}{pst-poly} &  \\ 
\hline 
\end{tabular} 

%\section{Non regular polygons}

\begin{tabular}{|c|c|c|} \hline \multicolumn{3}{|c|}{\BS{PstStarFive}[\RDD{xunit}=1.5]  \RDI{xunit}{pst-poly}  } \\
 \hline  
\begin{psgraph}[axesstyle=none,xticksize=-2 2,yticksize=-2 2,subticks=0](0,0)(-2,-2)(2,2){4cm}{4cm} 
\PstStarFive[xunit=1.5]
\end{psgraph}
&  
\begin{psgraph}[axesstyle=none,xticksize=-2 2,yticksize=-2 2,subticks=0](0,0)(-2,-2)(2,2){4cm}{4cm} 
\PstStarFive[yunit=1.5]
\end{psgraph}
&  
\begin{psgraph}[axesstyle=none,xticksize=-2 2,yticksize=-2 2,subticks=0](0,0)(-2,-2)(2,2){4cm}{4cm} 
\PstStarFive[xunit=0.5,yunit=1.5]
\end{psgraph}
\\ \hline  
\RDD{xunit}=1.5
&
\RDD{yunit}=1.5   \RDI{yunit}{pst-poly} 
&
\RDD{xunit}=0.5,\RDD{yunit}=1.5
\\ \hline 
\end{tabular}

 %>>>voir conflit PstPicture pst-gantt

\newpage

\label{bezier}
%\section[Courbes de Bezier]{Courbes de Bezier  \cite{pst-bezier}}
\SSCT{Courbes de Bezier}{Bezier Curves}


\TFRGB{utilisation du module}{Package} \og \textbf{pst-bezier} \fg{} \label{pst-bezier}

%\subsection{Commande psbcurve}
\SbSSCT{Commande psbcurve}{psbcurve command}

\begin{tabular}{|c|}	\hline 
\begin{psgraph}[axesstyle=none,xticksize= 0 3 ,yticksize=0 5, subticks=0](0,0)(5,3){5cm}{3cm } 
\psbcurve(1,1)(2,2)(3,1)(4,2)
\end{psgraph}
\\ 	\hline  
\BSS{psbcurve}(1,1)(2,2)(3,1)(4,2)   \BSI{psbcurve}{pst-bezier}
\\ 	\hline 
\end{tabular} 

\bigskip
\begin{tabular}{|c|}	\hline 
\begin{psgraph}[axesstyle=none,xticksize= 0 3 ,yticksize=0 5, subticks=0](0,0)(5,3){5cm}{3cm } 
\psbcurve[showpoints=true](1,1)(2,2)(3,1)(4,2)
\end{psgraph}
\\ 	\hline  
\BS{psbcurve}{\red [showpoints=true]}(1,1)(2,2)(3,1)(4,2)
\\ 	\hline 
\end{tabular} 


%\subsection{Modificateurs }
\SbSSCT{Modificateurs}{Modificators}

\begin{tabular}{|c|c|} 	\hline  
\begin{psgraph}[axesstyle=none,xticksize= 0 3 ,yticksize=0 5, subticks=0](0,0)(5,3){5cm}{3cm } 
\psset{showpoints=true}
\psbcurve[linecolor=blue,linestyle=dashed](1,1)(2,2)(3,1)(4,2)
\psbcurve[linecolor=red](1,1)(2,2)l(2,1)(3,1)(4,2)
\end{psgraph}
&  
\begin{psgraph}[axesstyle=none,xticksize= 0 3 ,yticksize=0 5, subticks=0](0,0)(5,3){5cm}{3cm } 
\psset{showpoints=true}
\psbcurve[linecolor=blue,linestyle=dashed](1,1)(2,2)(3,1)(4,2)
\psbcurve[linecolor=red](1,1)(2,2)r(2,1)(3,1)(4,2)
\end{psgraph}
\\ 	\hline  
\BS{psbcurve}(1,1)(2,2){\red l(2,1)}(3,1)(4,2)
&  
\BS{psbcurve}(1,1)(2,2){\red r(2,1)}(3,1)(4,2)
\\ 	\hline 
\end{tabular} 


\bigskip

\begin{tabular}{|c|c|} 	\hline  
\begin{psgraph}[axesstyle=none,xticksize= 0 3 ,yticksize=0 5, subticks=0](0,0)(5,3){5cm}{3cm } 
\psset{showpoints=true}
\psbcurve[linecolor=blue,linestyle=dashed](1,1)(2,2)(3,1)(4,2)
\psbcurve[linecolor=red](1,1)(2,2)L(2,1)(3,1)(4,2)
\end{psgraph}
&  
\begin{psgraph}[axesstyle=none,xticksize= 0 3 ,yticksize=0 5, subticks=0](0,0)(5,3){5cm}{3cm } 
\psset{showpoints=true}
\psbcurve[linecolor=blue,linestyle=dashed](1,1)(2,2)(3,1)(4,2)
\psbcurve[linecolor=red](1,1)(2,2)L(3,2.5)(3,1)(4,2)
\end{psgraph}
\\ 	\hline  
\BS{psbcurve}(1,1)(2,2){\red L(2,1)}(3,1)(4,2)
&  
\BS{psbcurve}(1,1)(2,2){\red L(3,2.5)}(3,1)(4,2)
\\ 	\hline 
\end{tabular}

\bigskip

\begin{tabular}{|c|c|} \hline 
\begin{psgraph}[axesstyle=none,xticksize= 0 3 ,yticksize=0 5, subticks=0](0,0)(5,3){5cm}{3cm }
\psset{showpoints=true}
v\psbcurve[linecolor=blue,linestyle=dashed](1,1)(2,2)(3,1)(4,2) \psbcurve[linecolor=red](1,1)(2,2)tl{0.5}(3,1)(4,2)

\end{psgraph}
 &   
\begin{psgraph}[axesstyle=none,xticksize= 0 3 ,yticksize=0 5, subticks=0](0,0)(5,3){5cm}{3cm }
\psset{showpoints=true}
\psbcurve[linecolor=blue,linestyle=dashed](1,1)(2,2)(3,1)(4,2)
\psbcurve[linecolor=red](1,1)(2,2)tl{1.5}(3,1)(4,2)
\end{psgraph}
 \\ \hline 
\BS{psbcurve}(1,1)(2,2){\red tl\AC{0.5} }(3,1)(4,2)
&  
\BS{psbcurve}(1,1)(2,2){\red tl\AC{1.5} }(3,1)(4,2)
\\ 	\hline 
\end{tabular}

\bigskip

\begin{tabular}{|c|c|} \hline 
\begin{psgraph}[axesstyle=none,xticksize= 0 3 ,yticksize=0 5, subticks=0](0,0)(5,3){5cm}{3cm }
\psset{showpoints=true}
\psbcurve[linecolor=blue,linestyle=dashed](1,1)(2,2)(3,1)(4,2) \psbcurve[linecolor=red](1,1)(2,2)tr{0.5}(3,1)(4,2)
\end{psgraph}
&   
\begin{psgraph}[axesstyle=none,xticksize= 0 3 ,yticksize=0 5, subticks=0](0,0)(5,3){5cm}{3cm }
\psset{showpoints=true}
\psbcurve[linecolor=blue,linestyle=dashed](1,1)(2,2)(3,1)(4,2)
\psbcurve[linecolor=red](1,1)(2,2)tr{1.5}(3,1)(4,2)
\end{psgraph}
\\ \hline 
\BS{psbcurve}(1,1)(2,2){\red tr\AC{0.5} }(3,1)(4,2)
&  
\BS{psbcurve}(1,1)(2,2){\red tr\AC{1.5} }(3,1)(4,2)
\\ 	\hline 
\end{tabular}

\bigskip

\begin{tabular}{|c|c|} \hline 
\begin{psgraph}[axesstyle=none,xticksize= 0 3 ,yticksize=0 5, subticks=0](0,0)(5,3){5cm}{3cm }
\psset{showpoints=true}
\psbcurve[linecolor=blue,linestyle=dashed](1,1)(2,2)(3,1)(4,2) \psbcurve[linecolor=red](1,1)(2,2)ts{0.5}(3,1)(4,2)

\end{psgraph}
&   
\begin{psgraph}[axesstyle=none,xticksize= 0 3 ,yticksize=-1 5, subticks=0](0,0)(-1,0)(5,3){6cm}{3cm }
\psset{showpoints=true}
\psbcurve[linecolor=blue,linestyle=dashed](1,1)(2,2)(3,1)(4,2)
\psbcurve[linecolor=red](1,1)(2,2)ts{1.5}(3,1)(4,2)
\end{psgraph}
\\ \hline 
\BS{psbcurve}(1,1)(2,2){\red ts\AC{0.5} }(3,1)(4,2)
&  
\BS{psbcurve}(1,1)(2,2){\red ts\AC{1.5} }(3,1)(4,2)
\\ 	\hline 
\end{tabular}

\bigskip


\begin{tabular}{|c|c|} \hline 
\begin{psgraph}[axesstyle=none,xticksize= 0 3 ,yticksize=0 5, subticks=0](0,0)(5,3){5cm}{3cm }
\psset{showpoints=true}
\psbcurve[linecolor=blue,linestyle=dashed](1,1)(2,2)(3,1)(4,2) \psbcurve[linecolor=red](1,1)(2,2)T{0.5}(3,1)(4,2)
\end{psgraph}
&   
\begin{psgraph}[axesstyle=none,xticksize= 0 3 ,yticksize=0 6, subticks=0](0,0)(0,0)(6,3){6cm}{3cm }
\psset{showpoints=true}
\psbcurve[linecolor=blue,linestyle=dashed](1,1)(2,2)(3,1)(4,2)
\psbcurve[linecolor=red](1,1)(2,2)T{1.5}(3,1)(4,2)
\end{psgraph}
\\ \hline 
\BS{psbcurve}(1,1)(2,2){\red T\AC{0.5} }(3,1)(4,2)
&  
\BS{psbcurve}(1,1)(2,2){\red T\AC{1.5} }(3,1)(4,2)
\\ 	\hline 
\end{tabular}
\bigskip

\begin{tabular}{|c|c|} \hline 
\begin{psgraph}[axesstyle=none,xticksize= 0 3 ,yticksize=0 5, subticks=0](0,0)(5,3){5cm}{3cm }
\psset{showpoints=true}
\psbcurve[linecolor=blue,linestyle=dashed](1,1)(2,2)(3,1)(4,2) \psbcurve[linecolor=red](1,1)(2,2)tr{-0.5}(3,1)(4,2)
 \end{psgraph}
&   
\begin{psgraph}[axesstyle=none,xticksize= 0 3 ,yticksize=-1 5, subticks=0](0,0)(-1,0)(5,3){6cm}{3cm }
\psset{showpoints=true}
\psbcurve[linecolor=blue,linestyle=dashed](1,1)(2,2)(3,1)(4,2)
\psbcurve[linecolor=red](1,1)(2,2)tl{-0.5}(3,1)(4,2)
\end{psgraph}
\\ \hline 
\BS{psbcurve}(1,1)(2,2){\red tr\AC{-0.5} }(3,1)(4,2)
&  
\BS{psbcurve}(1,1)(2,2){\red tl\AC{-0.5} }(3,1)(4,2)
\\ 	\hline 
\begin{psgraph}[axesstyle=none,xticksize= 0 3 ,yticksize=0 5, subticks=0](0,0)(5,3){5cm}{3cm }
	\psset{showpoints=true}
	\psbcurve[linecolor=blue,linestyle=dashed](1,1)(2,2)(3,1)(4,2) \psbcurve[linecolor=red](1,1)(2,2)ts{-0.5}(3,1)(4,2)
	
\end{psgraph}
&   
\\ \hline 
\BS{psbcurve}(1,1)(2,2){\red ts\AC{-0.5} }(3,1)(4,2)
&  
\\ 	\hline
\end{tabular}
\bigskip

%\subsection{Paramètre bcurveTension}
\SbSSCT{Paramètre bcurveTension}{bcurveTension parameter}

\begin{tabular}{|c|c|} \hline
\multicolumn{2}{|c|}{\BS{psbcurve}[\RDD{bcurveTension}=1.5,linecolor=red](1,1)(2,2){\red ts\AC{0.5} }(3,1)(4,2) \RDI{bcurveTension}{pst-bezier}}  \\ \hline  
\begin{psgraph}[axesstyle=none,xticksize= 0 3 ,yticksize=0 5, subticks=0](0,0)(5,3){5cm}{3cm }
\psset{showpoints=true}
\psbcurve[linecolor=blue,linestyle=dashed](1,1)(2,2)(3,1)(4,2) \psbcurve[bcurveTension=0.5,linecolor=red](1,1)(2,2)(3,1)(4,2)

\end{psgraph}
&   
\begin{psgraph}[axesstyle=none,xticksize= 0 3 ,yticksize=-1 6, subticks=0](0,0)(-1,0)(6,3){7cm}{3cm }
\psset{showpoints=true}
\psbcurve[linecolor=blue,linestyle=dashed](1,1)(2,2)(3,1)(4,2)
\psbcurve[bcurveTension=1.5,linecolor=red](1,1)(2,2)(3,1)(4,2)
\end{psgraph}
\\ \hline 
bcurveTension=0.5
&  
bcurveTension=1.5
\\ 	\hline 
\end{tabular}



\newpage

%\section{Notion de chemin  PSTricks }
\SSCT{Notion de chemin  PSTricks }{Path PSTricks}

\label{chem}

\begin{tabular}{|c|c|} \hline  
\begin{pspicture}*[shift=*](-1,-1)(5,4)
\psgrid[gridcolor=lightgray,subgridcolor=white](0,0)(0,0)(9,4)
\pscustom[linecolor=blue]{ \psline[linecolor=red](0,0)(1,2)(2,2)  
\psarc[linestyle=dashed](2,2){2}{-60}{60} 
 \fill[fillstyle=solid,fillcolor=green] }
\end{pspicture}
&  
\parbox[c]{8cm}{
\BSS{pscustom}[{\red linecolor=blue}]\{ \\

 \BS{psline}[{\red linecolor=red} ](0,0)(1,2)\\

\BS{psarc}[{\red linestyle=dashed}](2,2)\AC{2}\{-60\}\{60\} \\ 


 \BS{fill}[fillstyle=solid,fillcolor=green] \\ 
 \} }
\\ \hline 
\end{tabular}
\bigskip


\begin{tabular}{|c|c|}\hline 
\multicolumn{2}{|l|}{\BSS{pscustom} \{ \BS{psline}(0,0)(1,2)(2,2) \BS{psarc}[\RDD{liftpen}=1](2,2)\AC{2}\AC{-60}\AC{60}  } \\
\multicolumn{2}{|l|}{
 \BS{fill}[fillstyle=solid,fillcolor=green] \} } \\ \hline
\begin{pspicture}*(-1,-1)(5,4)
\psgrid[gridcolor=lightgray,subgridcolor=white](0,0)(0,0)(9,4)
\pscustom{\psline(0,0)(1,2)(2,2)  
\psarc[liftpen=1](2,2){2}{-60}{60} 
 \fill[fillstyle=solid,fillcolor=green]}
\end{pspicture}
 &  
\begin{pspicture}*(-1,-1)(5,4)
\psgrid[gridcolor=lightgray,subgridcolor=white](0,0)(0,0)(9,4)
\pscustom[linecolor=blue]{ \psline[linecolor=red](0,0)(1,2)(2,2)  
\psarc[liftpen=2](2,2){2}{-60}{60} 
 \fill[fillstyle=solid,fillcolor=green]}
\end{pspicture} 
 \\ 
\hline \RDD{liftpen}=1 &  \RDD{liftpen}=2\\ 
\hline 
\end{tabular} 

\bigskip


\begin{tabular}{|c|c|} \hline 
\multicolumn{2}{|l|}{\BSS{pscustom}[fillstyle=\RDD{eofill},fillcolor=green]\{} \\
\multicolumn{2}{|l|}{\BS{pscircle}(2,2)\AC{2}
\BS{pscircle}(4,2)\AC{2}\}} \\ \hline
 \begin{pspicture}(6,4)
\pscustom[fillstyle=eofill,fillcolor=green]{
\pscircle(2,2){2}
\pscircle(4,2){2}}
\end{pspicture}
&  
 \begin{pspicture}(6,4)
\pscustom[fillstyle=oefill,fillcolor=green]{
\pscircle(2,2){2}
\pscircle(4,2){2}}
\end{pspicture}
\\ \hline 
fillstyle= \RDDD{eofill}{fillstyle} & fillstyle= \RDDD{oefill}{fillstyle}\\ 
\hline 
\end{tabular} 




%%%%% % % %===================================

\newpage
%\section{Les coordonnées }
\SSCT{Les coordonnées }{coordinates}
 
%%\section{Syntaxe}
%
%\subsection{Quadrillage avec psgrid}
\SbSSCT{Quadrillage avec psgrid}{Grids}
%\BS{psgrid}(x0,y0)(x1,y1)(x2,y2)

\psset{unit=1cm,linewidth=1pt}

\BSS{psgrid}\rnode{AA}{(2,1)}\rnode{BB}{(0,0)}\rnode{CC}{(4,3)}

\begin{center}
\begin{pspicture}(-.5,-.5)(4.5,3.5)
\psgrid(2,1)(0,0)(4,3)
\dotnode[dotstyle=*](2,1){A}
\dotnode[dotstyle=*](0,0){B}
\dotnode[dotstyle=*](4,3){C}
\ncarc{->}{AA}{A}
\ncbar[angle=-90]{->}{BB}{B}
\ncline{->}{CC}{C}
\end{pspicture}
\end{center}


\begin{tabular}{|c|c|c|c|c|}  \hline
\multicolumn{5}{|c|}{  \TFRGB{Quadrillage principal}{Main grid}} \\ \hline
\begin{pspicture}(-.5,-.5)(2.5,2.5)
\psgrid[gridwidth=2pt](2,2)
\end{pspicture}
&
\begin{pspicture}(-.5,-.5)(2.5,2.5)
\psgrid[griddots=3](2,2)
\end{pspicture}
&
\begin{pspicture}(-.5,-.5)(2.5,2.5)
\psgrid[gridcolor=red](2,2)
\end{pspicture}
&
\begin{pspicture}(-.5,-.5)(2.5,2.5)
\psgrid[gridlabels=5pt](2,2)
\end{pspicture}
&
\begin{pspicture}(-.5,-.5)(2.5,2.5)
\psgrid[gridlabelcolor=red](2,2)
\end{pspicture}



\\ \hline
\RDD{gridwidth}= 2pt & \RDD{griddots}=3 & \RDD{gridcolor}=red & \RDD{gridlabels}=5pt & \RDD{gridlabelcolor}=red  \\ \hline
{\blue \dft{} : .8pt } & {\blue \dft{} : 0}  & {\blue \dft{} : black} & {\blue \dft{}: 10pt} & {\blue \dft{} : black}\\ \hline
\end{tabular}


\TFRGB{Exemple}{Example} : 


\begin{tabular}{|c|} \hline  
\begin{pspicture}(-.5,-.5)(6.5,2.5)
\psgrid[gridwidth=4pt,gridcolor=red,griddots=3,gridlabelcolor=blue,gridlabels=15pt](6,2)
\end{pspicture}
\\ \hline  
\BS{psgrid}[{\red gridwidth=4pt,gridcolor=red,griddots=3,gridlabelcolor=blue,gridlabels=15pt}](6,2)
\\ \hline 
\end{tabular} 

\smallskip

\begin{tabular}{|l|l|l|l|}  \hline
\multicolumn{4}{|c|}{ \TFRGB{Quadrillage secondaire}{secondary grid}} \\ \hline
\begin{pspicture}(-.5,-.5)(2.5,2.5)
\psgrid[subgriddiv=3](2,2)
\end{pspicture}
&
\begin{pspicture}(-.5,-.5)(2.5,2.5)
\psgrid[subgridwidth=1pt](2,2)
\end{pspicture}
&
\begin{pspicture}(-.5,-.5)(2.5,2.5)
\psgrid[subgridcolor=red](2,2)
\end{pspicture}
&
\begin{pspicture}(-.5,-.5)(2.5,2.5)
\psgrid[subgriddots=3](2,2)
\end{pspicture}\\ \hline
\RDD{subgriddiv}=3 & \RDD{subgridwidth}=1pt & \RDD{subgridcolor}=red & \RDD{subgriddots}=3 \\ \hline
{\blue  \dft{} : 5 }& {\blue  \dft{} : .4pt } & {\blue \dft{} : gray} & {\blue \dft{}: 0 }\\ \hline

\end{tabular}


\smallskip
\begin{tabular}{|c|} \hline  
\begin{pspicture}(-.5,-.5)(6.5,2.5)
\psgrid[subgriddiv=2,subgridwidth=2pt,subgridcolor=red,subgriddots=4](6,2)
\end{pspicture}
\\ \hline  
\BS{psgrid}[{\red subgriddiv=2,subgridwidth=1pt,subgridcolor=red,subgriddots=3}](6,3)
\\ \hline 
\end{tabular} 






\input{Coord}

%%%%%======================================================== 
\newpage

%\section[Les n\oe uds]{Les n\oe uds    \cite{pst-node}}
\SSCT{Les n\oe uds}{Nodes}

Utilisation du module \textbf{pst-node} \label{pst-node}

\SbSSCT{Les types de  n\oe uds }{Types of nodes}
\label{nodes}

\psset{fillcolor=yellow,fillstyle=solid,linecolor=blue,unit=1cm}


\SbSbSSCTTC{Avec coordonnées}{Avec coordonnées   \protect \footnote{fillcolor=yellow,linecolor=blue} }{ With coordinates}{With coordinates  \protect \footnote{fillcolor=yellow,linecolor=blue} }

\begin{tabular}{|c|c|c|} \hline  
\begin{psgraph*}[axesstyle=none,xticksize= 0 1 ,yticksize=0 1 , subticks=0, dx=.5,Dx=.5, dy=.5,Dy=.5 ](0,0)(0,0)(1,1){2cm}{2cm }
\dotnode(.5,0.5){A}
\end{psgraph*}
&  
\begin{psgraph*}[axesstyle=none,xticksize= 0 1 ,yticksize=0 1 , subticks=0, dx=.5,Dx=.5, dy=.5,Dy=.5 ](0,0)(0,0)(1,1){2cm}{2cm }
\fnode(.5,0.5){B}
\end{psgraph*}
&  
\begin{psgraph*}[axesstyle=none,xticksize= 0 1 ,yticksize=0 1 , subticks=0, dx=.5,Dx=.5, dy=.5,Dy=.5 ](0,0)(0,0)(1,1){2cm}{2cm }
\fnode[framesize=.5 5pt](.5,0.5){B} 
\end{psgraph*}
\\ \hline  
\BSS{dotnode}(.5,0.5)\AC{A} \BSI{dotnode}{pst-node} 
& 
\BSS{fnode}(.5,0.5)\AC{B} \BSI{fnode}{pst-node} 
& 
\BSS{fnode}[\RDD{framesize}=.5 5pt](.5,0.5)\AC{B} 
\RDI{framesize}{pst-node}
 \\ \hline 
\begin{psgraph*}[axesstyle=none,xticksize= 0 1 ,yticksize=0 1 , subticks=0, dx=.5,Dx=.5, dy=.5,Dy=.5 ](0,0)(0,0)(1,1){2cm}{2cm }
\dotnode*(.5,0.5){A}
\end{psgraph*}
&  
\begin{psgraph*}[axesstyle=none,xticksize= 0 1 ,yticksize=0 1 , subticks=0, dx=.5,Dx=.5, dy=.5,Dy=.5 ](0,0)(0,0)(1,1){2cm}{2cm }
\fnode*(.5,0.5){B}
\end{psgraph*}
&  
\begin{psgraph*}[axesstyle=none,xticksize= 0 1 ,yticksize=0 1 , subticks=0, dx=.5,Dx=.5, dy=.5,Dy=.5 ](0,0)(0,0)(1,1){2cm}{2cm }
\fnode*[framesize=.5 5pt](.5,0.5){B} 
\end{psgraph*}
\\ \hline  
\BSS{dotnode}*(.5,0.5)\AC{A} \BSI{dotnode}{pst-node} 
& 
\BSS{fnode}*(.5,0.5)\AC{B} \BSI{fnode}{pst-node} 
& 
\BSS{fnode}*[\RDD{framesize}=.5 5pt](.5,0.5)\AC{B} 
\RDI{framesize}{pst-node} \\ \hline 
\end{tabular} 

\bigskip

\begin{tabular}{|c|c|c|} \hline  
\begin{psgraph*}[axesstyle=none,xticksize= 0 1 ,yticksize=0 1 , subticks=0, dx=.5,Dx=.5, dy=.5,Dy=.5 ](0,0)(0,0)(1,1){2cm}{2cm }
\pnode(.5,0.5){A}
\end{psgraph*}
&  
\begin{psgraph*}[axesstyle=none,xticksize= 0 1 ,yticksize=0 1 , subticks=0, dx=.5,Dx=.5, dy=.5,Dy=.5 ](0,0)(0,0)(1,1){2cm}{2cm }
\cnode(.5,0.5){.2cm}{A}
\end{psgraph*}
&
\begin{psgraph*}[axesstyle=none,xticksize= 0 1 ,yticksize=0 1 , subticks=0, dx=.5,Dx=.5, dy=.5,Dy=.5 ](0,0)(0,0)(1,1){2cm}{2cm }
\Cnode[radius=.2cm](.5,0.5){A}
\end{psgraph*}  
\\ \hline 
\BSS{pnode}(.5,0.5)\AC{A} & 
\BSS{cnode}(.5,0.5)\AC{.2cm}\AC{A} & 
\BSS{Cnode}[radius=.2cm](.5,0.5)\AC{A} 
\\ \hline 
%\begin{psgraph*}[axesstyle=none,xticksize= 0 1 ,yticksize=0 1 , subticks=0, dx=.5,Dx=.5, dy=.5,Dy=.5 ](0,0)(0,0)(1,1){2cm}{2cm }
%\pnode(.5,0.5){A}
%\end{psgraph*}
&  
\begin{psgraph*}[axesstyle=none,xticksize= 0 1 ,yticksize=0 1 , subticks=0, dx=.5,Dx=.5, dy=.5,Dy=.5 ](0,0)(0,0)(1,1){2cm}{2cm }
\cnode*(.5,0.5){.2cm}{A}
\end{psgraph*}
&
\begin{psgraph*}[axesstyle=none,xticksize= 0 1 ,yticksize=0 1 , subticks=0, dx=.5,Dx=.5, dy=.5,Dy=.5 ](0,0)(0,0)(1,1){2cm}{2cm }
\Cnode*[radius=.2cm](.5,0.5){A}
\end{psgraph*}  
\\ \hline 
%\BSS{pnode}*(.5,0.5)\AC{A} 
& 
\BSS{cnode}*(.5,0.5)\AC{.2cm}\AC{A} & 
\BSS{Cnode}*[radius=.2cm](.5,0.5)\AC{A} 
\\ \hline 
\end{tabular} 

\bigskip

\begin{tabular}{|c|c|c|} \hline  
\begin{psgraph*}[axesstyle=none,xticksize= 0 1 ,yticksize=0 1 , subticks=0, dx=.5,Dx=.5, dy=.5,Dy=.5 ](0,0)(0,0)(1,1){2cm}{2cm }
\psnode(.5,0.5){A}{contenu}
\end{psgraph*}
&  
\begin{psgraph*}[axesstyle=none,xticksize= 0 1 ,yticksize=0 1 , subticks=0, dx=.5,Dx=.5, dy=.5,Dy=.5 ](0,0)(0,0)(1,1){2cm}{2cm }
\cnodeput{45}(.5,0.5){M}{contenu}
\end{psgraph*}
&  
\begin{psgraph*}[axesstyle=none,xticksize= 0 1 ,yticksize=0 1 , subticks=0, dx=.5,Dx=.5, dy=.5,Dy=.5 ](0,0)(0,0)(1,1){2cm}{2cm }
\Cnodeput[radius=1cm]{45}(.5,0.5){M}{contenu}
\end{psgraph*}
\\ \hline 
\BSS{psnode}(.5,0.5)\AC{A}\AC{contenu} \BSI{psnode}{pst-node}
&
 \BSS{cnodeput}\AC{45}(.5,0.5) \BSI{cnodeput}{pst-node}  
& 
\BSS{Cnodeput}[radius=1cm]  \BSI{Cnodeput}{pst-node} \\
 & \AC{M}\AC{contenu}  & \AC{45}(2,0)\AC{M}\AC{contenu}
\\ \hline 
  
\begin{psgraph*}[axesstyle=none,xticksize= 0 1 ,yticksize=0 1 , subticks=0, dx=.5,Dx=.5, dy=.5,Dy=.5 ](0,0)(0,0)(1,1){2cm}{2cm }
\psnode*(.5,0.5){A}{contenu}
\end{psgraph*}
&  
\begin{psgraph*}[axesstyle=none,xticksize= 0 1 ,yticksize=0 1 , subticks=0, dx=.5,Dx=.5, dy=.5,Dy=.5 ](0,0)(0,0)(1,1){2cm}{2cm }
\cnodeput*{45}(.5,0.5){M}{contenu}
\end{psgraph*}
&  
\begin{psgraph*}[axesstyle=none,xticksize= 0 1 ,yticksize=0 1 , subticks=0, dx=.5,Dx=.5, dy=.5,Dy=.5 ](0,0)(0,0)(1,1){2cm}{2cm }
\Cnodeput*[radius=1cm]{45}(.5,0.5){M}{contenu}
\end{psgraph*}
\\ \hline 
\BSS{psnode}*(.5,0.5)\AC{A}\AC{contenu} \BSI{psnode}{pst-node}
&
 \BSS{cnodeput}*\AC{45}(.5,0.5) \BSI{cnodeput}{pst-node}  
& 
\BSS{Cnodeput}*[radius=1cm]  \BSI{Cnodeput}{pst-node} \\
 & \AC{M}\AC{contenu}  & \AC{45}(2,0)\AC{M}\AC{contenu}
\\ \hline 
\end{tabular} 

\vspace{1cm}

\begin{tabular}{|c|}\hline  
Texte avant \Cnodeput*[radius=1cm]{45}(0,0){M}{contenu} texte près
\\ \hline   \\ \\
Texte avant \BS{Cnodeput}*[radius=1cm]\AC{45}(0,0)\AC{M}\AC{contenu} texte près \\
{\red \TFRGB{Ces n\oe uds n'ont pas de dimension}{These nodes do not have dimension} !}
\\ \hline 
\end{tabular} 

 %-----------------------------------------------------------------
\SbSSCT{Sans coordonnées}{Without coordinates}

\begin{tabular}{|c|c|c|} \hline  
 \rnode{A}{contenu}
&  
 \Rnode{B}{contenu} 
&  
 \rnode{C}{\psframebox{contenu}}
\\ \hline 
 \BSS{rnode}\AC{A}\AC{contenu} \BSI{rnode}{pst-node}
& 
\BSS{Rnode}\AC{B}\AC{contenu} \BSI{Rnode}{pst-node}
& 
\BSS{rnode}\AC{C}\AC{\BS{psframebox}\AC{contenu}} \BSI{rnode}{pst-node}      
\\ \hline 

&  
 \Rnode*{B}{contenu} 
&  
 \rnode{C}{\psframebox*{contenu}}
\\ \hline 

& 
\BSS{Rnode}*\AC{B}\AC{contenu} \BSI{Rnode}{pst-node}
& 
\BSS{rnode}\AC{C}\AC{\BS{}psframebox*\AC{contenu}} \BSI{rnode}{pst-node}      
\\ \hline 
\end{tabular} 

\bigskip

\begin{tabular}{|c|c|} \hline  
 \trinode{A}{contenu}
&  
 \trinode*{B}{contenu} 
\\ \hline 
 \BSS{trinode}\AC{A}\AC{contenu} \BSI{trinode}{pst-node}
& 
\BSS{trinode}*\AC{B}\AC{contenu}       
\\ \hline 
\end{tabular} 

\bigskip

\begin{tabular}{|c|}\hline  
Texte avant \dianode{A}{contenu} texte près
\\ \hline  
Texte avant \BS{dianode}\AC{A}\AC{contenu} texte près
\\ \hline 
\end{tabular} 

\bigskip

\begin{tabular}{|c|}\hline
\begin{psgraph*}[axesstyle=none,xticksize= 0 2 ,yticksize=0 2 , subticks=0, dx=.5,Dx=.5, dy=.5,Dy=.5 ](0,0)(0,0)(2,2){4cm}{4cm }
\rput(1,1){\dianode{A}{contenu}}
\end{psgraph*}
\\ \hline  
\BSS{rput}(1,1)\AC{\BS{dianode}\AC{A}\AC{contenu}}
\\ \hline 
\end{tabular} 


%===================
\subsubsection{Création de n\oe uds multiples}

\begin{tabular}{|c|c|} 	\hline 
	\multicolumn{2}{|c|}{\BSS{dotnodes}[dotstyle=o](1, 0.5)\AC{A}(2, 2.5)\AC{B}(3, 1.5)\AC{C}}\\ \hline 
	\begin{psgraph}[axesstyle=none,xticksize=0 4,yticksize=0 4,subticks=0](0,0)(4,4){4cm}{4cm}  
		\dotnodes[dotstyle=o](1, 0.5){A}(2, 2.5){B}(3, 1.5){C}
	\end{psgraph}
	&
	\begin{psgraph}[axesstyle=none,xticksize=0 4,yticksize=0 4,subticks=0](0,0)(4,4){4cm}{4cm}  
		\dotnodes[dotstyle=o](1, 0.5){A}(2, 2.5){B}(3, 1.5){C}
	\end{psgraph} \\ 	\hline 
	\BS{dotnodes} & \BS{dotnodes*}  \\ 
	\hline  
\end{tabular} 
%%================================================
\newpage
\SbSSCT{Les connexions entre les n\oe uds}{Connections between nodes}
\label{liaisons}

\SbSbSSCTTC{Les types de connexions disponibles}{Les types de connexions disponibles \protect \footnote{fillcolor=white,linecolor=blue} }{Types of connections available} {Types of connections available \protect \footnote{fillcolor=white,linecolor=blue}}
 
\psset{fillcolor=white,fillstyle=none,linewidth=1pt}
\label{node}

\begin{tabular}{|l|c|c|} \hline

 &\textbf{ \TFRGB{sans astérisque}{without asterisk}} & \textbf{\TFRGB{avec astérisque}{with asterisk}}\\  \hline
 \BSS{ncline}\AC{->}\AC{A}\AC{B} \BSI{ncline}{pst-node}	
& \begin{pspicture}(0,-.5)(3 ,.5) \dotnode[dotstyle=*](1,-.5){A} \dotnode[dotstyle=*](2,0){B} \ncline{->}{A}{B}  \end{pspicture}	
&  \begin{pspicture}(0,-.5)(3 ,.5) \dotnode[dotstyle=*](1,-.5){A} \dotnode[dotstyle=*](2,0){B} \ncline*{->}{A}{B}  \end{pspicture}	\\ \hline

 \BSS{nccurve}\AC{->}\AC{A}\AC{B} \BSI{nccurve}{pst-node}	
& \begin{pspicture}(0,-.5)(3 ,.5) \dotnode[dotstyle=*](1,-.5){A} \dotnode[dotstyle=*](2,0){B} \nccurve{->}{A}{B}  \end{pspicture}	
&  \begin{pspicture}(0,-.5)(3 ,.5) \dotnode[dotstyle=*](1,-.5){A} \dotnode[dotstyle=*](2,0){B} \nccurve*{->}{A}{B}  \end{pspicture}	\\ \hline

	 \BSS{ncarc}\AC{->}\AC{A}\AC{B}	\BSI{ncarc}{pst-node}
& \begin{pspicture}(0,-.5)(3 ,.5) \dotnode[dotstyle=*](1,-.5){A} \dotnode[dotstyle=*](2,0){B} \ncarc{->}{A}{B} \end{pspicture} 	
& \begin{pspicture}(0,-.5)(3 ,.5) \dotnode[dotstyle=*](1,-.5){A} \dotnode[dotstyle=*](2,0){B} \ncarc*{->}{A}{B} \end{pspicture}  \\ \hline

 \BSS{ncbar}\AC{->}\AC{A}\AC{B} \BSI{ncbar}{pst-node}
& \begin{pspicture}(0,-.5)(3 ,.5) \dotnode[dotstyle=*](1,-.5){A} \dotnode[dotstyle=*](2,0){B}
 \ncbar{->}{A}{B} \end{pspicture} 	
 & \begin{pspicture}(0,-.5)(3 ,.5) \dotnode[dotstyle=*](1,-.5){A} \dotnode[dotstyle=*](2,0){B}
  \ncbar*{->}{A}{B} \end{pspicture}  \\ \hline

 \BSS{ncdiag}\AC{->}\AC{A}\AC{B} \BSI{ncdiag}{pst-node}
& \begin{pspicture}(0,-.5)(3 ,.5) \dotnode[dotstyle=*](1,-.5){A} \dotnode[dotstyle=*](2,0){B}
 \ncdiag{->}{A}{B} \end{pspicture} 	
 & \begin{pspicture}(0,-.5)(3 ,.5) \dotnode[dotstyle=*](1,-.5){A} \dotnode[dotstyle=*](2,0){B}
  \ncdiag*{->}{A}{B} \end{pspicture} \\ \hline

 \BSS{ncdiagg}\AC{->}\AC{A}\AC{B} \BSI{ncdiagg}{pst-node}	
& \begin{pspicture}(0,-.5)(3 ,.5) \dotnode[dotstyle=*](1,-.5){A} \dotnode[dotstyle=*](2,0){B}
 \ncdiagg{->}{A}{B} \end{pspicture} 	
 & \begin{pspicture}(0,-.5)(3 ,.5) \dotnode[dotstyle=*](1,-.5){A} \dotnode[dotstyle=*](2,0){B}
  \ncdiagg*{->}{A}{B} \end{pspicture}  \\ \hline

\BSS{ncangle}\AC{->}\AC{A}\AC{B} \BSI{ncangle}{pst-node}	
& \begin{pspicture}(0,-.5)(3 ,.5) \dotnode[dotstyle=*](1,-.5){A} \dotnode[dotstyle=*](2,0){B}
 \ncangle{->}{A}{B} \end{pspicture} 	
 & \begin{pspicture}(0,-.5)(3 ,.5) \dotnode[dotstyle=*](1,-.5){A} \dotnode[dotstyle=*](2,0){B}
  \ncangle*{->}{A}{B} \end{pspicture}  \\ \hline
 

 \BSS{ncangles}\AC{->}\AC{A}\AC{B} \BSI{ncangles}{pst-node}
& \begin{pspicture}(0,-.5)(3 ,.5) \dotnode[dotstyle=*](1,-.5){A} \dotnode[dotstyle=*](2,0){B}
 \ncangles{->}{A}{B} \end{pspicture} 	
 &  \begin{pspicture}(0,-.5)(3 ,.5) \dotnode[dotstyle=*](1,-.5){A} \dotnode[dotstyle=*](2,0){B}
  \ncangles*{->}{A}{B} \end{pspicture}\\ \hline


	\BSS{nccircle}\AC{->}\AC{A}\AC{.3cm} \BSI{nccircle}{pst-node}
& \begin{pspicture}(0,-.5)(3 ,.5) \dotnode[dotstyle=*](1,-.5){A} \dotnode[dotstyle=*](2,0){B}
 \nccircle{->}{A}{.3cm} \end{pspicture} 	
 & \begin{pspicture}(0,-.5)(3 ,.5) \dotnode[dotstyle=*](1,-.5){A} \dotnode[dotstyle=*](2,0){B}
  \nccircle*{->}{A}{.3cm} \end{pspicture} \\ \hline
 
	 \BSS{ncbox}\AC{->}\AC{A}\AC{B} \BSI{ncbox}{pst-node}
& boxsize 	
 & \begin{pspicture}(0,-.5)(3 ,.5) \dotnode[dotstyle=*](1,0){A} \dotnode[dotstyle=*](2,0){B}
  \ncbox*{A}{B} \end{pspicture} \\ \hline 
  
 \BSS{ncarcbox}\AC{->}\AC{A}\AC{B} \BSI{ncarcbox}{pst-node}
& \begin{pspicture}(0,-.5)(3 ,.5) \dotnode[dotstyle=*](1,0){A} \dotnode[dotstyle=*](2,0){B}
 \ncarcbox{A}{B} \end{pspicture} 	
 & \begin{pspicture}(0,-.5)(3 ,.5) \dotnode[dotstyle=*](1,0){A} \dotnode[dotstyle=*](2,0){B}
  \ncarcbox*{A}{B} \end{pspicture}\\ \hline 
 
	\BSS{ncloop}\AC{->}\AC{A}\AC{B} \BSI{ncloop}{pst-node}
& \begin{pspicture}(0,-.5)(3 ,.5) \dotnode[dotstyle=*](1,-.5){A} \dotnode[dotstyle=*](2,-.5){B}
 \ncloop[loopsize=.7cm]{A}{B} \end{pspicture} 	
 &  \begin{pspicture}(0,-.5)(3 ,.5) \dotnode[dotstyle=*](1,-.5){A} \dotnode[dotstyle=*](2,-.5){B}
  \ncloop*[loopsize=.7cm]{A}{B} \end{pspicture}   \\ \hline 

\end{tabular} 

\SbSbSSCT{Les noeuds comme des points de dessin }{Nodes as points}

\psset{fillcolor=white,fillstyle=none,linewidth=1pt}
%
\begin{tabular}{|l|c|c|} \hline
 & \textbf{ \TFRGB{sans astérisque}{without asterisk}} & \textbf{\TFRGB{avec astérisque}{with asterisk}}\\  \hline
\BSS{pcline}\AC{->}(A)(B) \BSI{pcline}{pst-node}	
&
 \begin{pspicture}(0,-.5)(3 ,.5)
  \dotnode[dotstyle=*](1,-.5){A} 
  \dotnode[dotstyle=*](2,0){B} 
  \pcline{->}(A)(B)
  \end{pspicture}	
& 
 \begin{pspicture}(0,-.5)(3 ,.5) \dotnode[dotstyle=*](1,-.5){A} \dotnode[dotstyle=*](2,0){B}
\pcline*{->}(A)(B)
\end{pspicture}
\\ \hline

  \BSS{pccurve}\AC{->}(A)(B) \BSI{pccurve}{pst-node}	
& \begin{pspicture}(0,-.5)(3 ,.5) \dotnode[dotstyle=*](1,-.5){A} \dotnode[dotstyle=*](2,0){B} \pccurve{->}(A)(B)  \end{pspicture}	
&  \begin{pspicture}(0,-.5)(3 ,.5) \dotnode[dotstyle=*](1,-.5){A} \dotnode[dotstyle=*](2,0){B} \pccurve*{->}(A)(B)  \end{pspicture}	\\ \hline

	 \BSS{pcarc }\AC{->}(A)(B) \BSI{pcarc}{pst-node}	
& \begin{pspicture}(0,-.5)(3 ,.5) \dotnode[dotstyle=*](1,-.5){A} \dotnode[dotstyle=*](2,0){B} \pcarc {->}(A)(B) \end{pspicture} 	
& \begin{pspicture}(0,-.5)(3 ,.5) \dotnode[dotstyle=*](1,-.5){A} \dotnode[dotstyle=*](2,0){B} \pcarc *{->}(A)(B) \end{pspicture}  \\ \hline

 \BSS{pcbar}\AC{->}(A)(B) \BSI{pcbar}{pst-node}
& \begin{pspicture}(0,-.5)(3 ,.5) \dotnode[dotstyle=*](1,-.5){A} \dotnode[dotstyle=*](2,0){B}
 \pcbar{->}(A)(B) \end{pspicture} 	
 & \begin{pspicture}(0,-.5)(3 ,.5) \dotnode[dotstyle=*](1,-.5){A} \dotnode[dotstyle=*](2,0){B}
  \pcbar*{->}(A)(B) \end{pspicture}  \\ \hline

 \BSS{pcdiag}\AC{->}(A)(B) \BSI{pcdiag}{pst-node}
& \begin{pspicture}(0,-.5)(3 ,.5) \dotnode[dotstyle=*](1,-.5){A} \dotnode[dotstyle=*](2,0){B}
 \pcdiag{->}(A)(B) \end{pspicture} 	
 & \begin{pspicture}(0,-.5)(3 ,.5) \dotnode[dotstyle=*](1,-.5){A} \dotnode[dotstyle=*](2,0){B}
  \ncdiag*{->}{A}{B} \end{pspicture} \\ \hline

 \BSS{pcdiagg }\AC{->}(A)(B) \BSI{pcdiagg}{pst-node}	
& \begin{pspicture}(0,-.5)(3 ,.5) \dotnode[dotstyle=*](1,-.5){A} \dotnode[dotstyle=*](2,0){B}
 \pcdiagg {->}(A)(B) \end{pspicture} 	
 & \begin{pspicture}(0,-.5)(3 ,.5) \dotnode[dotstyle=*](1,-.5){A} \dotnode[dotstyle=*](2,0){B}
  \pcdiagg *{->}(A)(B) \end{pspicture}  \\ \hline

\BSS{pcangle}\AC{->}(A)(B)	\BSI{pcangle}{pst-node}
& \begin{pspicture}(0,-.5)(3 ,.5) \dotnode[dotstyle=*](1,-.5){A} \dotnode[dotstyle=*](2,0){B}
 \pcangle{->}(A)(B) \end{pspicture} 	
 & \begin{pspicture}(0,-.5)(3 ,.5) \dotnode[dotstyle=*](1,-.5){A} \dotnode[dotstyle=*](2,0){B}
  \pcangle*{->}(A)(B) \end{pspicture}  \\ \hline
 

 \BSS{pcangles}\AC{->}(A)(B) \BSI{pcangles}{pst-node}
& \begin{pspicture}(0,-.5)(3 ,.5) \dotnode[dotstyle=*](1,-.5){A} \dotnode[dotstyle=*](2,0){B}
 \pcangles{->}(A)(B) \end{pspicture} 	
 &  \begin{pspicture}(0,-.5)(3 ,.5) \dotnode[dotstyle=*](1,-.5){A} \dotnode[dotstyle=*](2,0){B}
  \pcangles*{->}(A)(B) \end{pspicture}\\ \hline

	 \BSS{pcbox}\AC{->}(A)(B) \BSI{pcbox}{pst-node}
& \begin{pspicture}(0,-.5)(3 ,.5) \dotnode[dotstyle=*](1,0){A} \dotnode[dotstyle=*](2,0){B}
 \pcbox(A)(B) \end{pspicture} 	
 & \begin{pspicture}(0,-.5)(3 ,.5) \dotnode[dotstyle=*](1,0){A} \dotnode[dotstyle=*](2,0){B}
  \pcbox*(A)(B) \end{pspicture} \\ \hline 
  
 \BSS{pcarcbox}\AC{->}(A)(B) \BSI{pcarcbox}{pst-node}
& \begin{pspicture}(0,-.5)(3 ,.5) \dotnode[dotstyle=*](1,0){A} \dotnode[dotstyle=*](2,0){B}
 \pcarcbox(A)(B) \end{pspicture} 	
 & \begin{pspicture}(0,-.5)(3 ,.5) \dotnode[dotstyle=*](1,0){A} \dotnode[dotstyle=*](2,0){B}
  \pcarcbox*(A)(B) \end{pspicture}\\ \hline 
 
	\BSS{pcloop }\AC{->}(A)(B) \BSI{pcloop}{pst-node}
& \begin{pspicture}(0,-.5)(3 ,.5) \dotnode[dotstyle=*](1,-.5){A} \dotnode[dotstyle=*](2,-.5){B}
 \pcloop [loopsize=.7cm](A)(B) \end{pspicture} 	
 &  \begin{pspicture}(0,-.5)(3 ,.5) \dotnode[dotstyle=*](1,-.5){A} \dotnode[dotstyle=*](2,-.5){B}
  \pcloop *[loopsize=.7cm](A)(B) \end{pspicture}   \\ \hline 

\end{tabular} 
\newpage
%==================================================================
\SbSbSSCT{Les options disponibles}{Options available}

\psset{fillcolor=white}
\begin{tabular}{|c| c |c|} \hline
\multicolumn{3}{|c|}{ \BS{}ncline[nodesep=.3cm]\AC{->}\AC{A}\AC{B} } \\ \hline
  &  &  \\

\begin{pspicture}(0,-.5)(3 ,.5) \dotnode[dotstyle=*](1,-.5){A} \dotnode[dotstyle=*](2,.5){B} \ncline[nodesep=.3cm]{->}{A}{B}  \end{pspicture}
&
\begin{pspicture}(0,-.5)(3 ,.5) \dotnode[dotstyle=*](1,-.5){A} \dotnode[dotstyle=*](2,.5){B} \ncline[nodesepA=.2cm]{->}{A}{B}  \end{pspicture}	
&
\begin{pspicture}(0,-.5)(3 ,.5) \dotnode[dotstyle=*](1,-.5){A} \dotnode[dotstyle=*](2,.5){B} \ncline[nodesepB=.4cm]{->}{A}{B}  \end{pspicture}	
\\ \hline
\RDD{nodesep}=0.3cm \RDI{nodesep}{pst-node} & \RDD{nodesepA}=0.2cm \RDI{nodesepA}{pst-node} & \RDD{nodesepB}=0.4cm \RDI{nodesepB}{pst-node} \\ \hline
{\blue \dft{} : 0pt }& {\blue \dft{} : 0pt} & {\blue \dft{}: 0pt} \\  \hline
\end{tabular}
\bigskip

%---------------------------------------------------------
\begin{tabular}{|c| c |c|} \hline
\multicolumn{3}{|c|}{ \BS{}ncarc[arcangle=90]\AC{->}\AC{A}\AC{B} } \\ \hline

\begin{pspicture}(0,-.5)(3 ,.5) \dotnode[dotstyle=*](1,-.5){A} \dotnode[dotstyle=*](2,0){B}
  \ncarc[linestyle=dotted]{A}{B} 
 \ncarc[arcangle=90]{->}{A}{B} 
\end{pspicture}
&
\begin{pspicture}(0,-.5)(3 ,.5) \dotnode[dotstyle=*](1,-.5){A} \dotnode[dotstyle=*](2,0){B}
  \ncarc[linestyle=dotted]{A}{B}  
\ncarc[arcangleA=90]{->}{A}{B} 
\end{pspicture}	
&
\begin{pspicture}(0,-.5)(3 ,.5) \dotnode[dotstyle=*](1,-.5){A} \dotnode[dotstyle=*](2,0){B}
  \ncarc[linestyle=dotted]{A}{B}  
\ncarc[arcangleB=90]{->}{A}{B}
\end{pspicture}	
\\ \hline
\RDD{arcangle}=90 \RDI{arcangle}{pst-node}  & \RDD{arcangleA}=90 \RDI{arcangleA}{pst-node}& \RDD{arcangleB}=90 \RDI{arcangleB}{pst-node} \\ \hline
{\blue \dft{}: 8} &{\blue \dft{}: 8} &{\blue \dft{} : 8}\\  \hline
\multicolumn{3}{|c|}{\cyan \TFRGB{seulement pour}{only for} \BS{ncarc} ! }  \\ \hline

\end{tabular}
\bigskip

%---------------------------------------------------------
\begin{tabular}{|c| c |c|} \hline
\multicolumn{3}{|c|}{ \BS{}ncdiag[angle=90]\AC{->}\AC{A}\AC{B} } \\ \hline
\begin{pspicture}(0,-.5)(3 ,.5) \dotnode[dotstyle=*](1,-.5){A} \dotnode[dotstyle=*](2,0){B} 
\ncdiag[linestyle=dotted]{A}{B} 
\ncdiag[angle=90]{->}{A}{B}  
\end{pspicture}
&
\begin{pspicture}(0,-.5)(3 ,.5) \dotnode[dotstyle=*](1,-.5){A} \dotnode[dotstyle=*](2,0){B} 
\ncdiag[linestyle=dotted]{A}{B} 
\ncdiag[angleA=15]{->}{A}{B} 
\end{pspicture}	
&
\begin{pspicture}(0,-.5)(3 ,.5) \dotnode[dotstyle=*](1,-.5){A} \dotnode[dotstyle=*](2,0){B} 
\ncdiag[linestyle=dotted]{A}{B} 
\ncdiag[angleB=180]{->}{A}{B} 
\end{pspicture}	
\\ \hline
\RDD{angle}=90 \RDI{angle}{pst-node} & 
\RDD{angleA}=15 \RDI{angleA}{pst-node} & 
\RDD{angleB}=180 \RDI{angleB}{pst-node}\\ \hline
{\blue \dft{}: 0} & {\blue \dft{} : 0} & {\blue \dft{} : 0}\\  \hline
\end{tabular}
\bigskip

%---------------------------------------------------------
\begin{tabular}{|c| c |c|} \hline
\multicolumn{3}{|c|}{ \BS{}ncdiag[arm=1cm]\AC{->}\AC{A}\AC{B} } \\ \hline
\begin{pspicture}(0,-.5)(3 ,.5) \dotnode[dotstyle=*](1,-.5){A} \dotnode[dotstyle=*](2,0){B} 
\ncdiag[linestyle=dotted]{A}{B}
\ncdiag[arm=1cm]{->}{A}{B}
\end{pspicture}
&
\begin{pspicture}(0,-.5)(3 ,.5) \dotnode[dotstyle=*](1,-.5){A} \dotnode[dotstyle=*](2,0){B}
\ncdiag[linestyle=dotted]{A}{B} 
\ncdiag[armA=1cm]{->}{A}{B}
\end{pspicture}	
&
\begin{pspicture}(0,-.5)(3 ,.5) \dotnode[dotstyle=*](1,-.5){A} \dotnode[dotstyle=*](2,0){B} 
\ncdiag[linestyle=dotted]{A}{B}
\ncdiag[armB=1cm]{->}{A}{B}
\end{pspicture}	
\\ \hline
\RDD{arm}=1cm \RDI{arm}{pst-node} & 
\RDD{armA}=1cm \RDI{armA}{pst-node} & 
\RDD{armB}=1cm \RDI{armB}{pst-node}\\ \hline
{\blue \dft{} : 10pt} & {\blue \dft{} : 10pt} & {\blue \dft{} : 10pt} \\  \hline
\end{tabular}
\bigskip

%-------------------------------------------------------------------
\begin{tabular}{|c|c|c|} \hline
 \begin{pspicture}(0,-1)(3 ,.5) \dotnode[dotstyle=*](1,-.5){A} \dotnode[dotstyle=*](2,0){B} \ncline[offset=5pt]{->}{A}{B}\ncline[offset=5pt]{->}{B}{A}    
 \end{pspicture}
&
 \begin{pspicture}(0,-1)(3 ,.5) \dotnode[dotstyle=*](1,-.5){A} \dotnode[dotstyle=*](2,0){B} \ncline[offsetA=5pt]{->}{A}{B}\ncline[linestyle=dotted]{A}{B}    
 \end{pspicture}
 &
 \begin{pspicture}(0,-1)(3 ,.5) \dotnode[dotstyle=*](1,-.5){A} \dotnode[dotstyle=*](2,0){B} \ncline[offsetB=5pt]{->}{A}{B}\ncline[linestyle=dotted]{A}{B}    
 \end{pspicture}
 \\  \hline	
	
\BS{}ncline[\RDD{offset}=5pt]\AC{->}\AC{A}\AC{B}  & \BS{}ncline[\RDD{offsetA}=5pt]\AC{->}\AC{A}\AC{B}  & \BS{}ncline[\RDD{offsetB}=5pt]\AC{->}\AC{A}\AC{B} \\
\BS{}ncline[{\red offset}=5pt]\AC{->}\AC{B}\AC{A} & \BS{}ncline[linestyle=dotted]\AC{A}\AC{B} & \BS{}ncline[linestyle=dotted]\AC{A}\AC{B} 
\\ \hline    
{\blue \dft{} : 0pt }	\RDI{offset}{pst-node} & {\blue \dft{} : 0pt}	\RDI{offsetA}{pst-node} & {\blue  \dft{} : 0pt} \RDI{offsetB}{pst-node}\\ \hline
\end{tabular}

%----------------------------------------------------------------
\bigskip
\begin{minipage}[c]{.5\linewidth}
\begin{tabular}{|c|} \hline
\begin{pspicture}(0,-1)(3 ,2) \dotnode[dotstyle=*](1,-.5){A} \dotnode[dotstyle=*](2,-.5){B}
 \ncloop[linestyle=dotted]{A}{B} 
 \ncloop[loopsize=2cm]{A}{B} 
 \end{pspicture} 
 \\  \hline	
 \BS{ncloop}[\RDD{loopsize}=2cm]\AC{A}\AC{B}  \RDI{loopsize}{pst-node} \\  \hline	
 {\blue \dft{} : 1 cm} \\ \hline

\end{tabular}

\end{minipage}\hfill 
\begin{minipage}[c]{.5\linewidth}
\begin{tabular}{|c|} \hline	
   \\

  \begin{pspicture}(0,-1)(3 ,.5) \dotnode[dotstyle=*](1,-.5){A} \dotnode[dotstyle=*](2,0){B}
  \ncarcbox[linestyle=dotted]{A}{B}  
 \ncarcbox[boxsize=.2]{A}{B} 
 \end{pspicture} 	\\  \hline
 
 \BS{ncarcbox}[\RDD{boxsize}=.2]\AC{A}\AC{B} \RDI{boxsize}{pst-node} \\ \hline 
{\blue  \dft{}: 0.4cm}	 \\ \hline
{\cyan \TFRGB{seulement pour}{only for}   \BS{ncbox}  et  \BS{ncarcbox}  ! }\\ \hline
\end{tabular}
\end{minipage}

%-------------------------------------------------------------------


\bigskip

\begin{tabular}{|c|c|c|}\hline	
\begin{pspicture}(0,-1)(3 ,2) \dotnode[dotstyle=*](1,-.5){A} \dotnode[dotstyle=*](2,1.5){B} 
\nccurve[linestyle=dotted]{A}{B}
\nccurve[ncurv=1]{->}{A}{B} 
\end{pspicture}
&  
\begin{pspicture}(0,-1)(3 ,2) \dotnode[dotstyle=*](1,-.5){A} \dotnode[dotstyle=*](2,1.5){B} 
\nccurve[linestyle=dotted]{A}{B}
\nccurve[ncurvA=1]{->}{A}{B}  

\end{pspicture}
&  
\begin{pspicture}(0,-1)(3 ,2) \dotnode[dotstyle=*](1,-.5){A} \dotnode[dotstyle=*](2,1.5){B} 
\nccurve[linestyle=dotted]{A}{B}
\nccurve[ncurvB=1]{->}{A}{B}

\end{pspicture}
\\ 
\hline 	\BS{}nccurve[\RDD{ncurv}=1]\AC{->}\AC{A}\AC{B}   	& 	\BS{}nccurve[\RDD{ncurvA}=1]\AC{->}\AC{A}\AC{B}  	& 	\BS{}nccurve[\RDD{ncurvB}=1]\AC{->}\AC{A}\AC{B}  \\ \hline  
{\blue \dft{} : 0.67} \RDI{ncurv}{pst-node} & {\blue \dft{} : 0.67} \RDI{ncurvA}{pst-node} & {\blue \dft{} : 0.67 } \RDI{ncurvB}{pst-node} \\ \hline
 
\multicolumn{3}{|c|}{ \cyan \TFRGB{seulement pour}{only for}   \BS{nccurve}  et  \BS{pccurve} !  } \\ 
\hline 
\end{tabular} 


%++++++++++++++++++++++++++++++++++++++++++++++++++++++++
\bigskip


\begin{tabular}{|c|c|}
\hline  
\multicolumn{2}{|c|}{ \TFRGB{personnalisation des liaisons}{personalization of the connections}} \\ \hline
 \begin{pspicture}(0,-.5)(3 ,2.5) \dotnode[dotstyle=*](.5,0){A} \dotnode[dotstyle=*](2,2){B}
\ncdiagg[linearc=.3cm,doubleline=true,arrowscale=2]{->}{A}{B} 
\end{pspicture}
&  
 \begin{pspicture}(0,-.5)(3 ,2.5) \dotnode[dotstyle=*](.5,0){A} \dotnode[dotstyle=*](2,2){B}
\ncbar[linestyle=dashed,linewidth=3pt,dashcolor=red]{)-[}{A}{B} 
\end{pspicture}
\\ \hline 
\parbox[b]{6cm}{\BS{}ncdiagg[linearc=.3cm,doubleline=true,
arrowscale=2]\AC{->}\AC{A}\AC{B} } &
\parbox[b]{6cm}{\BS{}ncbar[linestyle=dashed,linewidth=3pt,
dashcolor=red]\AC{)-[}\AC{A}\AC{B}   }\\ 
\hline 
\end{tabular} 




%=======================================================================
\SbSSCT{Les étiquettes}{Labels}


\SbSbSSCTTC{Les étiquettes sur les n\oe uds}{Les étiquettes sur les n\oe uds \protect \footnote{fillcolor=yellow,linecolor=blue} }{Labels on the nodes}{ Labels on the nodes \protect \footnote{fillcolor=yellow,linecolor=blue}}

\psset{fillcolor=yellow}
 \begin{tabular}{|c| p{3cm}|p{8cm}|} \hline
 \multicolumn{3}{|c|}{syntaxe :  $\backslash$nput*[paramètres]\{position=angle\}\{nom\}\{texte\}  } \\  
\hline
\BSS{nput} \BSI{nput}{pst-node}& \rule[-1cm]{0pt}{2cm} \cnode(1,0){.2cm}{A} \nput{0}{A}{texte} & \verb \nput{0}{A}{texte} \\ \hline
\BS{nput}* & \rule[-1cm]{0pt}{2cm} \cnode(1,0){.2cm}{A} \nput*{0}{A}{texte} & \verb \nput*{0}{A}{texte} \\ \hline
 position=angle & \rule[-1cm]{0pt}{2cm} \cnode(1,0){.2cm}{A} \nput*{45}{A}{à 45}\nput*{-45}{A}{à -45} & \verb \nput*{45}{A} \{à 45\}\\ \hline
\RDD{labelsep} \RDI{labelsep}{pst-node} & \rule[-1cm]{0pt}{2cm} \cnode(1,0){.2cm}{A} \nput*[labelsep=0.5cm]{0}{A}{texte} & \verb \nput*[labelsep=0.5cm]{0}{A}{texte} \\ \hline

\RDD{labelsep} \RDI{labelsep}{pst-node} & \rule[-1cm]{0pt}{2cm} \cnode(1,0){.2cm}{A} \nput*[labelsep=-0.1cm]{0}{A}{texte} & \verb \nput*[labelsep=-0.1cm]{0}{A}{texte} \\ \hline

\RDD{rot} \RDI{rot}{pst-node} & \rule[-1cm]{0pt}{2cm} \cnode(1,0){.2cm}{A} \nput*[rot=45]{0}{A}{rot=45} & \verb \nput*[rot=45]{0}{A}{rot=45} \\ \hline
\end{tabular} 

 
\SbSbSSCT{Les étiquettes sur les connexions}{Labels on the connections} 

 \psset{shortput=nab,fillcolor=yellow}
\begin{tabular}{|p{3cm}| p{3cm}|p{8cm}|} \hline
\rput[l](0,1){ncput} \rput[l](0,0.25){nbput} \rput[l](0,1.75){naput} 	& \begin{pspicture}(0,-1)(3 ,1) \dotnode[dotstyle=*](0,0){A} \dotnode[dotstyle=*](3,0){B} \ncline{->}{A}{B}\naput[npos=.3]{naput} \ncput{ncput} \nbput[npos=.7]{nbput} \end{pspicture}	&
 \rput[l](0,.25){\BS{ncline}\AC{->}\AC{A}\AC{B}\BSS{nbput}[\RDD{npos}{\red =.7}]\AC{nbput} } \BSI{nbput}{pst-node} \RDI{npos}{pst-node}
 \rput[l](0,1){\BS{ncline}\AC{->}\AC{A}\AC{B}\BSS{ncput}\AC{ncput} } \BSI{ncput}{pst-node}
  \rput[l](0,1.75){\BS{ncline}\AC{->}\AC{A}\AC{B}\BSS{naput}[{\red npos=.3}]\AC{naput} } \BSI{naput}{pst-node}
 \\ \hline

 
\rput[l](0,1){ncput*} \rput[l](0,0.25){nbput*} \rput[l](0,1.75){naput*} 	& \begin{pspicture}(0,-1)(3 ,1) \dotnode[dotstyle=*](0,0){A} \dotnode[dotstyle=*](3,0){B} \ncline{->}{B}{A}\naput*[npos=.3]{naput*} \ncput*{ncput*} \nbput*[npos=.7]{nbput*} \end{pspicture}	&  \rput[l](0,.25){\BS{ncline}\AC{->}\AC{B}\AC{A}\BS{nbput}*[npos=.7]\AC{nbput} } 
\rput[l](0,1){\BS{ncline}\AC{->}\AC{B}\AC{A}\BS{ncput}*\AC{ncput} } 
\rput[l](0,1.75){\BS{ncline}\AC{->}\AC{B}\AC{A}\BS{naput}*[npos=.3]\AC{naput} } 
 \\ \hline

 \\ \hline
 
 \rput[l](0,1.5){[\RDD{nrot}=90]} \RDI{nrot}{pst-node}	& \begin{pspicture}(0,-1.5)(3 ,1.5) \dotnode[dotstyle=*](1.5,-1.2){A} \dotnode[dotstyle=*](1.5,1.2){B} \ncline{->}{B}{A}\naput*[nrot=90]{naput*} \ncput*[nrot=90]{ncput*} \nbput*[nrot=90]{nbput*} \end{pspicture}	&
 
   \rput[l](0,.5){\BS{ncline}\{->\}\{B\}\{A\}\BS{nbput}*[nrot=90]\AC{nbput} } 
  \rput[l](0,1.5){\BS{ncline}\AC{->}\AC{B}\AC{A}\BS{naput}*[nrot=90]\AC{naput} } 
   \rput[l](0,2.5){\BS{ncline}\AC{->}\AC{B}\AC{A}\BS{naput}*[nrot=90]\AC{naput} } 
  \\ \hline



\end{tabular} 





\newpage


\SbSSCT{Les Matrices de n\oe uds}{Mathematicals and graphs}

 \psset{fillcolor=yellow,fillstyle=solid}

\SbSbSSCT{Création de la matrice de noeuds}{Creation of the  diagram}

\begin{center}
\begin{tabular}{|c|c|}
\hline  
\psmatrix
 A &  B \\
 C &  D\\
\endpsmatrix 
& \parbox[b]{5cm}{ 
\BSS{psmatrix} \BSI{psmatrix}{pst-node}

 A \&  B \textbackslash{}\textbackslash{} 
 
 C \&  D \textbackslash{}\textbackslash{} 
 
\BSS{endpsmatrix}  \BSI{endpsmatrix}{pst-node}}
\\ 
\hline 
\end{tabular}
\end{center} 

	

\SbSbSSCT{Les 10 types de noeuds}{10 types of nodes}


\begin{tabular}{|c|c|c|c|c|}
\hline 
\psmatrix  [mnode= R] R \\ \endpsmatrix
&
\psmatrix [mnode= r] r \\ \endpsmatrix
&
\psmatrix [mnode= C] \\ \endpsmatrix
&
\psmatrix [mnode= f]  \\ \endpsmatrix
&
\psmatrix [mnode= p] p \\ \endpsmatrix
\\\hline 
 [\RDD{mnode}{\red = R}]  R \RDI{mnode}{pst-node} & [{\red mnode= r}] r & [{\red mnode= C}]  C & [{\red mnode= f}]  f & [{\red mnode= p}]  p  \\
  
\hline 
\psmatrix  [mnode= circle] circle \\ \endpsmatrix
&
\psmatrix [mnode= oval]  oval \\ \endpsmatrix
&
\psmatrix [mnode=dia] dia \\ \endpsmatrix
&
\psmatrix  [mnode= tri]  tri   \\ \endpsmatrix
&
\psmatrix [mnode= dot] dot \\ \endpsmatrix
\\\hline
 [{\red mnode= circle}] circle  & [{\red mnode= oval}]  oval & [{\red mnode=dia}] dia & [{\red mnode= tri}]  tri & [{\red mnode= dot}] dot 
\\\hline 
\end{tabular} 

\bigskip
\begin{tabular}{|c|c|} 	\hline  
\BS{psmatrix}[\RDD{mnode}=tri] & \BS{psmatrix}[\RDD{emnode}=tri] \RDI{emnode}{pst-node}\\
A \& B \& \& D \BS{}\BS{} & A \& B \& \& D \BS{}\BS{}\\ 
  \&  \& C \& E \BS{}\BS{} &    \&  \& C \& E  \BS{}\BS{} \\
\BS{endpsmatrix} & \BS{endpsmatrix} 	\\ 
\hline 
\psmatrix[mnode=tri]
A & B & & D\\
&  & C & E\\
\endpsmatrix 
& 
\psmatrix[emnode=tri]
A & B & & D\\
&  & C & E\\
\endpsmatrix  \\ 
	\hline 
\end{tabular} 



%-----------------------------------------------------
\SbSbSSCT{Liaison des noeuds}{Connection of the nodes}

\begin{center}
\begin{tabular}{|c|}
\hline 
\psmatrix[rowsep=.2cm]
%mnode= R & mnode= r & mnode= C & mnode= f & mnode= p \\
[mnode= R]  R & [mnode= r] r &[mnode= C]  C & [mnode= f]  f & [mnode= p]  p  \\
%\hline
%mnode= circle &  mnode= oval & mnode=dia & mnode= tri & mnode= dot \\
[mnode= circle] circle  & [mnode= oval]  oval & [mnode=dia] dia & [mnode= tri]  tri & [mnode= dot] dot \\
\ncline[linecolor=red]{1,2}{2,1}
\ncline[linestyle=dotted]{1,3}{2,2}
\ncline[linestyle=dotted]{1,4}{2,3}
\ncline[linestyle=dotted]{1,5}{2,4}
%\ncline[linecolor=red]{1,6}{2,5}

\ncline[linecolor=blue]{1,1}{2,1}
\ncline[linestyle=dotted]{1,2}{2,2}
\ncline[linestyle=dotted]{1,3}{2,3}
\ncline[linestyle=dotted]{1,4}{2,4}
\ncline[linestyle=dotted]{1,5}{2,5}

\ncline[linecolor=cyan]{1,1}{2,2}
\ncline[linestyle=dotted]{1,2}{2,3}
\ncline[linestyle=dotted]{1,3}{2,4}
\ncline[linestyle=dotted]{1,4}{2,5}
\endpsmatrix
 \\ \hline 
 \parbox[b]{8cm}{ 
\BS{}ncline[linecolor=red]\AC{1,2}\AC{2,1} 

\BS{}ncline[linecolor=blue]\AC{1,1}\AC{2,1}

\BS{}ncline[linecolor=cyan]\AC{1,1}\AC{2,2} }
 \\ 
\hline 
\end{tabular} 
\end{center}



%\newpage
%--------------------------------------------------------------
\SbSbSSCT{\'Etiquettes sur les liaisons}{Labels on connections}

\begin{center}
\begin{tabular}{|c|c|c|c|}  \hline 
& & & \\
\psmatrix[rowsep=1cm]
[mnode= dot]  R & [mnode= dot] r  \\
[mnode= dot] circle  & [mnode= dot]  oval \\
\ncline[linecolor=red]{1,2}{2,1} <{\red A} 
\ncline[linecolor=blue]{1,2}{2,2}<{\blue B} 
\ncline[linecolor=green]{2,1}{2,2}<{\green C} 
\endpsmatrix
&
\psmatrix[rowsep=1cm]
[mnode= dot]  R & [mnode= dot] r  \\
[mnode= dot] circle  & [mnode= dot]  oval \\
 \ncline[linecolor=red]{1,2}{2,1}_{\red A}
 \ncline[linecolor=blue]{1,2}{2,2}_{\blue B} 
 \ncline[linecolor=green]{2,1}{2,2}_{\green C} 
\endpsmatrix
&
\psmatrix[rowsep=1cm]
[mnode= dot]  R & [mnode= dot] r  \\
[mnode= dot] circle  & [mnode= dot]  oval \\
  \ncline[linecolor=red]{1,2}{2,1}>{\red A}
  \ncline[linecolor=blue]{1,2}{2,2}>{\blue B}
  \ncline[linecolor=green]{2,1}{2,2}>{\green C} 
%\ncline[linecolor=red]{1,2}{2,1}>{C}
\endpsmatrix
&
\psmatrix[rowsep=1cm]
[mnode= dot]  R & [mnode= dot] r  \\
[mnode= dot] circle  & [mnode= dot]  oval \\
   \ncline[linecolor=red]{1,2}{2,1}^{\red A}
   \ncline[linecolor=blue]{1,2}{2,2}^{\blue B}
   \ncline[linecolor=green]{2,1}{2,2}^{\green C}
%\ncline[linecolor=red]{1,2}{2,1}^{D}
\endpsmatrix\\
\hline
\BS{}ncline\AC{1,2}\AC{2,1}<\AC{A} & \BS{}ncline\AC{1,2}\AC{2,1}\_\AC{A} & \BS{}ncline\AC{1,2}\AC{2,1}>\AC{A} & \BS{}ncline\AC{1,2}\AC{2,1}$\wedge$\AC{A}\\
\BS{}ncline\AC{1,2}\AC{2,2}<\AC{B} & \BS{}ncline\AC{1,2}\AC{2,2}\_\AC{B} & \BS{}ncline\AC{1,2}\AC{2,2}>\AC{B} & \BS{}ncline\AC{1,2}\AC{2,2}$\wedge$\AC{B}\\
\BS{}ncline\AC{2,1}\AC{2,2}<\AC{C} & \BS{}ncline\AC{2,1}\AC{2,2}\_\AC{C}& \BS{}ncline\AC{2,1}\AC{2,2}>\AC{C} & \BS{}ncline\AC{2,1}\AC{2,2}$\wedge$\AC{C}\\
\hline

\end{tabular}
\end{center}



\SbSbSSCT{Les autres paramètres}{Other parameters }


%\psset{fillcolor=green}



\begin{center}
\begin{tabular}{|c|c|}
\hline 
\multicolumn{2}{|c|}{name}\\ \hline 
\psmatrix[mnode= oval]
[name=A] A &  [name=B] B\\
 [name=C] C &  [name=D] D\\
\endpsmatrix 
\ncline[linecolor=red]{A}{D}
\pcline[linecolor=blue](B)(C)
 &
\parbox[b]{8cm}{
\BS{}psmatrix[mnode= oval] 

[{\red \RDD{name}=A}] A \&  [{\red name=B}] B \BS{}\BS{} %\textbackslash{}\textbackslash{} \\
\RDI{name}{pst-node}

[{\red name=C}] C \&  [{\red name=D}] D \textbackslash{}\textbackslash{}

\BS{}endpsmatrix  

\BS{}ncline[linecolor=red]{\red \AC{A}\AC{D}} 

\BS{}pcline[linecolor=blue]{\red (B)(C)}}
\\ \hline 
\end{tabular}                                                                              
\end{center}


%-------------------------------------------------------------------

\bigskip

\begin{center}
\begin{tabular}{|c|c|}\hline 
\multicolumn{2}{|c|}{mcol   \hspace{1cm}   {\blue \dft{}:  mcol=c } }\\ \hline 

\psmatrix[rowsep=.2cm,colsep=.2cm]
paramètres  &  Position du noeud  \\
mcol=l & [mnode= oval,mcol=l] A \\ 
 mcol=c & [mnode= oval,mcol=c] B \\ 
 mcol=r & [mnode= oval,mcol=r] C \\
\endpsmatrix
 & 
 \parbox[b]{6cm}{
 
\BS{}psmatrix[rowsep=.2cm,colsep=.2cm]

paramètres  \&  Position du noeud   \textbackslash{}\textbackslash{}

mcol=l \& [mnode= oval,\RDD{mcol}{\red =l}] A \textbackslash{}\textbackslash{}  \RDI{mcol}{pst-node}

mcol=c \& [mnode= oval,{\red mcol=c}] B \textbackslash{}\textbackslash{} 

mcol=r \& [mnode= oval,{\red mcol=r}] C \textbackslash{}\textbackslash{} 

\BS{}endpsmatrix }  \\ \hline 
\end{tabular} 
\end{center}

%=========================================================

\bigskip


\begin{center}
\begin{tabular}{|c|}
\hline  
\multicolumn{1}{|c|}{radius  }\\ \hline 
\psmatrix
[mnode=C] A &  [mnode=C, radius=1cm] B
\endpsmatrix 
\\  \hline 
\BS{}psmatrix

 [mnode=C]  \& [mnode=C, \RDD{radius}{\red =1cm}] \RDI{radius}{pst-node}

\BS{}endpsmatrix 
\\ 
\hline 
\end{tabular} 
\end{center}



%-------------------------------------------------------------------------
 \bigskip

\begin{center}
\begin{tabular}{|c|}
\hline
\multicolumn{1}{|c|}{mnodesize \hspace{1cm} {\blue \dft{} : mnodesize= -1pt } }
\\ \hline 
\psmatrix[mnode=oval,rowsep=.2cm,colsep=.2cm]
 A & B & [mnodesize=4cm]  C & D & E\\
\endpsmatrix 
\\ \hline  
\parbox[b]{10cm}{
\BS{}psmatrix[mnode=oval,rowsep=.2cm,colsep=.2cm] 

 A \& B \& [ \RDD{mnodesize}{\red =4cm}]  C \& D \& E 
 \RDI{mnodesize}{pst-node}
 
\BS{}endpsmatrix }
\\ 
\hline 
\end{tabular} 
\end{center}

%-------------------------------------------------------------
 \bigskip


\begin{center}
\begin{tabular}{|c|}
\hline
colsep   \hspace{1cm}   {\blue \dft{} :  colsep= 1.5cm}
\\ \hline  
\psmatrix[mnode=oval]
A & [colsep=0cm] B & [colsep=4cm]  C & D & E\\
\endpsmatrix 
\\ \hline  
A \& [ \RDD{colsep}{\red =0cm}] B \& [{\red colsep=4cm}]  C \& D \& E  \textbackslash{}\textbackslash{} \RDI{colsep}{pst-node}
\\ 
\hline 
\end{tabular}
\end{center} 


%-------------------------------------------------------

 \bigskip

\begin{center}
\begin{tabular}{|c|c|c|}
\hline 
\multicolumn{3}{|c|}{rowsep \hspace{.5cm} {\blue \dft{}  : rowsep=  1.5cm } }
\\ \hline
\psmatrix[mnode=oval] 
A \\
B \\
C \\
\endpsmatrix 
&   
\psmatrix[mnode=oval,rowsep=0cm] 
A \\
B \\
C \\
\endpsmatrix 
& 
\psmatrix[mnode=oval,rowsep=1cm] 
A \\
B \\
C \\
\endpsmatrix
\\ 
\hline 
 \RDD{rowsep}{\red =0cm} \RDI{rowsep}{pst-node}
 &  
{\red rowsep=1cm} 
 &
\dft{} 
 \\ 
\hline 
\end{tabular} 
\end{center}

%-------------------------------------------------------
 \bigskip
 
\begin{center}
\begin{tabular}{|c|c|}
\hline  
\multicolumn{2}{|c|}{\BSS{psspan} \BSI{pspan}{pst-node} }\\ \hline 
\psmatrix[mnode=oval,rowsep=.2cm,colsep=.2cm]
A &  B &  C & D & E\\
A &  B &   C  \psspan{2} & D\\
\endpsmatrix 
&  
\parbox[b]{6cm}{
A \&  B \&  C \& D \& E \BS{} \BS{}

A \&  B \&   C  {\red \BS{}psspan\AC{2}} \& D \BS{} \BS{} }
\\ 

\hline 
\end{tabular} 
\end{center}


 
 \psset{fillcolor=yellow}



 \newpage
 
\SbSSCT{Commenter un élément}{comments to a graphic}
\label{com}
\psset{fillstyle=none}

\begin{tabular}{|c|c|}	\hline 
\begin{psgraph}[axesstyle=none,xticksize=0 3cm,yticksize=0 4cm,subticks=0](0,0)(4,3){4cm}{3cm} 
\psComment(1,2)(3,1){Texte}
\end{psgraph}
&
\begin{psgraph}[axesstyle=none,xticksize=0 3cm,yticksize=0 4cm,subticks=0](0,0)(4,3){4cm}{3cm} 
\psComment{->}(1,2)(3,1){Texte}
\end{psgraph}
\\ \hline 
\BSS{psComment}(1,2)(3,1)\AC{Texte} \BSI{psComment}{pstricks-add}& 
\BSS{psComment}{\red \AC{->} }(1,2)(3,1)\AC{Texte}
\\ 	\hline 
\end{tabular} 

\bigskip

\begin{tabular}{|c|c|c|} 	\hline  
\multicolumn{3}{|c|}{\BS{psComment}{\red[ref=r]} (1,2)(3,1)\AC{Texte}  } \\ \hline
\begin{psgraph}[axesstyle=none,xticksize=0 3cm,yticksize=0 4cm,subticks=0](0,0)(4,3){4cm}{3cm} 
\psComment[ref=r](1,2)(3,1){Texte}
\end{psgraph}
&
\begin{psgraph}[axesstyle=none,xticksize=0 3cm,yticksize=0 4cm,subticks=0](0,0)(4,3){4cm}{3cm}  
	\psComment*[ref=l](1,2)(3,1){Texte}
\end{psgraph}
&
\begin{psgraph}[axesstyle=none,xticksize=0 3cm,yticksize=0 4cm,subticks=0](0,0)(4,3){4cm}{3cm} 
\psComment[ref=B](1,2)(3,1){Texte}
\end{psgraph}
 \\ 	\hline  
[ref=r]  &  [ref=l] & [ref=B] \\ 
	\hline 
\end{tabular} 





\bigskip

\begin{tabular}{|c|c|c|} 	\hline  
\multicolumn{3}{|c|}{\BS{psComment}[ref=r](1,2)(3,1)\AC{Texte}{\red [\BS{ncarc}]} } \\ \hline
\begin{psgraph}[axesstyle=none,xticksize=0 3cm,yticksize=0 4cm,subticks=0](0,0)(4,3){4cm}{3cm} 
\psComment[ref=r](1,2)(3,1){Texte}[\ncarc] 
\end{psgraph}
&
\begin{psgraph}[axesstyle=none,xticksize=0 3cm,yticksize=0 4cm,subticks=0](0,0)(4,3){4cm}{3cm}  
\psComment*[ref=r](1,2)(3,1){Texte}[\ncdiagg]
\end{psgraph}
&
\begin{psgraph}[axesstyle=none,xticksize=0 3cm,yticksize=0 4cm,subticks=0](0,0)(4,3){4cm}{3cm} 
\psComment[ref=r](1,2)(3,1){Texte}[\ncbar]
\end{psgraph}
\\ 	\hline  
[\BS{ncarc}]  & [\BS{ncdiagg}]  & [\BS{ncbar}]  \\ 
\hline 
\end{tabular}

\bigskip
A voir : problème avec le deuxième paramètre final [\BS{ncput}]

 \newpage
%%======================================================

\SSCT{Constructions particulières}{Particular constructions}

%\section[Constructions particulières]{Constructions particulières \cite{pstricks-add} }


\TFRGB{Voir aussi le module  de géométrie page}{See also the package of geometry on page} \pageref{geom}

\SbSSCT{Création de n\oe uds multiples}{Multiples nodes creation}

\psset{linewidth=1pt,tickcolor=black!10}
\psset{unit=1cm,fillcolor=white,fillstyle=none,linecolor=blue}

\begin{tabular}{|l|l|} \hline  
\begin{psgraph}[axesstyle=none,xticksize=0 4,yticksize=0 4,subticks=0](0,0)(4,4){3cm}{3cm} 
\pnodes(3,1){A}(2,3){B} (1,2){C}
\psdots[dotstyle=*,linecolor=blue](A) \nput{45}{A}{A} 
\psdots[dotstyle=*,linecolor=blue](B) \nput{45}{B}{B} 
\psdots[dotstyle=*,linecolor=blue](C) \nput{135}{C}{C} 
\psline (A) (B) (C)
 \end{psgraph}
& 
\begin{psgraph}[axesstyle=none,xticksize=0 4,yticksize=0 4,subticks=0](0,0)(4,4){3cm}{3cm} 
\pnodes{A} (3,1) (2,3) (1,2)
\psdots[dotstyle=*,linecolor=blue](A0) \nput{45}{A0}{A0} 
\psdots[dotstyle=*,linecolor=blue](A1) \nput{45}{A1}{A1} 
\psdots[dotstyle=*,linecolor=blue](A2) \nput{135}{A2}{A2} 
\psline (A0) (A1) (A2)
 \end{psgraph}
\\ \hline 

\BSS{pnodes}(3,1)\AC{A}(2,3)\AC{B} (1,2)\AC{C} \BSI{pnodes}{pst-node} 
&
\BSS{pnodes}\AC{A} (3,1)(2,3) (1,2)  \\
\BS{psline} (A) (B) (C)
&  
\BS{psline} (A0) (A1) (A2)
\\ \hline 
\end{tabular} 

%-----------------------------------------------------------------------------
\SbSSCT{Positionement calculé de n\oe uds}{Node positions calculated }

\SbSbSSCT{Positions relatives avec psLNode}{Relative position width psLNode}

\begin{tabular}{|c|c|}\hline  
\begin{psgraph}[axesstyle=none,xticksize=0 2cm,yticksize=0 6,subticks=0](0,0)(6,2){6cm}{2cm}
\pnode(1,1){B}\pnode(5,1){C}
  \psLNode(B)(C){0.75}{A}
\psdots[dotstyle=*,linecolor=red](A)  \nput{45}{A}{A}%
\psdots[dotstyle=*,linecolor=blue](B) \nput{45}{B}{B} 
\psdots[dotstyle=*,linecolor=blue](C) \nput{45}{C}{C}%
 \end{psgraph}
&  
\begin{psgraph}[axesstyle=none,xticksize=0 2cm,yticksize=0 6cm,subticks=0](0,0)(6,2){6cm}{2cm}
\pnode(1,1){B}\pnode(5,1){C}
  \psLNode(B)(C){-0.25}{A}
\psdots[dotstyle=*,linecolor=red](A)  \nput{45}{A}{A}%
\psdots[dotstyle=*,linecolor=blue](B) \nput{45}{B}{B} 
\psdots[dotstyle=*,linecolor=blue](C) \nput{45}{C}{C}%
 \end{psgraph} 
\\ \hline  
\BSS{psLNode}(B)(C)\AC{{\red 0.75}}\AC{A} \BSI{pslNode}{pst-node} 
&  
\BSS{psLNode}(B)(C)\AC{{\red -0.25}}\AC{A}
\\ \hline 
\end{tabular}

\SbSbSSCT{Positions relatives avec midAB}{Relative position width midAB}

\begin{tabular}{|c|} \hline  
\begin{psgraph}[axesstyle=none,xticksize=0 2cm,yticksize=0 6cm,subticks=0](0,0)(6,2){6cm}{2cm}
\pnode(1,1){B}\pnode(5,1){C}
\midAB(B)(C){A}
\psdots[dotstyle=*,linecolor=red](A)  \nput{45}{A}{A}%
\psdots[dotstyle=*,linecolor=blue](B) \nput{45}{B}{B} 
\psdots[dotstyle=*,linecolor=blue](C) \nput{45}{C}{C}%
 \end{psgraph}
\\ \hline  
\BSS{midAB}(B)(C)\AC{A} \BSI{midAB}{pst-node} 
\\ \hline 
\end{tabular}

%-------------------------------
\SbSbSSCT{Positions  avec psLDNode}{Position width psLDNode}

\begin{tabular}{|c|} \hline  
\begin{psgraph}[axesstyle=none,xticksize=0 2cm,yticksize=0 6cm,subticks=0](0,0)(6,2){6cm}{2cm}
\pnode(1,1){B}\pnode(5,1){C}
\psLDNode(B)(C){1cm}{A} 
\psdots[dotstyle=*,linecolor=red](A)  \nput{45}{A}{A}%
\psdots[dotstyle=*,linecolor=blue](B) \nput{45}{B}{B} 
\psdots[dotstyle=*,linecolor=blue](C) \nput{45}{C}{C}%
\end{psgraph} 
\\ \hline  
\BSS{psLDNode}(B)(C)\AC{{\red 1cm}}\AC{A} \BSI{psLNDode}{pst-node}  
\\ \hline 
\end{tabular} 

%-------------------------------------

\SbSbSSCT{Positions relatives avec psLCNode}{Relative position width psLCNode}

\begin{tabular}{|c|c|c|} \hline  
\begin{psgraph}[axesstyle=none,xticksize=0 5,yticksize=0 5,subticks=0](0,0)(5,5){3cm}{3cm}
\psline{-D>}(4,1)
\psline{-D>}(1,2)
\psLCNode(4,1){1}(1,2){1}{A}
\psdots[dotstyle=*,linecolor=red](A) \nput{135}{A}{A} 
\psline[linestyle=dashed](4,1)(A)
\psline[linestyle=dashed](1,2)(A)
\end{psgraph}
&  
\begin{psgraph}[axesstyle=none,xticksize=0 5,yticksize=0 5,subticks=0](0,0)(5,5){3cm}{3cm}
\psline{-D>}(4,1)
\psline{-D>}(1,2)
\psLCNode(4,1){.5}(1,2){1}{A}
\psdots[dotstyle=*,linecolor=red](A) \nput{135}{A}{A} 
\psline[linestyle=dashed](2,.5)(A)
\psline[linestyle=dashed](1,2)(A)
\end{psgraph}
&  
\begin{psgraph}[axesstyle=none,xticksize=0 5,yticksize=0 6,subticks=0](0,0)(6,5){3cm}{3cm}
\psline[linestyle=dashed](1.5,3)
\psline{-D>}(4,1)
\psline{-D>}(1,2)
\psLCNode(4,1){1}(1,2){1.5}{A}
\psdots[dotstyle=*,linecolor=red](A) \nput{135}{A}{A} 
\psline[linestyle=dashed](4,1)(A)
\psline[linestyle=dashed](1.5,3)(A)
\end{psgraph}
\\ \hline  
\BSS{psLCNode}(4,1){\red \AC{1}}(1,2){\red \AC{1}}\AC{A} \BSI{psLCNode}{pst-node} & 
\BSS{psLCNode}(4,1){\red \AC{.5}}(1,2){\red \AC{1}}\AC{A} &
\BSS{psLCNode}(4,1){\red \AC{1}}(1,2){\red \AC{1.5}}\AC{A}  \\ 
\hline 
\end{tabular} 

%-----------------------------------

\SbSbSSCT{Positions relatives avec psLCNodeVar}{Relative position width psLCNodeVar}

\begin{tabular}{|c|c|} \hline  
\begin{psgraph}[axesstyle=none,xticksize=0 5,yticksize=0 4,subticks=0](0,0)(5,4){3cm}{3cm}
\psline{-D>}(4,1)
\psline{-D>}(1,2)
\psLCNodeVar(4,1)(1,2)(1,1){A}
\psdots[dotstyle=*,linecolor=red](A) \nput{135}{A}{A} 
\psline[linestyle=dashed](4,1)(A)
\psline[linestyle=dashed](1,2)(A)
\end{psgraph}
&
\begin{psgraph}[axesstyle=none,xticksize=0 4,yticksize=0 5,subticks=0](0,0)(5,4){3cm}{3cm}
\psline{-D>}(4,1)
\psline{-D>}(1,2)
\psLCNodeVar(4,1)(1,2)(0.5,1){A}
\psdots[dotstyle=*,linecolor=red](A) \nput{135}{A}{A} 
\psline[linestyle=dashed](2,.5)(A)
\psline[linestyle=dashed](1,2)(A)
\end{psgraph}
\\ \hline 
\BSS{psLCNodeVar}(4,1)(1,2)({\red 1,1})\AC{A} \BSI{psLCNodeVar}{pst-node}  &
\BSS{psLCNodeVar}(4,1)(1,2)({\red 0.5,1})\AC{A} 
\\ \hline  
\end{tabular} 

%----------------------------------

\SbSbSSCT{Positions relatives avec rhombus}{Relative position width rhombus}

\begin{tabular}{|c|c|} \hline  
\begin{psgraph}[axesstyle=none,xticksize=0 6,yticksize=0 6,subticks=0](0,0)(6,6){3cm}{3cm}
\pnodes(5,2){A}(2,4){B} 
\uput[-135](B){B}
\uput[-45](A){A}
\psdots[dotstyle=*,linecolor=blue](B)
\psdots[dotstyle=*,linecolor=blue](A)
\rhombus{2}(A)(B){C}{D}
\psdots[dotstyle=*,linecolor=red](C)
\psdots[dotstyle=*,linecolor=red](D)
\uput[45](C){C}
\uput[45](D){D}
\psline[linestyle=dotted ](A)(B)
\psline[linestyle=dotted ](A)(C)
\psline[linestyle=dotted ](A)(D)
\psline[linestyle=dotted ](B)(C)
\psline[linestyle=dotted ](B)(D)
\psarc[linestyle=dashed,linewidth=1pt](A){2}{80}{210}
\end{psgraph}
&  
\begin{psgraph}[axesstyle=none,xticksize=0 6,yticksize=0 6,subticks=0](0,0)(6,6){3cm}{3cm}
\pnodes(5,2){A}(2,4){B} 
\uput[-135](B){B}
\uput[-45](A){A}
\psdots[dotstyle=*,linecolor=blue](B)
\psdots[dotstyle=*,linecolor=blue](A)
\rhombus{3}(A)(B){X}{Y}
\psdots[dotstyle=*,linecolor=red](X)
\psdots[dotstyle=*,linecolor=red](Y)
\psline[linestyle=dotted ](A)(B)
\psline[linestyle=dotted ](A)(X)
\psline[linestyle=dotted ](A)(Y)
\psline[linestyle=dotted ](B)(X)
\psline[linestyle=dotted ](B)(Y)
\uput[45](X){X}
\uput[45](Y){Y}
\psline[linestyle=dotted ](A)(B)
\psarc[linestyle=dashed,linewidth=1pt](A){3}{80}{210}
\end{psgraph}
\\ \hline  
\BSS{rhombus}\AC{{\red 2}}(A)(B)\AC{C}\AC{D} \BSI{rhombus}{pst-node} 
&  
\BSS{rhombus}\AC{{\red 3}}(A)(B)\AC{X}\AC{Y}
\\ \hline 
\end{tabular}
 
%-----------------------------------

\SbSbSSCT{Positions relatives avec psRelNodeVar}{Relative position width psRelNodeVar}

\begin{tabular}{|c|c|c|} \hline  
\begin{psgraph}[axesstyle=none,xticksize=0 5,yticksize=0 5,subticks=0](0,0)(5,5){3cm}{3cm}
 \pnode(0,1){B}
 \pnode(2,1){C}
\psRelNodeVar(B)(C)(1;45){A}
 \psline(B)(C)
\psline[linestyle=dotted](B)(A)
\psline[linestyle=dotted](C)(A)
\psarc[linestyle=dashed](B){2}{-20}{60}
\uput[-45](B){B}
\uput[45](A){A}
\uput[-45](C){C}
\psdots[dotstyle=*,linecolor=blue](B)
\psdots[dotstyle=*,linecolor=red](A)
\psdots[dotstyle=*,linecolor=blue](C)
 \end{psgraph}
&  
\begin{psgraph}[axesstyle=none,xticksize=0 5,yticksize=0 5,subticks=0](0,0)(5,5){3cm}{3cm}
\pnode(0,1){B}
\pnode(2,1){C}
\psRelNodeVar(B)(C)(2;45){A}
\psline(B)(C)
\psline[linestyle=dotted](B)(A)
\psarc[linestyle=dashed](B){4}{-20}{60}
\uput[-45](B){B}
\uput[45](A){A}
\uput[-45](C){C}
\psdots[dotstyle=*,linecolor=blue](B)
\psdots[dotstyle=*,linecolor=red](A)
\psdots[dotstyle=*,linecolor=blue](C)
 \end{psgraph}
&  
\begin{psgraph}[axesstyle=none,xticksize=0 5,yticksize=0 5,subticks=0](0,0)(5,5){3cm}{3cm}
\pnode(0,1){B}
\pnode(2,1){C}
\psRelNodeVar(B)(C)(2;30){A}
\psline(B)(C)
\psline[linestyle=dotted](B)(A)
\psarc[linestyle=dashed](B){4}{-20}{60}
\uput[-45](B){B}
\uput[45](A){A}
\uput[-45](C){C}
\psdots[dotstyle=*,linecolor=blue](B)
\psdots[dotstyle=*,linecolor=red](A)
\psdots[dotstyle=*,linecolor=blue](C)
\end{psgraph}
\\ \hline
\BSS{psRelNodeVar}(B)(C)({\red 1;45})\AC{A} \BSI{psRelNodeVar}{pst-node} 
&
\BSS{psRelNodeVar}(B)(C)({\red 2;45})\AC{A}
&  
\BSS{psRelNodeVar}(B)(C)({\red 2;30})\AC{A}
\\ \hline 
\end{tabular} 

%-----------------------------------

\SbSbSSCT{Positions relatives avec AplusB}{Relative position width AplusB} 

\begin{tabular}{|c|} \hline  
\begin{psgraph}[axesstyle=none,xticksize=0 4,yticksize=0 4,subticks=0](0,0)(4,4){3cm}{3cm}
\pnode(1,2){B}\pnode(2,1){C}
\pnode(0,0){O}
\AplusB(B)(C){A}
\psdots[dotstyle=*,linecolor=blue](B)
\psdots[dotstyle=*,linecolor=blue](C)
\psdots[dotstyle=*,linecolor=red](A)
\nput{45}{A}{A}%
\nput{135}{B}{B} 
\nput{135}{C}{C}%
\ncline {O} {B}
\ncline{O} {C}
\ncline {O} {B}
\ncline[linestyle=dashed]{->} {B} {A}
\ncline[linestyle=dashed]{->} {C} {A}
 \end{psgraph}
\\ \hline  
\BSS{AplusB}(B)(C)\AC{A} \BSI{AplusB}{pst-node} 
\\ \hline 
\end{tabular} 

%-----------------------------------

\SbSbSSCT{Positions relatives avec AtoB}{Relative position width AtoB} 

\begin{tabular}{|c|}\hline 
\begin{psgraph}[axesstyle=none,xticksize=-2 3,yticksize=-1 3,subticks=0](0,0)(-1,-2)(3,3){4cm}{4cm}
\pnode(1,2){B}\pnode(2,1){C}
\pnode(0,0){O}
\AtoB(B)(C){A}
\psdots[dotstyle=*,linecolor=blue](B)
\psdots[dotstyle=*,linecolor=blue](C)
\psdots[dotstyle=*,linecolor=red](A)
\nput{45}{A}{A}%
\nput{45}{B}{B} 
\nput{45}{C}{C}%
\ncline {O} {B}
\ncline {O} {C}
\ncline {O} {B}
\ncline[linestyle=dashed]{->} {B} {C}
\ncline[linestyle=dashed]{->} {O} {A}
\end{psgraph}
\\ \hline  
\BSS{AtoB}(B)(C)\AC{A} \BSI{AtoB}{pst-node} 
\\ \hline 
\end{tabular} 



%==================================================

\SbSSCT{N\oe ud sur une courbe}{Node on a curve} 

\SbSbSSCT{N\oe ud sur une courbe avec fnpnode  }{Node on a curve with fnpnode } 

\begin{tabular}{|l|} \hline  
\begin{psgraph}[axesstyle=frame,xticksize=-1.5 1.5cm,yticksize=0 8cm,subticks=0](0,0)(0,-1.5)(13,1.5){8cm}{3cm}
\psset{algebraic}
\psplot[plotpoints=200,linewidth=2pt]{0}{13}{sin(x)}
\fnpnode{2}{sin(x)}{A}
\fnpnode{10}{sin(x)}{B}
\psdots[dotstyle=*,linecolor=red](A) \uput[-135](A){A}
\psdots[dotstyle=*,linecolor=red](B) \uput[45](B){B}
\psline[linestyle=dashed] (A) (B)
 \end{psgraph}
\\ \hline  
\BS{psplot}[plotpoints=200,linewidth=2pt]{0}\AC{13}{sin(x)}\\
\BSS{fnpnode}\AC{2}\AC{sin(x)}\AC{A} \BSI{fnpnode}{pst-node} \\
\BSS{fnpnode}\AC{10}\AC{sin(x)}\AC{B}\\
\BS{psline}[linestyle=dashed] (A) (B)
\\ \hline 
\end{tabular} 

%--------------------------------------

\SbSbSSCT{N\oe uds sur une courbe avec fnpnodes  }{Nodes on a curve with fnpnodes } 

\begin{tabular}{|l|} \hline  
\begin{psgraph}[axesstyle=none,xticksize=-1.5 1.5cm,yticksize=0 8cm,subticks=0](0,0)(0,-1.5)(13,1.5){8cm}{3cm}
\psset{algebraic}
\psplot[plotpoints=200,linewidth=2pt]{0}{13}{sin(x)}
\fnpnodes[plotpoints=14]{0}{13}{sin(x)}{A}
\psdots[dotstyle=*,linecolor=red](A3)
\psdots[dotstyle=*,linecolor=red](A4)
\psdots[dotstyle=*,linecolor=red](A5)
\psdots[dotstyle=*,linecolor=red](A6)
\psdots[dotstyle=*,linecolor=red](A7)
\psdots[dotstyle=*,linecolor=red](A8)
\psdots[dotstyle=*,linecolor=red](A9)
\psdots[dotstyle=*,linecolor=red](A11)
\psdots[dotstyle=*,linecolor=red](A12)
\psdots[dotstyle=*,linecolor=red](A13)
\psdots[dotstyle=*,linecolor=red](A14)
\psdots[dotstyle=*,linecolor=red](A1)
\psdots[dotstyle=*,linecolor=red](A2) \uput[45](A2){A2}
\psdots[dotstyle=*,linecolor=red](A10) \uput[45](A10){A10}
\psline[linestyle=dashed]  (A2) (A10)
 \end{psgraph}
\\ \hline  
\BS{psplot}[plotpoints=200,linewidth=2pt]{0}\AC{13}{sin(x)} \\
\BSS{fnpnodes}[plotpoints=14]\AC{0}\AC{13}\AC{sin(x)}\AC{A} \BSI{fnpnodes}{pst-node}  \\
\BS{psline}[linestyle=dashed]  (A2) (A10)
\\ \hline 
\end{tabular} 
 
 
%--------------------------------------

\SbSbSSCT{N\oe ud sur une courbe paramétrique avec curvepnode  }{Node on a parametric curve with curvepnode } 

\begin{tabular}{|l|} \hline  
\begin{psgraph}[axesstyle=none,xticksize=-1.2 1.2 ,yticksize=-1.2 1.2 , subticks=0](0,0)(-1.2,-1.2)(1.2,1.2){4cm}{4cm}
\psset{algebraic}
\parametricplot[plotpoints=200,linewidth=2pt]{0}{6.28}{sin(t)|sin(2*t)}
\curvepnode{2}{sin(t)|sin(2*t)}{A}
\curvepnode{6}{sin(t)|sin(2*t)}{B}
\psdots[dotstyle=*,linecolor=red](A)  \uput[-45](A){A}
\psdots[dotstyle=*,linecolor=red](B)  \uput[135](B){B}
\psline[linestyle=dashed] (A) (B)
\end{psgraph}
\\ \hline  
\BS{parametricplot}[plotpoints=200]\AC{0}\AC{6.28}\AC{sin(t)|sin(2*t)}\\
\BSS{curvepnode}\AC{2}\AC{sin(t)|sin(2*t)}\AC{A} \BSI{curvepnode}{pst-node} \\
\BSS{curvepnode}\AC{6}\AC{sin(t)|sin(2*t)}\AC{B}\\
\BS{psline}[linestyle=dashed] (A) (B)
\\ \hline 
\end{tabular} 

\bigskip 

\paragraph{Création automatique d'un n\oe ud pour la tangente} :

\begin{tabular}{|c|c|} \hline  
\begin{psgraph}[axesstyle=none,xticksize=-1.5 1.5 ,yticksize=-1.5 1.5 , subticks=0](0,0)(-1.5,-1.5)(1.5,1.5){4cm}{4cm}
\psset{algebraic}
\parametricplot[plotpoints=200]{0}{6.28}{sin(t)|sin(2*t)}
\curvepnode{2}{sin(t)|sin(2*t)}{A}
\psxline[linewidth=2pt,linecolor=red]{<->}(A){-(Atang)}{(Atang)}
\uput[-45](A){A} 
\psline[linewidth=2pt,linecolor=red] (0,0) (Atang)
\psdots[dotstyle=*,linecolor=blue](Atang)
\psdots[dotstyle=*,linecolor=blue](A) 
\uput[-90](Atang){Atang}
\end{psgraph}
&  
\begin{psgraph}[axesstyle=none,xticksize=-1.5 1.5 ,yticksize=-1.5 1.5 , subticks=0](0,0)(-1.5,-1.5)(1.5,1.5){4cm}{4cm}
\psset{algebraic}
\parametricplot[plotpoints=200]{0}{6.28}{sin(t)|sin(2*t)}
\curvepnode{4}{sin(t)|sin(2*t)}{X}
\psxline[linewidth=2pt,linecolor=red]{<->}(X){-1.5(Xtang)}{0.5(Xtang)}
\uput[-90](X){X} 
\psline[linewidth=2pt,linecolor=red] (0,0) (Xtang)
\psdots[dotstyle=*,linecolor=blue](Xtang)
\psdots[dotstyle=*,linecolor=blue](X) \uput[-90](Xtang){Xtang}
\end{psgraph}
\\ \hline  
\BS{curvepnode}\AC{2}\AC{sin(t)|sin(2*t)}\AC{{\red A}}
&  
\BS{curvepnode}\AC{4}\AC{sin(t)|sin(2*t)}\AC{{\red X} }
\\
\BS{psxline}\AC{<->}(X)\AC{-({\red Atang})}\AC{({\red Atang})}
&
\BS{psxline}\AC{<->}(A)\AC{-0.5({\red Xtang})}\AC{1.5({\red Xtang})}
\\ \hline
\end{tabular} 


%--------------------------------------

\SbSbSSCT{N\oe uds sur une courbe paramétrique avec curvepnodes  }{Nodes on a parametric curve with curvepnodes }

\begin{tabular}{|c|c|} \hline 
\multicolumn{2}{|c|}{ \BSS{curvepnodes}[plotpoints=100]\AC{1}\AC{5}\AC{sin(t)|sin(2*t)}\AC{A} \BSI{curvepnodes}{pst-node}  } \\ \hline 
\begin{psgraph}[axesstyle=none,xticksize=-1.2 1.2 ,yticksize=-1.2 1.2 , subticks=0](0,0)(-1.2,-1.2)(1.2,1.2){4cm}{4cm}
\psset{algebraic}
\parametricplot[plotpoints=200,linecolor=red]{1}{5}{sin(t)|sin(2*t)}
\curvepnodes[plotpoints=100]{1}{5}{sin(t)|sin(2*t)}{A}
\cnodeput[fillstyle=solid](A20){A}{20}
\cnodeput[fillstyle=solid](A80){B}{80}
\psline[linestyle=dotted] (A) (B)
\end{psgraph}
&  
\begin{psgraph}[axesstyle=none,xticksize=-1.5 1.5 ,yticksize=-1.5 1.5 , subticks=0](0,0)(-1.5,-1.5)(1.5,1.5){4cm}{4cm}
\psset{algebraic}
\parametricplot[plotpoints=200,linecolor=red]{1}{5}{sin(t)|sin(2*t)}
\curvepnodes[plotpoints=100]{1}{5}{sin(t)|sin(2*t)}{A}
\psline[linewidth=2pt] (A20) (A80)
\end{psgraph}
\\ \hline  
\BS{cnodeput}(A20)\AC{A}\AC{20}

&  
\BS{psline}[linewidth=2pt] (A20) (A80)
\\
\BS{cnodeput}(A80)\AC{B}\AC{80}
&

\\ \hline 
\end{tabular} 

%--------------------------------------

\SbSSCT{Lignes relatives }{ Relative line }

\SbSbSSCT{Lignes relatives avec psRelNode  }{Relative line width psRelNode }

\begin{tabular}{|l|l|} \hline  
\begin{psgraph}[axesstyle=none,xticksize=0 5,yticksize=0 6,subticks=0](0,0)(0,0)(6,5){3cm}{2.5cm} 
\psline[linewidth=4pt](1,1)(3,2)
\psRelLine[linecolor=red](1,1)(3,2){2}{A}
\psdots[dotstyle=*,linecolor=red](A)  \nput{45}{A}{A}%
\psRelLine[linecolor=green,angle=30](1,1)(3,2){2}{B}
\psdots[dotstyle=*,linecolor=green](B)  \nput{0}{B}{B}%
\psdots[dotstyle=*,](1,1)
\psdots[dotstyle=*,](3,2)
\end{psgraph}
&  
\begin{psgraph}[axesstyle=none,xticksize=0 5,yticksize=0 6,subticks=0](0,0)(0,0)(6,5){3cm}{2.5cm} 
%\psline[linewidth=3pt](1,1)(3,2)
\pnode(1,1){B}\pnode(3,2){C}
\psRelNode[linecolor=red,angle=30](B)(C){2}{A}
\psdots[dotstyle=*,linecolor=red](A)  \nput{-45}{A}{A}%
\psdots[dotstyle=*,linecolor=blue](B) \nput{135}{B}{B} 
\psdots[dotstyle=*,linecolor=blue](C) \nput{45}{C}{C}%
\psline[linewidth=2pt](B)(C)
\psline[linewidth=2pt,linecolor=red](B)(A)
\end{psgraph}
\\ \hline  
\BSS{psRelLine}(1,1)(3,2)\AC{2}\AC{{\red A}}
\BSI{psRelLine}{pst-node}   \BSI{psRelLine}{pstricks-add} 
&  

\\
\BS{}psRelLine[{\red angle=30}](0,0)(2,1)\AC{2}\AC{{\red B}} 
&
\BSS{psRelNode}[linecolor=red,angle=30](B)(C)\AC{2}\AC{{\red A}}   \BSI{psRelNode}{pst-node}  \BSI{psRelNode}{pstricks-add} 
\\ \hline 
\end{tabular} 

\bigskip

\paragraph{Paramètre trueAngle} :
 
\begin{tabular}{|l|}  \hline  
\begin{psgraph}[axesstyle=none,xticksize=0 5,yticksize=0 4,subticks=0](0,0)(0,0)(4,6){6cm}{3cm} 
\psline[linewidth=3pt,linecolor=blue](1,1)(3,1)
\psRelLine[linecolor=red,angle=45](1,1)(3,1){1}{A}
\psRelLine[linecolor=red,angle=45,trueAngle](1,1)(3,1){1}{B}
\psdots[dotstyle=*,linecolor=red](A)   \nput{90}{A}{A}%
\psdots[dotstyle=*,linecolor=red](B)   \nput{90}{B}{B}%
\end{psgraph} 
\\  \hline  
\BS{psRelLine}[angle=45](1,1)(3,1)\AC{1}\AC{{\red A}} 
\\
\BS{psRelLine}[angle=45,\RDD{trueAngle}](1,1)(3,1)\AC{1}\AC{{\red B}} 
\\  \hline 
\end{tabular} 

%--------------------------------------

 
\SbSbSSCT{Lignes relatives avec psRelLineVar }{Relative line width psRelLineVar }

\begin{tabular}{|c|c|} \hline  
\begin{psgraph}[axesstyle=none,xticksize=0 6,yticksize=0 6,subticks=0](0,0)(6,6){3cm}{3cm}
\pnodes(4,1){A}(2,3){B} (1,2){C}
\uput[-135](B){B}
\uput[-45](A){A}
\psdots[dotstyle=*,linecolor=blue](B)
\psdots[dotstyle=*,linecolor=blue](A)
\psRelLineVar[linecolor=red](B)(A)(1;90){X}
\psarc[linestyle=dashed](B){2.828}{-60}{60}
\uput[45](X){X}
\psdots[dotstyle=*,linecolor=red](X) 
\psline(A)(B)
\end{psgraph}
&
\begin{psgraph}[axesstyle=none,xticksize=0 6,yticksize=0 6,subticks=0](0,0)(6,6){3cm}{3cm}
\pnodes(4,1){A}(2,3){B} (1,2){C}
\uput[-135](B){B}
\uput[-45](A){A}
\psdots[dotstyle=*,linecolor=blue](B)
\psdots[dotstyle=*,linecolor=blue](A)
\psRelLineVar[linecolor=red](B)(A)(0.5;135){Y} 
\uput[45](Y){Y}
\psarc[linestyle=dashed](B){1.414}{-70}{120}
\psdots[dotstyle=*,linecolor=red](Y)
\psline(A)(B)
\end{psgraph}
\\ \hline 
\BSS{psRelLineVar}(B)(A)({\red 1;90})\AC{X}  \BSI{psRelLineVar}{pst-node} &
\BSS{psRelLineVar}(B)(A)({\red 0.5;135})\AC{Y} 
\\ \hline 
\end{tabular} 
 
%--------------------------------------

 
\SbSbSSCT{Ligne par une série de points  avec psnline  }{ Line from seval points width psnline  }

\begin{tabular}{|l|l|} \hline  
\begin{psgraph}[axesstyle=none,xticksize=0 3cm,yticksize=0 5cm,subticks=0](0,0)(5,3){5cm}{3cm}
\pnodes{A}(1,1)(3,0.5)(4,2)(2,3)(1,2)

\psdots[dotstyle=*,linecolor=blue](A0) \uput[-135](A0){A0}
\psdots[dotstyle=*,linecolor=blue](A1) \uput[-45](A1){A1}
\psdots[dotstyle=*,linecolor=blue](A2) \uput[-45](A2){A2}
\psdots[dotstyle=*,linecolor=blue](A3) \uput[-90](A3){A3}
\psdots[dotstyle=*,linecolor=blue](A4) \uput[-45](A4){A4}
\psnline[linecolor=red](0,3){A}
\end{psgraph}
 &
\begin{psgraph}[axesstyle=none,xticksize=0 3cm,yticksize=0 5cm,subticks=0](0,0)(5,3){5cm}{3cm}
\pnodes{A}(1,1)(3,0.5)(4,2)(2,3)(1,2)
\psdots[dotstyle=*,linecolor=blue](A0) \uput[-135](A0){A0}
\psdots[dotstyle=*,linecolor=blue](A1) \uput[-45](A1){A1}
\psdots[dotstyle=*,linecolor=blue](A2) \uput[-45](A2){A2}
\psdots[dotstyle=*,linecolor=blue](A3) \uput[-90](A3){A3}
\psdots[dotstyle=*,linecolor=blue](A4) \uput[-45](A4){A4}
\psnline[linecolor=red](2,2){A}
\end{psgraph}
\\ \hline  
 \BS{pnodes}\AC{A}(1,1)(3,0.5)(4,2)(2,3)(1,1) 
 &
  \BS{pnodes}\AC{A}(1,1)(3,0.5)(4,2)(2,3)(1,1) 
 \\
 \BSS{psnline}({ \red 0,3})\AC{A} \BSI{psnline}{pst-node} 
 &
  \BSS{psnline}({ \red 2,2})\AC{A}
\\ \hline 
\end{tabular} 

%--------------------------------------

 
\SbSbSSCT{Courbe par une série de points  avec psncurve }{ Curve from seval points width psncurve }


\begin{tabular}{|l|l|} \hline  
\begin{psgraph}[axesstyle=none,xticksize=0 3cm,yticksize=0 5cm,subticks=0](0,0)(5,3){5cm}{3cm}
\pnodes{A}(1,1)(3,0.5)(4,2)(2,3)(1,2)

\psdots[dotstyle=*,linecolor=blue](A0) \uput[-135](A0){A0}
\psdots[dotstyle=*,linecolor=blue](A1) \uput[-45](A1){A1}
\psdots[dotstyle=*,linecolor=blue](A2) \uput[-45](A2){A2}
\psdots[dotstyle=*,linecolor=blue](A3) \uput[-90](A3){A3}
\psdots[dotstyle=*,linecolor=blue](A4) \uput[-45](A4){A4}
\psncurve[linecolor=red](0,3){A}
\end{psgraph}
 &
\begin{psgraph}[axesstyle=none,xticksize=0 3cm,yticksize=0 5cm,subticks=0](0,0)(5,3){5cm}{3cm}
\pnodes{A}(1,1)(3,0.5)(4,2)(2,3)(1,2)
\psdots[dotstyle=*,linecolor=blue](A0) \uput[-135](A0){A0}
\psdots[dotstyle=*,linecolor=blue](A1) \uput[-45](A1){A1}
\psdots[dotstyle=*,linecolor=blue](A2) \uput[-45](A2){A2}
\psdots[dotstyle=*,linecolor=blue](A3) \uput[-90](A3){A3}
\psdots[dotstyle=*,linecolor=blue](A4) \uput[-45](A4){A4}
\psncurve[linecolor=red](2,2){A}
\end{psgraph}
\\ \hline  
 \BS{pnodes}\AC{A}(1,1)(3,0.5)(4,2)(2,3)(1,1) 
 &
  \BS{pnodes}\AC{A}(1,1)(3,0.5)(4,2)(2,3)(1,1) 
 \\
 \BSS{psncurve}({\red 0,3})\AC{A} \BSI{psncurve}{pst-node} 
 &
  \BSS{psncurve}({\red 2,2})\AC{A}
\\ \hline 
\end{tabular}

%--------------------------------------
 
\SbSbSSCT{ligne par pas succesifs  avec psrline }{ Line by succesive step  width psrline }

\begin{tabular}{|c|} \hline  
\begin{psgraph}[axesstyle=none,xticksize=0 2cm,yticksize=0 5cm,subticks=0](0,0)(5,2){5cm}{2cm}
 \psrline[linewidth=2pt,linecolor=red](0,.5)(1,1)(1,-1)(2,1)
\psline[linestyle=dashed,arrowscale=2]{->}(0,.5)(1,.5)(1,1.5)\psline[linestyle=dashed,arrowscale=2]{->}(1,1.5)(2,1.5)(2,0.5)
\psline[linestyle=dashed,arrowscale=2]{->}(2,.5)(4,.5)(4,1.5)
 \end{psgraph}
\\ \hline  
\BSS{psrline}(0,0.5)(1,1)(1,-1)(2,1) \BSI{psrline}{pst-node} 

\\ \hline 
\end{tabular} 

%--------------------------------------
 
\SbSbSSCT{Lignes par rapport à un point avec psxline }{ Lines relative at a point  width psxline }

\begin{tabular}{|c|c|} \hline  
\begin{psgraph}[axesstyle=none,xticksize=0 6,yticksize=0 6,subticks=0](0,0)(6,6){4cm}{4cm}
\pnodes(3,1){A}(2,3){B} (1,2){C}
\uput[-45](B){B}
 \psxline[linecolor=red]{<->}(B){1,2}{3,1}
\psdots[dotstyle=*,dotscale=2](B)
 \psline (B)(3,5)
 \psline (B)(5,4) 
\end{psgraph}
&  
\begin{psgraph}[axesstyle=none,xticksize=0 6,yticksize=0 6,subticks=0](0,0)(6,6){4cm}{4cm}
\pnodes(3,1){A}(2,3){B} (1,2){C}
\uput[-45](A){A}
\uput[-45](B){B}
\uput[-45](C){C}
 \psxline[linecolor=red]{<->}(B){A}{C}
\psdots[dotstyle=*,linecolor=blue](A)
\psdots[dotstyle=*,linecolor=blue](B)
\psdots[dotstyle=*,linecolor=blue](C)
 \psline(0,0)(C)
  \psline (0,0)(A)
 \psline (B)(3,5)
 \psline (B)(5,4) 
\end{psgraph}
\\ \hline  
\BSS{psxline}\AC{<->}(B)\AC{1,2}\AC{3,1} \BSI{psxline}{pst-node} 
&  
\BSS{psxline}\AC{<->}(B)\AC{A}\AC{C}
\\ \hline 
\end{tabular}

\psset{unit=1}

 %-----------------------------------------------------------
\SbSSCT{Lignes parallèles et leur noeud final}{Parallel lines and their endpoint} 

\TFRGB{ Syntaxe : \\
\BS{}psParallelLine(Point 1)(point 2 )(point 3)\AC{longueur}\AC{nom extrémité} }{Syntax : 
\BS{}psParallelLine(Point 1)(point 2 )(point 3)\AC{lengh}\AC{end name}}

\bigskip


\begin{tabular}{|l|} \hline  
\begin{psgraph}[axesstyle=none,xticksize=0 5,yticksize=0 6,subticks=0](0,0)(0,0)(6,5){3cm}{2.5cm} 
\psline[linewidth=3pt,linecolor=blue](2,1)(4,2)
\psParallelLine[linecolor=red](2,1)(4,2)(1,2){2}{A}
\psdots[dotstyle=*,linecolor=red](A)   \nput{45}{A}{A}%
\psParallelLine[linecolor=green](2,1)(4,2)(1,3){1}{B} 
\psdots[dotstyle=*,linecolor=green](B)   \nput{45}{B}{B}%
\psParallelLine[linecolor=cyan](2,1)(4,2)(1,4){.5}{C}
\psdots[dotstyle=*,linecolor=cyan](C)   \nput{45}{C}{C}%
\end{psgraph} 
\\ \hline  
\BSS{psParallelLine}(2,1)(4,2))(1,2)\AC{{\red 2}}\AC{{\red A}}  \BSI{psParallelLine}{pstricks-add} 
\\
\BS{psParallelLine}(2,1)(4,2)(1,3)\AC{{\green 1}}\AC{{\green B}} 
\\
\BS{psParallelLine}(2,1)(4,2)(1,4)\AC{{\cyan 0.5}}\AC{{\cyan C}}
\\ \hline 
\end{tabular} 

%---------------------------

\SbSSCT{Lignes perpendiculaires une droite}{Perpendiculars to a lines}

\begin{tabular}{|l|} \hline  
\begin{psgraph}[axesstyle=none,xticksize=0 5,yticksize=0 6,subticks=0](0,0)(0,0)(6,5){3cm}{2.5cm} 
\psline(5,5)(3,0)
\psPline[linecolor=red,arrowscale=2]{->}(0,3)(5,5)(3,0)
\psPline[linecolor=green,arrowscale=2](1,4)(5,5)(3,0)
\end{psgraph} 
\\ \hline  
\BS{}psline(5,5)(3,0)
\\
\BSS{psPline}[linecolor=red]\AC{->}(0,3)(5,5)(3,0)
\\ 
\BS{}psPline[linecolor=green](1,4)(5,5)(3,0)
\\ \hline 
\end{tabular} 

%---------------------------------------

\subsection{Vecteur normal }

\begin{tabular}{|c|c|} \hline  
\begin{psgraph}[axesstyle=none,xticksize=0 3cm,yticksize=-2 4cm,subticks=0](0,0)(-2,0)(4,3){6cm}{3cm}
\pnode(2,1){B}
\normalvec(B){A}
\psline[linecolor=red]{->}(0,0) (A)
\psline{->}(0,0) (B)
\nput{90}{B}{B} 
\nput{90}{A}{A}%
\end{psgraph}
&
\begin{psgraph}[axesstyle=none,xticksize=-1.5 1.5 ,yticksize=-1.5 1.5 , subticks=0](0,0)(-1.5,-1.5)(1.5,1.5){2.5cm}{2.5cm}
\psset{algebraic}
\parametricplot[plotpoints=200]{0}{6.28}{sin(t)|sin(2*t)}
\curvepnode{2}{sin(t)|sin(2*t)}{A}
\psxline[linewidth=2pt,linestyle=dashed]{<->}(A){-(Atang)}{(Atang)}
\uput[-45](A){A}
\normalvec(Atang){B} 
\psdots[dotstyle=*,linecolor=blue](A) 
\psxline[linewidth=2pt,linecolor=red]{->}(A){}{-2(B)}
\end{psgraph}
\\ \hline  
\BSS{normalvec}(B)\AC{A} \BSI{normalvec}{pst-node} 
&
\BSS{normalvec}(Atang)\AC{B}
\\
&
\BS{psxline}\AC{->}(A)\AC{} \AC{-2(B)}
\\ \hline 
\end{tabular}

%-----------------------------------------

\SbSSCT{Tangentes}{Tangents}

\SbSbSSCT{Tangentes à un cercle par rapport à un point}{Tangent lines of a circle}

\begin{tabular}{|l|} \hline 
\begin{psgraph}[axesstyle=none,xticksize=-2 2,yticksize=-4 2,subticks=0](0,0)(-4,-2)(5,2){7cm}{3cm}  
\pscircle(03, 0){1}
\psdots[dotstyle=*,linecolor=blue,dotscale=2](-3,0)
\psCircleTangents(-3, 0)(3, 0){1}
\pcline[nodesep=-1cm,linestyle=dashed](-3, 0)(CircleT1)
\pcline[nodesep=-1cm,linestyle=dashed](-3, 0)(CircleT2)
\psdots[dotstyle=*,linecolor=red,dotscale=2](CircleT1)
\psdots[dotstyle=*,linecolor=red,dotscale=2](CircleT2)
\nput{-90}{CircleT1}{CircleT1}%
\nput{90}{CircleT2}{CircleT2}%
\psdots[dotstyle=*](3,0)
\psrline[linestyle=dotted]{->}(3,0)(1;45)
\end{psgraph}
\\ \hline  
\BS{pscircle}(3,0)\AC{1} \\
\BSS{psCircleTangents}(-3,0)(3,0)\AC{1}  \BSI{psCircleTangents}{pstricks-add} \\
\BS{pcline}[nodesep=-1cm,linecolor=blue](-3,0)({\red CircleT1}) \\
\BS{nput}\AC{-90}\AC{{\red CircleT1}}\AC{CircleT1}%
\\ \hline 
\end{tabular} 
 

%------------------------------------------------------------------------------
\SbSbSSCT{Tangentes à une ellipse par rapport à un point}{Tangent lines of an ellipse}

\begin{tabular}{|l|} \hline  
 
\begin{psgraph}[axesstyle=none,xticksize=-2 2,yticksize=-5 6,subticks=0](0,0)(-5,-2)(6,2){7cm}{3cm}  
\psellipse(3, 0)(2, 1)
\psdots[dotstyle=*,linecolor=blue,dotscale=2](-3,0)
\psEllipseTangents(3, 0)(2, 1)(-3,0)
\pcline[nodesep=-1cm,linestyle=dashed](-3, 0)(EllipseT1)
\pcline[nodesep=-1cm,linestyle=dashed](-3, 0)(EllipseT2)
\psdots[dotstyle=*,linecolor=red,dotscale=2](EllipseT1)
\psdots[dotstyle=*,linecolor=red,dotscale=2](EllipseT2)
\nput{-90}{EllipseT1}{EllipseT1}%
\nput{90}{EllipseT2}{EllipseT2}%
\psline[linestyle=dotted]{<->}(1,0)(3,0)
\psline[linestyle=dotted]{<->}(3,0)(3,1)
\end{psgraph}
\\ \hline  
\BS{}psellipse(3,0)(2,1) \\
\BSS{psEllipseTangents}(3,0)(2,1)(-3,0) \BSI{psEllipseTangents}{pstricks-add} \\
\BS{}pcline[nodesep=-1cm](-3,0)({\red EllipseT1}) \\
\BS{}nput\AC{90}\AC{{\red EllipseT1}}\AC{EllipseT1}%
\\ \hline 
\end{tabular} 

 
 %-----------------------------------------------------------\BSS{psCircleTangents}(-3,0)(3,0)\AC{1}-------------------------------------
\SbSbSSCT{Tangentes à deux cercles}{Tangent lines of circles}
 \psset{unit=.5cm}


\begin{tabular}{|l|} \hline  
\begin{psgraph}[axesstyle=none,xticksize=-3 3,yticksize=-4 14,subticks=0](0,0)(-4,-3)(14,3){9cm}{3cm}  
\pscircle(-1, 0){2}
\pscircle(5, 0){1}
\psCircleTangents(-1, 0){2}(5, 0){1}
\pcline[nodesep=-1cm,linestyle=dashed](CircleTO1)(CircleTO2)
\pcline[nodesep=-1cm,linestyle=dashed](CircleTO3)(CircleTO4)
\pcline[nodesep=-1cm,linestyle=dashed](CircleTO3)(CircleTC2)
\pcline[nodesep=-1cm,linestyle=dashed](CircleTO2)(CircleTC2)
\psdots[dotstyle=*,linecolor=red,dotscale=2](CircleTC2)
\psdots[dotstyle=*,linecolor=red,dotscale=2](CircleTO1)
\psdots[dotstyle=*,linecolor=red,dotscale=2](CircleTO2)
\psdots[dotstyle=*,linecolor=red,dotscale=2](CircleTO3) 
\psdots[dotstyle=*,linecolor=red,dotscale=2](CircleTO4)
\nput{90}{CircleTC2}{CircleTC2}% 
\nput{45}{CircleTO1}{CircleTO1}% 
\nput{45}{CircleTO2}{CircleTO2}% 
\nput{-45}{CircleTO3}{CircleTO3}% 
\nput{-45}{CircleTO4}{CircleTO4}% 
\end{psgraph}
\\ \hline  
\BSS{psCircleTangents}(-1, 0)\AC{2}(5,0)\AC{1} \\

\BS{}psdots[dotstyle=*,linecolor=red,dotscale=2]({\red CircleTC2})

\\ \hline  
\end{tabular} 
 

 
\bigskip
 \psset{xunit=1cm}


\begin{tabular}{|l|} \hline  
\begin{psgraph}[axesstyle=none,xticksize=-3 3,yticksize=-3 6,subticks=0](0,0)(-3,-3)(6,3){9cm}{3cm} 
\pscircle(-1, 0){2}
\pscircle(5, 0){1}
\psCircleTangents(-1, 0){2}(5, 0){1}
\pcline[nodesep=-1cm,linestyle=dashed](CircleTI1)(CircleTI2)
\pcline[nodesep=-1cm,linestyle=dashed](CircleTI3)(CircleTI4)
\psdots[dotstyle=*,linecolor=red,dotscale=2](CircleTC1)
\psdots[dotstyle=*,linecolor=red,dotscale=2](CircleTI1)
\psdots[dotstyle=*,linecolor=red,dotscale=2](CircleTI2)
\psdots[dotstyle=*,linecolor=red,dotscale=2](CircleTI3)
\psdots[dotstyle=*,linecolor=red,dotscale=2](CircleTI4)
\nput{90}{CircleTC1}{CircleTC1}% 
\nput{45}{CircleTI1}{CircleTI1}% 
\nput{-135}{CircleTI2}{CircleTI2}% 
\nput{-45}{CircleTI3}{CircleTI3}% 
\nput{135}{CircleTI4}{CircleTI4}% 
\end{psgraph}
\\ \hline  
\BS{}psdots[dotstyle=*,linecolor=red,dotscale=2]({\red CircleTC1}) \\

\BS{}nput\AC{90}\AC{{\red CircleTC1}}\AC{CircleTC1}% 
\\ \hline 
\end{tabular}  
 

\psset{unit=1cm,fillcolor=white,fillstyle=none,linecolor=blue}

%------------------------------------------------------------
\subsection{Intersections}

\SbSbSSCT{Point d'intersection avec psIntersectionPoint}{Intersection point of two lines width psIntersectionPoint}
\TFRGB{Syntaxe : \\
\BS{}psIntersectionPoint(point 1)(point 2)(point 3)(point 4)\AC{nom}}{Syntax : \BS{}psIntersectionPoint(point 1)(point 2)(point 3)(point 4)\AC{name}}

\bigskip
\psset{unit=1cm}


\begin{tabular}{|c|} \hline  
\begin{psgraph}[axesstyle=none,xticksize=-1 4,yticksize=0 5,subticks=0](0,0)(0,-1)(5,4){3cm}{3cm} 
\psdots[dotstyle=*,dotsize=6pt](1,3)
\psdots[dotstyle=*,dotsize=6pt](4,0)
\pnode(1,0){B} \pscircle*(B){3pt} \nput{135}{B}{B}%
\pnode(4,3){C} \pscircle*(C){3pt} \nput{45}{C}{C}%
\psline[linecolor=blue](B)(C)
\psline[linecolor=blue](1,3)(4,0)
\psIntersectionPoint(B)(C)(1,3)(4,0){A}
\psdots[dotstyle=*,linecolor=red,dotscale=2](A)   \nput{90}{A}{A}%
\end{psgraph}
\\ \hline  
\BSS{psIntersectionPoint}(B)(C)(1,3)(4,0)\AC{{\red A}}  \BSI{psIntersectionPoint}{pstricks-add} 
\\ \hline 
\end{tabular} 
  
%----------------------------------------

\SbSbSSCT{Points d'intersection avec polyIntersections} {Intersection points with polyIntersections}
%\subsubsection{polyIntersections}

\begin{tabular}{|c|c|} \hline  
\begin{psgraph}[axesstyle=none,xticksize=0 3cm,yticksize=0 4cm,subticks=0](0,0)(0,0)(4,3){4cm}{3cm}
\psline(0.5,.5)(3,.5)(2.5,2)(1,2.5)(.5,1)
\pnodes(1,1.5){A}(1.5,1){B} 
\polyIntersections{X}{Y}(A)(B)(0.5,.5)(3,.5)(2.5,2)(1,2.5)(.5,1)
\psdots(A)(B)\psline(X)(Y)
\psdots[dotstyle=*,linecolor=red](X)
\psdots[dotstyle=*,linecolor=red](Y)
\uput[0](X){X}\uput[-180](Y){Y}
\uput[0](A){A}\uput[-180](B){B}
\end{psgraph}
&  
\begin{psgraph}[axesstyle=none,xticksize=0 3cm,yticksize=0 4cm,subticks=0](0,0)(0,0)(4,3){4cm}{3cm}
\pnodes{P}(0.5,.5)(3,.5)(2.5,2)(1,2.5)(.5,1)
\pnode(1,1.5){A}\pnode(1.5,1){B}
\polyIntersections{X}{Y}(A)(B){P}{4}
\psnline(0,4){P}
\psdots(A)(B)\psline(X)(Y)
\psdots[dotstyle=*,linecolor=red](X)
\psdots[dotstyle=*,linecolor=red](Y)
\uput[-45](X){X}\uput[-180](Y){Y}
\uput[0](A){A}\uput[-180](B){B}
\end{psgraph}
\\ \hline  
\BSS{polyIntersections}\AC{X}\AC{Y}(A)(B) \BSI{polyIntersections}{pst-node} 
&  
\BS{pnodes}\AC{P}(0.5,.5)(3,.5)(2.5,2)(1,2.5)(.5,1)
\\
(0.5,0.5)(3,0.5)(2.5,2)(1,2.5)(0.5,1)
&
\BSS{polyIntersections}\AC{X}\AC{Y}(A)(B)\AC{P}\AC{4}
\\ \hline 
\end{tabular} 


\bigskip

\begin{tabular}{|l|} \hline  
\begin{psgraph}[axesstyle=none,xticksize=-1.5 1.5cm,yticksize=0 8cm,subticks=0](0,0)(0,-1.5)(13,1.5){12cm}{3cm}
\psset{algebraic}
\psplot[plotpoints=200]{0}{13}{sin(x)}
\fnpnodes[plotpoints=100]{0}{13}{sin(x)}{A}
\pnode(4,0.5){C} \pnode(6,.5){B}
\polyIntersections{X}{Y}(C)(B){A}{60}
\psline(X)(Y)
\psdots[dotstyle=*](B)
\psdots[dotstyle=*](C) 
\psdots[dotstyle=*,linecolor=red](X)
\psdots[dotstyle=*,linecolor=red](Y)
\psdots[dotstyle=*,linecolor=green](A60)
\uput[90](C){C}\uput[90](B){B}
\uput[0](X){X}\uput[180](Y){Y}\uput[90](A60){A60}
\end{psgraph}
\\ \hline  
\BS{fnpnodes}[plotpoints=100]\AC{0}\AC{13}\AC{sin(x)}\AC{A} \\
\BSS{polyIntersections}\AC{X}\AC{Y}(C)(B)\AC{A}\AC{60}
\\ \hline 
\end{tabular} 

\bigskip

\begin{tabular}{|l|} \hline  
\begin{psgraph}[axesstyle=none,xticksize=-1.2 1.2 ,yticksize=-1.5 1.5 , subticks=0](0,0)(-1.5,-1.2)(1.5,1.2){8cm}{3cm}
\psset{algebraic}
\parametricplot[plotpoints=200]{1}{5}{sin(t)|sin(2*t)}
\curvepnodes[plotpoints=100]{1}{5}{sin(t)|sin(2*t)}{A}
\pnode(0,0.5){C} \pnode(.5,.5){B}
\polyIntersections{X}{Y}(C)(B){A}{75}
\psline(X)(Y)
\psdots[dotstyle=*](B)
\psdots[dotstyle=*](C) 
\psdots[dotstyle=*,linecolor=red](X)
\psdots[dotstyle=*,linecolor=red](Y)
\psdots[dotstyle=*,linecolor=green](A75)
\uput[90](C){C}\uput[90](B){B}
\uput[0](X){X}\uput[180](Y){Y}\uput[90](A75){A75}
\end{psgraph}
\\ \hline 
\BS{curvepnodes}[plotpoints=100]\AC{1}\AC{5}\AC{sin(t)|sin(2*t)}\AC{A} \\
\BSS{polyIntersections}\AC{X}\AC{Y}(C)(B)\AC{A}\AC{75}
\\ \hline 
\end{tabular} 

 
 

\newpage
%-----------------------------------------------------------------------
\SbSSCT{Les 9 positions d'une figure par \BS{}psDefPSPNodes}{The 9 positions with \BS{}psDefPSPNodes}
\bigskip

\begin{tabular}{l}  
\begin{pspicture}(6,4)
\psaxes[xticksize=4,yticksize=6,axesstyle=none](6,4)
\psDefPSPNodes
\psdots[linecolor=red,dotscale=2](PSPbl)(PSPbc)(PSPbr)(PSPcl)(PSPcc)(PSPcr)(PSPtl)(PSPtc)(PSPtr)
 \uput[45](PSPbl){PSPbl} \uput[90](PSPbc){PSPbc}
\uput[90](PSPbr){PSPbr} \uput[45](PSPcl){PSPcl}
\uput[90](PSPcc){PSPcc} \uput[90](PSPcr){PSPcr}
\uput[45](PSPtl){PSPtl} \uput[90](PSPtc){PSPtc}
 \uput[90](PSPtr){PSPtr}
 \end{pspicture}
\\ 
\\ 
\\ \hline 
\BS{}begin{pspicture}(6,4)\\
 
\BSS{psDefPSPNodes}  \BSI{psDefPSPNodes}{pstricks-add} \\

\BS{}psdots({\red PSPbl})\\

\BS{}uput[45](PSPbl)\AC{{\red PSPbl}}
\\ 
\end{tabular}  

%--------------------------------

\SbSSCT{N\oe uds sur du texte avec  \BS{}psDefBoxNodes}{Nodes on text with \BS{}psDefBoxNodes}
\bigskip

\begin{tabular}{|l|} \hline  
 \psscalebox{8}{\psDefBoxNodes{nom}{\color{red!20} abcdefghij}}%
\shorthandoff{ :}
\uput[90](nom:tl){tl} \qdisk(nom:tl){3pt}
\uput[90](nom:tC){tC} \qdisk(nom:tC){3pt}
\uput[90](nom:tr){tr} \qdisk(nom:tr){3pt}
\uput[90](nom:Cl){Cl} \qdisk(nom:Cl){3pt}
\uput[90](nom:C){C} \qdisk(nom:C){3pt}
\uput[90](nom:Cr){Cr} \qdisk(nom:Cr){3pt}
\uput[90](nom:Bl){Bl} \qdisk(nom:Bl){3pt}
\uput[90](nom:BC){BC} \qdisk(nom:BC){3pt}
\uput[90](nom:Br){Br} \qdisk(nom:Br){3pt}
\uput[90](nom:bl){bl} \qdisk(nom:bl){3pt}
\uput[90](nom:bC){bC} \qdisk(nom:bC){3pt}
\uput[90](nom:br){br} \qdisk(nom:br){3pt}
\shorthandon{ :}
\\ \hline 
\\ 
 \BS{psscalebox}\AC{15}\AC{ \BSS{psDefBoxNodes}\AC{{\red nom}} \AC{ \BS{color}\AC{red!20} abcdefghij}}\\
 \BSS{shorthandoff}\AC{:} \footnotemark[1]  \\
\BS{uput}[90]({\red nom\xx tl})\AC{tl} \hspace{1cm} \BS{qdisk}({\red nom\xx tl})\AC{3pt} \\ \hspace{1cm} $\vdots $\\
 \BSS{shorthandon}\AC{:}
\\ \hline 
\end{tabular} 
 \footnotetext[1]{désactivation et ré-activation de \og : \fg  conflit entre ce module et Babel en français }


\subsection{ArrowNotch}

\begin{tabular}{|l|l|} \hline 
\multicolumn{2}{|c|}{ \BS{curvepnodes}[plotpoints=100]\AC{1}\AC{1.1}\AC{sin(t)|sin(2*t)}\AC{{\red A}} } 
\\ \hline
\begin{psgraph}[axesstyle=none,xticksize=0.5 0.6 ,yticksize=0.5 0.6 , subticks=0](0.8,0.8)(.9,.95){5cm}{5cm}
 \psset{algebraic}
\parametricplot[plotpoints=200]{1}{1.1}{sin(t)|sin(2*t)}
\curvepnodes[plotpoints=100]{1}{1.3}{sin(t)|sin(2*t)}{A} 
\ArrowNotch[arrowscale=10]{A}{0}{>}{X}
\psline[linecolor=red,arrowscale=5]{-D>}(X)(A0)
\psdots[dotstyle=*](X) \uput[-135](X){X}
\psdots[dotstyle=*](A0) \uput[-135](A0){A0}
\end{psgraph}
&  
\begin{psgraph}[axesstyle=none,xticksize=0.5 0.6 ,yticksize=0.5 0.6 , subticks=0](0.8,0.8)(.9,.95){5cm}{5cm}
\psset{algebraic}
\parametricplot[plotpoints=200]{1}{1.1}{sin(t)|sin(2*t)}
\curvepnodes[plotpoints=100]{1}{1.5}{sin(t)|sin(2*t)}{A} 
\ArrowNotch[arrowscale=10]{A}{0}{<}{V}
\psline[linecolor=red,arrowscale=5]{-D>}(V)(A0)
\psdots[dotstyle=*](V) \uput[-135](V){V}
\psdots[dotstyle=*](A0) \uput[-135](A0){A0}
\end{psgraph}
\\ \hline  
\BSS{ArrowNotch}[arrowscale=10]\AC{{\red A}}\AC{0}\AC{>}\AC{X} \BSI{ArrowNotch}{pst-node} 
&  
\BSS{ArrowNotch}[arrowscale=10]\AC{{\red A}}\AC{0}\AC{<}\AC{V}
\\
\BS{psline}[arrowscale=5]\AC{-D>}(X)(A0)
&
\BS{psline}[arrowscale=5]\AC{-D>}(V)(A0)

\\ \hline
\begin{psgraph}[axesstyle=none,xticksize=0.5 0.6 ,yticksize=0.5 0.6 , subticks=0](0.8,0.8)(.9,.95){5cm}{5cm}
 \psset{algebraic}
\parametricplot[plotpoints=200]{1}{1.1}{sin(t)|sin(2*t)}
\curvepnodes[plotpoints=100]{1}{1.3}{sin(t)|sin(2*t)}{A} 
\ArrowNotch[arrowscale=10]{A}{20}{>}{X}
\psline[linecolor=red,arrowscale=5]{-D>}(X)(A20)
\psdots[dotstyle=*](X) \uput[-135](X){X}
\psdots[dotstyle=*](A20) \uput[-135](A20){A20}
\end{psgraph}
&  
\begin{psgraph}[axesstyle=none,xticksize=0.5 0.6 ,yticksize=0.5 0.6 , subticks=0](0.8,0.8)(.9,.95){5cm}{5cm}
\psset{algebraic}
\parametricplot[plotpoints=200]{1}{1.1}{sin(t)|sin(2*t)}
\curvepnodes[plotpoints=100]{1}{1.3}{sin(t)|sin(2*t)}{A} 
\ArrowNotch[arrowscale=10]{A}{20}{<}{V}
\psline[linecolor=red,arrowscale=5]{-D>}(V)(A20)
\psdots[dotstyle=*](V) \uput[-135](V){V}
\psdots[dotstyle=*](A20) \uput[-135](A20){A20}
\end{psgraph}
\\ \hline  
\BSS{ArrowNotch}[arrowscale=10]\AC{{\red A}}\AC{0}\AC{>}\AC{X}
&  
\BSS{ArrowNotch}[arrowscale=10]\AC{{\red A}}\AC{0}\AC{<}\AC{V}
\\
\BS{psline}[arrowscale=5]\AC{-D>}(X)(A0)
&
\BS{psline}[arrowscale=5]\AC{-D>}(V)(A0)
\\ \hline 
\end{tabular} 



\subsection{Placement d'une étiquette à une distance donnée avec nlput }
\SbSSCT{Placement d'une étiquette à une distance donnée avec nlput }{nlput}

\begin{tabular}{|c|} \hline  
\begin{psgraph}[axesstyle=none,xticksize=0 2cm,yticksize=0 6cm,subticks=0](0,0)(6,2){6cm}{2cm}
\pnode(1,1){B}\pnode(5,1){C}
 \nlput(B)(C){1cm}{Texte}
\psdots[dotstyle=*,linecolor=blue](B) \nput{90}{B}{B} 
\psdots[dotstyle=*,linecolor=blue](C) \nput{90}{C}{C}%
 \end{psgraph}
\\ \hline  
\BSS{nlput}(B)(C)\AC{{\red 1cm}}\AC{Texte} \BSI{nlput}{pst-node} 
\\ \hline 
\end{tabular}  

\bigskip


\begin{tabular}{|c|c|c|} \hline 
\multicolumn{3}{|c|}{  \BSS{nlput}[\RDD{nrot}=\xx U](B)(C)\AC{1cm}\AC{\BS{red} Texte} \RDI{nrot}{pst-node}} \\ \hline  
\begin{pspicture}(5,2)
\pnode(0,0){A}
 \pnode(4,2){B}
 \ncline{A}{B}
 \nlput[nrot=:U](A)(B){1cm}{\red Texte}
 \end{pspicture}
&  
\begin{pspicture}(5,2)
\pnode(0,0){A}
 \pnode(4,2){B}
 \ncline{A}{B}
 \nlput[nrot=:D](A)(B){1cm}{\red Texte}
 \end{pspicture}
&  
\begin{pspicture}(5,2)
\pnode(0,0){A}
 \pnode(4,2){B}
 \ncline{A}{B}
 \nlput[nrot=90](A)(B){1cm}{\red Texte}
 \end{pspicture}
\\ \hline 
 nrot=\xx U  &  nrot=\xx U  & nrot=90
\\ \hline 
\begin{pspicture}(5,2)
\pnode(0,0){A}
 \pnode(4,2){B}
 \ncline{A}{B}
 \nlput[nrot=:L](A)(B){1cm}{\red Texte}
 \end{pspicture}
 &  
 \begin{pspicture}(5,2)
 \pnode(0,0){A}
 \pnode(4,2){B}
  \ncline{A}{B}
  \nlput[nrot=:R](A)(B){1cm}{\red Texte}
  \end{pspicture}
 &  
  \begin{pspicture}(5,2)
  \pnode(0,0){A}
  \pnode(4,2){B}
   \ncline{A}{B}
   \nlput(A)(B){1cm}{\red Texte}
\end{pspicture}
\\ \hline 
 nrot=\xx L  &  nrot=\xx R  & sans paramètre
 \\ \hline 
\end{tabular} 

  

\newpage

%\subsection{Homothétie}
\SSCT{Homothétie}{Homothety}

\psset{unit=0.8cm}
\label{hom}

\begin{tabular}{|c|} \hline
\parbox{10cm}{  
\BS{pscircle}[linestyle=dotted]\AC{1}  
\\
\BSS{psHomothetie}[linecolor=magenta](2,1)\AC{.5}\AC{\BS{pscircle}\AC{1}}  \BSI{psHomothetie}{pstricks-add}
\\
\BSS{psHomothetie}[linecolor=red](2,1)\AC{2}\AC{\BS{pscircle}\AC{1}}
\\
\BSS{psHomothetie}[linecolor=green](2,1)\AC{3}\AC{\BS{pscircle}\AC{1}}  }
\\ 	\hline 
\begin{pspicture}(-8.3,-6.3)(2.3,2.5)
\psgrid[subgriddiv=0](-8,-6)(2,2)
\pscircle[linecolor=blue,linestyle=dotted,linewidth=2pt]{1}
\psline(-8,-4)(2,1)
\psHomothetie[linecolor=magenta](2,1){.5}{\pscircle{1}}
\psHomothetie[linecolor=red](2,1){2}{\pscircle{1}}
\psHomothetie[linecolor=green](2,1){3}{\pscircle{1}}
\psdots[linecolor=red,dotsize=8pt](2,1)
\psline[linewidth=3pt,linecolor=red]{->}(2,1)(0,0)	\psdots(-4,-2)(-2,-1)(0,0)(1,.5)
%\end{psgraph}
\end{pspicture}
	\\ 	\hline 
\end{tabular} 


\bigskip


\begin{tabular}{|c|} \hline
\parbox{10cm}{  
\BS{pscircle}[linestyle=dotted]\AC{1}  
\\
\BSS{psHomothetie}[linecolor=magenta](-2,-1)\AC{.5}\AC{\BS{pscircle}\AC{1}}
\\
\BSS{psHomothetie}[linecolor=red](-2,-1)\AC{2}\AC{\BS{pscircle}\AC{1}}
\\
\BSS{psHomothetie}[linecolor=green](-2,-1)\AC{3}\AC{\BS{pscircle}\AC{1}} }
\\ 	\hline 
\begin{pspicture}(-2.3,-2.3)(8.3,6.5)
\psgrid[subgriddiv=0](-2,-2)(8,6)
\pscircle[linecolor=blue,linestyle=dotted,linewidth=2pt]{1}
\psline(8,4)(-2,-1)
\psHomothetie[linecolor=magenta](-2,-1){.5}{\pscircle{1}}
\psHomothetie[linecolor=red](-2,-1){2}{\pscircle{1}}
\psHomothetie[linecolor=green](-2,-1){3}{\pscircle{1}}
\psdots[linecolor=red,dotsize=8pt](2,1)
\psline[linewidth=3pt,linecolor=red]{->}(-2,-1)(0,0)	\psdots(4,2)(2,1)(0,0)(-1,-.5)
\end{pspicture}
	\\ 	\hline 
\end{tabular}


\begin{tabular}{|c|} \hline 
\parbox{10cm}{ 
\BS{pscircle}[linestyle=dotted]\AC{1}  
\\
\BSS{psHomothetie}[linecolor=magenta](2,1)\AC{-.5}\AC{\BS{pscircle}\AC{1}}
\\
\BSS{psHomothetie}[linecolor=black](2,1)\AC{-1}\AC{\BS{pscircle}\AC{1}}
\\
\BSS{psHomothetie}[linecolor=red](2,1)\AC{-2}\AC{\BS{pscircle}\AC{1}}
\\
\BSS{psHomothetie}[linecolor=green](2,1)\AC{-3}\AC{\BS{pscircle}\AC{1}} }
\\ 	\hline 
\psset{unit=0.8cm}
\begin{pspicture}(-2.5,-2.5)(12.5,8.7)
\psgrid[subgriddiv=0](-2,-2)(12,8)
%\begin{psgraph}[axesstyle=none,xticksize=-3 12 ,yticksize=-2 15, subticks=0,gridcolor=lightgray](0,0)(-2,-3)(15,12){8cm}{8cm} 
\pscircle[linecolor=blue,linestyle=dotted,linewidth=2pt]{1}
\psline(-2,-1)(12,6)
\psHomothetie[linecolor=magenta](2,1){-.5}{\pscircle{1}}
\psHomothetie[linecolor=black](2,1){-1}{\pscircle{1}}
\psHomothetie[linecolor=red](2,1){-2}{\pscircle{1}}
\psHomothetie[linecolor=green](2,1){-3}{\pscircle{1}}
\psdots[linecolor=red,dotsize=8pt](2,1)
\psline[linewidth=3pt,linecolor=red]{->}(2,1)(0,0)	\psdots(8,4)(6,3)(0,0)(3,1.5)(4,2)
%	\end{psgraph}
	\end{pspicture}
\psset{unit=1cm}
	\\ 	\hline 
\end{tabular} 

\psset{unit=1cm}

\newpage
%\section{Placer son dessin}
\SSCT{Placer son dessin}{Placing the picture}

\SbSSCT{Dans le texte}{In the text}

\psset{fillcolor=yellow,fillstyle=none,linecolor=blue}

\bigskip
\TFRGB{dessin directement  dans le texte}{picture in the text}  \psline[linecolor=red](0,0)(4,4) \psline[linecolor=blue](0,0)(4,2) \pscircle[linecolor=green]{2} \TFRGB{ici est inclus le code suivant}{here is the following code} : \BS{}psline[linecolor=red](0,0)(4,4) \BS{}psline[linecolor=blue](0,0)(4,2)  \BS{}pscircle[linecolor=green]\AC{2}
 
% les lignes sont tracées ici .\psline[linecolor=red](0,0)(4,4) \psline[linecolor=blue](0,0)(4,2)

\bigskip
\TFRGB{Le dessin se superpose au texte , Il n'a pas de dimension !}{
The drawing is superimposed on the text, it has no dimension!}

%----------------------------------------------------------
\SbSSCT{Dans un environnement pspicture}{In a pspicture environment}

\begin{tabular}{|l|l|}
\hline  
\multicolumn{2}{|c|}{2 syntaxes}\\ \hline  
\BSS{pspicture}(4,4)
&
\BSS{begin\AC{pspicture}}(4,4)
\\
\BS{}psframe(4,4)
&
\BS{}psframe(4,4)
\\
\BS{}pscircle[linecolor=red](2,2)\AC{1cm}
&
\BS{}pscircle[linecolor=red](2,2)\AC{1cm}
\\
\BSS{endpspicture}
&  
\BSS{end\AC{pspicture}}
\\ 
\hline 

\TFRGB{texte avant}{text before} 
\pspicture(4,4)
\psframe(4,4)
\pscircle[linecolor=red](2,2){1cm}
\endpspicture
\TFRGB{texte après}{text after} 

&  
\begin{pspicture}(4,4)
\TFRGB{texte avant}{text before} 
\psframe(4,4)
\pscircle[linecolor=red](2,2){1cm}
\end{pspicture}
\TFRGB{texte après}{text after} 
\\ \hline 
\end{tabular} 


	
	
	
\SbSSCT{Coupure de l'image}{Clipping the picture}	

\begin{tabular}{|c|p{2cm}|c|}
\hline 
\BS{}begin\AC{pspicture}(4,4) &  &\BS{}begin\AC{pspicture}{\red *}(4,4) \\
\BS{}pscircle[linecolor=red](2,2)\AC{2.5} & & \BS{}pscircle[linecolor=red](2,2)\AC{2.5}
\\ \hline 

\begin{pspicture}(4,4) 
\psframe(4,4)
\pscircle[linecolor=red](2,2){2.5}
\pscircle[linecolor=green](2,2){2}
\pscircle[linecolor=blue](2,2){1}
\end{pspicture}
&  &
\begin{pspicture}*(4,4) 	
\psframe(4,4)
\pscircle[linecolor=red](2,2){2.5}
\pscircle[linecolor=green](2,2){2}
\pscircle[linecolor=blue](2,2){1}
\end{pspicture}
\\  \hline 
\end{tabular} 

\SbSSCT{Rognage partiel}{Partial clipping}

\begin{tabular}{|c|c|}
\hline  
\begin{pspicture}*[shift=*](-2,-2)(3,2) %\psgrid
\psclip {\psdiamond(.5,.5)(2,1)}
\pscircle[linecolor=red]{.5}
\pscircle[linecolor=red]{1}
\endpsclip
\pscircle[linecolor=green]{1.5}
\end{pspicture}
& 
\parbox{8cm}{
\BS{}begin\AC{pspicture}*(-2,-2)(3,2) \\
\BSS{psclip} {\red \AC{\BS{}psdiamond(.5,.5)(2,1)}} \\ 
\BS{}pscircle[linecolor=red]\AC{.5} \\
\BS{}pscircle[linecolor=red]\AC{1} \\
\BSS{endpsclip} \\
\BS{}pscircle[linecolor=green]\AC{1.5} \\
\BS{}end\AC{pspicture}}
\\ \hline 
\end{tabular} 

%---------------------------------------------------------------

\SbSSCTTC{Positionnement  par rapport au texte}{Relative to the text line} %{Relative to the text line \cite{pst-news05}  \cite{pst-news08} }


\begin{tabular}{|c|c|c|c|}
\hline 
\multicolumn{4}{|c|}{ \TFRGB{avant}{before} \BS{}begin\AC{pspicture}[\RDD{shift}=*](1,1)  \BS{}psframe(1,1) \BS{end\AC{pspicture} \TFRGB{après}{after}  }}\\ \hline   
\TFRGB{avant}{before}  \begin{pspicture}(1,1) \psframe(1,1) \end{pspicture} \TFRGB{après}{after} 
&  
\TFRGB{avant}{before}  \begin{pspicture}[shift=*](1,1) \psframe(1,1) \end{pspicture} \TFRGB{après}{after} 
&  
\TFRGB{avant}{before}   \begin{pspicture}[shift=.5cm](1,1) \psframe(1,1) \end{pspicture} \TFRGB{après}{after} 
&  
\TFRGB{avant}{before}  \begin{pspicture}[shift=-.75cm](1,1) \psframe(1,1) \end{pspicture} \TFRGB{après}{after} 
\\ 
\hline {\blue \dft{}} & shift=* & shift=.5cm & shift=-.75cm \\ 
\hline 
\end{tabular}



\newpage 
%\section{Placer des objets}
\SSCT{Placer des objets}{Placing objects}

%======================================================================

\SbSSCT{Commande rput}{Macro rput}


\psset{fillcolor=yellow,fillstyle=solid,linecolor=blue}


\TFRGB{syntaxe : \BS{rput}*[point de référence]\AC{rotation}(coordonnées)\AC{contenu}}{syntax : \BS{rput}*[reference point]\AC{rotation}(coordinates)\AC{contents}}

\bigskip



\SbSbSSCTTC{Rôle de l'astérisque}{Rôle de l'astérisque \protect \footnote{La couleur de fond est en jaune et le point de référence en bleu}}{Role of the asterisk}{Role of the asterisk  \protect \footnote{Fillcolor=yellow and Reference point = blue disk}}

\begin{tabular}{|p{3cm}|p{3cm}|}\hline
\rput(1,0){objet}  &  \rput*(1,0){objet}  \\
\textbf{\BS{}rput}(1,0)\AC{objet} & \textbf{\BS{}rput*}(1,0)\AC{objet} \\ \hline
\end{tabular}
%-------------------------------------

\SbSbSSCT{Point de référence }{Reference point} % \protect \footnote{La couleur de fond est en jaune et le point de référence en bleu} 


\begin{tabular}{|c|l|p{3cm}|c|}

\hline \multicolumn{4}{|c|}{Horizontal}   								\\ \hline
 
\rule[-.35cm]{0mm}{.7cm} l	& \TFRGB{à gauche}{left}		
&  \rput*[l](1,0){objet}\qdisk(1,0){3pt} 			
&  \BS{rput*}[{\red l}](1,0)\AC{objet}\BS{qdisk}(1,0)\AC{3pt} \\ \hline

\rule[-.35cm]{0mm}{.7cm} r 	& \TFRGB{à droite}{right}		
&   \rput*[r](1,0){objet}\qdisk(1,0){3pt}				
& \BS{rput*}[{\red r}](1,0)\AC{objet}\BS{qdisk}(1,0)\AC{3pt}  \\  \hline
 \multicolumn{4}{|c|}{vertical}   								\\ \hline

\rule[-.5cm]{0mm}{.7cm} t	
& \TFRGB{en haut}{top}		
&  \rput*[t](1,0){objet}	\qdisk(1,0){3pt}		
& \BS{rput*}[{\red t}](1,0)\AC{objet}\BS{qdisk}(1,0)\AC{3pt}  \\ \hline


\rule[-.2cm]{0mm}{.7cm }b 	& \TFRGB{en bas}{bottom}		
&  \rput*[b ](1,0){objet}	 \qdisk(1,0){3pt}		&
\BS{rput*} [{\red b}](1,0)\AC{objet}\BS{qdisk}(1,0)\AC{3pt}  \\ \hline


\rule[-.25cm]{0mm}{.7cm} B 	& \TFRGB{sur la ligne d'écriture}{baseline}		& \rput*[B](1,0){objet}	 \qdisk(1,0){3pt} 			& \BS{rput*}[{\red B]}(1,0)\AC{objet}\BS{qdisk}(1,0)\AC{3pt} \\ \hline

 \multicolumn{4}{|c|}{horizontal \TFRGB{et}{and} vertical}   								\\ \hline

\rule[-.5cm]{0mm}{.7cm} rt	
& \TFRGB{à droite et en haut}{right and top}		
& \rput*[rt](1,0){objet}	\qdisk(1,0){3pt} 
& \BS{rput}*[{\red rt}](1,0)\AC{objet}\BS{qdisk}(1,0)\AC{3pt}  \\  \hline 
 \end{tabular} 

%------------------------------------------------------

\SbSbSSCT{Angle de rotation de l'objet}{Rotation angle of the object}


\begin{tabular}{|p{2cm}|p{2cm}|p{2cm}|p{2cm}|p{2cm}|p{2cm}|p{2cm}|}\hline
   \BS{rput*}[t]\AC{45} 
&  \BS{rput*}[t]\AC{90} 
&  \BS{rput*}[b]\AC{90}
&  \BS{rput*}[B]\AC{90} 
&  \BS{rput*}[l]\AC{90}
&  \BS{rput*}[r]\AC{90}  \\ 
\hline \rule[-1cm]{0mm}{2cm}  
\rput*[t]{45}(1,0){objet}\qdisk(1,0){3pt} & 
\rput*[t]{90}(1,0){objet}\qdisk(1,0){3pt} & 
\rput*[b]{90}(1,0){objet}\qdisk(1,0){3pt} & 
\rput*[B]{90}(1,0){objet}\qdisk(1,0){3pt} & 
\rput*[l]{90}(1,0){objet}\qdisk(1,0){3pt} & 
\rput*[r]{90}(1,0){objet}\qdisk(1,0){3pt}  \\ 
\hline 
\end{tabular} 
%
%----------------------------------------------------------

\SbSbSSCT{Angles de rotation  en points cardinaux}{Rotation angle in cardinal points}

\begin{tabular}{|p{2cm}|p{2cm}|p{2cm}|p{2cm}|p{2cm}|p{2cm}|p{2cm}|}
\hline  \TFRGB{haut}{top} 	&  \TFRGB{haut}{top}  	&  \TFRGB{haut}{top}  	& \TFRGB{haut}{top}   & \TFRGB{gauche}{left}   &  \TFRGB{droite}{right}  \\
\TFRGB{ et Est}{and east} 		& \TFRGB{ et Ouest}{and west} & \TFRGB{et Nord}{and north}  & \TFRGB{ et Sud}{and south} & \TFRGB{ et Est}{and east}  & \TFRGB{ et Est}{and east}  \\
\hline  
   \BS{}rput*[t]\AC{E} 
&  \BS{}rput*[t]\AC{W} 
& \BS{}rput*[t]\AC{N}  
&  \BS{}rput*[t]\AC{S}   
&  \BS{}rput*[l]\AC{W}  
&  \BS{}rput*[r]\AC{W}  \\ 
\hline \rule[-1cm]{0mm}{2cm}  \rput*[t]{E}(1,0){objet}\qdisk(1,0){3pt} & \rput*[t]{W}(1,0){objet}\qdisk(1,0){3pt} & \rput*[t]{N}(1,0){objet}\qdisk(1,0){3pt} & \rput*[t]{S}(1,0){objet}\qdisk(1,0){3pt} & \rput*[l]{W}(1,0){objet}\qdisk(1,0){3pt} & \rput*[r]{W}(1,0){objet}\qdisk(1,0){3pt}  \\ 
\hline 
\end{tabular} 



\newpage
%%---------------------------------------------------------------

\SbSSCT{Commande uput}{Macro uput}
\psset{fillcolor=yellow,fillstyle=solid,linecolor=blue}


\TFRGB{ syntaxe :\BSS{uput}*\AC{écartement}[point de référence]\AC{rotation}(coordonnées)\AC{contenu} }{syntax  : \BSS{uput}*\AC{spacing}[Reference point]\AC{rotation}(coordinates)\AC{content}}
\smallskip



\SbSbSSCTTC{Rôle de l'astérisque}{Rôle de l'astérisque \protect \footnote{La couleur de fond est en jaune et le point de référence en bleu}}{Role of the asterisk}{Role of the asterisk  \protect \footnote{Fillcolor=yellow and Reference point = blue disk}}

\begin{tabular}{|p{3cm}|p{3cm}|}\hline
 \uput[r](1,0){objet} &  \uput*[r](1,0){objet}  \\ \hline
\textbf{\BS{}uput}(1,0)\AC{objet} & \textbf{\BS{}uput*}(1,0)\AC{objet} \\ \hline
\end{tabular}
 
\SbSbSSCT{Point de référence  : angle}{Reference point : angle}

\begin{center}
\begin{tabular}{|c|p{3cm}|c|} \hline

 \rule[-1cm]{0mm}{2cm} à 45°	&  \uput*[45](1,0){objet}\qdisk(1,0){3pt} 			&  \BS{uput*}[{\red 45}](1,0)\AC{objet}\BS{qdisk}(1,0)\AC{3pt}\\ \hline

 \rule[-1cm]{0mm}{2cm} à 90° 			&   \uput*[90](1,0){objet}\qdisk(1,0){3pt}				& \BS{uput*}[{\red 90}](1,0)\AC{objet}\BS{qdisk}(1,0)\AC{3pt}\\  \hline

 \rule[-1cm]{0mm}{2cm} à 120°		&  \uput*[120](1,0){objet}	\qdisk(1,0){3pt}		& \BS{uput*}[{\red 120}](1,0)\AC{objet}\BS{qdisk}(1,0)\AC{3pt} \\ \hline
 \end{tabular} 
 \end{center}

\bigskip
%%-----------------------------------------------------------------

\SbSbSSCT{Point de référence  : points cardinaux }{Reference point : letter} 

\begin{tabular}{|p{2cm}|p{2cm}|p{2cm}|p{2cm}|p{2cm}|p{2cm}|p{2cm}|}
\hline 
  \BS{uput*}[{\red u}] 
&  \BS{uput*}[{\red r}]
&  \BS{uput*}[{\red d}]
& \BS{uput*}[{\red l}]   
& \BS{uput*}[{\red ul}]  
&  \BS{uput*}[{\red ur}]  \\ 
\hline  \rule[-1cm]{0mm}{2cm} 
  \uput*[u](1,0){objet}\qdisk(1,0){3pt} 
& \uput*[r](1,0){objet}\qdisk(1,0){3pt} 
& \uput*[d](1,0){objet}\qdisk(1,0){3pt} 
& \uput*[l](1,0){objet}\qdisk(1,0){3pt} 
& \uput*[ul](1,0){objet}\qdisk(1,0){3pt} 
& \uput*[ur](1,0){objet}\qdisk(1,0){3pt}  \\ \hline 
\end{tabular} 



\bigskip
\SbSbSSCT{Angle de rotation de l'objet}{Rotation angle of the object} 


\begin{tabular}{|p{2cm}|p{2cm}|p{2cm}|p{2cm}|p{2cm}|p{2cm}|p{2cm}|}
\hline  
  \BS{uput*}[{\red u}]\AC{{\red 45}} 
& \BS{uput*}[{\red u}]\AC{{\red 90}}
& \BS{uput*}[{\red d}]\AC{{\red 90}}
& \BS{uput*}[{\red l}]\AC{{\red 90}}
& \BS{uput*}[{\red r}]\AC{{\red 90}}  
& \BS{uput*}[{\red ur}]\AC{{\red 90}}  \\ \hline  
\rule[-1.5cm]{0mm}{3cm} 
\uput*[u]{45}(1,0){objet}\qdisk(1,0){3pt} 
& \uput*[u]{90}(1,0){objet}\qdisk(1,0){3pt} 
& \uput*[d]{90}(1,0){objet}\qdisk(1,0){3pt} 
& \uput*[l]{90}(1,0){objet}\qdisk(1,0){3pt} 
& \uput*[r]{90}(1,0){objet}\qdisk(1,0){3pt} 
& \uput*[ur]{90}(1,0){objet}\qdisk(1,0){3pt}  \\ \hline 
\end{tabular} 
%

 
\bigskip
\SbSbSSCT{\'Ecartement de l'objet par rapport au point de référence}{Spacing between object and reference point}

\dft: \RDD{labelsep}= 0.5 pt 
\bigskip

\TFRGB{Exemple}{E\emph{}xample} : 

\begin{tabular}{ll}
\BS{}psset\AC{\red labelsep=1cm } &  \% \TFRGB{nouveau écartement par défaut}{new default spacing} \\ 
\BS{}uput(1,0)\{ à 1cm \}& \%  \TFRGB{utilisation  nouveau écartement par défaut}{use of the new default spacing} \\ 
\BS{}uput {\red \AC{3cm}}(1,0)\AC{à 3cm} & \% \TFRGB{écartement spécifié à 3 cm}{spacing = 3cm}  \\ 
\BS{}uput{\red \AC{3cm}[-30]}(1,0)\AC{à 3cm et à -30°} & \% \TFRGB{écartement spécifié à 3  et à un angle de -30°}{spacing = 3cm angle= -30°} \\  
\BS{}qdisk(1,0)\AC{3pt}  & \%\TFRGB{ point de référence}{Reference point} \\ 
 
\end{tabular} 
\bigskip


{\blue

\psset{labelsep=1cm} 
\uput[0](1,0){ à 1cm }      \uput{3cm}[0](1,0){ à 3cm}  \uput{3cm}[-30](1,0){ à 3cm et à -30°}\qdisk(1,0){3pt}  }

\vspace{2cm}


\SbSSCT{Commande psrotate}{Macro psrotate}
\label{rot}
\psset{fillstyle=none                                                        }
\begin{tabular}{|c|}\hline  
\BSS{psrotate}(2,1)\AC{45}\AC{\BS{psline}(0, 1)(1, 2)(2, 2)(3, 4)}  \BSI{psrotate}{pstricks-add} 
\\ \hline  
\begin{psgraph}[axesstyle=none,xticksize=0 4cm,yticksize=0 4cm,subticks=0](0,0)(4,4){4cm}{4cm} 
\psline[linestyle=dotted](0, 1)(1, 2)(2, 2)(3, 4)
\psrotate(2,2){45}{\psline(0, 1)(1, 2)(2, 2)(3, 4)}
\qdisk(2,2){3pt}

 \end{psgraph}
\\ 
\hline 
\end{tabular} 


%\essais{pstplacess.tex}

\newpage

%%%===================================================
\SSCT{Créer ses couleurs}{Creating color}
%\section{Créer ses couleurs}

\input{Coul}

%\essais{pstcouless.tex}
%
%

%\subsection[Opacité des couleurs ]{Opacité des couleurs \cite{pst-news08}}
\SbSSCT{Opacité des couleurs }{Opacity} 

\input{Coul1}
\newpage

%\subsection[Transparence des couleurs]{Transparence des couleurs \cite{pst-news08}} 
\SbSSCT{Transparence des couleurs}{Tranparency} 

\input{Coul2}


%\subsection[en noir et blanc , en niveaux de gris ou en couleur]{en noir et blanc , en niveaux de gris ou en couleur \cite{pst-news13}}
 
\SbSSCT{en noir et blanc , en niveaux de gris ou en couleur}{Monochrome,Grayscale \& resetColor }
\input{Coul3}

\newpage
%\section{Créer ses commandes}
\SSCT{Créer ses commandes}{Own commands}

%\section{Syntaxe}

\maboite{ \TFRGB{Atention : la création de la commande doit être placée avant}{Warning: the creation of the command must be placed before} \BS{begin}\AC{document} ! }

\TFRGB{syntaxe :\BSS{newcommand}\AC{\BS{}nom}[nombre de variables]\AC{Description}}{syntax :\BSS{newcommand}\AC{\BS{}name}[
number of variables]\AC{Description}}
\bigskip


\textbf{\TFRGB{Exemple : commande avec une variable}{Example : command with one variable} :}

 \emph{\TFRGB{Création}{Creation}}



\begin{tabular}{ll}
\BS{newcommand} 									& \\ %\% début de la nouvelle commande \\ 
\AC{\BS{}maboite}[1]{\red\{	}	 					& \% \TFRGB{commande nommée ma boite et 1 seul d'argument}{command   named \og maboite \fg with one variable } \\ 
\BS{}begin\AC{center}									& \% \TFRGB{centrage sur la ligne}{centering the box  } \\ 
\BS{}psframebox[fillcolor=yellow,fillstyle=solid]{\magenta\{}  & \% \TFRGB{une boite de texte de couleur jaune}{a yellow text box} \\ 
\BS{}parbox\{  .5\BS{}linewidth \} 			& \% \TFRGB{parbox pour limiter la largeur de la boite}{use of \BS{parbox} to set the width of the box} \\ 
\{\BS{}centering  									& \% \TFRGB{centrage du texte dans la boite}{centering the text in the box} \\ 
\#1\} {\magenta\}}\BS{}end\AC{center}							& \%  \TFRGB{\#1 correspond à l'argument}{\#1 will be replaced by the variable} \\ 
{\red\}	} 														&  \\ %\% fin de la nouvelle commande \\ 
\end{tabular} 	

\bigskip
\emph{Utilisation} : \BS{maboite}\AC{contenu}

\maboite{contenu}


\bigskip
\textbf{\TFRGB{Exemple : commande sans variable}{Example : command without variable} :}
 
 \emph{\TFRGB{Création}{ creation}}
 
\BS{}newcommand\AC{\BS{}DFR}{\BS{}psset\AC{unit=.25cm,fillstyle=solid,linewidth=0pt} \BS{}begin\AC{pspicture*}(3,1.5) \BS{}psframe[fillcolor=blue](1,1.5)
\BS{}psframe[fillcolor=white](1,0)(2,1.5)\BS{}psframe[fillcolor=red](2,0)(3,1.5)
\BS{}end\AC{pspicture*} 


\emph{Utilisation}  :  \BS{DFR} \DFR

%
%\essais{pstcdeess.tex}
%
 
 
%\section[Créer ses styles]{Créer ses styles \cite{pst-user}}
\SSCT{Créer ses styles}{Own styles}


\TFRGB{syntaxe : \BSS{newpsstyle}\AC{nom}\AC{paramètres}}{syntax : \BSS{newpsstyle}\AC{name}\AC{parameters}}
\bigskip


\TFRGB{Exemple}{Example} :

 \emph{\TFRGB{Définition du nouveau style}{Definition of the new style}} :

\textbf{\BS{}newpsstyle}\AC{{\red mafleche}}\AC{arrowsize=4pt 6,arrowlength=2,doubleline=true,linewidth=1pt}

\newpsstyle{mafleche}{arrowsize=4pt 6,arrowlength=2,doubleline=true,linewidth=1pt} 
%
\bigskip

 \emph{\TFRGB{Utilisation  du nouveau style}{Using the new style} }:
%
\BS{}psline[{\red style=mafleche}]\AC{->}(0,0)(3,0) \hspace{1cm} 
\begin{pspicture}(4,.5)
\psline[style=mafleche]{->}(0,0)(4,0)
\end{pspicture}

\bigskip

\bigskip
\textbf{\TFRGB{Ajout ou modification d'un paramètre du style}{Adding or changing a parameter style \cite{pst-news08}}}

\bigskip
\addtopsstyle{mafleche}{linecolor=red}
\BSS{addtopsstyle}\AC{mafleche}\AC{linecolor=red}
\begin{pspicture}(4,.5)
\psline[style=mafleche]{->}(0,0)(4,0)
\end{pspicture}

\bigskip
\addtopsstyle{mafleche}{linestyle=dashed}
\BS{}addtopsstyle\AC{mafleche}\AC{linestyle=dashed}
\begin{pspicture}(4,.5)
\psline[style=mafleche]{->}(0,0)(4,0)
\end{pspicture}
%
%\essais{pststyless.tex}

 
%\section[Créer ses objets]{Créer ses objets \cite{pst-user}}
\SSCT{Créer ses objets}{Own objects}



\TFRGB{syntaxe : \BSS{newpsobject}\AC{nom}\AC{objet}\AC{paramètres}}{syntaxe : \BSS{newpsobject}\AC{name}\AC{object}\AC{paramèters}}:
\smallskip

\TFRGB{Exemple}{Example} :

\BS{}newpsobject\AC{maboite}\AC{psframebox}\AC{fillstyle=solid,fillcolor=yellow,linewidth=2pt,linecolor=red}
\smallskip

\newpsobject{maboite}{psframebox}{fillstyle=solid,fillcolor=yellow,linewidth=2pt,linecolor=red}

\bigskip

\BS{}maboite\AC{\TFRGB{ma boite personnalisée}{my custom box}} \maboite{\TFRGB{ma boite personnalisée}{my custom box}} 


%%%%%%%======================================================================
%\section{Mettre des objets en boîte}
\SSCT{Mettre des objets en boîte}{Boxed objects}

%\subsection[Les différentes boites]{Les différentes boites  \cite{pst-user}}


\psset{fillcolor=yellow,fillstyle=solid,linecolor=blue}


\psset{linecolor=blue,linewidth=2pt,fillcolor=yellow}

\begin{center}
\begin{tabular}{|c|c|c|}\hline
\multicolumn{3}{|c|}{ \BS{psframebox*}\AC{objet}} \\ \hline 
& \textbf{\TFRGB{sans astérisque}{without asterisk}} & \textbf{\TFRGB{avec astérisque}{with asterisk}} \\
\hline \BSS{psframebox}* 	&\rule[-.5cm]{0pt}{1cm}  \psframebox{objet} 	& \psframebox*{objet} \\ 

\hline \BSS{psdblframebox}*	&\rule[-.5cm]{0pt}{1cm}  \psdblframebox{objet} 	& \psdblframebox*{objet} \\ 

\hline \BSS{psshadowbox}* &\rule[-.5cm]{0pt}{1cm}\psshadowbox{objet} 	& \psshadowbox*{objet} \\

\hline \BSS{pscirclebox}*	&\rule[-.5cm]{0pt}{1.3cm}  \pscirclebox{objet}  & \pscirclebox*{objet} \\

\hline \BSS{psovalbox}*		&\rule[-.5cm]{0pt}{1cm}  \psovalbox{objet} 	& \psovalbox*{objet} \\

\hline \BSS{psdiabox}*		&\rule[-.5cm]{0pt}{1.3cm}  \psdiabox{objet}  	& \psdiabox*{objet} \\

\hline \BSS{pstribox}*		&\rule[-.3cm]{0pt}{1.3cm}  \pstribox{objet}  	& \pstribox*{objet} \\

\hline 
\end{tabular} 
\end{center}

\TFRGB{Exemple}{Example} :  \BS{}psdiabox\AC{\BS{}DFR}  \psdiabox{\DFR} 


\subsection{Options }



%\subsubsection{framesep=dimension}

\begin{tabular}{|c|c|c|}
\hline 
\multicolumn{3}{|c|}{ \BS{psframebox} \RDD{framesep}=.5cm]\AC{framesep=.5cm}} \\ \hline 
\psframebox{\dft{}}
&  
\psframebox[framesep=0cm]{framesep=0cm}
&  
\psframebox[framesep=.5cm]{framesep=.5cm}
\\ 
\hline {\blue \dft{} :  framesep=3pt }& framesep=0cm & framesep=.5cm \\ 
\hline 
\end{tabular} 

 
\bigskip


%\subsubsection{boxsep=true/false}

\begin{tabular}{|c|} \hline
\RDD{boxsep} {\blue  \dft : true} (\TFRGB{Ce paramètre ne s'applique qu'à}{Apply only to}  \BS{psframebox}, \BS{pscirclebox} \TFRGB{et}{and} \BS{psovalbox}) 
\\ \hline 
\TFRGB{texte avant}{text before} \psframebox[boxsep=true]{boxsep=true} \TFRGB{texte entre les 2 boites}{text between} \psframebox[boxsep=false]{boxsep=false} \TFRGB{texte après}{text after}
 \\ \hline 
\end{tabular} 



\bigskip
%\subsubsection{Options trimode pour  \BS{}pstribox}


\begin{tabular}{|c|c|c|} \hline
\multicolumn{3}{|c|}{ Option \RDD{trimode}  \TFRGB{pour}{only for}  \BS{pstribox} } \\ \hline 
& \textbf{\TFRGB{sans astérisque}{without asterisk}} & \textbf{\TFRGB{avec astérisque}{with asterisk}}  \\
\hline \BS{}pstribox*[{\red trimode=U}] &\rule[-.3cm]{0pt}{1.2cm} \pstribox[trimode=U]{objet}  & \pstribox*[trimode=U]{objet} \\
\hline \BS{}pstribox*[{\red trimode=D}] & \rule[-.5cm]{0pt}{1cm}\pstribox[trimode=D]{objet} & \pstribox*[trimode=D]{objet} \\
\hline \BS{}pstribox*[{\red trimode=R}] & \rule[-.5cm]{0pt}{1.2cm}\pstribox[trimode=R]{objet}  & \pstribox*[trimode=R]{objet} \\
\hline \BS{}pstribox*[{\red trimode=L}] & \rule[-.5cm]{0pt}{1.2cm}\pstribox[trimode=L]{objet}  & \pstribox*[trimode=L]{objet} \\
\hline 
\end{tabular} 

\bigskip
%\subsubsection{Limitation de la taille de la boite}

\begin{tabular}{|c|}\hline    
\BS{psframebox}\AC{\BS{parbox}[l]\AC{3cm}\AC{\TFRGB{utilisation de \BSS{parbox} pour limiter la largeur de la boite à 3cm}{use of \BSS{parbox} to limit the width of the framebox}}}
\\ \hline 
\psframebox{\parbox[l]{4cm}{\TFRGB{utilisation de \BSS{parbox} pour limiter la largeur de la boite à 3cm}{use of \BSS{parbox} to limit the width of the framebox}}}
 \\ \hline 
\end{tabular} 




\begin{center}


\end{center}

\psset{fillcolor=white}
%
%\essais{pstboit1ess.tex}

%\section[Mettre des objets en cadre]{Mettre des objets en cadre \cite{pst-news08}}
\SSCT{Mettre des objets en cadre}{Framed objects}

\SbSSCT{Texte dans un cadre]}{Text in a frame} 

\begin{tabular}{|c|c|}
\hline 
\begin{pspicture}(-0.2,-0.2)(4.2,1.7)
\psTextFrame(0,0)(4,1.5){\TFRGB{texte}{text}}
\end{pspicture}
&  
\begin{pspicture}(-0.2,-0.2)(4.2,1.7)
\psTextFrame*[linecolor=yellow](0,0)(4,1.5){\TFRGB{texte}{text}}
\end{pspicture}  
\\ 
\hline  
\BSS{psTextFrame}(0,0)(4,2)\AC{\TFRGB{texte}{text}}
&
\BS{}psTextFrame{\red *}[linecolor=yellow](0,0)(4,2)\AC{\TFRGB{texte}{text}}
\\ 
\hline 
\end{tabular} 

\SbSbSSCT{Problème de dépassement du cadre}{Problem : text too long for the frame}
 
 \begin{tabular}{|c|} \hline  
 \BS{psTextFrame}(0,0)(4,1)\AC{\TFRGB{Problème de dépassement du cadre}{Problem : text too long for the frame}}
\\ \hline
\begin{pspicture}(-.2,-0.2)(4.5,1.2)
\psTextFrame(0,0)(4,1){\TFRGB{Problème de dépassement du cadre}{Problem : text too long for the frame}}
\end{pspicture}
\\  \hline 
 \end{tabular} 


\bigskip

\textbf{Solutions}

\begin{tabular}{|c|c|} \hline 
  \begin{pspicture}(-.2,-0.2)(4.2,1.2)
 \psTextFrame(0,0)(4,1){\parbox{3.5cm}{\TFRGB{Problème de dépassement du cadre : problème résolu}{ text too long for the frame : Problem solved}}}
 \end{pspicture}
&
  \begin{pspicture}(-.2,-0.2)(4.2,1.2)
  \psTextFrame(0,0)(4,1){\begin{minipage}[c]{3.5cm}
 \TFRGB{Problème de dépassement du cadre : problème résolu}{ text too long for the frame : Problem solved} \end{minipage}}
  \end{pspicture}
\\ \hline  
\parbox{6.5cm}{ 
\BS{psTextFrame}(0,0)(4,2)\AC{ {\red\BS{parbox}}\AC{3.5cm}\AC{\TFRGB{Problème de dépassement du cadre : problème résolu}{ text too long for the frame : Problem solved}  }}}
&
\parbox{6.5cm}{ \BS{}psTextFrame(0,0)(4,2)\AC{  {\red \BS{}begin\AC{minipage}[c]\AC{3.5cm}}

\TFRGB{Problème de dépassement du cadre : problème résolu}{ text too long for the frame : Problem solved}
 {\red  \BS{end\AC{minipage}}}}}
\\ \hline 
\end{tabular} 

\SbSbSSCT{Rotation du texte}{Text rotation in the frame}

\begin{tabular}{|c|l|}
\hline  \begin{pspicture}[shift=*](-.2,-0.2)(4.2,2.2)
 \psTextFrame[rot=90](0,0)(4,2){\TFRGB{texte}{text}}
 \end{pspicture} & 
 \BS{}psTextFrame[\RDD{rot}=90](0,0.5)(4,2)\AC{\TFRGB{texte}{text}} \\ 
\hline 
\end{tabular}  
%
 \SbSbSSCT{Position du texte}{Position of the text in the frame}

 \begin{tabular}{|c|c|c|c|c|}
 \hline
 \multicolumn{5}{|c|}{ \BS{}psTextFrame[\RDD{ref=l}](0,0)(2,2)\AC{\TFRGB{texte}{text}}} \\
 \hline  
\begin{pspicture}(-.2,-0.2)(2.2,2.2)
\psTextFrame[ref=l](0,0)(2,2){\TFRGB{texte}{text}}
\end{pspicture} 
& 
\begin{pspicture}(-.2,-0.2)(2.2,2.2)
\psTextFrame[ref=r](0,0)(2,2){\TFRGB{texte}{text}}
\end{pspicture}   
&
\begin{pspicture}(-.2,-0.2)(2.2,2.2)
\psTextFrame[ref=t](0,0)(2,2){\TFRGB{texte}{text}}
\end{pspicture}
&
\begin{pspicture}(-.2,-0.2)(2.2,2.2)
 \psTextFrame[ref=b](0,0)(2,2){\TFRGB{texte}{text}}
 \end{pspicture}
&
\begin{pspicture}(-.2,-0.2)(2.2,2.2)
\psTextFrame[ref=B](0,0)(2,2){\TFRGB{texte}{text}}
\end{pspicture}
  \\ 
 \hline ref=l & ref=r  & ref=t  & ref=b & ref=B \\ 
 \hline 
 \end{tabular} 
\smallskip



\newpage
%\section[Mettre des objets en bouton]{Mettre des objets en bouton  \cite{pst-fr3d}}
\SSCT{Mettre des objets en bouton}{Buttoned objects}


\psset{fillcolor=yellow,fillstyle=solid,linecolor=blue}


\TFRGB{utilisation du module}{Package} \og \textbf{pst-fr3d} \fg{} \label{pst-fr3d}

\TFRGB{syntaxe : \BSS{PstFrameBoxThreeD}[paramètres]\AC{Contenu}}{syntax :\BSS{PstFrameBoxThreeD}[paramèters]\AC{Content}}  \BSI{PstFrameBoxThreeD}{pst-fr3d} 

\SbSSCT{Sans options}{Without option}


\begin{tabular}{|c|c|}
\hline  
\PstFrameBoxThreeD{Button}
&  
\PstFrameBoxThreeD{\shortstack{Un!\\Deux!\\Trois!}}
\\ \hline 

\BS{PstFrameBoxThreeD}\AC{Button}
 &
 \BS{PstFrameBoxThreeD}\AC{\BS{shortstack}\AC{Un!\BS{}\BS{}Deux!\BS{}\BS{}Trois!}}
\\ \hline

\PstFrameBoxThreeD{\DFR}
& 
\PstFrameBoxThreeD{$\dfrac{1}{\sqrt{3}}$} 
\\ \hline
  
\BS{PstFrameBoxThreeD}\AC{\BS{}DFR}
&  
\BS{PstFrameBoxThreeD}\AC{\$\BS{dfrac}\AC{1}\AC{\BS{sqrt}\AC{3}}\$}
\\ \hline 
\end{tabular} 


\SbSSCT{Dimensionnement}{Sizing}

\begin{tabular}{|l|c|c|c|} \hline
\RDD{doublesep} \rule[-.5cm]{0pt}{1.2cm} \RDI{doublesep}{pst-fr3d} &
\PstFrameBoxThreeD[doublesep=0]{doublesep=0} &
\PstFrameBoxThreeD[doublesep=0.2]{doublesep=0.2}&
\PstFrameBoxThreeD[doublesep=0.4]{doublesep=0.4}\\ \hline
\RDD{framesep} \rule[-.5cm]{0pt}{1.2cm} \RDI{framesep}{pst-fr3d}  &
\PstFrameBoxThreeD[framesep=0]{framesep=0}&
\PstFrameBoxThreeD[framesep=0.3]{framesep=0.3}&
\PstFrameBoxThreeD[framesep=0.5]{framesep=0.5}\\ \hline
\RDD{linewidth}\rule[-.5cm]{0pt}{1.2cm} \RDI{linewidth}{pst-fr3d} &
\PstFrameBoxThreeD[linewidth=0]{linewidth=0}&
\PstFrameBoxThreeD[linewidth=0.1]{linewidth=0.1}&
\PstFrameBoxThreeD[linewidth=0.2]{linewidth=0.2}\\ \hline
\RDD{framearc} \rule[-.5cm]{0pt}{1.2cm} \RDI{framearc}{pst-fr3d}  &
 \PstFrameBoxThreeD[framearc=0.3]{framearc=0.3}&
 \PstFrameBoxThreeD[framearc=0.6]{framearc=0.6}&
 \PstFrameBoxThreeD[framearc=1]{framearc=1}\\ \hline
\end{tabular}



\subsection{Aspect}
\psset{doublesep=0.1}

\begin{tabular}{|l|c|c|c|c|} \hline
\multicolumn{5}{|c|}{\BS{}PstFrameBoxThreeD[{\red FrameBoxThreeDColorHSB =0 0.3 1}]\AC{0 0.3 1} } \\ \hline 
\RDD{FrameBoxThreeDColorHSB} \rule[-.5cm]{0pt}{1.2cm} \RDI{FrameBoxThreeDColorHSB}{pst-fr3d} &
\PstFrameBoxThreeD[FrameBoxThreeDColorHSB =0 0.3 1]{0 0.3 1}&
\PstFrameBoxThreeD[FrameBoxThreeDColorHSB =0.3 0.3 1]{0.3 0.3 1}&
\PstFrameBoxThreeD[FrameBoxThreeDColorHSB =0.6 0.3 1]{0.6 0.3 1}&
\PstFrameBoxThreeD[FrameBoxThreeDColorHSB =1 0.3 1]{1 0.3 1}\\\hline 
\RDD{FrameBoxThreeDOn}=true/false \rule[-.5cm]{0pt}{1.2cm} \RDI{FrameBoxThreeDOn}{pst-fr3d} &
\PstFrameBoxThreeD[FrameBoxThreeDOn =true]{true}&
\PstFrameBoxThreeD[ FrameBoxThreeDOn =false]{false}&
\PstFrameBoxThreeD[FrameBoxThreeDColorHSB =0.6 0.6 0.8, FrameBoxThreeDOn =true]{true}&
 \PstFrameBoxThreeD[FrameBoxThreeDColorHSB =0.6 0.6 0.8, FrameBoxThreeDOn =false]{false}\\\hline 
\psset{linewidth=0.2} 
\RDD{FrameBoxThreeDOpposite}=true/false \RDI{FrameBoxThreeDOpposite}{pst-fr3d}  \rule[-.5cm]{0pt}{1.2cm}&
\PstFrameBoxThreeD[FrameBoxThreeDColorHSB =0.3 0.3 0.8, FrameBoxThreeDOn =false]{=false}&
\PstFrameBoxThreeD[FrameBoxThreeDColorHSB =0.3 0.3 0.8, FrameBoxThreeDOpposite =true]{=true}&
\PstFrameBoxThreeD[FrameBoxThreeDColorHSB =0.6 0.6 0.8, FrameBoxThreeDOn =false, FrameBoxThreeDBrightnessDistance =0.5]{=false}&
\PstFrameBoxThreeD[FrameBoxThreeDColorHSB =0.6 0.6 0.8, FrameBoxThreeDOpposite =true, FrameBoxThreeDBrightnessDistance =0.5]{=true}\\ \hline 

\end{tabular}
\psset{doublesep=0.2}



\begin{tabular}{|l|c|c|c|c|c|} \hline
\RDD{FrameBoxThreeDBrightnessDistance} \RDI{FrameBoxThreeDBrightnessDistance}{pst-fr3d}  \rule[-.5cm]{0pt}{1.2cm}&

\PstFrameBoxThreeD[FrameBoxThreeDColorHSB =0.6 0.3 0.5, FrameBoxThreeDBrightnessDistance =-0.5]{-0.5}&
 \PstFrameBoxThreeD[FrameBoxThreeDColorHSB =0.6 0.3 0.5, FrameBoxThreeDBrightnessDistance =-0.2]{-0.2}&
\PstFrameBoxThreeD[FrameBoxThreeDColorHSB =0.6 0.3 0.5,FrameBoxThreeDBrightnessDistance =0]{0}&
 \PstFrameBoxThreeD[FrameBoxThreeDColorHSB =0.6 0.3 0.5, FrameBoxThreeDBrightnessDistance =0.3]{d\'efaut}&
 \PstFrameBoxThreeD[FrameBoxThreeDColorHSB =0.6 0.3 0.5, FrameBoxThreeDBrightnessDistance =0.5]{0.5}\\
 & mini & & nul & défaut & maxi\\ \hline
\end{tabular}

\psset{doublesep=1.25\pslinewidth}

\newpage
%\section{Annuler des objets}
\SSCT{Annuler des objets}{Canceling objects}


\label{ann}

\begin{tabular}{|c|c|}	\hline  
\psCancel{Objet}
&
\psCancel{$\dfrac{1}{2}$}
\\ \hline 
\BSS{psCancel}\AC{Objet} \BSI{psCancel}{pstricks-add}& \BSS{psCancel}\AC{\$\BS{dfrac}\AC{1}\AC{2}\$} 
\\ \hline
& \\ 
\psCancel*{Objet}
& 
\psCancel*[opacity=0.5]{Objet} \\
& \\ 	\hline  
\BSS{psCancel*}\AC{Objet} & \BSS{psCancel*}[opacity=0.5]\AC{Objet} 
\\ \hline 
\end{tabular} 


\bigskip


\begin{tabular}{|c|c|c|}	\hline 
\multicolumn{3}{|c|}{\BS{psCancel}[\RDD{cancelType}=x]\AC{Objet} \RDI{cancelType}{pstricks-add} } \\ \hline
& &\\
\psCancel[cancelType=x]{Objet}
 &  
\psCancel[cancelType=s]{Objet} 
 &  
\psCancel[cancelType=b]{Objet} \\
& &
 \\ \hline  
[cancelType=x] & [cancelType=s] & [cancelType=b] \\ \hline 
\end{tabular} 


\newpage

%%%%%%======================================================================
%\section{Des lignes et liaisons spéciales}
\SSCT{Des lignes et liaisons spéciales}{
Lines and special connections }

%\subsection[Des lignes et liaisons spéciales]{Trait à main levé  \cite{pstricks-add}}
\SbSSCT{Trait à main levé }{Line by hand}

%\subsection{Trait à main levé  \cite{pst-add}}

\begin{center}
\begin{tabular}{|c|c|c|}
\hline  
\multicolumn{3}{|c|}{\BSS{pslineByHand}(0,0)(4,0)  \BSI{pslineByHand}{pstricks-add}} \\\hline  
%\rule[-1cm]{0pt}{2cm} 
\begin{pspicture}(-.15,-.5)(4.15 ,0.5)
\pslineByHand(0,0)(4,0) 
 \end{pspicture}
& 
\begin{pspicture}(-.15,-.5)(4.15 ,0.5)
\pslineByHand[varsteptol=5](0,0)(4,0)
 \end{pspicture}
  &
  \begin{pspicture}(-.15,-.5)(4.15 ,0.5)
  \pslineByHand[ VarStepEpsilon=.4](0,0)(4,0)
   \end{pspicture}
\\ \hline
{\blue \dft }	& 
\RDD{varsteptol}=5 \RDI{varsteptol}{pstricks-add}& \RDD{VarStepEpsilon}=.4 \RDI{VarStepEpsilon}{pstricks-add} 
\\ \hline
		& {\blue \dft : 2 } 			& {\blue \dft : .8  }				\\ \hline
\end{tabular}
\end{center}

\newpage
\SbSSCT{Symboles sur ligne}{Symbol on the line  }
\begin{center}
\begin{tabular}{|p{6cm}|p{6cm}|}
\hline  
%\multicolumn{2}{|c|}{paramètres \textbf{ArrowInside=->}} \\ \hline

\rule[-.5cm]{0pt}{1cm}	\psline[ArrowInside=->](0.5,0)(5,0)  		&
\rule[-.5cm]{0pt}{1cm} 	\psline[ArrowInside=-<](0.5,0)(5,0) 	\\ \hline
\BS{}psline[{\RDD{ArrowInside}=->}](0.5,0)(5,0) \RDI{ArrowInside}{pstricks-add} & 
\BS{}psline[{\red ArrowInside=-<}](0.5,0)(5,0)\\ \hline

\rule[-.5cm]{0pt}{1cm}	\psline[ ArrowInside=->>](0.5,0)(5,0)  		&
\rule[-.5cm]{0pt}{1cm} 	\psline[ArrowInside=-<<](0.5,0)(5,0) 	\\ \hline
\BS{}psline[{\red ArrowInside=->>}](0.5,0)(5,0) & \BS{}psline[{\red ArrowInside=-<<}](0.5,0)(5,0)\\ \hline

\rule[-.5cm]{0pt}{1cm}	\psline[ ArrowInside=-|](0.5,0)(5,0)  		&
\rule[-.5cm]{0pt}{1cm} 	\psline[ArrowInside=-|*](0.5,0)(5,0) 	\\ \hline
\BS{}psline[{\red ArrowInside=-|}](0.5,0)(5,0) & \BS{}psline[{\red ArrowInside=-|*}](0.5,0)(5,0)\\ \hline
\rule[-.5cm]{0pt}{1cm}	\psline[ ArrowInside=-[ ](0.5,0)(5,0)  		&
\rule[-.5cm]{0pt}{1cm}	\psline[ ArrowInside=-( ](0.5,0)(5,0)  \psline[ArrowInside=-) ](0.5,0)(5,0)		
\\ \hline
\BS{}psline[{\red ArrowInside=-[}](0.5,0)(5,0) & \BS{}psline[{\red ArrowInside=-( }](0.5,0)(5,0)\\ \hline


\rule[-.5cm]{0pt}{1cm}	\psline[ ArrowInside=-o](0.5,0)(5,0)  		&
\rule[-.5cm]{0pt}{1cm} 	\psline[ArrowInside=-*](0.5,0)(5,0) 	\\ \hline
\BS{}psline[{\red ArrowInside=-o}](0.5,0)(5,0) & \BS{}psline[{\red ArrowInside=-*}](0.5,0)(5,0)\\ \hline

\rule[-.5cm]{0pt}{1cm}	\psline[ ArrowInside=->|](0.5,0)(5,0)  		&
\rule[-.5cm]{0pt}{1cm} 	\psline[ArrowInside=-<|](0.5,0)(5,0) 	\\ \hline
\BS{}psline[{\red ArrowInside=->|}](0.5,0)(5,0) & \BS{}psline[{\red ArrowInside=-<|}](0.5,0)(5,0)\\ \hline

\rule[-.5cm]{0pt}{1cm}	\psline[ ArrowInside=-h](0.5,0)(5,0)  		&
\rule[-.5cm]{0pt}{1cm} 	\psline[ArrowInside=-H](0.5,0)(5,0) 	\\ \hline
\BS{}psline[{\red ArrowInside=-h}](0.5,0)(5,0) & \BS{}psline[{\red ArrowInside=-H}](0.5,0)(5,0)\\ \hline

\rule[-.5cm]{0pt}{1cm}	\psline[ ArrowInside=-v](0.5,0)(5,0)  		&
\rule[-.5cm]{0pt}{1cm} 	\psline[ArrowInside=-V](0.5,0)(5,0) 	\\ \hline
\BS{}psline[{\red ArrowInside=-v}](0.5,0)(5,0) & \BS{}psline[{\red ArrowInside=-V}](0.5,0)(5,0)\\ \hline

\rule[-.5cm]{0pt}{1cm}	\psline[ ArrowInside=-f](0.5,0)(5,0)  		&
\rule[-.5cm]{0pt}{1cm} 	\psline[ArrowInside=-F](0.5,0)(5,0) 	\\ \hline
\BS{}psline[{\red ArrowInside=-f}](0.5,0)(5,0) & \BS{}psline[{\red ArrowInside=-F}](0.5,0)(5,0)\\ \hline

\rule[-.5cm]{0pt}{1cm}	\psline[ ArrowInside=-t](0.5,0)(5,0)  		&
\rule[-.5cm]{0pt}{1cm} 	\psline[ArrowInside=-T](0.5,0)(5,0) 	\\ \hline
\BS{}psline[{\red ArrowInside=-t}](0.5,0)(5,0) & \BS{}psline[{\red ArrowInside=-T}](0.5,0)(5,0)\\ \hline
\end{tabular}
\end{center}

\begin{center}
\begin{tabular}{|p{6cm}|p{6cm}|}
\hline  
\multicolumn{2}{|c|}{\TFRGB{paramètres supplémentaires}{more option} \footnotemark[1] } \\ \hline

\rule[-.5cm]{0pt}{1cm} 	\psline[ArrowInsidePos=.3,ArrowInside=->](0.5,0)(5,0) 	& 
\rule[-.5cm]{0pt}{1cm}	\psline[ArrowInsidePos=20,ArrowInside=->](0.5,0)(5,0)  	\\ \hline
\RDD{ArrowInsidePos}=.3 (soit à 30\%) \RDI{ArrowInsidePos}{pstricks-add} & 
{\red ArrowInsidePos}=20 (soit à 20 pt)\\  \hline
 
\rule[-.5cm]{0pt}{1cm} 	\psline[ArrowInsideNo=5,ArrowInside=->](0.5,0)(5,0) 	& 
\rule[-.5cm]{0pt}{1cm} 	\psline[ArrowInsideNo=3,ArrowInside=-t](0.5,0)(5,0) 	\\ \hline
\RDD{ArrowInsideNo}=5,ArrowInside=-> \RDI{ArrowInsideNo}{pstricks-add} & 
{\red ArrowInsideNo}=3,ArrowInside=-t 	\\ \hline

\rule[-.5cm]{0pt}{1cm} 	\psline[ArrowInsideNo=5,ArrowInside=->,ArrowInsideOffset=0.1](0.5,0)(5,0) 	& 
\rule[-.5cm]{0pt}{1cm} 	\psline[ArrowInsideNo=3,ArrowInside=-t,ArrowInsideOffset=-.2](0.5,0)(5,0) 	\\ \hline
\RDD{ArrowInsideOffset}=0.1 \RDI{ArrowInsideOffset}{pstricks-add}  & {\red ArrowInsideOffset}=-0.2 	\\ \hline
\end{tabular}
\end{center}

\footnotetext[1]{\TFRGB{pour d'autres paramètres voir}{for other parameters see } page \pageref{parametres}}


%\subsection[Tracer avec des symboles]{Tracer avec des symboles  \protect \footnote{ valable seulement pour \BS{}psline, \BS{pspolygon},\BS{pscurve}
%\BS{psccurve} et \BS{psbezier}} \cite{pst-news10}} 
\SbSSCT{Tracer avec des symboles}{Drawing with symbols} \protect \footnote{ valable seulement pour \BS{}psline, \BS{pspolygon},\BS{pscurve}
\BS{psccurve} et \BS{psbezier}}

%\SbSSCT{Tracer avec des symboles}{Drawing with symbols}


\psset{fillstyle=none,linewidth=1pt} 

\psset{unit=.75cm}
\begin{tabular}{|c|c|c|c|}\hline 
\multicolumn{4}{|c|}{\BS{}psline[{\red linestyle=}\RDDD{symbol}{linestyle}](-2,0)(2,0)} \\
\hline   

\begin{pspicture}(-2.5,-.5)(2.5,2.5)
 \psline[linestyle=dotted,linecolor=red](-2,0)(1,0)(2,2)
 \psline[linestyle=symbol](-2,0)(1,0)(2,2)
\end{pspicture} 
 &
\begin{pspicture}(-2.5,-.5)(2.5,2.5)
 \psline[linestyle=dotted,linecolor=red](-2,0)(1,0)(2,2)
 \psline[linestyle=symbol,symbolStep=.5](-2,0)(1,0)(2,2)
\end{pspicture}
 &
\begin{pspicture}(-2.5,-.5)(2.5,2.5)
 \psline[linestyle=dotted,linecolor=red](-2,0)(1,0)(2,2)
 \psline[linestyle=symbol,symbolWidth=.5cm](-2,0)(1,0)(2,2)
\end{pspicture}
 &
\begin{pspicture}(-2.5,-.5)(2.5,2.5)
 \psline[linestyle=dotted,linecolor=red](-2,0)(1,0)(2,2)
 \psline[linestyle=symbol,rotateSymbol=true](-2,0)(1,0)(2,2)
\end{pspicture} 
 \\
  \hline 
\dft & \RDD{symbolStep}=.5 & \RDD{symbolWidth}=.5cm &  \RDD{rotateSymbol}=true \\
& {\blue \dft : 20pt} & {\blue \dft : 10pt } & {\blue \dft : false}\\
\hline 
	\end{tabular}
\bigskip
	
\begin{tabular}{|c|c|c|c|}\hline 	

\multicolumn{4}{|c|}{ \BS{pscurve[linestyle=symbol,\RDD{symbolFont}{\red =PSTricksDotFont}](-2,1)(0,-1)(2,1.5)} }\\
\hline  
\begin{pspicture}(-2.5,-2)(2.5,2)
\pscurve[linestyle=dotted,linecolor=red](-2,1)(0,-1)(2,1.5)
\pscurve[linestyle=symbol,symbolFont=PSTricksDotFont](-2,1)(0,-1)(2,1.5)
\end{pspicture}
 &
\begin{pspicture}(-2.5,-2)(2.5,2)
\pscurve[linestyle=dotted,linecolor=red](-2,1)(0,-1)(2,1.5)
\pscurve[linestyle=symbol,symbolFont=PSTricksDotFont,symbolWidth=1](-2,1)(0,-1)(2,1.5)
\end{pspicture}
 &
\begin{pspicture}(-2.5,-2)(2.5,2)
\pscurve[linestyle=dotted,linecolor=red](-2,1)(0,-1)(2,1.5)
\pscurve[linestyle=symbol,symbolFont=PSTricksDotFont,rotateSymbol=true](-2,1)(0,-1)(2,1.5)
\end{pspicture}
&
\begin{pspicture}(-2.5,-2)(2.5,2)
\pscurve[linestyle=dotted,linecolor=red](-2,1)(0,-1)(2,1.5)
\pscurve[linestyle=symbol,symbolFont=PSTricksDotFont,rotateSymbol=true,startAngle=45](-2,1)(0,-1)(2,1.5)
\end{pspicture} 
 \\ \hline 
{\blue \dft} & \red{symbolWidth} =1cm & \red{rotateSymbol} =true &  rotateSymbol=true \\
  & {\blue \dft : 10pt} & {\blue \dft : false }&  \RDD{startAngle=45} 
\\ 
  \hline  
	\end{tabular}
  \bigskip
  	
  \begin{tabular}{|c|c|c|c|}\hline 
\multicolumn{2}{|c|}{ symbolFont=Dingbats ( \dft) } &                                                                      
\multicolumn{2}{|c|}{ symbolFont=PSTricksDotFont } \\
\hline   

\begin{pspicture}(-2.5,-.5)(2.5,2.5)
 \psline[linestyle=dotted,linecolor=red](-2,0)(1,0)(2,2)
 \psline[linestyle=symbol,symbol=u,symbolFont=Dingbats](-2,0)(1,0)(2,2)
\end{pspicture} 
 &
\begin{pspicture}(-2.5,-.5)(2.5,2.5)
 \psline[linestyle=dotted,linecolor=red](-2,0)(1,0)(2,2)
 \psline[linestyle=symbol,symbol=120,symbolFont=Dingbats](-2,0)(1,0)(2,2)
\end{pspicture}
 &
\begin{pspicture}(-2.5,-.5)(2.5,2.5)
 \psline[linestyle=dotted,linecolor=red](-2,0)(1,0)(2,2)
 \psline[linestyle=symbol,symbol=u,symbolFont=PSTricksDotFont](-2,0)(1,0)(2,2)
\end{pspicture}
 &
\begin{pspicture}(-2.5,-.5)(2.5,2.5)
 \psline[linestyle=dotted,linecolor=red](-2,0)(1,0)(2,2)
 \psline[linestyle=symbol,symbol=120,symbolFont=PSTricksDotFont](-2,0)(1,0)(2,2)
\end{pspicture} 
 \\
  \hline 
\RDD{symbol}=u & symbol=120 & symbol=u & symbol=120
\\  \hline 
\end{tabular} 



 
\newpage
\SbSbSSCT{Symboles disponibles avec le clavier}{Symbols available with the keyboard}
\psset[pstricks]{symbolFont=Dingbats}
\begin{center}
	\begin{tabular}{| c @{:} c |c @{:} c |c @{:} c |c @{:} c |c @{:} c |c @{:} c |c @{:} c | }  \hline
\multicolumn{14}{|c|}{symbolFont=Dingbats ( \dft)} \\  \hline
	A &  \pspicture[shift=*](-1,-.5)(1,.5) \psline[linestyle=symbol,symbol=A](-.5,.2)(.5,.2)\endpspicture &
	B &  \pspicture[shift=*](-1,-.5)(1,.5) \psline[linestyle=symbol,symbol=B](-.5,.2)(.5,.2)\endpspicture &
	C &  \pspicture[shift=*](-1,-.5)(1,.5) \psline[linestyle=symbol,symbol=C](-.5,.2)(.5,.2)\endpspicture &
	D &  \pspicture[shift=*](-1,-.5)(1,.5) \psline[linestyle=symbol,symbol=D](-.5,.2)(.5,.2)\endpspicture &
	E &  \pspicture[shift=*](-1,-.5)(1,.5) \psline[linestyle=symbol,symbol=E](-.5,.2)(.5,.2)\endpspicture &
	F &  \pspicture[shift=*](-1,-.5)(1,.5) \psline[linestyle=symbol,symbol=F](-.5,.2)(.5,.2)\endpspicture &
	G &  \pspicture[shift=*](-1,-.5)(1,.5) \psline[linestyle=symbol,symbol=G](-.5,.2)(.5,.2)\endpspicture \\
	
	H &  \pspicture[shift=*](-1,-.5)(1,.5) \psline[linestyle=symbol,symbol=H](-.5,.2)(.5,.2)\endpspicture &
	I &  \pspicture[shift=*](-1,-.5)(1,.5) \psline[linestyle=symbol,symbol=I](-.5,.2)(.5,.2)\endpspicture &
	J &  \pspicture[shift=*](-1,-.5)(1,.5) \psline[linestyle=symbol,symbol=J](-.5,.2)(.5,.2)\endpspicture &
	K &  \pspicture[shift=*](-1,-.5)(1,.5) \psline[linestyle=symbol,symbol=K](-.5,.2)(.5,.2)\endpspicture &
	L &  \pspicture[shift=*](-1,-.5)(1,.5) \psline[linestyle=symbol,symbol=L](-.5,.2)(.5,.2)\endpspicture &
	M &  \pspicture[shift=*](-1,-.5)(1,.5) \psline[linestyle=symbol,symbol=M](-.5,.2)(.5,.2)\endpspicture &
	N &  \pspicture[shift=*](-1,-.5)(1,.5) \psline[linestyle=symbol,symbol=N](-.5,.2)(.5,.2)\endpspicture \\
	
	O &  \pspicture[shift=*](-1,-.5)(1,.5) \psline[linestyle=symbol,symbol=O](-.5,.2)(.5,.2)\endpspicture &
	P &  \pspicture[shift=*](-1,-.5)(1,.5) \psline[linestyle=symbol,symbol=P](-.5,.2)(.5,.2)\endpspicture &
	Q &  \pspicture[shift=*](-1,-.5)(1,.5) \psline[linestyle=symbol,symbol=Q](-.5,.2)(.5,.2)\endpspicture &
	R &  \pspicture[shift=*](-1,-.5)(1,.5) \psline[linestyle=symbol,symbol=R](-.5,.2)(.5,.2)\endpspicture &
	S &  \pspicture[shift=*](-1,-.5)(1,.5) \psline[linestyle=symbol,symbol=S](-.5,.2)(.5,.2)\endpspicture &
	T &  \pspicture[shift=*](-1,-.5)(1,.5) \psline[linestyle=symbol,symbol=T](-.5,.2)(.5,.2)\endpspicture &
	U &  \pspicture[shift=*](-1,-.5)(1,.5) \psline[linestyle=symbol,symbol=U](-.5,.2)(.5,.2)\endpspicture \\
	
	V &  \pspicture[shift=*](-1,-.5)(1,.5) \psline[linestyle=symbol,symbol=V](-.5,.2)(.5,.2)\endpspicture &
	W &  \pspicture[shift=*](-1,-.5)(1,.5) \psline[linestyle=symbol,symbol=W](-.5,.2)(.5,.2)\endpspicture &
	X &  \pspicture[shift=*](-1,-.5)(1,.5) \psline[linestyle=symbol,symbol=X](-.5,.2)(.5,.2)\endpspicture &
	Y &  \pspicture[shift=*](-1,-.5)(1,.5) \psline[linestyle=symbol,symbol=Y](-.5,.2)(.5,.2)\endpspicture &
	Z &  \pspicture[shift=*](-1,-.5)(1,.5) \psline[linestyle=symbol,symbol=Z](-.5,.2)(.5,.2)\endpspicture &
	1 &  \pspicture[shift=*](-1,-.5)(1,.5) \psline[linestyle=symbol,symbol=1](-.5,.2)(.5,.2)\endpspicture &
	2 &  \pspicture[shift=*](-1,-.5)(1,.5) \psline[linestyle=symbol,symbol=2](-.5,.2)(.5,.2)\endpspicture \\	
	
	
	3 &  \pspicture[shift=*](-1,-.5)(1,.5) \psline[linestyle=symbol,symbol=3](-.5,.2)(.5,.2)\endpspicture &
	4 &  \pspicture[shift=*](-1,-.5)(1,.5) \psline[linestyle=symbol,symbol=4](-.5,.2)(.5,.2)\endpspicture &
	5 &  \pspicture[shift=*](-1,-.5)(1,.5) \psline[linestyle=symbol,symbol=5](-.5,.2)(.5,.2)\endpspicture &
	6 &  \pspicture[shift=*](-1,-.5)(1,.5) \psline[linestyle=symbol,symbol=6](-.5,.2)(.5,.2)\endpspicture &
	7 &  \pspicture[shift=*](-1,-.5)(1,.5) \psline[linestyle=symbol,symbol=7](-.5,.2)(.5,.2)\endpspicture &
	8 &  \pspicture[shift=*](-1,-.5)(1,.5) \psline[linestyle=symbol,symbol=8](-.5,.2)(.5,.2)\endpspicture &
	9 &  \pspicture[shift=*](-1,-.5)(1,.5) \psline[linestyle=symbol,symbol=9](-.5,.2)(.5,.2)\endpspicture \\
	
	
	a &  \pspicture[shift=*](-1,-.5)(1,.5) \psline[linestyle=symbol,symbol=a](-.5,.2)(.5,.2)\endpspicture &
	b &  \pspicture[shift=*](-1,-.5)(1,.5) \psline[linestyle=symbol,symbol=b](-.5,.2)(.5,.2)\endpspicture &
	c &  \pspicture[shift=*](-1,-.5)(1,.5) \psline[linestyle=symbol,symbol=c](-.5,.2)(.5,.2)\endpspicture &
	d &  \pspicture[shift=*](-1,-.5)(1,.5) \psline[linestyle=symbol,symbol=d](-.5,.2)(.5,.2)\endpspicture &
	e &  \pspicture[shift=*](-1,-.5)(1,.5) \psline[linestyle=symbol,symbol=e](-.5,.2)(.5,.2)\endpspicture &
	f &  \pspicture[shift=*](-1,-.5)(1,.5) \psline[linestyle=symbol,symbol=f](-.5,.2)(.5,.2)\endpspicture &
	g &  \pspicture[shift=*](-1,-.5)(1,.5) \psline[linestyle=symbol,symbol=g](-.5,.2)(.5,.2)\endpspicture \\
		
	h &  \pspicture[shift=*](-1,-.5)(1,.5) \psline[linestyle=symbol,symbol=h](-.5,.2)(.5,.2)\endpspicture &
	i &  \pspicture[shift=*](-1,-.5)(1,.5) \psline[linestyle=symbol,symbol=i](-.5,.2)(.5,.2)\endpspicture &
	j &  \pspicture[shift=*](-1,-.5)(1,.5) \psline[linestyle=symbol,symbol=j](-.5,.2)(.5,.2)\endpspicture &
	k &  \pspicture[shift=*](-1,-.5)(1,.5) \psline[linestyle=symbol,symbol=k](-.5,.2)(.5,.2)\endpspicture &
	l &  \pspicture[shift=*](-1,-.5)(1,.5) \psline[linestyle=symbol,symbol=l](-.5,.2)(.5,.2)\endpspicture &
	m &  \pspicture[shift=*](-1,-.5)(1,.5) \psline[linestyle=symbol,symbol=m](-.5,.2)(.5,.2)\endpspicture &
	n &  \pspicture[shift=*](-1,-.5)(1,.5) \psline[linestyle=symbol,symbol=n](-.5,.2)(.5,.2)\endpspicture \\
		
	o &  \pspicture[shift=*](-1,-.5)(1,.5) \psline[linestyle=symbol,symbol=o](-.5,.2)(.5,.2)\endpspicture &
	p &  \pspicture[shift=*](-1,-.5)(1,.5) \psline[linestyle=symbol,symbol=p](-.5,.2)(.5,.2)\endpspicture &
	q &  \pspicture[shift=*](-1,-.5)(1,.5) \psline[linestyle=symbol,symbol=q](-.5,.2)(.5,.2)\endpspicture &
	r &  \pspicture[shift=*](-1,-.5)(1,.5) \psline[linestyle=symbol,symbol=r](-.5,.2)(.5,.2)\endpspicture &
	s &  \pspicture[shift=*](-1,-.5)(1,.5) \psline[linestyle=symbol,symbol=s](-.5,.2)(.5,.2)\endpspicture &
	t &  \pspicture[shift=*](-1,-.5)(1,.5) \psline[linestyle=symbol,symbol=t](-.5,.2)(.5,.2)\endpspicture &
	u &  \pspicture[shift=*](-1,-.5)(1,.5) \psline[linestyle=symbol,symbol=u](-.5,.2)(.5,.2)\endpspicture \\
		
	v &  \pspicture[shift=*](-1,-.5)(1,.5) \psline[linestyle=symbol,symbol=v](-.5,.2)(.5,.2)\endpspicture &
	w &  \pspicture[shift=*](-1,-.5)(1,.5) \psline[linestyle=symbol,symbol=w](-.5,.2)(.5,.2)\endpspicture &
	x &  \pspicture[shift=*](-1,-.5)(1,.5) \psline[linestyle=symbol,symbol=x](-.5,.2)(.5,.2)\endpspicture &
	y &  \pspicture[shift=*](-1,-.5)(1,.5) \psline[linestyle=symbol,symbol=y](-.5,.2)(.5,.2)\endpspicture &
	z &  \pspicture[shift=*](-1,-.5)(1,.5) \psline[linestyle=symbol,symbol=z](-.5,.2)(.5,.2)\endpspicture &
	+ &  \pspicture[shift=*](-1,-.5)(1,.5) \psline[linestyle=symbol,symbol=+](-.5,.2)(.5,.2)\endpspicture &
	- &  \pspicture[shift=*](-1,-.5)(1,.5) \psline[linestyle=symbol,symbol=-](-.5,.2)(.5,.2)\endpspicture \\	
		
	* &  \pspicture[shift=*](-1,-.5)(1,.5) \psline[linestyle=symbol,symbol=*](-.5,.2)(.5,.2)\endpspicture &
	' &  \pspicture[shift=*](-1,-.5)(1,.5) \psline[linestyle=symbol,symbol='](-.5,.2)(.5,.2)\endpspicture &	
	> &  \pspicture[shift=*](-1,-.5)(1,.5) \psline[linestyle=symbol,symbol=>](-.5,.2)(.5,.2)\endpspicture &
	< &  \pspicture[shift=*](-1,-.5)(1,.5) \psline[linestyle=symbol,symbol=<](-.5,.2)(.5,.2)\endpspicture &
	0 &  \pspicture[shift=*](-1,-.5)(1,.5) \psline[linestyle=symbol,symbol=0](-.5,.2)(.5,.2)\endpspicture &
	/ &  \pspicture[shift=*](-1,-.5)(1,.5) \psline[linestyle=symbol,symbol=/](-.5,.2)(.5,.2)\endpspicture &
	. &  \pspicture[shift=*](-1,-.5)(1,.5) \psline[linestyle=symbol,symbol=.](-.5,.2)(.5,.2)\endpspicture \\ \hline
	\end{tabular}
\end{center}




\psset[pstricks]{symbolFont=PSTricksDotFont}
\begin{center}
	\begin{tabular}{| c @{:} c |c @{:} c |c @{:} c |c @{:} c |c @{:} c |c @{:} c |c @{:} c | }  \hline
\multicolumn{14}{|c|}{symbolFont=PSTricksDotFont} \\  \hline
	A &  \pspicture[shift=*](-1,-.5)(1,.5) \psline[linestyle=symbol,symbol=A](-.5,.2)(.5,.2)\endpspicture &
	B &  \pspicture[shift=*](-1,-.5)(1,.5) \psline[linestyle=symbol,symbol=B](-.5,.2)(.5,.2)\endpspicture &
	C &  \pspicture[shift=*](-1,-.5)(1,.5) \psline[linestyle=symbol,symbol=C](-.5,.2)(.5,.2)\endpspicture &
	D &  \pspicture[shift=*](-1,-.5)(1,.5) \psline[linestyle=symbol,symbol=D](-.5,.2)(.5,.2)\endpspicture &
	E &  \pspicture[shift=*](-1,-.5)(1,.5) \psline[linestyle=symbol,symbol=E](-.5,.2)(.5,.2)\endpspicture &
	F &  \pspicture[shift=*](-1,-.5)(1,.5) \psline[linestyle=symbol,symbol=F](-.5,.2)(.5,.2)\endpspicture &
	G &  \pspicture[shift=*](-1,-.5)(1,.5) \psline[linestyle=symbol,symbol=G](-.5,.2)(.5,.2)\endpspicture \\
	
	H &  \pspicture[shift=*](-1,-.5)(1,.5) \psline[linestyle=symbol,symbol=H](-.5,.2)(.5,.2)\endpspicture &
	I &  \pspicture[shift=*](-1,-.5)(1,.5) \psline[linestyle=symbol,symbol=I](-.5,.2)(.5,.2)\endpspicture &
	J &  \pspicture[shift=*](-1,-.5)(1,.5) \psline[linestyle=symbol,symbol=J](-.5,.2)(.5,.2)\endpspicture &
	K &  \pspicture[shift=*](-1,-.5)(1,.5) \psline[linestyle=symbol,symbol=K](-.5,.2)(.5,.2)\endpspicture &
	L &  \pspicture[shift=*](-1,-.5)(1,.5) \psline[linestyle=symbol,symbol=L](-.5,.2)(.5,.2)\endpspicture &
	M &  \pspicture[shift=*](-1,-.5)(1,.5) \psline[linestyle=symbol,symbol=M](-.5,.2)(.5,.2)\endpspicture &
	N &  \pspicture[shift=*](-1,-.5)(1,.5) \psline[linestyle=symbol,symbol=N](-.5,.2)(.5,.2)\endpspicture \\
	
	O &  \pspicture[shift=*](-1,-.5)(1,.5) \psline[linestyle=symbol,symbol=O](-.5,.2)(.5,.2)\endpspicture &
	P &  \pspicture[shift=*](-1,-.5)(1,.5) \psline[linestyle=symbol,symbol=P](-.5,.2)(.5,.2)\endpspicture &
	Q &  \pspicture[shift=*](-1,-.5)(1,.5) \psline[linestyle=symbol,symbol=Q](-.5,.2)(.5,.2)\endpspicture &
	R &  \pspicture[shift=*](-1,-.5)(1,.5) \psline[linestyle=symbol,symbol=R](-.5,.2)(.5,.2)\endpspicture &
	S &  \pspicture[shift=*](-1,-.5)(1,.5) \psline[linestyle=symbol,symbol=S](-.5,.2)(.5,.2)\endpspicture &
	T &  \pspicture[shift=*](-1,-.5)(1,.5) \psline[linestyle=symbol,symbol=T](-.5,.2)(.5,.2)\endpspicture &
	U &  \pspicture[shift=*](-1,-.5)(1,.5) \psline[linestyle=symbol,symbol=U](-.5,.2)(.5,.2)\endpspicture \\
	
	V &  \pspicture[shift=*](-1,-.5)(1,.5) \psline[linestyle=symbol,symbol=V](-.5,.2)(.5,.2)\endpspicture &
	W &  \pspicture[shift=*](-1,-.5)(1,.5) \psline[linestyle=symbol,symbol=W](-.5,.2)(.5,.2)\endpspicture &
	X &  \pspicture[shift=*](-1,-.5)(1,.5) \psline[linestyle=symbol,symbol=X](-.5,.2)(.5,.2)\endpspicture &
	Y &  \pspicture[shift=*](-1,-.5)(1,.5) \psline[linestyle=symbol,symbol=Y](-.5,.2)(.5,.2)\endpspicture &
	Z &  \pspicture[shift=*](-1,-.5)(1,.5) \psline[linestyle=symbol,symbol=Z](-.5,.2)(.5,.2)\endpspicture &
	1 &  \pspicture[shift=*](-1,-.5)(1,.5) \psline[linestyle=symbol,symbol=1](-.5,.2)(.5,.2)\endpspicture &
	2 &  \pspicture[shift=*](-1,-.5)(1,.5) \psline[linestyle=symbol,symbol=2](-.5,.2)(.5,.2)\endpspicture \\	
	
	
	3 &  \pspicture[shift=*](-1,-.5)(1,.5) \psline[linestyle=symbol,symbol=3](-.5,.2)(.5,.2)\endpspicture &
	4 &  \pspicture[shift=*](-1,-.5)(1,.5) \psline[linestyle=symbol,symbol=4](-.5,.2)(.5,.2)\endpspicture &
	5 &  \pspicture[shift=*](-1,-.5)(1,.5) \psline[linestyle=symbol,symbol=5](-.5,.2)(.5,.2)\endpspicture &
	6 &  \pspicture[shift=*](-1,-.5)(1,.5) \psline[linestyle=symbol,symbol=6](-.5,.2)(.5,.2)\endpspicture &
	7 &  \pspicture[shift=*](-1,-.5)(1,.5) \psline[linestyle=symbol,symbol=7](-.5,.2)(.5,.2)\endpspicture &
	8 &  \pspicture[shift=*](-1,-.5)(1,.5) \psline[linestyle=symbol,symbol=8](-.5,.2)(.5,.2)\endpspicture &
	9 &  \pspicture[shift=*](-1,-.5)(1,.5) \psline[linestyle=symbol,symbol=9](-.5,.2)(.5,.2)\endpspicture \\
	
	
	a &  \pspicture[shift=*](-1,-.5)(1,.5) \psline[linestyle=symbol,symbol=a](-.5,.2)(.5,.2)\endpspicture &
	b &  \pspicture[shift=*](-1,-.5)(1,.5) \psline[linestyle=symbol,symbol=b](-.5,.2)(.5,.2)\endpspicture &
	c &  \pspicture[shift=*](-1,-.5)(1,.5) \psline[linestyle=symbol,symbol=c](-.5,.2)(.5,.2)\endpspicture &
	d &  \pspicture[shift=*](-1,-.5)(1,.5) \psline[linestyle=symbol,symbol=d](-.5,.2)(.5,.2)\endpspicture &
	e &  \pspicture[shift=*](-1,-.5)(1,.5) \psline[linestyle=symbol,symbol=e](-.5,.2)(.5,.2)\endpspicture &
	f &  \pspicture[shift=*](-1,-.5)(1,.5) \psline[linestyle=symbol,symbol=f](-.5,.2)(.5,.2)\endpspicture &
	g &  \pspicture[shift=*](-1,-.5)(1,.5) \psline[linestyle=symbol,symbol=g](-.5,.2)(.5,.2)\endpspicture \\
		
	h &  \pspicture[shift=*](-1,-.5)(1,.5) \psline[linestyle=symbol,symbol=h](-.5,.2)(.5,.2)\endpspicture &
	i &  \pspicture[shift=*](-1,-.5)(1,.5) \psline[linestyle=symbol,symbol=i](-.5,.2)(.5,.2)\endpspicture &
	j &  \pspicture[shift=*](-1,-.5)(1,.5) \psline[linestyle=symbol,symbol=j](-.5,.2)(.5,.2)\endpspicture &
	k &  \pspicture[shift=*](-1,-.5)(1,.5) \psline[linestyle=symbol,symbol=k](-.5,.2)(.5,.2)\endpspicture &
	l &  \pspicture[shift=*](-1,-.5)(1,.5) \psline[linestyle=symbol,symbol=l](-.5,.2)(.5,.2)\endpspicture &
	m &  \pspicture[shift=*](-1,-.5)(1,.5) \psline[linestyle=symbol,symbol=m](-.5,.2)(.5,.2)\endpspicture &
	n &  \pspicture[shift=*](-1,-.5)(1,.5) \psline[linestyle=symbol,symbol=n](-.5,.2)(.5,.2)\endpspicture \\
		
	o &  \pspicture[shift=*](-1,-.5)(1,.5) \psline[linestyle=symbol,symbol=o](-.5,.2)(.5,.2)\endpspicture &
	p &  \pspicture[shift=*](-1,-.5)(1,.5) \psline[linestyle=symbol,symbol=p](-.5,.2)(.5,.2)\endpspicture &
	q &  \pspicture[shift=*](-1,-.5)(1,.5) \psline[linestyle=symbol,symbol=q](-.5,.2)(.5,.2)\endpspicture &
	r &  \pspicture[shift=*](-1,-.5)(1,.5) \psline[linestyle=symbol,symbol=r](-.5,.2)(.5,.2)\endpspicture &
	s &  \pspicture[shift=*](-1,-.5)(1,.5) \psline[linestyle=symbol,symbol=s](-.5,.2)(.5,.2)\endpspicture &
	t &  \pspicture[shift=*](-1,-.5)(1,.5) \psline[linestyle=symbol,symbol=t](-.5,.2)(.5,.2)\endpspicture &
	u &  \pspicture[shift=*](-1,-.5)(1,.5) \psline[linestyle=symbol,symbol=u](-.5,.2)(.5,.2)\endpspicture \\
		
	v &  \pspicture[shift=*](-1,-.5)(1,.5) \psline[linestyle=symbol,symbol=v](-.5,.2)(.5,.2)\endpspicture &
	w &  \pspicture[shift=*](-1,-.5)(1,.5) \psline[linestyle=symbol,symbol=w](-.5,.2)(.5,.2)\endpspicture &
	x &  \pspicture[shift=*](-1,-.5)(1,.5) \psline[linestyle=symbol,symbol=x](-.5,.2)(.5,.2)\endpspicture &
	y &  \pspicture[shift=*](-1,-.5)(1,.5) \psline[linestyle=symbol,symbol=y](-.5,.2)(.5,.2)\endpspicture &
	z &  \pspicture[shift=*](-1,-.5)(1,.5) \psline[linestyle=symbol,symbol=z](-.5,.2)(.5,.2)\endpspicture &
	+ &  \pspicture[shift=*](-1,-.5)(1,.5) \psline[linestyle=symbol,symbol=+](-.5,.2)(.5,.2)\endpspicture &
	- &  \pspicture[shift=*](-1,-.5)(1,.5) \psline[linestyle=symbol,symbol=-](-.5,.2)(.5,.2)\endpspicture \\	
		
	* &  \pspicture[shift=*](-1,-.5)(1,.5) \psline[linestyle=symbol,symbol=*](-.5,.2)(.5,.2)\endpspicture &
	' &  \pspicture[shift=*](-1,-.5)(1,.5) \psline[linestyle=symbol,symbol='](-.5,.2)(.5,.2)\endpspicture &	
	> &  \pspicture[shift=*](-1,-.5)(1,.5) \psline[linestyle=symbol,symbol=>](-.5,.2)(.5,.2)\endpspicture &
	< &  \pspicture[shift=*](-1,-.5)(1,.5) \psline[linestyle=symbol,symbol=<](-.5,.2)(.5,.2)\endpspicture &
	0 &  \pspicture[shift=*](-1,-.5)(1,.5) \psline[linestyle=symbol,symbol=0](-.5,.2)(.5,.2)\endpspicture &
	/ &  \pspicture[shift=*](-1,-.5)(1,.5) \psline[linestyle=symbol,symbol=/](-.5,.2)(.5,.2)\endpspicture &
	. &  \pspicture[shift=*](-1,-.5)(1,.5) \psline[linestyle=symbol,symbol=.](-.5,.2)(.5,.2)\endpspicture \\ \hline
	\end{tabular}
\end{center}

\newpage
%\subsection[Les bobines]{Les bobines \cite{pst-user} \cite{pst-coil}}
\SbSSCT{Les bobines}{Coils}


\TFRGB{Utilisation du module}{Package} \textbf{pst-coil} \label{pst-coil}

\SbSbSSCT{Les 3 types de bobines}{The 3 types of coils}
\psset{fillstyle=none,unit=1cm}

\begin{center}
\begin{tabular}{|p{4cm}|p{4cm}|p{4cm}|}  \hline  
\rule[-1.25cm]{0pt}{2.5cm}	\pscoil(0.5,0)(4,0)  & \pszigzag(0.5,0)(4,0) & \pssin(0.5,0)(4,0) \\ \hline
\BSS{pscoil}(0.5,0)(4,0) \BSI{pscoil}{pst-coil} & \BSS{pszigzag}(0.5,0)(4,0)  \BSI{pszigzag}{pst-coil} & \BSS{pssin}(0.5,0)(4,0)   \BSI{pssin}{pst-coil}
\\ \hline	
\end{tabular}
\end{center}

\begin{center}
\begin{tabular}{|p{4cm}|p{4cm}|p{4cm}|}  \hline  
\rule[-1.25cm]{0pt}{2.5cm}	\pscoil*(0.5,0)(4,0)  & \pszigzag*(0.5,0)(4,0) & \pssin*(0.5,0)(4,0) \\ \hline
\textbf{\BS{}pscoil*}(0.5,0)(4,0)  & \textbf{\BS{}pszigzag*}(0.5,0)(4,0)  & \textbf{\BS{}pssin*}(0.5,0)(4,0)  \\ \hline	
\end{tabular}
\end{center}

\SbSbSSCT{Paramètres des bobines}{Parameters of coils}
\begin{center}
\begin{tabular}{|p{6cm}|p{6cm}|}  \hline  
\rule[-.7cm]{0pt}{1.4cm}	\pscoil[coilwidth=0.5cm](0.5,0)(5,0)  & \pszigzag[coilwidth=0.5cm](0.5,0)(5,0)  \\ \hline
\BS{pscoil}[\RDD{coilwidth}=0.5cm](0.5,0)(5,0)  \RDI{coilwidth}{pst-coil} & \BS{pszigzag}[{\red coilwidth=0.5cm}](0.5,0)(5,0)  \\ \hline
\multicolumn{2}{|c|}{\blue \dft: 1cm}\\ \hline 	
\end{tabular}
\end{center}

\begin{center}
\begin{tabular}{|p{6cm}|p{6cm}|}  \hline  
\rule[-.7cm]{0pt}{1.4cm}	\pscoil[coilheight=0.5](0.5,0)(5,0)  & \pszigzag[coilheight=0.5](0.5,0)(5,0)  \\ \hline
\BS{pscoil}[\RDD{coilheight}=0.5](0.5,0)(5,0) \RDI{coilheight}{pst-coil} &
 \BS{}pszigzag[{\red coilheight0.5}](0.5,0)(5,0)  \\ \hline	
\multicolumn{2}{|c|}{\blue  \dft : 1}\\ \hline 	
\end{tabular}
\end{center}

\begin{center}
\begin{tabular}{|p{6cm}|p{6cm}|}  \hline  
\rule[-.7cm]{0pt}{1.4cm}	\pscoil[coilarm=1](0.5,0)(5,0)  & \pszigzag[coilarm=1](0.5,0)(5,0)  \\ \hline
\BS{}pscoil[\RDD{coilarm}=1](0.5,0)(5,0) \RDI{coilarm}{pst-coil} & \BS{}pszigzag[{\red coilarm=1}](0.5,0)(5,0)  \\ \hline	

\rule[-.7cm]{0pt}{1.4cm}	\pscoil[coilarmA=1](0.5,0)(5,0)  & \pszigzag[coilarmB=1](0.5,0)(5,0)  \\ \hline
\BS{}pscoil[\RDD{coilarmA}=1](0.5,0)(5,0) \RDI{coilarmA}{pst-coil} & \BS{}pszigzag[\RDD{coilarmB}=1](0.5,0)(5,0)  \RDI{coilarmB}{pst-coil} \\ \hline
\multicolumn{2}{|c|}{\blue  \dft : 0.5cm}\\ \hline 
\end{tabular}
\end{center}

\begin{center}
\begin{tabular}{|p{6cm}|p{6cm}|}  \hline  
\rule[-.7cm]{0pt}{1.4cm}	\pscoil[coilaspect=0](0.5,0)(5,0)  & \pscoil[coilaspect=30](0.5,0)(5,0)  \\ \hline
\BS{}pscoil[\RDD{coilaspect}=0](0.5,0)(5,0) \RDI{coilaspect}{pst-coil} & \BS{}pscoil[{\red coilaspect=30}](0.5,0)(5,0)  \\ \hline	
\multicolumn{2}{|c|}{ \blue \dft :45}\\ \hline 
\end{tabular}
\end{center}

\begin{center}
\begin{tabular}{|p{6cm}|p{6cm}|}  \hline  
\rule[-.7cm]{0pt}{1.4cm}	\pscoil[coilinc=1](0.5,0)(5,0)  & \pscoil[coilinc=60](0.5,0)(5,0)  \\ \hline
\BS{}pscoil[\RDD{coilinc}=1](0.5,0)(5,0) \RDI{coilinc}{pst-coil}  & \BS{}pscoil[{\red  coilinc=30 }](0.5,0)(5,0)  \\ \hline
\multicolumn{2}{|c|}{\blue \dft : 10}\\ \hline 	
\end{tabular}
\end{center}

\begin{center}
\begin{tabular}{|p{6cm}|p{6cm}|}  \hline  
\rule[-2cm]{0pt}{4cm}	\pszigzag[bow=1cm](0.5,0)(5,0)  & \pszigzag[bow=-1cm](0.5,0)(5,0)  \\ \hline
\BS{pszigzag}[\RDD{bow}=1cm](0.5,0)(5,0) \RDI{bow}{pst-coil}  & 
\BS{}pszigzag[{\red bow=-1cm}](0.5,0)(5,0)  \\ \hline
\multicolumn{2}{|c|}{\blue \dft : 0}\\ \hline 	
\end{tabular}
\end{center}


\begin{center}
\begin{tabular}{|p{6cm}|p{6cm}|}  \hline  
\rule[-1.5cm]{0pt}{3cm}	\pssin(0.5,0)(5,0)  & \pssin[periods=5](0.5,0)(5,0)  \\ \hline
\textbf{\BS{}pssin}(0.5,0)(5,0)  & \BS{}pssin[\RDD{periods}=5(0.5,0)(5,0) \RDI{periods}{pst-coil}  \\ \hline

\rule[-1.5cm]{0pt}{3cm}	\pssin[periods=2cm ](0.5,0)(5,0)  & \pssin[amplitude=0.5](0.5,0)(5,0)  \\ \hline
\BS{}pssin[{\red periods=2cm}](0.5,0)(5,0)  & \BS{}pssin[\RDD{amplitude}=0.5](0.5,0)(5,0)  \RDI{amplitude}{pst-coil} \\ \hline

\rule[-1.5cm]{0pt}{3cm}	\pssin[ppoints=5 ](0.5,0)(5,0)  & \pssin[ppoints=2000](0.5,0)(5,0)  \\ \hline
\BS{}pssin[\RDD{ppoints}=5 ](0.5,0)(5,0) \RDI{ppoints}{pst-coil}  & 
\BS{}pssin[{\red ppoints=2000}]](0.5,0)(5,0)  \\ \hline
\multicolumn{2}{|c|}{\blue \dft:  periods = 1 , amplitude = 1 , ppoints= 360}\\ \hline 	
\end{tabular}
\end{center}

\begin{center}
\begin{tabular}{|l|}  \hline 

\multicolumn{1}{|c|}{
\begin{pspicture}(0,-4)(11,2)
\psaxes[yticksize=11,linewidth=.2mm,linecolor=black,linestyle=dashed,xticksize=-2.5 1.5]{->}(0,0)(0,-3.5)(11,1.5)
	\pssin[linecolor=red,linewidth=2pt,function=dup sin mul .01 mul](1,0)(10,0) 
	\pssin[linecolor=green,linewidth=2pt,function=cos](1,0)(10,0) 
	\pssin[linecolor=cyan,linewidth=2pt,function=2 mul sin](1,0)(10,0)
	\pssin[linecolor=blue,linewidth=2pt,function=360 div](1,0)(10,0)	
\end{pspicture}}\\\hline
	\BS{}pssin[linecolor=blue,linewidth=2pt,\RDD{function=360} div](1,0)(10,0) \RDI{function}{pst-coil} \\
	\BS{}pssin[linecolor=green,linewidth=2pt,{\green function=cos}](1,0)(10,0)  \\ 
	\BS{}pssin[linecolor=cyan,linewidth=2pt,{\cyan function=2 mul sin}](1,0)(10,0) \\ 
	\BS{}pssin[linecolor=red,linewidth=2pt,{\red function=dup sin mul .01 mul}](1,0)(10,0)  \\	\hline
\end{tabular}
\end{center}


\SbSbSSCT{Liaison de n\oe uds en bobine}{Connecting nodes with coils}

\begin{center}
\begin{tabular}{|c|c|c|} \hline
\multicolumn{3}{|c|}{\BS{}dotnode[dotstyle=*](.5,-.5)\AC{{\red A}} \BS{}dotnode[dotstyle=*](3.5,0)\AC{{\red B}}} \\ \hline
\begin{pspicture}(0,-1)(4 ,1) \dotnode[dotstyle=*](.5,-.5){A} \dotnode[dotstyle=*](3.5,0){B} \nccoil{->}{A}{B}  
\end{pspicture}	
&
\begin{pspicture}(0,-1)(4 ,1)\dotnode[dotstyle=*](.5,-.5){A} \dotnode[dotstyle=*](3.5,0){B} \nczigzag{->}{A}{B}  
\end{pspicture}	
&  
\begin{pspicture}(0,-1)(4 ,1)\dotnode[dotstyle=*](.5,-.5){A} \dotnode[dotstyle=*](3.5,0){B} \ncsin{->}{A}{B}  
\end{pspicture}	
\\ \hline
\BSS{nccoil}\AC{->}\AC{A}\AC{B}  \BSI{nccoil}{pst-coil} & \BSS{nczigzag}\AC{->}\AC{A}\AC{B}  \BSI{nczigzag}{pst-coil} & \BSS{ncsin}\AC{->}\AC{A}\AC{B}   \BSI{ncsin}{pst-coil}
\\ \hline
\begin{pspicture}(0,-1)(4 ,1) \dotnode[dotstyle=*](.5,-.5){A} \dotnode[dotstyle=*](3.5,0){B} \pccoil{->}(A)(B)   
\end{pspicture}	
&
\begin{pspicture}(0,-1)(4 ,1)\dotnode[dotstyle=*](.5,-.5){A} \dotnode[dotstyle=*](3.5,0){B}
 \pczigzag {->}(A)(B)  
\end{pspicture}	
&  
\begin{pspicture}(0,-1)(4 ,1)\dotnode[dotstyle=*](.5,-.5){A} \dotnode[dotstyle=*](3.5,0){B}
 \pcsin{->}(A)(B)  
\end{pspicture}
\\ \hline
\BSS{pccoil}\AC{->}(A)(B)  \BSI{pccoil}{pst-coil} & \BSS{pczigzag} \AC{->}(A)(B)  \BSI{pczigzag}{pst-coil} &  \BSS{pcsin}\AC{->}(A)(B)   \BSI{pcsin}{pst-coil}\\ \hline
\begin{pspicture}(0,-1)(4 ,1) \dotnode[dotstyle=*](.5,-.5){A} \dotnode[dotstyle=*](3.5,0){B} \nccoil*{->}{A}{B}  
\end{pspicture}	
&
\begin{pspicture}(0,-1)(4 ,1)\dotnode[dotstyle=*](.5,-.5){A} \dotnode[dotstyle=*](3.5,0){B} \nczigzag*{->}{A}{B}  
\end{pspicture}	
&  
\begin{pspicture}(0,-1)(4 ,1)\dotnode[dotstyle=*](.5,-.5){A} \dotnode[dotstyle=*](3.5,0){B} \ncsin*{->}{A}{B}  
\end{pspicture}	
\\ \hline
\textbf{\BS{}nccoil*}\AC{->}\AC{A}\AC{B}  & \textbf{\BS{}nczigzag*}\AC{->}\AC{A}\AC{B}  & \textbf{\BS{}ncsin*}\AC{->}\AC{A}\AC{B}  
\\ \hline
\end{tabular}
\end{center}

\psset{fillcolor=yellow}  % pst-coil

\newpage

%\subsection[Les accolades]{Les accolades  \cite{pstricks-add}}
\SbSSCT{Les accolades}{Braces}


%Syntaxe : 
%\BSS{psbrace}[Options](A)(B)\AC{\TFRGB{texte}{text}}

\SbSbSSCT{Dans un environnement pspicture}{Braces in  pspicture}


\begin{tabular}{|c|c|} \hline  
 \begin{pspicture}[shift=*](-.5,-0.5)(5,4.5)
\psaxes[xticksize=4,yticksize=4,axesstyle=frame](5,4)
\pnode(1,1){B}\pnode(3,2){C}
\psbrace(B)(C){\TFRGB{texte}{text}}
\psbrace(4,3)(1,3){}
\psdots[dotstyle=*,linecolor=blue,dotscale=2](B) \nput{45}{B}{B} 
\psdots[dotstyle=*,linecolor=blue,dotscale=2](C) \nput{45}{C}{C}%
 \end{pspicture} 
& 
\parbox{3.5cm}{ 
\BS{}pnode(1,1)\AC{B} \BS{}pnode(3,2)\AC{C}

\BSS{psbrace}(B)(C)\AC{\TFRGB{texte}{text}} \BSI{psbrace}{pstricks-add} 

\bigskip
\textbf{\BS{psbrace}}(4,3)(1,3)\AC{} }
\\ \hline 
\end{tabular} 

%\begin{center}
%
%\end{center}
 
%\bigskip

 
% \psbrace[bracePos=0.25,nodesepB=10 pt,rot=90](0,2)(\linewidth,2){\fbox{Text}}
\bigskip
\begin{tabular}{|c|c|}\hline 
\multicolumn{2}{|c|}{\TFRGB{Rôle de l'astérisque}{Role of the asterisk} \footnotemark[1] } \\
\hline 
 \begin{pspicture}(-.5,-.5)(4,3)
 \psaxes[xticksize=2.5,yticksize=4,axesstyle=frame](4,2.5)
 \psbrace[braceWidth=.5cm,fillcolor=yellow](1,1)(3,2){}
 \end{pspicture}  &  
 \begin{pspicture}(-.5,-.5)(4,3)
 \psaxes[xticksize=2.5,yticksize=5,axesstyle=frame](4,2.5)
 \psbrace*[braceWidth=.5cm,fillcolor=yellow](1,1)(3,2){}
\end{pspicture} \\ \hline 
 \BS{}psbrace(1,1)(3,2)\AC{} &
  \BS{}psbrace*(1,1)(3,2)\AC{}  \\
\hline 
\end{tabular} 

\footnotetext[1]{braceWidth=.5cm,fillcolor=yellow}

\SbSbSSCT{Dans le texte}{The brace in the text}
 \TFRGB{le noeud  A est ici \pnode(0,0){A} et le noeud B est ici}{the node  A is here \pnode(0,0){A} and the node  B is here} \pnode(0,0){B} \BS{}psbrace(A)(B)\AC{\TFRGB{texte}{text}}
\psbrace(A)(B){\TFRGB{texte}{text}}
\smallskip

{\red \TFRGB{L'accolage n'a pas de dimension}{The brace has no dimension !}}

\vspace{1cm}


\pnode(0,0){A} \TFRGB{ici, se trouve le noeud}{here is the node}  A 

\BS{}vspace\AC{1cm}
\vspace{1cm}

 \pnode(0,0){B} \TFRGB{ici, se trouve le noeud}{here is the node}  B   \BS{}psbrace(A)(B)\AC{}
\psbrace[rot=90](A)(B){}

\vspace{1cm}




\subsubsection{Options}
\begin{tabular}{|c|c|c|}
\hline 
 \begin{pspicture}(-.5,-.5)(4.2,2.5)
\psaxes[xticksize=2.2,yticksize=4,axesstyle=frame](4,2.2)
\pnode(1,1){B}\pnode(3,2){C}
\psbrace[braceWidth=5pt](B)(C){\TFRGB{texte}{text}}
 \end{pspicture} 
&
 \begin{pspicture}(-.5,-.5)(4.2,2.5)
\psaxes[xticksize=2.2,yticksize=4,axesstyle=frame](4,2.2)
\pnode(1,1){B}\pnode(3,2){C}
\psbrace[braceWidthInner=.5cm](B)(C){\TFRGB{texte}{text}}
 \end{pspicture}
 &
 \begin{pspicture}(-.5,-.5)(4.2,2.5)
\psaxes[xticksize=2.2,yticksize=4,axesstyle=frame](4,2.2)
\pnode(1,1){B}\pnode(3,2){C}
\psbrace[braceWidthOuter=.5cm](B)(C){\TFRGB{texte}{text}}
 \end{pspicture}
 \\ \hline
\RDD{braceWidth}=5pt \RDI{braceWidth}{pstricks-add} & \RDD{braceWidthInner}=.5cm \RDI{braceWidthInner}{pstricks-add} &
\RDD{braceWidthOuter}=.5cm \RDI{braceWidthOuter}{pstricks-add}
\\ \hline
{\blue \dft : \BS{pslinewidth}}  & {\blue  \dft : 10\BS{pslinewidth}} & {\blue \dft : 10\BS{pslinewidth} }
  \\ \hline
 \begin{pspicture}(-.5,-.5)(4.2,2.5)
\psaxes[xticksize=2.2,yticksize=4,axesstyle=frame](4,2.2)
\pnode(1,1){B}\pnode(3,2){C}
\psbrace[bracePos=.25](B)(C){\TFRGB{texte}{text}}
 \end{pspicture}
&
 \begin{pspicture}(-.5,-.5)(4.2,2.5)
\psaxes[xticksize=2.2,yticksize=4,axesstyle=frame](4,2.2)
\pnode(1,1){B}\pnode(3,2){C}
\psbrace[nodesepA=.5cm](B)(C){\TFRGB{texte}{text}}
 \end{pspicture}
&
 \begin{pspicture}(-.5,-.5)(4.2,2.5)
\psaxes[xticksize=2.2,yticksize=4,axesstyle=frame](4,2.2)
 \pnode(1,1){B}\pnode(3,2){C}
 \psbrace[nodesepB=.5cm](B)(C){\TFRGB{texte}{text}}
  \end{pspicture}
 \\ \hline
\RDD{bracePos}=.25 \RDI{bracePos}{pstricks-add} & 
\RDD{nodesepA}=5pt \RDI{nodesepA}{pstricks-add} & \RDD{nodesepB}=5pt \RDI{nodesepB}{pstricks-add}\\
Position (\%) & \TFRGB{décalage horizontal}{horizontal  offset} & \TFRGB{décalage vertical}{vertical offset}
 \\ \hline
{\blue \dft : .5} &  {\blue\dft : 0pt  }& {\blue \dft : 0pt }
  \\ \hline
 \begin{pspicture}(-.5,-.5)(4.2,2.5)
\psaxes[xticksize=2.2,yticksize=4,axesstyle=frame](4,2.2)
 \pnode(1,1){B}\pnode(3,2){C}
 \psbrace[rot=90](B)(C){\TFRGB{texte}{text}} 
  \end{pspicture} 
&
 \begin{pspicture}(-.5,-.5)(4.2,2.5)
\psaxes[xticksize=2.2,yticksize=4,axesstyle=frame](4,2.2)
 \pnode(1,1){B}\pnode(3,2){C}
 \psbrace[rot=90,ref=r](B)(C){\TFRGB{texte}{text}}
  \end{pspicture} 
&
 \begin{pspicture}(-.5,-.5)(4.2,2.5)
\psaxes[xticksize=2.2,yticksize=4,axesstyle=frame](4,2.2)
 \pnode(1,1){B}\pnode(3,2){C}
 \psbrace[rot=90,ref=l](B)(C){\TFRGB{texte}{text}}
  \end{pspicture} 
   \\ \hline 
\RDD{rot}=90  \RDI{rot}{pstricks-add} & 
rot=90,\RDD{ref} = {\red r}  \RDI{ref}{pstricks-add} & rot=90,ref={\red l }
   \\ \hline 
 \begin{pspicture}(-.5,-.5)(4.2,2.5)
\psaxes[xticksize=2.2,yticksize=4,axesstyle=frame](4,2.2)
 \pnode(1,1){B}\pnode(3,2){C}
 \psbrace[rot=90,ref=b](B)(C){\TFRGB{texte}{text}}
  \end{pspicture} 
 &
 \begin{pspicture}(-.5,-.5)(4.2,2.5)
\psaxes[xticksize=2.2,yticksize=4,axesstyle=frame](4,2.2)
  \pnode(1,1){B}\pnode(3,2){C}
  \psbrace[rot=90,ref=t](B)(C){\TFRGB{texte}{text}}
    \end{pspicture} 
  &
 \begin{pspicture}(-.5,-.5)(4.2,2.5)
\psaxes[xticksize=2.2,yticksize=4,axesstyle=frame](4,2.2)
 \pnode(1,1){B}\pnode(3,2){C}
 \psbrace[rot=90,ref=C](B)(C){\TFRGB{texte}{text}}
  \end{pspicture}
   \\ \hline 
rot=90,ref={\red b} & rot=90,ref={\red t} & rot=90,ref={\red C}
 \\ \hline
 \begin{pspicture}(-.5,-.5)(4.2,2.5)
\psaxes[xticksize=2.2,yticksize=4,axesstyle=frame](4,2.2)
 \pnode(1,1){B}\pnode(3,2){C}
 \psbrace[rot=90,ref=B](B)(C){\TFRGB{texte}{text}}
  \end{pspicture} 
 &
 \begin{pspicture}(-.5,-.5)(4.2,2.5)
\psaxes[xticksize=2.2,yticksize=4,axesstyle=frame](4,2.2)
  \pnode(1,1){B}\pnode(3,2){C}
  \psbrace[rot=90,ref=lC](B)(C){\TFRGB{texte}{text}}
    \end{pspicture} 
  &
 \begin{pspicture}(-.5,-.5)(4.2,2.5)
\psaxes[xticksize=2.2,yticksize=4,axesstyle=frame](4,2.2)
 \pnode(1,1){B}\pnode(3,2){C}
 \psbrace[fillcolor=green](B)(C){\TFRGB{texte}{text}}
  \end{pspicture}
   \\ \hline 
rot=90,ref={\red B} & rot=90,ref={\red lC}& \RDD{fillcolor}=green  \RDI{fillcolor}{pstricks-add}\\ \hline
  \end{tabular} 
  
\bigskip
 

%%%%%======================================================================
\SSCT{Des remplissages spéciaux}{Special fillings}
%\section{Des remplissages spéciaux}
%\subsection{Des gradients de couleurs}
\SbSSCT{Des gradients de couleurs}{Color gradient}

\subsubsection[Module pst-grad]{Module pst-grad \cite{pst-user} \cite{pst-grad}}


%\TFRGB{Utilisation du module}{Package}  \textbf{pst-grad} 
\label{pst-grad}

\begin{center}
\begin{tabular}{|c|c|c|c|} \hline 
\multicolumn{4}{|c|}{\BS{}psframe[fillstyle=\RDDD{gradient}{fillstyle}](0.5,.5)(2.5,2.5)   } \\ \hline
\begin{pspicture}(3,3)
\psframe[fillstyle=gradient](0.5,.5)(2.5,2.5)
\end{pspicture}
&
\begin{pspicture}(3,3)
\psframe[fillstyle=gradient,gradbegin=green](0.5,.5)(2.5,2.5)
\end{pspicture}
&
\begin{pspicture}(3,3)
\psframe[fillstyle=gradient,gradend=green](0.5,.5)(2.5,2.5)
\end{pspicture}
&
\begin{pspicture}(3,3)
\psframe[fillstyle=gradient,gradbegin=red,gradend=green](0.5,.5)(2.5,2.5)
\end{pspicture} \\ \hline
{\blue \dft}  & 
\RDD{gradbegin}=green \RDI{gradbegin}{pst-grad} & 
\RDD{gradend}=green \RDI{gradend}{pst-grad} & 
{\red gradbegin}=red \\
& & & {\red gradend}=green \\ \hline \hline

\begin{pspicture}(3,3)
\psframe[fillstyle=gradient,gradlines=5](0.5,.5)(2.5,2.5)
\end{pspicture}
&
\begin{pspicture}(3,3)
\psframe[fillstyle=gradient,gradmidpoint=0.7](0.5,.5)(2.5,2.5)
\end{pspicture}
&
\begin{pspicture}(3,3)
\psframe[fillstyle=gradient,gradangle=45](0.5,.5)(2.5,2.5)
\end{pspicture}
&
\begin{pspicture}(3,3)
\psframe[fillstyle=gradient,gradangle=90](0.5,.5)(2.5,2.5)
\end{pspicture} \\ \hline
\RDD{gradlines}=5  \RDI{gradlines}{pst-grad} & \RDD{gradmidpoint}=0.7  \RDI{gradmidpoint}{pst-grad} & \RDD{gradangle}=45  \RDI{gradangle}{pst-grad} & \RDD{gradangle}=90  \RDI{gradangle}{pst-grad} \\ \hline
{\blue \dft  : 500} & {\blue \dft  : 0.9} &  \multicolumn{2}{c|}{\blue \dft  : 0}\\ \hline 
\end{tabular}
\end{center}


\begin{center}
\begin{tabular}{|c|c|c|c|} \hline 
\multicolumn{4}{|c|}{\BS{}psframe[  fillstyle=gradient,\RDD{GradientCircle}=true ](0.5,.5)(2.5,2.5)  \RDI{GradientCircle}{pst-grad} } \\ \hline
\begin{pspicture}(3,3)
\psframe[fillstyle=gradient,GradientCircle=true](0.5,.5)(2.5,2.5)
\end{pspicture}
&
\begin{pspicture}(3,3)
\psframe[fillstyle=gradient,GradientCircle=true,GradientScale=.5](0.5,.5)(2.5,2.5)
\end{pspicture}
&
\begin{pspicture}(3,3)
\psframe[fillstyle=gradient,GradientCircle=true,GradientScale=2](0.5,.5)(2.5,2.5)
\end{pspicture}
&
\begin{pspicture}(3,3)
\psframe[fillstyle=gradient,GradientCircle=true,GradientPos={(1,1)}](0.5,.5)(2.5,2.5)
\end{pspicture}\\ \hline
  & \RDD{GradientScale}=.5  \RDI{GradientScale}{pst-grad} & \RDD{GradientScale}=2  \RDI{GradientScale}{pst-grad} &  \RDD{GradientPos}=\AC{(1,1)} \RDI{GradientPos}{pst-grad} \\ \hline
\end{tabular}
\end{center} % pst-grad
 
\newpage
\subsubsection[Module pst-slpe]{Module pst-slpe  \cite{pst-slpe}}




%\TFRGB{Utilisation du module}{Package}  \textbf{pst-slpe}
\label{pst-slpe}



\begin{center}
\begin{tabular}{|c|c|c|c|c|c|} \hline 
\multicolumn{6}{|c|}{\BS{psframe}[\RDD{fillstyle=slope} ](0.5,0.5)(2.5,2.5) \RDI{fillstyle}{pst-slpe} } \\ \hline

\begin{pspicture}(2,2)
\psframe[fillstyle=slope](0.3,.3)(1.7,1.7)
\end{pspicture}
&
\begin{pspicture}(2,2)
\psframe[fillstyle=ccslope](0.3,.3)(1.7,1.7)
\end{pspicture} 
&
\begin{pspicture}(2,2)
\psframe[fillstyle=radslope](0.3,.3)(1.7,1.7)
\end{pspicture}
&
\begin{pspicture}(2,2)
\psframe[fillstyle=slopes](0.3,.3)(1.7,1.7)
\end{pspicture}
&
\begin{pspicture}(2,2)
\psframe[fillstyle=ccslopes](0.3,.3)(1.7,1.7)
\end{pspicture}
&
\begin{pspicture}(2,2)
\psframe[fillstyle=radslopes](0.3,.3)(1.7,1.7)
\end{pspicture} \\ \hline
slope \BDI{slope}{pst-slpe} & 
ccslope \BDI{ccslope}{pst-slpe}&  
radslope \BDI{radslope}{pst-slpe}&
slopes  \BDI{slopes}{pst-slpe}& 
ccslopes \BDI{ccslopes}{pst-slpe}&  
radslopes  \BDI{radslopes}{pst-slpe} \\ \hline
\end{tabular}
\end{center}
%
%%-------------------------------------------------------------



\begin{center}
\begin{tabular}{|c|c|c|c|} \hline 
\multicolumn{4}{|c|}{\BS{}psframe[fillstyle=slope](0.5,0.5)(2.5,1.5) } \\ \hline
\begin{pspicture}(0,.3)(3,1.7)
\psframe[fillstyle=slope](0.5,.5)(2.5,1.5)
\end{pspicture}
&
\begin{pspicture}(0,.3)(3,1.7)
\psframe[fillstyle=slope,slopebegin=green](0.5,.5)(2.5,1.5)
\end{pspicture}
&
\begin{pspicture}(0,.3)(3,1.7)
\psframe[fillstyle=slope,slopeend=green](0.5,.5)(2.5,1.5)
\end{pspicture}
&
\begin{pspicture}(0,.3)(3,1.7)
\psframe[fillstyle=slope,slopebegin=red,slopeend=green](0.5,.5)(2.5,1.5)
\end{pspicture} \\ \hline
{\blue \dft} 	& 
\RDD{slopebegin}=green  \RDI{slopebegin}{pst-slpe}	&
& 
\RDD{slopebegin}=red \\
& 			
& 
\RDD{slopeend}=green  \RDI{slopeend}{pst-slpe}	& \RDD{slopeend}=green \\ \hline
\end{tabular}
\end{center}
%
%%---------------------------------------------------

\begin{center}
\begin{tabular}{|c|c|c|c|} \hline 
\begin{pspicture}(0,-.3)(10,1.8)
\psframe[fillstyle=slopes,slopecolors=0 1 0 0 4 0 1 0 7 0 0 1 3](1,0.5)(9,1.5)
\psaxes[linecolor=black,xsubticks=10,linewidth=1pt,linestyle=dashed,yAxis=false]{->}(0,0)(10,0)

\end{pspicture} \\ \hline

\BS{psframe}[fillstyle=slopes,\RDD{slopecolors}= \psframebox{\red 0} \psovalbox{1 0 0} \psframebox{\red 4} \psovalbox{0 1 0} \psframebox{\red 7} \psovalbox{0 0 1} \psframebox{\blue 3}](1,.5)(9,2.5)  \RDI{slopecolors}{pst-slpe}\\ \hline
\psframebox{\red Position} \psovalbox{couleur en RGB} \psframebox{\blue  nombre de couleurs}\\ \hline
\end{tabular}
\end{center}
%
%%---------------------------------------------------

\begin{center}
\begin{tabular}{|c|c|c|c|c|c|} \hline 
\multicolumn{6}{|c|}{\BS{}psframe[ fillstyle=slope,\RDD{slopesteps}=5 ](0.3,0.3)(1.7,1.7) {\blue  (\dft : 100)} \RDI{slopesteps}{pst-slpe}  } \\ \hline

\begin{pspicture}(2,2)
\psframe[fillstyle=slope,slopesteps=5](0.3,.3)(1.7,1.7)
\end{pspicture}
&
\begin{pspicture}(2,2)
\psframe[fillstyle=ccslope,slopesteps=5](0.3,.3)(1.7,1.7)
\end{pspicture} 
&
\begin{pspicture}(2,2)
\psframe[fillstyle=radslope,slopesteps=5](0.3,.3)(1.7,1.7)
\end{pspicture}
&
\begin{pspicture}(2,2)
\psframe[fillstyle=slopes,slopesteps=5](0.3,.3)(1.7,1.7)
\end{pspicture}
&
\begin{pspicture}(2,2)
\psframe[fillstyle=ccslopes,slopesteps=5](0.3,.3)(1.7,1.7)
\end{pspicture}
&
\begin{pspicture}(2,2)
\psframe[fillstyle=radslopes,slopesteps=5](0.3,.3)(1.7,1.7)
\end{pspicture} \\ \hline
slope & ccslope &  radslope &
 slopes  & ccslopes &  radslopes   \\ \hline
\end{tabular}
\end{center}
%
%%-------------------------------------------------------------
%

\begin{center}
\begin{tabular}{|c|c|c|c|c|c|} \hline 
\multicolumn{6}{|c|}{\BS{}psframe[ fillstyle=slope,\RDD{slopeangle}=45 ](0.5,0.5)(2.5,2.5) {\blue ( \dft 0)} \RDI{slopeangle}{pst-slpe} } \\ \hline

\begin{pspicture}(2,2)
\psframe[fillstyle=slope,slopeangle=45](0.3,.3)(1.7,1.7)
\end{pspicture}
&
\begin{pspicture}(2,2)
\psframe[fillstyle=ccslope,slopeangle=45](0.3,.3)(1.7,1.7)
\end{pspicture} 
&
\begin{pspicture}(2,2)
\psframe[fillstyle=radslope,slopeangle=45](0.3,.3)(1.7,1.7)
\end{pspicture}
&
\begin{pspicture}(2,2)
\psframe[fillstyle=slopes,slopeangle=45](0.3,.3)(1.7,1.7)
\end{pspicture}
&
\begin{pspicture}(2,2)
\psframe[fillstyle=ccslopes,slopeangle=45](0.3,.3)(1.7,1.7)
\end{pspicture}
&
\begin{pspicture}(2,2)
\psframe[fillstyle=radslopes,slopeangle=45](0.3,.3)(1.7,1.7)
\end{pspicture} \\ \hline
slope & ccslope &  radslope &
 slopes  & ccslopes &  radslopes   \\ \hline
\end{tabular}
\end{center}
%
%%-------------------------------------------------------------


\begin{center}
\begin{tabular}{|c|c|c|c|c|c|} \hline 
\multicolumn{6}{|c|}{\BS{}psframe[ fillstyle=slope,\RDD{slopecenter}= .25 .25](0.5,0.5)(2.5,2.5) {\blue (\dft .5 .5)} \RDI{slopecenter}{pst-slpe} } \\ \hline

\begin{pspicture}(2,2)
\psframe[fillstyle=slope,slopecenter= .25 .25](0.3,.3)(1.7,1.7)
\end{pspicture}
&
\begin{pspicture}(2,2)
\psframe[fillstyle=ccslope,slopecenter= .25 .25](0.3,.3)(1.7,1.7)
\end{pspicture} 
&
\begin{pspicture}(2,2)
\psframe[fillstyle=radslope,slopecenter= .25 .25](0.3,.3)(1.7,1.7)
\end{pspicture}
&
\begin{pspicture}(2,2)
\psframe[fillstyle=slopes,slopecenter= .25 .25](0.3,.3)(1.7,1.7)
\end{pspicture}
&
\begin{pspicture}(2,2)
\psframe[fillstyle=ccslopes,slopecenter= .25 .25](0.3,.3)(1.7,1.7)
\end{pspicture}
&
\begin{pspicture}(2,2)
\psframe[fillstyle=radslopes,slopecenter= .25 .25](0.3,.3)(1.7,1.7)
\end{pspicture} \\ \hline
slope & ccslope &  radslope &
 slopes  & ccslopes &  radslopes   \\ \hline
\end{tabular}
\end{center}

%%-------------------------------------------------------------

\begin{center}
\begin{tabular}{|c|c|c|c|c|c|} \hline 
\multicolumn{6}{|c|}{\BS{}psframe[ fillstyle=slope,\RDD{sloperadius} =.75 ](0.5,0.5)(2.5,2.5) {\blue (\dft 0.5cm) }  \RDI{sloperadius}{pst-slpe} } \\ \hline

\begin{pspicture}(2,2)
\psframe[fillstyle=slope,sloperadius =.75](0.3,.3)(1.7,1.7)
\end{pspicture}
&
\begin{pspicture}(2,2)
\psframe[fillstyle=ccslope,sloperadius =.75](0.3,.3)(1.7,1.7)
\end{pspicture} 
&
\begin{pspicture}(2,2)
\psframe[fillstyle=radslope,sloperadius =.75](0.3,.3)(1.7,1.7)
\end{pspicture}
&
\begin{pspicture}(2,2)
\psframe[fillstyle=slopes,sloperadius =.75](0.3,.3)(1.7,1.7)
\end{pspicture}
&
\begin{pspicture}(2,2)
\psframe[fillstyle=ccslopes,sloperadius =.75](0.3,.3)(1.7,1.7)
\end{pspicture}
&
\begin{pspicture}(2,2)
\psframe[fillstyle=radslopes,sloperadius =.75](0.3,.3)(1.7,1.7)
\end{pspicture} \\ \hline
slope & ccslope &  radslope &
 slopes  & ccslopes &  radslopes   \\ \hline
\end{tabular}
\end{center}
%
%%-------------------------------------------------------------

\begin{center}
\begin{tabular}{|c|c|c|c|c|c|} \hline 
\multicolumn{6}{|c|}{\BS{}psframe[\RDD{fading},  fillstyle=slope](0.5,0.5)(2.5,2.5) \RDI{fading}{pst-slpe} } \\ \hline

\begin{pspicture}(2,2)
\psframe[fading,fillstyle=slope](0.3,.3)(1.7,1.7)
\end{pspicture}
&
\begin{pspicture}(2,2)
\psframe[fading,fillstyle=ccslope](0.3,.3)(1.7,1.7)
\end{pspicture} 
&
\begin{pspicture}(2,2)
\psframe[fading,fillstyle=radslope](0.3,.3)(1.7,1.7)
\end{pspicture}
&
\begin{pspicture}(2,2)
\psframe[fading,fillstyle=slopes](0.3,.3)(1.7,1.7)
\end{pspicture}
&
\begin{pspicture}(2,2)
\psframe[fading,fillstyle=ccslopes](0.3,.3)(1.7,1.7)
\end{pspicture}
&
\begin{pspicture}(2,2)
\psframe[fading,fillstyle=radslopes](0.3,.3)(1.7,1.7)
\end{pspicture} \\ \hline
slope & ccslope &  radslope &
 slopes  & ccslopes &  radslopes   \\ \hline
\end{tabular}
\end{center}
%
%%-------------------------------------------------------------


\begin{center}
\begin{tabular}{|c|c|c|c|c|c|} \hline 
\multicolumn{6}{|c|}{\BS{}psframe[fading, \RDD{startfading}=0.5,  fillstyle=slope](0.5,0.5)(2.5,2.5)  \RDI{startfading}{pst-slpe}  } \\ \hline

\begin{pspicture}(2,2)
\psframe[fading,startfading=0.5,fillstyle=slope](0.3,.3)(1.7,1.7)
\end{pspicture}
&
\begin{pspicture}(2,2)
\psframe[fading,startfading=0.5,fillstyle=ccslope](0.3,.3)(1.7,1.7)
\end{pspicture} 
&
\begin{pspicture}(2,2)
\psframe[fading,startfading=0.5,fillstyle=radslope](0.3,.3)(1.7,1.7)
\end{pspicture}
&
\begin{pspicture}(2,2)
\psframe[fading,startfading=0.5,,fillstyle=slopes](0.3,.3)(1.7,1.7)
\end{pspicture}
&
\begin{pspicture}(2,2)
\psframe[fading,startfading=0.5,fillstyle=ccslopes](0.3,.3)(1.7,1.7)
\end{pspicture}
&
\begin{pspicture}(2,2)
\psframe[fading,startfading=0.5,fillstyle=radslopes](0.3,.3)(1.7,1.7)
\end{pspicture} \\ \hline
slope & ccslope &  radslope &
 slopes  & ccslopes &  radslopes   \\ \hline
\end{tabular}
\end{center}
%
%%-------------------------------------------------------------

\begin{center}
\begin{tabular}{|c|c|c|c|c|c|} \hline 
\multicolumn{6}{|c|}{\BS{}psframe[fading, \RDD{endfading}=0.5,  fillstyle=slope](0.5,0.5)(2.5,2.5) \RDI{endfading}{pst-slpe}  } \\ \hline

\begin{pspicture}(2,2)
\psframe[fading,endfading=0.5,fillstyle=slope](0.3,.3)(1.7,1.7)
\end{pspicture}
&
\begin{pspicture}(2,2)
\psframe[fading,endfading=0.5,fillstyle=ccslope](0.3,.3)(1.7,1.7)
\end{pspicture} 
&
\begin{pspicture}(2,2)
\psframe[fading,endfading=0.5,fillstyle=radslope](0.3,.3)(1.7,1.7)
\end{pspicture}
&
\begin{pspicture}(2,2)
\psframe[fading,endfading=0.5,,fillstyle=slopes](0.3,.3)(1.7,1.7)
\end{pspicture}
&
\begin{pspicture}(2,2)
\psframe[fading,endfading=0.5,fillstyle=ccslopes](0.3,.3)(1.7,1.7)
\end{pspicture}
&
\begin{pspicture}(2,2)
\psframe[fading,endfading=0.5,fillstyle=radslopes](0.3,.3)(1.7,1.7)
\end{pspicture} \\ \hline
slope & ccslope &  radslope &  slopes  & ccslopes &  radslopes   \\ \hline
\end{tabular}
\end{center}
%
%%-------------------------------------------------------------

\begin{center}
\begin{tabular}{|c|c|c|c|c|c|} \hline 
\multicolumn{6}{|c|}{\BSS{psBall} [option](1,1)\AC{blue}\AC{.8} \BSI{psBall}{pst-slpe} }\\ \hline

%\begin{pspicture}(2,2)
%\psBall(1,1){black}{.8}
%\end{pspicture}
%&
\begin{pspicture}(2,2)
\psBall(1,1){blue}{.8}
\end{pspicture} 
&
\begin{pspicture}(2,2)
\psBall[sloperadius=10pt](1,1){blue}{.8}

\end{pspicture}
&
\begin{pspicture}(2,2)
\psBall[slopebegin=red](1,1){blue}{.8}
\end{pspicture}
&
\begin{pspicture}(2,2)
\psBall[slopeend=red](1,1){blue}{.8}
\end{pspicture}
&
\begin{pspicture}(2,2)
\psBall[fading](1,1){blue}{.8}
\end{pspicture}
&
\begin{pspicture}(2,2)
\psBall[slopesteps=5](1,1){blue}{.8}
\end{pspicture}
 \\ \hline
sans option & sloperadius=10pt &  slopebegin=red  & slopeend=red &  fading & slopesteps=5 \\ \hline
\end{tabular}
\end{center}
 % pst-slpe

\newpage
%\subsection[Remplissage par des motifs]{Remplissage par des motifs \cite{pst-fill}}
\SbSSCT{Remplissage par des motifs}{Filling with pattern}

\TFRGB{Utilisation du module}{Package} \textbf{pst-fill} \label{pst-fill}

\psboxfill{\MonMotif}

\begin{center}
\begin{tabular}{|p{14cm}|} \hline
\multicolumn{1}{|c|}{ \TFRGB{Création du motif}{Creating the pattern} : \begin{pspicture}(-.5,-.27)(.5,.27)
 \psframe[dimen=middle,fillcolor=yellow,fillstyle=solid,linecolor=blue](-.5,-.25)(.5,.25)
 \pscircle[dimen=middle,fillcolor=magenta,fillstyle=solid,linecolor=magenta]{.1}
\end{pspicture} 
} \\ \hline 
\BS{}newcommand\AC{{\red \BS{}MonMotif}}\{

\BS{}begin\AC{pspicture}(1,.5)

\BS{}psframe[dimen=middle,fillcolor=yellow,fillstyle=solid,linecolor=blue](1,.5)

\BS{}pscircle[dimen=middle,fillcolor=magenta,fillstyle=solid,linecolor=magenta](.5,.25)\AC{.1}

\BS{}end\AC{pspicture} \} \\ \hline 
\multicolumn{1}{|c|}{\TFRGB{Utilisation du motif}{Using the pattern} : \BSS{psboxfill}\AC{{\red\BS{}MonMotif}}  \BSI{psboxfill}{pst-fill} }\\   \hline

\multicolumn{1}{|l|}{\BS{}pspolygon[fillstyle=\RDDD{boxfill}{fillstyle}](0,.5)(1,2)(3,2)(4,0)(2,.5)}\\  
\multicolumn{1}{|c|}{\begin{pspicture}(2.5,2.5) \pspolygon[fillstyle=boxfill](0,.5)(1,2)(3,2)(4,0)(2,.5) \end{pspicture} }\\ \hline
\end{tabular}
\end{center}

%---------------------------------------------------
 
 

\psset{arrowscale=2}
\psset{linecolor=blue,fillcolor=yellow,linewidth=1pt}


 
 
 \begin{center}
 \begin{tabular}{|c|c|c|} \hline 
\begin{pspicture}(2.5,2.5)
\psframe[fillstyle=boxfill,fillangle=90](0,0)(2,2)
\end{pspicture}
&
\begin{pspicture}(4,2.5)
\pspolygon[fillstyle=boxfill,fillangle=45](0,.5)(1,2)(3,2)(4,0)(2,.5)
\end{pspicture}
&
\begin{pspicture}(2.5,2.5)
\pscircle[fillstyle=boxfill,fillangle=-30](1,1){1}
\end{pspicture}   \\ \hline
 \RDD{fillangle}=90  \RDI{fillangle}{pst-fill} & {\red fillangle}=45 & {\red fillangle}=-30 \\ \hline
\multicolumn{3}{|c|}{ \blue \dft  : 0 }\\ \hline

 \end{tabular}
 \end{center}
 
 
\begin{center}
\begin{tabular}{|c|c|c|} \hline 
\begin{pspicture}(4,2.5)
\pspolygon[fillstyle=boxfill,fillsepx=.25cm](0,.5)(1,2)(3,2)(4,0)(2,.5)
\end{pspicture}
&
\begin{pspicture}(4,2.5)
\pspolygon[fillstyle=boxfill,fillsepy=.25cm](0,.5)(1,2)(3,2)(4,0)(2,.5)
\end{pspicture}
&
\begin{pspicture}(4,2.5)
\pspolygon[fillstyle=boxfill,fillsep=.25cm](0,.5)(1,2)(3,2)(4,0)(2,.5)
\end{pspicture} \\ \hline
\RDD{fillsepx}=.25cm  \RDI{fillsepx}{pst-fill} &
\RDD{fillsepy}=.25cm  \RDI{fillsepy}{pst-fill} &  \RDD{fillsep}=.25cm  \RDI{fillsep}{pst-fill}\\ \hline
\end{tabular}
\end{center}

 
\begin{center}
\begin{tabular}{|c|c|c|} \hline 
\begin{pspicture}(4,2.5)
\pspolygon[fillstyle=boxfill,fillcyclex=3](0,.5)(1,2)(3,2)(4,0)(2,.5)
\end{pspicture}
&
\begin{pspicture}(4,2.5)
\pspolygon[fillstyle=boxfill,fillcycley=3](0,.5)(1,2)(3,2)(4,0)(2,.5)
\end{pspicture}
&
\begin{pspicture}(4,2.5)
\pspolygon[fillstyle=boxfill,fillcycle=3](0,.5)(1,2)(3,2)(4,0)(2,.5)
\end{pspicture} \\ \hline
\RDD{fillcyclex}=3  \RDI{fillcyclex}{pst-fill} & 
\RDD{fillcycley}=3  \RDI{fillcycley}{pst-fill} &  
\RDD{fillcycle}=3  \RDI{fillcycle}{pst-fill} \\ \hline
\multicolumn{3}{|c|}{ 3  \TFRGB{correspond à}{corresponds to} $1/3$, {\blue \dft  : 0} }\\ \hline
\end{tabular}
\end{center}

\begin{center}
\begin{tabular}{|c|c|c|} \hline 
\begin{pspicture}(4,2.5)
\pspolygon[fillstyle=boxfill,fillmovex=.1](0,.5)(1,2)(3,2)(4,0)(2,.5)
\end{pspicture}
&
\begin{pspicture}(4,2.5)
\pspolygon[fillstyle=boxfill,fillmovey=.1](0,.5)(1,2)(3,2)(4,0)(2,.5)
\end{pspicture}
&
\begin{pspicture}(4,2.5)
\pspolygon[fillstyle=boxfill,fillmove=.1](0,.5)(1,2)(3,2)(4,0)(2,.5)
\end{pspicture} \\ \hline
\RDD{fillmovex}=.1 \RDI{fillmovex}{pst-fill}  & 
\RDD{fillmovey}=.1 \RDI{fillmovey}{pst-fill}  &  \RDD{fillmove}=.1 \RDI{fillmove}{pst-fill}  \\ \hline
\multicolumn{3}{|c|}{ .1  \TFRGB{correspond à}{corresponds to} 0,1 cm , {\blue \dft  : 0 }}\\ \hline
\end{tabular}
\end{center}


\begin{center}
\begin{tabular}{|c|c|c|c|} \hline 
\multicolumn{4}{|c|}{\BS{}pspolygon[fillstyle=boxfill](0.6,.7)(1.8,2.4)(2.3,.6)(4,1.2)(2.5,.3) }\\ \hline
\begin{pspicture}(.5,0)(4,3)
\pspolygon[fillstyle=boxfill](0.6,.7)(1.8,2.4)(2.3,1.6)(3.8,1.2)(2.5,.3)
\end{pspicture}
&
\begin{pspicture}(.5,0)(4,3)
\pspolygon[fillstyle=boxfill,fillloopaddx=1](0.6,.7)(1.8,2.4)(2.3,1.6)(3.8,1.2)(2.5,.3)
\end{pspicture}
&
\begin{pspicture}(.5,0)(4,3)
\pspolygon[fillstyle=boxfill,fillloopaddy=1](0.6,.7)(1.8,2.4)(2.3,1.6)(3.8,1.2)(2.5,.3)
\end{pspicture} 
&
\begin{pspicture}(.5,0)(4,3)
\pspolygon[fillstyle=boxfill,fillloopadd=1](0.6,.7)(1.8,2.4)(2.3,1.6)(3.8,1.2)(2.5,.3)
\end{pspicture} \\ \hline
  & 
  \RDD{fillloopaddx}=1  \RDI{fillloopaddx}{pst-fill}  & \RDD{fillloopaddy}=1  \RDI{fillloopaddy}{pst-fill}  & 
  \RDD{fillloopadd}=1  \RDI{fillloopadd}{pst-fill}  \\ \hline
\multicolumn{4}{|c|}{\blue \dft : 0} \\ \hline
\end{tabular}
\end{center} % pst-fill

%\subsection[Remplissage par des points aléatoires]{Remplissage par des points aléatoires \cite{pstricks-add}}
\SbSSCT{Remplissage par des points aléatoires}{random points filling}


\psset{unit=1cm,fillcolor=yellow,fillcolor=white}


\begin{center}
\begin{tabular}{|c|c|c|} \hline 
\multicolumn{3}{|c|}{\BSS{psRandom}\AC{\BS{}pspolygon(0,.5)(1,2)(3,2)(4,0)(2,.5)} \BSI{psRandom}{pstricks-add}  }\\ \hline
\begin{pspicture}(4,2.5)
\psRandom{\pspolygon(0,.5)(1,2)(3,2)(4,0)(2,.5)}
\end{pspicture}
&
\begin{pspicture}(4,2.5)
\psRandom(1,1)(2,2){\pspolygon(0,.5)(1,2)(3,2)(4,0)(2,.5)}
\end{pspicture}
&
\begin{pspicture}(4,2.5)
\psRandom(0,0)(4,2){\pspolygon(0,.5)(1,2)(3,2)(4,0)(2,.5)}
\end{pspicture} \\ \hline
\BS{psRandom}\AC{\BS{}pspolygon \dots }&
\BS{psRandom}{\red (1,1)(2,2)}\AC{\dots }&
\BS{psRandom}{\red (0,0)(4,2)}\AC{\dots }\\ \hline
\end{tabular}
\end{center}

\begin{center}
\begin{tabular}{|c|c|c|} \hline 
\multicolumn{3}{|c|}{\BS{psRandom}[{\red options}](0,0)(4,2)\AC{\BS{}pspolygon(0,.5)(1,2)(3,2)(4,0)(2,.5)}}\\ \hline
\begin{pspicture}(4,2.5)
\psRandom[randomPoints=100](0,0)(4,2){\pspolygon(0,.5)(1,2)(3,2)(4,0)(2,.5)}
\end{pspicture}
&
\begin{pspicture}(4,2.5)
\psRandom[color](0,0)(4,2){\pspolygon(0,.5)(1,2)(3,2)(4,0)(2,.5)}
\end{pspicture}
&
\begin{pspicture}(4,2.5)
\psRandom[dotstyle=+](0,0)(4,2){\pspolygon(0,.5)(1,2)(3,2)(4,0)(2,.5)}
\end{pspicture} \\ \hline
\RDD{randomPoints}=100 \RDI{randomPoints}{pstricks-add} &
\RDD{color}  \RDI{color}{pstricks-add}&
\RDD{dotstyle}=+  \RDI{dotstyle}{pstricks-add}
 \\ \hline  
{\blue \dft :  1000} & &   \\ \hline  
\end{tabular}
\end{center} % random dots
\newpage

% ========================================================================
%\section[Effets spéciaux avec du texte ]{Effets spéciaux avec du texte  \cite{pst-user}  \cite{pst-text}}
\SSCT{Effets spéciaux avec du texte}{Special effects}

\subsection{pstextpath}
\label{pst-text}

%\psset{doublesep=3pt}

\begin{tabular}{|c|} \hline  
\rnode{A}{position} \hspace{.5cm} \rnode{B}{\TFRGB{décalage}{offset}} \hspace{.5cm} \rnode{C}{\TFRGB{support graphique}{graphic support}}\\
\\
\BS{pstextpath}\rnode{AA}{\red[r]} \rnode{BB}{\cyan (0,0)}\Rnode{CC}{\Vert \AC{\BS{psline}(0, 0)(5, 1)}}\AC{\BS{red} \TFRGB{texte }{text}}
\\ \hline 
\begin{pspicture}(-1,-.15)(6,1.5) 
 \pstextpath[r](0,0){\psline(0, 0)(5, 1)}{\red \TFRGB{texte}{text}}
\end{pspicture} 
\\ \hline 
\end{tabular} 
\psset{linestyle = solid}
\ncline{A}{AA}
\ncline{B}{BB}
\ncline{C}{CC}

\SbSbSSCT{Positionnement sur différents objets graphiques}{Positioning on different graphic objects}


\psset{linestyle = dotted}
\begin{tabular}{|c|c|c|} \hline  
 \multicolumn{3}{|l|}{ \hspace{1cm} \BS{pstextpath}{\red[r]} (0,0)\AC{\BS{psline}(0, 0)(1, 1)(2, 1)(3.5, 3.5)}\AC{\BS{red} \TFRGB{à droite}{on the right}} } \\ 
 \multicolumn{3}{|l|}{  \hspace{1cm} \BS{pstextpath}{\red[l]} (0,0)\AC{\BS{psline}(0, 0)(1, 1)(2, 1)(3.5, 3.5)}\AC{\BS{cyan} \TFRGB{à gauche}{on the left}} } \\ 
 \multicolumn{3}{|l|}{ \hspace{1cm} \BS{pstextpath}{\red[c]} (0,0)\AC{\BS{psline}(0, 0)(1, 1)(2, 1)(3.5, 3.5)}\AC{ \TFRGB{centré}{centered} }} \\  
 \hline

\begin{psgraph}[axesstyle=none,ticks=none,labels=none](0,0)(4,4){4cm}{4cm} 
 \pstextpath[r](0,0){\psline(0, 0)(1, 1)(2, 1)(3.5, 3.5)}{\red \TFRGB{à droite}{on the right}}
 \pstextpath[l](0,0){\psline(0, 0)(1, 1)(2, 1)(3.5, 3.5)}{\cyan \TFRGB{à gauche}{on the left}}
 \pstextpath[c](0,0){\psline(0, 0)(1, 1)(2, 1)(3.5, 3.5)}{\TFRGB{centré}{centered}}
 \end{psgraph}
&  
\begin{psgraph}[axesstyle=none,ticks=none,labels=none](0,0)(4,4){4cm}{4cm} 
\pstextpath[r](0,0){\pspolygon(1, 0.5)(2, 3)(3.5, 1.5)(2.5,1)}{\red \TFRGB{à droite}{on the right}}
\pstextpath[l](0,0){\pspolygon(1, 0.5)(2, 3)(3.5, 1.5)(2.5,1)}{\cyan \TFRGB{à gauche}{on the left}}
\pstextpath[c](0,0){\pspolygon(1, 0.5)(2, 3)(3.5, 1.5)(2.5,1)}{\TFRGB{centré}{centered}}
 \end{psgraph}
&  
\begin{psgraph}[axesstyle=none,ticks=none,labels=none](0,0)(4,4){4cm}{4cm} 
\pstextpath[r](0,0){\psframe(.5,.5)(3.5,3.5)}{\red \TFRGB{à droite}{on the right}}
\pstextpath[l](0,0){\psframe(.5,.5)(3.5,3.5)}{\cyan \TFRGB{à gauche}{on the left}}
\pstextpath[c](0,0){\psframe(.5,.5)(3.5,3.5)}{\TFRGB{centré}{centered}}
 \end{psgraph}

\\ \hline  
\BS{psline} & \BS{pspolygon} & \BS{psframe}
\\ \hline  
\begin{psgraph}[axesstyle=none,ticks=none,labels=none](0,0)(4,4){4cm}{4cm} 
\pstextpath[r](0,0){\psdiamond(2,2)(1.75,1.75)}{\red \TFRGB{à droite}{on the right}}
\pstextpath[l](0,0){\psdiamond(2,2)(1.75,1.75)}{\cyan \TFRGB{à gauche}{on the left}}
\pstextpath[c](0,0){\psdiamond(2,2)(1.75,1.75)}{\TFRGB{centré}{centered}}
 \end{psgraph}
&  
\begin{psgraph}[axesstyle=none,ticks=none,labels=none](0,0)(4,4){4cm}{4cm} 
\pstextpath[r](0,0){\pstriangle(2,1)(3.5,2)}{\red \TFRGB{à droite}{on theright}}
\pstextpath[l](0,0){\pstriangle(2,1)(3.5,2)}{\cyan \TFRGB{à gauche}{on the left}}
\pstextpath[c](0,0){\pstriangle(2,1)(3.5,2)}{\TFRGB{centré}{centered}}
 \end{psgraph}
&  
\begin{psgraph}[axesstyle=none,ticks=none,labels=none](0,0)(4,4){4cm}{4cm}  
\pstextpath[r](0,0){\psarc(1.5,2){2}{-120}{120}}{\red \TFRGB{à droite}{on the right}}
\pstextpath[l](0,0){\psarc(1.5,2){2}{-120}{120}}{\cyan \TFRGB{à gauche}{on the left}}
\pstextpath[c](0,0){\psarc(1.5,2){2}{-120}{120}}{\TFRGB{centré}{centered}}
\end{psgraph}
\\ \hline 
\BS{psdiamond}& \BS{pstriangle} & \BS{psarc}
\\ \hline 


\end{tabular} 

\ncline{A}{AA}
\ncline{B}{BB}
\ncline{C}{CC}

%
%
%
\begin{tabular}{|c|c|c|} \hline 
\begin{psgraph}[axesstyle=none,ticks=none,labels=none](0,0)(4,4){4cm}{4cm} 
 \pstextpath[r](0,0){\pscircle(2,2){1.5}}{\red \TFRGB{à droite}{on the right}}
 \pstextpath[l](0,0){\pscircle(2,2){1.5}}{\cyan \TFRGB{à gauche}{on the left}}
 \pstextpath[c](0,0){\pscircle(2,2){1.5}}{\TFRGB{centré}{centered}}
\end{psgraph}
&
\begin{psgraph}[axesstyle=none,ticks=none,labels=none](0,0)(4,4){4cm}{4cm} 
\pstextpath[r](0,0){\psellipse(2, 2)(1.75,1)}{\red \TFRGB{à droite}{on the right}}
\pstextpath[l](0,0){\psellipse(2, 2)(1.75,1)}{\cyan \TFRGB{à gauche}{on the left}}
\pstextpath[c](0,0){\psellipse(2, 2)(1.75,1)}{\TFRGB{centré}{centered}}
\end{psgraph}
&
\begin{psgraph}[axesstyle=none,ticks=none,labels=none](0,0)(4,4){4cm}{4cm} 
\pstextpath[r](0,0){\psellipticarc(2, 2)(1.75,1){150}{30}}{\red \TFRGB{à droite}{on the right}}
\pstextpath[l](0,0){\psellipticarc(2, 2)(1.75,1){150}{30}}{\cyan \TFRGB{à gauche}{on the left}}
\pstextpath[c](0,0){\psellipticarc(2, 2)(1.75,1){150}{30}}{\TFRGB{centré}{centered}}
 \end{psgraph}
\\ \hline 
\BS{pscircle} & \BS{psellipse} & \BS{psellipticarc} \\ \hline 
\begin{psgraph}[axesstyle=none,ticks=none,labels=none](0,0)(4,4){4cm}{4cm} 
\pstextpath[r](0,0){\pscurve[showpoints=true](1,1)(2,3)(3,2)(2,1)(3,1)}{\red \TFRGB{à droite}{on the right}}
\pstextpath[l](0,0){\pscurve[showpoints=true](1,1)(2,3)(3,2)(2,1)(3,1)}{\cyan \TFRGB{à gauche}{on the left}}
\pstextpath[c](0,0){\pscurve[showpoints=true](1,1)(2,3)(3,2)(2,1)(3,1)}{\TFRGB{centré}{centered}}
 \end{psgraph}
&
\begin{psgraph}[axesstyle=none,ticks=none,labels=none](0,0)(4,4){4cm}{4cm} 
\pstextpath[r](0,0){\psccurve[showpoints=true](1,1)(2,3)(3,2)(2,1)(3,1)}{\red \TFRGB{à droite}{on the right}}
\pstextpath[l](0,0){\psccurve[showpoints=true](1,1)(2,3)(3,2)(2,1)(3,1)}{\cyan \TFRGB{à gauche}{on the left}}
\pstextpath[c](0,0){\psccurve[showpoints=true](1,1)(2,3)(3,2)(2,1)(3,1)}{\TFRGB{centré}{centered}}
 \end{psgraph}
&
\begin{psgraph}[axesstyle=none,ticks=none,labels=none](0,0)(4,4){4cm}{4cm} 
\pstextpath[r](0,0){\psecurve[showpoints=true](1,1)(2,3)(3,2)(2,1)(3,1)}{\red \TFRGB{à droite}{on the right}}
\pstextpath[l](0,0){\psecurve[showpoints=true](1,1)(2,3)(3,2)(2,1)(3,1)}{\cyan \TFRGB{à gauche}{on the left}}
\pstextpath[c](0,0){\psecurve[showpoints=true](1,1)(2,3)(3,2)(2,1)(3,1)}{\TFRGB{centré}{centered}}
\end{psgraph}
\\ \hline  
\BS{pscurve} & \BS{psccurve} & \BS{psecurve}
\\ \hline  
\begin{psgraph}[axesstyle=none,ticks=none,labels=none](0,0)(4,4){4cm}{4cm} 
 \pstextpath[r](0,0){\pswedge(1,2){2}{-60}{60}}{\red \TFRGB{à droite}{on the right}}
 \pstextpath[l](0,0){\pswedge(1,2){2}{-60}{60}}{\cyan \TFRGB{à gauche}{on the left}}
 \pstextpath[c](0,0){\pswedge(1,2){2}{-60}{60}}{\TFRGB{centré}{centered}}
 \end{psgraph}
&  
\begin{psgraph}[axesstyle=none,ticks=none,labels=none](0,0)(4,4){4cm}{4cm} 
\pstextpath[r](0,0){\psbezier(1,0)(8,4)(-4,4)(3,0)}{\red \TFRGB{à droite}{on the right}}
\pstextpath[l](0,0){\psbezier(1,0)(8,4)(-4,4)(3,0)}{\cyan \TFRGB{à gauche}{on the left}}
\pstextpath[c](0,0){\psbezier(1,0)(8,4)(-4,4)(3,0)}{\TFRGB{centré}{centered}}
 \end{psgraph}
&  
\begin{psgraph}[axesstyle=none,ticks=none,labels=none,dx=\psPiH](0,-2)(13,2){4cm}{4cm} 
\pstextpath[r](0,0){\psplot[algebraic,plotpoints=200]{0}{12.56}{sin(x)}}{\red \TFRGB{à droite}{on the right}}
\pstextpath[l](0,0){\psplot[algebraic,plotpoints=200]{0}{12.56}{sin(x)}}{\cyan \TFRGB{à gauche}{on the left}}
\pstextpath[c](0,0){\psplot[algebraic,plotpoints=200]{0}{12.56}{sin(x)}}{\TFRGB{centré}{centered}}
 \end{psgraph}
\\ \hline 
\BS{pswedge} & \BS{psbezier} & \BS{psplot}[algebraic]\AC{0}\AC{12.56}\AC{sin(x)}
\\ \hline 
\end{tabular}

%------------------------------------------------------------------------
\bigskip
\SbSbSSCT{Décalage}{Offset}
\psset{linestyle = solid}

\begin{tabular}{|c|c|c|c|}\hline  
 \multicolumn{4}{|c|}{ \BS{pstextpath}[l]{\red (0,.5)}\AC{\BS{psline}(0, 0)(1, 1)}\AC{\TFRGB{texte}{text}} } \\ \hline
\begin{pspicture}(-1.5,-.5)(1.5,1.5)
 \psset{linewidth=0.2pt}
\pstextpath[l](0,.5){\psline(0, 0)(1, 1)} {\TFRGB{texte}{text}}
%\psdots(-.75,0)
 \end{pspicture}
&  
 \begin{pspicture}(-1.5,-.5)(1.5,1.5)
  \psset{linewidth=0.2pt}
 \pstextpath[l](0,-.5){\psline(0, 0)(1, 1)}{\TFRGB{texte}{text}}
%  \psdots(-.75,0)
  \end{pspicture}
&  
  \begin{pspicture}(-1.5,-.5)(1.5,1.5)
   \psset{linewidth=0.2pt}
  \pstextpath[l](.5,0){\psline(0, 0)(1, 1)}{\TFRGB{texte}{text}}
%    \psdots(-.75,0)
   \end{pspicture}
&  
  \begin{pspicture}(-1.5,-.5)(1.5,1.5)
   \psset{linewidth=0.2pt}
  \pstextpath[l](.5,.5){\psline(0, 0)(1, 1)}{\TFRGB{texte}{text}}
%    \psdots(-.75,0)
   \end{pspicture}
\\ \hline 
(0,0.5) & (0,-0.5) & (0.5,0) & (0.5,0.5)
\\ \hline 
 \multicolumn{4}{|c|}{ \blue \dft : (0,\BDD{\BS{TPoffset}}) \hspace{.5cm}  \BS{TPoffset}= -0.7ex. }
\\ \hline 
\end{tabular} 



\subsection{pscharpath}

\psset{linestyle = solid}
\TFRGB{
\BS{DeclareFixedFont}\AC{\BS[nom]}\AC{\BS{encodage}}
\AC{\BS{famille}}\AC{\BS{Séries}}\AC{\BS{forme}}\AC{taille}
}{
\BS{DeclareFixedFont}\AC{\BS[nom]}\AC{\BS{encodingdefault}}
\AC{\BS{familydefault}}\AC{\BS{seriesdefault}}\AC{\BS{shapedefault}}\AC{taille}
}

\bigskip

\rnode{E}{\TFRGB{nom de la fonte}{name}} \hspace{.5cm}
\rnode{A}{\TFRGB{encodage}{encoding }: T1} \hspace{.5cm}
\rnode{B}{\TFRGB{famille}{family} : Times} \hspace{.5cm} 
\rnode{C}{\TFRGB{séries}{series} : bold} \hspace{.5cm}
\rnode{D}{\TFRGB{forme}{shape} :  normale}

\bigskip
\begin{tabular}{|c|} \hline  
\BSS{DeclareFixedFont}\AC{\rnode{EE}{\BS{Font}}}\AC{\rnode{AA}{T1}}\AC{\rnode{BB}{ptm}}\AC{\rnode{CC}{b}}\AC{\rnode{DD}{n}}\AC{2cm} \BSI{DeclareFixedFont}{pst-text} \\
%
\BSS{pscharpath}\AC{\BS{Font} PSTricks} \BSI{pscharpath}{pst-text}
\\ \hline  
\DeclareFixedFont{\Font}{T1}{ptm}{b}{n}{2cm}

\pscharpath{\Font PSTricks}
\\ \hline 
\end{tabular} 
\ncline{A}{AA}
\ncline{B}{BB}
\ncline{C}{CC}
\ncline{D}{DD}
\ncline{E}{EE}

%\bigskip
%
%\begin{tabular}{|c|} \hline  
%\BSS{DeclareFixedFont}\AC{\BS{RM}}\AC{T1}\AC{pcr}\AC{b}\AC{n}\AC{2cm} \\
%%
%\BSS{pscharpath}\AC{\BS{RM} PSTricks}
%\\ \hline  
%\DeclareFixedFont{\RM}{T1}{pcr}{b}{n}{2cm}
%
%\pscharpath{\RM PSTricks}
%\\ \hline 
%\end{tabular} 
%\bigskip
%\BSS{DeclareFixedFont}\AC{\BS{RM}}\AC{T1}\AC{ppl}\AC{b}\AC{n}\AC{2cm} \\
%\BSS{pscharpath}\AC{\BS{RM} PSTricks}

\SbSbSSCT{Quelques familles de fonte}{Some families}

\begin{tabular}{|c|c|} \hline  
\DeclareFixedFont{\Font}{T1}{ppl}{b}{n}{1.5cm}

\pscharpath{\Font PSTricks}
&
\DeclareFixedFont{\Font}{T1}{pag}{b}{n}{1.5cm}

\pscharpath{\Font PSTricks}
\\ \hline  
famille : ppl (Palatino)
&
famille : pag (AvantGarde)
\\ \hline  
\DeclareFixedFont{\Font}{T1}{pcr}{b}{n}{1.5cm}

\pscharpath{\Font PSTricks}

&
\DeclareFixedFont{\Font}{T1}{pnc}{b}{n}{1.5cm}

\pscharpath{\Font PSTricks}
\\ \hline 
famille : pcr (Courier)
&
famille : pnc (NewCenturySchoolbook)
\\ \hline 

\DeclareFixedFont{\Font}{T1}{psy}{b}{n}{1.5cm}

\pscharpath{\Font PSTricks}
&
\DeclareFixedFont{\Font}{T1}{pzc}{b}{n}{1.5cm}

\pscharpath{\Font PSTricks}
\\ \hline 
famille : psy (Symbol)
&
famille : pzc (ZapfChancery)
\\ \hline 

\DeclareFixedFont{\Font}{T1}{phv}{b}{n}{1.5cm}

\pscharpath{\Font PSTricks}
&
\DeclareFixedFont{\Font}{T1}{pzd}{b}{n}{1.5cm}

\pscharpath{\Font PSTricks}
\\ \hline 
famille : phv (Helvetica)
&
famille : pzd (ZapfDingbats)
\\ \hline 

\end{tabular} 
\bigskip


%---------------------------------------------------------------------------
\SbSbSSCT{Mise en forme}{Formatting}

\DeclareFixedFont{\Font}{T1}{ptm}{b}{n}{2cm}

\begin{tabular}{|c|}\hline  
 
\BS{pscharpath}[linecolor=lightgray]\AC{\BS{Font} PsTricks}
\\ \hline
\pscharpath[linecolor=lightgray]{\Font PsTricks}
\\ \hline 
\BS{pscharpath}[fillstyle=gradient,gradbegin=red,gradend=cyan,shadow=true]\AC{\BS{Font} PsTricks}
\\ \hline
\pscharpath[fillstyle=gradient,gradbegin=red,gradend=cyan,shadow=true]{\Font PsTricks}
\\ \hline 
\BS{pscharpath}[doubleline=true]\AC{\BS{Font} PsTricks}
\\ \hline
\pscharpath[doubleline=true]{\Font PsTricks}
\\ \hline 
\BS{pscharpath}[shadow=true]\AC{\BS{Font} PsTricks}
\\ \hline
\pscharpath[shadow=true]{\Font PsTricks}
\\ \hline 
\end{tabular}

\bigskip
\begin{tabular}{|c|}\hline  
\BSS{pscharpath} \TFRGB{avec astérisque}{ with asterisk}
\\ \hline
 \BSS{pscharpath*}\AC{\BS{Font} PsTricks} \BSI{pscharpath*}{pst-text} 
\\ \hline
\pscharpath*{\Font PsTricks}
\\ \hline 
\BS{pscharpath*}[linecolor=cyan]\AC{\BS{Font} PsTricks}
\\ \hline
\pscharpath*[linecolor=cyan]{\Font PsTricks}
\\ \hline 
\BS{pscharpath}[doubleline=true,linecolor=magenta]\AC{\BS{Font} PsTricks}
\\ \hline
\pscharpath[doubleline=true,linecolor=magenta]{\Font PsTricks}
\\ \hline 

%\BS{pscharpath*}[shadow=true]\AC{\BS{Font} PsTricks}
%\\ \hline
%\pscharpath*[shadow=true]{\Font PsTricks}
%\\ \hline 
\end{tabular}

%--------------------------------------------------------------------------------
\SbSbSSCT{Effets spéciaux}{Special effects}
\DeclareFixedFont{\Font}{T1}{ptm}{b}{n}{2.5cm}

\begin{tabular}{|l|} \hline  
 \BSS{psboxfill}\AC{\BS{tiny} pstricks} \BSI{psboxfill}{pst-text} \\
\BS{pscharpath}[{\red fillstyle=boxfill,fillangle=45}] 
\AC{\BS{Font} PsTricks}
\\ \hline  
 \psboxfill{\tiny pstricks}
\pscharpath[fillstyle=boxfill,fillangle=45]{\Font PsTricks}
\\ \hline 
\end{tabular} 

\bigskip



\begin{tabular}{|l|} \hline  
 \BS{DeclareFixedFont}\AC{\BS{Font}}\AC{T1}\AC{phv}\AC{b}\AC{n}\AC{2cm} \BSI{DeclareFixedFont}{pst-text}\\
\BSS{pstextpath}(0,0)\AC{% \\
\BSS{pscharpath*}[linestyle=none]\AC{\BS{Font} PsTricks}} \BSI{pstextpath}{pst-text}\\
\AC{\BS{tiny}  PsTricks PsTricks PsTricks ...}
\\ \hline  
 \DeclareFixedFont{\Font}{T1}{phv}{b}{n}{3cm}
\pstextpath(0,0){%
\pscharpath*[linestyle=none]{\Font PsTricks}}{\tiny  PsTricks PsTricks PsTricks PsTricks PsTricks PsTricks PsTricks PsTricks PsTricks PsTricks PsTricks PsTricks PsTricks PsTricks PsTricks PsTricks PsTricks PsTricks PsTricks  PsTricks PsTricks PsTricks PsTricks PsTricks PsTricks PsTricks PsTricks PsTricks PsTricks PsTricks PsTricks PsTricks PsTricks PsTricks PsTricks PsTricks PsTricks PsTricks PsTricks PsTricks PsTricks PsTricks PsTricks PsTricks PsTricks PsTricks PsTricks PsTricks PsTricks PsTricks PsTricks PsTricks PsTricks PsTricks PsTricks PsTricks PsTricks PsTricks PsTricks PsTricks PsTricks PsTricks PsTricks PsTricks PsTricks
PsTricks PsTricks PsTricks PsTricks PsTricks PsTricks PsTricks PsTricks PsTricks PsTricks PsTricks PsTricks PsTricks PsTricks
PsTricks PsTricks PsTricks PsTricks PsTricks PsTricks PsTricks PsTricks PsTricks PsTricks PsTricks PsTricks PsTricks PsTricks
PsTricks PsTricks PsTricks PsTricks PsTricks PsTricks PsTricks PsTricks PsTricks PsTricks PsTricks PsTricks PsTricks PsTricks
 PsTricks}
\\ \hline 
\end{tabular} 

%--------------------------------------------------------------------------------------
%\newpage
\subsection{pscharclip}

\begin{tabular}{|l|} \hline  
 \BS{DeclareFixedFont}\AC{\BS{Font}}\AC{T1}\AC{pcr}\AC{b}\AC{n}\AC{2cm} \\
 \BS{begin}\AC{pspicture*}(12,3) \\
 \BSS{begin\AC{pscharclip}}[doubleline=true]\AC{ \\
 \BS{rput}(6,1.5)\AC{\BS{Font} PSTricks} \\
 } \\
\BSS{end\AC{pscharclip}} \\
\BS{end}\AC{pspicture*}  \BSI{pscharclip}{pst-text} \\
\\ \hline  
 \DeclareFixedFont{\Font}{T1}{pcr}{b}{n}{2cm}
 \begin{pspicture*}(12,2.5)

\begin{pscharclip}[doubleline=true]{ \rput(6,1.25){\Font PSTricks}}
\end{pscharclip}
\end{pspicture*}
\\ \hline 
\end{tabular} 





% \DeclareFixedFont{\RM}{T1}{ptm}{b}{n}{2cm}
% \begin{pspicture*}(10,3)
%\begin{pscharclip*}[linewidth=0.1pt]{ \rput(4,1.5){\RM PSTricks}}
%\end{pscharclip}
%\end{pspicture*}

%
%\newpage
% \pscharpath{\Font PsTricks}\\
%\pscharpath[linecolor=lightgray]{\Font PsTricks}\\
%
%\pscharpath*[linecolor=lightgray]{\Font PsTricks}\\
%
%\pscharpath[fillstyle=gradient,gradbegin=red,gradend=cyan,shadow=true]{\Font PsTricks}\\
%\pscharpath*[fillstyle=gradient,gradbegin=red,gradend=cyan,shadow=true]{\Font PsTricks}\\
%
%
%\pscharpath[fillstyle=penrose*]{\Font PsTricks}\\
%\pscharpath*[fillstyle=penrose*]{\Font PsTricks}\\
%
%\pscharpath[doubleline=true]{\Font PsTricks}\\
%\pscharpath[fillstyle=gradient,gradbegin=red,gradend=cyan]{\Font PsTricks}\\
%
% \pscharpath*{\Font PsTricks}\\
%\pscharpath*[linecolor=lightgray]{\Font PsTricks}\\
%\pscharpath*[fillstyle=gradient,gradbegin=red,gradend=cyan]{\Font PsTricks}\\
%\pscharpath*[fillstyle=gradient,gradbegin=red,gradend=cyan,shadow=true]{\Font PsTricks}\\
%\pscharpath*[fillstyle=penrose*]{\Font PsTricks}\\
%\pscharpath*[doubleline=true]{\Font PsTricks}\\
%\newpage
%
%
%\DeclareFixedFont{\Font}{T1}{ptm}{b}{n}{3cm}
%%\begin{pspicture*}[showgrid=true](\linewidth,3cm)
%%\begin{pscharclip}[linewidth=1pt]{%
%%\rput(0.5\linewidth,1){\Font PSTricks}}%
%% \end{pscharclip}
%%\end{pspicture*}
%
%\newpage
%\RDD{clipping path}, 8, 9 \\
%\BSS{DeclareFixedFont}, 7 \\
%\textbf{Dimension} \\
%\BSS{linewidth}, 8, 9 \\
%\BSS{TPoffset}, 4 \\
%
%
%dvips, 3, 7 \\
%\BSS{endpscharclip}, 8 \\
%
%\textbf{Environment} \\
%minipage, 8, 9 \\
%pscharclip, 8 \\
%pscharclip*, 8 \\
%pspicture, 4, 8 \\
%pspicture*, 8 \\
%line style, 4 \\
%\BSS{linewidth}, 8, 9 \\
%
%\textbf{Macro} \\
%\BSS{DeclareFixedFont}, 7 \\
%\BSS{endpscharclip}, 8 \\
%\BSS{parbox}, 4 \\
%\BSS{psboxfill}, 10 \\
%\BSS{pscharclip}, 8, 10 \\
%\BSS{pscharclip*}, 8 \\
%\BSS{pscharpath}, 7, 8, 10 \\
%\BSS{pscharpath*}, 7 \\
%\BSS{pscustom}, 5 \\
%\BSS{pstextpath}, 4, 7 \\
%\BSS{rput}, 8 \\
%\BSS{vspace}, 4 \\
%
%minipage, 8, 9 \\
%none, 4 \\
%offset, 4 \\
%outline font, 7 \\
%Package \\
%pst-text, 3, 4 \\
%setspace, 9 \\
%\BSS{parbox}, 4 \\
%path, 3 \\
%
%\textbf{Program} \\
%dvips, 3, 7 \\
%\BSS{psboxfill}, 10 \\
%\BSS{pscharclip}, 8, 10 \\
%pscharclip, 8 \\
%\BSS{pscharclip*}, 8 \\
%pscharclip*, 8 \\
%\BSS{pscharpath}, 7, 8, 10 \\
%\BSS{pscharpath*}, 7 \\
%\BSS{pscustom}, 5 \\
%pspicture, 4, 8 \\
%pspicture*, 8 \\
%pst-text, 3, 4 \\
%\BSS{pstextpath}, 4, 7 \\
%\RDD{Rokicki}, 3, 7 \\
%%\rput, 8
%\RDD{setspace}, 9 \\
%\BSS{TPoffset}, 4 \\
%
%Value \\
%none, 4 \\
%\BSS{vspace}, 4 \\
%11

\newpage

 ========================================================================
%\section[Objets divers ]{Objets divers \cite{pstricks-add}}
\SSCT{Objets divers}{Various objects}

%\subsection{Des dés}
\SbSSCT{Des dés}{Dices}
\label{des}

\begin{tabular}{|c|c|c|c|c|c|}
\hline 
\begin{pspicture}(-.5,-.5)(.5,.7)
\psdice{1}
\end{pspicture}
 & 
\begin{pspicture}(-.5,-.5)(.5,.7)
\psdice{2}
\end{pspicture} 
 & 
\begin{pspicture}(-.5,-.5)(.5,.7)
\psdice{3}
\end{pspicture}  
& 
\begin{pspicture}(-.5,-.5)(.5,.7)
\psdice{4}
\end{pspicture}
& 
\begin{pspicture}(-.5,-.5)(.5,.7)
\psdice{5}
\end{pspicture}
&
\begin{pspicture}(-.5,-.5)(.5,.7)
\psdice{6}
\end{pspicture}
 \\ 
\hline \BSS{psdice}\AC{1}  \BSI{psdice}{pstricks-add}  &  \BSS{psdice}\AC{2} &  \BSS{psdice}\AC{3} &  \BSS{psdice}\AC{4} &  \BSS{psdice}\AC{5} &  \BSS{psdice}\AC{6} \\ 
\hline 
\end{tabular} 

\bigskip
\begin{tabular}{|c|}\hline  
\begin{pspicture}(-1,-1.2)(1,1.2)
\psdice[unit=2]{5}
\end{pspicture}
\\ \hline  
\BS{psdice}[\RDD{unit}=2]\AC{5}  \RDI{unit}{pstricks-add} 
\\ \hline 
\end{tabular} 


 


%\section{Dessins humouristiques}
\SbSSCT{Dessins humouristiques}{Fun drawing}

\label{fun}

\TFRGB{utilisation du module}{Package} \og \textbf{pst-fun} \fg{} \label{pst-fun}

\subsubsection{Commandes brutes}

\begin{tabular}{|c|c|} 	\hline  
\begin{psgraph}[axesstyle=none,xticksize=-2.5 2.5, yticksize=-1 2, subticks=0](0,0)(-1,-2.5)(2,2.5){3cm}{5cm} 
	\psBill
\end{psgraph}
& 
\begin{psgraph}[axesstyle=none,xticksize= -.5 4.5,yticksize=0 7, subticks=0](0,0)(0,-.5)(7,4.5){7cm}{5cm }
 \psFish
\end{psgraph}
\\ \hline  
\BSS{psBill}  \BSI{psBill}{pst-fun}
&  
 \BSS{psFish}  \BSI{psFish}{pst-fun}
\\ 	\hline 
\end{tabular} 



\bigskip

\begin{tabular}{|c|c|c|} 	\hline 
\begin{psgraph}[axesstyle=none,xticksize= -2.5 2.5,yticksize=-2 2, subticks=0](0,0)(-2,-2.5)(2,2.5){4cm}{5cm } 
 \psLouisXIII
\end{psgraph}
&
\begin{psgraph}[axesstyle=none,xticksize= -1.5 3.5,yticksize=0 4, subticks=0](0,0)(0,-1.5)(4,3.5){4cm}{5cm } 
\psBird 
\end{psgraph}
&
\begin{psgraph}[axesstyle=none,xticksize= -2.5 2.5,yticksize=-2 2, subticks=0](0,0)(-2,-2.5)(2,2.5){4cm}{5cm } 
 \psAnt 
\end{psgraph}
\\ 	\hline  
\BSS{psLouisXIII}  \BSI{psLouisXIII}{pst-fun}
&  
\BSS{psBird} 
&  \BSS{psAnt}  \BSI{psAnt}{pst-fun}
\\  \hline 
\end{tabular} 

\bigskip

\begin{tabular}{|c|} 	\hline 
\begin{psgraph}[axesstyle=none,xticksize= 0 6 ,yticksize=0 11, subticks=0](0,0)(0,0)(11,6){11cm}{6cm } 
\psPulpo 
\end{psgraph}
\\  \hline  
\BSS{psPulpo}   \BSI{psPulpo}{pst-fun}
\\ 	\hline 
\end{tabular} 

\bigskip

\begin{tabular}{|c|} 	\hline  
\begin{psgraph}[axesstyle=none,xticksize= -.5 10.5 ,yticksize=0 8, subticks=0](0,0)(0,-.5)(8,10.5){8cm}{11cm } 
\psLuke
\end{psgraph}
 \\ 	\hline  
  \BSS{psLuke}  \BSI{psLuke}{pst-fun}
 \\ 	\hline 
\end{tabular}
 
\bigskip
\begin{tabular}{|c|c|} 	\hline  
\begin{psgraph}[axesstyle=none,xticksize= 0 7 ,yticksize=0 6, subticks=0](0,0)(0,0)(6,7){6cm}{7cm }
 \psParrot{1}
\end{psgraph}
&
\begin{psgraph}[axesstyle=none,xticksize= 0 7 ,yticksize=0 6, subticks=0](0,0)(0,0)(6,7){6cm}{7cm }
\psParrot{.5}
\end{psgraph}
 \\ 	\hline  
 \BSS{psParrot}\AC{1}  \BSI{psParrot}{pst-fun}
 &  
\BSS{psParrot}\AC{.5}
 \\ 	\hline 
\end{tabular} 
 
\bigskip
\begin{tabular}{|c|c|} 	\hline  
\begin{psgraph}[axesstyle=none,xticksize= 0 5 ,yticksize=-3 3, subticks=0](0,0)(-3,0)(3,5){6cm}{5cm }
\psKangaroo{1}
\end{psgraph}
&
\begin{psgraph}[axesstyle=none,xticksize= 0 5 ,yticksize=-3 3, subticks=0](0,0)(-3,0)(3,5){6cm}{5cm }
\psKangaroo{5}
\end{psgraph}
\\ 	\hline  
\BSS{psKangaroo}\AC{1}  \BSI{psKangaroo}{pst-fun}
&  
\BSS{psKangaroo}\AC{5}
\\ 	\hline 
\end{tabular}
 
  
\bigskip
\begin{tabular}{|c|c|} 	\hline  
\begin{psgraph}[axesstyle=none,xticksize= -2 3 ,yticksize=-2 3, subticks=0](0,0)(-2,-2)(3,3){5cm}{5cm }
\psPig(0,0)
\end{psgraph}
&
\begin{psgraph}[axesstyle=none,xticksize= -2 3 ,yticksize=-2 3, subticks=0](0,0)(-2,-2)(3,3){5cm}{5cm }
 \psPig(1,1)
\end{psgraph}
\\ 	\hline  
  	\BSS{psPig}(0,0) \BSI{psPig}{pst-fun}
  	&  
  	\BSS{psPig}(1,1)
  	\\ 	\hline 
  \end{tabular}
 
 

  
 
\subsubsection{options}

\begin{tabular}{|c|c|} \hline  
\begin{psgraph}[axesstyle=none,xticksize= -2.5 2.5,yticksize=-2 3, subticks=0](0,0)(-2,-2.5)(3,2.5){5cm}{5cm } 
	\psAnt[fillcolor=red] 
\end{psgraph}
& 
\begin{psgraph}[axesstyle=none,xticksize= -2.5 2.5,yticksize=-2 3, subticks=0](0,0)(-2,-2.5)(3,2.5){5cm}{5cm } 
	\psAnt[fillstyle=slope] 
\end{psgraph}
\\ 	\hline  
\BS{psAnt}[{\red fillcolor}=red] &  \BS{psAnt}[{\red fillstyle}=slope]
\\ 	\hline
\end{tabular} 


\bigskip
\begin{tabular}{|c|c|} 	\hline 
\begin{psgraph}[axesstyle=none,xticksize= -2.5 3.5,yticksize=-2 3, subticks=0](0,0)(-2,-2.5)(3,3.5){5cm}{6cm } 
	\rput(1,1){\psAnt} 
\end{psgraph}
&  
\begin{psgraph}[axesstyle=none,xticksize= -2.5 3.5,yticksize=-2 3, subticks=0](0,0)(-2,-2.5)(3,3.5){5cm}{6cm }  
	\rput{-60}{\psAnt} 
\end{psgraph}
\\ 	\hline
\BSS{rput} {\red (1,1)}\AC{\BS{psAnt}} &  \BSS{rput}{\red \AC{-60}}\AC{\BS{psAnt}}
\\ 	\hline 
\end{tabular} 

\bigskip
\begin{tabular}{|c|c|} 	\hline
\begin{psgraph}[axesstyle=none,xticksize= -1.5 2.5,yticksize=-1 2, subticks=0](0,0)(-1,-1.5)(2,2.5){3cm}{4cm } 
	\psscalebox{0.5}{\psAnt} 
\end{psgraph}
&  
\begin{psgraph}[axesstyle=none,xticksize= -1.5 2.5,yticksize=-1 2, subticks=0](0,0)(-1,-1.5)(2,2.5){3cm}{4cm }  
	\rput(1,1){\psscalebox{0.5}{\psAnt}} 
\end{psgraph}
\\ 	\hline
\BSS{psscalebox}\AC{0.5}\AC{\BS{psAnt}}  &  \BS{rput}\AC{\BSS{psscalebox}\AC{0.5}\AC{\BS{psAnt}} }
\\ 	\hline 
\end{tabular}

\bigskip
\begin{tabular}{|c|c|} 	\hline  
	\begin{psgraph}[axesstyle=none,xticksize= -2 3 ,yticksize=0 10, subticks=0](0,0)(0,-2)(10,3){10cm}{5cm}
%\begin{pspicture}[showgrid](0,-1.3)(10,3)
 \psBird[Branch] % \rput{-20}(4,1.8){\psBird}
\end{psgraph}
\\ 	\hline 
 \BS{psBird}[\RDD{Branch}]  \RDI{Branch}{pst-fun}
\\ 	\hline 
\end{tabular}


  
\bigskip
\begin{tabular}{|c|c|} 	\hline  
\begin{psgraph}[axesstyle=none,xticksize= -2 2 ,yticksize=-2 2, subticks=0](0,0)(-2,-2)(2,2){4cm}{4cm}
\psPig[ eyeColor=red](0,0)
\end{psgraph}
&
\begin{psgraph}[axesstyle=none,xticksize= -2 2 ,yticksize=-2 2, subticks=0](0,0)(-2,-2)(2,2){4cm}{4cm }
\psPig[noseColor=yellow](0,0)
\end{psgraph}
\\ 	\hline  
\BS{psPig}[\RDD{eyeColor}=red](0,0) \RDI{eyeColor}{pst-fun}
&  
\BS{psPig}[\RDD{noseColor}=yellow](0,0) \RDI{noseColor}{pst-fun}
\\ 	\hline 
  \end{tabular}

\newpage
%======================================================================

%\section[Créer un graphe]{Créer un graphe \cite{pst-user}  \cite{pst-plot}}
\SSCT{Créer un graphe}{Create a graph}

Utilisation du module \textbf{pst-plot} \label{pst-plot}

%\section{Syntaxe}
\psset{fillstyle=none,unit=1cm}

\subsection{Environnement}

\SbSbSSCT{Dans un environnement classique }{pspicture}
 
\begin{itemize}
\item  Axes : Macro \BSS{psaxes} \BSI{psaxes}{pst-plot} 
\item  Quadrillages : Macro  \BSS{psgrid}  \BSI{psgrid}{pst-plot} 
\end{itemize}

\SbSbSSCT{Dans un environnement psgraph}{psgraph}

\psset{linecolor=red,arrowscale=2}



\psset{linecolor=blue}
\TFRGB{Deux syntaxes}{Two syntaxes} :

\BSS{psgraph}[Options]  \AC{\TFRGB{flèches}{arrows}}(xOrig,yOrig)(xMin,yMin)(xMax,yMax)\AC{\TFRGB{largeur graphe}{graph width}}
\AC{\TFRGB{hauteur graphe}{graph height}} 
\BSS{endpsgraph} 
\BSI{psgraph}{pst-plot}  \BSI{endpsgraph}{pst-plot} 

\TFRGB{ou}{or}

\BSS{begin\AC{psgraph}} [Options]\AC{\TFRGB{flèches}{arrows}}(xOrig,yOrig)(xMin,yMin)(xMax,yMax)
\AC{\TFRGB{largeur graphe}{graph width}}\AC{\TFRGB{hauteur graphe}{graph height}}
. . .
\BSS{end\AC{psgraph}} 

\bigskip

\TFRGB{
Remarque : 
\begin{itemize}
\item  L'indication de la largeur et de la hauteur du graphe permettent la mise à l'échelle automatique
\item Si hauteur graphe = ! , les deux axes ont la même unité
\end{itemize}}{
Remarks : 
\begin{itemize}
\item  The indication of the width and height of the graph allow automatic scaling
\item If graph  height =! , Both axes have the same unit
\end{itemize}}

%======================================

\SbSSCT{Type de tracé}{Type of graph}

%Option : \RDD{plotstyle} 

\readdata{\dat}{table.dat}
\psset{xunit=.4cm,yunit=.3cm}
\psset{linecolor=blue,linewidth=2pt}


\begin{tabular}{|c|c|c|} \hline  
\begin{pspicture}(-1.5,-1.5)(8.5,11)
\psaxes(0,0) (8,10)
\listplot[xticksize=0 10,yticksize= 0 8,Dx=1,dy=1,Dy=1,linecolor=red]{\dat}
\end{pspicture} 
& 
\begin{pspicture}(-1.5,-1.5)(8.5,11)
\psaxes(0,0) (8,10)
\listplot[plotstyle=dots,xticksize=0 10,yticksize= 0 8,Dx=1,dy=1,Dy=1,linecolor=red]{\dat}
\end{pspicture}  
&
\begin{pspicture}(-1.5,-1.5)(8.5,11)
\psaxes(0,0) (8,10)
\listplot[plotstyle=polygon,xticksize=0 10,yticksize= 0 8,Dx=1,dy=1,Dy=1,linecolor=red]{\dat}
\end{pspicture} 
 \\  \hline 
{\blue \dft }: \RDD{plotstyle}=  \BDD{line} \BDI{line}{pst-plot}  \RDI{plotstyle}{pst-plot} & 
{\red plotstyle}= \BDD{dots} \BDI{dots}{pst-plot} & 
{\red plotstyle}=  \BDD{polygon} \BDI{polygon }{pst-plot} \\ 
\hline 
\end{tabular}

\bigskip
\begin{tabular}{|c|c|c|} \hline  
\begin{pspicture}(-1.5,-1.5)(8.5,11)
\psaxes(0,0) (8,10) 
\listplot[plotstyle=curve,xticksize=0 10,yticksize= 0 8,Dx=1,dy=1,Dy=1,linecolor=red]{\dat}
\end{pspicture} 
& 
\begin{pspicture}(-1.5,-1.5)(8.5,11)
\psaxes(0,0) (8,10) 
\listplot[plotstyle=ecurve,xticksize=0 10,yticksize= 0 8,Dx=1,dy=1,Dy=1,linecolor=red]{\dat}
\end{pspicture} 
&
\begin{pspicture}(-1.5,-1.5)(8.5,11)
\psaxes(0,0) (8,10) 
\listplot[plotstyle=ccurve,xticksize=0 10,yticksize= 0 8,Dx=1,dy=1,Dy=1,linecolor=red]{\dat}
\end{pspicture}   
 \\  \hline 
{\red plotstyle}= \BDD{curve} \BDI{curve}{pst-plot} &
{\red plotstyle}=\BDD{ecurve} \BDI{ecurve}{pst-plot} & 
 {\red plotstyle}= \BDD{ccurve} \BDI{ccurve}{pst-plot} \\ 
\hline 
\end{tabular}

\bigskip
\begin{tabular}{|c|c|c|} \hline  
\begin{pspicture}(-1.5,-1.5)(8.5,11)
\psaxes(0,0) (8,10)
\listplot[plotstyle=bar,linecolor=red]{\dat}
\end{pspicture} 
& 
\begin{pspicture}(-1.5,-1.5)(8.5,11)
\psaxes(0,0) (8,10)
\listplot[plotstyle=ybar,linecolor=red]{\dat}
\end{pspicture}  
&
\begin{pspicture}(-1.5,-1.5)(8.5,11)
\psaxes(0,0) (8,10)
\listplot[plotstyle=dots,linecolor=red]{\dat} 
\listplot[plotstyle=xvalues]{\dat}
\end{pspicture} 
 \\  \hline 
{\red plotstyle}= \BDD{bar} \BDI{bar}{pst-plot} &
 {\red plotstyle}= \BDD{ybar} \BDI{ybar}{pst-plot} & 
 {\red plotstyle}= \BDD{xvalues } \BDI{xvalues }{pst-plot}\\ 
\hline 
\end{tabular}

\bigskip
\begin{tabular}{|c|c|c|} \hline  
\begin{pspicture}(-1.5,-1.5)(8.5,11)
\psaxes(0,0) (8,10) 
\listplot[plotstyle=LineToXAxis,linecolor=red]{\dat}
\end{pspicture} 
& 
\begin{pspicture}(-1.5,-1.5)(8.5,11)
\psaxes(0,0) (8,10) 
\listplot[plotstyle=LineToYAxis,linecolor=red]{\dat}
\end{pspicture} 
&
\begin{pspicture}(-1.5,-1.5)(8.5,11)
\psaxes(0,0) (8,10)
\listplot[plotstyle=dots,linecolor=red]{\dat} 
\listplot[plotstyle=values]{\dat}
\end{pspicture}   
 \\  \hline 
{\red plotstyle}= \BDD{LineToXAxis} \BDI{LineToXAxis}{pst-plot} & 
{\red plotstyle}= \BDD{LineToYAxis} \BDI{LineToYAxis}{pst-plot} & 
{\red plotstyle}=  \BDD{values} \BDI{values}{pst-plot} \\ 
\hline 
\end{tabular}

\bigskip
\begin{tabular}{|c|c|c|} \hline  
\begin{pspicture}(-1.5,-1.5)(8.5,11)
\psaxes(0,0) (8,10)

\listplot[Hue=300,plotstyle=colordots]{\dat}
\end{pspicture} 
& 
\begin{pspicture}(-1.5,-1.5)(8.5,11)
\psaxes(0,0) (8,10) 
\listplot[plotstyle=dots,linecolor=red]{\dat} 
\listplot[plotstyle=LSM,linecolor=magenta]{\dat}
\end{pspicture} 
&
\begin{pspicture}(-1.5,-1.5)(8.5,11)
\psaxes(0,0) (8,10)
\listplot[shadow=true,plotstyle=bar]{\dat}
\end{pspicture}   
 \\  \hline 
Hue=100,{\red plotstyle}= \BDD{colordots} \BDI{colordots}{pst-plot} & 
{\red plotstyle}= \BDD{LSM} \BDI{LSM}{pst-plot}  & shadow=true,{\red plotstyle}=\BDD{bar} \BDI{bar}{pst-plot} \\ 
\hline 
\end{tabular}


%=================================================

\SbSSCT{Les axes }{Axes}
\SbSbSSCT{Dimensionnement}{Dimensioning}
 
\psset{unit=1cm,linewidth=1pt}

\begin{tabular}{|c|c|} \hline 
\BSS{psaxes}\AC{<->}\rnode[fillcolor=yellow]{AA}{(0,0)}\rnode[fillcolor=green]{BB}{(-1,-2)}\rnode[fillcolor=cyan]{CC}{(3,3)}
&
\textbf{\BS{psaxes}}\AC{->}\rnode[fillcolor=cyan]{DD}{(4,2)} \\
& \\ \hline
 
\begin{pspicture}[shift=*](-1.5,-3)(4.5,3.5)
\psaxes[linecolor=black,linewidth=2pt]{<->}(0,0)(-1,-2)(3,3)
\dotnode[dotstyle=*](0,0){A}
\dotnode[dotstyle=*](-1,-2){B}
\dotnode[dotstyle=*](3,3){C}
\ncarc[angleB=135]{->}{AA}{A}
\ncbar[angleA=-90,angleB=135]{->}{BB}{B}
\ncline{->}{CC}{C}
\end{pspicture}
&  
\begin{pspicture}[shift=*](-.5,-.5)(5,2.5)
\psaxes[linecolor=black,linewidth=2pt]{->}(4.5,2.5)
\dotnode[dotstyle=*](4,2){D}
\ncline{->}{DD}{D}
\end{pspicture}
\\ \hline  
\end{tabular} 

%----------------------------------------------------------
\psset{unit=.7cm}
\subsubsection{Types d'axes}


\psset{unit=.7cm}

\begin{tabular}{|c|c|c|}  \hline 
\begin{pspicture}(-1,-1)(4,4)
\psaxes[axesstyle=none]{->}(3,3)
\end{pspicture} 
&
 \begin{pspicture}(-1,-1)(4,4)  
 \psaxes[axesstyle=frame]{->}(3,3)
\end{pspicture} 
&  
\begin{pspicture}(-3,-3)(3,3) 
 \psaxes[axesstyle=polar]{->}(2,2)
\end{pspicture} \\ 
 \hline 
\RDD{axesstyle}=none \RDI{axesstyle}{pst-plot} & {\red axesstyle}=frame &  {\red axesstyle}=polar\\ \hline 
 \multicolumn{3}{|c|}{ \blue \dft : {\red axesstyle}=axes  } \\  \hline  

\end{tabular} 


%--------------------------------------------------
\SbSbSSCT{Choix des axes}{choice of axes}


\psset{unit=0.5cm}
\begin{tabular}{|c|c|c|c|} \hline
 \begin{pspicture}(-1,-1)(4,4)  
 \psaxes[xyAxes=true]{->}(3,3)
\end{pspicture}  
& 
 \begin{pspicture}(-1,-1)(4,4)  
 \psaxes[xyAxes=false]{->}(3,3)
\end{pspicture}  
&  \begin{pspicture}(-1,-1)(4,4)  
 \psaxes[xAxis=false]{->}(3,3)
\end{pspicture} 
&  \begin{pspicture}(-1,-1)(4,4) 
 \psaxes[yAxis=false]{->}(3,3)
\end{pspicture} \\ \hline 
 \RDD{xyAxes}=true  \RDI{xyAxes}{pst-plot}  &  xyAxes=false &  \RDD{xAxis}=false  \RDI{xAxis}{pst-plot} &  \RDD{yAxis}=false  \RDI{yAxis}{pst-plot} \\ \hline 
 \multicolumn{4}{|c|}{\blue \dft : xAxis = yAxis = xyAxes = true } \\  \hline  
\end{tabular} 

%----------------------------------------------------------------

\SbSbSSCT{Espacement des graduations}{Units of the axis}


\psset{xunit=.9cm,yunit=.9cm}
\begin{tabular}{|c|c|c|} \hline 

\begin{pspicture}(-1,-1)(3.5,3.5) 
 \psaxes[Ox=2,Dx=2.0,Dy=.5]{->}(3,3)
\end{pspicture}  
&  
\begin{pspicture}(-1,-1)(3.5,3.5)
 \psaxes[Oy=2,dx=.5,dy=1.5]{->}(3,3)
\end{pspicture} 
&  
\begin{pspicture}(-1,-1)(3.5,3.5) 
 \psaxes[dx=.5,Dx=2,dy=.5,Dy=1.5]{->}(3,3) 
\end{pspicture} \\  \hline  
\RDD{Ox}=2 	 \RDI{Ox}{pst-plot} 	&  
\RDD{Oy}=2   \RDI{Oy}{pst-plot} 	& {\red dx}=.5 {\red Dx}=2   \\ 
\RDD{Dx}=2.0  \RDI{Dx}{pst-plot} 	&  
\RDD{dx}=.5  \RDI{dx}{pst-plot} 	&  
{\red dy}=.5 {\red Dy}= 1.5 \\ 
\RDD{Dy}=0.5  \RDI{Dy}{pst-plot} 	&  
\RDD{dy}=1.5 \RDI{dy}{pst-plot} 	&  \\
\hline 
\multicolumn{3}{|c|}{ \blue \dft{}:  Ox = Oy= 0 \hspace{1cm} Dx= Dy = 1 }
\\ \hline

\end{tabular} 


%====================================================
\SbSbSSCT{Origine}{Origin}


\begin{tabular}{|c|c|}\hline 
\begin{pspicture}(-1,-1)(4,4) 
\psaxes[showorigin=true]{->}(3,3)
\end{pspicture} 
&
\begin{pspicture}(-1,-1)(4,4) 
\psaxes[showorigin=false]{->}(3,3)
\end{pspicture} \\
\hline  
\RDD{showorigin}=true  \RDI{showorigin}{pst-plot} 	{\blue (\dft)}	& {\red showorigin}=false  	 \\ 
\hline 
\end{tabular}


%-----------------------------------------------------------------------
\SbSbSSCT{Titres des axes}{Titles on axes}

%\emph{Commande simple :}

%syntaxe :\BS{}psaxes\AC{flèche}(coordonnée origine)(coordonnée maximale)[titre axe X,position angulaire][titre axe Y,position angulaire]
%\smallskip




\begin{tabular}{|c|} \hline 

\BS{}psaxes\AC{->}(0,0)(6,5)[$X$,-90][$Y$,180]
\\ \hline 
\begin{pspicture*}(-1,-1)(7,3)
\psaxes{->}(0,0)(6,2.5)[{\red $X$},-90][{\red$Y$},180]
\end{pspicture*}
 
 \\  \hline 
\end{tabular} 


%\bigskip
%\emph{Options avancées}

%\smallskip
%\begin{tabular}{|c|c|c|l|} \hline
%nom 			& type						& \dft 	&   Rôle \\\hline 
%\RDD{xAxisLabel}		& texte						&  				& titre axe X \\
%\RDD{yAxisLabel}		& texte						& 				& titre axe Y 	 \\
%\RDD{xAxisLabelPos}	& (x,y) (centré si x= c)	&  				& position du titre axe X \\
%\RDD{yAxisLabelPos}	& (x,y) (centré si y= c)	& 				&  position du titre axe Y \\ \hline
%\multicolumn{4}{|c|}{correction taille de la boîte contenant le graphique} \\
%\multicolumn{4}{|c|}{A définir avant la commande \BS{psgraph} grâce à la commande \BS{psset}}
% \\\hline
%\RDD{urx}				& dimension					& 0pt			& coin supérieur droit en Y \\
%\RDD{ury}				& dimension					& 0pt			& coin supérieur droit en X  \\
%\RDD{llx}				& dimension					& 0pt 			& coin inférieur gauche en X \\
%\RDD{lly}				& dimension	 				& 0pt			& coin inférieur gauche en Y \\
%\hline 
%\end{tabular} 
%
%\bigskip

\begin{tabular}{|c|c|} \hline 
\multicolumn{2}{|c|}{\BS{psset}\AC{\RDD{llx}=0,\RDD{lly}=0,\RDD{urx}=0,\RDD{ury}=0,\RDD{xAxisLabel}=X,\RDD{yAxisLabel}=titre axe Y,\RDD{yAxisLabelPos}=\AC{-1cm,c}}   
}\\
\multicolumn{2}{|c|}{\RDI{llx}{pst-plot}  \RDI{lly}{pst-plot} \RDI{urx}{pst-plot}  \RDI{ury}{pst-plot}
\RDI{xAxisLabel}{pst-plot}  \RDI{yAxisLabel}{pst-plot}  
\RDI{yAxisLabelPos}{pst-plot}}
\\
& \\ \hline
\rule{6cm}{0pt} & \rule{6cm}{0pt} \\


\psframebox[linestyle=dashed,linecolor=red]{
\psset{llx=0,lly=0,urx=0,ury=0,xAxisLabel=X,yAxisLabel=titre axe Y,yAxisLabelPos={-1cm,c}}
\begin{psgraph}[axesstyle=frame,](0,0)(3,3){4cm}{3cm}%
\end{psgraph}} 
& 
\psframebox[linestyle=dashed,linecolor=red]{
\psset{llx=-1cm,lly=-1.25cm,urx=.5cm,ury=.51cm,xAxisLabel=titre axe X, yAxisLabel=Y,xAxisLabelPos={c,-1cm}}
\begin{psgraph}[axesstyle=frame,](0,0)(3,3){4 cm}{3cm}%
\end{psgraph}} \\ 
\rule[-1cm]{6cm}{0pt} 		& \rule{6cm}{0pt} 
\\ \hline 
{\red xAxisLabel}=X 					& {\red xAxisLabel}=titre axe X  \\
{\red yAxisLabel}=titre axe Y 		& {\red yAxisLabel}= Y \\
{\red llx}=0 						& {\red llx}=-1cm   \\
{\red lly}=0 						& {\red lly}=-1.25cm  \\
{\red urx}=0 						& {\red urx}=.5cm   \\
{\red ury}=0 						& {\red ury}=.5cm  \\ 
{\red yAxisLabelPos}=\AC{-1cm,c} 	& \RDD{xAxisLabelPos}=\AC{c,-1cm} \RDI{xAxisLabelPos}{pst-plot} \\
\hline
\end{tabular}



\psset{xAxisLabel={} ,yAxisLabel={}}



%-----------------------------------------------------------
\SbSSCT{Marques de graduations}{Ticks marks}

\SbSbSSCT{Style des marques de graduation}{Style of the tick marks}


\psset{yunit=0.7cm,xunit=0.7cm}
 
\begin{tabular}{|c|c|c|c|} \hline  
  \begin{pspicture}(-1,-1)(3.5,3.5)
  \psaxes[tickstyle=full]{->}(3,3)
  \end{pspicture}
& 
\begin{pspicture}(-1,-1)(3.5,3.5) 
\psaxes[tickstyle=top]{->}(3,3)
\end{pspicture}
& 
\begin{pspicture}(-1,-1)(3.5,3.5)
\psaxes[tickstyle=bottom]{->}(3,3)
\end{pspicture} 
& 
\begin{pspicture}(-1,-1)(3.5,3.5)  
\psaxes[axesstyle=frame,tickstyle=inner]{->}(3,3)
\end{pspicture} 
 \\  \hline 
\RDD{tickstyle}=full {\blue (\dft)} \RDI{tickstyle}{pst-plot} & 
{\red tickstyle}=top  & {\red tickstyle}=bottom 	& {\red tickstyle}=inner 	\\
							& 						&							& axesstyle=frame 			\\ \hline
\end{tabular} 



\SbSbSSCT{Présence des marques de graduation}{Ticks on axes}
%--------------------------------------------------
 
\psset{yunit=0.7cm,xunit=0.7cm}

\begin{tabular}{|c|c|c|c|} \hline 
\begin{pspicture}(-1,-1)(3.5,3.5) 
\psaxes[ticks=all]{->}(3,3)
\end{pspicture} 
& 
\begin{pspicture}(-1,-1)(3.5,3.5) 
\psaxes[ticks=x]{->}(3,3)
\end{pspicture} 
& 
  \begin{pspicture}(-1,-1)(3.5,3.5)  
\psaxes[ticks=y]{->}(3,3)
\end{pspicture} 
& 
\begin{pspicture}(-1,-1)(3.5,3.5) 
\psaxes[ticks=none]{->}(3,3)
\end{pspicture}  \\  \hline 
 \RDD{ticks}=all {\blue ( \dft )} \RDI{ticks}{pst-plot} & {\red ticks}=x & {\red ticks}=y & {\red ticks}=none \\ \hline 
\end{tabular} 


%--------------------------------------------------

\SbSbSSCT{Taille des graduations}{Size of ticks marks}
 


\begin{tabular}{|c|c|c|} \hline  
  \begin{pspicture}[ticksize=-5pt 10pt](-1,-1)(4,4) 
 \psaxes{->}(3,3)
 \end{pspicture} 
& 
  \begin{pspicture}(-1,-1)(4,4) 
\psaxes[xticksize=2,yticksize=2](3,3)
\end{pspicture} 
&
  \begin{pspicture}(-1,-1)(4,4) 
\psaxes[xticksize=5pt 0pt,yticksize=10pt 0pt](3,3)
\end{pspicture} 
 \\  \hline 
\RDD{ticksize}=5pt 10pt  \RDI{ticksize}{pst-plot} 	& \RDD{xticksize}=2  \RDI{xticksize}{pst-plot} 	& {\red xticksize}=5pt 0pt  \\ \hline 
							& {\red yticksize}=2	& \RDD{yticksize}=10pt 0pt  \RDI{yticksize}{pst-plot} \\ \hline 
\multicolumn{3}{|c|}{\blue  \dft : ticksize = xticksize = yticksize = -4pt 4pt} \\  \hline  
\end{tabular}


%--------------------------------------------------

\psset{unit=.7cm}
\SbSbSSCT{\'Epaisseur des graduations}{Width of ticks marks}
 
\begin{tabular}{|c|c|c|} \hline  
  \begin{pspicture}[subticks=3,tickwidth=10pt](-1,-1)(4,4) 
 \psaxes{->}(3,3)
 \end{pspicture} 
& 
  \begin{pspicture}(-1,-1)(4,4) 
\psaxes[subticks=3,subtickwidth=5pt]{->}(3,3)
\end{pspicture} 
&
  \begin{pspicture}(-1,-1)(4,4) 
\psaxes[subticks=3,tickwidth=1em,subtickwidth=1ex]{->}(3,3)
\end{pspicture} 
 \\  \hline 
\RDD{tickwidth}=10pt  \RDI{tickwidth}{pst-plot} & \RDD{subtickwidth}=5pt  \RDI{subtickwidth}{pst-plot} & 
{\red tickwidth}=1em \\ 
  &  & {\red subtickwidth}=1ex \\ \hline 
\multicolumn{3}{|c|}{\blue \dft : tickwidth = subtickwidth = 0.5\BS{pslinewidth}} \\  \hline  
\end{tabular}



\SbSbSSCT{Nombre de graduations secondaires }{Number of subticks}
%--------------------------------------------------
 
\begin{tabular}{|c|c|c|} \hline  
  \begin{pspicture}[subticks=3](-1,-1)(4,4) 
 \psaxes{->}(3,3)
 \end{pspicture} 
& 
  \begin{pspicture}(-1,-1)(4,4) 
\psaxes[xsubticks=3]{->}(3,3)
\end{pspicture} 
&
  \begin{pspicture}(-1,-1)(4,4) 
\psaxes[ysubticks=3]{->}(3,3)
\end{pspicture} 
 \\  \hline 
\RDD{subticks}=2  \RDI{subticks}{pst-plot} & \RDD{xsubticks}=2  \RDI{xsubticks}{pst-plot} & \RDD{ysubticks}=2  \RDI{ysubticks}{pst-plot} \\ \hline 
\multicolumn{3}{|c|}{\blue \dft  subticks = xsubticks = ysubticks = 0 } \\  \hline 

\end{tabular}


%--------------------------------------------------
\SbSbSSCT{Tailles des marques de graduation secondaires / principales }{Size of subticks}
 
\begin{tabular}{|c|c|c|} \hline  
  \begin{pspicture}[subticks=3,subticksize	=1](-1,-1)(4,4) 
 \psaxes{->}(3,3)
 \end{pspicture} 
& 
  \begin{pspicture}(-1,-1)(4,4) 
\psaxes[xsubticks=3,xsubticksize=.5]{->}(3,3)
\end{pspicture} 
&
  \begin{pspicture}(-1,-1)(4,4) 
\psaxes[ysubticks=3,ysubticksize=2 ]{->}(3,3)
\end{pspicture} 
 \\  \hline 
\RDD{subticksize}=1   \RDI{subticksize}{pst-plot} & \RDD{xsubticksize}=.5  \RDI{xsubticksize}{pst-plot}  & \RDD{ysubticksize}=2  \RDI{ysubticksize}{pst-plot} \\ \hline 
\multicolumn{3}{|c|}{\blue \dft{} : subticksize = xsubticksize = subticksize = 0.75 } \\  \hline 

\end{tabular}



%--------------------------------------------------
\SbSbSSCT{Couleurs des marques de graduation}{Color of tick marks}
 
\begin{tabular}{|c|c|c|} \hline  
  \begin{pspicture}[subticks=3,tickcolor=red,subtickcolor=green](-1,-1)(4,4) 
 \psaxes{->}(3,3)
 \end{pspicture} 
& 
  \begin{pspicture}(-1,-1)(4,4) 
\psaxes[xsubticks=3,xtickcolor=red,xsubtickcolor=green]{->}(3,3)
\end{pspicture} 
&
  \begin{pspicture}(-1,-1)(4,4) 
\psaxes[ysubticks=3,ytickcolor=red,ysubtickcolor=green]{->}(3,3)
\end{pspicture} 
 \\  \hline 
\RDD{tickcolor}=red 	 \RDI{tickcolor}{pst-plot} 	& \RDD{xtickcolor}=red 	 \RDI{xtickcolor}{pst-plot} 	& \RDD{ytickcolor}=red 	 \RDI{ytickcolor}{pst-plot} 	\\ 
\RDD{subtickcolor}=green   \RDI{subtickcolor}{pst-plot} 	& \RDD{xsubtickcolor}=green \RDI{xsubtickcolor}{pst-plot}  & \RDD{ysubtickcolor}=green  \RDI{ysubtickcolor}{pst-plot} \\ \hline 
\multicolumn{3}{|c|}{ \blue \dft{} : tickcolor = xtickcolor = ytickcolor = black } \\ 
\multicolumn{3}{|c|}{ \blue subtickcolor = xsubtickcolor = ysubtickcolor = darkgray } \\  \hline 
\end{tabular}

%--------------------------------------------------
\SbSbSSCT{Style des marques de graduation}{Style of ticks}

\psset{unit=.7cm}

\begin{tabular}{|c|c|c|} \hline  
  \begin{pspicture}(-1,-1)(4,4) 
 \psaxes[ticksize=3,subticksize=1,subticks=3,yticklinestyle=dashed,xticklinestyle	=dotted]
(3,3)
 \end{pspicture} 
& 
  \begin{pspicture}(-1,-1)(4,4) 
\psaxes[ticksize=3,subticksize=1,subticks=3,xsubticklinestyle=solid,ysubticklinestyle=none ](3,3)
\end{pspicture} 
&
  \begin{pspicture}(-1,-1)(4,4) 
\psaxes[ticksize=3,,subticksize=1,subticks=3,ticklinestyle= dotted,subticklinestyle=dashed	 ](3,3)
\end{pspicture} 
 \\  \hline 
\RDD{yticklinestyle}=dashed 
\RDI{yticklinestyle}{pst-plot} &

 \RDD{xsubticklinestyle}=solid \RDI{xsubticklinestyle}{pst-plot} &
  
\RDD{ticklinestyle}= dotted  
\RDI{ticklinestyle}{pst-plot} \\
 
\RDD{xticklinestyle}=dotted 
\RDI{xticklinestyle}{pst-plot} &

\RDD{ysubticklinestyle}=none 
\RDI{ysubticklinestyle}{pst-plot} &

 \RDD{subticklinestyle}=dashed  \RDI{subticklinestyle}{pst-plot} \\ 
\hline 
\multicolumn{3}{|c|}{\blue \dft{} : ticklinestyle = xticklinestyle = yticklinestyle = solid }\\
\multicolumn{3}{|c|}{\blue subticklinestyle = xsubticklinestyle =  ysubticklinestyle = solid } \\ \hline 
\multicolumn{3}{|c|}{ \textbf{Option} : solid/dashed/dotted/none } \\ \hline 
\end{tabular}




%--------------------------------------------------
\SbSSCT{\'Etiquettes de graduation}{Labels on axis}
\SbSbSSCT{\'Etiquettes}{Choice of axis}


\psset{unit=0.5cm}

\begin{tabular}{|c|c|c|c|}
\hline
 \begin{pspicture}(-1,-1)(4,4)  
 \psaxes[labels=all]{->}(3,3)
\end{pspicture}  
& \begin{pspicture}(-1,-1)(4,4)  
 \psaxes[labels=x]{->}(3,3)
\end{pspicture}  
&  \begin{pspicture}(-1,-1)(4,4)  
 \psaxes[labels=y]{->}(3,3)
\end{pspicture} 
&  \begin{pspicture}(-1,-1)(4,4) 
 \psaxes[labels=none]{->}(3,3)
\end{pspicture} \\ 
\hline  
\RDD{labels}= all {\blue (\dft)} \RDI{labels}{pst-plot} & 
{\red labels}=x  & {\red labels}=y  & {\red labels}=none \\  \hline 
\end{tabular} 


%--------------------------------------------------
\SbSSCT{Position des étiquettes}{Position of labels}


\psset{unit=0.7cm}
\begin{tabular}{|c|c|c|}
\hline \begin{pspicture}(-1,-1)(4,4)  
 \psaxes[ylabelPos=axis, xlabelPos=axis,tickcolor=red]{->}(3,3)
\end{pspicture}  
&  \begin{pspicture}(-1,-1)(4,4)  
 \psaxes[ylabelPos=right, xlabelPos=bottom,tickcolor=red]{->}(3,3)
\end{pspicture} 
&  \begin{pspicture}(-1,-1)(4,4) 
 \psaxes[ ylabelPos=left, xlabelPos=top,tickcolor=red]{->}(3,3)
\end{pspicture} \\ 
\hline  \RDD{xlabelPos}=axis  \RDI{xlabelPos}{pst-plot} 
&  {\red xlabelPos}=bottom {\blue (\dft)} 
&  {\red xlabelPos=top} \\ 
   \RDD{ylabelPos}=axis  \RDI{ylabelPos}{pst-plot} 
   &  {\red ylabelPos}=right 
   &  {\red ylabelPos}=left {\blue (\dft)}  \\
\hline 
\end{tabular} 



\bigskip
\psset{unit=0.7cm}
\begin{tabular}{|c|c|c|}
\hline \begin{pspicture}(-1,-1)(4,4)  
 \psaxes[labelsep=.5cm,tickcolor=red]{->}(3,3)
\end{pspicture}  
&  \begin{pspicture}(-1,-1)(4,4)  
 \psaxes[xlabelsep=-.5cm, ylabelsep=.5cm,,tickcolor=red]{->}(3,3)
\end{pspicture} 
&  \begin{pspicture}(-1,-1)(4,4) 
 \psaxes[ylabelsep=-.5cm, xlabelsep=.5cm,tickcolor=red]{->}(3,3)
\end{pspicture} \\ 
\hline  
\RDD{labelsep}= .5cm  \RDI{labelsep}{pst-plot} 
&  
\RDD{xlabelsep}= -.5cm \RDI{xlabelsep}{pst-plot} 
&  
{\red  xlabelsep}= .5cm 
\\ 

& 
\RDD{ylabelsep}= .5cm  \RDI{ylabelsep}{pst-plot} 
&  
{\red ylabelsep}=-.5cm 
\\ \hline 
\multicolumn{3}{|c|}{ \blue \dft{} :  labelsep = 5pt, xlabelsep = 5pt, ylabelsep =5pt } \\ \hline 
\end{tabular} 
\bigskip

\psset{unit=0.7cm}
\begin{tabular}{|c|c|c|}
\hline \begin{pspicture}(-1,-1)(4,4)  
 \psaxes[xlabelOffset=0.5,tickcolor=red]{->}(3,3)
\end{pspicture}  
&  \begin{pspicture}(-1,-1)(4,4)  
 \psaxes[ylabelOffset=0.5,tickcolor=red]{->}(3,3)
\end{pspicture} 
&  \begin{pspicture}(-1,-1)(4,4) 
 \psaxes[ xlabelOffset=-0.5,tickcolor=red]{->}(3,3)
\end{pspicture} \\ \hline  
\RDD{xlabelOffset}=0.5  \RDI{xlabelOffset}{pst-plot} 
&  
\RDD{ylabelOffset}=0.5  \RDI{ylabelOffset}{pst-plot} 
&  
\RDD{xlabelOffset}= -0.5  \RDI{xlabelOffset}{pst-plot}  
\\ \hline 
\multicolumn{3}{|c|}{ \blue \dft{} :  xlabelOffset =0 , xlabelOffset = 0 } \\ \hline 
\end{tabular} 


% % % % % % % % % % % % A voir
%%--------------------------------------------------
\SbSbSSCT{Taille des étiquettes }{Size of labels}

\psset{unit=0.5cm}

\begin{tabular}{|c|c|c|}
\hline \begin{pspicture}(-1,-1)(4,4)  
 \psaxes[labelFontSize=\scriptstyle]{->}(3,3)
\end{pspicture}  
& % \psset{mathLabel=false}
\begin{pspicture}(-1,-1)(4,4)  
 \psaxes[xlabelFontSize=\footnotesize]{->}(3,3)
\end{pspicture} 
& 
 %\psset{mathLabel=true}
\begin{pspicture}(-1,-1)(4,4) 
 \psaxes[ylabelFontSize=\small]{->}(3,3)
\end{pspicture} \\ 
\hline 
\RDD{labelFontSize}=\BS{scriptstyle}  \RDI{labelFontSize}{pst-plot} &

 \RDD{xlabelFontSize}=\BS{}footnotesize \RDI{xlabelFontSize}{pst-plot}  
& 

\RDD{ylabelFontSize}=\BS{}small
\RDI{ylabelFontSize}{pst-plot} 
\\ \hline 
\end{tabular} 
%
%
%%\psset{mathLabel=false}
%%--------------------------------------------------
\SbSbSSCT{\'Etiquette avec extension}{Labels with extra}

\psset{yunit=0.5cm,xunit=1.2cm}

\begin{tabular}{|c|c|}
\hline 
\begin{pspicture}(-1.5,-1.5)(3.5,3.5)  
 \psaxes[xlabelFactor=$$\cdot 10^3 $$,ylabelFactor= V]{->}(3,3)
\end{pspicture}  
& 
\begin{pspicture}(-1.5,-1.5)(3.5,3.5)  
 \psaxes[ylabelFactor=$$\cdot 10^6$$,xlabelFactor=s]{->}(3,3)
\end{pspicture} \\  \hline 
\RDD{xlabelFactor}=\BS{cdot} $10^3$  \RDI{xlabelFactor}{pst-plot} 	& {\red  xlabelFactor}= s \\ 
{\red  ylabelFactor}= V 				&  \RDD{ylabelFactor}=$\cdot 10^6$ \RDI{ylabelFactor}{pst-plot} 
\\  \hline 

\end{tabular} 



%--------------------------------------------------
\SbSbSSCT{Les décimales dans les étiquettes}{Decimals in labels}


\psset{yunit=0.5cm,xunit=0.65cm}

\begin{tabular}{|c|c|c|}
\hline \begin{pspicture}(-1.5,-1)(4,4)  
 \psaxes[comma=true,xyDecimals=1]{->}(3,3)
\end{pspicture}  
& 
\begin{pspicture}(-1.5,-1)(4,4)  
 \psaxes[comma=false,xDecimals=1]{->}(3,3)
\end{pspicture} 
& 
\begin{pspicture}(-1.5,-1)(4,4) 
 \psaxes[decimalSeparator=h,yDecimals=2]{->}(3,3)
\end{pspicture} \\ 
\hline 
\RDD{comma}=true \RDI{comma}{pst-plot} 
& {\red comma}= false {\blue (\dft )}
& \RDD{decimalSeparator}=h   \RDI{decimalSeparator}{pst-plot} \\
\RDD{xyDecimals}=1 \RDI{xyDecimals}{pst-plot}
&  \RDD{xDecimals}=1  \RDI{xDecimals}{pst-plot}
& \RDD{yDecimals}=2  \RDI{yDecimals}{pst-plot}\\ 

\hline 
\end{tabular} 

%------------------------------------------------------------------------------
\SbSbSSCT{Liste comme étiquettes de graduations}{List of labels as  axis} 


\begin{center}
\begin{tabular}{|c|} \hline
 \begin{psgraph}[labels=none](5,4){5cm}{3cm} 
\TFRGB{\psaxes[xLabels={,un,deux,trois,quatre},xLabelsRot=45,yLabels={,petit,moyen,grand},yLabelsRot=30](5,4)}{\psaxes[xLabels={,one,two,three,four},xLabelsRot=45,yLabels={,small,big,huge},yLabelsRot=30](5,4)}
\end{psgraph} \\ 
\\ \hline
% \BS{}begin\AC{psgraph}[labels=none](5,4)\AC{5cm}\AC{3cm} \\
\TFRGB{\BS{}psaxes[\RDD{xLabels}=\AC{,un,deux,trois,quatre},\RDD{xLabelsRot}=45,}{\BS{}psaxes[\RDD{xLabels}=\AC{,one,two,three,four},\RDD{xLabelsRot}=45,} 
 \RDI{xLabels}{pst-plot} \RDI{xLabelsRot}{pst-plot}  \\
\TFRGB{\RDD{yLabels}=\AC{,petit,moyen,grand},\RDD{yLabelsRot}=30](5,4)}{\RDD{yLabels}=\AC{,small,big,huge},\RDD{yLabelsRot}=30](5,4)}  \RDI{yLabelsRot}{pst-plot} \\ \hline
\end{tabular}
\end{center}



\newpage
%-----------------------------------------------------------------------
\SbSSCT{Légende}{Legend}
\SbSbSSCT{Position de la légende}{Legend position}

%\TFRGB{
%Syntaxe :  
%\BS{pslegend}[position] (décalage en x,décalage `en y ) \AC{Légende }}
%{Syntax :  
%\BS{pslegend}[position] ( x offset, y offset ) \AC{légend }} 

\begin{tabular}{|c|c|}
\hline 
\pslegend[lb]{left bottom}
  \begin{psgraph}[axesstyle=frame,](0,0)(6,6){6cm}{3cm}
 \end{psgraph} 
& 
\pslegend[lt]{left top}
\begin{psgraph}[axesstyle=frame,](0,0)(6,6){6cm}{3cm}
\end{psgraph} 
\\  \hline 
\BSS{pslegend}[\RDD{lb}]\AC{left bottom} \BSI{pslegend}{pst-plot} \RDI{lb}{pst-plot}& \BSS{pslegend}[\RDD{ lt}]\AC{left top} \RDI{lt}{pst-plot}
\\  \hline  
\pslegend[rb]{right bottom}
  \begin{psgraph}[axesstyle=frame,](0,0)(6,6){6cm}{3cm}
 \end{psgraph} 
&  
\pslegend[rt]{right top}
  \begin{psgraph}[axesstyle=frame,](0,0)(6,6){6cm}{3cm}
 \end{psgraph}
\\ \hline 
\BSS{pslegend}[\RDD{rb}]\AC{right bottom}  \RDI{rb}{pst-plot} & 
\BSS{pslegend}[\RDD{rt}]\AC{right top} {\blue (\dft)}  \RDI{rt}{pst-plot}
\\  \hline 
\end{tabular} 
\bigskip

\begin{tabular}{|c|c|}
\hline 
\pslegend[lb](20,10){left bottom}
  \begin{psgraph}[axesstyle=frame](0,0)(6,6){6cm}{3cm}
 \end{psgraph} 
 &
 \pslegend[lb](10,20){left bottom}
   \begin{psgraph}[axesstyle=frame](0,0)(6,6){6cm}{3cm}
  \end{psgraph} 
  \\ \hline  
 \BS{pslegend}[lb]{\red (20,10)}\AC{left bottom} &  \BS{pslegend}[lb]{\red (10,20)}\AC{left bottom}
 \\ \hline  
\end{tabular} 
%------------------------------------------------------------------------------

\SbSbSSCT{Aspect de la légende}{Legend style}



\newpsstyle{legendstyle}{fillstyle=solid,fillcolor=cyan,shadow=true}

\pslegend[lt]{\red\rule[1ex]{2em}{1pt} & courbe 1\\
\blue\rule[1ex]{2em}{1pt} &  courbe 2\\
 \green\rule[1ex]{2em}{1pt} & courbe 3}
 
 \begin{center}
 \begin{tabular}{|c|} \hline
  \begin{psgraph}[axesstyle=frame,](0,0)(12,6){0.6\linewidth}{3cm}
 \end{psgraph} \\  \hline
\BSS{newpsstyle}\AC{\RDD{legendstyle}} \AC{fillstyle=solid,fillcolor=cyan,shadow=true} \\
 \BS{pslegend}[lt]\AC{\BS{red} \BS{rule}[1ex]\AC{2em} \AC{1pt} \& courbe 1 \BS{}\BS{}  \\
\BS{blue}  \BS{rule}[1ex]\AC{2em}\AC{1pt} \&  courbe 2 \BS{}\BS{}\\ 
\BS{green} \BS{rule}[1ex]\AC{2em}\AC{1pt} \& courbe 3
}  \RDI{legendstyle}{pst-plot} \\  \hline 
 \end{tabular}
 \end{center}
 
 
 %-----------------------------------------------------------
\SbSSCT{Points particuliers sur les axes}{Point on axes}
 
 \TFRGB{syntaxe}{syntax} :\\
\BSS{psxTick} [Options]\AC{rotation}(x position)\AC{label} 
\BSI{psxTick}{pst-plot}\\
\BSS{psyTick} [Options]\AC{rotation}(y position)\AC{label} 
\BSI{psyTick}{pst-plot}\\
\smallskip


\begin{center}
\begin{tabular}{|l|} \hline
\psset{llx=-1.5cm,lly=-.5cm,urx=.5cm  ,ury=.5cm  ,dotscale=2}
 \begin{psgraph}[axesstyle=frame,](0,0)(6,6){0.8\linewidth}{3cm}
 \psxTick[linecolor=red,labelsep=-20pt ]{45}(3.25){\red X=3.25}
\psyTick[linecolor=magenta](1.5){\magenta Y=1.5}
\end{psgraph} \\ \hline
\textbf{\BS{}psxTick}[linecolor=red,labelsep=-20pt ]\AC{45}(3.25)\AC{\BS{}red X=3.25}\\
\textbf{\BS{}psyTick}[linecolor=magenta](1.5)\AC{\BS{}magenta Y=1.5} \\
\hline
\end{tabular}
\end{center}

%---------------------------------------------------------

\SbSSCT{Portion de courbe}{Portion of curve }



\psset{llx=-1cm,lly=-.5cm,urx=.5cm  ,ury=.5cm  ,dotscale=2}
\begin{center}
\begin{tabular}{|l|} \hline
\readdata{\dat}{mesdata.dat}
 \begin{psgraph}[axesstyle=frame,xticksize=0 4cm,yticksize=0 12cm,subticks=0,Dx=100,Dy=.01](0,0)(750,.12){12cm}{4cm} 
\listplot[linecolor=black,linewidth=1pt]{\dat}
\listplot[xStart=200,xEnd=300,linecolor=blue,linewidth=5pt]{\dat}
\listplot[yStart=0,yEnd=.05,linecolor=red,linewidth=5pt]{\dat}
\listplot[nStart=200,nEnd=300,linecolor=magenta,linewidth=5pt]{\dat}
\end{psgraph}\\ \hline
\BS{}listplot[{\red \RDD{xStart}=200,\RDD{xEnd}=300},linecolor=blue,linewidth=5pt]\AC{\BS{}dat} 
\RDI{xStart}{pst-plot}   \RDI{xEnd}{pst-plot} \\
\BS{}listplot[{\red \RDD{yStart}=0,\RDD{yEnd}=.05},linecolor=red,linewidth=5pt]\AC{\BS{}dat} 
\RDI{yStart}{pst-plot}  \RDI{yEnd}{pst-plot} \\
\BS{}listplot[{\red \RDD{nStar}t=200,\RDD{nEnd}=300},linecolor=magenta,linewidth=5pt]\AC{\BS{}dat}
 \RDI{nStar}{pst-plot}  \RDI{nEnd}{pst-plot} \\
\hline

\end{tabular}
\end{center}

%--------------------------------------------------
\SbSSCT{Option yMaxValue et yMinValue}{yMaxValue and yMinValue}


\begin{center}
\begin{tabular}{|l|} \hline

 \begin{psgraph}[axesstyle=frame,dx=100](0,0)(0,-1.2)(13,1.2){12cm}{3cm} 
 \psset{algebraic=true}
 \psplot[plotpoints=2000,linecolor=blue,linestyle=dotted]{0}{12.56}{sin(x)}
\psplot[yMaxValue=.7,yMinValue=-.7,plotpoints=2000,linecolor=red,linewidth=5pt]{0}{12.56}{sin(x)}
\end{psgraph} \\ \hline
\BS{}psplot[{\red \RDD{yMaxValue}=.7,\RDD{yMinValue}=-.7},plotpoints=2000,linecolor=red,linewidth=5pt]\AC{0}\AC{12.56}\AC{sin(x)}
 \RDI{yMaxValue}{pst-plot}  \RDI{yMinValue}{pst-plot}\\
\hline
\multicolumn{1}{|c|}{ \blue \dft{}  yMaxValue=  1.e30 \hspace{1cm} yMinValue =  -1.e30 }
\\ \hline

\end{tabular}
\end{center}

%--------------------------------------------------
\SbSSCT{\'Echelle trigonométrique}{Axis with trigonmetrical units}


\psset{yunit=0.45cm,xunit=0.45cm}

\begin{tabular}{|c|c|c|}
\hline \begin{pspicture}(-1.5,-1.5)(7,4)  
 \psaxes[trigLabels=true,trigLabelBase=3]{->}(\psPiTwo,3)
\end{pspicture}  
& 
\begin{pspicture}(-1.5,-1.5)(7,4)   
 \psaxes[xtrigLabels=true,dx=\pstRadUnit]{->}(\psPiTwo,3)
\end{pspicture} 
& 
\begin{pspicture}(-1.5,-1.5)(7,4)  
 \psaxes[ytrigLabels=true,xunit=\pstRadUnit]{->}(\psPi,3)
\end{pspicture} \\ 
\hline 
\RDD{trigLabels}=true  \RDI{trigLabels}{pst-plot} & \RDD{xtrigLabels}=true  \RDI{xtrigLabels}{pst-plot} & \RDD{ytrigLabels}=true  \RDI{ytrigLabels}{pst-plot} \\
\RDD{trigLabelBase}=3  \RDI{trigLabelBase}{pst-plot} 
& dx=\BSS{pstRadUnit} \BSI{pstRadUnit}{pst-plot} 
&xunit=\BS{}pstRadUnit \\ \hline 
\multicolumn{3}{|c|}{ \blue \dft : {\red trigLabelBase} = 0 , {\red trigLabels} = false , {\red xtrigLabels} = false , {\red ytrigLabels} = false}
\\ \hline
\end{tabular} 

\bigskip
\TFRGB{Constantes prédéfinies}{Trigonometrical constants} 


\begin{tabular}{|c|c|c|}
\hline nom					& valeur 		& math \\ 
\hline \BS{}psPiFour 		& 12.566371 	& $4 \pi$ \\ 
\hline \BS{}psPiTwo 		& 6.283185 		& $2 \pi$  \\ 
\hline \BS{}psPi			& 3.14159265		&  $ \pi$ \\ 
\hline \BS{}psPiH 			& 1.570796327	& $\pi/2$  \\ 
\hline \BS{}pstRadUnit 		& 1.047198cm 	& $\pi/3$  \\ 
\hline \BS{}pstRadUnitInv 	& 0.95493cm 	& $3/\pi$  \\ 
\hline 
\end{tabular} 
 

%==================================================================
\SbSSCT{\'Echelle logarithmique}{Logarithmic axis}
\psset{yunit=.7cm,xunit=.7cm}


\begin{tabular}{|c|c|c|} \hline
\pspicture(-1,-1)(3.5,3.5)
 \psaxes[subticks=5,xylogBase=10,logLines=all](3,3)
\endpspicture\hspace{1cm}
&
\pspicture(-1,-1)(3.5,3.5)
 \psaxes[subticks=10,axesstyle=frame,ylogBase=2,xlogBase=10,tickstyle=inner](3,3)
 \endpspicture
 & 
 \pspicture(-1,-1)(3.5,3.5)
 \psaxes[ylogBase={},logLines=y,xticksize=max,xsubticksize=1,tickcolor=red,subtickcolor=blue,tickwidth=1pt,subticks=5,xsubticks =10](3,3)
\endpspicture\\  \hline
\RDD{xylogBase}=10   \RDI{xylogBase}{pst-plot} 	& \RDD{xlogBase}=10    \RDI{xlogBase}{pst-plot} 	& \RDD{ylogBase}=\AC{} \RDI{ylogBase}{pst-plot} \\
\RDD{logLines}=all   \RDI{logLines}{pst-plot} 	& 
{\red ylogBase}=2 	& 
{\red logLines}=all\\
subticks=5 		& subticks=10,		& subticks=5\\
 				&  tickstyle=inner 	& xsubticks =10\\ \hline 
\end{tabular}

%------------------------------
\SbSSCT{Coordonnées de l'environnement psgraph}{Coordinates of psgraph }

\begin{tabular}{|c|}
\hline  
\psset{llx=-5mm,lly=-.5cm,urx=.5cm  ,ury=.5cm  ,dotscale=2}
\begin{psgraph}[axesstyle=none,xticksize=2,yticksize=6](0,0)(6,2){6cm}{2cm}
\psdot[linecolor=red](\psgraphLLx,\psgraphLLy)
\psdot[linecolor=blue](\psgraphLLx,\psgraphURy)
\psdot[linecolor=cyan](\psgraphURx,\psgraphLLy)
 \psdot[linecolor=green](\psgraphURx,\psgraphURy)
\end{psgraph} 
\\  \hline 
\BS{psdot}[linecolor=red](\BS{\BDD{psgraphLLx}},\BS{\BDD{psgraphLLy}}) \BDI{psgraphLLx}{pst-plot} \BDI{psgraphLLy}{pst-plot} \\
\BS{psdot}[linecolor=blue]({\blue \BS{psgraphLLx}},\BS{\BDD{psgraphURy}})  \BDI{psgraphURy}{pst-plot} \\
\BS{psdot}[linecolor=cyan](\BDD{psgraphURx},{\blue \BS{psgraphLLy}}) \BDI{psgraphURx}{pst-plot}  \\
\BS{psdot}[linecolor=green]({\blue \BS{psgraphURx},\BS{psgraphURy}}) \\
\hline 
\end{tabular} 

%\subsection{paramètres d'un graphe en barres }
\SbSSCT{paramètres d'un graphe en barres }{Parameters of a bar graph }

\label{bar}
\begin{tabular}{|c|} 	\hline  
\BS{listplot}[plotstyle=bar,\RDD{barwidth}=1]\AC{\BS{dat}} 
\RDI{barwidth}{pst-plot} \\ \hline  
\readdata{\dat}{table.dat}
\begin{pspicture}(-1.5,-1.5)(8.5,11)
\psaxes(0,0) (8,10) 
\listplot[plotstyle=bar,barwidth=1]{\dat} 
\end{pspicture}
\\ \hline 
\dft : barwidth = 0.25cm
\\ \hline 
\end{tabular} 


\begin{tabular}{|c|} 	\hline  
\BS{listplot}[plotstyle=bar,\RDD{interrupt}=\AC{7,1,5}] \RDI{interrupt}{pst-plot}\\
\AC{ 0 5 1 17 2 15 3 20 4 1 5 3 6 22 7 1 8 18 } 
\\ \hline 
\psset{yunit=.5cm}
\begin{pspicture}(-1.5,-1.5)(10,19)
\psaxes(0,0) (10,18) 
\listplot[plotstyle=bar,barwidth=0.3cm,interrupt={7,1,5}]{
	0 5 1 17 2 15 3 20 4 1 5 3 6 22 7 1 8 18 }
\end{pspicture}
\\ \hline 
\end{tabular}




\newpage

%\section[Créer un graphe d'après un fichier de données]{Créer un graphe d'après un fichier de données  \cite{pst-user}  \cite{pst-plot}}
\SSCT{Créer un graphe d'après un fichier de données}{Data graph}


\psset{llx=-1.5cm,lly=-.5cm,urx=.5cm,ury=0.5cm,fillstyle=none}

\subsection{Macro  fileplot , psfileplot \cite{pst-user}  \cite{pst-plot}}

\TFRGB{Syntaxe}{syntax} : \BSS{fileplot} [Options] \AC{\TFRGB{fichier}{file}} ou  \BSS{psfileplot} [Options] \AC{\TFRGB{fichier}{file}} \BSI{psfileplot}{pst-plot} 

\bigskip


\begin{tabular}{|l|}
\hline  
\psset{xunit=.015cm,yunit=30cm}

\begin{pspicture}(-1cm,-1cm)(748,.13)
\psaxes[Dx=100,Dy=.02]{->}(748,.13)
\fileplot[linecolor=red,linewidth=2pt]{mesdata.dat}
\end{pspicture}
\\ \hline 

\BSS{fileplot}[linecolor=red,linewidth=2pt]\AC{mesdata.dat} \BSI{fileplot}{pst-plot} 

 \\ \hline 
option plotstyle : \TFRGB{seulement}{only} \hspace{1cm}  \og line \fg \hspace{1cm} \og polygon  \fg  \hspace{1cm}\og dots \fg \\
\TFRGB{Séparateurs de données}{Separators of data} : \hspace{1cm}  \og \AC{}  \fg  \hspace{1cm}  \og ()  \fg  \hspace{1cm} \og ,  \fg  \hspace{1cm} \og \TFRGB{espace}{space}  \fg 
 \\ \hline 
\end{tabular} 

%------------------

\subsection{Macro  dataplot , psdataplot}

\TFRGB{Syntaxe}{syntax} :  \BSS{dataplot} [Options] \AC{\BS{}macro} ou \BSS{psdataplot} [Options] \AC{\BS{}macro}

\BSI{dataplot}{pst-plot} \BSI{psdataplot}{pst-plot} 
 
\TFRGB{Elle doit être précédé de}{It must be preceded by }  : \BSS{readdata}\AC{\BS{}macro}\AC{nomfichier}
\BSI{readdata}{pst-plot}
\bigskip
 



\begin{center}
\begin{tabular}{|l|} \hline

 \begin{psgraph}[axesstyle=frame,xticksize=0 4cm,yticksize=0 9cm,subticks=0,Dx=100,dy=.01,Dy=.2](0,0)(750,.12){9cm}{4cm} 
\readdata{\dat}{mesdata.dat}
\dataplot[linecolor=red,linewidth=2pt]{\dat}
\end{psgraph}\\ \hline
\textbf{\BS{}readdata}\AC{{\red \BS{}dat}}\AC{mesdata.dat} \\
\textbf{\BS{dataplot}}[linecolor=red,linewidth=2pt]\AC{{\red \BS{}dat}} 

\\\hline
\end{tabular}
\end{center}


\subsection{Macro  savedata}


\TFRGB{Syntaxe}{syntax} : \BSS{savedata}\AC{\BS{}macro}[données en XY]
\BSI{savedata}{pst-plot}

\bigskip

\textbf{\BS{}savedata}\AC{\BS{}mydata}[\AC{x0, y0}, \AC{x1., y1}, .... \AC{xn., yn}]


\newpage
\subsection{Macro listplot , pslistplot}


\TFRGB{Syntaxe}{syntax} : \BSS{listplot}\AC{data}   \BSS{pslistplot}\AC{data}
\BSI{listplot}{pst-plot} \BSI{pslistplot}{pst-plot}

\bigskip


\begin{center}
\begin{tabular}{|l|} \hline
\psset{llx=-1cm,lly=-1cm,urx=.5cm,ury=0.5cm}
 \begin{psgraph}[axesstyle=frame,xticksize=0 4cm,yticksize=0 9cm,subticks=0,Dx=100,Dy=.02](0,0)(750,.12){9cm}{4cm} 
\readdata{\dat}{mesdata.dat}
\listplot[plotstyle=curve,showpoints=true,dotstyle=triangle]{\dat}
\end{psgraph} \\ \hline
\textbf{\BS{}listplot}[plotstyle=curve,showpoints=true,dotstyle=triangle]\AC{\BS{}dat} \\ \hline
\TFRGB{liste des coordonnées séparées que par des espaces blancs}{separator of data  :  space only} !
 \\ \hline
\end{tabular}
\end{center}


\bigskip
%=====================================================
\SbSSCT{\'Echelle des données}{Scale factor}

\BSS{pstScalePoints}(\TFRGB{facteur échelle X,facteur échelle Y}{X scale factor , Y scale factor})\AC{\TFRGB{code calcul postscript sur X}{PostScript code applied to the x values}}\AC{\TFRGB{code calcul postscript sur Y}{PostScript code applied to the Y values} }
\BSI{pstScalePoints}{pst-plot} 
\bigskip


\begin{center}
\begin{tabular}{|l|} \hline
\readdata{\dat}{mesdata.dat}
\psset{llx=-.5cm,lly=-.5cm,urx=.5cm,ury=0.5cm}
 \begin{psgraph}[axesstyle=frame,xticksize=0 4cm,yticksize=0 9cm,subticks=0,Dx=100,dy=1,Dy=2](0,0)(750,12){9cm}{4cm} 
\pstScalePoints(1,100){}{}
\listplot[linecolor=red,linewidth=2pt]{\dat}
\end{psgraph} \\\hline
\BSS{pstScalePoints}({\red 1,100})\AC{}\AC{} \\ \hline
 \multicolumn{1}{|c|}{\cyan \TFRGB{ne fonctionne qu'avec \BS{listplot} et \BS{pslistplot}{Only work with \BS{listplot} and \BS{pslistplot}! }}} \\  \hline  
\end{tabular}
\end{center}
 
%----------------------------------------------------------------------------------------
\SbSSCT{Options de lecture de fichier}{Options  reading the file of data}



\begin{tabular}{|c|c|} \hline
\psset{llx=-1cm}
 \begin{psgraph}[axesstyle=frame,xticksize=0 3cm,yticksize=0 5cm,subticks=0,Dx=100,Dy=.02](0,0)(750,.12){5cm}{3cm} 
 \readdata[ignoreLines=100]{\dat}{mesdata.dat}
\listplot[plotstyle=line,linecolor=blue,linewidth=2pt ]{\dat}
\end{psgraph}
&
 \readdata{\dat}{mesdata.dat}
 \psset{llx=-1cm}
 \begin{psgraph}[axesstyle=frame,xticksize=0 3cm,yticksize=0 5cm,subticks=0,Dx=100,Dy=.02](0,0)(750,.12){5cm}{3cm}
\listplot[plotstyle=curve,linestyle=dotted,linewidth=1pt ]{\dat}
\listplot[plotstyle=dots,linecolor=red,nStep=100,dotscale=2]{\dat}
\end{psgraph} \\  \hline
\RDD{ignoreLines}=100  \RDI{ignoreLines}{pst-plot} & \RDD{nStep}=100  \RDI{nStep}{pst-plot} \\ \hline
 

\end{tabular}


\bigskip
\begin{tabular}{|c|c|} \hline
\psset{llx=-1cm}
 \begin{psgraph}[axesstyle=frame,xticksize=0 3cm,yticksize=0 5cm,subticks=0,Dx=100,Dy=.02](0,0)(750,.12){5cm}{3cm} 
\listplot[plotstyle=curve,linestyle=dotted,linewidth=1pt ]{\dat}
\listplot[plotstyle=dots,linecolor=red,xStep=100,dotscale=2]{\dat}
\end{psgraph}
&
\psset{llx=-1cm}
 \begin{psgraph}[axesstyle=frame,xticksize=0 3cm,yticksize=0 5cm,subticks=0,Dx=100,Dy=.02](0,0)(350,.12){5cm}{3cm}
  
\listplot[plotstyle=curve,linestyle=dotted,linewidth=1pt ,xEnd=300]{\dat}
\listplot[plotstyle=dots,linecolor=red,yStep=0.02,dotscale=2,xEnd=300]{\dat}
\end{psgraph} \\  \hline
\RDD{xStep}=100  \RDI{xStep}{pst-plot}& \RDD{yStep}=0.02,xEnd=300  \RDI{yStep}{pst-plot} \\ \hline
\end{tabular}

%----------------------------------------------------------------------------------

\SbSSCT{Table de données multiples}{Multiple data table}


\TFRGB{Soit une table de données est organisée ainsi}{The data table has 4 columns of data } :

\smallskip

\begin{center}
\begin{tabular}{|c|c|c|c|}
\hline  A & B & C & B \\ 
\hline  &  &  &  \\ 
\hline  &  &  &  \\ 
\hline 
\end{tabular} 
\end{center}
\bigskip

\begin{tabular}{|l @{:} l |}
\hline 
\multicolumn{2}{|c|}{ \BS{listplot}[\RDD{plotNoMax}=3,\RDD{plotNoX=2},\RDD{plotNo}=2]\AC{\BS{data}}  
\RDI{plotNoMax}{pst-plot}  \RDI{plotNoX}{pst-plot} 
 \RDI{plotNo}{pst-plot} 
} 

\\ \hline 
plotNoX=2 & \TFRGB{ la colonne B correspond à X}{ X values on  column B} \\
plotNoMax=3 & \TFRGB{ soit 2 colonnes y + 1 colonne x}{ 1 columm with x values + 2 columms with y values} \\ 
plotNo=2  & \TFRGB{ la colonne C correspond à Y}{ Y values on  column C} \\  \hline 
\end{tabular} 

\newpage
\subsection{Macro sur Excel}
\TFRGB{Voici un programme en Visual Basic permettant de créer son fichier de données d'après une feuille Excel}{Here is a Visual Basic program to create a data file from an Excel spreadsheet}


\bigskip

\psframebox{\parbox[c]{\linewidth}{

Sub mesdata() \\
deb = 8      \hspace{2cm}        ' \TFRGB{première ligne de données}{first line of data} \\
fin = 382     \hspace{2cm}         ' \TFRGB{dernière ligne de données}{last line of data} \\
colX = 5       \hspace{2cm}        ' \TFRGB{colonne des valeurs de X}{column of the values  X} \\
colY = 6        \hspace{2cm}       ' \TFRGB{colonne des valeurs de Y}{column of the values  Y} \\
nom = "mesdata.dat"  \hspace{1cm}  ' \TFRGB{nom du fichier}{name of the file} \\

Dim valX, valY As Double \\

'\TFRGB{pour effacer le fichier}{to erase the file} \\
Open nom For Output Access Write As \#1 \\
Close \#1 \\

'\TFRGB{création du fichier}{creation of the file} \\
For i = deb To fin \\
Open nom For Append As \#1 \\
valX = Cells(i, colX) \\
valY = Cells(i, colY) \\

Write \#1, valX \\
Write \#1, valY \\
Close \#1 \\
Next \\

End Sub}} \\

\TFRGB{A copier dans un module Excel et modifier les paramètres}{copy  this code in a module Excel and  modify the parameters}  deb, fin , colX, colY et nom

 

%%\essais{pstgraph2ess.tex}

\newpage

%\section[Créer un graphe d'après une équation]{Créer un graphe d'après une équation  \cite{pst-user} \cite{pst-plot}}
\SSCT{Créer un graphe d'après une équation}{Equation graph}
 


\subsection{Macro psplot }

%\TFRGB{syntaxe}{syntax} :  \BSS{psplot} [Options] \AC{x min}\AC{x max}\AC{fonction}
% 
%
%
% \bigskip


\psset{yunit=1cm}
\begin{center}
\begin{tabular}{|l|} \hline

 \multicolumn{1}{|c|}{
 \begin{psgraph*}[,xticksize= -1.5 1.5 ,yticksize=12cm , subticks=0, dx=90,Dx=90 , dy=.5,Dy=.5](0,0)(0,-1.5)(720,1.5){12cm}{5cm } 
 \psplot[plotpoints=200,linecolor=red]{0}{720}{x sin}
 \end{psgraph*} 
%\pspicture(0,-1.5)(8,1.50)
%\psset{xunit=0.01cm}
%\psplot[plotpoints=200,linecolor=red]{0}{720}{x sin}
%\psline{->}(800,0)
%\endpspicture 
}\\ \hline
% \BS{}psset\AC{xunit=0.01cm}\\
\textbf{\BS{psplot}}[\rnode{A}{\RDD{plotpoints}=200},linecolor=red]\rnode{B}{\AC{0}\AC{720}}\rnode{C}{\AC{x sin}}  \RDI{plotpoints}{pst-plot}  \\ \hline
\\
\rnode{AA}{\TFRGB{ nombre de points utilisés}{number of point used}}  \hspace{.5cm}
\rnode{BB}{\TFRGB{unité de x en degré}{x unit : degree}}  \hspace{0.5cm}
\rnode{CC}{\TFRGB{fonction  en langage PostScript  \footnotemark[1]}{Function in postscript code }} 
 \\ \hline
{\blue \dft{} :  plotpoints = 50 }
\\ \hline 


\end{tabular}
\end{center}
\footnotetext[1]{\TFRGB{formule de calcul en langage PostScript (voir \pageref{postcript})}{formula in the PostScript language} }

 \ncline[linecolor=blue]{A}{AA}  \ncline[linecolor=blue]{B}{BB}  \ncline[linecolor=blue]{C}{CC}


 %-----------------------------------------------------
 \subsection{Macro  parametricplot}
 
%\TFRGB{syntaxe}{syntax} :  \BSS{parametricplot} [Options]\AC{t min}\AC{t max}\AC{x(t) y(t) }
%\BSI{parametricplot}{pst-plot} 
%  \bigskip

  
%\BS{}psset\AC{xunit=2cm}\\


%
%\begin{psgraph*}[,xticksize= -2 2 ,yticksize=-2 2, subticks=0, dx=.5,Dx=.5, dy=.5,Dy=.5](0,0)(-2,-2)(2,2){5cm}{5cm } 
%\parametricplot[linewidth=1.2pt,plotstyle=ccurve,linecolor=red]{0}{360}{t sin t 2 mul sin}
%\end{psgraph*}
 

\begin{center}
\begin{tabular}{|c|} \hline
\begin{psgraph*}[,xticksize= -1.5 1.5 ,yticksize=-1.5 1.5 , subticks=0, dx=.5,Dx=.5, dy=.5,Dy=.5](0,0)(-1.5,-1.5)(1.5,1.5){5cm}{5cm } 
\parametricplot[linewidth=1.2pt,plotstyle=ccurve,linecolor=red]{0}{360}{t sin t 2 mul sin}
\end{psgraph*}
%
%\begin{pspicture}(-5cm,-1.5)(5cm,1.5)
%
%\psset{xunit=2cm}
%\parametricplot[linewidth=1.2pt,plotstyle=ccurve,linecolor=red]{0}{360}{t sin t 2 mul sin}
%\psline{<->}(0,-1.2)(0,1.2)
%\psline{<->}(-1.2,0)(1.2,0)
%\end{pspicture}
\\ \hline
\textbf{\BS{}parametricplot}[linewidth=1.2pt,plotstyle=ccurve,linecolor=red]\\
\rnode{A}{\AC{0}\AC{360}}\Rnode*[fillcolor=yellow,fillstyle=solid]{B}{\AC{t sin t 2 mul sin} }\hspace{1cm} (functions : $sin(t)$ et  $sin(2t)$)\\
%\BS{psline}\AC{<->}(0,-1.2)(0,1.2)\\
%\BS{psline}\AC{<->}(-1.2,0)(1.2,0)\\ 
\hline
\\ \\
 \rnode{AA}{\TFRGB{L'unité  de t est le degré }{ t unit in degree}} \hspace{1cm} 
\rnode{BB}{\TFRGB{Les deux fonctions doivent être écrites en langage PostScript }{the 2 functions in postcript code } !}
\\ \hline
\end{tabular}
\end{center}
 \ncline[linecolor=blue]{A}{AA}  \ncline[linecolor=blue,armA=.5cm,angleA=90]{B}{BB}  
 
\newpage

%=====================================================
\SbSSCT{Graphe polaire}{Polar graph}

%\TFRGB{syntaxe}{syntax} :  
%\BS{}psplot [\RDD{polarplot}=true]\AC{angle début}\AC{angle fin}[commande postscript]\AC{équation} \RDI{polarplot}{pst-plot} 

%Exemple :



%\AC{}psplot[plotstyle=curve,polarplot=true,linecolor=red,algebraic=true]\AC{0}\AC{\BS{}psPiTwo}\AC{6*sin(2*x)}

%\begin{psgraph*}[axesstyle=none,xticksize= -5 5 ,yticksize=-5 5, subticks=0](0,0)(-5,-5)(5,5){5cm}{5cm } 
%\psplot[plotstyle=curve,polarplot=true,linecolor=red]{0}{360}{ x 2 mul sin 6 mul }
%\end{psgraph*}

\begin{center}
\begin{tabular}{|c|} \hline
\begin{psgraph*}[xticksize= -5 5 ,yticksize=-5 5, subticks=0](0,0)(-5,-5)(5,5){5cm}{5cm } 
\psplot[plotstyle=curve,polarplot=true,linecolor=red]{0}{360}{ x 2 mul sin 6 mul }
\end{psgraph*}
\\ \hline
\BS{}psplot[plotstyle=curve,{\red polarplot}=true,linecolor=red]\\
\AC{0}\AC{360} \AC{ x 2 mul sin 6 mul }  \\  \hline ($6*sin(2*x)$)\\ \hline
\end{tabular}
\end{center}

%\psplot[plotstyle=curve,polarplot=true,linecolor=red,algebraic=true]{0}{\psPiTwo}{6*sin(2*x)}


\subsection{Modules infix-RPN et pst-infixplot  \cite{pst-infix}}
% \begin{itemize}
% \item  Ces Modules permettent de s'affranchir des équations en langage PostScript
% \item Le module  infix-RPN utilise \BS{RPN} 
% \item Le module  pst-infixplot donne les commandes \BS{}psPlot et \BS{}parametricPlot 
%\end{itemize}
%
%
%\bigskip
%Exemples :
\label{infix-RPN}
\label{pst-infixplot}

%% \multicolumn{1}{|c|}{
% \begin{psgraph*}[,xticksize= -1.5 1.5 ,yticksize=12cm , subticks=0, dx=90,Dx=90 , dy=.5,Dy=.5](0,0)(0,-1.5)(720,1.5){12cm}{5cm } 
%\infixtoRPN{sin(x)} 
%\psplot[plotpoints=200,linecolor=red]{0}{720}{\RPN}
% \end{psgraph*} 
%%\pspicture(0,-1.5)(8,1.50)
%%\psset{xunit=0.01cm}
%%\psplot[plotpoints=200,linecolor=red]{0}{720}{x sin}
%%\psline{->}(800,0)
%%\endpspicture 
%%}


\begin{center}
\begin{tabular}{|l|} \hline
 \begin{psgraph*}[,xticksize= -1.5 1.5 ,yticksize=12cm , subticks=0, dx=90,Dx=90 , dy=.5,Dy=.5](0,0)(0,-1.5)(720,1.5){10cm}{3cm } 
\infixtoRPN{sin(x)} 
\psplot[plotpoints=200,linecolor=red]{0}{720}{\RPN}
 \end{psgraph*} 
%\pspicture(0,-1.5)(9cm,1.50)
%\psset{xunit=0.01cm}
%\infixtoRPN{sin(x)} 
%\psplot[plotpoints=200,linecolor=red]{0}{720}{\RPN}
%\psline{->}(800,0)
%\endpspicture 
\\ \hline
 \BSS{infixtoRPN}\AC{sin(x)} \BSI{infixtoRPN}{pst-plot}  \\
\BS{}psplot[plotpoints=200]\AC{0}\AC{720}\AC{{\red \BS{}RPN}}
\\ \hline
\end{tabular} 
\end{center}



\smallskip
\begin{center}
\begin{tabular}{|l|} \hline
 \begin{psgraph*}[,xticksize= -1.5 1.5 ,yticksize=12cm , subticks=0, dx=90,Dx=90 , dy=.5,Dy=.5](0,0)(0,-1.5)(720,1.5){10cm}{3cm } 
\psPlot[linecolor=red]{0}{720}{sin(x)}
 \end{psgraph*} 
 
%\pspicture(0,-1.1)(9cm,1.1)
%\psset{xunit=0.01cm}
%\psset{plotpoints=200}
%\psPlot[linecolor=red]{0}{720}{sin(x)}
%\psline{->}(8cm,0)
%\endpspicture 
\\ \hline
%\BS{}psset\AC{plotpoints=200} \\
 \BSS{psPlot}\AC{0}\AC{720}\AC{sin(x)}
\\ \hline
\end{tabular} 
\end{center}




\smallskip
\begin{center}
\begin{tabular}{|c|} \hline
\begin{psgraph*}[,xticksize= -1.5 1.5 ,yticksize=-1.5 1.5 , subticks=0, dx=.5,Dx=.5, dy=.5,Dy=.5](0,0)(-1.5,-1.5)(1.5,1.5){5cm}{5cm } 
\parametricPlot[linecolor=red,plotpoints=200]{0}{360}{sin(t)}{sin(2*t)}
\end{psgraph*}
%\psset{plotpoints=200}
%\pspicture(-2,-1.5)(2,1.5)
%\psset{xunit=1.7cm}
%\parametricPlot[linecolor=red,plotpoints=200]{0}{360}{sin(t)}{sin(2*t)}
%\psline{<->}(0,-1.2)(0,1.2)
%\psline{<->}(-1.2,0)(1.2,0)
%\endpspicture
\\ \hline
 \BSS{parametricPlot}[linecolor=red,plotpoints=200]\AC{0}\AC{360}\AC{sin(t)}\AC{sin(2*t)} \\ \hline
\end{tabular} 
\end{center}

%\newpage
\subsection{Option algebraic}
%\begin{itemize}
%\item  Cette option permet de s'affranchir des équations en langage PostScript
%\item L'unité de x ou de t est le\emph{ radian }
%\end{itemize}
%
%
%
%\bigskip
%%
%\BS{}psplot[algebraic,plotpoints=200]\AC{0}\AC{12.56}\AC{ sin(x)}
%\smallskip

\begin{center}
\begin{tabular}{|c|} \hline
 \begin{psgraph*}[,xticksize= -1.5 1.5 ,yticksize=13	 , subticks=0, dx=1,Dx=1, dy=.5,Dy=.5](0,0)(0,-1.5)(13,1.5){10cm}{3cm } 
\psplot[algebraic,plotpoints=200,linecolor=red]{0}{12.56}{ sin(x)}
 \end{psgraph*} 
%\pspicture(0,-1.1)(8,1.1)
%\psset{xunit=.5cm}
%\psplot[algebraic,plotpoints=200,linecolor=red]{0}{12.56}{ sin(x)}
%\psline{->}(15,0)
%\endpspicture
\\ \hline
\BS{}psplot[\RDD{algebraic},plotpoints=200]\AC{0}\AC{12.56}\AC{{\red sin(x)}} \RDI{algebraic}{pst-plot} \\ \hline
\TFRGB{L'unité de x  est le}{x unit in }\emph{ radian }
 \\ \hline
\end{tabular} 
\end{center}
\bigskip
%

\smallskip


\begin{center}
\begin{tabular}{|c|} \hline
\begin{psgraph*}[,xticksize= -1.5 1.5 ,yticksize=-1.5 1.5 , subticks=0, dx=.5,Dx=.5, dy=.5,Dy=.5](0,0)(-1.5,-1.5)(1.5,1.5){5cm}{5cm } 
\parametricplot[algebraic,plotpoints=200,linecolor=red]{0}{6.28}{sin(t)|sin(2*t)}
\end{psgraph*}
%\pspicture(-2,-1.5)(2,1.5)
%\psset{xunit=1.7cm}
%\parametricplot[algebraic,plotpoints=200,linecolor=red]{0}{6.28}{sin(t)|sin(2*t)}
%\psline{<->}(0,-1.2)(0,1.2)
%\psline{<->}(-1.2,0)(1.2,0)
%\endpspicture 
\\ \hline
\BSS{parametricplot}[{\red algebraic},plotpoints=200]\AC{0}\AC{6.28}\AC{{\red sin(t)|sin(2*t)}}
\\ \hline
\TFRGB{L'unité de  t est le}{t unit in }\emph{ radian }
\\ \hline
\end{tabular} 
\end{center}

\subsection{Options VarStep et VarStepEpsilon }
\label{var}

\begin{tabular}{|c|} \hline 
 \begin{psgraph*}[,xticksize= -1.5 1.5 ,yticksize=13	 , subticks=0, dx=1,Dx=1, dy=.5,Dy=.5](0,0)(0,-1.5)(13,1.5){10cm}{3cm } 
\psplot[algebraic,VarStep=true,linecolor=red,showpoints=true,VarStepEpsilon=1]{0}{12.56}{ sin(x)}
 \end{psgraph*} 
 \\ \hline   
\BS{psplot}[algebraic,\RDD{VarStep}=true,showpoints=true,\RDD{VarStepEpsilon}=1]{0}\AC{12.56}\AC{ sin(x)}  \RDI{VarStep}{pstricks-add}   \RDI{VarStepEpsilon}{pstricks-add} 
\\ \hline  
 \begin{psgraph*}[,xticksize= -1.5 1.5 ,yticksize=13	 , subticks=0, dx=1,Dx=1, dy=.5,Dy=.5](0,0)(0,-1.5)(13,1.5){10cm}{3cm } 
\psplot[algebraic,VarStep=true,linecolor=red,showpoints=true,VarStepEpsilon=.1]{0}{12.56}{ sin(x)}
 \end{psgraph*} 
 \\ \hline  
\BS{psplot}[algebraic,\RDD{VarStep}=true,showpoints=true,\RDD{VarStepEpsilon}=0.1]{0}\AC{12.56}\AC{ sin(x)}
\\ \hline  
 \begin{psgraph*}[,xticksize= -1.5 1.5 ,yticksize=13	 , subticks=0, dx=1,Dx=1, dy=.5,Dy=.5](0,0)(0,-1.5)(13,1.5){10cm}{3cm } 
\psplot[algebraic,VarStep=true,linecolor=red,showpoints=true,VarStepEpsilon=0.01]{0}{12.56}{ sin(x)}
 \end{psgraph*} 
 \\ \hline  
\BS{psplot}[algebraic,\RDD{VarStep}=true,showpoints=true,\RDD{VarStepEpsilon}=.001]{0}\AC{12.56}\AC{ sin(x)}
\\ \hline  

\end{tabular} 






\newpage

% \section[Des outils pour les graphes]{Des outils pour les graphes \cite{pstricks-add} }
\SSCT{Des outils pour les graphes}{Tools for graph} 
%\SbSSCTTC{Coordonnées d'un point}{Coordonnées d'un point \cite{pst-plot}}{Coordinates of a point}{Coordinates of a point \cite{pst-plot}}
\SbSSCT{Coordonnées d'un point}{Coordinates of a point}
 

\begin{tabular}{|c|} \hline
 \begin{psgraph}[axesstyle=frame,xticksize=-1.5 1.5 , yticksize=0 13,subticks=0,Dx=1 , dy=.5,Dy=.5] (0,0)(0,-1.5)(13,1.5){12cm}{4cm} 
\psplot[algebraic,plotpoints=200]{0}{12.56}{ sin(x)}
\psCoordinates[linecolor=red,linestyle=dashed,dotstyle=square,dotscale=2](*4 { sin(x)})
\end{psgraph}
%\pspicture(0,-1.1)(8,1.1)
%\psset{xunit=.5cm}
%\psplot[algebraic,plotpoints=200]{0}{12.56}{ sin(x)}
%\psline{->}(15,0)
%\psline{->}(0,-1.1)(0,1.1)
%\psCoordinates[linecolor=red,linestyle=dashed,dotstyle=square,dotscale=2](*4 { sin(x)})
%\endpspicture
\\ \hline
\BSS{psCoordinates}[linecolor=red,linestyle=dashed,dotstyle=square,dotscale=2](*4 \AC{sin(x)})  \BSI{psCoordinates}{pst-plot} \\ \hline
\end{tabular} 

%---------------------------------------------
\SbSSCTTC{Tangente}{Tangente \cite{pstricks-add}}{Tangent}{Tangent \cite{pstricks-add}}
\SbSSCT{Tangente}{Tangent}

\SbSbSSCT{Tangente à une courbe d'après un fichier de données}{Tangent to a data file curve }
\BSS{psTangentLine}[Options] (x1,y1)(x2,y2)(x3,y3)\AC{x}\AC{dx}
\psset{llx=-.7cm,lly=-.5cm,urx=.5cm,ury=0.5cm,fillstyle=none}
\BSI{psTangentLine}{pst-plot} 
\begin{center}
\begin{tabular}{|c|} \hline
 \begin{psgraph}[axesstyle=frame,xticksize=0 4cm,yticksize=0 12cm,subticks=0,Dx=100,dy=.01,Dy=.2](0,0)(750,.12){12cm}{4cm} 
\fileplot[linecolor=blue,linewidth=1pt]{mesdata.dat}
\psTangentLine[linecolor=red,arrows=<->,arrowscale=2] (198,0.0824)(200,0.0811)(202,0.07962){200}{30}
\psTangentLine[linecolor=magenta,arrows=->,arrowscale=2](118,0.0465)(120,0.0445)(122,0.0428){120}{30}
\end{psgraph} 
\\\hline
\BSS{psTangentLine}[linecolor=magenta,{\red arrows=->}](118,0.0465)(120,0.0445)(122,0.0428)\AC{120}\AC{30}  \BSI{psTangentLine}{pstricks-add} \\
\BS{}psTangentLine[linecolor=red,{\red arrows=<->}] (198,0.0824)(200,0.0811)(202,0.07962)\AC{200}\AC{30}

 \\\hline
\end{tabular}
\end{center}


\newpage
\SbSbSSCTTC{Tangente à une fonction}{Tangente à une fonction \cite{pstricks-add}}{Tangent to a function curve}{Tangent to a function curve \cite{pstricks-add}}
\psset{xunit=1cm,yunit=.8cm}



\TFRGB{syntaxe}{syntax} : \BSS{psplotTangent} * [Options] \AC{x}\AC{dx}\AC{function}
\BSI{psplotTangent}{pst-plot} 


\begin{center}
\begin{tabular}{|c|} \hline
\TFRGB{Commande sans astérisque}{Command without asterisk}  \\ \hline
 \begin{psgraph*}[,xticksize= -1.5 1.5 ,yticksize=13	 , subticks=0, dx=1,Dx=1, dy=.5,Dy=.5](0,0)(0,-1.5)(13,1.5){10cm}{3cm } 
\psplot[algebraic,plotpoints=200,linecolor=blue]{0}{12.56}{ sin(x)}
\psplotTangent[linecolor=red,arrows=<->,arrowscale=2,algebraic=true,linewidth=2pt]{\psPiH}{2}{sin(x)}
\psplotTangent[linecolor=magenta,arrows=<-,arrowscale=2,algebraic=true,linewidth=2pt]{\psPi}{2}{sin(x)}
\psplotTangent[linecolor=green,arrows=->,arrowscale=2,algebraic=true,linewidth=2pt]{\psPiTwo}{3}{sin(x)}
 \end{psgraph*} 
\\\hline
\BS{psplotTangent}[linecolor=red,arrows=<->]\AC{\BS{}psPiH}\AC{2}\AC{sin(x)}  \footnotemark[1]  \BSI{psplotTangent}{pstricks-add} \\
\BS{}psplotTangent[linecolor=magenta,arrows=<-]\AC{\BS{}psPi}\AC{2}\AC{sin(x)}\\
\BS{}psplotTangent[linecolor=green,arrows=->]\AC{\BS{}psPiTwo}\AC{3}\AC{sin(x)}
\\\hline
\TFRGB{Commande avec astérisque}{Command with asterisk}  \\ \hline
 \begin{psgraph*}[,xticksize= -1.5 1.5 ,yticksize=13	 , subticks=0, dx=1,Dx=1, dy=.5,Dy=.5](0,0)(0,-1.5)(13,1.5){10cm}{3cm } 
\psplot[algebraic,plotpoints=200,linecolor=blue]{0}{12.56}{ sin(x)}
\psplotTangent*[linecolor=red,arrows=<->,arrowscale=2,algebraic=true,linewidth=2pt]{\psPiH}{2}{sin(x)}
\psplotTangent*[linecolor=magenta,arrows=<-,arrowscale=2,algebraic=true,linewidth=2pt]{\psPi}{2}{sin(x)}
\psplotTangent*[linecolor=green,arrows=->,arrowscale=2,algebraic=true,linewidth=2pt]{\psPiTwo}{3}{sin(x)}
 \end{psgraph*} %\\
 \\\hline
 \BS{}psplotTangent*[linecolor=red,arrows=<->]\AC{\BS{}psPiH}\AC{2}\AC{sin(x)}\\
 \BS{}psplotTangent*[linecolor=magenta,arrows=<-]\AC{\BS{}psPi}\AC{2}\AC{sin(x)}\\
 \BS{}psplotTangent*[linecolor=green,arrows=->]\AC{\BS{}psPiTwo}\AC{3}\AC{sin(x)}
 \\\hline

\end{tabular}
\end{center}
%\footnotetext[1]{arrowscale=2,algebraic=true,linewidth=2pt} 

%=====================================================================================

\SbSbSSCTTC{Tangente à une courbe polaire}{Tangente à une courbe polaire \cite{pstricks-add}}{Tangent to a polar curve}{Tangent to a polar curve \cite{pstricks-add}}
\psset{unit=0.4cm}



\begin{center}
\begin{tabular}{|c|c|} \hline
\TFRGB{Commande sans astérisque}{Command without asterisk}   & \TFRGB{Commande avec astérisque}{Command with asterisk} 
 \\\hline
 \begin{psgraph*}[,xticksize= -6 6,yticksize=-6 6 , subticks=0, dx=1,Dx=1, dy=1,Dy=1 ](0,0)(-6,-6)(6,6){4cm}{4cm }
 \psplot[plotstyle=curve,polarplot=true,linecolor=blue,algebraic=true]{0}{\psPiTwo}{6*sin(2*x)}
 \psplotTangent[polarplot,linecolor=red,arrows=->,arrowscale=2,algebraic=true,linewidth=2pt]{2}{3}{6*sin(2*x)}
 \end{psgraph*}
 &
 \begin{psgraph*}[,xticksize= -6 6,yticksize=-6 6 , subticks=0, dx=1,Dx=1, dy=1,Dy=1 ](0,0)(-6,-6)(6,6){4cm}{4cm }
 \psplot[plotstyle=curve,polarplot=true,linecolor=blue,algebraic=true]{0}{\psPiTwo}{6*sin(2*x)}
 \psplotTangent*[polarplot,linecolor=red,arrows=->,arrowscale=2,algebraic=true,linewidth=2pt]{2}{3}{6*sin(2*x)}
 \end{psgraph*}
 \\\hline
\multicolumn{2}{|c|}{\BS{}psplotTangent[polarplot,linecolor=red,arrows=->]\AC{2}\AC{3}\AC{6*sin(2*x)} \footnotemark[1]}
 \\\hline
\end{tabular}
\end{center}

\footnotetext[1]{arrowscale=2,algebraic=true,linewidth=2pt}
%===========================================================

\SbSbSSCTTC{Normale à une courbe}{Normale à une courbe \cite{pstricks-add}}{Normal of the tangent line}{Normal of the tangent line \cite{pstricks-add}}

\begin{center}
\begin{tabular}{|c|c|} \hline
\begin{psgraph*}[,xticksize=  0 4,yticksize=0 4 , subticks=0, dx=1,Dx=1, dy=1,Dy=1 ](0,0)(0,0)(4,4){4cm}{4cm }
\pscurve[showpoints=true](1,1)(2,3)(3,2)
\psTangentLine[linecolor=blue,arrows=<->,arrowscale=2,algebraic=true](1,1)(2,3)(3,2){2}{1}
\psTangentLine[linecolor=red,arrows=->,arrowscale=2,Tnormal](1,1)(2,3)(3,2){2}{1}
\end{psgraph*}
&
 \begin{psgraph*}[,xticksize= -2 1.5 ,yticksize=13	 , subticks=0, dx=1,Dx=1, dy=.5,Dy=.5](0,0)(0,-2)(13,1.5){7cm}{4cm } 
\psplot[algebraic,plotpoints=200,linecolor=blue]{0}{12.56}{ sin(x)}
\psplotTangent[linecolor=blue,arrows=<->,arrowscale=2,algebraic=true]{5}{3}{sin(x)}
\psplotTangent[linecolor=red,algebraic=true,Tnormal,arrows=->,arrowscale=2]{5}{2}{sin(x)}
 \end{psgraph*} 
 \\ \hline
\BS{psTangentLine}[\RDD{Tnormal}](1,1)(2,3)(3,2)\AC{2}\AC{1}  \RDI{Tnormal}{pstricks-add} 
&
\BS{}psplotTangent[{\red Tnormal}]{5}\AC{2}\AC{sin(x)}
 \\\hline
\end{tabular}
\end{center}



%================================================================
\SbSbSSCTTC{Dérivée}{Dérivée \cite{pstricks-add}}{Derivatives of a function}{Derivatives of a function \cite{pstricks-add}}
\psset{unit=1.5cm}


 
\begin{center}
\begin{tabular}{|c|} \hline
 \begin{psgraph*}[,xticksize= -1.5 1.5,yticksize=0 13 , subticks=0, dx=1,Dx=1, dy=.5,Dy=.5 ](0,0)(0,-1.5)(13,1.5){12cm}{3cm }
 \psplot[algebraic,plotpoints=200,linecolor=blue]{0}{12.56}{sin(.75*x)}
  \psplot[algebraic,plotpoints=200,linecolor=red,linewidth=2pt]{0}{12.56}{Derive(1,sin(.75*x))}
   \psplot[algebraic,plotpoints=200,linecolor=green,linewidth=2pt]{0}{12.56}{Derive(2,sin(.75*x))}
 \end{psgraph*}
%\pspicture(0,-1.1)(8,1.1)
%\psset{xunit=.75cm}
%\psplot[algebraic,plotpoints=200,linecolor=blue]{0}{12.56}{sin(.75*x)}
% \psplot[algebraic,plotpoints=200,linecolor=red,linewidth=2pt]{0}{12.56}{Derive(1,sin(.75*x))}
%  \psplot[algebraic,plotpoints=200,linecolor=green,linewidth=2pt]{0}{12.56}{Derive(2,sin(.75*x))}
%\psline{->}(15,0)
%\endpspicture

 \\\hline
%\BS{}psplot[algebraic,plotpoints=200,linecolor=blue]\AC{0}\AC{12.56}\AC{sin(.75*x)} \\
\BS{}psplot[algebraic,plotpoints=200,linecolor=red]\AC{0}\AC{12.56}\AC{{\red \RDD{Derive}(\textbf{1},sin(.75*x))}}  \RDI{Derive}{pstricks-add}  \\
\BS{}psplot[algebraic,plotpoints=200,linecolor=green]\AC{0}\AC{12.56}\AC{{\red Derive(\textbf{2},sin(.75*x))}}
 \\\hline
\end{tabular}
\end{center}
%====================================================
\newpage
\SbSbSSCTTC{Intégrale de Riemann}{Intégrale de Riemann \cite{pstricks-add}}{Riemann integral}{Riemann integral \cite{pstricks-add}}
 

\begin{center}
\begin{tabular}{|c|c|} \hline
\psset{xunit=.5cm}
\pspicture(0,-1.1)(13.5,1.1)
\psStep[algebraic,linecolor=magenta,StepType=upper,fillcolor=yellow,fillstyle=solid](0,12.56){24}{ sin(x)}
\psplot[algebraic,plotpoints=200,linecolor=blue]{0}{12.56}{ sin(x)}
\psline{->}(13,0)
\endpspicture
&
\psset{xunit=.5cm}
\pspicture(0,-1.1)(13.5,1.1)
\psStep[algebraic,linecolor=magenta,StepType=u,fillcolor=yellow,fillstyle=solid](0,12.56){24}{ sin(x)}
\psplot[algebraic,plotpoints=200,linecolor=blue]{0}{12.56}{ sin(x)}
\psline{->}(13,0)
\endpspicture
 \\\hline
 \BSS{psStep}[\RDD{StepType}=upper](0,12.56)\AC{24}\AC{sin(x)} \BSI{psStep}{pst-plot}
  \BSI{psStep}{pstricks-add}   \RDI{StepType}{pstricks-add} 
 &
 \BSS{psStep}[\RDD{StepType}=u](0,12.56)\AC{24}\AC{sin(x)}
 \\\hline

\psset{xunit=.5cm}
\pspicture(0,-1.1)(13.5,1.1)
\psStep[algebraic,linecolor=magenta,StepType=lower,fillcolor=yellow,fillstyle=solid](0,12.56){24}{ sin(x)}
\psplot[algebraic,plotpoints=200,linecolor=blue]{0}{12.56}{ sin(x)}
\psline{->}(15,0)
\endpspicture
&
\psset{xunit=.5cm}
\pspicture(0,-1.1)(13.5,1.1)
\psStep[algebraic,linecolor=magenta,StepType=l,fillcolor=yellow,fillstyle=solid](0,12.56){24}{ sin(x)}
\psplot[algebraic,plotpoints=200,linecolor=blue]{0}{12.56}{ sin(x)}
\psline{->}(13,0)
\endpspicture
 \\\hline
 \BS{psStep}[\RDD{StepType}=lower](0,12.56)\AC{24}\AC{sin(x)}&
 \BS{}psStep[\RDD{StepType}=l](0,12.56)\AC{24}\AC{sin(x)}
 \\\hline
\psset{xunit=.5cm}
\pspicture(0,-1.1)(13.5,1.1)
\psStep[algebraic,linecolor=magenta,StepType=Riemann,fillcolor=yellow,fillstyle=solid](0,12.56){24}{ sin(x)}
\psplot[algebraic,plotpoints=200,linecolor=blue]{0}{12.56}{ sin(x)}
\psline{->}(15,0)
\endpspicture
&
\psset{xunit=.5cm}
\pspicture(0,-1.1)(13.5,1.1)
\psStep[algebraic,linecolor=magenta,StepType=R,fillcolor=yellow,fillstyle=solid](0,12.56){24}{ sin(x)}
\psplot[algebraic,plotpoints=200,linecolor=blue]{0}{12.56}{ sin(x)}
\psline{->}(13,0)
\endpspicture
 \\\hline
 \BS{}psStep[\RDD{StepType}=Riemann](0,12.56)\AC{24}\AC{sin(x)}&
 \BS{}psStep[\RDD{StepType}=R](0,12.56)\AC{24}\AC{sin(x)}
 \\\hline

\psset{xunit=.5cm}
\pspicture(0,-1.1)(13.5,1.1)
%\psaxes[Dx=100,Dy=.02]{->}(13,1.5)


\psStep[algebraic,linecolor=magenta,StepType=infimum,fillcolor=yellow,fillstyle=solid](0,12.56){24}{ sin(x)}
\psplot[algebraic,plotpoints=200,linecolor=blue]{0}{12.56}{ sin(x)}
\psline{->}(15,0)
\endpspicture
&
\psset{xunit=.5cm}
\pspicture(0,-1.1)(13.5,1.1)
\psStep[algebraic,linecolor=magenta,StepType=i,fillcolor=yellow,fillstyle=solid](0,12.56){24}{ sin(x)}
\psplot[algebraic,plotpoints=200,linecolor=blue]{0}{12.56}{ sin(x)}
\psline{->}(13,0)
\endpspicture
 \\\hline
 \BS{}psStep[\RDD{StepType}=infimum](0,12.56)\AC{24}\AC{sin(x)}&
 \BS{}psStep[\RDD{StepType}=i](0,12.56)\AC{24}\AC{sin(x)}
 \\\hline
 
\psset{xunit=.5cm}
\pspicture(0,-1.1)(13.5,1.1)
%\psaxes[Dx=100,Dy=.02]{->})(13,1.5)

\psStep[algebraic,linecolor=magenta,StepType=supremum,fillcolor=yellow,fillstyle=solid](0,12.56){24}{ sin(x)}
\psplot[algebraic,plotpoints=200,linecolor=blue]{0}{12.56}{ sin(x)}
\psline{->}(15,0)
\endpspicture
&
\psset{xunit=.5cm}
\pspicture(0,-1.1)(13.5,1.1)
\psStep[algebraic,linecolor=magenta,StepType=s,fillcolor=yellow,fillstyle=solid](0,12.56){24}{ sin(x)}
\psplot[algebraic,plotpoints=200,linecolor=blue]{0}{12.56}{ sin(x)}
\psline{->}(13,0)
\endpspicture
 \\\hline
 \BS{}psStep[\RDD{StepType}=supremum](0,12.56)\AC{24}\AC{sin(x)}&
 \BS{}psStep[\RDD{StepType}=s](0,12.56)\AC{24}\AC{sin(x)}
 \\\hline

\end{tabular}
\end{center}

\SbSbSSCTTC{Méthode de Newton}{Méthode de Newton \cite{pst-plot}}{Newton method}{Newton method \cite{pst-plot}}

\TFRGB{syntaxe}{syntax} : \BS{psNewton} [Options] \AC{$x0$} \AC{f(x)} \AC{\TFRGB{nombre d'itération}{number of iteration}}

\bigskip
\psset{yunit=.5cm}
\begin{tabular}{|c|c|}
\hline 
\multicolumn{2}{|c|}{\BS{}psplot[algebraic,linestyle=dotted]\AC{0}\AC{12.56}\AC{0.5*x $\hat{}$  2-2} }\\ 
\multicolumn{2}{|c|}{\BSS{psNewton}[linecolor=red]\AC{4}\AC{0.5*x $\hat{}$ 2-2}\AC{20}  \BSI{psNewton}{pst-plot} }\\ \hline 	 
\begin{pspicture}*[algebraic,shift=*](-1,-2)(5,7)
%\psframe(-1,-2)(5,7)
\psaxes{->}(4.5,6)
\psplot[algebraic,plotpoints=200,linestyle=dotted]{0}{12.56}{0.5*x^2-2}
\psframe[linestyle=dashed,linecolor=green](1.9,-.2)(2.2,.5)
\pnode(2.2,0.2){A}
\psNewton[linecolor=red,linewidth=0.5pt]{4}{0.5*x^2-2}{20}
\end{pspicture}
&  
\psset{unit=15cm,yunit=6cm}
\begin{pspicture}*[algebraic,shift=*](1.9,-.25)(2.2,.55)
\psaxes[Dx=.1]{->}(4,4)
\psplot[algebraic,plotpoints=200,linestyle=dotted]{1.5}{2.5}{0.5*x^2-2}
\psframe[linestyle=dashed,linecolor=green](1.9,-.2)(2.2,.5)
\pnode(1.9,0.2){B}

\psNewton[linecolor=red,linewidth=0.5pt,arrowscale=3]{2.1}{0.5*x^2-2}{2}
\end{pspicture}
\\ 
\hline 
\end{tabular} 
%------------------------------
\ncline[linestyle=dashed,linecolor=green]{->}{A}{B}
\bigskip


\begin{tabular}{|c|c|}
\hline 
\multicolumn{2}{|c|}{\BSS{psNewton}[linecolor=red,\RDD{plotstyle=xvalues}]\AC{4}\AC{0.5*x $\hat{}$ 2-2}\AC{1}  \BSI{psNewton}{pst-plot} }\\ \hline 	 
\begin{pspicture}*[algebraic,shift=*](-1,-2)(5,7)
%\psframe(-1,-2)(5,7)
\psaxes{->}(4.5,6)
\psplot[algebraic,plotpoints=200,linestyle=dotted]{0}{12.56}{0.5*x^2-2}
\psframe[linestyle=dashed,linecolor=green](1.9,-.2)(2.2,.5)
\psNewton[linecolor=red,linewidth=0.5pt]{4}{0.5*x^2-2}{20}
\psNewton[linecolor=red,linewidth=0.5pt,plotstyle=xvalues]{4}{0.5*x^2-2}{1}
\end{pspicture}
&  
\psset{unit=15cm,yunit=6cm}
\begin{pspicture}*[algebraic,shift=*](1.9,-.25)(2.2,.55)
\psaxes[Dx=.1]{->}(4,4)
\psplot[algebraic,plotpoints=200,linestyle=dotted]{1.5}{2.5}{0.5*x^2-2}
\psframe[linestyle=dashed,linecolor=green](1.9,-.2)(2.2,.5)
\psNewton[linecolor=red,linewidth=0.5pt]{2.1}{0.5*x^2-2}{2}
\psNewton[linecolor=red,linewidth=0.5pt,plotstyle=xvalues]{2.1}{0.5*x^2-2}{1}
\end{pspicture}
\\ \hline 
\end{tabular} 
%---------------------------------------------------------

\bigskip
\begin{tabular}{|c|c|}
\hline 
\multicolumn{2}{|c|}{\BSS{psNewton}[linecolor=red,{\red showDerivation=false}]\AC{4}\AC{0.5*x $\hat{}$ 2-2}\AC{1} }\\ \hline  
\begin{pspicture}*[algebraic,shift=*](-.5,-3)(5,7)
\psaxes{->}(10,10)
\psplot[algebraic,plotpoints=200,linestyle=dotted]{0}{12.56}{0.5*x^2-2}
\psNewton[linecolor=red,linewidth=0.5pt,showDerivation=false]{4}{0.5*x^2-2}{20}
\end{pspicture}
&  
\begin{pspicture}*[algebraic,shift=*](-.5,-3)(5,7)
\psaxes{->}(10,10)
\psplot[algebraic,plotpoints=200,linestyle=dotted]{0}{12.56}{0.5*x^2-2}
\psNewton[linecolor=red,linewidth=0.5pt,showDerivation=true]{4}{0.5*x^2-2}{20}
\end{pspicture}
\\ 
\hline \RDD{showDerivation} {\red =false } &   {\red  showDerivation=true} (par défaut) \\ 
\hline 
\end{tabular} 

\newpage
\subsection[Macro psFixpoint]{Macro psFixpoint \cite{pst-plot}}
\TFRGB{syntaxe}{syntax} : \BS{psFixpoint} [Options] \AC{$x_0$}\AC{f(x)}\AC{\TFRGB{nombre d'itération}{number of iteration}}

\bigskip
\begin{tabular}{|c|}
\hline  
\psset{unit=.5cm}

\begin{pspicture}*[algebraic](-2.5,-2)(10,10.5)
%\pframe(-2.5,-2)(10,10.5)
\psaxes{->}(10,10)
\psplot[algebraic,plotpoints=200,linestyle=dotted,linewidth=2pt]{1}{10}{0.2*x^2-2}
\psline[linecolor=red,linestyle=dashed](10,10)
\psFixpoint[linecolor=red]{6}{0.2*x^2-2}{3}
\end{pspicture}
\psset{unit=1cm}
\\  \hline 
\BS{}psplot[algebraic,linestyle=dotted]\AC{1}\AC{10}\AC{0.5*x $\hat{}$ 2-2} \\
\BS{}psline[linecolor=red,linestyle=dashed](10,10) \\
\BSS{psFixpoint}[linecolor=red]\AC{6}\AC{0.5*x $\hat{}$ 2-2}\AC{3} \BSI{psFixpoint}{pst-plot}  \\ \hline 
\end{tabular} 

\psset{unit=1cm,yunit=1cm}


\newpage
\subsection[Macro psVectorfield]{Macro psVectorfield \cite{pst-plot}}

\label{vec}
\begin{tabular}{|c|} \hline  
Solutions de $ \dfrac{dy}{dx}=x+y+1 $
\\ \hline  
\BSS{psVectorfield}[algebraic](-2,-2)(2,2)\AC{ x+y+1} \BSI{psVectorfield}{pst-plot} 
\\ \hline  
\psset{unit=1.5cm}
\begin{pspicture}(-2.5,-2.2)(2.5,2.2)
\psVectorfield[algebraic](-2,-2)(2,2){ x+y+1}
\end{pspicture}
\\ \hline 
\end{tabular} 




\bigskip

\begin{tabular}{|c|} \hline 
\BSS{psVectorfield}[algebraic,\RDD{Dx}=0.3,\RDD{Dy}=0.3](-2,-2)(2,2)\AC{ x+y+1}  \RDI{Dx}{pst-plot}   \RDI{Dy}{pst-plot} 
\\ \hline  
\psset{unit=1.5cm}
\begin{pspicture}(-2.2,-2.5)(2.2,2.5)
%\psaxes[ticksize=0 4pt,axesstyle=frame,tickstyle=inner,subticks=20,
%Ox=-1,Oy=-1](-1,-1)(1,1)
%\psset{arrows=->,algebraic}
\psVectorfield[algebraic,Dx=.3,Dy=.3](-2,-2)(2,2){ x+y+1}
\end{pspicture}
\\ \hline  
\dft : Dx= 0.1 , Dy= 0.1
\\ \hline 
\end{tabular}  

\newpage

%\section{Tracé de fonctions mathématiques \cite{pst-func} }
\SSCT{Tracé de fonctions mathématiques}{mathematical functions}
\label{func}

%\subsection{Courbe de Bezier}
\SbSSCT{Courbe de Bezier}{Bezier curve}

\begin{tabular}{|c|c|c|} \hline  
\begin{psgraph*}[axesstyle=none,xticksize= -2 2 ,yticksize=-2 2, subticks=0](0,0)(-2,-2)(2,2){2.5cm}{2.5cm } 
 \psBezier1[showpoints=true]{<->}(-2,-1)(0,2)
\end{psgraph*}
&  
\begin{psgraph*}[axesstyle=none,xticksize= -2 2 ,yticksize=-2 2, subticks=0](0,0)(-2,-2)(2,2){2.5cm}{2.5cm } 
 \psBezier2[showpoints=true]{<->}(-2,-1)(0,2)(2,1)
\end{psgraph*}
&  
\begin{psgraph*}[axesstyle=none,xticksize= -2 2 ,yticksize=-2 2, subticks=0](0,0)(-2,-2)(2,2){2.5cm}{2.5cm } 
 \psBezier3[showpoints=true]{<->}(-2,-1)(0,2)(2,1)(1,-1)
\end{psgraph*}
\\ \hline 
 \BSS{psBezier1} &   \BSS{psBezier2} & \BSS{psBezier3}  
\\ \hline 
\begin{psgraph*}[axesstyle=none,xticksize= -2 2 ,yticksize=-2 2, subticks=0](0,0)(-2,-2)(2,2){2.5cm}{2.5cm } 
 \psBezier4[showpoints=true]{<->}(-2,-1)(0,2)(2,1)(2,-1)(0,-2)
\end{psgraph*}
&  
\begin{psgraph*}[axesstyle=none,xticksize= -2 2 ,yticksize=-2 2, subticks=0](0,0)(-2,-2)(2,2){2.5cm}{2.5cm } 
 \psBezier5[showpoints=true]{<->}(-2,-1)(0,2)(2,1)(2,-1)(0,-2)(-1,-1)
\end{psgraph*}
&  
\begin{psgraph*}[axesstyle=none,xticksize= -2 2 ,yticksize=-2 2, subticks=0](0,0)(-2,-2)(2,2){2.5cm}{2.5cm } 
 \psBezier6[showpoints=true]{<->}(-2,-1)(0,2)(2,1)(2,-1)(0,-2)(-1,-1)(-2,1)
\end{psgraph*}
\\ \hline
 \BSS{psBezier4} &   \BSS{psBezier5} & \BSS{psBezier6} 
\\ \hline 
\begin{psgraph*}[axesstyle=none,xticksize= -2 2 ,yticksize=-2 2, subticks=0](0,0)(-2,-2)(2,2){2.5cm}{2.5cm } 
 \psBezier7[showpoints=true]{<->}(-2,-1)(0,2)(2,1)(2,-1)(0,-2)(-1,-1)(-2,1)(-1,2)
\end{psgraph*}
&  
\begin{psgraph*}[axesstyle=none,xticksize= -2 2 ,yticksize=-2 2, subticks=0](0,0)(-2,-2)(2,2){2.5cm}{2.5cm } 
 \psBezier8[showpoints=true]{<->}(-2,-1)(0,2)(2,1)(2,-1)(0,-2)(-1,-1)(-2,1)(-1,2)(1,1)
\end{psgraph*}
&  
\begin{psgraph*}[axesstyle=none,xticksize= -2 2 ,yticksize=-2 2, subticks=0](0,0)(-2,-2)(2,2){2.5cm}{2.5cm } 
 \psBezier9[showpoints=true]{<->}(-2,-1)(0,2)(2,1)(2,-1)(0,-2)(-1,-1)(-2,1)(-1,2)(1,1)(1,0)
\end{psgraph*}
\\ \hline
 \BSS{psBezier7} &   \BSS{psBezier8} & \BSS{psBezier9} 
 \\ \hline  
\end{tabular} 

\newpage
%============================
%\subsection{Polynôme de Chebyshev }
\SbSSCT{Polynôme de Chebyshev }{Chebyshev polynomial}

\subsubsection{Polynôme de première espèce}

\begin{tabular}{|c|c|c|} \hline
\multicolumn{3}{|c|}{ \BS{psplot}\AC{-1}\AC{1}\AC{1 x \BSS{ChebyshevT} }  \BSI{ChebyshevT}{pst-func} } 
\\ \hline   
\begin{psgraph*}[axesstyle=none,xticksize= -1.5 1.5 ,yticksize=-1.5 1.5 , subticks=0 ](0,0)(-1.5,-1.5)(1.5,1.5){3.5cm}{3.5cm}  
 \psplot{-1}{1}{1 x \ChebyshevT}
\end{psgraph*}
&  
\begin{psgraph*}[axesstyle=none,xticksize= -1.5 1.5 ,yticksize=-1.5 1.5 , subticks=0](0,0)(-1.5,-1.5)(1.5,1.5){3.5cm}{3.5cm} 
 \psplot{-1}{1}{3 x \ChebyshevT}
\end{psgraph*}
&  
\begin{psgraph*}[axesstyle=none,xticksize= -1.5 1.5 ,yticksize=-1.5 1.5 , subticks=0](0,0)(-1.5,-1.5)(1.5,1.5){3.5cm}{3.5cm}   
 \psplot{-1}{1}{6 x \ChebyshevT}
\end{psgraph*}
\\ \hline  
1 x \BSS{ChebyshevT} & 3 x \BSS{ChebyshevT} & 6 x \BSS{ChebyshevT} \\ 
\hline 
\end{tabular} 

\subsubsection{Polynôme de deuxième espèce }

\begin{tabular}{|c|c|c|} \hline
\multicolumn{3}{|c|}{ \BS{psplot}\AC{-1}\AC{1}\AC{1 x \BSS{ChebyshevU} }  \BSI{ChebyshevU}{pst-func} } 
\\ \hline   
\begin{psgraph*}[axesstyle=none,xticksize= -1.5 1.5 ,yticksize=-1.5 1.5 , subticks=0 ](0,0)(-1.5,-1.5)(1.5,1.5){3.5cm}{3.5cm} 
 \psplot{-1}{1}{1 x \ChebyshevU}
\end{psgraph*}
&  
\begin{psgraph*}[axesstyle=none,xticksize= -1.5 1.5 ,yticksize=-1.5 1.5 , subticks=0](0,0)(-1.5,-1.5)(1.5,1.5){3.5cm}{3.5cm} 
 \psplot{-1}{1}{3 x \ChebyshevU}
\end{psgraph*}
&  
\begin{psgraph*}[axesstyle=none,xticksize= -1.5 1.5 ,yticksize=-1.5 1.5 , subticks=0](0,0)(-1.5,-1.5)(1.5,1.5){3.5cm}{3.5cm}  
 \psplot{-1}{1}{6 x \ChebyshevU}
\end{psgraph*}
\\ \hline  
1 x \BSS{ChebyshevU} & 3 x \BSS{ChebyshevU} & 6 x \BSS{ChebyshevU} \\ 
\hline 
\end{tabular}

%ù+++++++++++++

\newpage

%\subsection{Fonction polynomiale}
\SbSSCT{Fonction polynomiale}{Function plynomial}

\begin{tabular}{|c|c|c|c|}	\hline 
\multicolumn{4}{|c|}{ \BSS{psPolynomial}[coeff= 1 ]\AC{-2}\AC{2} \BSI{psPolynomial}{pst-func} }
\\ \hline 
\begin{psgraph*}[axesstyle=none,xticksize= -2 2 ,yticksize=-2 2, subticks=0](0,0)(-2,-2)(2,2){2.5cm}{2.5cm } 
  \psPolynomial[coeff= 1 ]{-2}{2}
\end{psgraph*}
&
\begin{psgraph*}[axesstyle=none,xticksize= -2 2 ,yticksize=-2 2, subticks=0](0,0)(-2,-2)(2,2){2.5cm}{2.5cm } 
  \psPolynomial[coeff= 0 1 ]{-2}{4}
\end{psgraph*}
&
\begin{psgraph*}[axesstyle=none,xticksize= -2 2 ,yticksize=-2 2, subticks=0](0,0)(-2,-2)(2,2){2.5cm}{2.5cm } 
  \psPolynomial[coeff=0 0 1 ]{-2}{4}
\end{psgraph*}
&
\begin{psgraph*}[axesstyle=none,xticksize= -2 2 ,yticksize=-2 2, subticks=0](0,0)(-2,-2)(2,2){2.5cm}{2.5cm } 
  \psPolynomial[coeff=0 0 0 1 ]{-2}{4}
\end{psgraph*}
\\ 	\hline  
\RDD{coeff}= 1 & \RDD{coeff}=0  1 & \RDD{coeff}=0  1 & \RDD{coeff}=0 0 01
 \RDI{coeff}{pst-func} \\
$f(x)=1$ & $f(x)=x$ & $f(x)=x^2$ & $f(x)=x^3 $
\\ 	\hline 
\end{tabular}


 
 \bigskip

\begin{tabular}{|c|} \hline  
\begin{psgraph*}[axesstyle=none,xticksize= -2 2 ,yticksize=-2 4, subticks=0](0,0)(-2,-2)(4,2){4cm}{3cm } 
  \psPolynomial[coeff=0 0 0 1 ]{-2}{4}
  \psPolynomial[coeff=0 0 0 1 ,linecolor=red,xShift=2 ]{-2}{4}
\end{psgraph*}
\\ \hline  
\BS{psPolynomial}[coeff=0 0 0 1 ,linecolor=red,\RDD{xShift}=2 ]\AC{-2}\AC{4}
\\ \hline 
\end{tabular}  

 \bigskip


\begin{tabular}{|c|c|c|} \hline 
\multicolumn{3}{|c|}{ \BSS{psPolynomial}[coeff=0 0 0 0 0 1 ,linecolor=red,\RDD{Derivation}=1 ]\AC{-2}\AC{2}  \RDI{Derivation}{pst-func} } \\ \hline 
\begin{psgraph*}[axesstyle=none,xticksize= -2 2 ,yticksize=-2 2, subticks=0](0,0)(-2,-2)(2,2){3.5cm}{3.5cm} 
  \psPolynomial[coeff=0 0 0 0 0 1 ]{-2}{4}
  \psPolynomial[coeff=0 0 0 0 0 1 ,linecolor=red,Derivation=1 ]{-2}{4}
\end{psgraph*}
&
\begin{psgraph*}[axesstyle=none,xticksize= -2 2 ,yticksize=-2 2, subticks=0](0,0)(-2,-2)(2,2){3.5cm}{3.5cm} 
  \psPolynomial[coeff=0 0 0 0 0 1 ]{-2}{4}
  \psPolynomial[coeff=0 0 0 0 0 1 ,linecolor=red,Derivation= 2 ]{-2}{4}
\end{psgraph*}
&
\begin{psgraph*}[axesstyle=none,xticksize= -2 2 ,yticksize=-2 2, subticks=0](0,0)(-2,-2)(2,2){3.5cm}{3.5cm} 
  \psPolynomial[coeff=0 0 0 0 0 1 ]{-2}{4}
  \psPolynomial[coeff=0 0 0 0 0 1 ,linecolor=red,Derivation= 3 ]{-2}{4}
\end{psgraph*}
\\ \hline 
\RDD{Derivation}= 1 & \RDD{Derivation}= 2 & \RDD{Derivation}= 3
\\ \hline 
\end{tabular}  



\bigskip
\begin{tabular}{|c|} \hline 
\begin{psgraph*}[axesstyle=none,xticksize= -2 2 ,yticksize=-3 3, subticks=0](0,0)(-3,-2)(3,2){6cm}{4cm } 
 \psPolynomial[markZeros,dotscale=3,coeff=1 1 -1 -.5 0.15]{-3}{3}%
\end{psgraph*}
\\ \hline
 \BS{psPolynomial}[\RDD{markZeros},dotscale=3,coeff=1 1 -1 -.5 0.15]\AC{-3}\AC{3}%
\\ \hline 
\end{tabular}

%----------------------

\bigskip
\begin{tabular}{|c|} \hline 
\begin{psgraph*}[axesstyle=none,xticksize= -3 2 ,yticksize=-3 3, subticks=0](0,0)(-3,-3)(3,2){8cm}{5cm } 
 \psPolynomial[markZeros,dotscale=2,zeroLineTo=1,coeff=1 1 -1 -.5 0.15]{-3}{3}%
 \psPolynomial[linestyle=dotted,Derivation=1,coeff=1 1 -1 -.5 0.15]{-3}{3}% 
\end{psgraph*}
\\ \hline
 \BS{psPolynomial}[markZeros,\RDD{zeroLineTo}=1,coeff=1 1 -1 -.5 0.15]\AC{-3}\AC{3} \\
  \BS{psPolynomial}[linestyle=dotted,Derivation=1,coeff=1 1 -1 -.5 0.15]\AC{-3}\AC{3}
\\ \hline 
\end{tabular}

\bigskip
\begin{tabular}{|c|} \hline 
\begin{psgraph*}[axesstyle=none,xticksize= -1 2 ,yticksize=-2 2, subticks=0](0,0)(-2,-1)(2,2){8cm}{4cm } 
 \psPolynomial[coeff=1 1 -1 -.5 0.15]{-3}{3}%
 \psPolynomial[markZeros,linestyle=dotted,Derivation=1,zeroLineTo=0,zeroLineTo=0,zeroLineStyle=solid,zeroLineColor=red,zeroLineWidth=3pt,coeff=1 1 -1 -.5 0.15]{-3}{3}% 
\end{psgraph*}
\\ \hline
 \BS{psPolynomial}[coeff=1 1 -1 -.5 0.15]\AC{-3}\AC{3} \\
  \BS{psPolynomial}[markZeros,linestyle=dotted,Derivation=1,zeroLineTo=0,\\
 \RDD{zeroLineStyle}=solid,\RDD{zeroLineColor}=red,\RDD{zeroLineWidth}=3pt,\\
 coeff=1 1 -1 -.5 0.15]\AC{-3}\AC{3}
\\ \hline 
\end{tabular}


\bigskip
\begin{tabular}{|c|} \hline 
\begin{psgraph*}[axesstyle=none,xticksize= -1 2 ,yticksize=-3 3, subticks=0](0,0)(-2,-1)(1,2){8cm}{4cm } 
 \psPolynomial[coeff=1 1 -1 -.5 0.15]{-2}{1}%
 \psPolynomial[markZeros,linestyle=dotted,Derivation=2,zeroLineTo=0,zeroLineStyle=solid,zeroLineColor=red,zeroLineWidth=3pt,coeff=1 1 -1 -.5 0.15]{-2}{1}% 
\end{psgraph*}
\\ \hline
 \BS{psPolynomial}[coeff=1 1 -1 -.5 0.15]\AC{-3}\AC{3} \\
  \BS{psPolynomial}[markZeros,linestyle=dotted,Derivation=2,zeroLineTo=0,\\
\hspace{1cm} \RDD{zeroLineStyle}=solid,\RDD{zeroLineColor}=red,\RDD{zeroLineWidth}=3pt,\\
 coeff=1 1 -1 -.5 0.15]\AC{-3}\AC{3}
\\ \hline 
\end{tabular}

\newpage

% \subsection{Polynôme de Bernstein}
 \SbSSCT{Polynôme de Bernstein}{Bernstein polynomial}
 
 
\begin{tabular}{|c|c|c|} \hline  
\begin{psgraph*}[axesstyle=none,xticksize= -.5 1.5 ,yticksize=-.5 1.5,xticksize= -.5  1.5 , dx=.5,Dx=.5, dy=.5,Dy=.5 , subticks=0] (0,0)(-.5,-.5)(1.5,1.5){3.5cm}{3.5cm }  
\psBernstein(0,0)
\end{psgraph*}
&  
\begin{psgraph*}[axesstyle=none,xticksize= -.5 1.5 ,yticksize=-.5 1.5,xticksize= -.5  1.5 , dx=.5,Dx=.5, dy=.5,Dy=.5 , subticks=0] (0,0)(-.5,-.5)(1.5,1.5){3.5cm}{3.5cm }    
\psBernstein(0,1)
\end{psgraph*}
&  
\begin{psgraph*}[axesstyle=none,xticksize= -.5 1.5 ,yticksize=-.5 1.5,xticksize= -.5  1.5 , dx=.5,Dx=.5, dy=.5,Dy=.5 , subticks=0] (0,0)(-.5,-.5)(1.5,1.5){3.5cm}{3.5cm }   
\psBernstein(1,1)
\end{psgraph*}
\\ \hline  
\BSS{psBernstein}(0,0) \BSI{psBernstein}{pst-func} & \BSS{psBernstein}(0,1) & \BSS{psBernstein}(1,1)
\\ \hline 
\begin{psgraph*}[axesstyle=none,xticksize= -.5 1.5 ,yticksize=-.5 1.5,xticksize= -.5  1.5 , dx=.5,Dx=.5, dy=.5,Dy=.5 , subticks=0] (0,0)(-.5,-.5)(1.5,1.5){3.5cm}{3.5cm }    
\psBernstein(0,2)
\end{psgraph*}
&  
\begin{psgraph*}[axesstyle=none,xticksize= -.5 1.5 ,yticksize=-.5 1.5,xticksize= -.5  1.5 , dx=.5,Dx=.5, dy=.5,Dy=.5 , subticks=0] (0,0)(-.5,-.5)(1.5,1.5){3.5cm}{3.5cm }   
\psBernstein(1,2)
\end{psgraph*}
&  
\begin{psgraph*}[axesstyle=none,xticksize= -.5 1.5 ,yticksize=-.5 1.5,xticksize= -.5  1.5 , dx=.5,Dx=.5, dy=.5,Dy=.5 , subticks=0] (0,0)(-.5,-.5)(1.5,1.5){3.5cm}{3.5cm }    
\psBernstein(2,2)
\end{psgraph*}  
\\ \hline
\BSS{psBernstein}(0,2) & \BSS{psBernstein}(1,2) & \BSS{psBernstein}(2,2)
\\ \hline   
\begin{psgraph*}[axesstyle=none,xticksize= -.5 1.5 ,yticksize=-.5 1.5,xticksize= -.5  1.5 , dx=.5,Dx=.5, dy=.5,Dy=.5 , subticks=0] (0,0)(-.5,-.5)(1.5,1.5){3.5cm}{3.5cm } 
\psBernstein(0,3) 
\end{psgraph*}
&  
\begin{psgraph*}[axesstyle=none,xticksize= -.5 1.5 ,yticksize=-.5 1.5,xticksize= -.5  1.5 , dx=.5,Dx=.5, dy=.5,Dy=.5 , subticks=0] (0,0)(-.5,-.5)(1.5,1.5){3.5cm}{3.5cm }     
\psBernstein(1,3)
\psBernstein[linestyle=dotted](2,3)
\end{psgraph*}
&  
\begin{psgraph*}[axesstyle=none,xticksize= -.5 1.5 ,yticksize=-.5 1.5,xticksize= -.5  1.5 , dx=.5,Dx=.5, dy=.5,Dy=.5 , subticks=0] (0,0)(-.5,-.5)(1.5,1.5){3.5cm}{3.5cm }   
\psBernstein(3,3)
\end{psgraph*} 
\\ \hline 
\BSS{psBernstein}(0,3) & \BSS{psBernstein}(1,3) & \BSS{psBernstein}(3,3)\\
& \BSS{psBernstein}(2,3) & 
\\ \hline
\end{tabular}

\bigskip 

\begin{tabular}{|c|c|c|} \hline  
\begin{psgraph*}[axesstyle=none,xticksize= 0 1 ,yticksize=0 1 , dx=.2,Dx=.2, dy=.2,Dy=.2 , subticks=0](0,0)(0,0)(1,1){12cm}{4cm }
\psBernstein[linestyle=dotted](5,5)
\psBernstein[linestyle=dotted](4,5)
\psBernstein[linestyle=dotted](3,5)
\psBernstein[linestyle=dotted](2,5)
\psBernstein[linestyle=dotted](1,5)  
\psBernstein[linestyle=dotted](0,5) 
\psBernstein[envelope](0.1,5)
\end{psgraph*}
\\ \hline 
\BS{psBernstein}[\RDD{envelope}](0,5) \RDI{envelope}{pst-func}  
\\ \hline
\end{tabular} 

\newpage

%\subseZeroction{Zéros d'une fonction ou  point d'intersection de deux
fonction}
 \SbSSCT{Zéros d'une fonction ou  point d'intersection de deux
 fonction}{Zeros or intersections}
 
\begin{tabular}{|c|} \hline  
\begin{psgraph*}[axesstyle=none,xticksize= -1 2 ,yticksize=0 10](0,0)(0,-1)(10,2){10cm}{3cm} 
 \psplot[plotpoints=500,algebraic,linewidth=0.8pt]{0.001}{9.75}{cos(x)+.5}
 \psZero[algebraic](0.5,5){cos(x)+.5}{A}
% \psComment{->}(3,1)(N1){noeud N1}
\rput[B](3,1.5){ \ovalnode{B}{n\oe ud A}}
\ncline{->}{B}{A}
\end{psgraph*}
\\ \hline  
 \BSS{psZero}[algebraic](0.5,5)\AC{cos(x)+.5}\AC{A}
\\ \hline 
\end{tabular} 

\bigskip

\begin{tabular}{|c|} \hline  
\begin{psgraph*}[axesstyle=none,xticksize= -1 2 ,yticksize=0 10](0,0)(0,-1)(10,2){10cm}{3cm} 
 \psplot[plotpoints=500,algebraic,linestyle=dotted]{0.001}{9.75}{sin(x)}
 \psplot[plotpoints=500,algebraic,linewidth=0.8pt]{0.001}{9.75}{cos(x)+.5} 
 \psZero[algebraic](3,5){cos(x)+.5}[sin(x)]{N1}
% \psComment{->}(3,1)(N1){noeud N1}
\rput[B](3,1.5){ \ovalnode{B}{n\oe ud N1}}
\ncline{->}{B}{N1}
\end{psgraph*}
\\ \hline  
 \BSS{psZero}[algebraic](0.5,5)\AC{cos(x)+.5}[sin(x)]\AC{N1}
\\ \hline 
\end{tabular}

\bigskip

\begin{tabular}{|c|c|} \hline  
\multicolumn{2}{|c|}{ \BSS{psZero}[algebraic,\RDD{markZeros}](0.5,5)\AC{cos(x)+.5[sin(x)]}\AC{A} } 
\\ \hline 
\begin{psgraph*}[axesstyle=none,xticksize= -1 2 ,yticksize=0 4](0,0)(0,-1)(4,2){4cm}{3cm}
 \psplot[plotpoints=500,algebraic,linewidth=.1pt]{0.001}{9.75}{sin(x)} 
 \psplot[plotpoints=500,algebraic,linewidth=0.8pt]{0.001}{9.75}{cos(x)+.5}
 \psZero[algebraic,markZeros](0.5,5){cos(x)+.5}[sin(x)]{A}
\end{psgraph*}
&
\begin{psgraph*}[axesstyle=none,xticksize= -1 2 ,yticksize=0 4](0,0)(0,-1)(4,2){4cm}{3cm} 
 \psplot[plotpoints=500,algebraic,linewidth=.1pt]{0.001}{9.75}{sin(x)}
 \psplot[plotpoints=500,algebraic,linewidth=0.8pt]{0.001}{9.75}{cos(x)+.5}
 \psZero[algebraic,markZeros,onlyNode](0.5,5){cos(x)+.5}[sin(x)]{A}
\end{psgraph*}
\\ \hline 
\RDD{ markZeros} \RDI{markZeros}{pst-func} & \RDD{onlyNode}  \RDI{onlyNode}{pst-func}
\\ \hline 
\end{tabular} 

\bigskip

\begin{tabular}{|c|c|} \hline  
%\multicolumn{2}{|c|}{ \BSS{psZero}[algebraic,\RDD{markZeros}](0.5,5)\AC{cos(x)+.5[sin(x)]}\AC{A} } 
%\\ \hline 
\begin{psgraph*}[axesstyle=none,xticksize= -1 2 ,yticksize=0 4](0,0)(0,-1)(4,2){4cm}{3cm}
 \psplot[plotpoints=500,algebraic,linewidth=.1pt]{0.001}{9.75}{sin(x)} 
 \psplot[plotpoints=500,algebraic,linewidth=0.8pt]{0.001}{9.75}{cos(x)+.5}
 \psZero[algebraic,PrintCoord](0.5,5){cos(x)+.5}[sin(x)]{A}
\end{psgraph*}
&
\begin{psgraph*}[axesstyle=none,xticksize= -1 2 ,yticksize=0 4](0,0)(0,-1)(4,2){4cm}{3cm} 
 \psplot[plotpoints=500,algebraic,linewidth=.1pt]{0.001}{9.75}{sin(x)}
 \psplot[plotpoints=500,algebraic,linewidth=0.8pt]{0.001}{9.75}{cos(x)+.5}
 \psZero[algebraic,markZeros,onlyYVal](0.5,5){cos(x)+.5}[sin(x)]{A}
\end{psgraph*}
\\ \hline 
\RDD{PrintCoord}  \RDI{PrintCoord}{pst-func} & \RDD{onlyYVal}  \RDI{onlyYVal}{pst-func}
\\ \hline 
\end{tabular} 

\bigskip

\begin{tabular}{|c|c|} \hline  
%\multicolumn{2}{|c|}{ \BSS{psZero}[algebraic,\RDD{markZeros}](0.5,5)\AC{cos(x)+.5[sin(x)]}\AC{A} } 
%\\ \hline 
\begin{psgraph*}[axesstyle=none,xticksize= -1 2 ,yticksize=0 4](0,0)(0,-1)(4,2){4cm}{3cm}
 \psplot[plotpoints=500,algebraic,linewidth=.1pt]{0.001}{9.75}{sin(x)} 
 \psplot[plotpoints=500,algebraic,linewidth=0.8pt]{0.001}{9.75}{cos(x)+.5}
 \psZero[algebraic,PointName=Point,PrintCoord](0.5,5){cos(x)+.5}[sin(x)]{A}
\end{psgraph*}
&
\begin{psgraph*}[axesstyle=none,xticksize= -1 2 ,yticksize=0 4](0,0)(0,-1)(4,2){4cm}{3cm} 
 \psplot[plotpoints=500,algebraic,linewidth=.1pt]{0.001}{9.75}{sin(x)}
 \psplot[plotpoints=500,algebraic,linewidth=0.8pt]{0.001}{9.75}{cos(x)+.5}
 \psZero[algebraic,markZeros,originV,PrintCoord](0.5,5){cos(x)+.5}[sin(x)]{A}
\end{psgraph*}
\\ \hline 
\RDD{PointName},PrintCoord  \RDI{PointName}{pst-func} & \RDD{originV},PrintCoord  \RDI{originV}{pst-func}
\\ \hline 
\dft: PointName= I &
\\ \hline 
\end{tabular} 

\bigskip

\begin{tabular}{|c|c|} \hline  
%\multicolumn{2}{|c|}{ \BSS{psZero}[algebraic,\RDD{markZeros}](0.5,5)\AC{cos(x)+.5[sin(x)]}\AC{A} } 
%\\ \hline 
\begin{psgraph*}[axesstyle=none,xticksize= -1 2 ,yticksize=0 4](0,0)(0,-1)(4,2){4cm}{3cm}
 \psplot[plotpoints=500,algebraic,linewidth=.1pt]{0.001}{9.75}{sin(x)} 
 \psplot[plotpoints=500,algebraic,linewidth=0.8pt]{0.001}{9.75}{cos(x)+.5}
 \psZero[algebraic,PrintCoord,decimals=3](0.5,5){cos(x)+.5}[sin(x)]{A}
\end{psgraph*}
&
\begin{psgraph*}[axesstyle=none,xticksize= -1 2 ,yticksize=0 4](0,0)(0,-1)(4,2){4cm}{3cm} 
 \psplot[plotpoints=500,algebraic,linewidth=.1pt]{0.001}{9.75}{sin(x)}
 \psplot[plotpoints=500,algebraic,linewidth=0.8pt]{0.001}{9.75}{cos(x)+.5}
 \psZero[algebraic,markZeros,PrintCoord,ydecimals=4](0.5,5){cos(x)+.5}[sin(x)]{A}
\end{psgraph*}
\\ \hline 
\RDD{decimals}=3,PrintCoord  \RDI{originV}{pst-func}& \RDD{ydecimals}=4,PrintCoord \RDI{ydecimals}{pst-func}
%\\ \hline 
%\dft: PointName= I &
\\ \hline 
\end{tabular}


\bigskip

\begin{tabular}{|c|c|} \hline  
%\multicolumn{2}{|c|}{ \BSS{psZero}[algebraic,\RDD{markZeros}](0.5,5)\AC{cos(x)+.5[sin(x)]}\AC{A} } 
%\\ \hline 
\begin{psgraph*}[axesstyle=none,xticksize= -1 2 ,yticksize=0 4](0,0)(0,-1)(4,2){4cm}{3cm}
 \psplot[plotpoints=500,algebraic,linewidth=.1pt]{0.001}{9.75}{sin(x)} 
 \psplot[plotpoints=500,algebraic,linewidth=0.8pt]{0.001}{9.75}{cos(x)+.5}
 \psZero[algebraic,PrintCoord,xShift=.5](0.5,5){cos(x)+.5}[sin(x)]{A}
\end{psgraph*}
&
\begin{psgraph*}[axesstyle=none,xticksize= -1 2 ,yticksize=0 4](0,0)(0,-1)(4,2){4cm}{3cm} 
 \psplot[plotpoints=500,algebraic,linewidth=.1pt]{0.001}{9.75}{sin(x)}
 \psplot[plotpoints=500,algebraic,linewidth=0.8pt]{0.001}{9.75}{cos(x)+.5}
 \psZero[algebraic,markZeros,PrintCoord,yShift=.5](0.5,5){cos(x)+.5}[sin(x)]{A}
\end{psgraph*}
\\ \hline 
\RDD{xShift}=.5,PrintCoord  \RDI{xShift}{pst-func} & \RDD{yShift}=.5,PrintCoord  \RDI{yShift}{pst-func}
%\\ \hline 
%\dft: PointName= I &
\\ \hline 
\end{tabular} 


\bigskip

\begin{tabular}{|c|c|} \hline  
%\multicolumn{2}{|c|}{ \BSS{psZero}[algebraic,\RDD{markZeros}](0.5,5)\AC{cos(x)+.5[sin(x)]}\AC{A} } 
%\\ \hline 
\begin{psgraph*}[axesstyle=none,xticksize= -1 2 ,yticksize=0 4](0,0)(0,-1)(4,2){4cm}{3cm}
 \psplot[plotpoints=500,algebraic,linewidth=.1pt]{0.001}{9.75}{sin(x)} 
 \psplot[plotpoints=500,algebraic,linewidth=0.8pt]{0.001}{9.75}{cos(x)+.5}
 \psZero[algebraic,PrintCoord,yShift=.5,,postString=123](0.5,5){cos(x)+.5}[sin(x)]{A}
\end{psgraph*}
&
%\begin{psgraph*}[axesstyle=none,xticksize= -1 2 ,yticksize=0 4](0,0)(0,-1)(4,2){4cm}{3cm} 
%% \psplot[plotpoints=500,algebraic,linewidth=.1pt]{0.001}{9.75}{sin(x)}
% \psplot[plotpoints=500,algebraic,linewidth=0.8pt]{0.001}{9.75}{cos(x)+.5}
% \psZero[xShift=-0.2,yShift=0.15,postString=1,Newton](0.5,5){ x cos }{A}
%\end{psgraph*}
\\ \hline 
\RDD{postString}=123,PrintCoord  \RDI{postString}{pst-func} &
% \RDD{yShift}=.5,PrintCoord  \RDI{yShift}{pst-func}
%\\ \hline 
%\dft: PointName= I &
\\ \hline 
\end{tabular} 

\newpage

%\subsection{Fonction de Fourier}
 \SbSSCT{Fonction de Fourier}{Fourrier}
 
\begin{tabular}{|c|} \hline  
\begin{psgraph*}[axesstyle=none,xticksize= -3 3 ,yticksize=-10 10 , subticks=0 ](0,0)(-10,-3)(10,3){10cm}{6cm } 
 \psFourier[cosCoeff=0 1 -1 ]{-10}{10}
\end{psgraph*}
\\ \hline  
\BSS{psFourier}[\RDD{cosCoeff}=0 1 -1 ]\AC{-10}\AC{10}
\BSI{psFourier}{pst-func} \RDI{cosCoeff}{pst-func}
\\ \hline 
\dft : cosCoeff =0
\\ \hline 

\end{tabular}   


\bigskip
\begin{tabular}{|c|} \hline  
\begin{psgraph*}[axesstyle=none,xticksize= -2 2 ,yticksize=-5 5 , subticks=0 ](0,0)(-5,-2)(5,2){10cm}{4cm } 
 \psFourier[sinCoeff=1 .5 .33 .25 .2 .165 .14  .125 ]{-5}{5}
\end{psgraph*}
\\ \hline  
\BSS{psFourier}[\RDD{sinCoeff}=1 .5 .33 .25 .2 .165 .14  .125 ]\AC{-5}\AC{5}  \RDI{sinCoeff}{pst-func}
\\ \hline 
\dft : sinCoeff =1
\\ \hline 
\end{tabular}

\newpage
%\subsection{Fonction de Bessel}
 \SbSSCT{Fonction de Bessel}{Bessel}
 
\begin{tabular}{|c|} \hline  
  $\displaystyle J_n(x)=\frac{1}{\pi} \int_0^\pi \cos  (x \sin t-nt)  dt$ 
\\ \hline  
\begin{psgraph*}[axesstyle=none,xticksize= -1.5 1.5 ,yticksize=-20 20 , subticks=0,Dx=5 ](0,0)(-20,-1.5)(20,1.5){10cm}{4cm } 
 \psBessel{0}{-20}{20}
\end{psgraph*}  
\\ \hline  
n= 0 \hspace{1cm} \BSS{psBessel}\AC{0}\AC{-20}\AC{20} 
\\ \hline 
\begin{psgraph*}[axesstyle=none,xticksize= -1.5 1.5 ,yticksize=-20 20 , subticks=0,Dx=5 ](0,0)(-20,-1.5)(20,1.5){10cm}{4cm } 
 \psBessel{2}{-20}{20}
\end{psgraph*}  
\\ \hline  
n= 2 \hspace{1cm} \BSS{psBessel}\AC{2}\AC{-20}\AC{20} 
\\ \hline 
\end{tabular} 

\bigskip

\begin{tabular}{|c|} \hline  
  $\displaystyle f(x)=2.5 J_0(x) +sin(t)$ 
\\ \hline  
\begin{psgraph*}[axesstyle=none,xticksize= -3 3 ,yticksize=-20 20 , subticks=0,Dx=5 ](0,0)(-20,-3)(20,3){10cm}{4cm } 
 \psBessel[constI=2.5,constII={ t k  sin }]{0}{-20}{20}%
\end{psgraph*} 
\\ \hline 
\BS{psBessel}[\RDD{constI}=2.5,\RDD{constII}=\AC{ t k  sin }]\AC{0}\AC{-20}\AC{20} \RDI{constI}{pst-func}  \RDI{constII}{pst-func}
\\ \hline 
\end{tabular} 


\newpage
%\subsection{Fonction de Bessel modifiée}
\SbSSCT{Fonction de Bessel modifiée}{modified Bessel}

\begin{tabular}{|c|c|c|} \hline 
\multicolumn{3}{|c|}{\BSS{psModBessel}[yMaxValue=5,\RDD{nue}=0]\AC{0}\AC{5} \BSI{psModBessel}{pst-func} \RDI{nue}{pst-func}} \\ \hline
\begin{psgraph*}[axesstyle=none,xticksize= 0 5 ,yticksize=0 5](0,0)(0,0)(4,5){3.5cm}{5cm} 
 \psModBessel[yMaxValue=5,nue=0]{0}{5}
\end{psgraph*}
&  
\begin{psgraph*}[axesstyle=none,xticksize= 0 5 ,yticksize=0 5](0,0)(0,0)(4,5){3.5cm}{5cm} 
 \psModBessel[yMaxValue=5,nue=1]{0}{5}
\end{psgraph*}
&  
\begin{psgraph*}[axesstyle=none,xticksize= 0 5 ,yticksize=0 5](0,0)(0,0)(4,5){3.5cm}{5cm} 
 \psModBessel[yMaxValue=5,nue=2]{0}{5}
\end{psgraph*}
\\ \hline 
\RDD{nue}=0 & \RDD{nue}=1 & \RDD{nue}= 2
\\ \hline 
\multicolumn{3}{|c|}{\dft : nue=0} \\ \hline
\end{tabular} 

\newpage


%\subsection{Sinus intégral}
\SbSSCT{Sinus intégral}{Integral sinus}

\begin{tabular}{|c|} \hline  
\begin{psgraph*}[axesstyle=none,xticksize= -2 2 ,yticksize=-14 14, dx=2,Dx=2](0,0)(-14,-2)(14,2){10cm}{3cm} 
 \psSi{-14.5}{14.5}
\end{psgraph*}
\\ \hline  
 \BSS{psSi}\AC{-14.5}\AC{14.5}  \BSI{psSi}{pst-func}
\\ \hline 
\end{tabular} 

\bigskip

\begin{tabular}{|c|} \hline  
\begin{psgraph*}[axesstyle=none,xticksize= -4 1 ,yticksize=-14 14, dx=2,Dx=2](0,0)(-14,-4)(14,1){10cm}{3cm} 
 \pssi{-14.5}{14.5}
\end{psgraph*}
\\ \hline  
 \BSS{pssi}\AC{-14.5}\AC{14.5}  \BSI{pssi}{pst-func}
\\ \hline 
\end{tabular} 

%\subsection{Cosinus intégral}
\SbSSCT{Cosinus intégral}{Integral cosinus}

\begin{tabular}{|c|} \hline  
\begin{psgraph*}[axesstyle=none,xticksize= -4 1 ,yticksize=-12 12 , dx=2,Dx=2](0,0)(-12,-4)(12,1){10cm}{3cm} 
 \psCi[plotpoints=500]{-11.5}{11.5}
\end{psgraph*}
\\ \hline  
 \BSS{psCi}\AC{-11.5}\AC{11.5}  \BSI{psCi}{pst-func}
\\ \hline 
\end{tabular} 

\bigskip

\begin{tabular}{|c|} \hline  
\begin{psgraph*}[axesstyle=none,xticksize= 0 4 ,yticksize=-12 12 , dx=2,Dx=2](0,0)(-12,0)(12,4){10cm}{3cm} 
 \psci[plotpoints=500]{-11.5}{11.5}
\end{psgraph*}
\\ \hline  
 \BSS{psci}\AC{-11.5}\AC{11.5} \BSI{psci}{pst-func}
\\ \hline 
\end{tabular}

\newpage

%\subsection{Intégration et Convolution}
\SbSSCT{Intégration et Convolution}{Integration and convolution }

\begin{tabular}{|c|} \hline  
\begin{psgraph*}[axesstyle=none,xticksize= 0 1 ,yticksize=-6 6 , dx=2,Dx=2, dy=.5,Dy=.5](0,0)(-6,0)(6,1){10cm}{3cm} 
 \psplot[linestyle=dotted]{-6}{6}{x 0 2 GAUSS}
 \psCumIntegral{-6}{6}{0 2 GAUSS}
\end{psgraph*}
\\ \hline 
 \BS{psplot}[linestyle=dotted]\AC{-6}\AC{6}\AC{x 0 2 GAUSS} \\
 \BSS{psCumIntegral}\AC{-10}\AC{10}\AC{0 2 GAUSS}
 \BSI{psCumIntegral}{pst-func} 
\\ \hline 
\end{tabular}

\bigskip


\begin{tabular}{|c|} \hline  
\begin{psgraph*}[axesstyle=none,xticksize= 0 1 ,yticksize=-6 6 , dx=2,Dx=2, dy=.5,Dy=.5](0,0)(-6,0)(6,1){10cm}{3cm} 
 \psplot[linestyle=dotted]{-6}{6}{x 0 2 GAUSS}
 \psCumIntegral{0}{6}{0 2 GAUSS}
\end{psgraph*}
\\ \hline  
 \BSS{psCumIntegral}\AC{0}\AC{6}\AC{0 2 GAUSS}
\\ \hline 
\end{tabular}

\bigskip


\begin{tabular}{|c|} \hline  
\begin{psgraph*}[axesstyle=none,xticksize= 0 1 ,yticksize=-6 6 , dx=2,Dx=2, dy=.5,Dy=.5](0,0)(-6,0)(6,1){10cm}{3cm}  
 \psplot[linestyle=dotted]{-6}{6}{x 0 .5 GAUSS}
 \psIntegral{-6}{6}(-2,4){.5 GAUSS}
\end{psgraph*}
\\ \hline  
 \BSS{psIntegral}\AC{-2}\AC{4}\AC{.5 GAUSS} 
  \BSI{psIntegral}{pst-func} 
\\ \hline 
\end{tabular}

\bigskip

\begin{tabular}{|c|} \hline  
\begin{psgraph*}[axesstyle=none,xticksize= 0 1 ,yticksize=-6 6 , dx=2,Dx=2, dy=.5,Dy=.5](0,0)(-6,0)(6,1){10cm}{3cm}  
 \psplot[linestyle=dotted]{-6}{6}{x 0 .5 GAUSS}
 \psIntegral[Simpson=10]{-6}{6}(-2,4){.5 GAUSS}
\end{psgraph*}
\\ \hline  
 \BS{psIntegral}[\RDD{Simpson}=10]\AC{-2}\AC{4}\AC{.5 GAUSS}
\\ \hline 
\end{tabular}

\bigskip


\begin{tabular}{|c|} \hline  
\begin{psgraph*}[axesstyle=none,xticksize= 0 1.5 ,yticksize=-6 6 , dx=2,Dx=2, dy=.5,Dy=.5](0,0)(-6,0)(6,1.5){10cm}{3cm}
 \psplot[linestyle=dashed]{-5}{5}{x abs 2 le {0.5}{0} ifelse}
  \psplot[linestyle=dotted]{-5}{5}{x abs 1 le {0.75}{0} ifelse}
 \psConv{-5}{5}(-6,6) {abs 2 le {0.5}{0} ifelse}{abs 2 le {0.75}{0} ifelse}
\end{psgraph*}
\\ \hline  
 \BS{psplot}[linestyle=dashed]\AC{-5}\AC{5}\AC{x abs 2 le {0.5}{0} ifelse} \\ 
 \BS{psplot}[linestyle=dotted]\AC{-5}\AC{5}\AC{x abs 1 le {0.75}{0} ifelse} \\
 \BSS{psConv}\AC{-5}\AC{5}\AC(-6,6) \AC{abs 2 le {0.5}{0} ifelse}\AC{abs 2 le {0.75}{0} ifelse}
   \BSI{psConv}{pst-func} 
\\ \hline 
\end{tabular}

%====================


\newpage

%\subsection{Loi de Gauss}
\SbSSCT{Loi de Gauss}{Gauss Distribution}

\begin{tabular}{|c|c|} \hline  
\begin{psgraph*}[axesstyle=none,xticksize= -5 1 ,yticksize=-12 12 , subticks=0 ](0,0)(-2,0)(2,1){6cm}{2cm } 
 \psGauss{-2}{2}%
\end{psgraph*}
&  
\begin{psgraph*}[axesstyle=none,xticksize= -5 1 ,yticksize=-12 12 , subticks=0 ](0,0)(-2,0)(2,1){6cm}{2cm } 
 \psGaussI{-2}{2}%
\end{psgraph*}
\\ \hline  
 \BSS{psGauss}\AC{-2}\AC{2} \BSI{psGauss}{pst-func}
&  
 \BSS{psGaussI}\AC{-2}\AC{2} \BSI{psGaussI}{pst-func}
\\ 
\hline 
\end{tabular} 
\bigskip

\begin{tabular}{|c|c|} \hline  
\begin{psgraph*}[axesstyle=none,xticksize= -5 1 ,yticksize=-12 12 , subticks=0 ](0,0)(-2,0)(2,1){6cm}{2cm } 
 \psGauss[linestyle=dotted]{-2}{2}%
 \psGauss[mue=0.5]{-2}{2}%
\end{psgraph*}
&  
\begin{psgraph*}[axesstyle=none,xticksize= -5 1 ,yticksize=-12 12 , subticks=0 ](0,0)(-2,0)(2,1){6cm}{2cm }
 \psGauss[linestyle=dotted]{-2}{2}% 
 \psGauss[mue=-.5]{-2}{2}%
\end{psgraph*}
\\ \hline  
 \BSS{psGauss}[\RDD{mue}=0.5]\AC{-2}\AC{2}  \RDI{mue}{pst-func}
&  
 \BSS{psGauss}[\RDD{mue}=0.5]\AC{-2}\AC{2}
\\ 
\hline 
\end{tabular}

\bigskip

\begin{tabular}{|c|c|} \hline  
\begin{psgraph*}[axesstyle=none,xticksize= 0 2 ,yticksize=-2 2 , subticks=0 ](0,0)(-2,0)(2,2){6cm}{2cm } 
 \psGauss[linestyle=dotted]{-2}{2}%
 \psGauss[sigma=.25]{-2}{2}%
\end{psgraph*}
&  
\begin{psgraph*}[axesstyle=none,xticksize= 0 2 ,yticksize=-2 2 , subticks=0 ](0,0)(-2,0)(2,2){6cm}{2cm }
 \psGauss[linestyle=dotted]{-2}{2}% 
 \psGauss[sigma=1]{-2}{2}%
\end{psgraph*}
\\ \hline 
 
 \BSS{psGauss}[\RDD{sigma}=0.25]\AC{-2}\AC{2} \RDI{sigma}{pst-func}
&  
 \BSS{psGauss}[\RDD{sigma}=1]\AC{-2}\AC{2}
\\ \hline 
\end{tabular}

\newpage 

%\subsection{Loi binomiale}
\SbSSCT{Loi binomiale}{Binomial Distribution}

\begin{tabular}{|c|c|c|} \hline  
\begin{psgraph*}[axesstyle=none,xticksize= 0 1 ,yticksize=-1 3 , subticks=0, dy=.2,Dy=.2](0,0)(-1,0)(3,1){3cm}{2cm }
\psBinomial{2}{0.5}
\end{psgraph*}
&  
\begin{psgraph*}[axesstyle=none,xticksize= 0 1 ,yticksize=-1 3 , subticks=0, dy=.2,Dy=.2 ](0,0)(-1,0)(3,1){3cm}{2cm }
\psBinomial{2}{0.25}
\end{psgraph*}
&  
\begin{psgraph*}[axesstyle=none,xticksize= 0 1 ,yticksize=-1 3 , subticks=0, dy=.2,Dy=.2 ](0,0)(-1,0)(3,1){3cm}{2cm }
\psBinomial{2}{0.75}
\end{psgraph*}
\\ \hline  
\BSS{psBinomial}\AC{2}\AC{0.5}  \BSI{psBinomial}{pst-func}
&  
\BSS{psBinomial}\AC{2}\AC{0.25}
&  
\BSS{psBinomial}\AC{2}\AC{0.75}
\\ \hline 
\end{tabular} 

\bigskip

\begin{tabular}{|c|c|} \hline  
  
\begin{psgraph*}[axesstyle=none,xticksize= 0 1 ,yticksize=-1 5 , subticks=0, dy=.2,Dy=.2 ](0,0)(-1,0)(5,1){5cm}{2cm }
\psBinomial{3}{0.5}
\end{psgraph*}
&  
\begin{psgraph*}[axesstyle=none,xticksize= 0 1 ,yticksize=-1 5 , subticks=0, dy=.2,Dy=.2 ](0,0)(-1,0)(5,1){5cm}{2cm }
\psBinomial{4}{0.5}
\end{psgraph*}
\\ \hline  
 
\BSS{psBinomial}\AC{3}\AC{0.5}
&  
\BSS{psBinomial}\AC{4}\AC{0.5}
\\ \hline 
\end{tabular}

\bigskip

\begin{tabular}{|c|c|} \hline  
  
\begin{psgraph*}[axesstyle=none,xticksize= 0 1 ,yticksize=-1 5 , subticks=0, dy=.2,Dy=.2 ](0,0)(-1,0)(5,1){5cm}{2cm }
\psBinomial[linestyle=dotted]{4}{0.5}
\psBinomial{2,4}{0.5}
\end{psgraph*}
&  
\begin{psgraph*}[axesstyle=none,xticksize= 0 1 ,yticksize=-1 5 , subticks=0, dy=.2,Dy=.2 ](0,0)(-1,0)(5,1){5cm}{2cm }
\psBinomial[linestyle=dotted]{4}{0.5}
\psBinomial{1,2,4}{0.5}
\end{psgraph*}
\\ \hline  
 
\BSS{psBinomial}\AC{2,4}\AC{0.5}
&  
\BSS{psBinomial}\AC{1,2,4}\AC{0.5}
\\ \hline 
\end{tabular}

\bigskip

\begin{tabular}{|c|c|} \hline  
  
\begin{psgraph*}[axesstyle=none,xticksize= 0 1 ,yticksize=-3 3 , subticks=0, dy=.2,Dy=.2 ](0,0)(-3,0)(3,1){5cm}{2cm }
\psBinomialN{3}{0.5}
\end{psgraph*}
&  
\begin{psgraph*}[axesstyle=none,xticksize= 0 1 ,yticksize=-3 3 , subticks=0, dy=.2,Dy=.2 ](0,0)(-3,0)(3,1){5cm}{2cm }
\psBinomialN{4}{0.5}
\end{psgraph*}
\\ \hline  
 
\BSS{psBinomialN}\AC{3}\AC{0.5} \BSI{psBinomialN}{pst-func}
&  
\BSS{psBinomialN}\AC{4}\AC{0.5}
\\ \hline 
\end{tabular}

\newpage

\subsubsection{paramètres}

\begin{tabular}{|c|c|} \hline  
  
\begin{psgraph*}[axesstyle=none,xticksize= 0 1 ,yticksize=-1 5 , subticks=0, dy=.2,Dy=.2 ](0,0)(-1,0)(5,1){5cm}{2cm }
\psBinomial[printValue]{3}{0.5}
\end{psgraph*}
&  
\begin{psgraph*}[axesstyle=none,xticksize= 0 1 ,yticksize=-1 5 , subticks=0, dy=.2,Dy=.2 ](0,0)(-1,0)(5,1){5cm}{2cm }
\psBinomial[markZeros]{4}{0.5}
\end{psgraph*}
\\ \hline  
 
\BSS{psBinomial}[\RDD{printValue}]\AC{3}\AC{0.5} \RDI{printValue}{pst-func}
&  
\BS{psBinomial}[\RDD{markZeros}]\AC{4}\AC{0.5}
\RDI{markZeros}{pst-func}
\\ \hline 
\end{tabular}


\bigskip

\begin{tabular}{|c|c|} \hline  
  
\begin{psgraph*}[axesstyle=none,xticksize= 0 1 ,yticksize=-1 5 , subticks=0, dy=.2,Dy=.2 ](0,0)(-1,0)(5,1){5cm}{2cm }
\psBinomial[fillcolor=yellow,fillstyle=solid,]{3}{0.5}
\end{psgraph*}
&  
\begin{psgraph*}[axesstyle=none,xticksize= 0 1 ,yticksize=-1 5 , subticks=0, dy=.2,Dy=.2 ](0,0)(-1,0)(5,1){5cm}{2cm }
\psBinomial[barwidth=0.5]{4}{0.5} 
\end{psgraph*}
\\ \hline  
 
\BSS{psBinomial}[fillcolor=yellow]\AC{3}\AC{0.5}
&  
\BS{psBinomial}[\RDD{barwidth}=0.5]\AC{4}\AC{0.5} \RDI{barwidth}{pst-func}
\\ \hline 
\end{tabular}

\bigskip

\begin{tabular}{|c|c|} \hline  
  
\begin{psgraph*}[axesstyle=none,xticksize= 0 1 ,yticksize=-1 5 , subticks=0, dy=.2,Dy=.2 ](0,0)(-1,0)(5,1){5cm}{2cm }
\psBinomial[fillstyle=vlines,markZeros]{3}{0.5}
\end{psgraph*}
&  
\begin{psgraph*}[axesstyle=none,xticksize= 0 1 ,yticksize=-1 5 , subticks=0, dy=.2,Dy=.2 ](0,0)(-1,0)(5,1){5cm}{2cm }
\psBinomial[barwidth=0.5,markZeros]{4}{0.5}
\end{psgraph*}
\\ \hline  
 
[fillstyle=vlines,,markZeros]
&  
[barwidth=0.5,markZeros]
\\ \hline 
\end{tabular}

\newpage
%\subsection{Loi de Poisson}
\SbSSCT{Loi de Poisson}{Poisson Distribution}

\begin{tabular}{|c|c|c|} \hline  
\begin{psgraph*}[axesstyle=none,xticksize= 0 1 ,yticksize=-1 5 , subticks=0, dy=.2,Dy=.2 ](0,0)(-1,0)(3,1){3cm}{2cm }
 \psPoisson{2}{1}
\end{psgraph*}
&  
\begin{psgraph*}[axesstyle=none,xticksize= 0 1 ,yticksize=-1 5 , subticks=0, dy=.2,Dy=.2 ](0,0)(-1,0)(5,1){3cm}{2cm }
 \psPoisson{3}{1}
\end{psgraph*}
&  
\begin{psgraph*}[axesstyle=none,xticksize= 0 1 ,yticksize=-1 5 , subticks=0, dy=.2,Dy=.2 ](0,0)(-1,0)(5,1){3cm}{2cm }
 \psPoisson{4}{1}
\end{psgraph*}
\\ \hline  
\BSS{psPoisson}\AC{2}\AC{1} \BSI{psPoisson}{pst-func}  & \BSS{psPoisson}\AC{3}\AC{1}   &  \BSS{psPoisson}\AC{4}\AC{1} \\ 
\hline 
\end{tabular} 

\bigskip

\begin{tabular}{|c|c|c|} \hline  
\begin{psgraph*}[axesstyle=none,xticksize= 0 1 ,yticksize=-1 5 , subticks=0, dy=.2,Dy=.2 ](0,0)(-1,0)(5,1){3cm}{2cm }
 \psPoisson{4}{2}
\end{psgraph*}
&  
\begin{psgraph*}[axesstyle=none,xticksize= 0 1 ,yticksize=-1 5 , subticks=0, dy=.2,Dy=.2 ](0,0)(-1,0)(5,1){3cm}{2cm }
 \psPoisson{4}{3}
\end{psgraph*}
&  
\begin{psgraph*}[axesstyle=none,xticksize= 0 1 ,yticksize=-1 5 , subticks=0, dy=.2,Dy=.2 ](0,0)(-1,0)(5,1){3cm}{2cm }
 \psPoisson{4}{4}
\end{psgraph*}
\\ \hline  
\BSS{psPoisson}\AC{4}\AC{2}   & \BSS{psPoisson}\AC{4}\AC{3}   &  \BSS{psPoisson}\AC{4}\AC{4} \\ 
\hline 
\end{tabular} 


\bigskip

\begin{tabular}{|c|c|c|} \hline  
\begin{psgraph*}[axesstyle=none,xticksize= 0 1 ,yticksize=-1 5 , subticks=0, dy=.2,Dy=.2 ](0,0)(-1,0)(5,1){3cm}{2cm }
 \psPoisson[linestyle=dotted]{4}{2}
 \psPoisson{1,4}{2}
\end{psgraph*}
&  
\begin{psgraph*}[axesstyle=none,xticksize= 0 1 ,yticksize=-1 5 , subticks=0, dy=.2,Dy=.2 ](0,0)(-1,0)(5,1){3cm}{2cm }
 \psPoisson[linestyle=dotted]{4}{2}
 \psPoisson{2,4}{2}
\end{psgraph*}
&  
\begin{psgraph*}[axesstyle=none,xticksize= 0 1 ,yticksize=-1 5 , subticks=0, dy=.2,Dy=.2 ](0,0)(-1,0)(5,1){3cm}{2cm }
 \psPoisson[linestyle=dotted]{4}{2}
 \psPoisson{3,4}{2}
\end{psgraph*}
\\ \hline  
\BSS{psPoisson}\AC{1,4}\AC{2}   & \BSS{psPoisson}\AC{2,4}\AC{2}   &  \BSS{psPoisson}\AC{3,4}\AC{2} \\ 
\hline 
\end{tabular}

\bigskip

\begin{tabular}{|c|c|c|} \hline  
\begin{psgraph*}[axesstyle=none,xticksize= 0 1 ,yticksize=-1 5 , subticks=0, dy=.2,Dy=.2 ](0,0)(-1,0)(5,1){3cm}{2cm }
 \psPoisson[markZeros]{4}{2}
\end{psgraph*}
&  
\begin{psgraph*}[axesstyle=none,xticksize= 0 1 ,yticksize=-1 5 , subticks=0, dy=.2,Dy=.2 ](0,0)(-1,0)(5,1){3cm}{2cm }
 \psPoisson[printValue]{4}{2}
\end{psgraph*}
&  
\begin{psgraph*}[axesstyle=none,xticksize= 0 1 ,yticksize=-1 5 , subticks=0, dy=.2,Dy=.2 ](0,0)(-1,0)(5,1){3cm}{2cm }
 \psPoisson[barwidth=0.5]{4}{2}
\end{psgraph*}
\\ \hline  
\BS{psPoisson}[\RDD{markZeros}]\AC{4}\AC{2}   & \BS{psPoisson}[\RDD{printValue}]\AC{4}\AC{2}   &  \BS{psPoisson}[\RDD{barwidth}=0.5]\AC{4}\AC{2} \\ 
\hline 
\end{tabular}

\newpage
%\subsection{Loi Gamma }
\SbSSCT{Loi Gamma}{Gamma Distribution}

\begin{center}
\begin{tabular}{|c|} \hline  
\begin{psgraph*}[axesstyle=none,xticksize= 0 1 ,yticksize=0 3.5 , subticks=0, dy=.2,Dy=.2 ](0,0)(0,0)(3.5,1){5cm}{3cm }
 \psGammaDist{0.1}{3}
\end{psgraph*}
\\ \hline  
\BSS{psGammaDist}\AC{0.1}\AC{3}  \BSI{psGammaDist}{pst-func}
\\ \hline 
\end{tabular} 
\end{center}



 
\bigskip
\begin{tabular}{|c|c|} \hline  
\begin{psgraph*}[axesstyle=none,xticksize= 0 1 ,yticksize=0 3.5 , subticks=0, dy=.2,Dy=.2 ](0,0)(0,0)(3.5,1){5cm}{3cm }
 \psGammaDist[linestyle=dotted]{0.1}{3}
 \psGammaDist[alpha=0.25]{0.1}{3}
\end{psgraph*}

&  
\begin{psgraph*}[axesstyle=none,xticksize= 0 1 ,yticksize=0 3.5 , subticks=0, dy=.2,Dy=.2 ](0,0)(0,0)(3.5,1){5cm}{3cm }
 \psGammaDist[linestyle=dotted]{0.1}{3}
 \psGammaDist[alpha=0.75]{0.1}{3}
\end{psgraph*}
\\ \hline 
 \BS{psGammaDist}[\RDD{alpha}=0.25]\AC{0.1}\AC{3} \RDI{alpha}{pst-func} &
  \BS{psGammaDist}[\RDD{alpha}=0.75]\AC{0.1}\AC{3}  \\ 
\hline 
\end{tabular} 

\bigskip
\begin{tabular}{|c|c|} \hline  
\begin{psgraph*}[axesstyle=none,xticksize= 0 1 ,yticksize=0 3.5 , subticks=0, dy=.2,Dy=.2 ](0,0)(0,0)(3.5,1){5cm}{3cm }
 \psGammaDist[linestyle=dotted]{0.1}{3}
 \psGammaDist[beta=0.25]{0.1}{3}
\end{psgraph*}

&  
\begin{psgraph*}[axesstyle=none,xticksize= 0 1 ,yticksize=0 3.5 , subticks=0, dy=.2,Dy=.2 ](0,0)(0,0)(3.5,1){5cm}{3cm }
 \psGammaDist[linestyle=dotted]{0.1}{3}
 \psGammaDist[beta=0.75]{0.1}{3}
\end{psgraph*}
\\ \hline 
 \BS{psGammaDist}[\RDD{beta}=0.25]\AC{0.1}\AC{3} \RDI{beta}{pst-func} &
  \BS{psGammaDist}[\RDD{beta}=0.75]\AC{0.1}\AC{3}  \\ 
\hline 
\end{tabular} 


\bigskip
\begin{tabular}{|c|c|} \hline  
\begin{psgraph*}[axesstyle=none,xticksize= 0 1 ,yticksize=0 3.5 , subticks=0, dy=.2,Dy=.2 ](0,0)(0,0)(3.5,1){5cm}{3cm }
 \psGammaDist[linestyle=dotted]{0.1}{3}
 \psGammaDist[alpha=0.25,beta=0.75]{0.1}{3}
\end{psgraph*}

&  
\begin{psgraph*}[axesstyle=none,xticksize= 0 1 ,yticksize=0 3.5 , subticks=0, dy=.2,Dy=.2 ](0,0)(0,0)(3.5,1){5cm}{3cm }
 \psGammaDist[linestyle=dotted]{0.1}{3}
 \psGammaDist[alpha=0.75,beta=0.25]{0.1}{3}
\end{psgraph*}
\\ \hline 
[alpha=0.25,beta=0.75] &
[alpha=0.75,beta=0.25]  \\ 
\hline 
\end{tabular} 

\newpage
%\subsection{Loi du $\chi^2$}
\SbSSCT{Loi du $\chi^2$}{$\chi^2$ Distribution}

\begin{center}
\begin{tabular}{|c|} \hline  
\begin{psgraph*}[axesstyle=none,xticksize= 0 1 ,yticksize=0 5 , subticks=0, dy=.2,Dy=.2 ](0,0)(0,0)(5,1){10cm}{3cm }
\psChiIIDist{0.01}{5}
\end{psgraph*}
\\ \hline  
\BSS{psChiIIDist}\AC{0.01}\AC{5} \BSI{psChiIIDist}{pst-func}
\\ \hline 
\end{tabular} 
\end{center}

\begin{tabular}{|c|} \hline  
\begin{psgraph*}[axesstyle=none,xticksize= 0 1 ,yticksize=0 5 , subticks=0, dy=.2,Dy=.2 ](0,0)(0,0)(5,1){10cm}{3cm }
\psChiIIDist[linestyle=dotted]{0.01}{5}
\psChiIIDist[nue=.5]{0.01}{5}
\end{psgraph*}
\\ \hline  
\BS{psChiIIDist}[\RDD{nue}=.5]\AC{0.01}\AC{5} \RDI{nue}{pst-func}
\\ \hline 
\end{tabular}

\bigskip
\begin{tabular}{|c|} \hline  
\begin{psgraph*}[axesstyle=none,xticksize= 0 1 ,yticksize=0 5 , subticks=0, dy=.2,Dy=.2 ](0,0)(0,0)(5,1){10cm}{3cm }
\psChiIIDist[linestyle=dotted]{0.01}{5}
\psChiIIDist[nue=2]{0.01}{5}
\end{psgraph*}
\\ \hline  
\BS{psChiIIDist}[\RDD{nue}=2]\AC{0.01}\AC{5}
\\ \hline 
\end{tabular}

\bigskip
\begin{tabular}{|c|} \hline  
\begin{psgraph*}[axesstyle=none,xticksize= 0 1 ,yticksize=0 5 , subticks=0, dy=.2,Dy=.2 ](0,0)(0,0)(5,1){10cm}{3cm }
\psChiIIDist[linestyle=dotted]{0.01}{5}
\psChiIIDist[nue=3]{0.01}{5}
\end{psgraph*}
\\ \hline  
\BS{psChiIIDist}[\RDD{nue}=3]\AC{0.01}\AC{5}
\\ \hline 
\end{tabular}

%\subsection{Loi de Student}
\SbSSCT{Loi de Student}{Student Distribution}

\begin{tabular}{|c|} \hline  
\begin{psgraph*}[axesstyle=none,xticksize= 0 1 ,yticksize=-4 4 , subticks=0, dy=.2,Dy=.2 ](0,0)(-4,0)(4,1){10cm}{3cm }
 \psTDist{-4}{4}
\end{psgraph*}
\\ \hline  
\BSS{psTDist}\AC{4}\AC{4} \BSI{psTDist}{pst-func}
\\ \hline 
\end{tabular} 

\bigskip
\begin{tabular}{|c|} \hline  
\begin{psgraph*}[axesstyle=none,xticksize= 0 1 ,yticksize=-4 4 , subticks=0, dy=.2,Dy=.2 ](0,0)(-4,0)(4,1){10cm}{3cm }
 \psTDist[linestyle=dotted]{-4}{4}
 \psTDist[nue=.5]{-4}{4}
\end{psgraph*}
\\ \hline  
\BSS{psTDist}[\RDD{nue}=.5]\AC{4}\AC{4} \RDI{nue}{pst-func}
\\ \hline 
\end{tabular}

\bigskip
\begin{tabular}{|c|} \hline  
\begin{psgraph*}[axesstyle=none,xticksize= 0 1 ,yticksize=-4 4 , subticks=0, dy=.2,Dy=.2 ](0,0)(-4,0)(4,1){10cm}{3cm }
 \psTDist[linestyle=dotted]{-4}{4}
 \psTDist[nue=10]{-4}{4}
\end{psgraph*}
\\ \hline  
\BSS{psTDist}[\RDD{nue}=10]\AC{4}\AC{4} \\ \hline 
\end{tabular}



\newpage
%\subsection{Loi de F}
\SbSSCT{Loi de F}{F Distribution}

\begin{tabular}{|c|} \hline  
\begin{psgraph*}[axesstyle=none,xticksize= 0 1 ,yticksize=0 5 , subticks=0, dy=.2,Dy=.2 ](0,0)(0,0)(5,1){10cm}{3cm }
\psFDist{0.1}{5}
\end{psgraph*}
\\ \hline  
\BSS{psFDist}\AC{0.1}\AC{5} \BSI{psFDist}{pst-func}
\\ \hline 
\end{tabular} 


\bigskip

\begin{tabular}{|c|} \hline  
\begin{psgraph*}[axesstyle=none,xticksize= 0 1 ,yticksize=0 5 , subticks=0, dy=.2,Dy=.2 ](0,0)(0,0)(5,1){10cm}{3cm }
\psFDist[linestyle=dotted]{0.1}{5}
\psFDist[nue=3]{0.01}{5}
\end{psgraph*}

\\ \hline  
\BSS{psFDist}[\RDD{nue}=3]\AC{0.1}\AC{5} \RDI{nue}{pst-func}
\\ \hline 
\end{tabular}



\bigskip

\begin{tabular}{|c|} \hline  
\begin{psgraph*}[axesstyle=none,xticksize= 0 1 ,yticksize=0 5 , subticks=0, dy=.2,Dy=.2 ](0,0)(0,0)(5,1){10cm}{3cm }
\psFDist[linestyle=dotted]{0.1}{5}
\psFDist[mue=12]{0.01}{5}
\end{psgraph*}
\\ \hline  
\BSS{psFDist}[\RDD{mue}=12]\AC{0.1}\AC{5} \RDI{mue}{pst-func}
\\ \hline 
\end{tabular}



\bigskip

\begin{tabular}{|c|} \hline  
\begin{psgraph*}[axesstyle=none,xticksize= 0 1 ,yticksize=0 5 , subticks=0, dy=.2,Dy=.2 ](0,0)(0,0)(5,1){10cm}{3cm }
\psFDist[linestyle=dotted]{0.1}{5}
\psFDist[nue=3,mue=12]{0.01}{5}
\end{psgraph*}
\\ \hline  
\BSS{psFDist}[nue=3,mue=12]\AC{0.1}\AC{5} \RDI{mue}{pst-func}
\\ \hline 
\end{tabular}

\newpage

%\subsection{Loi de Beta}
\SbSSCT{Loi de Beta}{Beta Distribution}

\begin{center}
\begin{tabular}{|c|} \hline  
\begin{psgraph*}[axesstyle=none,xticksize= 0 2 ,yticksize=0 1 , subticks=0, dy=.2,Dy=.2 ](0,0)(0,0)(1,2){5cm}{3cm }
 \psBetaDist{0.01}{0.99}
\end{psgraph*}
\\ \hline  
\BSS{psBetaDist}\AC{0.01}\AC{0.99} \BSI{psBetaDist}{pst-func}
\\ \hline 
\end{tabular}
\end{center}



\bigskip

\begin{tabular}{|c|c|c|} \hline 
\multicolumn{3}{|c|}{ \BSS{psBetaDist}[alpha=0.1]\AC{0.01}\AC{0.99} }
\\ \hline 
\begin{psgraph*}[axesstyle=none,xticksize= 0 2 ,yticksize=0 1 , subticks=0, dy=.2,Dy=.2 ](0,0)(0,0)(1,2){3cm}{3cm }
 \psBetaDist[alpha=0.1]{0.01}{0.99}
\end{psgraph*}
&  
\begin{psgraph*}[axesstyle=none,xticksize= 0 2 ,yticksize=0 1 , subticks=0, dy=.2,Dy=.2 ](0,0)(0,0)(1,2){3cm}{3cm }
 \psBetaDist[alpha=.5]{0.01}{0.99}
\end{psgraph*}
&  
\begin{psgraph*}[axesstyle=none,xticksize= 0 2 ,yticksize=0 1 , subticks=0, dy=.2,Dy=.2 ](0,0)(0,0)(1,2){3cm}{3cm }
 \psBetaDist[alpha=.9]{0.01}{0.99}
\end{psgraph*}
\\ \hline 
[alpha=0.1] & [alpha=0.5] & [alpha=0.9] \\ 
\hline 
\multicolumn{3}{|c|}{ \dft : alpha= 1 }
\\ \hline
\end{tabular} 

\bigskip

\begin{tabular}{|c|c|c|} \hline 
\multicolumn{3}{|c|}{ \BSS{psBetaDist}[beta=0.1]\AC{0.01}\AC{0.99} }
\\ \hline 
\begin{psgraph*}[axesstyle=none,xticksize= 0 2 ,yticksize=0 1 , subticks=0, dy=.2,Dy=.2 ](0,0)(0,0)(1,2){3cm}{3cm }
 \psBetaDist[beta=0.1]{0.01}{0.99}
\end{psgraph*}
&  
\begin{psgraph*}[axesstyle=none,xticksize= 0 2 ,yticksize=0 1 , subticks=0, dy=.2,Dy=.2 ](0,0)(0,0)(1,2){3cm}{3cm }
 \psBetaDist[beta=.5]{0.01}{0.99}
\end{psgraph*}
&  
\begin{psgraph*}[axesstyle=none,xticksize= 0 2 ,yticksize=0 1 , subticks=0, dy=.2,Dy=.2 ](0,0)(0,0)(1,2){3cm}{3cm }
 \psBetaDist[beta=.9]{0.01}{0.99}
\end{psgraph*}
\\ \hline 
[beta=0.1] & [beta=0.5] & [beta=0.9] \\ 
\hline 
\multicolumn{3}{|c|}{ \dft : beta= 1 }
\\ \hline
\end{tabular} 

\bigskip

\begin{tabular}{|c|c|c|} \hline 
\multicolumn{3}{|c|}{ \BSS{psBetaDist}[beta=0.1]\AC{0.01}\AC{0.99} }
\\ \hline 
\begin{psgraph*}[axesstyle=none,xticksize= 0 2 ,yticksize=0 1 , subticks=0, dy=.2,Dy=.2 ](0,0)(0,0)(1,2){3cm}{3cm }
 \psBetaDist[alpha=.1,beta=0.1]{0.01}{0.99}
\end{psgraph*}
&  
\begin{psgraph*}[axesstyle=none,xticksize= 0 2 ,yticksize=0 1 , subticks=0, dy=.2,Dy=.2 ](0,0)(0,0)(1,2){3cm}{3cm }
 \psBetaDist[alpha=.1,beta=.5]{0.01}{0.99}
\end{psgraph*}
&  
\begin{psgraph*}[axesstyle=none,xticksize= 0 2 ,yticksize=0 1 , subticks=0, dy=.2,Dy=.2 ](0,0)(0,0)(1,2){3cm}{3cm }
 \psBetaDist[alpha=.1,beta=.9]{0.01}{0.99}
\end{psgraph*}
\\ \hline 
[alpha=.1,beta=0.1] & [alpha=.1,beta=0.5] & [alpha=.1,beta=0.9] \\ 
\hline 
\end{tabular}

\newpage

%\subsection{Loi de Cauchy}
\SbSSCT{Loi de Cauchy}{Cauchy Distribution}

\begin{center}
\begin{tabular}{|c|} \hline  
\begin{psgraph*}[axesstyle=none,xticksize= 0 1 ,yticksize=-3 3 , subticks=0, dy=.2,Dy=.2 ](0,0)(-3,0)(3,1){6cm}{2cm }
\psCauchy{-3}{3}
\end{psgraph*}
\\ \hline  
\BSS{psCauchy}\AC{-3}\AC{3} \BSI{psCauchy}{pst-func}
\\ \hline 
\end{tabular}
\end{center}

\bigskip

\begin{tabular}{|c|c|c|} \hline 
\multicolumn{3}{|c|}{ \BSS{psCauchy}[b=0.1]\AC{-3}\AC{3} }
\\ \hline 
\begin{psgraph*}[axesstyle=none,xticksize= 0 1 ,yticksize=-3  3 , subticks=0, dy=.2,Dy=.2 ](0,0)(-3,0)(3,1){3cm}{3cm }
\psCauchy[b=0.1]{-3}{3}
\end{psgraph*}
&  
\begin{psgraph*}[axesstyle=none,xticksize= 0 1 ,yticksize=-3  3 , subticks=0, dy=.2,Dy=.2 ](0,0)(-3,0)(3,1){3cm}{3cm }
\psCauchy[b=.5]{-3}{3}
\end{psgraph*}
&  
\begin{psgraph*}[axesstyle=none,xticksize= 0 1 ,yticksize=-3  3 , subticks=0, dy=.2,Dy=.2 ](0,0)(-3,0)(3,1){3cm}{3cm }
\psCauchy[b=1]{-3}{3}
\end{psgraph*}
\\ \hline 
[b=0.1]] & [b=0.5] & [b=1] \\ 
\hline 
\multicolumn{3}{|c|}{ \dft : b = 1 }
\\ \hline
\end{tabular}


\bigskip

\begin{tabular}{|c|c|c|} \hline 
\multicolumn{3}{|c|}{ \BSS{psCauchy}[m=0.1]\AC{-3}\AC{3} }
\\ \hline 
\begin{psgraph*}[axesstyle=none,xticksize= 0 1 ,yticksize=-3  3 , subticks=0, dy=.2,Dy=.2 ](0,0)(-3,0)(3,1){3cm}{3cm }
\psCauchy[m=-1]{-3}{3}
\end{psgraph*}
&  
\begin{psgraph*}[axesstyle=none,xticksize= 0 1 ,yticksize=-3  3 , subticks=0, dy=.2,Dy=.2 ](0,0)(-3,0)(3,1){3cm}{3cm }
\psCauchy[m=0]{-3}{3}
\end{psgraph*}
&  
\begin{psgraph*}[axesstyle=none,xticksize= 0 1 ,yticksize=-3  3 , subticks=0, dy=.2,Dy=.2 ](0,0)(-3,0)(3,1){3cm}{3cm }
\psCauchy[m=1]{-3}{3}
\end{psgraph*}
\\ \hline 
[m=-1]] & [m=0] & [m=1] \\ 
\hline 
\multicolumn{3}{|c|}{ \dft : m = 0 }
\\ \hline
\end{tabular}

\newpage

\begin{center}
\begin{tabular}{|c|} \hline  
\begin{psgraph*}[axesstyle=none,xticksize= 0 1 ,yticksize=-3 3 , subticks=0, dy=.2,Dy=.2 ](0,0)(-3,0)(3,1){6cm}{2cm }
\psCauchyI{-3}{3}
\end{psgraph*}
\\ \hline  
\BSS{psCauchyI}\AC{-3}\AC{3} \BSI{psCauchyI}{pst-func}
\\ \hline 
\end{tabular}
\end{center}




\bigskip

\begin{tabular}{|c|c|c|} \hline 
\multicolumn{3}{|c|}{ \BSS{psCauchyI}[b=0.1]\AC{-3}\AC{3} }
\\ \hline 
\begin{psgraph*}[axesstyle=none,xticksize= 0 1 ,yticksize=-3  3 , subticks=0, dy=.2,Dy=.2 ](0,0)(-3,0)(3,1){3cm}{3cm }
\psCauchyI[linestyle=dotted]{-2.5}{2.5}
\psCauchyI[b=0.1]{-3}{3}
\end{psgraph*}
&  
\begin{psgraph*}[axesstyle=none,xticksize= 0 1 ,yticksize=-3  3 , subticks=0, dy=.2,Dy=.2 ](0,0)(-3,0)(3,1){3cm}{3cm }
\psCauchyI[linestyle=dotted]{-3}{3}
\psCauchyI[b=.5]{-3}{3}
\end{psgraph*}
&  
\begin{psgraph*}[axesstyle=none,xticksize= 0 1 ,yticksize=-3  3 , subticks=0, dy=.2,Dy=.2 ](0,0)(-3,0)(3,1){3cm}{3cm }
\psCauchyI[linestyle=dotted]{-2.5}{2.5}
\psCauchyI[b=1]{-3}{3}
\end{psgraph*}
\\ \hline 
[b=0.1]] & [b=0.5] & [b=1] \\ 
\hline 
\multicolumn{3}{|c|}{ \dft : b = 1 }
\\ \hline
\end{tabular}


\bigskip

\begin{tabular}{|c|c|c|} \hline 
\multicolumn{3}{|c|}{ \BSS{psCauchyI}[m=0.1]\AC{-3}\AC{3} }
\\ \hline 
\begin{psgraph*}[axesstyle=none,xticksize= 0 1 ,yticksize=-3  3 , subticks=0, dy=.2,Dy=.2 ](0,0)(-3,0)(3,1){3cm}{3cm }
\psCauchyI[linestyle=dotted]{-3}{3}
\psCauchyI[m=-1]{-3}{3}
\end{psgraph*}
&  
\begin{psgraph*}[axesstyle=none,xticksize= 0 1 ,yticksize=-3  3 , subticks=0, dy=.2,Dy=.2 ](0,0)(-3,0)(3,1){3cm}{3cm }
\psCauchyI[linestyle=dotted]{-2.5}{2.5}
\psCauchyI[m=0]{-3}{3}
\end{psgraph*}
&  
\begin{psgraph*}[axesstyle=none,xticksize= 0 1 ,yticksize=-3  3 , subticks=0, dy=.2,Dy=.2 ](0,0)(-3,0)(3,1){3cm}{3cm }
\psCauchyI[linestyle=dotted]{-3}{3}
\psCauchyI[m=1]{-3}{3}
\end{psgraph*}
\\ \hline 
[m=-1] & [m=0] & [m=1] \\ 
\hline 
\multicolumn{3}{|c|}{ \dft : m = 0 }
\\ \hline
\end{tabular}

\newpage

%\subsection{Loi de Weibull}
\SbSSCT{Loi de Weibull}{Weibull Distribution}

\begin{center}
\begin{tabular}{|c|} \hline  
\begin{psgraph*}[axesstyle=none,xticksize= 0 1 ,yticksize=0 3 , subticks=0, dy=.2,Dy=.2 ](0,0)(0,0)(3,1){6cm}{2cm }
\psWeibull{0}{3}
\end{psgraph*}
\\ \hline  
\BSS{psWeibull}\AC{0}\AC{3} \BSI{psWeibull}{pst-func} 
\\ \hline 
\end{tabular}
\end{center}

\bigskip

\begin{tabular}{|c|c|} \hline  
\begin{psgraph*}[axesstyle=none,xticksize= 0 1 ,yticksize=0 3 , subticks=0, dy=.2,Dy=.2 ](0,0)(0,0)(3,1){6cm}{2cm }
\psWeibull[linestyle=dotted]{0}{3}
\psWeibull[alpha=.5]{0}{3}
\end{psgraph*}
& 
\begin{psgraph*}[axesstyle=none,xticksize= 0 1 ,yticksize=0 3 , subticks=0, dy=.2,Dy=.2 ](0,0)(0,0)(3,1){6cm}{2cm }
\psWeibull[linestyle=dotted]{0}{3}
\psWeibull[alpha=2]{0}{3}
\end{psgraph*}
 \\ \hline  
\BS{psWeibull}[\RDD{alpha}=.5]\AC{0}\AC{3} 
&  \BS{psWeibull}\RDD{alpha}=2]\AC{0}\AC{3} \\ 
\hline 
\multicolumn{2}{|c|}{ \dft : alpha=1 }
\\ \hline 
\end{tabular} 

\bigskip

\begin{tabular}{|c|c|} \hline  
\begin{psgraph*}[axesstyle=none,xticksize= 0 1 ,yticksize=0 3 , subticks=0, dy=.2,Dy=.2 ](0,0)(0,0)(3,1){6cm}{2cm }
\psWeibull[linestyle=dotted]{0}{3}
\psWeibull[beta=.5]{0}{3}
\end{psgraph*}
& 
\begin{psgraph*}[axesstyle=none,xticksize= 0 1 ,yticksize=0 3 , subticks=0, dy=.2,Dy=.2 ](0,0)(0,0)(3,1){6cm}{2cm }
\psWeibull[linestyle=dotted]{0}{3}
\psWeibull[beta=2]{0}{3}
\end{psgraph*}
 \\ \hline  
\BS{psWeibull}[\RDD{beta}=.5]\AC{0}\AC{3} 
&  \BS{psWeibull}\RDD{beta}=2]\AC{0}\AC{3} \\ 
\hline 
\multicolumn{2}{|c|}{ \dft :beta=1 }
\\ \hline
\end{tabular} 

\bigskip

\begin{tabular}{|c|c|} \hline  
\begin{psgraph*}[axesstyle=none,xticksize= 0 2 ,yticksize=0 3 , subticks=0, dy=.2,Dy=.2 ](0,0)(0,0)(3,2){6cm}{4cm }
\psWeibull[linestyle=dotted]{0}{3}
\psWeibull[alpha=2,beta=.5]{0}{3}
\end{psgraph*}
& 
\begin{psgraph*}[axesstyle=none,xticksize= 0 2 ,yticksize=0 3 , subticks=0, dy=.2,Dy=.2 ](0,0)(0,0)(3,2){6cm}{4cm }
\psWeibull[linestyle=dotted]{0}{3}
\psWeibull[alpha=2,beta=2]{0}{3}
\end{psgraph*}
 \\ \hline  
\BS{psWeibull}[alpha=2,beta=.5]\AC{0}\AC{3} 
&  \BS{psWeibull}[alpha=2,beta=2]\AC{0}\AC{3} \\ 
\hline 
\end{tabular} 


\newpage

\begin{center}
\begin{tabular}{|c|} \hline  
\begin{psgraph*}[axesstyle=none,xticksize= 0 1 ,yticksize=0 3 , subticks=0, dy=.2,Dy=.2 ](0,0)(0,0)(3,1){6cm}{2cm }
\psWeibullI{0}{3}
\end{psgraph*}
\\ \hline  
\BSS{psWeibullI}\AC{0}\AC{3}  \BSI{psWeibullI}{pst-func}
\\ \hline 
\end{tabular}
\end{center}

\bigskip


\begin{tabular}{|c|c|} \hline  
\begin{psgraph*}[axesstyle=none,xticksize= 0 1 ,yticksize=0 3 , subticks=0, dy=.2,Dy=.2 ](0,0)(0,0)(3,1){6cm}{2cm }
\psWeibullI[linestyle=dotted]{0}{3}
\psWeibullI[alpha=.5]{0}{3}
\end{psgraph*}
& 
\begin{psgraph*}[axesstyle=none,xticksize= 0 1 ,yticksize=0 3 , subticks=0, dy=.2,Dy=.2 ](0,0)(0,0)(3,1){6cm}{2cm }
\psWeibullI[linestyle=dotted]{0}{3}
\psWeibullI[alpha=2]{0}{3}
\end{psgraph*}
 \\ \hline  
\BS{psWeibullI}[\RDD{alpha}=.5]\AC{0}\AC{3}  \RDI{alpha}{pst-func}
&  \BS{psWeibullI}\RDD{alpha}=2]\AC{0}\AC{3} \\ 
\hline 
\multicolumn{2}{|c|}{ \dft : alpha=1 }
\\ \hline
\end{tabular} 


\bigskip


\begin{tabular}{|c|c|} \hline  
\begin{psgraph*}[axesstyle=none,xticksize= 0 1 ,yticksize=0 3 , subticks=0, dy=.2,Dy=.2 ](0,0)(0,0)(3,1){6cm}{2cm }
\psWeibullI[linestyle=dotted]{0}{3}
\psWeibullI[beta=.5]{0}{3}
\end{psgraph*}
& 
\begin{psgraph*}[axesstyle=none,xticksize= 0 1 ,yticksize=0 3 , subticks=0, dy=.2,Dy=.2 ](0,0)(0,0)(3,1){6cm}{2cm }
\psWeibullI[linestyle=dotted]{0}{3}
\psWeibullI[beta=2]{0}{3}
\end{psgraph*}
 \\ \hline  
\BS{psWeibullI}[\RDD{beta}=.5]\AC{0}\AC{3} \RDI{beta}{pst-func}
&  \BS{psWeibullI}\RDD{beta}=2]\AC{0}\AC{3} \\ 
\hline
\multicolumn{2}{|c|}{ \dft : beta=1 }
\\ \hline 
\end{tabular} 


\bigskip

\begin{tabular}{|c|c|} \hline  
\begin{psgraph*}[axesstyle=none,xticksize= 0 1 ,yticksize=0 3 , subticks=0, dy=.2,Dy=.2 ](0,0)(0,0)(3,1){6cm}{2cm }
\psWeibullI[linestyle=dotted]{0}{3}
\psWeibullI[alpha=2,beta=.5]{0}{3}
\end{psgraph*}
& 
\begin{psgraph*}[axesstyle=none,xticksize= 0 1 ,yticksize=0 3 , subticks=0, dy=.2,Dy=.2 ](0,0)(0,0)(3,1){6cm}{2cm }
\psWeibullI[linestyle=dotted]{0}{3}
\psWeibullI[alpha=2,beta=2]{0}{3}
\end{psgraph*}
 \\ \hline  
\BS{psWeibullI}[alpha=2,beta=.5]\AC{0}\AC{3} 
&  \BS{psWeibullI}[alpha=2,beta=2]\AC{0}\AC{3} \\ 
\hline 
\end{tabular}

\newpage
%\subsection{Loi de Vasicek}
\SbSSCT{Loi de Vasicek}{Vasicek Distribution}

\begin{center}
\begin{tabular}{|c|} \hline  
\begin{psgraph*}[axesstyle=none,xticksize= 0 5, dx=.2,Dx=.2 ,yticksize=0 1, subticks=0 ](0,0)(0,0)(1,5){6cm}{3cm }
 \psVasicek{0}{0.9999}
\end{psgraph*}
\\ \hline  
\BSS{psVasicek}\AC{0}\AC{3}  \BSI{psVasicek}{pst-func}
\\ \hline 
\end{tabular}
\end{center}

\bigskip


\begin{tabular}{|c|c|} \hline  
\begin{psgraph*}[axesstyle=none,xticksize= 0 10, dx=.2,Dx=.2 ,yticksize=0 1, subticks=0 ](0,0)(0,0)(1,10){6cm}{3cm }
\psVasicek[linestyle=dotted]{0}{0.9999}
\psVasicek[pd=.1]{0}{0.9999}
\end{psgraph*}
& 
\begin{psgraph*}[axesstyle=none,xticksize= 0 5, dx=.2,Dx=.2 ,yticksize=0 1, subticks=0 ](0,0)(0,0)(1,5){6cm}{3cm }
\psVasicek[linestyle=dotted]{0}{0.9999}
\psVasicek[pd=.5]{0}{0.9999}
\end{psgraph*}
 \\ \hline  
\BS{psVasicek}[\RDD{pd}=.1]\AC{0}\AC{3}  \RDI{pd}{pst-func}
&  \BS{psVasicek}[\RDD{pd}=.5]\AC{0}\AC{3} \\ 
\hline 
\multicolumn{2}{|c|}{ \dft : pd = 0.22 }
\\ \hline
\end{tabular} 

\bigskip


\begin{tabular}{|c|c|} \hline  
\begin{psgraph*}[axesstyle=none,xticksize= 0 10, dx=.2,Dx=.2 ,yticksize=0 1, subticks=0 ](0,0)(0,0)(1,10){6cm}{3cm }
\psVasicek[linestyle=dotted]{0}{0.9999}
\psVasicek[R2=.05]{0}{0.9999}
\end{psgraph*}
& 
\begin{psgraph*}[axesstyle=none,xticksize= 0 5, dx=.2,Dx=.2 ,yticksize=0 1, subticks=0 ](0,0)(0,0)(1,5){6cm}{3cm }
\psVasicek[linestyle=dotted]{0}{0.9999}
\psVasicek[R2=.2]{0}{0.9999}
\end{psgraph*}
 \\ \hline  
\BS{psVasicek}[\RDD{R2}=.05]\AC{0}\AC{3}  \RDI{R2}{pst-func}
&  \BS{psVasicek}[\RDD{R2}=.2]\AC{0}\AC{3} \\ 
\hline 
\multicolumn{2}{|c|}{ \dft : R2 = 0.11 }
\\ \hline
\end{tabular}


\bigskip


\begin{tabular}{|c|c|} \hline  
\begin{psgraph*}[axesstyle=none,xticksize= 0 5, dx=.2,Dx=.2 ,yticksize=0 1, subticks=0 ](0,0)(0,0)(1,5){6cm}{3cm }
\psVasicek[linestyle=dotted,pd=.5]{0}{0.9999}
\psVasicek[pd=.5,R2=.05]{0}{0.9999}
\end{psgraph*}
& 
\begin{psgraph*}[axesstyle=none,xticksize= 0 5, dx=.2,Dx=.2 ,yticksize=0 1, subticks=0 ](0,0)(0,0)(1,5){6cm}{3cm }
\psVasicek[linestyle=dotted,pd=.5]{0}{0.9999}
\psVasicek[pd=.5,R2=.2]{0}{0.9999}
\end{psgraph*}
 \\ \hline  
\BS{psVasicek}[pd=.5,R2=.05]\AC{0}\AC{3}  \RDI{R2}{pst-func}
&  \BS{psVasicek}[pd=.5,R2=.2]\AC{0}\AC{3} \\ 
\hline 
\end{tabular}

\newpage
%\subsection{Courbe de Lorenz}
\SbSSCT{Courbe de Lorenz}{Lorenz curve}

\begin{tabular}{|c|c|} \hline  
\begin{psgraph*}[axesstyle=none,xticksize= 0 1,dx=.2,Dx=.2 , dy=.2,Dy=.2 ,yticksize=0 1, subticks=0 ](0,0)(0,0)(1,1){5cm}{3cm }
 \psLorenz{0.1 0.2 0.3 }
\end{psgraph*}
&
\begin{psgraph*}[axesstyle=none,xticksize= 0 1,dx=.2,Dx=.2 , dy=.2,Dy=.2 ,yticksize=0 1, subticks=0 ](0,0)(0,0)(1,1){5cm}{3cm }
 \psLorenz*{.1 .2 .3 }
\end{psgraph*}
\\ \hline  
\BSS{psLorenz}\AC{0.1 0.2 0.3} \BSI{psLorenz}{pst-func}
&
\BSS{psLorenz}*\AC{0.1 0.2 0.3 }
\\ \hline 
\end{tabular}


\bigskip

\begin{tabular}{|c|c|} \hline  
\begin{psgraph*}[axesstyle=none,xticksize= 0 1, dx=.2,Dx=.2 ,yticksize=0 1, subticks=0,lly=-13mm ](0,0)(0,0)(1,1){5cm}{3cm }
 \psLorenz[linestyle=dotted]{.1 .2 .3 }
 \psLorenz[plotstyle=bezier]{.1 .2 .3 }
\end{psgraph*}
&  
\begin{psgraph*}[axesstyle=none,xticksize= 0 1, dx=.2,Dx=.2 ,yticksize=0 1, subticks=0,lly=-13mm ](0,0)(0,0)(1,1){5cm}{3cm }
 \psLorenz[linestyle=dotted]{.1 .2 .3 }
 \psLorenz[Gini]{.1 .2 .3 }
\end{psgraph*}
\\ \hline  
\BS{psLorenz}[plotstyle=bezier]\AC{.1 .2 .3}  
&
\BS{psLorenz}[\RDD{Gini}]\AC{.1 .2 .3 }
 \\ \hline 
\end{tabular} 

\newpage
%\subsection{Courbe de Lamé : superellipses}
\SbSSCT{Courbe de Lamé : superellipses}{Lame curve}


 \begin{tabular}{|c|c|c|c|}  \hline  
\begin{psgraph*}[axesstyle=none,xticksize= -1 1 , yticksize=-1 1, subticks=0 ](0,0)(-1,-1)(1,1){2cm}{2cm}
\psLame{.5}
\end{psgraph*} 
 &  
\begin{psgraph*}[axesstyle=none,xticksize= -1 1 , yticksize=-1 1, subticks=0 ](0,0)(-1,-1)(1,1){2cm}{2cm}
\psLame{.75}
\end{psgraph*} 
 &  
\begin{psgraph*}[axesstyle=none,xticksize= -1 1 , yticksize=-1 1, subticks=0 ](0,0)(-1,-1)(1,1){2cm}{2cm}
\psLame{2}
\end{psgraph*} 
 &  
\begin{psgraph*}[axesstyle=none,xticksize= -1 1 , yticksize=-1 1, subticks=0 ](0,0)(-1,-1)(1,1){2cm}{2cm}
\psLame{5}
\end{psgraph*} 
 \\  \hline 
 \BSS{psLame}\AC{.5} \BSI{psLame}{pst-func} &  \BSS{psLame}\AC{.75} &  \BSS{psLame}\AC{2} &   \BSS{psLame}\AC{5}\\ 
 \hline 
 \end{tabular} 

\bigskip

\begin{tabular}{|c|c|} \hline  
\begin{psgraph*}[axesstyle=none,xticksize= -2 2 ,yticksize=-2 2,subticks=0 ](0,0)(-2,-2)(2,2){4cm}{4cm }
\psLame[radiusA=2]{.5}
\end{psgraph*}
&  
\begin{psgraph*}[axesstyle=none,xticksize= -2 2 ,yticksize=-2 2,subticks=0 ](0,0)(-2,-2)(2,2){4cm}{4cm }
\psLame[radiusB=2]{.5}
\end{psgraph*}
\\ \hline 
\BS{psLame}[\RDD{radiusA}=2]\AC{.5} \RDI{radiusA}{pst-func}&
\BS{psLame}[\RDD{radiusB}=2]\AC{.5}  \RDI{radiusB}{pst-func}\\ 
\hline 
\end{tabular} 


%\subsection{Fonction de Thomae}
\SbSSCT{Fonction de Thomae}{Thomae curve}

 \psset{unit=4cm}


\begin{tabular}{|c|c|c|c|} \hline  
\begin{psgraph*}[axesstyle=none,xticksize= 0 1, dx=.2,Dx=.2, dy=.2,Dy=.2  ,yticksize=0 1, subticks=0,lly=-13mm ](0,0)(0,0)(1,1){2.5cm}{2.5cm }
 \psThomae[dotsize=5pt](0,1){1}
\end{psgraph*}
&  
\begin{psgraph*}[axesstyle=none,xticksize= 0 1, dx=.2,Dx=.2, dy=.2,Dy=.2  ,yticksize=0 1, subticks=0,lly=-13mm ](0,0)(0,0)(1,1){2.5cm}{2.5cm }
 \psThomae[dotsize=5pt](0,1){2}
\end{psgraph*}
&  
\begin{psgraph*}[axesstyle=none,xticksize= 0 1, dx=.2,Dx=.2, dy=.2,Dy=.2  ,yticksize=0 1, subticks=0,lly=-13mm ](0,0)(0,0)(1,1){2.5cm}{2.5cm }
 \psThomae[dotsize=5pt](0,1){3}
\end{psgraph*}
&  
\begin{psgraph*}[axesstyle=none,xticksize= 0 1, dx=.2,Dx=.2, dy=.2,Dy=.2  ,yticksize=0 1, subticks=0,lly=-13mm ](0,0)(0,0)(1,1){2.5cm}{2.5cm }
 \psThomae[dotsize=5pt](0,1){10}
\end{psgraph*}
\\ 
\hline 
\BSS{psThomae}(0,1)\AC{1}  \BSI{psThomae}{pst-func} & \BSS{psThomae}(0,1)\AC{2}  & \BSS{psThomae}(0,1)\AC{3} & \BSS{psThomae}(0,1)\AC{10}\\ 
\hline 
\end{tabular} 

\bigskip

\begin{tabular}{|c|c|} \hline  
\begin{psgraph*}[axesstyle=none,xticksize= 0 1, dx=.5,Dx=.5, dy=.5,Dy=.5  ,yticksize=0 2, subticks=0 ](0,0)(0,0)(2,1){5cm}{3cm }
 \psThomae[dotsize=5pt](0,2){10}
\end{psgraph*}
&  
\begin{psgraph*}[axesstyle=none,xticksize= 0 1, dx=.5,Dx=.5, dy=.5,Dy=.5  ,yticksize=0 2.5, subticks=0 ](0,0)(0,0)(2.5,1){7cm}{3cm }
 \psThomae[dotsize=5pt](0.5,2.5){10}
\end{psgraph*}
\\ \hline 
\BSS{psThomae}(0,2){10}(0,2)\AC{10} & \BSS{psThomae}(0,2){10}(0.5,2)\AC{10} \\ 
\hline 
\end{tabular}


%\subsection{Fonction de Weierstrass}
\SbSSCT{Fonction de Weierstrass}{Weierstrass curve}

\begin{tabular}{|c|c|} \hline  
\begin{psgraph*}[axesstyle=none,xticksize= -.5  .5, dx=.5,Dx=.5, dy=.5,Dy=.5  ,yticksize=0 2, subticks=0 ](0,0)(0,-.5)(2,.5){5cm}{3cm }
 \psWeierstrass(0,2){2}
\end{psgraph*}
&  
\begin{psgraph*}[axesstyle=none,xticksize= -.5  .5, dx=.5,Dx=.5, dy=.5,Dy=.5  ,yticksize=0 2, subticks=0 ](0,0)(0,-.5)(2,.5){5cm}{3cm }
% \psWeierstrass[linestyle=dotted](0,2){2}
 \psWeierstrass(0,2){5}
\end{psgraph*}
\\ \hline  
\BSS{psWeierstrass}(0,2)\AC{2} & \BSS{psWeierstrass}(0,2)\AC{5}  \BSI{psWeierstrass}{pst-func}\\ 
\hline 
\end{tabular} 

\bigskip

 \begin{tabular}{|c|c|} \hline  
 \begin{psgraph*}[axesstyle=none,xticksize= 0  .5, dx=.1,Dx=.1, dy=.1,Dy=.1  ,yticksize=.5 1, subticks=0 ](0,0)(0.5,0)(1,.5){5cm}{3cm }
  \psWeierstrass(.5,1){2}
 \end{psgraph*}
 &  
 \begin{psgraph*}[axesstyle=none,xticksize= 0  .05, dx=.01,Dx=.01, dy=.01,Dy=.01  ,yticksize=.95 1, subticks=0 ](0,0)(0.95,0)(1,.05){5cm}{3cm }
  \psWeierstrass(.95,1){2}
 \end{psgraph*}

 \\ \hline  
 \BSS{psWeierstrass}(.5,1)\AC{2} & \BSS{psWeierstrass}(.95,1)\AC{5}  \BSI{psWeierstrass}{pst-func}\\ 
 \hline 
 \end{tabular} 
 
 \bigskip
 
\begin{tabular}{|c|c|} \hline  
\begin{psgraph*}[axesstyle=none,xticksize= -.5  .5, dx=.5,Dx=.5, dy=.5,Dy=.5  ,yticksize=0 2, subticks=0 ](0,0)(0,-.5)(2,.5){5cm}{3cm }
 \psWeierstrass[linestyle=dotted](0,2){2}
 \psWeierstrass[epsilon=1.e-1](0,5){2}
\end{psgraph*}
&  
\begin{psgraph*}[axesstyle=none,xticksize= -.5  .5, dx=.5,Dx=.5, dy=.5,Dy=.5  ,yticksize=0 2, subticks=0 ](0,0)(0,-.5)(2,.5){5cm}{3cm }
 \psWeierstrass[linestyle=dotted](0,2){2}
 \psWeierstrass[epsilon=1](0,5){2}
\end{psgraph*}
\\ \hline  
 \BSS{psWeierstrass}[\RDD{epsilon}=1.e-1](0,5)\AC{2} &   \BSS{psWeierstrass}[\RDD{epsilon}=1](0,5)\AC{2}
 \RDI{epsilon}{pst-func} \\ \hline 
\multicolumn{2}{|c|}{ \dft : epsilon=1.e-18  } 
\\ \hline 
 
\end{tabular} 
 
\newpage
%\subsection{Fonction définie implicitement}
\SbSSCT{Fonction définie implicitement}{implicit defined functions}

\begin{tabular}{|c|} \hline  
\begin{psgraph*}[axesstyle=none,xticksize= -6  3 ,yticksize=-5 5, subticks=0 ](0,0)(-5,-6)(5,3){10cm}{5cm }
 \psplotImp[linewidth=2pt](-6,-7)(4,3){4 x 3 exp y 3 exp add 4 x y mul mul sub }
\end{psgraph*}
\\ \hline  
 \BSS{psplotImp}(-6,-7)(4,3)\AC{4 x 3 exp y 3 exp add 4 x y mul mul sub }
\\ \hline 
\end{tabular}

\bigskip

\begin{tabular}{|c|} \hline   
\begin{psgraph*}[axesstyle=none,xticksize= -6  3 ,yticksize=-5 5, subticks=0 ](0,0)(-5,-6)(5,3){10cm}{5cm }
 \psplotImp[algebraic,linewidth=2pt](-6,-7)(4,3){x^3 +y^3 -4*x*y}
\end{psgraph*}
\\ \hline  
 \BSS{psplotImp}[\RDD{algebraic}](-6,-7)(4,3)\AC{x\^{}3 +y\^{}3 -4*x*y }
\\ \hline 
\end{tabular}

\bigskip
\begin{tabular}{|c|} \hline   
\begin{psgraph*}[axesstyle=none,xticksize= -6  3 ,yticksize=-5 5, subticks=0 ](0,0)(-5,-6)(5,3){10cm}{5cm }
 \psplotImp[algebraic,linewidth=2pt,stepFactor=2](-6,-7)(4,3){x^3 +y^3 -4*x*y}
\end{psgraph*}
\\ \hline  
 \BSS{psplotImp}[algebraic,\RDD{stepFactor}=2](-6,-7)(4,3)\AC{x\^{}3 +y\^{}3 -4*x*y }
\\ \hline 
\end{tabular}

\bigskip
\begin{tabular}{|c|} \hline   
\begin{psgraph*}[axesstyle=none,xticksize= -1  1 ,yticksize=-1 1, subticks=0 ](0,0)(-1,-1)(1,1){5cm}{5cm }
 \psplotImp[algebraic,polarplot,linewidth=1pt](-1,-1)(1,1){ r + cos(10*phi) }
\end{psgraph*}
\\ \hline  
 \BSS{psplotImp}[algebraic,\RDD{polarplot}](-1,-1)(1,1)\AC{r + cos(10*phi) }
\\ \hline 
\end{tabular}

\bigskip

\begin{tabular}{|c|} \hline   
\begin{psgraph*}[axesstyle=none,xticksize= -1  1 ,yticksize=-1 1, subticks=0 ](0,0)(-1,-1)(1,1){5cm}{5cm }
 \psplotImp[algebraic,polarplot,linewidth=1pt,stepFactor=1](-1,-1)(1,1){ r + cos(10*phi) }
\end{psgraph*}
\\ \hline  
 \BSS{psplotImp}[algebraic,polarplot,\RDD{stepFactor}=1](-1,-1)(1,1)\AC{r + cos(10*phi) }
\\ \hline 
\end{tabular}

\newpage

%\subsection{Fonction de rotation}
\SbSSCT{Fonction de rotation}{Rotating functions}

 
\begin{tabular}{|c|}\hline  
\begin{psgraph*}[axesstyle=none,xticksize= -2  2 ,yticksize=0 5, subticks=0 ](0,0)(0,-2)(5,2){10cm}{5cm }
 \psVolume[fillstyle=solid,fillcolor=blue!40](0,4){4}{x sqrt}
\end{psgraph*}

\\ \hline  
 \BSS{psVolume}[fillstyle=solid,fillcolor=blue!40](0,4)\AC{4}\AC{x sqrt}  \BSI{psVolume}{pst-func}
\\ \hline 
\end{tabular} 

\bigskip

\begin{tabular}{|c|}\hline  
\begin{psgraph*}[axesstyle=none,xticksize= -1  1 ,yticksize=0 7, subticks=0 ](0,0)(0,-1)(7,1){10cm}{5cm }
 \psVolume[fillstyle=solid,fillcolor=yellow,algebraic](0,6.28){20}{cos(x)}
\end{psgraph*}
\\ \hline  
 \BSS{psVolume}[fillstyle=solid,fillcolor=yellow,algebraic](0,6.28)\AC{20}\AC{cos(x)}
\\ \hline 
\end{tabular} 

\newpage

% \section[Créer un graphe en camembert]{Créer un graphe en camembert \cite{pstricks-add} }
\SSCT{Créer un graphe en camembert}{Pie chart} 

 \TFRGB{
 Syntaxe :  
 \BS{psChart}[options]\AC{liste de valeurs }\AC{liste des valeurs décalées }\AC{rayon}}
 {Syntax :  
 \BS{psChart}[options]{value list }{list shifted values }{radius}}
 
 
 
\psset{fillstyle=solid,unit=1cm}
 
\begin{center}
\begin{tabular}{|c|c|}  \hline 
\begin{pspicture}(-2.5,-1.5)(2.5,1.5)
\psChart{1,2,3,4,5}{}{1cm}
\end{pspicture} 
&
\begin{pspicture}(-2.5,-1.5)(2.5,1.5)
\psChart{1,2,3,4,5}{2,5}{1cm}
\end{pspicture}
\\ \hline
\BSS{psChart}\AC{1,2,3,4,5}\AC{}\AC{1cm} \BSI{psChart}{pstricks-add} &
\textbf{\BS{}psChart}\AC{1,2,3,4,5}{\red \AC{2,5}}\AC{1cm}
\\ \hline
\end{tabular}
\end{center}

% \subsubsection{Paramètres}
\begin{center}
\begin{tabular}{|c|c|c|}  \hline 
\begin{pspicture}(-2,-1.5)(2,1.5)
\psChart[chartColor=color]{1,1,1,1,1,1,1,1,1,1,1 }{1}{1cm}
\end{pspicture}
&
\begin{pspicture}(-2.5,-1.5)(2.5,1.5)
\psChart[userColor={orange,teal,red!20}]{1,2,3}{1}{1cm}
\end{pspicture}
&
\begin{pspicture}(-2,-1.5)(2,1.5)
\psChart[chartColor=color,shadow=true]{1,2,3}{1}{1cm}
\end{pspicture}

\\ \hline
\RDD{chartColor}=color \RDI{chartColor}{pstricks-add} & \RDD{userColor}=\AC{orange,teal,red!20}  \RDI{userColor}{pstricks-add} & 
\RDD{shadow}=true  \RDI{shadow}{pstricks-add}
\\ \hline
\end{tabular}
\end{center}

\begin{center}
\begin{tabular}{|c|c|c|}  \hline 
\begin{pspicture}(-2,-2)(2,2)
\psChart[chartColor=color,chartSep=5pt]{1,2,3}{1}{1cm}
\end{pspicture}
&
\begin{pspicture}(-2,-2)(2,2)
\psChart[chartColor=color,chartSep=5pt]{1,2,3}{1,2}{1cm}
\end{pspicture}
&
\begin{pspicture}(-2,-2)(2,2)
\psChart[chartColor=color,chartSep=5pt]{1,2,3}{1,2,3}{1cm}
\end{pspicture}
\\ \hline
\multicolumn{3}{|c|}{\RDD{chartSep}=5pt  par défaut : 10pt  \RDI{chartSep}{pstricks-add} } \\ \hline
%chartColor=color & userColor={blue,green,red,cyan,orange} 
%\\ \hline
%{\red chartSep}=.5cm & 
%\\ \hline
\end{tabular}
\end{center}
%------------------------------------------------------------------
\SbSSCT{Etiquettes}{Labels}
\begin{center}
\begin{tabular}{|c|c|c|}  \hline 
\multicolumn{3}{|c|}{\TFRGB{3 positions possibles}{Three possible positions}} \\ \hline
\begin{pspicture}(-2,-1.5)(2,1.5)
\psChart[chartColor=color]{1,2}{}{1cm}
\psdots[dotstyle=*,dotscale=2,linecolor=green](psChartO1)
\psdots[dotstyle=*,dotscale=2,linecolor=green](psChartO2)
\end{pspicture} 
&
\begin{pspicture}(-2,-1.5)(2,1.5)
\psChart[chartColor=color]{1,2}{}{1cm}
\psdots[dotstyle=*,dotscale=2,linecolor=green](psChartI1)
\psdots[dotstyle=*,dotscale=2,linecolor=green](psChartI2)
\end{pspicture}
&
\begin{pspicture}(-2,-1.5)(2,1.5)
\psChart[chartColor=color]{1,2}{}{1cm}
\psdots[dotstyle=*,dotscale=2,linecolor=green](psChart1)
\psdots[dotstyle=*,dotscale=2,linecolor=green](psChart2)
\end{pspicture}
\\ \hline
\BS{}psdots(\BDD{psChartO1})  & \BS{}psdots(\BDD{psChartI1})  & \BS{}psdots(\BDD{psChart1}) \\ 
\BS{}psdots(\BDD{psChartO2})  & \BS{}psdots(\BDD{psChartI2})  & \BS{}psdots(\BDD{psChart2}) \\ \hline
\end{tabular}
\end{center}



%---------------------------------------------------------

\begin{center}
\begin{tabular}{|c|c|}  \hline 
\multicolumn{2}{|c|}{\TFRGB{Liaison des points}{Connection of the points}} \\ \hline
\begin{pspicture}(-2,-2)(2,2)
\psChart[chartColor=color,chartNodeO=2]{1,2}{}{1cm}
\pcline[linecolor=green](psChartO1)(psChartI1)
\pcline[linecolor=green](psChartO2)(psChart2)
\end{pspicture} 
&
\begin{pspicture}(-3,-1.5)(3,1.5)
\psChart[chartColor=color,chartNodeI=.5]{1,2}{}{1cm}
\ncline[linecolor=green]{psChartO1}{psChartI1}
\ncline[linecolor=green]{psChartO2}{psChart2}
\end{pspicture}
\\ \hline
\BS{}pcline(psChartO1)(psChartI1)
 &
\BS{}ncline\AC{psChartO1}\AC{psChartI1}
\\ 
\BS{}pcline(psChartO2)(psChart2)
 &
\BS{}ncline\AC{psChartO2}\AC{psChart2} \\ 
\hline
\end{tabular}
\end{center}

%----------------------------------------------------------
\begin{center}
\begin{tabular}{|c|c|}  \hline 
\multicolumn{2}{|c|}{\TFRGB{Attachement des étiquettes aux points}{labels attachment points}} \\ \hline
\begin{pspicture}(-3,-2.5)(3,2.5)
\psChart[chartColor=color]{1,2}{}{1cm}
\rput*[l](psChartO1){psChartO1}
\rput*[l](psChartO2){O2}
\end{pspicture} 
&
\begin{pspicture}(-3,-2.5)(3,2.5)
\psChart[chartColor=color]{1,2}{}{1cm}
\rput*[l](psChartI1){psChartI1}
\rput*[l](psChartI2){I2}
\end{pspicture}
\\ \hline
\BS{}rput*[l](\BDD{psChartO1)}\AC{psChartO1} &
\BS{}rput*[l](\BDD{psChartI1})\AC{psChartI1} 
\\ 
\BS{}rput*[l](\BDD{psChartO2})\AC{O2} &
\BS{}rput*[l](\BDD{psChartI2})\AC{I2} 
\\ 
\hline
\end{tabular}
\end{center}
\newpage
\begin{center}
\begin{tabular}{|c|c|}  \hline 
\begin{pspicture}(-2,-2.5)(2,2.5)
\psChart[chartColor=color,chartNodeO=2]{1,2}{}{1cm}
\psdots[dotstyle=*,dotscale=2,linecolor=green](psChartO1)
\psdots[dotstyle=*,dotscale=2,linecolor=green](psChartO2)
\end{pspicture} 
&
\begin{pspicture}(-2,-2.5)(2,2.5)
\psChart[chartColor=color,chartNodeI=.5]{1,2}{}{1cm}
\psdots[dotstyle=*,dotscale=2,linecolor=green](psChartI1)
\psdots[dotstyle=*,dotscale=2,linecolor=green](psChartI2)
\end{pspicture}
\\ \hline
\RDD{chartNodeO}=2 \RDI{chartNodeO}{pstricks-add} & \RDD{chartNodeI}=.5 \RDI{chartNodeI}{pstricks-add}\\ 
par défaut : 1.5 & par défaut : .75 \\ 
\hline
\end{tabular}
\end{center}
%--------------------------------------------------------
\begin{center}
\begin{tabular}{|c|c|c|c|c|c|c|c|c|c|}  \hline
\multicolumn{10}{|c|}{\TFRGB{Couleurs de remplissage}{Filling color }}\\
\multicolumn{10}{|c|}{\BS{}psframe[fillcolor=\BDD{chartFillColor1}(1,1) à \BS{}psframe[fillcolor=\BDD{chartFillColor10}(1,1)}\\ \hline
\begin{pspicture}(1,1.2)
\psframe[fillcolor=chartFillColor1](1,1)
\end{pspicture}
&
\begin{pspicture}(1,1.2)
\psframe[fillcolor=chartFillColor2](1,1)
\end{pspicture}
&
\begin{pspicture}(1,1.2)
\psframe[fillcolor=chartFillColor3](1,1)
\end{pspicture}
&
\begin{pspicture}(1,1.2)
\psframe[fillcolor=chartFillColor4](1,1)
\end{pspicture}
&
\begin{pspicture}(1,1.2)
\psframe[fillcolor=chartFillColor5](1,1)
\end{pspicture}
&
\begin{pspicture}(1,1.2)
\psframe[fillcolor=chartFillColor6](1,1)
\end{pspicture}
&
\begin{pspicture}(1,1.2)
\psframe[fillcolor=chartFillColor7](1,1)
\end{pspicture}
&
\begin{pspicture}(1,1.2)
\psframe[fillcolor=chartFillColor8](1,1)
\end{pspicture}
&
\begin{pspicture}(1,1.2)
\psframe[fillcolor=chartFillColor9](1,1)
\end{pspicture}
&
\begin{pspicture}(1,1.2)
\psframe[fillcolor=chartFillColor10](1,1)
\end{pspicture}
\\ \hline

\end{tabular}
\end{center}
%------------------------------------------

\begin{center}
\begin{tabular}{|c|c|c|c|c|c|c|c|c|c|}  \hline
\multicolumn{10}{|c|}{\TFRGB{Couleurs complémentaires}{More colors}}\\
\multicolumn{10}{|c|}{\BS{}psframe[fillcolor=\BDD{-chartFillColor1}(1,1) à \BS{}psframe[fillcolor=\BDD{-chartFillColor10}(1,1)}\\ \hline
\begin{pspicture}(1,1.2)
\psframe[fillcolor=-chartFillColor1](1,1)
\end{pspicture}
&
\begin{pspicture}(1,1.2)
\psframe[fillcolor=-chartFillColor2](1,1)
\end{pspicture}
&
\begin{pspicture}(1,1.2)
\psframe[fillcolor=-chartFillColor3](1,1)
\end{pspicture}
&
\begin{pspicture}(1,1.2)
\psframe[fillcolor=-chartFillColor4](1,1)
\end{pspicture}
&
\begin{pspicture}(1,1.2)
\psframe[fillcolor=-chartFillColor5](1,1)
\end{pspicture}
&
\begin{pspicture}(1,1.2)
\psframe[fillcolor=-chartFillColor6](1,1)
\end{pspicture}
&
\begin{pspicture}(1,1.2)
\psframe[fillcolor=-chartFillColor7](1,1)
\end{pspicture}
&
\begin{pspicture}(1,1.2)
\psframe[fillcolor=-chartFillColor8](1,1)
\end{pspicture}
&
\begin{pspicture}(1,1.2)
\psframe[fillcolor=-chartFillColor9](1,1)
\end{pspicture}
&
\begin{pspicture}(1,1.2)
\psframe[fillcolor=-chartFillColor10](1,1)
\end{pspicture}
\\ \hline

\end{tabular}
\end{center}

%\begin{pspicture}(12,1)
%\psframe[fillcolor=chartFillColor1](1,1)
%\psframe[fillcolor=chartFillColor2](1,0)(2,1)
%\psframe[fillcolor=chartFillColor3](2,0)(3,1)
%\psframe[fillcolor=chartFillColor4](3,0)(4,1)
%\psframe[fillcolor=chartFillColor5](4,0)(5,1)
%\psframe[fillcolor=chartFillColor6](5,0)(6,1)
%\psframe[fillcolor=chartFillColor7](6,0)(7,1)
%\psframe[fillcolor=chartFillColor8](7,0)(8,1)
%\psframe[fillcolor=chartFillColor9](8,0)(9,1)
%\psframe[fillcolor=chartFillColor10](9,0)(10,1)
%
%
%\end{pspicture} 

%\begin{pspicture}(12,1)
%\psframe[fillcolor=-chartFillColor1](1,1)
%\psframe[fillcolor=-chartFillColor2](1,0)(2,1)
%\psframe[fillcolor=-chartFillColor3](2,0)(3,1)
%\psframe[fillcolor=-chartFillColor4](3,0)(4,1)
%\psframe[fillcolor=-chartFillColor5](4,0)(5,1)
%\psframe[fillcolor=-chartFillColor6](5,0)(6,1)
%\psframe[fillcolor=-chartFillColor7](6,0)(7,1)
%\psframe[fillcolor=-chartFillColor8](7,0)(8,1)
%\psframe[fillcolor=-chartFillColor9](8,0)(9,1)
%\psframe[fillcolor=-chartFillColor10](9,0)(10,1)
%\end{pspicture} 
%
%\begin{center}
%\begin{tabular}{|c|c|}  \hline 
%\begin{pspicture}(-3,-3)(3,3)
%\psChart[chartColor=color,chartNodeO=1.2]{1,2,3,4,5 }{3}{2cm}
%\rput[l](psChartO1){O1}
%\rput[l](psChartO2){\textcolor{chartFillColor2}{O2}}
%\rput[l](psChartO3){\textcolor{-chartFillColor3}{O3}}
%\rput[l](psChartO4){\textcolor{-chartFillColor4}{O4}}
%\rput[l](psChartO5){\textcolor{chartFillColor5}{O5}}
%\end{pspicture} 
%&
%\begin{pspicture}(-3,-3)(3,3)
%\psChart[chartColor=color,shadow=true,shadowsize=5pt,chartNodeI=.5]{1,2,3,4,5 }{}{2cm}
%\rput*[l](psChartI1){I1}
%\rput[l](psChartI2){I2}
%\rput*[c](psChartI3){I3}
%\rput*[c](psChart4){4}
%\rput*[l](psChart5){5}
%%\psdots[dotstyle=*,dotscale=2](psChartI4)
%  \end{pspicture} 
%  \\ \hline
%\BS{}rput[l]\AC(psChartO1){O1} &  \BS{}rput*[l](psChartI1)\AC{I1} \\  
%\BS{}rput[l](psChartO2){\BS{}textcolor\AC{chartFillColor2}\AC{O2}} & 
%\BS{}rput[l](psChartI2)\AC{I2} \\
%
%\hline
%  \end{tabular}
%  \end{center} % camembert


\newpage
%%%=============================================

%\section{Les répétitions}
\SSCT{Les répétitions}{Repetitions}

\subsection[Multirput]{Multirput \cite{pst-user} }


\psset{fillcolor=yellow,fillstyle=solid,linecolor=red,unit=1cm}




\begin{tabular}{|c|c|}\hline 
\begin{pspicture}(-.5,-.5)(6,3)
%\psframe[fillstyle=none](-.5,-.5)(6,3) 
 \multirput(.5,0)(.5,.25){10}{A}
 \end{pspicture}
& 
\begin{pspicture}(-.5,-.5)(6,3)
%\psframe[fillstyle=none](-.5,-.5)(6,3) 
 \multirput*(.5,0)(.5,.25){10}{A}
 \end{pspicture}
\\ \hline  
\BSS{multirput}( \rnode{A}{0.5,0})( \rnode{B}{0.5,0.25})\AC{ \rnode{C}{10}}\AC{A}
&  
\textbf{\BS{multirput}}{\red *}(0.5,0)(0.5,0.25)\AC{10}\AC{A}
\\ \hline 
& \\
 \rnode{AA}{\TFRGB{origine}{origin}} \hspace{1cm}  \rnode{BB}{\TFRGB{décalage}{shift}} \hspace{1cm}\rnode{CC}{10 \TFRGB{fois}{times}}
 & 
 \\ \hline 
\end{tabular}
 \ncline[linecolor=blue]{A}{AA}  \ncline[linecolor=blue]{B}{BB}  \ncline[linecolor=blue]{C}{CC}
\bigskip


\begin{tabular}{|c|c|}\hline 
\begin{pspicture}(-.5,-.5)(6,3)
%\psframe[fillstyle=none](-.5,-.5)(6,3) 
 \multirput{45}(.5,0)(.5,.25){10}{A}
 \end{pspicture}
& 
\begin{pspicture}(-.5,-.5)(6,3)
%\psframe[fillstyle=none](-.5,-.5)(6,3) 
 \multirput*{45}(.5,0)(.5,.25){10}{A}
 \end{pspicture}
\\ \hline  
\BSS{multirput}\AC{{\red 45}}(0.5,0)(0.5,0.25)\AC{10}\AC{A}
&  
\textbf{\BS{multirput}}*\AC{{\red 45}}(0.5,0)(0.5,0.25)\AC{10}\AC{A}
\\ \hline 
\end{tabular} 


\subsection[multips]{multips \cite{pst-user} }


\begin{tabular}{|c|c|}
\hline  
\begin{pspicture}(-.5,-.5)(6,3)
%\psframe[fillstyle=none](-.5,-.5)(6,3)
\multips(.5,0)(.5,.25){10}{\psframe(1,.5) }
\end{pspicture}
&
\begin{pspicture}(-.5,-.5)(6,3)
%\psframe[fillstyle=none](-.5,-.5)(6,3)  
\multips(.5,0)(.5;15){10}{\psframe(1,.5) }
\end{pspicture}
\\ \hline 
\BS{multips}(\rnode{A}{0.5,0})(\rnode{B}{0.5,0.25})\AC{\rnode{C}{10}}\AC{\BS{psframe}(1,.5) }
 &  
\BS{multips}(\rnode{D}{0.5;0})(\rnode{E}{0.5;15})\AC{10}\AC{\BS{psframe}(1,.5) }  
 \\ \hline 
& \\
 \rnode{AA}{\TFRGB{origine}{origin}} \hspace{1cm}  \rnode{BB}{\TFRGB{décalage}{shift}} \hspace{1cm}\rnode{CC}{10 \TFRGB{fois}{times}}
 &  \rnode{DD}{\TFRGB{coordonnées polaires}{polar coordinates}}
 \\ \hline 
\end{tabular}  
\ncline[linecolor=blue]{A}{AA}  \ncline[linecolor=blue]{B}{BB}  \ncline[linecolor=blue]{C}{CC} 
 \ncline[linecolor=blue]{D}{DD}  \ncline[linecolor=blue]{E}{DD} 

\bigskip
\begin{tabular}{|c|c|}
\hline  
\begin{pspicture}(-.5,-.5)(6,4)
%\psframe[fillstyle=none](-.5,-.5)(6,4)
\multips{45}(.5,0)(.5,.25){10}{\psframe(1,.5) }
\end{pspicture}
&
\begin{pspicture}(-.5,-.5)(6,4)
%\psframe[fillstyle=none](-.5,-.5)(6,4)  
\multips{45}(.5,0)(.5;15){10}{\psframe(1,.5)}
\end{pspicture}
\\ \hline 
\BS{multips}\AC{{\red 45}}(.5,0)(.5,.25)\AC{10}\AC{\BS{psframe}(1,.5)}
 &  
\BS{multips}\AC{{\red 45}}(.5;0)(.5;15)\AC{10}\AC{\BS{psframe}(1,.5)}  
 \\ \hline 
\end{tabular}  
%======================================================================

\subsection[rmultiput]{rmultiput \cite{pstricks-add} }

%Répétitions à pas variables

\psset{fillcolor=yellow,fillstyle=solid,linecolor=red}


%\TFRGB{syntaxe}{syntax} :

% \BSS{rmultiput}*[Options] \AC{objet à répéter }(1° position)(2° position). . . (n position )


%
%\vspace{1cm}
%Exemple :

\begin{tabular}{|c|c|}
\hline  
\begin{pspicture}(-1,-2)(7,2)
%\psframe[fillstyle=none](-1,-2)(7,2)
\rmultiput{\DFR }(0,0)(2,1)(4,0)(6,-1)
\end{pspicture}
&  
\begin{pspicture}(-1,-2)(7,2)
%\psframe[fillstyle=none](-1,-2)(7,2)
\rmultiput*[rot=45] {\DFR }(0,0)(2,1)(4,0)(6,-1)
\end{pspicture}
\\ 
\hline 
\BSS{rmultiput}\AC{\BS{DFR}}(0,0)(2,1)(4,0)(6,-1) \BSI{rmultiput}{pstricks-add} & 
\BS{rmultiput}*[{\red rot=45}]\AC{\BS{DFR}}(0,0)(2,1)(4,0)(6,-1) \\ 
\hline 
\end{tabular} 



%
%\essais{pstrep1ess.tex}

\subsection[Multido]{Multido \cite{pst-user} \cite{multido} }


%\subsection{$\backslash$multirput}
Utilisation du module \textbf{multido} \label{multido}

\psset{fillcolor=yellow,fillstyle=solid,linecolor=blue,unit=1cm}


%Syntaxe : 

%\begin{tabular}{ll}
% \BSS{multido} 				&   \\ 
%\AC{variables} 						& \% variable = valeur initial + incrément \\
%\AC{nombre de répétitions} 			& \% nombre entier\\
%\AC{élément à répéter} 				& \% on peut utiliser les variables \\
%\end{tabular}
%
%
%
%
%\bigskip
%
%Exemples : \BS{}multido\AC{\BS{}i=2+-3}\AC{10}\AC{\BS{}i, } donne : \multido{\i=2+-3}{10}{\i, }

%\multido{\i=2+1, \n=5.0-0.5}{10}{\psframe[fillsyle=none](\i,\n)}
\begin{tabular}{|c|}\hline
\begin{pspicture}(6,3.5)
%   \put(0,0){\vector(1,0){8}}
   \multido{\i=1+1, \n=3.+-0.5}{5}{\psframe[fillstyle=none](\i,\n) }
\end{pspicture}   \\  \hline 

\BSS{multido} \Rnode*{A}{\AC{\textbf{\BS{i}}=1+1,\BS{n}=3.+-0.5}} \Rnode*[fillcolor=green]{B}{\AC{5}} \AC{\BS{psframe} \Rnode*[fillcolor=cyan]{C}{(\BS{i},\BS{n})} }
 \\ \hline 
\\
\rnode{AA}{\TFRGB{variable = valeur initiale+incrément}{variable = initial value + increment}} \hspace{1cm} \rnode{BB}{5 \TFRGB{fois}{times}} \hspace{1cm} \rnode{CC}{utilisation}
 \\ \hline
\end{tabular} 
 \ncline[linecolor=blue]{A}{AA}  \ncline[linecolor=blue]{B}{BB}  \ncline[linecolor=blue]{C}{CC}
 
\bigskip

\begin{tabular}{|l|l|} \hline
 \multicolumn{2}{|c|}{ \TFRGB{Types de variables}{variables types}  } \\  \hline  
\TFRGB{initiale}{initial} & dimension \\ \hline
d ou D				&  \TFRGB{longueur}{lenght}  \\  \hline
i ou I				&  \TFRGB{nombre entier}{integer}\\ \hline
n ou N 				& \TFRGB{nombre réel (même nombre de décimales)}{real} \\ \hline
r ou R 				&  \TFRGB{Réel (4 chiffres maxima de part et d'autre)}{real}\\ \hline
\end{tabular}


%
%\essais{pstrep2ess.tex}

\subsection[Psforeach]{  Commande psforeach \cite{pst-news10} }


 


\psset{fillcolor=yellow,fillstyle=solid,linecolor=blue,unit=1cm}


%
%\begin{tabular}{ll}
% \BSS{psforeach}				&   \\ 
%\AC{variable} 						& \% nom de la variable \\
%\AC{liste des valeurs} 				& \% liste des valeurs de la variable\\
%\AC{action} 						& \% qui utilise les valeurs de la variable \\
%\end{tabular}
%
%\bigskip
%\emph{Exemple: liste de valeurs}

\begin{tabular}{|l|} \hline  
%\begin{pspicture}[0,-.5](11,0.5)

\psforeach{\nA}{0, 1, 1.5, 3, 5,10}{\psdot[dotscale=2](\nA,0)}

%\end{pspicture}
\\
\\ \hline  
\BSS{psforeach}\Rnode{A}{\AC{{\red \BS{nA}}}}\Rnode{B}{\AC{0, 1, 1.5, 3, 5,10}}\Rnode{C}{\AC{\BS{}psdot[dotscale=2]({\red \BS{nA}},0)}}  \hspace{1cm}
\\
\\
 \hspace{1cm}  \Rnode{AA}{variable}  \hspace{1cm} \Rnode{BB}{\TFRGB{liste des valeurs}{list of values}} \hspace{1cm}  \Rnode{CC}{action}
\\ \hline 
\end{tabular} 
 \ncline[linecolor=blue]{A}{AA}  \ncline[linecolor=blue]{B}{BB}  \ncline[linecolor=blue]{C}{CC}
 
 
%\begin{tabular}{ll}
%$\backslash$psforeach\AC{$\backslash$nA}\AC{0, 1, 1.5, 3, 5,10}	& \% variable $\backslash$nA et ses 6 valeurs  \\  
%\AC{$\backslash$psdot[dotscale=2]($\backslash$nA,0)}			& \% un point à chaque abscisse $\backslash$nA  \\ \
%\end{tabular}
%
%\smallskip 
%
%\psforeach{\nA}{0, 1, 1.5, 3, 5,10}{\psdot[dotscale=2](\nA,0)}

\bigskip

%\emph{Exemple: }

\begin{tabular}{|l|} \hline

 \multicolumn{1}{|c|}{ \TFRGB{liste de valeurs avec pas régulier}{list of values at a regular step}  } \\  \hline  
\psforeach{\nA}{0, 1,..,10}{\psdot[dotscale=2](\nA,0)}
\\
\\  \hline  
 \hspace{1cm}  \BS{psforeach}\AC{\BS{nA}}\AC{{\red 0, 1,..,10}}\AC{\BS{psdot}[dotscale=2](\BS{nA},0)}  \hspace{1cm} 
\\ \hline 
\end{tabular} 


%\begin{tabular}{ll}
%$\backslash$psforeach\AC{$\backslash$nA}\AC{0, 1 ,..,10}	& \% variable $\backslash$nA de 0 à 10 par pas de 1  \\  
%\AC{$\backslash$psdot[dotscale=2]($\backslash$nA,0)}		& \% un point à chaque abscisse $\backslash$nA  \\ \
%\end{tabular}
% 
%\smallskip 



\bigskip


 
\begin{tabular}{|l|} \hline 

 \multicolumn{1}{|c|}{ \TFRGB{utilisation du numéro d'index}{use of the index number}  } \\  \hline  

\psforeach{\nA}{0, 1, 1.5, 2.25, 5,10}{\rput(\nA,0){\the\psLoopIndex}}
\\
\\ \hline  
\BS{psforeach}\AC{\BS{A}}\AC{0, 1, 1.5, 2.25, 5,10}\AC{\BS{rput}(\BS{nA},0)\AC{{\red\BS{the}\BS{psLoopIndex}}}}
\\ \hline  
\end{tabular} 

%$\backslash$psforeach\AC{$\backslash$nA}\AC{0, 1, 1.5, 2.25, 5,10}\AC{$\backslash$rput($\backslash$nA,0)\AC{$\backslash$the$\backslash$psLoopIndex}}
%\smallskip 
%
%\psforeach{\nA}{0, 1, 1.5, 2.25, 5,10}{\rput(\nA,0){\the\psLoopIndex}}



%\essais{pstrep3ess.tex}


\newpage


 %======================================================================
 
%\section[La géométrie]{La géométrie  \cite{pst-eucl} }
\SSCT{La géométrie}{Geometry}

Utilisation du module \textbf{pst-eucl} \label{pst-eucl}(consultez le fichier\textbf{ pst-eucl-doc.pdf} )




%\subsection{\'Elements de base}
\SbSSCT{\'Elements de base}{Basic elements}

\psset{fillcolor=yellow,linecolor=blue,dotscale=2}

\subsubsection{Points}
\psset{fillstyle=none,unit=1cm,arrowscale=1}
\newrgbcolor{Vert}{0 .8 0}
\label{geom}
%
\begin{tabular}{|c|c|} \hline
\TFRGB{axes par défaut}{default axes} & \TFRGB{axes personnalisés}{personalized axes} \\ \hline
\begin{pspicture}(-.5,-.5)(6,5.5)
\psaxes[ticksize=5,axesstyle=frame](5,5)
\pstGeonode(1,2){A}(3,1){A_1}(4,4){C}
\cnodeput{0}(2,4){D}{D}
\end{pspicture}
&
\begin{pspicture}(-.5,-.5)(6,5.5)
\psaxes[ticksize=5,axesstyle=frame](5,5)
\pstGeonode[linecolor=red](3,1){A}(2,2){B}(4,2){C}
\pnode(1,3){H}
\ncline[nodesepB=-2]{->}{A}{B}
\ncline[nodesepB=-1]{->}{A}{C}
\pstOIJGeonode(1,1){E}{A}{B}{C}(2,1){D}
\ncline[linestyle=dashed,linewidth=1pt,linecolor=red]{B}{E}
\ncline[linestyle=dashed,linewidth=1pt,linecolor=red]{C}{D}
\ncline[linestyle=dashed,linewidth=1pt,linecolor=red]{H}{D}
\end{pspicture}\\ \hline

\parbox[c]{.48\linewidth}{
 \BSS{pstGeonode}(1,2)\AC{A}(3,1)\AC{A\_1}(4,4)\AC{C} \BSI{pstGeonode}{pst-eucl}}
&
\parbox[c]{.48\linewidth}{
\BS{pstGeonode(3,1)\AC{A}(2,2)\AC{B}(4,2)\AC{C}}}
\\
 \BSS{cnodeput}\AC{0}(2,4)\AC{D}\AC{D}  \footnotemark[1]
&
\parbox[c]{.48\linewidth}{
\BSS{pstOIJGeonode}(1,1)\AC{E}{\red\AC{A}\AC{B}\AC{C}} (2,1)\AC{D} \BSI{pstOIJGeonode}{pst-eucl}}
\\ \hline
\end{tabular}

\footnotetext[1]{\TFRGB{On peut aussi utiliser les n\oe uds du module \textbf{pstnode}}{other possible nodes see} page \pageref{nodes}}

%\bigskip
%\begin{tabular}{|c|c|} \hline
%\begin{pspicture}(-.5,-.5)(6,5.5)
%\psaxes[ticksize=5,axesstyle=frame](5,5)
%\pstGeonode[PointName=none](1,2){A}(3,1){A_1}(4,4){C}
%\cnodeput{0}(2,4){D}{D}
%\end{pspicture}
%&
%\begin{pspicture}(-.5,-.5)(6,5.5)
%\psaxes[ticksize=5,axesstyle=frame](5,5)
%\pstGeonode[PosAngle=45](1,2){A}(3,1){A_1}(4,4){C}
%\cnodeput{0}(2,4){D}{D}
%\end{pspicture}
%\\ \hline
%\BS{pstGeonode}[\RDD{PointName}=none]& \BS{pstGeonode}[\RDD{PosAngle}=45] \\
%(1,2)\AC{A}(3,1)\AC{B}(4,4)\AC{C} \RDI{PointName}{pst-eucl} & 
%(1,2)\AC{A}(3,1)\AC{B}(4,4)\AC{C} \RDI{PosAngle}{pst-eucl}
%\\ \hline
%\end{tabular}
%--------------------------------------------------
\bigskip

\begin{center}
\begin{tabular}{|c|p{3cm}|c|p{3cm}|} \hline
\multicolumn{4}{|c|}{\TFRGB{Types de points}{Point type}}\\ \hline
\TFRGB{paramètre}{parameter} & \TFRGB{exemple}{sample} \footnotemark[2]  &\TFRGB{paramètre}{parameter} & \TFRGB{exemple}{sample}   \\ \hline
* & \pstGeonode[PointSymbol=*](1,0){A} &  o & \pstGeonode[PointSymbol=o](1,0){A}    \\
+ & \pstGeonode[PointSymbol=+](1,0){A} &  x & \pstGeonode[PointSymbol=x](1,0){A}  \\
asterisk & \pstGeonode[PointSymbol=asterisk](1,0){A} &  oplus & \pstGeonode[PointSymbol=oplus](1,0){A}   \\
otimes & \pstGeonode[PointSymbol=otimes](1,0){A} & | & \pstGeonode[PointSymbol=|](1,0){A}   \\
triangle & \pstGeonode[PointSymbol=triangle](1,0){A} &  triangle* & \pstGeonode[PointSymbol=triangle*](1,0){A}  \\
square & \pstGeonode[PointSymbol=square](1,0){A} &  
square* & \pstGeonode[PointSymbol=square*](1,0){A} \\  
diamond & \pstGeonode[PointSymbol=diamond](1,0){A}    &
diamond* & \pstGeonode[PointSymbol=diamond*](1,0){A} \\ 
pentagon & \pstGeonode[PointSymbol=pentagon](1,0){A}  &
pentagon* & \pstGeonode[PointSymbol=pentagon*](1,0){A} \\ \hline

\end{tabular}
\end{center}
\footnotetext[2]{linecolor=blue,fillcolor=yellow,dotscale=2}

%\newpage
%Exemple : 
\begin{tabular}{|c|} \hline  
\psset{unit=.5cm}

\begin{pspicture}(-.5,-.5)(5,5.5)
\psaxes[ticksize=5,axesstyle=frame](5,5)
\pstGeonode[linecolor=red,PointSymbol={otimes,asterisk,diamond*}](1,2){A}(3,1){B}(4,4){C}(3,3){D}(1,4){E}
\end{pspicture}

\psset{unit=1cm}
\\ \hline  
\BS{pstGeonode}{\red [PointSymbol=\AC{otimes,asterisk,diamond*}]} \\ (1,2)\AC{A}(3,1)\AC{B}(4,4)\AC{C}(3,3)\AC{D}(1,4)\AC{E}

\\ \hline 
\end{tabular} 



%------------------------------------------------------------
 

\begin{center}
\begin{tabular}{|c|c|c|c|} \hline
\multicolumn{4}{|c|}{\BS{pstGeonode}[PointNameSep=.7cm](1,1)\AC{A}}\\ \hline
\begin{pspicture}(-.5,-.5)(2,2.5)
\psaxes[ticksize=2,axesstyle=frame](2,2)
\pstGeonode(1,1){A}
\end{pspicture}
&
\begin{pspicture}(-.5,-.5)(2,2.5)
\psaxes[ticksize=2,axesstyle=frame](2,2)
\pstGeonode[PointNameSep=.7cm ](1,1){A}
\end{pspicture}
&
\begin{pspicture}(-.5,-.5)(2,2.5)
\psaxes[ticksize=2,axesstyle=frame](2,2)
\pstGeonode[PosAngle=45](1,1){A}
\end{pspicture}
&
\begin{pspicture}(-.5,-.5)(2,2.5)
\psaxes[ticksize=2,axesstyle=frame](2,2)
\pstGeonode[PointName=none](1,1){A}
\end{pspicture}
\\ \hline
{\blue\dft  }	&\RDD{PointNameSep}=.7cm \RDI{PointNameSep}{pst-eucl} &  \RDD{PosAngle}=45 \RDI{PosAngle}{pst-eucl}		&  \RDD{PointName}=none  \RDI{PointName}{pst-eucl} \\ \hline
			& {\blue \dft = 1em  }  & {\blue \dft = 0  } 	&      \\ \hline
\end{tabular}
\end{center}

%-----------------------------------------


\begin{center}
\begin{tabular}{|c|c|c|c|} \hline
\multicolumn{4}{|c|}{ \BS{pstGeonode}[\RDD{CurveType}=polyline](0,1)\AC{A}(1,0)\AC{B}(1.5,1.5)\AC{C} \RDI{CurveType}{pst-eucl} }\\ \hline
\begin{pspicture}(-.5,-.5)(2,2.5)
\psaxes[ticksize=2,axesstyle=frame](2,2)
\pstGeonode[linecolor=red,CurveType=polyline](0,1){A}(1,0){B}(1.5,1.5){C}
\end{pspicture}
&
\begin{pspicture}(-.5,-.5)(2,2.5)
\psaxes[ticksize=2,axesstyle=frame](2,2)
\pstGeonode[linecolor=red,CurveType=polygon](0,1){A}(1,0){B}(1.5,1.5){C}
\end{pspicture}
&
\begin{pspicture}(-.5,-.5)(2,2.5)
\psaxes[ticksize=2,axesstyle=frame](2,2)
\pstGeonode[linecolor=red,CurveType=curve](0,1){A}(1,0){B}(1.5,1.5){C}
\end{pspicture}
&
\begin{pspicture}(-.5,-.5)(2,2.5)
\psaxes[ticksize=2,axesstyle=frame](2,2)
\pstGeonode(0,1){A}(1,0){B}(1.5,1.5){C}
\ncline[linecolor=red]{A}{B}
\end{pspicture}
\\ \hline
\RDD{CurveType}=polyline \RDI{CurveType}{pst-eucl} &\RDD{CurveType}=polygon \RDI{CurveType}{pst-eucl} & \RDD{CurveType}=curve \RDI{CurveType}{pst-eucl}&  \BSS{ncline}\AC{A}\AC{B} \BSI{ncline}{pst-eucl} \footnotemark[1]\\ \hline
\end{tabular}
\end{center}

\footnotetext[1]{\TFRGB{On peut aussi utiliser les liaisons des noeuds voir}{other possibilities see} page \pageref{liaisons}}

%============================================================
\SbSbSSCT{Droites et segments de droite}{Lines}
 
%syntaxe :  \BSS{pstLineAB}[paramètres]\AC{A}\AC{B} 

\begin{tabular}{|c|c|c|} \hline
\multicolumn{3}{|c|}{\BSS{pstLineAB}[nodesepA=.5]\AC{A}\AC{B} \BSI{pstLineAB}{pst-eucl} }\\ \hline
\begin{pspicture}(-.5,-.5)(2,2.5)
\psaxes[ticksize=2,axesstyle=frame](2,2)
\pstGeonode(0.5,1.5){A}(1.5,0.5){B}
\pstLineAB[linecolor=red]{A}{B}
\end{pspicture}
&
\begin{pspicture}(-.5,-.5)(2,2.5)
\psaxes[ticksize=2,axesstyle=frame](2,2)
\pstGeonode(0.5,1.5){A}(1.5,0.5){B}
\pstLineAB[linecolor=red,nodesepA=.5]{A}{B}
\end{pspicture}
&
\begin{pspicture}(-.5,-.5)(2,2.5)
\psaxes[ticksize=2,axesstyle=frame](2,2)
\pstGeonode(0.5,1.5){A}(1.5,0.5){B}
\pstLineAB[linecolor=red,nodesepB=.5]{A}{B}
\end{pspicture}
\\  \hline
{\blue \dft{}} & [\RDD{nodesepA}=0.5] \RDI{nodesepA}{pst-eucl} & [\RDD{nodesepB}=0.5
\\  \hline
\begin{pspicture}(-.5,-.5)(2,2.5)
\psaxes[ticksize=2,axesstyle=frame](2,2)
\pstGeonode(0.5,1.5){A}(1.5,0.5){B}
\ncline[linecolor=red,nodesepA=-1]{A}{B}
\end{pspicture}
&
\begin{pspicture}(-.5,-.5)(2,2.5)
\psaxes[ticksize=2,axesstyle=frame](2,2)
\pstGeonode(0.5,1.5){A}(1.5,0.5){B}
\ncline[linecolor=red,nodesepB=-1]{A}{B}
\end{pspicture}
&
\begin{pspicture}(-.5,-.5)(2,2.5)
\psaxes[ticksize=2,axesstyle=frame](2,2)
\pstGeonode(0.5,1.5){A}(1.5,0.5){B}
\ncline[linecolor=red,nodesep=-1]{A}{B}
\end{pspicture} 
\\ \hline
[\RDD{nodesepA}=-1] \RDI{nodesepB}{pst-eucl} 
 &  [\RDD{nodesepB}=-1] \RDI{nodesepA}{pst-eucl} &
[\RDD{nodesep}=-1] \RDI{nodesep}{pst-eucl} \\ \hline
\end{tabular}


%===============================================




\SbSbSSCT{Marquage des droites}{Marks}
 
%syntaxe :   \BSS{pstMarkSegment}[paramètres]\AC{point A}\AC{point B}



\begin{tabular}{|c|c|c|c|} \hline
\multicolumn{4}{|c|}{\BSS{pstSegmentMark}[\RDD{SegmentSymbol}=pstslash]\AC{A}\AC{B} \BSI{pstSegmentMark}{pst-eucl} \RDI{SegmentSymbol}{pst-eucl}}\\ \hline
\begin{pspicture}(-.5,-.5)(2,2.5)
\psaxes[axesstyle=frame](2,2)
\pstGeonode(0.5,1.5){A}(1.5,0.5){B}
\pstSegmentMark[linecolor=red,SegmentSymbol=pstslash]{A}{B}
\end{pspicture}
&
\begin{pspicture}(-.5,-.5)(2,2.5)
\psaxes[axesstyle=frame](2,2)
\pstGeonode(0.5,1.5){A}(1.5,0.5){B}
\pstSegmentMark[linecolor=red,SegmentSymbol=pstslashh]{A}{B}
\end{pspicture}
&
\begin{pspicture}(-.5,-.5)(2,2.5)
\psaxes[axesstyle=frame](2,2)
\pstGeonode(0.5,1.5){A}(1.5,0.5){B}
\pstSegmentMark[linecolor=red,SegmentSymbol=pstslashhh]{A}{B}
\end{pspicture}
&
\begin{pspicture}(-.5,-.5)(2,2.5)
\psaxes[axesstyle=frame](2,2)
\pstGeonode(0.5,1.5){A}(1.5,0.5){B}
\pstSegmentMark[linecolor=red,SegmentSymbol=MarkCros]{A}{B}
\end{pspicture}

\\ \hline
pstslash \BDI{pstslash}{pst-eucl} &
pstslashh \BDI{pstslashh}{pst-eucl} & 
pstslashhh  \BDI{pstslashhh}{pst-eucl} & 
MarkCros \BDI{MarkCros}{pst-eucl} \\ \hline
%\end{tabular}
%
%\begin{tabular}{|c|c|c|c|} \hline
\begin{pspicture}(-.5,-.5)(2,2.5)
\psaxes[axesstyle=frame](2,2)
\pstGeonode(0.5,1.5){A}(1.5,0.5){B}
\pstSegmentMark[linecolor=red,SegmentSymbol=MarkHash]{A}{B}
\end{pspicture}
&
\begin{pspicture}(-.5,-.5)(2,2.5)
\psaxes[axesstyle=frame](2,2)
\pstGeonode(0.5,1.5){A}(1.5,0.5){B}
\pstSegmentMark[linecolor=red,SegmentSymbol=MarkHashh]{A}{B}
\end{pspicture}
&
\begin{pspicture}(-.5,-.5)(2,2.5)
\psaxes[axesstyle=frame](2,2)
\pstGeonode(0.5,1.5){A}(1.5,0.5){B}
\pstSegmentMark[linecolor=red,SegmentSymbol=MarkHashhh]{A}{B}
\end{pspicture}
&
\begin{pspicture}(-.5,-.5)(2,2.5)
\psaxes[axesstyle=frame](2,2)
\pstGeonode(0.5,1.5){A}(1.5,0.5){B}
\pstSegmentMark[linecolor=red,SegmentSymbol=MarkCross]{A}{B}
\end{pspicture}
\\ \hline
MarkHash  \BDI{MarkHash}{pst-eucl} & 
MarkHashh  \BDI{MarkHashh}{pst-eucl} & 
MarkHashhh  \BDI{MarkHashhh}{pst-eucl} &  
MarkCross  \BDI{MarkCross }{pst-eucl} \\ \hline
\end{tabular}
\bigskip



\begin{tabular}{|c|c|c|} \hline
\multicolumn{3}{|c|}{\BSS{pstSegmentMark}[\RDD{MarkAngle}=90]\AC{A}\AC{B} }\\ \hline
\begin{pspicture}(-.5,-.5)(2,2.5)
\psaxes[axesstyle=frame](2,2)
\pstGeonode(0.5,1.5){A}(1.5,0.5){B}
\pstSegmentMark[linecolor=red,MarkAngle=90]{A}{B}
\end{pspicture}
&
\begin{pspicture}(-.5,-.5)(2,2.5)
\psaxes[axesstyle=frame](2,2)
\pstGeonode(0.5,1.5){A}(1.5,0.5){B}
\pstSegmentMark[linecolor=red,MarkHashLength=.5]{A}{B}
\end{pspicture}
&
\begin{pspicture}(-.5,-.5)(2,2.5)
\psaxes[axesstyle=frame](2,2)
\pstGeonode(0.5,1.5){A}(1.5,0.5){B}
\pstSegmentMark[linecolor=red,MarkHashSep=.5]{A}{B}
\end{pspicture}
\\ \hline
\RDD{MarkAngle}=90  \RDI{MarkAngle}{pst-eucl}&
\RDD{MarkHashLength}=.5  \RDI{MarkHashLength}{pst-eucl}&
\RDD{MarkHashSep}=.5  \RDI{MarkHashSep}{pst-eucl}
\\ \hline
\dft : 45 & \dft : 1.25mm  &  \dft : .625mm
\\ \hline
\end{tabular}

%===============================================================

\subsubsection{Triangles}
 
%syntaxe : \textbf{\BS{}pstTriangle}[paramètres](xA,yA){A}(xB,yB){B}(xC,yC){C}

\begin{center}
\begin{tabular}{|c|} \hline
\begin{pspicture}(-.5,-.5)(5,5.5)
\psaxes[ticksize=5,axesstyle=frame](5,5)
 \pstTriangle[linecolor=red](1,2){A}(3,1){P}(4,4){Z} 
\end{pspicture}\\  \hline
 \BSS{pstTriangle}(1,2)\AC{A}(3,1)\AC{P}(4,4)\AC{Z} \BSI{pstTriangle}{pst-eucl} \\  \hline
\end{tabular}
\end{center}

%\subsubsection{Paramètres des triangles }
\begin{tabular}{|c|c|c|c|c|} \hline
\multicolumn{3}{|c|}{ \BS{pstTriangle}[\RDD{PointName}=none](0.5,1)\AC{A}(1.5,0.5)\AC{B}(1,1.5)\AC{C} \RDI{PointName}{pst-eucl}}\\ \hline
\begin{pspicture}(-.5,-.5)(2,2.5)
%\psaxes[axesstyle=frame](2,2)
 \pstTriangle[PointName=none](0.5,1){A}(1.5,0.5){B}(1,1.5){C}
\end{pspicture}
&
\begin{pspicture}(-.5,-.5)(2,2.5)
%\psaxes[axesstyle=frame](2,2)
 \pstTriangle[PointName=sommet](0.5,1){A}(1.5,0.5){B}(1,1.5){C}
\end{pspicture}
&
\begin{pspicture}(-.5,-.5)(2,2.5)
%\psaxes[axesstyle=frame](2,2)
\pstTriangle[PointName={X,Y,Z}](0.5,1){A}(1.5,0.5){B}(1,1.5){C}
\end{pspicture}
\\ \hline
\RDD{PointName}=none   & \RDD{PointName}=sommet &  >> A REVOIR <<
\\ \hline
\begin{pspicture}(-.5,-.5)(2,2.5)
%\psaxes[axesstyle=frame](2,2)
 \pstTriangle[PointNameA=none](0.5,1){A}(1.5,0.5){B}(1,1.5){C}
\end{pspicture}
&
\begin{pspicture}(-.5,-.5)(2,2.5)
%\psaxes[axesstyle=frame](2,2)
 \pstTriangle[PointNameB=X](0.5,1){A}(1.5,0.5){B}(1,1.5){C}
\end{pspicture}
&
\begin{pspicture}(-.5,-.5)(2,2.5)
%\psaxes[axesstyle=frame](2,2)
 \pstTriangle[PointNameC=sommet](0.5,1){A}(1.5,0.5){B}(1,1.5){C}
\end{pspicture}
\\ \hline

\RDD{PointNameA}=none  \RDI{PointNameA}{pst-eucl} &  \RDD{PointNameB}=X  \RDI{PointNameB}{pst-eucl} & \RDD{PointNameC}=sommet  \RDI{PointNameC}{pst-eucl}
\\ \hline 
\end{tabular}

\bigskip

\begin{tabular}{|c|c|c|c|} \hline
\multicolumn{4}{|c|}{\BS{pstTriangle}[\RDD{PosAngle}=45](0.5,1)\AC{A}(1.5,0.5)\AC{B}(1,1.5)\AC{C} \RDI{PosAngle}{pst-eucl} }
\\ \hline
\begin{pspicture}(-.5,-.5)(2,2.5)
\pstTriangle[PosAngle={45,30,40} ] (0.5,1){A}(1.5,0.5){B}(1,1.5){C}
\end{pspicture}
&
\begin{pspicture}(-.5,-.5)(2,2.5)
\pstTriangle[PosAngleA=90](0.5,1){A}(1.5,0.5){B}(1,1.5){C}
\end{pspicture}
&
\begin{pspicture}(-.5,-.5)(2,2.5)
\pstTriangle[PosAngleB=90](0.5,1){A}(1.5,0.5){B}(1,1.5){C}
\end{pspicture}
&
\begin{pspicture}(-.5,-.5)(2,2.5)
 \pstTriangle[PosAngleC=0](0.5,1){A}(1.5,0.5){B}(1,1.5){C}
\end{pspicture}
\\ \hline
\RDD{PosAngle}=180 &
\RDD{PosAngleA}=90 \RDI{PosAngleA}{pst-eucl}& \RDD{PosAngleB}=90 \RDI{PosAngleB}{pst-eucl}& \RDD{PosAngleC}=0 \RDI{PosAngleC}{pst-eucl} 
\\ \hline
\multicolumn{4}{|c|}{\dft : sur la bissectrice }
\\ \hline
\end{tabular}

\bigskip



\begin{tabular}{|c|c|c|c|} \hline
\multicolumn{4}{|c|}{\BS{pstTriangle}[\RDD{PointSymbolA}=o](0.5,1)\AC{A}(1.5,0.5)\AC{B}(1,1.5)\AC{C} }
\\ \hline
\begin{pspicture}(-.5,-.5)(2,2.5)
%\psaxes[axesstyle=frame](2,2)
 \pstTriangle[PointSymbolA=o](0.5,1){A}(1.5,0.5){B}(1,1.5){C}
\end{pspicture}
&
\begin{pspicture}(-.5,-.5)(2,2.5)
%\psaxes[axesstyle=frame](2,2)
 \pstTriangle[PointSymbolB=o](0.5,1){A}(1.5,0.5){B}(1,1.5){C}
\end{pspicture}
&
\begin{pspicture}(-.5,-.5)(2,2.5)
%\psaxes[axesstyle=frame](2,2)
 \pstTriangle[PointSymbolC=o](0.5,1){A}(1.5,0.5){B}(1,1.5){C}
\end{pspicture}
&
\begin{pspicture}(-.5,-.5)(2,2.5)
\psaxes[axesstyle=frame](2,2)
 \pstTriangle[PointSymbol=o](0.5,1){A}(1.5,0.5){B}(1,1.5){C}
\end{pspicture}
\\ \hline
\RDD{PointSymbolA}=o \RDI{PointSymbolA}{pst-eucl} & \RDD{PointSymbolB}=o \RDI{PointSymbolB}{pst-eucl} & \RDD{PointSymbolC}=o \RDI{PointSymbolC}{pst-eucl} & \RDD{PointSymbol}=o \RDI{PointSymbol}{pst-eucl}\\ \hline
\end{tabular}




%==================================================================

\subsubsection{Angles}
 


\begin{tabular}{|c|c|c|} \hline
\begin{pspicture}(-.5,-.5)(5,4.5)
\psaxes[xticksize=4,yticksize=5,axesstyle=frame](5,4)
\pstGeonode[CurveType=polyline,linewidth=2pt](1,1){A}(4,1){B}(4,3){C}
\pstRightAngle[linecolor=red]{A}{B}{C}
%\psline(A)(B)(C)
%\pstMarkAngle[linecolor=red]{A}{C}{B}{$\alpha$}
\end{pspicture}
&
\begin{pspicture}(-.5,-.5)(5,4.5)
\psaxes[xticksize=4,yticksize=5,axesstyle=frame](5,4)
\pstGeonode[CurveType=polyline,linewidth=2pt](1,1){A}(4,3){C}(4,1){B}
%\psline(A)(C)(B)
%\pstRightAngle[linecolor=red]{A}{B}{C}
\pstMarkAngle[linecolor=red]{A}{C}{B}{$\alpha$}
\end{pspicture}\\ \hline
 \BSS{pstRightAngle}\AC{A}\AC{B}\AC{C} \BSI{pstRightAngle}{pst-eucl}
&
 \BSS{pstMarkAngle}\AC{A}\AC{C}\AC{B}\AC{$\alpha$} \BSI{pstMarkAngle}{pst-eucl}
\\ \hline
\end{tabular}


\bigskip
%---------------------------------------------
%\subsubsection{}

\psset{linewidth=2pt}

\begin{tabular}{|c|c|c|} \hline
\multicolumn{3}{|c|}{\BS{pstRightAngle}[linecolor=red,\RDD{RightAngleType}=german]\AC{C}\AC{B}\AC{A}}\\ \hline
\begin{pspicture}(-.5,-.5)(2,2.5)
%\psaxes[axesstyle=frame](2,2)
\pnode(0,0){A}\pnode(2,0){B}\pnode(2,2){C}
\psline(A)(B)(C)
\pstRightAngle[linecolor=red,RightAngleType=german]{C}{B}{A}
\end{pspicture}
&
\begin{pspicture}(-.5,-.5)(2,2.5)
%\psaxes[axesstyle=frame](2,2)
\pnode(0,0){A}\pnode(2,0){B}\pnode(2,2){C}
\psline(A)(B)(C)
\pstRightAngle[linecolor=red,RightAngleType=suisseromand]{C}{B}{A}
\end{pspicture}
&
\begin{pspicture}(-.5,-.5)(2,2.5)
%\psaxes[axesstyle=frame](2,2)
\pnode(0,0){A}\pnode(2,0){B}\pnode(2,2){C}
\psline(A)(B)(C)
\pstRightAngle[linecolor=red]{C}{B}{A} 
\end{pspicture}
\\  \hline
\RDD{RightAngleType}= \BDD{german} \RDI{RightAngleType}{pst-eucl} \BDI{german}{pst-eucl} & {\red RightAngleType}=\BDD{suisseromand} \BDI{suisseromand}{pst-eucl}  &  {\blue\dft  }\\ \hline
\end{tabular}

\bigskip

\begin{tabular}{|c|c|c|} \hline
\multicolumn{2}{|c|}{\BS{pstRightAngle}[linecolor=red,\RDD{RightAngleSize}=1]\AC{C}\AC{B}\AC{A} \RDI{RightAngleSize}{pst-eucl} }\\ \hline
\begin{pspicture}(-.5,-.5)(2,2.5)
%\psaxes[axesstyle=frame](2,2)
\pnode(0,0){A}\pnode(2,0){B}\pnode(2,2){C}
\psline(A)(B)(C)
\pstRightAngle[linecolor=red,RightAngleSize=1]{C}{B}{A} 
\end{pspicture}
&
\begin{pspicture}(-.5,-.5)(2,2.5)
%\psaxes[axesstyle=frame](2,2)
\pnode(0,0){A}\pnode(2,0){B}\pnode(2,2){C}
\psline(A)(B)(C)
\pstRightAngle[linecolor=red,RightAngleSize=-.3]{C}{B}{A} 
\end{pspicture}
\\ \hline
RightAngleSize=1 & RightAngleSize=-.3
\\ \hline
\multicolumn{2}{|c|}{ \dft : RightAngleSize = 0.28 unit }
\\ \hline
\end{tabular}

\bigskip
%\subsubsection{Marquage des angles quelconques}




\begin{tabular}{|c|c|c|c|} \hline
\multicolumn{4}{|c|}{\BS{pstMarkAngle}[\RDD{LabelSep}=.5]\AC{A}\AC{C}\AC{B}\AC{\$\BS{alpha}\$} }\\ \hline
\begin{pspicture}(-.5,-.5)(2,2.5)
%\psaxes[axesstyle=frame](2,2)
\pnode(0,0){A} \pnode(2,0){B} \pnode(1,2){C}
\psline(A)(C)(B)
\pstMarkAngle[linecolor=red]{A}{C}{B}{$\alpha$}
%\pstMarkAngle[linecolor=red,LabelSep=.3cm]{C}{B}{A}{$\alpha$}
\end{pspicture}
&
\begin{pspicture}(-.5,-.5)(2,2.5)
%\psaxes[axesstyle=frame](2,2)
\pnode(0,0){A} \pnode(2,0){B} \pnode(1,2){C}
\psline(A)(C)(B)
\pstMarkAngle[linecolor=red,LabelSep=.5]{A}{C}{B}{$\alpha$}
\end{pspicture}
&
\begin{pspicture}(-.5,-.5)(2,2.5)
%\psaxes[axesstyle=frame](2,2)
\pnode(0,0){A} \pnode(2,0){B} \pnode(1,2){C}
\psline(A)(C)(B)
\pstMarkAngle[linecolor=red,LabelAngleOffset=10]{A}{C}{B}{$\alpha$}
%\pstMarkAngle[linecolor=red,LabelSep=.3cm]{C}{B}{A}{$\alpha$}
\end{pspicture}
&
\begin{pspicture}(-.5,-.5)(2,2.5)
%\psaxes[axesstyle=frame](2,2)
\pnode(0,0){A} \pnode(2,0){B} \pnode(1,2){C}
\psline(A)(C)(B)
\pstMarkAngle[linecolor=red,LabelAngleOffset=-10]{A}{C}{B}{$\alpha$}
%\pstMarkAngle[linecolor=red,LabelSep=.3cm]{C}{B}{A}{$\alpha$}
\end{pspicture}
\\ \hline
{\blue \dft } & \RDD{LabelSep}=.3cm \RDI{LabelSep}{pst-eucl} & \RDD{LabelAngleOffset}=10 \RDI{LabelAngleOffset}{pst-eucl} & \RDD{LabelAngleOffset}=-10 \RDI{LabelAngleOffset}{pst-eucl} \\ \hline
 			& {\blue \dft : 1} &  {\blue\dft : 0 }  & {\blue \dft : 0 }  \\ \hline
%\end{tabular}
%
%\begin{tabular}{|c|c|c|c|} \hline
\begin{pspicture}(-.5,-.5)(2,2.5)
%\psaxes[axesstyle=frame](2,2)
\pnode(0,0){A} \pnode(2,0){B} \pnode(1,2){C}
\psline(A)(C)(B)
\pstMarkAngle[linecolor=red,LabelRefPt=l]{A}{C}{B}{$\alpha$}

\end{pspicture}
&
\begin{pspicture}(-.5,-.5)(2,2.5)
%\psaxes[axesstyle=frame](2,2)
\pnode(0,0){A} \pnode(2,0){B} \pnode(1,2){C}
\psline(A)(C)(B)
\pstMarkAngle[linecolor=red,Mark=MarkCros]{A}{C}{B}{$\alpha$}
%\pstMarkAngle[linecolor=red,LabelSep=.3cm]{C}{B}{A}{$\alpha$}
\end{pspicture}
&
\begin{pspicture}(-.5,-.5)(2,2.5)
%\psaxes[axesstyle=frame](2,2)
\pnode(0,0){A} \pnode(2,0){B} \pnode(1,2){C}
\psline(A)(C)(B)
\pstMarkAngle[linecolor=red,MarkAngleRadius=.8]{A}{C}{B}{$\alpha$}

\end{pspicture}
&
\begin{pspicture}(-.5,-.5)(2,2.5)
%\psaxes[axesstyle=frame](2,2)
\pnode(0,0){A} \pnode(2,0){B} \pnode(1,2){C}
\psline(A)(C)(B)
\pstMarkAngle[linecolor=red,arrows=->,MarkAngleRadius=.8]{A}{C}{B}{$\alpha$}

\end{pspicture}
\\ \hline
\RDD{LabelRefPt}=l \RDI{LabelRefPt}{pst-eucl} & \RDD{Mark}=MarkCros  \RDI{Mark}{pst-eucl} &  \RDD{MarkAngleRadius}=.8 \RDI{MarkAngleRadius}{pst-eucl}& \RDD{arrows}=-> \RDI{arrows}{pst-eucl} \\ \hline
{\blue \dft : c } &  & {\blue \dft : .4  } &  ,MarkAngleRadius=.8 \\ \hline
\end{tabular}

\psset{linewidth=1pt}
%===========================================================


\SbSbSSCT{Cercles}{circles}

\begin{center}
\begin{tabular}{|c|c|} \hline
\begin{pspicture}(-.5,-.5)(4,4.5)
\psaxes[ticksize=4,axesstyle=frame](4,4)
\pstGeonode(2,2){A}(3.5,2){B}
\pstCircleOA[linecolor=red]{A}{B}
\end{pspicture}
&
\begin{pspicture}(-.5,-.5)(4,4.5)
\psaxes[ticksize=4,axesstyle=frame](4,4)
\pstGeonode(0.5,2){A}(3.5,2){B}
\pstCircleAB[linecolor=red]{A}{B}
\end{pspicture}
\\  \hline
 \BSS{pstCircleOA}\AC{A}\AC{B} \BSI{pstCircleOA}{pst-eucl} &  \BSS{pstCircleAB}\AC{A}\AC{B} \BSI{pstCircleAB}{pst-eucl}\\ \hline
\end{tabular}
\end{center}


%
%\BS{}pstCircleOA[linecolor=blue, Radius=\BS{}pstDistAB\AC{B}\AC{C}]\AC{A}\AC{}\\
%
%\BS{}pstCircleOA[linecolor=vert, Diameter=\BS{}pstDistAB\AC{A}\AC{C}]\AC{B}\AC{}\\

\begin{tabular}{|c|c|} \hline
\begin{pspicture}(-.5,-.5)(4,4.5)
\psaxes[ticksize=4,axesstyle=frame](4,4)
\pstGeonode(2,2){A}(1,1){B}(3,1){C}
\pstCircleOA[linecolor=red,Radius=\pstDistAB{B}{C}]{A}{}
\end{pspicture}
&
\begin{pspicture}(-.5,-.5)(4,4.5)
\psaxes[ticksize=4,axesstyle=frame](4,4)
\pstGeonode(2,2){A}(1,1){B}(3,1){C}
\pstCircleOA[linecolor=red,Diameter=\pstDistAB{B}{C}]{A}{}
\end{pspicture}
\\  \hline
\BS{pstCircleOA}[ & 
\BS{pstCircleOA}[ \\ 
\RDD{Radius}=\BDD{pstDistAB}\AC{B}\AC{C}] \AC{A}\AC{} \BSI{pstDistAB}{pst-eucl} \RDI{Radius}{pst-eucl}
& 
\RDD{Diameter}=\BDD{pstDistAB}\AC{B}\AC{C}] \AC{A}\AC{} \RDI{Diameter}{pst-eucl}\\ \hline
%\end{tabular}
%\bigskip
%
%\begin{tabular}{|c|c|} \hline
\begin{pspicture}(-.5,-.5)(4,4.5)
\psaxes[ticksize=4,axesstyle=frame](4,4)
\pstGeonode(2,2){A}(1,1){B}(3,1){C}
\pstCircleOA[linecolor=red,DistCoef=.5,Radius=\pstDistAB{B}{C}]{A}{}
\end{pspicture}
&
\begin{pspicture}(-.5,-.5)(4,4.5)
\psaxes[ticksize=4,axesstyle=frame](4,4)
\pstGeonode(2,2){A}(1,1){B}(3,1){C}
\pstCircleOA[linecolor=red,DistCoef=2,Diameter=\pstDistAB{B}{C}]{A}{}
\end{pspicture}
\\  \hline
\BS{pstCircleOA}[\RDD{DistCoef}=.5 & 
\BS{pstCircleOA}[\RDD{DistCoef}=2 \\ 
Radius=pstDistAB\AC{B}\AC{C}] \AC{A}\AC{} 
& 
Diameter=pstDistAB \AC{B}\AC{C}] \AC{A}\AC{} \RDI{DistCoef}{pst-eucl}\\ \hline
\end{tabular}

\bigskip
\begin{tabular}{|c|c|} \hline
\begin{pspicture}(-.5,-.5)(4,4.5)
\psaxes[ticksize=4,axesstyle=frame](4,4)
\pstGeonode(2,2){A} %(1,1){B}(3,1){C}
\pstCircleOA[linecolor=red,Radius=\pstDistVal{2}]{A}{}
\end{pspicture}
&
\begin{pspicture}(-.5,-.5)(4,4.5)
\psaxes[ticksize=4,axesstyle=frame](4,4)
\pstGeonode(2,2){A} %(1,1){B}(3,1){C}
\pstCircleOA[linecolor=red,Diameter=\pstDistVal{2}]{A}{}
\end{pspicture}
\\  \hline
\BS{}pstCircleOA{\red[Radius=\BDD{pstDistVal}\AC{2}]}\AC{A}\AC{} \BDI{pstDistVal}{pst-eucl}& 
\BS{}pstCircleOA{\red[Diameter=\BDD{pstDistVal}\AC{2}]}\AC{A}\AC{} \\ \hline
\end{tabular}



%-------------------------------------------------------------------------------
\SbSbSSCT{Arcs de cercle}{Arcs}


\begin{tabular}{|c|c|} \hline
\begin{pspicture}(-.5,-.5)(4,4.5)
\psaxes[ticksize=4,axesstyle=frame](4,4)
\pstGeonode(2,2){A}(3,1){B}(3,3){C}
\pstArcOAB[linecolor=red]{A}{B}{C}
\end{pspicture}
&
\begin{pspicture}(-.5,-.5)(4,4.5)
\psaxes[ticksize=4,axesstyle=frame](4,4)
\pstGeonode(2,2){A}(3,1){B}(3,3){C}
\pstArcnOAB[linecolor=red]{A}{B}{C}
\end{pspicture}
\\  \hline
 \BSS{pstArcOAB}\AC{A}\AC{B}\AC{C} \BSI{pstArcOAB}{pst-eucl} & 
 \BSS{pstArcnOAB}\AC{A}\AC{B}\AC{C} \BSI{pstArcnOAB}{pst-eucl}\\ \hline
\end{tabular}

%\newpage %<<<<<<<<<<<<<<<<<<<<<<<<<<<<<<<<<<<<<<<<<<<<<<<<

\SbSSCT{Point sur cercle}{Point on circle}

\begin{tabular}{|c|c|} \hline
%\multicolumn{2}{|c|}{ \BSS{pstCurvAbsNode}\AC{A}\AC{B}\AC{C}\AC{\BSS{pstDistVal}\AC{1}}
%\BSI{pstCurvAbsNode}{pst-eucl} 
% }
%\\ \hline
\begin{pspicture}(-.5,-.5)(4,4.5)
\psaxes[ticksize=4,axesstyle=frame](4,4)
\pstGeonode(2,2){A}(3.5,2){B}
\pstCircleOA{A}{B}
\pstCurvAbsNode[linecolor=red]{A}{B}{C}{10}
\end{pspicture}
&
\begin{pspicture}(-.5,-.5)(4,4.5)
\psaxes[ticksize=4,axesstyle=frame](4,4)
\pstGeonode(2,2){A}(3.5,2){B}
\pstCircleOA{A}{B}
\pstCurvAbsNode[linecolor=red]{A}{B}{C}{\pstDistVal{1}}
\end{pspicture}
\\  \hline
\BSS{pstCurvAbsNode}\AC{A}\AC{B}\AC{C} \AC{10} &
\BSS{pstCurvAbsNode}\AC{A}\AC{B}\AC{C}   \\
& \AC{\BSS{pstDistVal}\AC{1}} \BSI{pstDistVal}{pst-eucl}
\\ \hline
\begin{pspicture}(-.5,-.5)(4,4.5)
\psaxes[ticksize=4,axesstyle=frame](4,4)
\pstGeonode(2,2){A}(3.5,2){B}
\pstCircleOA{A}{B}
\pstCurvAbsNode[CurvAbsNeg=true,linecolor=red]{A}{B}{C}{10}
\end{pspicture}
&
\begin{pspicture}(-.5,-.5)(4,4.5)
\psaxes[ticksize=4,axesstyle=frame](4,4)
\pstGeonode(2,2){A}(3.5,2){B}
\pstCircleOA{A}{B}
\pstCurvAbsNode[CurvAbsNeg=true,linecolor=red]{A}{B}{C}{\pstDistVal{1}}
\end{pspicture}
\\  \hline
\BSS{pstCurvAbsNode} [\RDD{CurvAbsNeg}=true]  &
\BSS{pstCurvAbsNode} [\RDD{CurvAbsNeg}=true]    \\
\AC{A}\AC{B}\AC{C} \AC{10}
& \AC{A}\AC{B}\AC{C} \AC{\BSS{pstDistVal}\AC{1}} \RDI{CurvAbsNeg}{pst-eucl}
\\ \hline
\end{tabular}
 

%------------------------------------------------------------
\SbSbSSCT{Courbe générique}{Generic curve}
 


\begin{tabular}{|c|c|} \hline
\multicolumn{2}{|c|}{ \BS{pstGeonode}(2,2)\AC{A} (3,1)\AC{B\_1} (3,3)\AC{B\_2} (1,3)\AC{B\_3} \AC(1,1){B\_4} 
 }
\\ \hline
\begin{pspicture}(-.5,-.5)(4,4.5)
\psaxes[ticksize=4,axesstyle=frame](4,4)
\pstGeonode(2,2){A}(3,1){B_1}(3,3){B_2}(1,3){B_3}(1,1){B_4} 
\pstGenericCurve[linecolor=red]{B_}{2}{4}
\end{pspicture}
&
\begin{pspicture}(-.5,-.5)(4,4.5)
\psaxes[ticksize=4,axesstyle=frame](4,4)
\pstGeonode(2,2){A}(3,1){B_1}(3,3){B_2}(1,3){B_3}(1,1){B_4} %(4,4){C}
\pstGenericCurve[GenCurvFirst=A,linecolor=red]{B_}{1}{4}
\end{pspicture}
\\ \hline
\BSS{pstGenericCurve}\AC{B\_}\AC{2}\AC{4} \BSI{pstGenericCurve}{pst-eucl}&
\BS{pstGenericCurve} [\RDD{GenCurvFirst}=A] \AC{B\_}\AC{1}\AC{4} \RDI{GenCurvFirst}{pst-eucl}
\\ \hline
\begin{pspicture}(-.5,-.5)(4,4.5)
\psaxes[ticksize=4,axesstyle=frame](4,4)
\pstGeonode(2,2){A}(3,1){B_1}(3,3){B_2}(1,3){B_3}(1,1){B_4} 
\pstGenericCurve[GenCurvLast=A,linecolor=red]{B_}{1}{4}
\end{pspicture}
&
\begin{pspicture}(-.5,-.5)(4,4.5)
\psaxes[ticksize=4,axesstyle=frame](4,4)
\pstGeonode(2,2){A}(3,1){B_1}(3,3){B_2}(2,3){B_3}(1,3){B_4}(1,1){B_5}
\pstGenericCurve[GenCurvInc=2,linecolor=red]{B_}{1}{5}
\end{pspicture}\\ \hline
 \BSS{pstGenericCurve}[\RDD{GenCurvLast}=A]&
 \BS{pstGenericCurve}[\RDD{GenCurvInc}=2] \\ 
\AC{B\_}\AC{1}\AC{4}  \RDI{GenCurvLast}{pst-eucl} &
\AC{B\_}\AC{1}\AC{5}
\RDI{GenCurvInc}{pst-eucl}
\\ \hline
\end{tabular}

%\subsection[Transformations géométriques]{Transformations géométriques \cite{pst-eucl} }
\SbSSCT{Transformations géométriques}{Geometric transformations}

\SbSbSSCT{Symétrie par rapport à un point}{Central symmetry}

\psset{unit=.7cm}
\begin{tabular}{|c|c|c|} \hline
%\multicolumn{3}{|c|}{ \BSS{pstSymO}[linecolor=red]\AC{A}\AC{B}  }\\ \hline
\begin{pspicture}(-.5,-.5)(5,4.3)
\psaxes[xticksize=4,yticksize=5,axesstyle=frame](5,4)
\pstGeonode(2,2){A}(3,1){B} %(3,3){C}
\pstSymO[linecolor=red]{A}{B}
%\psline[linestyle=dotted](B)(B')
%\psline[linestyle=dashed](A)(B)
\end{pspicture}
&
\begin{pspicture}(-.5,-.5)(5,4.3)
\psaxes[xticksize=4,yticksize=5,axesstyle=frame](5,4)
\pstGeonode(2,2){A}(3,1){B} %(3,3){C}
\pstSymO[linecolor=Vert]{A}{B}[D]
%\psline[linestyle=dashed](A)(B)
%\psline[linestyle=dotted](B)(D)
\end{pspicture}
&
\begin{pspicture}(-.5,-.5)(5,4.3)
\psaxes[xticksize=4,yticksize=5,axesstyle=frame](5,4)
\pstGeonode(2,2){A}(4,1){B}(3,3){C}
\pstSymO[linecolor=red]{A}{B,C}[D,E]
\end{pspicture}
\\ \hline

\BSS{pstSymO}[linecolor=red]  \BSI{pstSymO}{pst-eucl} &  \BSS{pstSymO}[linecolor=Vert]  
&\BSS{pstSymO}[linecolor=red]    \\
\AC{A}\AC{B}  & \AC{A}\AC{B}[D] & \AC{A}\AC{B,C}[D,E]
 \\ \hline
\end{tabular}

\bigskip
\begin{tabular}{|c|c|c|} \hline
\begin{pspicture}(-.5,-.5)(5,4.3)
\psaxes[xticksize=4,yticksize=5,axesstyle=frame](5,4)
\pstGeonode(2,2){A}(4,1){B}(3,3){C}
\pstSymO[CodeFig=true]{A}{B,C}[D,E]
\end{pspicture}
&
\begin{pspicture}(-.5,-.5)(5,4.3)
\psaxes[xticksize=4,yticksize=5,axesstyle=frame](5,4)
\pstGeonode(2,2){A}(4,1){B}(3,3){C}
\pstSymO[CodeFig=true,CodeFigColor=red]{A}{B,C}[D,E]
\end{pspicture}
 \\ \hline
\BS{pstSymO}[\RDD{CodeFig}=true] & \BS{pstSymO}[CodeFig=true,\RDD{CodeFigColor}=red]   \\
 \AC{A}\AC{B,C}[D,E] \RDI{CodeFig}{pst-eucl} &  \AC{A}\AC{B,C}[D,E] \RDI{CodeFigColor}{pst-eucl}
 \\ \hline
 \dft : CodeFig = false & \dft : CodeFigColor = cyan
  \\ \hline
  


\begin{pspicture}(-.5,-.5)(5,4.3)
\psaxes[xticksize=4,yticksize=5,axesstyle=frame](5,4)
\pstGeonode(2,2){A}(4,1){B}(3,3){C}
\pstSymO[CodeFig=true,CodeFigStyle=dotted,CodeFigColor=red]{A}{B,C}[D,E]
\end{pspicture}
&
\begin{pspicture}(-.5,-.5)(5,4.3)
\psaxes[xticksize=4,yticksize=5,axesstyle=frame](5,4)
\pstGeonode(2,2){A}(4,1){B}(3,3){C}
\pstSymO[CodeFig=true,CodeFigStyle=solid,CodeFigColor=red]{A}{B,C}[D,E]
\end{pspicture}
\\ \hline
\BS{pstSymO}[CodeFig=true,\RDD{CodeFigStyle}=dotted]  & \BS{pstSymO}[CodeFig=true,\RDD{CodeFigStyle}=solid]   \\
 \AC{A}\AC{B,C}[D,E] \RDI{CodeFigStyle}{pst-eucl} &  \AC{A}\AC{B,C}[D,E] 
 \\ \hline
 \dft : CodeFigStyle = dashed & 
  \\ \hline
  

\end{tabular}

Autres options possibles : PointSymbol PosAngle PointName PointNameSep PtNameMath 


% <<<<<<<<<<<<<<<<<<<<<<<<<<<<<<<<<<<<<<<<<<<<<<<<<<<<<<<<<<<<

\SbSbSSCT{Symétrie par rapport à une droite}{Orthogonal symmetry}

\begin{tabular}{|c|c|} \hline
\multicolumn{2}{|c|}{ \BSS{pstOrtSym}[options]\AC{A}\AC{B}\AC{C}  \BSI{pstOrtSym}{pst-eucl} }\\ \hline
\begin{pspicture}(-.5,-.5)(5,5.3)
\psaxes[ticksize=5,axesstyle=frame](5,5)
\pstGeonode(1,1){A}(4,1){C}(3,3){B}
\pstOrtSym [linecolor=red]{A}{B}{C}
\psline[linestyle=dashed](B)(A)
\psline[linestyle=dotted](C)(C')
\end{pspicture}
&
\begin{pspicture}(-.5,-.5)(5,5.3)
\psaxes[ticksize=5,axesstyle=frame](5,5)
\pstGeonode(1,1){A}(4,1){C}(3,3){B}
\pstOrtSym [CodeFig=true,CodeFigColor=red]{A}{B}{C} %  <<<<<<<<<<<<<<<<<<<<<<<<<<<<<<<
\psline[linestyle=dashed](B)(A)
\end{pspicture}
\\ \hline
[linecolor=red] & [CodeFig=true,\RDD{CodeFigColor}=red] \RDI{CodeFigColor}{pst-eucl} \\ \hline
 &  {\blue \dft :  CodeFigColor=cyan } \\ \hline
\end{tabular}


%%--------------------------------------------------
\subsubsection{Rotation}

\begin{tabular}{|c|c|c|} \hline
\begin{pspicture}(-.5,-.5)(5,4.3)
\psaxes[xticksize=4,yticksize=5,axesstyle=frame](5,4)
\pstGeonode(2,2){A}(4,1){B} %(1,3){C}
\pstRotation[linecolor=red,CodeFig=true]{A}{B}
%\pstRotation[linecolor=Vert]{A}{C}[D]
\psline[linestyle=dotted](A)(B')
\psline(A)(B)
\end{pspicture}
&
\begin{pspicture}(-.5,-.5)(5,4.3)
\psaxes[xticksize=4,yticksize=5,axesstyle=frame](5,4)
\pstGeonode(1,3){A}(2,1){B} %(1,3){C}
\pstRotation[linecolor=red,RotAngle=45,CodeFig=true]{A}{B}[D]
\psline(B)(A)
\psline[linestyle=dotted](A)(D)
\end{pspicture}
&
\begin{pspicture}(-.5,-.5)(5,4.3)
\psaxes[xticksize=4,yticksize=5,axesstyle=frame](5,4)
\pstGeonode(2,2){A}(4,1){B}(1,3){C}
\pstRotation[linecolor=red,RotAngle=45,CodeFig=true]{A}{B,C}[D,E]
\psline  (B)(A) (C) %(D)
%\psline[linestyle=dashed](C)(A) %(E)
\psline[linestyle=dotted](A)(D)
\psline[linestyle=dotted](A)(E)
\end{pspicture}
\\ \hline
\BSS{pstRotation}[linecolor=red]  & \textbf{\BS{pstRotation}}[\RDD{RotAngle}=45]  & \textbf{\BS{pstRotation}}[\RDD{RotAngle}=45] \RDI{RotAngle}{pst-eucl} \\
\AC{A}\AC{B}  \BSI{pstRotation}{pst-eucl} & \AC{A}\AC{C}[D] \RDI{RotAngle}{pst-eucl}& \AC{A}\AC{B,C}[D,E] \\\hline
 & \multicolumn{2}{|c|} {\blue \dft : RotAngle=60 }\\
\hline
\end{tabular}




\bigskip
\begin{tabular}{|c|c|c|} \hline
\multicolumn{2}{|c|} {\BS{pstRotation}[CodeFig=true,CodeFigColor=red,\RDD{TransformLabel}=\BS{alpha}]\AC{A}\AC{B} \RDI{TransformLabel}{pst-eucl}  }
\\ \hline
\begin{pspicture}(-.5,-.5)(5,4.3)
\psaxes[xticksize=4,yticksize=5,axesstyle=frame](5,4)
\pstGeonode(2,2){A}(4,1){B} 
\pstRotation[CodeFig=true,CodeFigColor=red,TransformLabel=\alpha]{A}{B} %<<<<<<<<<<<<<<<<<<<<<<<<<<<<<<<<<<<<<<<<<<<<
\psline[linestyle=dashed](B)(A)
\end{pspicture}
&
\begin{pspicture}(-.5,-.5)(5,4.3)
\psaxes[xticksize=4,yticksize=5,axesstyle=frame](5,4)
\pstGeonode(2,2){A}(4,1){B}
\pstRotation[CodeFig=true,CodeFigColor=red,TransformLabel=\frac{\pi}{3}]{A}{B}[C] 
\psline[linestyle=dashed](B)(A)
\end{pspicture}
\\ \hline
\RDD{TransformLabel}=\BS{alpha} \RDI{TransformLabel}{pst-eucl} & \RDD{TransformLabel}=\BS{frac}\AC{\BS{pi}}\AC{3}  \RDI{TransformLabel}{pst-eucl} \\ \hline
\end{tabular}



\bigskip
%%
%%%>>>>>>>>>>>>>>>>>>> \BS{pstAngleAOB} incompatible avec pst-3dplot <<<<<<<<<<<<<<<<<<<<<<<
%%


\begin{tabular}{|c|c|c|} \hline
\multicolumn{3}{|c|} {\BS{pstRotation}[CodeFig=true,RotAngle=\RDD{pstAngleAOB}\AC{C}\AC{A}\AC{B}] \AC{A}\AC{B} \AC{D} \RDI{pstAngleAOB}{pst-eucl}  }
\\ \hline
\parbox{4cm}{
\begin{pspicture}(-.5,-.5)(5,5.3)
\psaxes[xticksize=5,yticksize=5,axesstyle=frame](5,5)
\pstGeonode(3,3){A}(1,1){B} (1,3){C}
\pstRotation[linecolor=red,RotAngle=\pstAngleAOB{C}{A}{B},CodeFig=true]{A}{B}[D]
\psline(B)(A)
%\rput (2,-1){\framebox{Rotation = \makebox[1cm][l]{\psPrintValue[decimals=2]{\pstAngleAOB{C}{A}{B}}}}}
%\rput (2,-1){Rotation = \psPrintValue[decimals=2] {\pstAngleAOB{C}{A}{B}}}
\psline[linestyle=dotted](A)(D)
\end{pspicture} \\
Rotation = \psPrintValue[decimals=2]{\pstAngleAOB{C}{A}{B}}}
&
\parbox{4cm}{
\begin{pspicture}(-.5,-.5)(5,5.3)
\psaxes[xticksize=5,yticksize=5,axesstyle=frame](5,5)
\pstGeonode(1,2){A}(4,1){B}(3,3){C}
\pstRotation[linecolor=red,RotAngle=\pstAngleAOB{C}{B}{A},CodeFig=true]{A}{B,C}[D,E]
%\rput* (2,-1){\framebox{Rotation = \makebox[1cm][l]{\psPrintValue[decimals=2]{\pstAngleAOB{C}{B}{A}}}}}
%\rput (2,-1){Rotation = \psPrintValue[decimals=2] {\pstAngleAOB{C}{B}{A}}}
\psline  (B)(A) (C) %(D)
\psline[linestyle=dashed](C)(A) %(E)
\psline[linestyle=dotted](A)(D)
\psline[linestyle=dotted](A)(E)
\end{pspicture}\\
Rotation = \psPrintValue[decimals=2]{\pstAngleAOB{C}{B}{A}}}
&
\parbox{4cm}{
\begin{pspicture}(-.5,-.5)(5,5.3)
\psaxes[xticksize=5,yticksize=5,axesstyle=frame](5,5)
\pstGeonode(2,2){A}(4,1){B}(3,3){C}
\pstRotation[linecolor=red,AngleCoef=3,RotAngle=\pstAngleAOB{C}{B}{A},CodeFig=true]{A}{B,C}[D,E]
%\psline  (B)(A) (C) %(D)

%\rput (2,-1){\framebox{Rotation = \makebox[1cm][l]{\psPrintValue[decimals=2] {  \pstAngleAOB{C}{B}{A} 3 mul}}}}
%\psline[linestyle=dashed](C)(A) %(E)
%\psline[linestyle=dotted](A)(D)
%\psline[linestyle=dotted](A)(E)
\end{pspicture}\\
Rotation = \psPrintValue[decimals=2] {  \pstAngleAOB{C}{B}{A} 3 mul}}
\\ \hline

\end{tabular}
%<<<<<<<<
%\bigskip
%%
%%%\begin{pspicture}(-.5,-.5)(5,4.3)
%%%\psaxes[xticksize=4,yticksize=5,axesstyle=frame](5,4)
%%%\pstGeonode(3,2){A}(3,3){B}(1,2){C}
%%%\pstRotation[RotAngle=45, CodeFig=true]{A}{D}[E] 
%%%%\pstRotation[RotAngle=\pstAngleAOB{C}{A}{B}, CodeFig=true]{A}{D}[E] 
%%%%\pstRotation[ RotAngle=\pstAngleAOB{B}{A}{C} ,CodeFigColor=red, CodeFig=true]{A}{B}[C]
%%%%\pstRotation[CodeFig=true,CodeFigColor=red,TransformLabel=\alpha]{A}{B}
%%%%\pstRotation[linecolor=Vert,RotAngle=\pstAngleAOB{B}{A}{C}]{A}{C}[D]
%%%
%%%\end{pspicture}
%%
%%%\begin{tabular}{|c|c|} \hline
%%%\begin{pspicture}(-.5,-.5)(5,4.3)
%%%\psaxes[xticksize=4,yticksize=5,axesstyle=frame](5,4)
%%%\pstGeonode(3,2){A}(3,3){B}(1,3){C}
%%%\pstRotation[RotAngle=\pstAngleAOB{B}{A}{C}]{A}{D}[E] %<<<<<<<<<<<<<<<<<<<<<<<<<<<<<<<<<<<<<
%%%\pstRotation[AngleCoef=.5, RotAngle=\pstAngleAOB{B}{A}{C},CodeFigColor=red, CodeFig=true]{A}{B}[C]
%%%\pstRotation[CodeFig=true,CodeFigColor=red,TransformLabel=\alpha]{A}{B}
%%%\pstRotation[linecolor=Vert,RotAngle=\pstAngleAOB{B}{A}{C}]{A}{C}[D]
%%%
%%%\end{pspicture}
%%%&
%%%\begin{pspicture}(-.5,-.5)(4,4.3)
%%%\psaxes[ticksize=4,axesstyle=frame](4,4)
%%%\pstGeonode(3,2){A}(3,3){B} %(1,3){C}
%%%\pstRotation[CodeFig=true,CodeFigColor=red,TransformLabel=\alpha]{A}{B}
%%%\pstRotation[RotAngle=\pstAngleAOB{B}{A}{C},CodeFig=true,CodeFigColor=red,TransformLabel=\widehat{BAC}]{A}{C}[D]
%%%\psline[linestyle=dashed](B)(A)(B')
%%%\psline[linestyle=dashed](C)(A)(D)
%%%\end{pspicture}
%%%\\ \hline
%%%RotAngle=\textbf{\BS{}pstAngleAOB}\AC{B}\AC{A}\AC{C} & {\red TransformLabel}=\BS{widehat}\AC{BAC}
%%%\\ \hline
%%%\end{tabular}
%%
%%
%%%------------------------------------------------------
\subsubsection{Translation}

\begin{tabular}{|c|c|c|} \hline
\BSS{pstTranslation}[options]\AC{B}\AC{A}\AC{C}  \BSI{pstTranslation}{pst-eucl} & 
\multicolumn{2}{|c|}{ \BSS{pstTranslation}[options]\AC{A}\AC{B}\AC{C}  \BSI{pstTranslation}{pst-eucl} }\\ \hline
\begin{pspicture}(-.5,-.5)(5,4.3)
\psaxes[xticksize=4,yticksize=5,axesstyle=frame](5,4)
\pstGeonode(1,2){A}(3,1){B}(3,2){C}
\pstTranslation[linecolor=red]{B}{A}{C}
\psline[linestyle=dashed](B)(A)
\psline[linestyle=dashed](C)(C')
\end{pspicture}
&
\begin{pspicture}(-.5,-.5)(5,4.3)
\psaxes[xticksize=4,yticksize=5,axesstyle=frame](5,4)
\pstGeonode(1,2){A}(3,1){B}(3,2){C}
\pstTranslation[linecolor=red,DistCoef=.5]{A}{B}{C}
\psline[linestyle=dashed](B)(A)
\psline[linestyle=dashed](C)(C')
\end{pspicture}
&
\begin{pspicture}(-.5,-.5)(6,4.3)
\psaxes[xticksize=4,yticksize=6,axesstyle=frame](6,4)
\pstGeonode(1,2){A}(3,1){B}(3,2){C}
\pstTranslation[CodeFig=true]{A}{B}{C}
\psline[linestyle=dashed](B)(A)
%\psline[linestyle=dashed](C)(C')
\end{pspicture}
\\ \hline
 & [\RDD{DistCoef}=0.5] \RDI{DistCoef}{pst-eucl} & [\RDD{CodeFig}=true] \RDI{CodeFig}{pst-eucl} 
\\ \hline

\end{tabular}


%
%%--------------------------------------------------------------------------
\SbSbSSCT{Homothétie}{Homothetie}

\begin{tabular}{|c|c|c|} \hline
\multicolumn{2}{|c|}{ \BSS{pstHomO}\AC{A}\AC{B,C}[D,E]  \BSI{pstHomO}{pst-eucl} }\\ \hline
\begin{pspicture}(-.5,-.5)(5,4.3)
\psaxes[xticksize=4,yticksize=5,axesstyle=frame](5,4)
\pstGeonode(0,2){A}(4,1){B}(4,3){C}
\pstLineAB[nodesepB=-1,linestyle=dashed]{A}{B}
\pstLineAB[nodesepB=-1,linestyle=dashed]{A}{C}
\pstHomO[linecolor=red]{A}{B,C}[D,E]
\end{pspicture}
&
\begin{pspicture}(-.5,-.5)(5,4.3)
\psaxes[xticksize=4,yticksize=5,axesstyle=frame](5,4)
\pstGeonode(0,2){A}(4,1){B}(4,3){C}
\pstLineAB[nodesepB=-1,linestyle=dashed]{A}{B}
\pstLineAB[nodesepB=-1,linestyle=dashed]{A}{C}
\pstHomO[linecolor=red,HomCoef=.75]{A}{B,C}[D,E]
\end{pspicture} 


\\ \hline
 &
[\RDD{HomCoef}=.75] \RDI{HomCoef}{pst-eucl} 
\\ \hline
\end{tabular}



%%-------------------------------------------------------------
\SbSbSSCT{Projection orthogonale}{Orthogonal projection}

\begin{tabular}{|c|} \hline  
\begin{pspicture}(-.5,-.5)(5,5.3)
\psaxes[ticksize=5,axesstyle=frame](5,5)
 \pstTriangle(1,3){A}(3,1){B}(4,4){C}
 \pstProjection[CodeFig=true, CodeFigColor=red]{A}{B}{C}[D]
\end{pspicture}
\\ \hline  
\BSS{pstProjection}[CodeFig=true,CodeFigColor=red]\AC{A}\AC{B}\AC{C}[D]  \BSI{pstProjection}{pst-eucl}
\\ \hline 
\end{tabular} 



%\subsection[Constructions particulières en géométrie ]{Constructions particulières en géométrie \cite{pst-eucl} }
\SbSSCT{Constructions particulières en géométrie}{Particular constructions}

\SbSbSSCT{Point milieu}{midpoint}

%\verb \pstMiddleAB %[hpari]{hAi}{hBi}{hIi}
\begin{tabular}{|c|c|}\hline 
\multicolumn{2}{|c|}{ \BSS{pstMiddleAB}[linecolor=red]\AC{A}\AC{B}\AC{C} \BSI{pstMiddleAB}{pst-eucl} }\\ \hline 
\begin{pspicture}(-.5,-.5)(4,3.5)
\psaxes[xticksize=3,yticksize=4,axesstyle=frame](4,3)
\pstGeonode(1,2){A}(3,1){B}
\pstMiddleAB[linecolor=red]{A}{B}{C}
\end{pspicture}
&  
\begin{pspicture}(-.5,-.5)(4,3.5)
\psaxes[xticksize=3,yticksize=4,axesstyle=frame](4,3)
\pstGeonode(1,2){A}(3,1){B}
\pstMiddleAB[CodeFig,CodeFigColor=green]{A}{B}{C}
\end{pspicture}
\\ \hline  
[linecolor=red]
&  
[CodeFig,CodeFigColor=green]
\\ \hline 
\end{tabular} 

%\begin{center}
%\begin{tabular}{|c|} \hline
%\begin{pspicture}(-.5,-.5)(4,3.5)
%\psaxes[xticksize=3,yticksize=4,axesstyle=frame](4,3)
%\pstGeonode(1,2){A}(3,1){B}
%\pstMiddleAB[linecolor=red]{A}{B}{C}
%\end{pspicture}
%\\ \hline
%\BSS{pstMiddleAB}[linecolor=red]\AC{A}\AC{B}\AC{C} \BSI{pstMiddleAB}{pst-eucl} \\ \hline
%\end{tabular}
%\end{center}





%--------------------

\SbSbSSCT{Centre de gravité d'un triangle}{Center of gravity of a triangle}

%\verb \pstCGravABC %[hpari]{hAi}{hBi}{hCi}{hGi}



\begin{center}
\begin{tabular}{|c|} \hline
\begin{pspicture}(-.5,-.5)(5,5.5)
\psaxes[ticksize=5,axesstyle=frame](5,5)
\pstTriangle(1,2){A}(3,1){B}(4,4){C}
\pstCGravABC[linecolor=red]{A}{B}{C}{G}
\end{pspicture}
\\ \hline
\BSS{pstCGravABC}[linecolor=red]\AC{A}\AC{B}\AC{C}\AC{G} \BSI{pstCGravABC}{pst-eucl} \\ \hline
\end{tabular}
\end{center}

%----------

\SbSbSSCT{Centre du cercle circonscrit d'un triangle}{Circumcenter of a triangle}


\begin{tabular}{|c|c|c|}\hline 
\multicolumn{3}{|c|}{ \BSS{pstCircleABC}\AC{A}\AC{B}\AC{C}\AC{O} \BSI{pstCircleABC}{pst-eucl} }
\\ \hline 
\begin{pspicture}(-.5,-.5)(5,5.5)
\psaxes[ticksize=5,axesstyle=frame](5,5)
\pstTriangle(1,2){A}(4,1){B}(3,4){C}
\pstCircleABC[linecolor=red,linewidth=1pt]{A}{B}{C}{O}
\end{pspicture}
&  
\begin{pspicture}(-.5,-.5)(5,5.5)
\psaxes[ticksize=5,axesstyle=frame](5,5)
\pstTriangle(1,2){A}(4,1){B}(3,4){C}
\pstCircleABC[CodeFig,CodeFigColor=red,linewidth=1pt]{A}{B}{C}{O}
\end{pspicture}
&  
\begin{pspicture}(-.5,-.5)(5,5.5)
\psaxes[ticksize=5,axesstyle=frame](5,5)
\pstTriangle(1,2){A}(4,1){B}(3,4){C}
\pstCircleABC[CodeFig,CodeFigColor=red,DrawCirABC=false,linewidth=1pt]{A}{B}{C}{O}
\end{pspicture}
\\ \hline  
[linecolor=red]
&  
[CodeFig,CodeFigColor=green]
&
[CodeFig,CodeFigColor=green \\
&& ,\RDD{DrawCirABC=false}  \RDI{DrawCirABC}{pst-eucl}
\\ \hline 
\end{tabular}
\bigskip

\begin{tabular}{|c|c|c|}\hline 
\multicolumn{3}{|c|}{ \BS{pstCircleABC}[CodeFig,CodeFigColor=red,\RDD{SegmentSymbolA}=MarkCros]\AC{A}\AC{B}\AC{C}\AC{O} \RDI{SegmentSymbolA}{pst-eucl} }
\\ \hline 
\begin{pspicture}(-.5,-.5)(5,5.5)
\psaxes[ticksize=5,axesstyle=frame](5,5)
\pstTriangle(1,2){A}(4,1){B}(3,4){C}
\pstCircleABC[CodeFig,CodeFigColor=red,linewidth=1pt,SegmentSymbolA=MarkCros]{A}{B}{C}{O}
\end{pspicture}
&  
\begin{pspicture}(-.5,-.5)(5,5.5)
\psaxes[ticksize=5,axesstyle=frame](5,5)
\pstTriangle(1,2){A}(4,1){B}(3,4){C}
\pstCircleABC[CodeFig,CodeFigColor=red,linewidth=1pt,SegmentSymbolB=MarkCros]{A}{B}{C}{O}
\end{pspicture}
&  
\begin{pspicture}(-.5,-.5)(5,5.5)
\psaxes[ticksize=5,axesstyle=frame](5,5)
\pstTriangle(1,2){A}(4,1){B}(3,4){C}
\pstCircleABC[CodeFig,CodeFigColor=red,linewidth=1pt,,SegmentSymbolC=MarkCros]{A}{B}{C}{O}
\end{pspicture}
\\ \hline  
\RDD{SegmentSymbolA}=MarkCros \RDI{SegmentSymbolA}{pst-eucl}
&  
\RDD{SegmentSymbolB}=MarkCros \RDI{SegmentSymbolB}{pst-eucl}
&
\RDD{SegmentSymbolC}=MarkCros \RDI{SegmentSymbolC}{pst-eucl}
\\ \hline 
\end{tabular}

%
%\begin{center}
%\begin{tabular}{|c|} \hline
%\begin{pspicture}(-.5,-.5)(5,5.5)
%\psaxes[ticksize=5,axesstyle=frame](5,5)
%\pstTriangle(1,2){A}(4,1){B}(3,4){C}
%\pstCircleABC[linecolor=red,linestyle=dashed]{A}{B}{C}{O}
%\end{pspicture}\\ \hline
%\BSS{pstCircleABC}\AC{A}\AC{B}\AC{C}\AC{O} \BSI{pstCircleABC}{pst-eucl} \\ \hline
%\end{tabular}
%\end{center}





%------------

\SbSbSSCT{Perpendiculaire par rapport à une droite}{Perpendicular to a  line}

\begin{tabular}{|c|c|}\hline 
\multicolumn{2}{|c|}{ \BSS{pstMediatorAB}\AC{A}\AC{B}\AC{C}\AC{D} \BSI{pstMediatorAB}{pst-eucl}  }\\ \hline 
\begin{pspicture}(-.5,-.5)(6,4.5)
\psaxes[xticksize=4,yticksize=6,axesstyle=frame](6,4)
\pstGeonode[CurveType=polyline](1,2){A}(5,1){B}

\pstMediatorAB[linecolor=red]{A}{B}{C}{D}
\end{pspicture}
&  
\begin{pspicture}(-.5,-.5)(6,4.5)
\psaxes[xticksize=4,yticksize=6,axesstyle=frame](6,4)
\pstGeonode[CurveType=polyline](1,2){A}(5,1){B}

\pstMediatorAB[CodeFig,CodeFigColor=red]{A}{B}{C}{D}
\end{pspicture}
\\ \hline  
[linecolor=red]
&  
[CodeFig,CodeFigColor=red]
\\ \hline 
\end{tabular} 
\bigskip

\begin{tabular}{|c|c|c|}\hline 
\multicolumn{3}{|c|}{ \BS{pstMediatorAB}[linecolor=red,PointName=none,\RDD{nodesep}=-1]\AC{A}\AC{B}\AC{C}\AC{D} \RDI{nodesep}{pst-eucl}  }\\ \hline 
\begin{pspicture}(-.5,-.5)(5,4.5)
\psaxes[xticksize=4,yticksize=4,axesstyle=frame](4,4)
\pstGeonode[CurveType=polyline](1,2){A}(3,1){B}
\pstMediatorAB[linecolor=red,PointName=none,nodesep=-1]{A}{B}{C}{D}
\end{pspicture}
&
\begin{pspicture}(-.5,-.5)(5,4.5)
\psaxes[xticksize=4,yticksize=4,axesstyle=frame](4,4)
\pstGeonode[CurveType=polyline](1,2){A}(3,1){B}
\pstMediatorAB[linecolor=red,PointName=none,nodesepA=-2]{A}{B}{C}{D}
\end{pspicture}
&
\begin{pspicture}(-.5,-.5)(5,4.5)
\psaxes[xticksize=4,yticksize=4,axesstyle=frame](4,4)
\pstGeonode[CurveType=polyline](1,2){A}(3,1){B}
\pstMediatorAB[linecolor=red,PointName=none,nodesepB=-0.5]{A}{B}{C}{D}
\end{pspicture}
\\ \hline 
\RDD{nodesep}=-1 \RDI{nodesep}{pst-eucl} & 
\RDD{nodesepA}=-2 \RDI{nodesepA}{pst-eucl} &
\RDD{nodesepB}=-0.5 \RDI{nodesepA}{pst-eucl}
\\ \hline 
\begin{pspicture}(-.5,-.5)(5,4.5)
\psaxes[xticksize=4,yticksize=4,axesstyle=frame](4,4)
\pstGeonode[CurveType=polyline](1,2){A}(3,1){B}
\pstMediatorAB[linecolor=red,PointName=none,nodesep=0.3]{A}{B}{C}{D}
\end{pspicture}
&
\begin{pspicture}(-.5,-.5)(5,4.5)
\psaxes[xticksize=4,yticksize=4,axesstyle=frame](4,4)
\pstGeonode[CurveType=polyline](1,2){A}(3,1){B}
\pstMediatorAB[linecolor=red,PointName=none,nodesepA=0.3]{A}{B}{C}{D}
\end{pspicture}
&
\begin{pspicture}(-.5,-.5)(5,4.5)
\psaxes[xticksize=4,yticksize=4,axesstyle=frame](4,4)
\pstGeonode[CurveType=polyline](1,2){A}(3,1){B}
\pstMediatorAB[linecolor=red,PointName=none,nodesepB=0.3]{A}{B}{C}{D}
\end{pspicture}
\\ \hline 
\RDD{nodesep}=0.3 & \RDD{nodesepA}=0.3 &\RDD{nodesepB}=0.3
\\ \hline 

\end{tabular}


%\begin{center}
%\begin{tabular}{|c|c|} \hline
%\begin{pspicture}(-.5,-.5)(4,3.5)
%\psaxes[xticksize=3,yticksize=4,axesstyle=frame](4,3)
%\pstGeonode[CurveType=polyline](1,2){A}(3,1){B}
%
%\pstMediatorAB[linecolor=red]{A}{B}{C}{D}
%\end{pspicture}
%\\ \hline
%\BSS{pstMediatorAB}\AC{A}\AC{B}\AC{C}\AC{D} \BSI{pstMediatorAB}{pst-eucl} \\ \hline
%\end{tabular}
%\end{center}

%---------

\SbSbSSCT{Bissectrice d'un angle}{Bissector}
%\verb \pstBissectBAC %[hpari]{hBi}{hAi}{hCi}{hNi}
%
%\verb \pstOutBissectBAC %[hpari]{hBi}{hAi}{hCi}{hNi}

\begin{center}
\begin{tabular}{|c|c|} \hline
\begin{pspicture}(-.5,-1.5)(5.5,3.5)
\psaxes[xticksize=4,yticksize=5,axesstyle=frame](0,-1)(5,3)
\pstGeonode[CurveType=polyline](1.5,2){A}(3,1){B}(4,2){C}
%\pstBissectBAC[linecolor=red]{C}{B}{A}{D}
\pstBissectBAC[linecolor=red]{A}{B}{C}{D}
\pstCircleOA[linestyle=dotted,linecolor=red]{B}{A}
%\pstCircleOA[linestyle=dotted,linecolor=red]{B}{C}
\end{pspicture}
&
\begin{pspicture}(-.5,-1.5)(5.5,3.5)
\psaxes[xticksize=4,yticksize=5,axesstyle=frame](0,-1)(5,3)
\pstGeonode[CurveType=polyline](1.5,2){A}(3,1){B}(4,2){C}
\pstOutBissectBAC[linecolor=red]{A}{B}{C}{D}
%\pstOutBissectBAC[linecolor=cyan]{C}{B}{A}{E}
\pstCircleOA[linestyle=dotted,linecolor=red]{B}{A}
%\pstCircleOA[linestyle=dotted,linecolor=cyan]{B}{C}
\end{pspicture}  
\\ \hline
\BSS{pstBissectBAC}\AC{A}\AC{B}\AC{C}\AC{D} \BSI{pstBissectBAC}{pst-eucl} & 
\BSS{pstOutBissectBAC}\AC{A}\AC{B}\AC{C}\AC{D} 
\BSI{pstOutBissectBAC}{pst-eucl}
%\BS{pstBissectBAC}\AC{A}\AC{B}\AC{C}\AC{E} & 
%\BS{pstOutBissectBAC}\AC{C}\AC{B}\AC{A}\AC{E} 
\\ \hline
\begin{pspicture}(-.5,-1.5)(5.5,3.5)
\psaxes[xticksize=4,yticksize=5,axesstyle=frame](0,-1)(5,3)
\pstGeonode[CurveType=polyline](1.5,2){A}(3,1){B}(4,2){C}
\pstBissectBAC[linecolor=red]{C}{B}{A}{D}
\pstCircleOA[linestyle=dotted,linecolor=red]{B}{C}
\end{pspicture}
&
\begin{pspicture}(-.5,-1.5)(5.5,3.5)
\psaxes[xticksize=4,yticksize=5,axesstyle=frame](0,-1)(5,3)
\pstGeonode[CurveType=polyline](1.5,2){A}(3,1){B}(4,2){C}
\pstOutBissectBAC[linecolor=red]{C}{B}{A}{D}
\pstCircleOA[linestyle=dotted,linecolor=cyan]{B}{C}
\end{pspicture}  
\\ \hline
\BS{pstBissectBAC}\AC{A}\AC{B}\AC{C}\AC{D} & 
\BS{pstOutBissectBAC}\AC{C}\AC{B}\AC{A}\AC{D} 
\\ \hline
\end{tabular}
\end{center}

\subsection[Intersections]{Intersections \cite{pst-eucl} }

\SbSbSSCT{Intersection de deux droites}{Intersection of two lines}
%\verb \pstInterLL %[hpari]{hAi}{hBi}{hCi}{hDi}{hMi}

\begin{center}
\begin{tabular}{|c|} \hline
\begin{pspicture}(-.5,-.5)(5,4.5)
\psaxes[xticksize=4,yticksize=5,axesstyle=frame](5,4)
\pstGeonode(1,3){A} (4,1){B} (2,1){C} (4,3.5){D}

\psline[linecolor=Vert](A)(B)
\psline[linecolor=cyan](C)(D)
\pstInterLL[linecolor=red]{A}{B}{C}{D}{E}
\end{pspicture}
\\ \hline
\BSS{pstInterLL}[linecolor=red]{\Vert\AC{A}\AC{B}}{\cyan\AC{C}\AC{D}}{\red\AC{E}}  \BSI{pstInterLL}{pst-eucl}\\ \hline
\end{tabular}
\end{center}

%-----------------

\SbSbSSCT{Intersection d'une droite et un cercle}{Intersection of a line and a circle}

\begin{center}
\begin{tabular}{|c|} \hline
\begin{pspicture}(-.5,-.5)(5,4.5)
\psaxes[xticksize=4,yticksize=5,axesstyle=frame](5,4)
\pstGeonode(1,3.5){A}(4,1){B}(2,2){C}(1,2){D}
\psline[linecolor=Vert](A)(B)
\pstCircleOA[linecolor=cyan]{C}{D}
\pstInterLC[linecolor=red]{A}{B}{C}{D}{E}{F}
\end{pspicture}
\\ \hline
\BSS{pstInterLC}{\Vert\AC{A}\AC{B}}{\cyan\AC{C}\AC{D}}{\red\AC{E}\AC{F}} \BSI{pstInterLC}{pst-eucl} \\ \hline
\end{tabular}
\end{center}

\bigskip
\begin{center}
	\begin{tabular}{|c|} \hline
\begin{pspicture}(-.5,-.5)(5,4.5)
\psaxes[xticksize=4,yticksize=5,axesstyle=frame](5,4)
\pstGeonode(1,3.5){A}(4,1){B}(2,2){C}
\psline[linecolor=Vert](A)(B)
\pscircle[linestyle=dotted](C){1}
\pstInterLC[linecolor=red,Radius=\pstDistVal{1}]{A}{B}{C}{}{E}{F}
\end{pspicture}
\\ \hline
\BSS{pstInterLC}[Radius=\BS{pstDistVal}\AC{1}]{\Vert\AC{A}\AC{B}}{\cyan\AC{C}\AC{}}{\red\AC{E}\AC{F}} \BSI{pstInterLC}{pst-eucl} \\ \hline

\end{tabular}
\end{center}


%------------

\SbSbSSCT{Intersection de deux cercles}{Intersection of two circles}


\begin{center}
\begin{tabular}{|c|} \hline
\begin{pspicture}(-.5,-.5)(5,5.5)
\psaxes[ticksize=5,axesstyle=frame](5,5)
\pstGeonode(2,3){A}(1,3){B}(3,2){C}(4,1){D}
\pstCircleOA[linecolor=Vert]{A}{B}
\pstCircleOA[linecolor=cyan]{C}{D}
\pstInterCC[linecolor=red,PosAngleA=45,PosAngleB=225,]{A}{B}{C}{D}{E}{F}
\end{pspicture}
\\ \hline
\BSS{pstInterCC}[linecolor=red,PosAngleA=45,PosAngleB=225]\\
{\Vert\AC{A}\AC{B}}{\cyan\AC{C}\AC{D}}{\red\AC{E}\AC{F}} \BSI{pstInterCC}{pst-eucl}
\\ \hline
\end{tabular}
\end{center}

\begin{center}
\begin{tabular}{|c|} \hline
\begin{pspicture}(-.5,-.5)(5,5.5)
\psaxes[ticksize=5,axesstyle=frame](5,5)
\pstGeonode(2,3){A}(1,3){B}(3,2){C}(4,1){D}
\pstInterCC[PosAngleA=135,PosAngleB=135, CodeFig=true, CodeFigColor=red]{A}{B}{C}{D}{E}{F}
\end{pspicture} \\ \hline
CodeFig=true  \\ \hline
\end{tabular}
\end{center}

\begin{tabular}{|c|c|c|} \hline
\begin{pspicture}(-.5,-.5)(5,5.5)
\psaxes[ticksize=5,axesstyle=frame](5,5)
\pstGeonode(2,3){A}(1,3){B}(3,2){C}(4,1){D}
\pstInterCC[PosAngleA=135,PosAngleB=135, CodeFig=true,CodeFigAarc=false, CodeFigColor=red]{A}{B}{C}{D}{E}{F}
\end{pspicture}
&
\begin{pspicture}(-.5,-.5)(5,5.5)
\psaxes[ticksize=5,axesstyle=frame](5,5)
\pstGeonode(2,3){A}(1,3){B}(3,2){C}(4,1){D}
\pstInterCC[PosAngleA=135,PosAngleB=135,CodeFig=true, CodeFigBarc=false, CodeFigColor=red]{A}{B}{C}{D}{E}{F}
\end{pspicture}
\\ \hline
 CodeFig=true & CodeFig=true \\
 \RDD{CodeFigAarc}=false  \RDI{CodeFigAarc}{pst-eucl} & \RDD{CodeFigBarc}=false  \RDI{CodeFigBarc}{pst-eucl}
\\ \hline
  \blue \dft : CodeFigAarc =true  & { \blue \dft : CodeFigBarc =true  } \\  \hline 
%& \dft : CodeFigAarc =true & CodeFigBarc =true 
%\\ \hline
\end{tabular}

%---------
\SbSbSSCT{Intersection de deux courbes}{Intersection of two  plots}


\psset{unit=.5cm} 
\begin{center}
\begin{tabular}{|c|} \hline
\begin{pspicture}*(-5,-2)(5,6.5)  
\psaxes[xticksize=-1 6,yticksize=-4 4,axesstyle=frame](0,0)(-4,-1)(4,6)
\def \F{x 2 exp neg 5 add }
\def \G{x 2 exp }
\psplot[linecolor=Vert]{-2.5}{2.5}{\F} 
\psplot[linecolor=cyan]{-2.5}{2.5}{\G} 
\pstInterFF[linecolor=red]{\F}{\G}{1}{A} 
\pstInterFF[linecolor=red]{\F}{\G}{-1}{B} 	
\end{pspicture} \\ \hline
\BS{}def {\Vert\BS{}F\AC{x 2 exp neg 5 add}}  \\
\BS{}def {\cyan\BS{}G\AC{x 2 exp}}\\
\BSS{pstInterFF}\AC{{\Vert\BS{}F}}\AC{{\cyan\BS{}G}}\AC{1}{\red\AC{A}} \BSI{pstInterFF}{pst-eucl} \\
\BSS{pstInterFF}\AC{{\Vert\BS{}F}}\AC{{\cyan\BS{}G}}\rnode{D}{\AC{-1}}{\red\AC{B}} \BSI{pstInterFF}{pst-eucl}
\\ \hline
\end{tabular}
\end{center}

\rnode{C}{\psframebox{\TFRGB{point d'intersection le plus proche de}{point of intersection nearest to x=-1}}} \ncdiagg{->}{C}{D}
\psset{unit=1cm}


%--------------------------------------------------
\SbSbSSCT{Intersection d'une droite et d'une courbe}{Intersection of a line and a curve}


\psset{unit=.5cm}  
\begin{center}
\begin{tabular}{|c|} \hline
\begin{pspicture}*(-5,-2)(5,6.5)  
\psaxes[xticksize=-1 6,yticksize=-4 4,axesstyle=frame](0,0)(-4,-1)(4,6)
\pstGeonode(-3,5){A}(3,2){B}
\pstLineAB[linecolor=cyan]{A}{B}	
\def \F{x 2 exp neg 5 add }
\psplot[linecolor=Vert]{-2.5}{2.5}{\F} 
\pstInterFL[linecolor=red]{\F}{A}{B}{1}{C}	
\end{pspicture} 
\\ \hline
{\cyan\BS{pstLineAB}\AC{A}\AC{B}} \\ 
\BS{def} {\Vert\BS{}F\AC{x 2 exp neg 5 add}} \\ 
\BS{pstInterFL}\AC{{\Vert\BS{}F}}{\cyan\AC{A}\AC{B}}\AC{1}{\red\AC{C}} \BSI{pstInterFL}{pst-eucl}	
\\ \hline
\end{tabular}
\end{center}

\psset{unit=1cm}
%----------------------------------------------------------------------
\SbSbSSCT{Intersection d'un cercle et d'une courbe}{Intersection of a circle and a curve}
%\verb \pstInterFC %[hpari]{hfi}{hOi}{hAi}{hx0 i}{hMi}


\psset{unit=.5cm}

\begin{center}
\begin{tabular}{|c|} \hline 
\begin{pspicture}*(-5,-5)(5.5,6.5)  
\psaxes[xticksize=-4 6,yticksize=-5 5,axesstyle=frame](0,0)(-5,-4)(5,6)
\pstGeonode[linecolor=cyan](0,0){A}(-3,2){B}
\pstCircleOA[linecolor=cyan]{A}{B}	 
	
\def \F{x 2 exp neg 5 add }
\psplot[linecolor=Vert]{-3}{3}{\F} 
%
\pstInterFC[linecolor=red]{\F}{A}{B}{-1}{C} 
\pstInterFC[linecolor=red]{\F}{A}{B}{1}{D}  
\pstInterFC[linecolor=red]{\F}{A}{B}{-3}{E}	
\end{pspicture}
\\ \hline
\BS{}def {\Vert  \BS{}F\AC{x 2 exp neg 5 add }} : fonction  $5 - x^2$\\
{\cyan\BS{}pstCircleOA\AC{A}\AC{B}}\\
\BSS{pstInterFC}\AC{{\Vert\BS{}F}}{\cyan\AC{A}\AC{B}}\AC{-1}{\red\AC{C}} \BSI{pstInterFC}{pst-eucl}\\
\BS{pstInterFC}\AC{{\Vert\BS{}F}}{\cyan\AC{A}\AC{B}}\AC{1}{\red\AC{D}}\\
\BS{pstInterFC}\AC{{\Vert\BS{}F}}{\cyan\AC{A}\AC{B}}\AC{3}{\red\AC{E}} \\
\BS{pstInterFC}{\blue\AC{{\Vert\BS{}F}}}{\cyan\AC{A}\AC{B}}{\red\AC{1}\AC{D}}\\
\BS{pstInterFC}{\blue\AC{{\Vert\BS{}F}}}{\cyan\AC{A}\AC{B}}{\red\AC{3}\AC{E}}
\\ \hline
\end{tabular}
\end{center}
\psset{unit=1cm}  

%%--------------------------------------------------------------

%\section[Les vecteurs]{Les vecteurs  \cite{pstricks-add}}
\SSCT{Les vecteurs }{Vectors}

\SbSSCT{Chaine de vecteurs}{Vectors chain }
\begin{center}
\begin{tabular}{|l|} \hline
\begin{pspicture}(-.5,-.5)(4.5,3.5)
\psaxes[xticksize=3,yticksize=4,axesstyle=frame](4,3)
\psStartPoint(1,1)
\psVector[linecolor=blue](1;30)
\psVector[linecolor=cyan](2;0)
\psVector[linecolor=red](1;120)
\psVector[linestyle=dashed](2;180)
 \end{pspicture}\\ \hline
\BSS{psStartPoint}(1,1) \BSI{psStartPoint}{pstricks-add} \\
\BSS{psVector}[linecolor=blue](1;30)       	\BSI{psVector}{pstricks-add}\\
\BS{psVector}[linecolor=cyan](2;0)\\
\BS{psVector}[linecolor=red](1;120)\\
\BS{psVector}[linestyle=dashed](2;180)\\ \hline
 \end{tabular}
\end{center}

\subsection{Options}

\begin{center}
\begin{tabular}{|l|} \hline
\psset{linecolor=blue}
\begin{pspicture}(-.5,-.5)(4.5,3.5)
\psaxes[xticksize=3,yticksize=4,axesstyle=frame](4,3)
 \psStartPoint(.5,.5)
\psVector[linecolor=blue,markAngle](2;30)
\psVector[linecolor=red,markAngle](2;150)
 \end{pspicture}\\ \hline
 \BS{}psStartPoint(.5,.5) \\
 \BS{}psVector[linecolor=blue,\RDD{markAngle}](2;30) \RDI{markAngle}{pstricks-add} \\
 \BS{}psVector[linecolor=red,{\red markAngle}](2;150)\\  \hline
 
 \end{tabular}
\end{center} 
 
 
 \begin{center}
 \begin{tabular}{|l|l|} \hline 
\begin{pspicture}(-.5,-.5)(5.5,4.5)
\psaxes[xticksize=4,yticksize=5,axesstyle=frame](5,4)
 \psStartPoint[A](1,1)
\psVector(1;30)
\psVector(2;0)
\psVector(2;120)

\uput[90](A1){A1}
\uput[-90](A2){A2}
\uput[90](A3){A3}
%\psline[linecolor=red]{->>}(A1)(A3)
 \end{pspicture}
&
\begin{pspicture}(-.5,-.5)(5.5,4.5)
\psaxes[xticksize=4,yticksize=5,axesstyle=frame](5,4)
 \psStartPoint[A](2,1)
\psVector[linecolor=blue](1;30)
\psVector[linecolor=blue](2;0)
 \psStartPoint[B](2,1)
\psVector[linecolor=cyan](2;120)

\psline[linecolor=red]{->>}(B1)(A2)
 \end{pspicture}\\  \hline
  \BS{}psStartPoint{\red[A]}(1,1) &  \BS{}psStartPoint{\red [A]}(2,1) \\
\BS{}uput[90](A1){\red \AC{A1}} & \BS{}psStartPoint{\red [B]}(2,1) \\
\BS{}uput[-90](A2){\red\AC{A2}} & \BS{}psline[linecolor=red]\AC{->>}{\red (B1)(A2)} \\ \hline
 \end{tabular}
\end{center} 


%%==============================================================
\newpage

%\section[Les diagrammes arborescents ]{Les diagrammes arborescents \cite{pst-tree}}
\SSCT{Les diagrammes arborescents}{Trees} 

\label{pst-tree}

\subsection{structure}



%-----------------------------------------------------------------------
\SbSSCT{les noeuds}{The nodes}

\begin{tabular}{|c|c|}
%\hline \multicolumn{2}{|c|}{ Sans contenu} \\ 
\hline  \textbf{ \TFRGB{sans astérisque}{without asterisk}} & \textbf{\TFRGB{avec astérisque}{with asterisk}} \\
\hline \pstree{\Tp}{\Tp \Tp \Tp \Tp} & \pstree{\Tp*}{\Tp* \Tp* \Tp* \Tp*} \\ 
\hline \multicolumn{2}{|c|}{ \BS{pstree}\AC{\BSS{Tp}*}\AC{\BS{}Tp* \BS{}Tp* \BS{}Tp* \BS{}Tp*} } \\
\hline \pstree{\Tc{3pt}}{\Tc{3pt} \Tc{3pt} \Tc{3pt} \Tc{3pt}} & \pstree{\Tc*{3pt}}{\Tc*{3pt} \Tc*{3pt} \Tc*{3pt} \Tc*{3pt}} \\ 
\hline \multicolumn{2}{|c|}{ \BS{pstree}\AC{\BSS{Tc}*}\AC{\BS{}Tc* \BS{}Tc* \BS{}Tc* \BS{Tc}=*} } \\ 
\hline \pstree{\TC}{\TC \TC \TC \TC} & \pstree{\TC*}{\TC* \TC* \TC* \TC*} \\ 
\hline \multicolumn{2}{|c|}{ \BS{}pstree\AC{\BSS{TC}*}\AC{\BS{}TC* \BS{}TC* \BS{}TC* \BS{}TC*}  \BSI{TC}{pst-tree} } \\ 
\hline \pstree{\Tf}{\Tf \Tf \Tf \Tf} & \pstree{\Tf*}{\Tf* \Tf* \Tf* \Tf*} \\ 
\hline \multicolumn{2}{|c|}{ \BS{}pstree\AC{\BSS{Tf}*}\AC{\BS{}Tf* \BS{}Tf* \BS{}Tf* \BS{}Tf*}  \BSI{TC}{pst-tree}} \\ 
\hline \pstree{\Tdot}{\Tdot \Tdot \Tdot \Tdot} & \pstree{\Tdot*}{\Tdot* \Tdot* \Tdot* \Tdot*} \\
\hline \multicolumn{2}{|c|}{ \BS{}pstree\AC{\BSS{Tdot}*}\AC{\BS{}Tdot* \BS{}Tdot* \BS{}Tdot* \BS{}Tdot*}  \BSI{Tdot}{pst-tree} } \\ 
\hline 
\end{tabular}


\begin{tabular}{|c|c|}
\hline  \textbf{ \TFRGB{sans astérisque}{without asterisk}} & \textbf{\TFRGB{avec astérisque}{with asterisk}}\\
\hline \pstree{\Tr{A1}}{\Tr{B1} \Tr{B2} \Tr{B3} \Tr{B4}} & \pstree{\Tr*{A1}}{\Tr*{B1} \Tr*{B2} \Tr*{B3} \Tr*{B4}} \\
\hline \multicolumn{2}{|c|}{ \BS{pstree}\AC{\BSS{Tr}\AC{A1}}\AC{\BS{Tr}\AC{B1} \BS{Tr}\AC{B2} \BS{Tr}\AC{B3} \BS{Tr}\AC{B4}}  \BSI{Tr}{pst-tree}  } \\ 
 
\hline \pstree{\TR{A1}}{\TR{B1} \TR{B2} \TR{B3} \TR{B4}} & \pstree{\TR*{A1}}{\TR*{B1} \TR*{B2} \TR*{B3} \TR*{B4}} \\ 
\hline \multicolumn{2}{|c|}{ \BS{pstree}\AC{\BSS{TR}\AC{A1}}\AC{\BS{TR}\AC{B1} \BS{TR}\AC{B2} \BS{TR}\AC{B3} \BS{TR}\AC{B4}}  \BSI{TR}{pst-tree}  }\\

\hline \pstree{\Tcircle{A1}}{\Tcircle{B1} \Tcircle{B2} \Tcircle{B3} \Tcircle{B4}} & \pstree{\Tcircle*{A1}}{\Tcircle*{B1} \Tcircle*{B2} \Tcircle*{B3} \Tcircle*{B4}} \\
\hline 
\multicolumn{2}{|c|}{ \BS{}pstree \AC{\BSS{Tcircle}\AC{A1}}\AC{\BS{Tcircle}\AC{B1} \BS{Tcircle}\AC{B2} \BS{Tcircle}\AC{B3} \BS{Tcircle}\AC{B4} }  \BSI{Tcircle}{pst-tree}  }\\
 

\hline \pstree{\TCircle{A1}}{\TCircle{B1} \TCircle{B2} \TCircle{B3} \TCircle{B4}} & \pstree{\TCircle*{A1}}{\TCircle*{B1} \TCircle*{B2} \TCircle*{B3} \TCircle*{B4}} \\
\hline 
\multicolumn{2}{|c|}{ \BS{}pstree \AC{\BSS{TCircle}\AC{A1}}\AC{\BS{TCircle}\AC{B1} \BS{TCircle}\AC{B2} \BS{TCircle}\AC{B3} \BS{TCircle}\AC{B4} }  \BSI{TCircle}{pst-tree}  }\\
 
\hline \pstree{\Toval{A1}}{\Toval{B1} \Toval{B2} \Toval{B3} \Toval{B4}} & \pstree{\Toval*{A1}}{\Toval*{B1} \Toval*{B2} \Toval*{B3} \Toval*{B4}} \\
\hline 
\multicolumn{2}{|c|}{ \BS{}pstree \AC{\BSS{Toval}\AC{A1}}\AC{\BS{Toval}\AC{B1} \BS{Toval}\AC{B2} \BS{Toval}\AC{B3} \BS{Toval}\AC{B4} }  \BSI{Toval}{pst-tree}  }\\


\hline \pstree{\Tdia{A1}}{\Tdia{B1} \Tdia{B2} \Tdia{B3} \Tdia{B4}} & \pstree{\Tdia*{A1}}{\Tdia*{B1} \Tdia*{B2} \Tdia*{B3} \Tdia*{B4}} \\
\hline 
\multicolumn{2}{|c|}{ \BS{}pstree \AC{\BSS{Tdia}\AC{A1}}\AC{\BS{Tdia}\AC{B1} \BS{Tdia}\AC{B2} \BS{Tdia}\AC{B3} \BS{Tdia}\AC{B4} }  \BSI{Tdia}{pst-tree} }\\

\hline \pstree{\Ttri{A1}}{\Ttri{B1} \Ttri{B2} \Ttri{B3} \Ttri{B4}} & \pstree{\Ttri*{A1}}{\Ttri*{B1} \Ttri*{B2} \Ttri*{B3} \Ttri*{B4}} \\
\hline 
\multicolumn{2}{|c|}{ \BS{}pstree \AC{\BSS{Ttri}\AC{A1}}\AC{\BS{Ttri}\AC{B1} \BS{Ttri}\AC{B2} \BS{Ttri}\AC{B3} \BS{Ttri}\AC{B4} }  \BSI{Ttri}{pst-tree}  }\\
\hline 
\end{tabular}


\begin{tabular}{|c|c|}
\hline  
\pstree{\Toval{A1}}{
\pstree{\Tfan}{\Toval{B1}}
\pstree{\Tfan[fansize=.5]}{\Toval{B2}}
\pstree{\Tfan}{\Toval{B3} \Toval{B4} }}
&   
\parbox[t]{8cm}{  
\BS{pstree}\AC{\BS{Toval}\AC{A1}}\AC{\%  \\
\BS{pstree}\AC{\BSS{Tfan}}\AC{\BS{Toval}\AC{B1}}  \BSI{Tfan}{pst-tree}  \\
\BS{pstree}\AC{\textbf{\BS{Tfan}[\RDD{fansize}=.5]}}\AC{\BS{Toval}\AC{B2}} \RDI{fansize}{pst-tree} \\
\BS{pstree}\AC{\textbf{\BS{Tfan}}}\AC{\BS{Toval}\AC{B3} \BS{Toval}\AC{B4} } \\
} 

\bigskip
{\blue \dft{} :   fansize= 1cm }\\
}
\\ \hline 
\end{tabular} 

%---------------------------------------------------
\subsection{Orientation}


\begin{tabular}{|c|c|c|c|}
\hline \multicolumn{4}{|c|}{  \BS{pstree}[\RDD{treemode}=R]\AC{\BS{Toval}\AC{A1}}\AC{\BS{Toval}\AC{B1} \BS{Toval}\AC{B2}}  \RDI{treemode}{pst-tree}  } \\
\hline 
 \pstree[treemode=R]{\Toval{A1}}{\Toval{B1} \Toval{B2}}
&  
 \pstree[treemode=D]{\Toval{A1}}{\Toval{B1} \Toval{B2} }
&  
\pstree[treemode=L]{\Toval{A1}}{\Toval{B1} \Toval{B2}}
&
\pstree[treemode=U]{\Toval{A1}}{\Toval{B1} \Toval{B2} }
\\ 
\hline treemode=R  & treemode=D & treemode=L & treemode=U \\ 
\hline 
\end{tabular} 

\bigskip
 \begin{tabular}{|c|c|}
\hline \multicolumn{2}{|c|}{\BS{pstree}[\RDD{treeflip}=true]\AC{ \BS{Toval}\AC{A1}} \{\BS{Toval}\AC{B1}  \RDI{treeflip}{pst-tree}  } \\
\multicolumn{2}{|c|}{  \BS{pstree}[treemode=R]\AC{\BS{Toval}\AC{B2}}\AC{\BS{Toval}\AC{C1}}\} } \\ 
\hline 
 \pstree[treeflip=true]{\Toval{A1}}{\Toval{B1} \pstree[treemode=R]{\Toval{B2}}{\Toval{C1}}}
 &  
 \pstree[treeflip=false]{\Toval{A1}}{\Toval{B1} \pstree[treemode=R]{\Toval{B2}}{\Toval{C1}}}
 \\ 
 \hline treeflip=true & treeflip=false {\blue (\dft)}\\ 
 \hline 
 \end{tabular} 
 
 %----------------------------------------------------------------------
\SbSSCT{Distance entre 2 noeuds de même niveau}{Distance between two nodes on the same level}

 
\begin{tabular}{|c|c|}    \hline 
\multicolumn{2}{|c|}{  \BS{pstree}[\RDD{treefit}=tight] \AC{ \BS{Toval}\AC{A1}} \{ \BS{Toval}\AC{B1} \RDI{treefit}{pst-tree}  }\\  
\multicolumn{2}{|c|}{  \BS{pstree}\AC{ \BS{Toval}\AC{B2}} \AC{ \BS{Toval}\AC{C1} \BS{Toval}\AC{C2} }\} } \\
\hline  
\pstree[treefit=tight]{\Toval{A1}} {\Toval{B1} \pstree{\Toval{B2}}{\Toval{C1} \Toval{C2}}}
&  
\pstree[treefit=loose]{\Toval{A1}}{\Toval{B1} \pstree{\Toval{B2}}{\Toval{C1} \Toval{C2}}}
\\ 
\hline {\red  treefit=tight} {\blue (\dft)} & {\red treefit=loose }  \\ 
\hline 
\end{tabular}

\bigskip
\begin{tabular}{|c|c|}
\hline 
\multicolumn{2}{|c|}{  \BS{pstree}[\RDD{treesep}=0cm]\AC{\BS{Toval}\AC{A1}}\AC{\BS{Toval}\AC{B1} \BS{Toval}\AC{B2}} \RDI{treesep}{pst-tree}  } \\ 
\hline  
\pstree[treesep=0cm]{\Toval{A1}}{\Toval{B1} \Toval{B2}}
&  
\pstree[treesep=2cm]{\Toval{A1}}{\Toval{B1} \Toval{B2}}
\\ 
\hline  {\red treesep}=0cm & {\red treesep}=2cm {\blue (\dft: treesep=0.75cm) }  \\ 
\hline 
\end{tabular} 

\bigskip
\begin{tabular}{|c|c|}    \hline 
\multicolumn{2}{|c|}{ \BS{pstree}[\RDD{treenodesize}=1]\AC{\BS{Toval}\AC{A1}} \AC{\BS{Toval}\AC{B1}  \BS{Toval}\AC{B2 B2 B2} \BS{Toval}\AC{B3 B3 B3} } \RDI{treenodesize}{pst-tree}  }\\  
\hline  
\pstree[treenodesize=1]{\Toval{A1}} {\Toval{B1} \Toval{B2 B2 B2} \Toval{B3 B3 B3} }
&  
\pstree[treenodesize=-1]{\Toval{A1}} {\Toval{B1} \Toval{B2 B2 B2} \Toval{B3 B3 B3} }
\\  \hline 
{\red treenodesize=1}
&
{\red treenodesize=-1}
 \\ \hline 
\end{tabular}

\bigskip
\begin{tabular}{|c|c|}    \hline 
\multicolumn{2}{|c|}{ \BS{pstree}\AC{\BS{Toval}\AC{A1}} \AC{\BS{Toval}\AC{B1} \BSS{tspace}\AC{-1cm} \BS{Toval}\AC{B2} \BS{Toval}\AC{B3} }  \BSI{Tspace}{pst-tree}  }\\  
\hline  
\pstree{\Toval{A1}} {\Toval{B1} \tspace{-1cm} \Toval{B2} \Toval{B3} }
&  
\pstree{\Toval{A1}}{\Toval{B1}\tspace{1cm} \Toval{B2} \Toval{B3} }
 \\ \hline 
{\red  \BS{tspace}\AC{-1cm}}
&
{\red  \BS{tspace}\AC{1cm}}
 \\ \hline 
\end{tabular}

\bigskip
\begin{tabular}{|c|c|} \hline 
\parbox[c]{4cm}{ 
\pstree[levelsep=1cm]{\Toval{A1}}{\Toval{B1} 
\pstree{\Toval{B2}}{\pstree[thistreefit=loose]{\Toval{C1}}{\red \Toval{D1} \pstree{\Toval{D2}}{\black \Toval{E1} \Toval{E2}}}\Toval{C2}}}
}
& 
\parbox[l]{8cm}{ 
\BS{pstree}[levelsep=1cm]\AC{\BS{Toval}\AC{A1}}\\
\AC{\BS{Toval}\AC{B1} \\
\BS{pstree}\AC{\BS{Toval}\AC{B2}}\AC{ \\
\BS{pstree}[\RDD{thistreefit}=loose]\AC{\BS{Toval}\AC{C1}}  \RDI{thistreefit}{pst-tree}  \\
\AC{\BS{Toval}\AC{D1} \\
\BS{pstree}\AC{\BS{Toval}\AC{D2}}\\
\AC{\BS{Toval}\AC{E1} \\
\BS{Toval}\AC{E2}}}\\
\BS{Toval}\AC{C2}}}
} 
\\ \hline 
\end{tabular}
\bigskip

\begin{tabular}{|c|c|} \hline 
\parbox[c]{4cm}{ 
\pstree[levelsep=1cm]{\Toval{A1}}{\Toval{B1} 
\pstree{\Toval{B2}}{ \pstree[thistreesep=0]{\Toval{C1}}{\red \Toval{D1} \pstree{\Toval{D2}}{\black \Toval{E1} \Toval{E2}}}\Toval{C2}}}
}
& 
\parbox[l]{8cm}{ 
\BS{pstree}[levelsep=1cm]\AC{\BS{Toval}\AC{A1}}\\
\AC{\BS{Toval}\AC{B1} \\
\BS{pstree}\AC{\BS{Toval}\AC{B2}}\AC{ \\
\BS{pstree}[\RDD{thistreesep}=0]\AC{\BS{Toval}\AC{C1}} \RDI{thistreesep}{pst-tree}  \\
\AC{\BS{Toval}\AC{D1} \\
\BS{pstree}\AC{\BS{Toval}\AC{D2}}\\
\AC{\BS{Toval}\AC{E1} \\
\BS{Toval}\AC{E2}}}\\
\BS{Toval}\AC{C2}}}
} 
\\ \hline 
\end{tabular}



%--------------------------------------------------------------------------
\SbSSCT{Distance entre noeuds successifs}{Distance between successive nodes}


\begin{tabular}{|c|c|c|c|}
\hline 
 \multicolumn{4}{|c|}{\BS{pstree}[\RDD{levelsep}=1cm] \AC{\BS{Toval}\AC{A1}} \AC{ \BS{Toval}\AC{B1} \BS{pstree}\AC{\BS{Toval}\AC{B2}} \AC{\BS{Toval}\AC{C1} \BS{Toval}\AC{C2}} } \RDI{levelsep}{pst-tree}  }\\  
\hline  
\pstree[levelsep=1cm]{\Toval{A1}} {\Toval{B1} \pstree{\Toval{B2}}{\Toval{C1} \Toval{C2}}}
& 
\pstree[levelsep=*1cm]{\Toval{A1}} {\Toval{B1} \pstree{\Toval{B2}}{\Toval{C1} \Toval{C2}}}
& 
\pstree[levelsep=.5cm]{\Toval{A1}}{\Toval{B1} \pstree{\Toval{B2}}{\Toval{C1} \Toval{C2}}}
& 
\pstree[levelsep=*.5cm]{\Toval{A1}}{\Toval{B1} \pstree{\Toval{B2}}{\Toval{C1} \Toval{C2}}}
\\ \hline 
\hline levelsep=1cm & levelsep=*1cm & levelsep=.5cm   &  levelsep=*.5cm   \\ \hline 
 \multicolumn{4}{|c|}{ \blue (\dft : 2cm)}
\\ \hline 
\end{tabular}

\bigskip
\begin{tabular}{|c|c|} \hline 
\parbox[c]{4cm}{ 
\pstree[levelsep=1cm]{\Toval{A1}}{\Toval{B1} 
\pstree{\Toval{B2}}{\pstree[levelsep=.5cm]{\Toval{C1}}{\red \Toval{D1} \pstree{\Toval{D2}}{\Toval{E1} \Toval{E2}}}\Toval{C2}}}
}
& 
\parbox[l]{8cm}{ 
\BS{pstree}[levelsep=1cm]\AC{\BS{Toval}\AC{A1}}\\
\AC{\BS{Toval}\AC{B1} \\
\BS{pstree}\AC{\BS{Toval}\AC{B2}}\AC{ \\
\BS{pstree}[ \RDD{levelsep}=0.5cm]\AC{\BS{Toval}\AC{C1}}  \RDI{levelsep}{pst-tree} \\
\AC{\BS{Toval}\AC{D1} \\
\BS{pstree}\AC{\BS{Toval}\AC{D2}}\\
\AC{\BS{Toval}\AC{E1} \\
\BS{Toval}\AC{E2}}}\\
\BS{Toval}\AC{C2}}}
 
} 
\\ \hline 
\end{tabular} 

\bigskip

\begin{tabular}{|c|c|} \hline 
\parbox[c]{4cm}{ 
\pstree[levelsep=1cm]{\Toval{A1}}{\Toval{B1} 
\pstree{\Toval{B2}}{\pstree[thislevelsep=.5cm]{\Toval{C1}}{\red  \Toval{D1} \pstree{\Toval{D2}}{\black \Toval{E1} \Toval{E2}}}\Toval{C2}}}
}
& 
\parbox[l]{8cm}{ 
\BS{pstree}[levelsep=1cm]\AC{\BS{Toval}\AC{A1}}\\
\AC{\BS{Toval}\AC{B1} \\
\BS{pstree}\AC{\BS{Toval}\AC{B2}}\AC{ \\
\BS{pstree}[\RDD{thislevelsep}=0.5cm]\AC{\BS{Toval}\AC{C1}}  \RDI{thislevelsep}{pst-tree}  \\
\AC{\BS{Toval}\AC{D1} \\
\BS{pstree}\AC{\BS{Toval}\AC{D2}}\\
\AC{\BS{Toval}\AC{E1} \\
\BS{Toval}\AC{E2}}}\\
\BS{Toval}\AC{C2}}}
 
} 
\\ \hline 
\end{tabular} 
%---------------------------------------------------------------------------------------------
\SbSSCT{Liaison des noeuds}{Connecting the nodes}



\bigskip
  \begin{tabular}{|c|c|l|}
  \hline 
  \parbox[l]{3.5cm}{ 
  \renewcommand{\psedge}{\ncdiag[angle=-90, armA=0, angleB=90,armB=1cm]}
  \pstree{\Toval{A1}} {\Toval{B1} \pstree{\Toval{B2}}{\Toval{C1} \Toval{C2}}}
  }
  &
    \parbox[l]{3.5cm}{   
  \def\psedge{\nccurve[angleA=-90,angleB=90, nodesepB=3pt]}
  
  \pstree{\Toval{A1}} {\Toval{B1} \pstree{\Toval{B2}}{\Toval{C1} \Toval{C2}}}
  }
  &
\parbox[l]{6cm}{
\emph {\TFRGB{Redéfinition du type de liaison  2 possibilités}{the 2 ways to set the  connection  type} :}
\bigskip

  \BSS{renewcommand}\AC{\BSS{psedge}} \AC{\BS{ncdiag}[angle=-90,armA=0,angleB=90,armB=1cm]}

\bigskip
  \BSS{def}\BSS{psedge}\AC{\BS{nccurve}[angleA=-90 ,angleB=90,nodesepB=3pt]}
}
  \\   \hline 
 \multicolumn{3}|c|{ {\blue  \dft : \BS{ncline}} \hspace{1cm} \TFRGB{Autres possibilités voir}{other connection see} page \pageref{node} } \\
% \multicolumn{3}{|c|}{ Autres possibilités voir page \pageref{node}} \\
  \hline 
  \end{tabular} 

\bigskip
\begin{tabular}{|c|c|}
\hline  
\pstree{\Toval{A1}} {\Toval[edge={\ncdiag[angleA=-90,angleB=90,armA=0,armB=1cm]}]{B1} \Toval[edge={\nccurve[angleA=-90,angleB=90]}]{B2}}
&  
 \parbox[t]{12cm}{ 
\smallskip
\BS{pstree}\AC{\BS{Toval}\AC{A1}} \\
\AC{\BS{Toval}[\RDD{edge}=\AC{\BS{ncdiag}[angleA=-90,angleB=90,armA=0,armB=1cm]}]\AC{B1}  \RDI{edge}{pst-tree}  \\ 
\BS{Toval}[{\red edge}=\AC{\BS{nccurve}[angleA=-90,angleB=90]}]\AC{B2}}
}
\\ 
\hline 
\end{tabular} 

\bigskip

\begin{tabular}{|c|c|}
\hline  
\pstree{\Toval{A1}} {\Toval[name=A]{B1} \Toval[name=B]{B2}}
\ncline[linestyle=dashed,linecolor=red]{A}{B}
& 
 \parbox[t]{12cm}{ 
\smallskip
\BS{pstree}\AC{\BS{Toval}\AC{A1}} \\
\AC{\BS{Toval}[\RDD{name}=A]\AC{B1}\BS{Toval[{\red name}=B]\AC{B2}}}  \RDI{name}{pst-tree} \\
\smallskip
\BS{ncline}[linestyle=dashed,linecolor=red]\AC{A}\AC{B} \\
}
\\ 
\hline 
\end{tabular} 


\SbSSCT{Etiquettes}{Labels}
\SbSbSSCT{Etiquettes sur les liaisons}{Labels on the connection}

\begin{tabular}{|c|c|c|c|}
\hline 
\multicolumn{4}{|c|}{ \BS{pstree}\AC{\BS{Toval}\AC{A1}} \AC{\BS{Toval}\AC{B1} {\red  \BS{tlput}\AC{\BS{}red l}}  \BS{Toval}\AC{B2}  } }\\  
\hline  
\pstree{\Toval{A1}} {\Toval{B1} \tlput{\red l} \Toval{B2}} 
& 
\pstree{\Toval{A1}} {\Toval{B1} \trput{\red r} \Toval{B2}} 
&  
\pstree{\Toval{A1}} {\Toval{B1} \taput{\red a} \Toval{B2}}
& 
\pstree{\Toval{A1}} {\Toval{B1} \tbput{\red b} \Toval{B2}}
 \\ 
   \hline \BSS{tlput}\AC{\BS{}red l}  \BSI{tlput}{pst-tree} & \BSS{trput}\AC{\BS{}red r}  \BSI{trput}{pst-tree}  & \BSS{taput}\AC{\BS{}red a}  \BSI{taput}{pst-tree}  &  \BSS{tbput}\AC{\BS{}red a}  \BSI{tbput}{pst-tree} \\ 
   \hline 
   \end{tabular} 

\bigskip
\begin{tabular}{|c|c|c|c|}
\hline 
 \multicolumn{4}{|c|}{ \BS{pstree}[{\red treemode=L}]\AC{\BS{Toval}\AC{A1}} \AC{\BS{Toval}\AC{B1} {\red  \BS{tlput}\AC{\BS{}red l}}  \BS{Toval}\AC{B2}  } }\\  
 \hline  
\pstree[treemode=L]{\Toval{A1}} {\Toval{B1} \tlput{\red l} \Toval{B2}}
&  
\pstree[treemode=L]{\Toval{A1}} {\Toval{B1} \trput{\red r} \Toval{B2}}
&  
\pstree[treemode=L]{\Toval{A1}} {\Toval{B1} \taput{\red a} \Toval{B2}}
&  
\pstree[treemode=L]{\Toval{A1}} {\Toval{B1} \tbput{\red b} \Toval{B2}}
\\ 
   \hline \BS{tlput}\AC{\BS{}red l} & \BS{trput}\AC{\BS{}red r}  & \BS{taput}\AC{\BS{}red a}  &  \BS{taput}\AC{\BS{}red a} \\ 
\hline 
\end{tabular} 
\bigskip

\begin{tabular}{|c|c|c|c|}
\hline 
\multicolumn{4}{|c|}{\BS{}psset\AC{tpos=.75}  \BS{pstree}\AC{\BS{Toval}\AC{A1}} \AC{\BS{Toval}\AC{B1} {\red  \BS{tlput}\AC{\BS{}red X}}  \BS{Toval}\AC{B2}  } }\\  
\hline 
\psset{tpos=0} 
\pstree{\Toval{A1}} {\Toval{B1} \tlput{\red X} \Toval{B2}} 
&
\psset{tpos=.75} 
\pstree{\Toval{A1}} {\Toval{B1} \tlput{\red X} \Toval{B2}} 
& 
\psset{tpos=1} 
\pstree{\Toval{A1}} {\Toval{B1} \tlput{\red X} \Toval{B2}}
& 
\pstree{\Toval{A1}} {\Toval{B1} \tlput{\red X} \Toval{B2}}
 \\ 
   \hline tpos=0 & tpos=.75 & tpos=1 & {\blue \dft }\\ 
   \hline 
   \end{tabular} 

%------------------------------------------------------------------------------------------------
\SbSbSSCT{Etiquettes sur les noeuds}{Labels on the nodes}

\begin{tabular}{|c|c|c|c|}
\hline  
 \multicolumn{4}{|c|}{ \BS{pstree}\AC{\BS{Toval}\AC{A1}} \AC{\BS{Toval}\AC{B1}  $\sim$[\RDD{tnpos}=l]\AC{\BS{}red l}}  \BS{Toval}\AC{B2}   \RDI{tnpos}{pst-tree} } \\  
 \hline  
\pstree{\Toval{A1}} {\Toval{B1}~[tnpos=l]{\red l} \Toval{B2}} 
& 
\pstree{\Toval{A1}} {\Toval{B1}~[tnpos=r]{\red r} \Toval{B2}} 
&  
\pstree{\Toval{A1}} {\Toval{B1}~[tnpos=a]{\red a} \Toval{B2}}
& 
\pstree{\Toval{A1}} {\Toval{B1}~[tnpos=b]{\red b}\Toval{B2}}
 \\ 
\hline tnpos=l & tnpos=r & tnpos=a & tnpos=b \\ 
\hline 
\end{tabular}

\bigskip
\begin{tabular}{|c|c|c|c|}
\hline 
 \multicolumn{4}{|c|}{ \BS{pstree}[{\red treemode=L}]\AC{\BS{Toval}\AC{A1}} \AC{\BS{Toval}\AC{B1} { $\sim$[{\red tnpos}=l]\AC{\BS{}red l}}  \BS{Toval}\AC{B2}  } }\\  
 \hline   
\pstree[treemode=L]{\Toval{A1}} {\Toval{B1} ~[tnpos=l]{\red l} \Toval{B2}}
&  
\pstree[treemode=L]{\Toval{A1}} {\Toval{B1}~[tnpos=r]{\red r} \Toval{B2}}
&  
\pstree[treemode=L]{\Toval{A1}} {\Toval{B1} ~[tnpos=a]{\red a} \Toval{B2}}
&  
\pstree[treemode=L]{\Toval{A1}} {\Toval{B1} ~[tnpos=b]{\red b} \Toval{B2}}
\\ 
\hline tnpos=l & tnpos=r & tnpos=a & tnpos=b \\ 
\hline 
\end{tabular}

\bigskip
\begin{tabular}{|c|c|c|c|}
\hline  
 \multicolumn{4}{|c|}{ \BS{pstree}\AC{\BS{Toval}\AC{A1}} \AC{\BS{Toval}\AC{B1} {$\sim$[tnpos=b,\RDD{tnsep}=1cm]\AC{\BS{}red  1cm}}  \BS{Toval}\AC{B2}  }  \RDI{tnsep}{pst-tree} }\\  
 \hline  
\pstree{\Toval{A1}} {\Toval{B1}~[tnpos=b,tnsep=1cm]{\red 1cm} \Toval{B2}} 
& 
\pstree{\Toval{A1}} {\Toval{B1}~[tnpos=b,tnsep=-1cm]{\red -1cm} \Toval{B2}} 
&  
\pstree{\Toval{A1}} {\Toval{B1}~[tnpos=b,tnsep=0cm]{\red 0cm} \Toval{B2}}
& 
\pstree{\Toval{A1}} {\Toval{B1}~[tnpos=b]{\blue\dft}\Toval{B2}}
 \\ 
\hline tnsep=1cm & tnsep=-1cm & tnsep=0cm & {\blue\dft }\\ 
\hline 
\end{tabular}


\bigskip
\begin{tabular}{|c|c|c|c|}
\hline  
 \multicolumn{4}{|c|}{ \BS{pstree}\AC{\BS{Toval}\AC{A1}} \AC{\BS{Toval}\AC{B1} {$\sim$[tnpos=b,\RDD{tnheight}=1cm]\AC{\BS{}red l}}  \BS{Toval}\AC{B2}  }  \RDI{tnheight}{pst-tree}  }\\  
 \hline  
\pstree{\Toval{A1}} {\Toval{B1}~[tnpos=b,tnheight=1cm]{\red 1cm} \Toval{B2}} 
& 
\pstree{\Toval{A1}} {\Toval{B1}~[tnpos=b,tnheight=-1cm]{\red -1cm} \Toval{B2}} 
&  
\pstree{\Toval{A1}} {\Toval{B1}~[tnpos=b,tnheight=0cm]{\red 0cm} \Toval{B2}}
& 
\pstree{\Toval{A1}} {\Toval{B1}~[tnpos=b]{\blue \dft}\Toval{B2}}
 \\ 
\hline tnheight=1cm & tnheight=-1cm & tnheight=0cm & {\blue \dft } \\ 
\hline 
\end{tabular}
\bigskip

\begin{tabular}{|c|c|c|c|}
\hline  
 \multicolumn{4}{|c|}{ \BS{pstree}\AC{\BS{Toval}\AC{A1}} \AC{\BS{Toval}\AC{B1} {$\sim$[tnpos=b,\RDD{tnyref}=1cm]\AC{\BS{}red l}}  \BS{Toval}\AC{B2}  }  \RDI{tnyref}{pst-tree}  }\\  
 \hline  
\pstree{\Toval{A1}} {\Toval{B1}~[tnpos=b,tnyref=1cm]{\red 1cm} \Toval{B2}} 
& 
\pstree{\Toval{A1}} {\Toval{B1}~[tnpos=b,tnyref=-1cm]{\red -1cm} \Toval{B2}} 
&  
\pstree{\Toval{A1}} {\Toval{B1}~[tnpos=b,tnyref=0cm]{\red 0cm} \Toval{B2}}
& 
\pstree{\Toval{A1}} {\Toval{B1}~[tnpos=b]{\blue \dft}\Toval{B2}}
 \\ 
\hline tnyref=1cm &tnyref=-1cm & tnyref=0cm & {\blue \dft }\\ 
\hline 
\end{tabular}


%------------------------------------------------------------------------------
\newpage
\SbSSCT{Showbbox}{Showbbox}

\begin{tabular}{|c|c|}
\hline 
 \multicolumn{2}{|l|}{   \BS{psset}\AC{\RDD{showbbox}=true}     }\\ 
  \multicolumn{2}{|l|}{ \BS{pstree}\AC{\BS{Toval}\AC{A1}} \AC{\BS{Toval}\AC{B1}  \BS{Toval}\AC{B2}}    \RDI{showbbox}{pst-tree}  }\\  
\hline 
&\\  
\psset{showbbox=true}
\pstree[treemode=L]{\Toval{A1}} {\Toval{B1}  \Toval{B2}}
&  
\psset{showbbox=false}
\pstree[treemode=L]{\Toval{A1}} {\Toval{B1} \Toval{B2}}
\\ 
&\\
\hline showbbox=true & showbbox=false \\ 
\hline 
\end{tabular} 

 
 \bigskip
 \begin{tabular}{|c|c|c|c|}
 \hline  
 \multicolumn{4}{|c|}{   \BS{psset}\AC{showbbox=true}   \BS{pstree}[bbl=1cm]\AC{\BS{Toval}\AC{A1}} \AC{\BS{Toval}\AC{B1}  \BS{Toval}\AC{B2}}    }\\  

 \psset{showbbox=true}
 \pstree[bbl=1cm]{\Toval{A1}} {\Toval{B1}  \Toval{B2}} 
 &  
 \psset{showbbox=true}
  \pstree[bbr=1cm]{\Toval{A1}} {\Toval{B1}  \Toval{B2}}  
 & 
 \psset{showbbox=true} 
 \pstree[bbh=1cm]{\Toval{A1}} {\Toval{B1}  \Toval{B2}}
 &
 \psset{showbbox=true}  
  \pstree[bbd=1cm]{\Toval{A1}} {\Toval{B1}  \Toval{B2}}
 \\ 
\RDD{bbl}=1cm \RDI{bbl}{pst-tree} &
 \RDD{bbr}=1cm  \RDI{bbr}{pst-tree}& 
 \RDD{bbh}=1cm  \RDI{bbh}{pst-tree}& 
 \RDD{bbd}=1cm  \RDI{bbd}{pst-tree}\\ 
 \hline 
 \end{tabular} 

   \bigskip 
 
  \begin{tabular}{|cccc|}
  \hline 
  \multicolumn{4}{|c|}{   \BS{psset}\AC{showbbox=true}   \BS{pstree}[\RDD{xbbl}=1cm]\AC{\BS{Toval}\AC{A1}} \AC{\BS{Toval}\AC{B1}  \BS{Toval}\AC{B2}}   \RDI{xbbl}{pst-tree}  }\\  

  &&&\\ 
  \psset{showbbox=true}
  \pstree[xbbl=1cm]{\Toval{A1}} {\Toval{B1}  \Toval{B2}} 
  &  
  \psset{showbbox=true}
   \pstree[xbbr=1cm]{\Toval{A1}} {\Toval{B1}  \Toval{B2}}  
  & 
  \psset{showbbox=true} 
  \pstree[xbbh=1cm]{\Toval{A1}} {\Toval{B1}  \Toval{B2}}
  & 
  \psset{showbbox=true} 
   \pstree[xbbd=1cm]{\Toval{A1}} {\Toval{B1}  \Toval{B2}}
  \\ 

 \RDD{xbbl}=1cm  \RDI{xbbl}{pst-tree}& 
 \RDD{xbbr}=1cm  \RDI{xbbr}{pst-tree}& 
 \RDD{xbbh}=1cm   \RDI{xbbh}{pst-tree}& 
 \RDD{xbbd}=1cm   \RDI{xbbd}{pst-tree}\\ 
\hline 
\end{tabular}  

\bigskip
\begin{tabular}{|c|c|c|}
\hline  
\pstree{\Tc{3pt}~[tnpos=a,tndepth=0pt,radius=4pt]{root}}{%
\TC*~[tnpos=l,tndepth=0pt]{$h$}
\TC*~[tnpos=r,tndepth=0pt]{$i$}}
&  
\pstree{\Tc{3pt}~[tnpos=a,tndepth=1cm,radius=4pt]{root}}{%
\TC*~[tnpos=l,tndepth=1cm]{$h$}
\TC*~[tnpos=r,tndepth=1cm]{$i$}}
&  
\pstree{\Tc{3pt}~[tnpos=a,tndepth=-1cm,radius=4pt]{root}}{%
\TC*~[tnpos=l,tndepth=-1cm]{$h$}
\TC*~[tnpos=r,tndepth=-1cm]{$i$}}
\\ 
\hline \RDD{tndepth}=0pt  \RDI{tndepth}{pst-tree} & \RDD{tndepth}=1cm  \RDI{tndepth}{pst-tree} &  
\RDD{tndepth}=-1cm  \RDI{tndepth}{pst-tree}\\ 
\hline 
\end{tabular} 

%----------------------------------------------------------
\SbSSCT{skiplevel}{skiplevel}
\begin{tabular}{|c|c|} \hline 
\parbox[l]{6cm}{ 
\pstree[levelsep=1cm,treesep=0cm]{\Toval{A1}}{%
\skiplevel{\Toval{\red B1}}
\pstree{\Toval{B2}}{%
\Toval{C1}
\skiplevels{2}
\pstree{\Toval{\red C2}}{\Toval{D1} \Toval{D2}}
\Toval{\red C3}
\endskiplevels
\pstree{\Toval{C4}}{\Toval{E1} \Toval{E2}}}}
}

& 
\parbox[l]{8cm}{ 
\BS{pstree}[levelsep=1cm,treesep=0cm]\AC{\BS{Toval}\AC{A1}}\AC{\% \\
\\
\BSS{skiplevel} \AC{\BS{Toval}{\red \AC{\BS{}red  B1}} }  \BSI{skiplevel}{pst-tree} \\
 
\BS{pstree}\AC{\BS{Toval}\AC{B2}}\AC{\% \\
\BS{Toval}\AC{C1} \\
\\
\BSS{skiplevels}\AC{2} \BSI{skiplevels}{pst-tree}\\
\BS{pstree}\AC{\BS{Toval}{\red \AC{\BS{}red C2}}}\AC{\BS{Toval}\AC{D1} \BS{Toval}\AC{D2}} \\
\BS{Toval}{\red \AC{\BS{}red C3}} \\
\BSS{endskiplevels} \BSI{endskiplevel}{pst-tree} \\
\\
\BS{pstree}\AC{\BS{Toval}\AC{C4}}\AC{\BS{Toval}\AC{E1} \BS{Toval}\AC{E2}}}} \\
} 
\\ \hline 
\end{tabular}



%%==============================================================
\newpage

%\section[Les animations ]{ Les animations \cite{anim}}
\SSCT{Les animations}{Animations}
\label{anim}
\SbSSCT{Animation à partir de fichiers d'image }{Animation from picture files}

%\begin{verbatim}\animategraphics[<options>]{<frame rate>}{<file basename>}{<first>}{<last>} \end{verbatim}


\begin{tabular}{|c|c|} \hline 
\TFRGB{première image}{first frame} & \TFRGB{second et dernière image}{second and last frame}
\\ \hline
\includegraphics{XXX1.ps}
&  
\includegraphics{XXX2.ps}
\\ 
\hline \BS{includegraphics}\AC{XXX1.ps} &  \BS{includegraphics}\AC{XXX2.ps}\\ 
\hline 
\end{tabular} 



%
\begin{minipage}{7cm}
\animategraphics[controls,loop,autoplay]{4}{XXX}{1}{2}  
\end{minipage}\hfill
\begin{minipage}{7cm}
\begin{tabular}{|l@{:}l|}
\hline \BSS{animategraphics} &  \\ 
\hline [ controls, & \TFRGB{boutons de contrôle}{Inserts control buttons} \\ 
\hline loop &  \TFRGB{en boucle}{animation restarts automatically}\\ 
\hline autoplay ] & \TFRGB{auto demarrage}{Start animation automatically } \\ 
\hline \AC{4} &  \TFRGB{4 fois par seconde}{4 frame per second}\\ 
\hline \AC{XXX} & \TFRGB{base du nom fichier}{file base name} \\ 
\hline \AC{1} & \TFRGB{numero de debut}{number of the first frame} \\ 
\hline \AC{2} & \TFRGB{numeo de fin}{number of the last frame} \\ 
\hline 
\end{tabular} 

\end{minipage}

%--------------------------------------------------------------
\subsection{Animateinline}


\begin{minipage}{7cm}
\begin{center}
\begin{animateinline}[ controls,loop,autoplay]{5}%
\begin{pspicture}(6,6)
\psdiamond*[gangle=45](3,3)(2,.5)
\psdiamond*[gangle=135](3,3)(2,.5)
\end{pspicture}
\newframe%
\begin{pspicture}(6,6)
\psdiamond*[gangle=0](3,3)(2,.5)
\psdiamond*[gangle=90](3,3)(2,.5)
\end{pspicture}
\end{animateinline}% 
\end{center}
\end{minipage}\hfill
\begin{minipage}{7cm}
\BSS{begin\AC{animateinline}}[controls,loop,autoplay]\AC{5} \\

\emph{\% \TFRGB{première image}{first frame }}\\
\BS{begin}\AC{pspicture}(6,6) \\
\BS{psdiamond*}[gangle=45](3,3)(2,.5) \\
\BS{psdiamond*}[gangle=135](3,3)(2,.5) \\
\BS{end}\AC{pspicture} \\

\emph{\% \TFRGB{deuxième}{second frame }}\\
\BSS{newframe} \\
\BS{begin}\AC{pspicture}(6,6) \\
\BS{psdiamond*}[gangle=0](3,3)(2,.5) \\
\BS{psdiamond*}[gangle=90](3,3)(2,.5) \\
\BS{end}\AC{pspicture} \\
\\
\BSS{end\AC{animateinline}} \\
\end{minipage}
%
%%------------------------------------------------------------------------
%\newpage
\subsection{Multiframe}

\begin{minipage}{7cm}
\begin{center}
\begin{animateinline}[poster=first, controls, palindrome,autoplay]{12}%
\multiframe{29}{iAngle=80+10,Rdim=2.0+-0.2}{\begin{pspicture}(6,6)
 \psdiamond*[gangle=\iAngle](3,3)(\Rdim,.5)
 \rput(1,1){\iAngle}
 \rput(5,1){\Rdim}
\end{pspicture} }%
\end{animateinline}%
\end{center}
\end{minipage}\hfill
\begin{minipage}{7cm}

\BS{begin}\AC{animateinline}[poster=first,controls, palindrome]\AC{12} \\
\BSS{multiframe}\AC{29}\AC{{\red iAngle}=80+10, {\red Rdim}=2.0+-0.2}\AC{ \\
\BS{begin}\AC{pspicture}(6,6) \\
\BS{psdiamond*}[gangle={\red \BS{iAngle}}](3,3)({\red \BS{Rdim}},.5)\\
\BS{rput}(1,1)\AC{{\red \BS{iAngle}}} \\
\BS{rput}(5,1)\AC{{\red \BS{Rdim}}} \\
\BS{end}\AC{pspicture} } \\
\BS{end}\AC{animateinline}%


\end{minipage}
\bigskip

\TFRGB{L'initiale de la variable définit son type}{The first  letter of the variable name determines his type }

\begin{tabular}{|c|l|}
\hline  entier &  initiale : i ou I \\ 
\hline  réelles &  initiale : n, N, r ou R \\ 
\hline  longueurs & initiale : d ou D \\ 
\hline 
\end{tabular} 
%%---------------------------------------------------------------------------------
\newpage
\subsection{Timeline}

\begin{minipage}{6.5cm}
\begin{animateinline}[controls,autoplay,timeline=xxx.txt]{5}%
\begin{pspicture}(6,6)
\pscircle[fillcolor=yellow,fillstyle=solid](3,3){2.5}
\end{pspicture} 
\newframe%
\begin{pspicture}(6,6)
\pscircle[linecolor=red,fillcolor=green,fillstyle=solid](3,3){2.5}
\end{pspicture} 
\newframe
\multiframe{10}{iAngle=60+10}{\begin{pspicture}(6,6)
\psdiamond*[gangle=\iAngle](3,3)(2,.5)
\end{pspicture} }%
\end{animateinline}%

\end{minipage}\hfill
\begin{minipage}{8cm}

\BS{begin}\AC{animateinline}

[controls,autoplay,\RDD{timeline}=xxx.txt]\AC{5} \\

\emph{\% \TFRGB{1 image de fond}{}first background image (image \No 0)} \\
\BS{begin}\AC{pspicture}(6,6) \\
\BS{pscircle}[fillcolor=yellow,fillstyle=solid](3,3)\AC{2.5} \\
\BS{end}\AC{pspicture}  \\

\textbf{\BS{newframe}} \emph{\%  \TFRGB{2 page de fond}{second background image } (image \No 1)} \\
\BS{begin}\AC{pspicture}(6,6) \\ 
\BS{pscircle}[linecolor=red,fillcolor=green,fillstyle=solid](3,3)\AC{2.5} \\
\BS{end}\AC{pspicture} \\

\textbf{\BS{newframe}} \emph{\% \TFRGB{animation}{animation frames} (images \No 2 - 11) } \\
\BS{multiframe}\AC{10}\AC{iAngle=60+10}\AC{ \\
\BS{begin}\AC{pspicture}(6,6) \\
\BS{psdiamond*}[gangle=\BS{iAngle}](3,3)(2,.5) \\
\BS{end}\AC{pspicture}  } \\
\BS{end}\AC{animateinline} \\
\end{minipage}

\bigskip
\SbSbSSCT{Création du fichier pour timeline}{ Creation of the file for timeline}

\TFRGB{Pour créer le fichier xxx.txt , en insérant le code suivant avant}{to create the file xxx.txt, insert the following code before } \BS{begin\AC{document}}
\bigskip

\begin{minipage}{7cm}
\BSS{begin\AC{filecontents}}\AC{xxx.txt} \\
::\Rnode{A}{0x0,8} \\
::2 \\
::7 \\
::3 \\
::6 \\
::\Rnode{B}{c,1x3,5} \\ 
::4 \\
::11 \\
::5 \\
::7 \\
::9 \\
\BSS{end\AC{filecontents}}
\end{minipage}\hfill
\begin{minipage}{7cm}
\rnode{AA}{0x0 : image \No 0 \TFRGB{sert de fond tout le temps}{= background image for all frame}} \\

 \vspace{.5cm}
\Rnode{BB}{c : \TFRGB{efface les images précédentes}{clear the background image} } \\

 \vspace{.5cm}
\Rnode{CC}{1x3 : image \No 1  \TFRGB{sert de fond 3 fois}{= background image for 3  frames}} \\

 \vspace{.5cm}
 \TFRGB{Ordre de passage des images}{ Order of frames} : 8,2,7,3,6,5,4,11,5,7,9
\end{minipage}
 \ncline[linecolor=blue]{A}{AA}  \ncline[linecolor=blue]{B}{BB}  \ncline[linecolor=blue]{B}{CC}
 
\SbSbSSCT{option pour le fichier xxx.txt}{option for the file xxx.txt}
 
\begin{tabular}{|c|l|}  \hline  
* : : 3
&
\TFRGB{pause à l'image}{pause at frame } \No 3    
\\   \hline  
: 10 : 3
&  
\TFRGB{vitesse 10 par seconde à l'image}{10 frames per second at frame} \No  3
\\   \hline  
: : 3 : code 
&  
\TFRGB{code java possible à l'image}{java code at frame} \No  3
\\   \hline 
 \end{tabular}  
 
\newpage
\SbSSCT{Animation d'un graphe}{Graph animation}
\readdata{\dat}{mesdata.dat}

\begin{animateinline}[poster=last,controls]{5}%
\multiframe{70}{ifin=10+10}{
\begin{psgraph}[axesstyle=frame,xticksize=0 4cm,yticksize=0 9cm,subticks=0,Dx=100,Dy=.02](0,0)(750,.12){9cm}{4cm} 
\listplot[xEnd=\ifin,linecolor=blue,linewidth=5pt]{\dat}
\end{psgraph} }%
\end{animateinline}

\bigskip
\noindent

\BS{readdata}\AC{\BS{dat}}\AC{mesdata.dat} \\
\BS{begin}\AC{animateinline}[poster=last,controls]\AC{5} \\
\BS{multiframe}\AC{70}\AC{{\red ifin}=10+10}\AC{ \\
\BS{begin}\AC{psgraph}[axesstyle=frame,xticksize=0 4cm,yticksize=0 9cm,subticks=0,Dx=100,Dy=.02](0,0)(750,.12)\AC{9cm}\AC{4cm} \\
\BS{listplot}[xEnd={\red \BS{ifin}},linecolor=blue,linewidth=5pt]\AC{\BS{dat}} \\
\BS{end}\AC{psgraph} } \\
\BS{end}\AC{animateinline}








%
%\newpage
%\subsection{Command options}
%
%
%\begin{tabular}{|c|c|} \hline  
%poster[=first | none | last] &   \\  \hline  
%every=<num>					& Build animation from every <num>th frame only \\
% autopause 					& \TFRGB{animation en pause à la fermeture de la page}{Pause animation when the page is closed }  \\  \hline  
% autoplay 					&  Start animation after the page has opened.\\  \hline  
% autoresume 				&  Resume previously paused animation \\  \hline  
% loop 						&  The animation restarts immediately after reaching the end. \\  \hline  
% palindrome 				&  The animation continuously plays forwards and backwards. \\  \hline  
% step 						&  Step through the animation one frame at a time per mouse-click. \\  \hline  
% width=<width> 				&  \TFRGB{redimensionne la largeur }{Resize the width of the animation widget} \\  \hline  
% height=<height> 			&  \TFRGB{redimensionne la hauteur}{Resize the height of the animation widget} \\  \hline  
% depth=<depth> 				&  Resize the animation widget \\  \hline  
% scale=<factor> 			&  Scales the animation widget \\  \hline  
% bb=<llx> <lly> <urx> <ury> &   \\  \hline  
% viewport=<llx> <lly> <urx> <ury> &   \\  \hline  
% trim=<left> <bottom> <right> <top> &   \\  \hline  
% controls 					&  Inserts control buttons \\  \hline  
% buttonsize=<size> 			&  Changes the control button \\  \hline  
% buttonbg=<colour> 			&  strokes on transparent background \\  \hline  
% buttonfg=<colour> 			& specifies the stroke colour  \\  \hline  
% draft 						& animation is not embedded.Instead, a box with the exact
% dimensions  \\  \hline  
% final 						&   \\  \hline  
% useocg 					&  alternative animation method based on Optional Content Groups \\  \hline  
% measure 					&  Measures the frame rate during one cycle of the animation. \\  \hline  
%\end{tabular}




\newpage

%\section[Créer un dessin en 3D]{Créer un dessin en 3D  \cite{pst-3dplot}}
\SSCT{Créer un dessin en 3D}{3D drawing}

Utilisation du module \textbf{pst-3dplot} \label{pst-3dplot}

\psset{unit=.5cm}
\psset{fillcolor=yellow,fillstyle=solid}
\SbSSCT{Les axes en 3 D}{3 D axis}
\label{3dplot}
%\subsubsection{pstThreeDCoor}

\begin{center}
\begin{tabular}{|c|c|c|}
\hline
 \multicolumn{2}{|c|}{ \BSS{pstThreeDCoor}  \BSI{pstThreeDCoor}{pst-3dplot} } \\ \hline
  

\begin{pspicture}(-3,-2.5)(3,4.25)
\pstThreeDCoor[drawing=true]
\end{pspicture}
&  
\begin{pspicture}(-3,-2.5)(3,4.25)
\pstThreeDCoor[drawing=false]
\end{pspicture}
\\ 
\hline  \RDD{drawing}=true (\dft)  \RDI{drawing}{pst-3dplot} &  \RDD{drawing}=false\\ 
\hline 
\end{tabular} 
\end{center}



\bigskip

\begin{tabular}{|c|c|c|}
\hline
 \multicolumn{3}{|c|}{ \BS{pstThreeDCoor}[\RDD{xMax}=2,\RDD{yMax}=2,\RDD{zMax}=2] } \\ \hline
\begin{pspicture}(-3,-2.5)(3,4.5)
\pstThreeDCoor[xMax=2,yMax=2,zMax=2]
\end{pspicture}
&  
\begin{pspicture}(-3,-2.5)(3,4.5)
\pstThreeDCoor[xMin=-2,yMin=-2,zMin=-2]
\end{pspicture}
&  
\begin{pspicture}(-3,-2.5)(3,4.5)
\pstThreeDCoor[nameX=A,nameY=B,nameZ=C]
\end{pspicture}

\\ 
\hline  \RDD{xMax}=2,\RDD{yMax}=2,\RDD{zMax}=2 &  \RDD{xMin}=-2,\RDD{yMin}=-2,\RDD{zMin}=-2 & \RDD{nameX}=A,\RDD{nameY}=B,\RDD{nameZ}=C  
\\
\RDI{xMax}{pst-3dplot} \RDI{yMax}{pst-3dplot} \RDI{zMax}{pst-3dplot} 
&
\RDI{xMin}{pst-3dplot} \RDI{yMin}{pst-3dplot} \RDI{zMin}{pst-3dplot}
&
\RDI{nameX}{pst-3dplot} \RDI{nameY}{pst-3dplot} \RDI{nameZ}{pst-3dplot}
\\ \hline 
{\blue \dft : xMax=yMax=zMax=4} & {\blue  \dft : xMin=yMin=zMin=-1}& \\ \hline 
\end{tabular} 

\subsubsection{Option spotX}
\begin{center}
\psframebox[fillstyle=none]{\BS{pstThreeDCoor}[\RDD{spotX}=60,\RDD{spotY}=60,\RDD{spotZ}=60] \RDI{nspotX}{pst-3dplot} \RDI{nspotY}{pst-3dplot} \RDI{nspotZ}{pst-3dplot}}


\begin{animateinline}[controls,autoplay,loop]{2}%
\multiframe{20}{iAngle=60+10}{
\begin{pspicture}(-4,-4.5)(4,4.25)
\psframe(-4,-4.5)(4,4.25)
\pstThreeDCoor[linecolor=blue,linestyle=dotted]
\pstThreeDCoor[spotX=\iAngle,spotY=\iAngle,spotZ=\iAngle]
\rput(-.5,-4){ spotX : } \rput(1,-4){ \iAngle}
\end{pspicture}}
\end{animateinline} 
\end{center}


\newpage
%-----------------------------------------------------------------------------------
\SbSbSSCT{Orientation des axes }{Axis orientation}
\begin{center}
\psframebox[fillstyle=none]{\BS{pstThreeDCoor}[linecolor=blue,linestyle=dotted] }
\end{center}


\begin{minipage}{7cm}

\begin{center}
\psframebox[fillstyle=none]{\BS{pstThreeDCoor}[\RDD{Alpha}=30] \RDI{Alpha}{pst-3dplot}}
\begin{animateinline}[controls,autoplay,loop,,palindrome]{2}%
\multiframe{20}{iAngle=60+10}{
\begin{pspicture}(-4,-4.5)(4,4.25)
\psframe(-4,-4.5)(4,4.25)
\pstThreeDCoor[linecolor=blue,linestyle=dotted]
\pstThreeDCoor[Alpha=\iAngle]
\rput(-.8,-4){ Alpha : } \rput(1,-4){ \iAngle}
\end{pspicture}}
\end{animateinline}
\end{center}
\end{minipage}\hfill
\begin{minipage}{7cm}
\begin{center}
\psframebox[fillstyle=none]{\BS{pstThreeDCoor}[\RDD{Beta}=30] \RDI{Beta}{pst-3dplot} }
\begin{animateinline}[controls,autoplay,loop,palindrome]{2}%
\multiframe{20}{iAngle=60+10}{
\begin{pspicture}(-4,-4.5)(4,4.25)
\psframe(-4,-4.5)(4,4.25)
\pstThreeDCoor[linecolor=blue,linestyle=dotted]
\pstThreeDCoor[Beta=\iAngle]
\rput(-.5,-4){ Beta : } \rput(1,-4){ \iAngle}
\end{pspicture}}
\end{animateinline}
\end{center}
\end{minipage}

\bigskip

\begin{tabular}{|c|c|c|}
\hline
 \multicolumn{3}{|c|}{ \BS{pstThreeDCoor}[linestyle=dotted,linecolor=blue] \BS{pstThreeDCoor}[\RDD{RotX}=30] \RDI{RotX}{pst-3dplot}} \\ \hline
\begin{pspicture}(-4,-2.5)(4,5.25)
\pstThreeDCoor[linestyle=dotted,linecolor=blue]
\pstThreeDCoor[RotX=30]
\end{pspicture}
&  
\begin{pspicture}(-4,-2.5)(4,5.25)
\pstThreeDCoor[linestyle=dotted,linecolor=blue]
\pstThreeDCoor[RotY=-30]
\end{pspicture}
&  
\begin{pspicture}(-4,-2.5)(4,5.25)
\pstThreeDCoor[linestyle=dotted,linecolor=blue]
\pstThreeDCoor[RotZ=30]
\end{pspicture}

\\ 
\hline  \RDD{RotX}=30 &  \RDD{RotY}=-30 \RDI{RotY}{pst-3dplot} & \RDD{RotZ}=30 \RDI{RotZ}{pst-3dplot} \\ \hline 
{\blue \dft : RotX=0} & {\blue \dft : RotY=0} & {\blue \dft : RotZ=0} \\ \hline 
\end{tabular} 
 
 \bigskip
%\subsubsection{RotAngle xRotVec}
\psframebox[fillstyle=none]{\parbox[l]{12cm}{
\BS{pstThreeDCoor}[\RDD{RotSequence}=quaternion,\RDD{RotAngle}=10, \RDI{RotSequence}{pst-3dplot} \RDI{RotAngle}{pst-3dplot} \\ \RDD{xRotVec}=3,\RDD{yRotVec}=0,\RDD{zRotVec}=3, \RDI{xRotVec}{pst-3dplot}  \RDI{yRotVec}{pst-3dplot} \RDI{zRotVec}{pst-3dplot}  \\
 xMin=0,xMax=3, yMin=0,yMax=3, zMin=0,zMax=3]\\
\\
\BS{pstThreeDLine}[linecolor=blue, linewidth=2pt, arrows=->](0,0,0)(3,0,3)}}

\begin{center}
\begin{animateinline}[controls,autoplay,loop]{2}%
\multiframe{10}{iA=0+10}{
\begin{pspicture}(-3,-3)(3,3.5)
\psframe(-3,-3)(3,3.5)
\pstThreeDCoor[linecolor=red, RotSequence=quaternion, RotAngle=\iA, xRotVec=3,yRotVec=0,zRotVec=3,
 xMin=0,xMax=3, yMin=0,yMax=3, zMin=0,zMax=3]
\pstThreeDLine[linecolor=blue, linewidth=2pt, arrows=->](0,0,0)(3,0,3)
\uput[0](-2.28,-2.5){RotAngle : \iA}
\end{pspicture}}
\end{animateinline}
\end{center}
%--------------------------------------------------------------------------
\subsubsection{Option Ticks}

\begin{tabular}{|c|c|c|}
\hline
 \multicolumn{3}{|c|}{ \BS{pstThreeDCoor}[\RDD{IIIDticks},\RDD{IIIDticksize}=.5pt] \RDI{IIIDticks}{pst-3dplot}  \RDI{IIIDticksize}{pst-3dplot}  \RDI{IIIDticks}{pst-3dplot} } \\ \hline
\begin{pspicture}(-4,-2.5)(4,4.25)

\pstThreeDCoor[IIIDticks,IIIDticksize=.5pt]
\end{pspicture}
&  
\begin{pspicture}(-4,-2.5)(4,4.25)
\pstThreeDCoor[IIIDticks,IIIDlabels]
\end{pspicture}
&  
\begin{pspicture}(-4,-2.5)(4,4.25)
\pstThreeDCoor[IIIDticks,IIIDlabels,Dx=.5,Dy=.5,Dz=.5]
\end{pspicture}
\\ 
\hline  \RDD{IIIDticks},\RDD{IIIDticksize}= .5pt 
\RDI{IIIDticks}{pst-3dplot} \RDI{IIIDticksize}{pst-3dplot} 
&  
\RDD{IIIDticks},\RDD{IIIDlabels} 
 \RDI{IIIDlabels}{pst-3dplot}
& 
\RDD{Dx}=.5,\RDD{Dy}=.5,\RDD{Dz}=.5 
 \RDI{Dx}{pst-3dplot} \RDI{Dy}{pst-3dplot} \RDI{Dz}{pst-3dplot}
\\ \hline
{\blue \dft : IIIDticksize=0.1} & {\blue  \dft : IIIDlabels=false}  & {\blue  \dft : Dx=Dy=Dz=1} \\ \hline 
\end{tabular} 


\bigskip

\begin{tabular}{|c|} \hline  
 \BS{pstThreeDCoor}[IIIDticks,IIIDlabels, yMin=-3,\RDD{IIIDOffset}=\AC{(1,-2,1)}]
  \RDI{IIIDOffset}{pst-3dplot}
\\ \hline  
 \begin{pspicture}(-4,-1.25)(1,4)
 \pstThreeDCoor[IIIDticks,IIIDlabels, yMin=-3,IIIDOffset={(1,-2,1)}]
 \end{pspicture}
\\ \hline 
\end{tabular} 


\subsubsection{Option pstThreeDPlaneGrid}
 
 \begin{tabular}{|c|c|c|} \hline 
  \multicolumn{3}{|c|}{  \BSS{pstThreeDPlaneGrid}[planeGrid=xz](0,0)(3,3) \BSI{pstThreeDPlaneGrid}{pst-3dplot} } \\ \hline
  \begin{pspicture}(-4,-3.5)(5,5)
  \pstThreeDCoor %[xMin=0,yMin=0,zMin=0,linewidth=2pt]
  \pstThreeDPlaneGrid[linecolor=blue](0,0)(3,3)
  \end{pspicture}
 &
  \begin{pspicture}(-4,-3.5)(5,5)
  \pstThreeDCoor %[xMin=0,yMin=0,zMin=0,linewidth=2pt]
  \pstThreeDPlaneGrid[linecolor=blue,planeGrid=xz](0,0)(3,3)
  \end{pspicture}
 &
  \begin{pspicture}(-4,-3.5)(5,5)
  \pstThreeDCoor %[xMin=0,yMin=0,zMin=0,linewidth=2pt]
 \pstThreeDPlaneGrid[linecolor=blue,planeGrid=yz](0,0)(3,3)
  \end{pspicture}
  \\ \hline 
 \dft ( \RDD{planeGrid}=xy )  \RDI{planeGrid}{pst-3dplot} 
 & \RDD{planeGrid}=xz  & \RDD{planeGrid}=yz \\ \hline 
  
 \end{tabular} 
 \bigskip
  
  \begin{tabular}{|c|c|c|} \hline  
  \multicolumn{2}{|c|}{  BS{pstThreeDPlaneGrid}[planeGrid=xz, \RDD{planeGridOffset}=-1](0,0)(3,3) \RDI{planeGridOffset}{pst-3dplot}} \\ \hline
  \begin{pspicture}(-4,-3.5)(5,5)
  \pstThreeDCoor %[xMin=0,yMin=0,zMin=0,linewidth=2pt]
  \pstThreeDPlaneGrid[linecolor=blue,planeGrid=xz,planeGridOffset=-1](0,0)(3,3)
  \end{pspicture} 
 &  
  \begin{pspicture}(-4,-3.5)(5,5)
  \pstThreeDCoor %[xMin=0,yMin=0,zMin=0,linewidth=2pt]
 \pstThreeDPlaneGrid[linecolor=blue,planeGrid=xz,planeGridOffset=1](0,0)(3,3)
  \end{pspicture}
 \\ \hline 
 \RDD{planeGridOffset}=-1 & \RDD{planeGridOffset}=1 
 \\  \hline 
 \end{tabular} 
 
 \bigskip
 \begin{tabular}{|c|c|c|} \hline  
 \multicolumn{3}{|c|}{  \BS{pstThreeDPlaneGrid}[planeGrid=xy,\RDD{subticks}=2](0,0)(3,3) } \\ \hline
 \begin{pspicture}(-4,-3.5)(5,5)
 \pstThreeDCoor %[xMin=0,yMin=0,zMin=0,linewidth=2pt]
 \pstThreeDPlaneGrid[linecolor=blue,planeGrid=xy,subticks=2](0,0)(3,3)
 \end{pspicture} 
 &  
 \begin{pspicture}(-4,-3.5)(5,5)
 \pstThreeDCoor %[xMin=0,yMin=0,zMin=0,linewidth=2pt]
 \pstThreeDPlaneGrid[linecolor=blue,planeGrid=xz,subticks=5](0,0)(3,3)
  \end{pspicture}
 &  
 \begin{pspicture}(-4,-3.5)(5,5)
 \pstThreeDCoor %[xMin=0,yMin=0,zMin=0,linewidth=2pt]
 \pstThreeDPlaneGrid[linecolor=blue,planeGrid=yz,subticks=20](0,0)(3,3)
  \end{pspicture}
 \\ \hline 
 planeGrid=xy &  planeGrid=xz & planeGrid=yz \\
 subticks=2 & subticks=5 & subticks=20 \\ 
 \hline 
 \end{tabular} 
  \bigskip
  
 \begin{tabular}{|c|c|c|} \hline  
 \multicolumn{3}{|c|}{  \BS{pstThreeDPlaneGrid}[planeGrid=xy,\RDD{xsubticks}=5](0,0)(3,3) } \\ \hline
 \begin{pspicture}(-4,-3.5)(5,5)
 \pstThreeDCoor %[xMin=0,yMin=0,zMin=0,linewidth=2pt]
 \pstThreeDPlaneGrid[linecolor=blue,planeGrid=xy,xsubticks=5](0,0)(3,3)
 \end{pspicture} 
 &  
 \begin{pspicture}(-4,-3.5)(5,5)
 \pstThreeDCoor %[xMin=0,yMin=0,zMin=0,linewidth=2pt]
 \pstThreeDPlaneGrid[linecolor=blue,planeGrid=xz,xsubticks=5](0,0)(3,3)
  \end{pspicture}
 &  
 \begin{pspicture}(-4,-3.5)(5,5)
 \pstThreeDCoor %[xMin=0,yMin=0,zMin=0,linewidth=2pt]
 \pstThreeDPlaneGrid[linecolor=blue,planeGrid=yz,xsubticks=5](0,0)(3,3)
  \end{pspicture}
 \\ \hline 
 planeGrid=xy &  planeGrid=xz & planeGrid=yz \\
\RDD{xsubticks}=5 & \RDD{xsubticks}=5 & \RDD{xsubticks}=5 \\ 
 \hline 
 \end{tabular} 
 \bigskip

\begin{tabular}{|c|c|c|} \hline  
\multicolumn{3}{|c|}{  \BS{pstThreeDPlaneGrid}[planeGrid=xy,\RDD{ysubticks}=2](0,0)(3,3) } \\ \hline
\begin{pspicture}(-4,-3.5)(5,5)
\pstThreeDCoor %[xMin=0,yMin=0,zMin=0,linewidth=2pt]
\pstThreeDPlaneGrid[linecolor=blue,planeGrid=xy,ysubticks=5](0,0)(3,3)
\end{pspicture} 
&  
\begin{pspicture}(-4,-3.5)(5,5)
\pstThreeDCoor %[xMin=0,yMin=0,zMin=0,linewidth=2pt]
\pstThreeDPlaneGrid[linecolor=blue,planeGrid=xz,ysubticks=5](0,0)(3,3)
 \end{pspicture}
&  
\begin{pspicture}(-4,-3.5)(5,5)
\pstThreeDCoor %[xMin=0,yMin=0,zMin=0,linewidth=2pt]
\pstThreeDPlaneGrid[linecolor=blue,planeGrid=yz,ysubticks=5](0,0)(3,3)
\end{pspicture}
\\ \hline 
planeGrid=xy &  planeGrid=xz & planeGrid=yz \\
\RDD{ysubticks}=5 & \RDD{ysubticks}=5 & \RDD{ysubticks}=5 \\ 
\hline 
\end{tabular} 


%------------------------------------------------------------------
 
 \subsubsection{Option coorType}
 \begin{tabular}{|c|c|c|c|c|}
 \hline
  \multicolumn{3}{|c|}{ \BS{pstThreeDCoor}[\RDD{coorType}=0] } \\ \hline
  \begin{pspicture}(-4,-2.5)(5,5)
  \pstThreeDCoor[coorType=0]
  \end{pspicture}
  & 
  \begin{pspicture}(-4,-3)(5,5)
 \pstThreeDCoor[coorType=1]
 \end{pspicture}
 &  
  \begin{pspicture}(-4,-2.5)(5,5)
 \pstThreeDCoor[coorType=2]
 \end{pspicture}
 \\  \hline 
  \RDD{coorType}=0 &  \RDD{coorType}=1 &  \RDD{coorType}=2 \\ \hline
 \begin{pspicture}(-4,-2.5)(5,5)
 \pstThreeDCoor[coorType=3]
 \end{pspicture}
 &  
  \begin{pspicture}(-4,-2.5)(5,5)
 \pstThreeDCoor[coorType=4]
 \end{pspicture}
 &
 
 \\  \hline 
 \RDD{coorType}=3 & \RDD{coorType}=4  & \\ \hline 
 \end{tabular} 
 
 
 


%\subsection{Les objets en 3D}
\SSCT{Les objets en 3D}{3D Objects}

%\subsection{Objets en 3D}
\psset{dotstyle=*,dotscale=2,drawCoor=true,fillstyle=none}

\begin{tabular}{|c|c|c|}
\hline  
\begin{pspicture}(-2,-2)(2,2)
\pstThreeDCoor[xMin=-2,xMax=2,yMin=-2,yMax=2,zMin=-2,zMax=2]
\pstThreeDDot(-1,1,1)
\pstThreeDDot(1.5,-1,3)
\end{pspicture}
&  
\begin{pspicture}(-2,-2.25)(2,2.25)
\pstThreeDCoor[xMin=-2,xMax=2,yMin=-2,yMax=2,zMin=-2,zMax=2]
\pstThreeDDot(-1,1,1)
\pstThreeDDot(1.5,-1,3)
\pstThreeDLine[linewidth=3pt,linecolor=blue](-1,1,1)(1.5,-1,3)
\end{pspicture}
&
\begin{pspicture}(-3,-4.25)(3,3.25)
\pstThreeDCoor[xMin=-4,xMax=4,yMin=-3,yMax=5,zMin=-4,zMax=3]
\pstThreeDTriangle[drawCoor=true,linecolor=blue,linewidth=2pt,fillstyle=none](3,1,2)(1,4,-1)(-2,2,0)
\end{pspicture}
\\ \hline 
\BSS{pstThreeDDot}(-1,1,1) & \BSS{pstThreeDLine} & \BSS{pstThreeDTriangle} \\
\BSS{pstThreeDDot}(1.5,-1,3) &(-1,1,1)(1.5,-1,-1) & (3,1,2)(1,4,-1)(-2,2,0) \\ 
\hline 
%------------------------------------------------------------------------------
\begin{pspicture}(-3,-2)(4,5)
\pstThreeDCoor[xMin=-3,xMax=3,yMin=-1,yMax=4,zMin=-1,zMax=3]
\pstThreeDSquare[linewidth=2pt,drawCoor=false](-2,2,3)(3,0,0)(0,1,-1)
\pstThreeDDot[linecolor=red](-2,2,3)
\end{pspicture}
& 
\begin{pspicture}(-2,-2.25)(2,5)
\pstThreeDCoor[xMax=3,yMax=2,zMax=4,yMin=-1.5]
%\pstThreeDDot[linecolor=red,drawCoor=true](1,0.5,0.5)
\pstThreeDEllipse(2,-1,2)(-1,1,0)(1,0,-1)
\pstThreeDDot[linecolor=red,drawCoor=true](2,-1,2)
\pstThreeDLine[linewidth=1pt,linecolor=red,linestyle=dashed](2,-1,2)(1,0,2)
\pstThreeDLine[linewidth=.5pt,linecolor=red,linestyle=dashed](2,-1,2)(1,-1,2)
\pstThreeDLine[linewidth=.5pt,linecolor=red,linestyle=dashed](1,-1,2)(1,0,2)

\pstThreeDLine[linewidth=1pt,linecolor=cyan,linestyle=dashed](2,-1,2)(3,-1,1)
\pstThreeDLine[linewidth=.5pt,linecolor=cyan,linestyle=dashed](2,-1,2)(3,-1,2)
\pstThreeDLine[linewidth=.5pt,linecolor=cyan,linestyle=dashed](3,-1,2)(3,-1,1)
\end{pspicture}
& 
\begin{pspicture}(-2,-2.25)(2,5)
\pstThreeDCoor[xMax=3,yMax=2,zMax=4,yMin=-1.5]
%\pstThreeDDot[linecolor=red,drawCoor=true](1,0.5,0.5)
\pstThreeDCircle(2,-1,2)(-1,1,0)(1,0,-1)
\pstThreeDDot[linecolor=red,drawCoor=true](2,-1,2)
\pstThreeDLine[linewidth=1pt,linecolor=red,linestyle=dashed](2,-1,2)(1,0,2)
\pstThreeDLine[linewidth=.5pt,linecolor=red,linestyle=dashed](2,-1,2)(1,-1,2)
\pstThreeDLine[linewidth=.5pt,linecolor=red,linestyle=dashed](1,-1,2)(1,0,2)

\pstThreeDLine[linewidth=1pt,linecolor=cyan,linestyle=dashed](2,-1,2)(3,-1,1)
\pstThreeDLine[linewidth=.5pt,linecolor=cyan,linestyle=dashed](2,-1,2)(3,-1,2)
\pstThreeDLine[linewidth=.5pt,linecolor=cyan,linestyle=dashed](3,-1,2)(3,-1,1)
\end{pspicture}
\\ \hline 
\BSS{pstThreeDSquare} 	& \BSS{pstThreeDEllipse} 	&  \\ 
\Rnode*[fillcolor=green]{E}{(-2,2,3)}  \Rnode*[fillcolor=yellow]{F}{(3,0,0)(0,1,-1)} 
& 
\Rnode*[fillcolor=green]{A}{(2,-1,2)} \Rnode*[fillcolor=yellow]{B}{(-1,1,0)(1,0,-1)} 
&   
\BSS{pstThreeDCircle} 
\Rnode*[fillcolor=green]{C}{(1,-1,2)} \Rnode*[fillcolor=yellow]{D}{\AC{2}}
\\
& & \\
 \rnode{EE}{position} \hspace{.5cm} \rnode{FF}{2 \TFRGB{vecteurs}{vectors}} 
& \rnode{AA}{\TFRGB{centre}{center}} \hspace{.5cm} \rnode{BB}{2 \TFRGB{vecteurs}{vectors}}    
& \rnode{CC}{\TFRGB{centre}{center}} \hspace{.5cm} \rnode{DD}{2 \TFRGB{vecteurs}{vectors}}
\\ \hline  
\ncline[linecolor=blue]{A}{AA}  \ncline[linecolor=blue]{B}{BB}
\ncline[linecolor=blue]{C}{CC}  \ncline[linecolor=blue]{D}{DD}
\ncline[linecolor=blue]{E}{EE}  \ncline[linecolor=blue]{F}{FF}
%----------------------------------------------------------------------------------------
\begin{pspicture}(-2.5,-2)(2,6)
\pstThreeDCoor[zMax=6]
\pstIIIDCylinder{1.5}{4}
\end{pspicture} 
&  
\begin{pspicture}(-2,-2)(2,6)
\pstThreeDCoor[zMax=6]
\psCylinder{1.5}{4}
\end{pspicture}
& 
\begin{pspicture}(-2,-1)(2,6)
\pstThreeDCoor[xMax=2,yMax=2,zMin=0,zMax=6,IIIDticks]%
\pstParaboloid{4}{2}
\end{pspicture}
\\ \hline 
\BSS{pstIIIDCylinder}\AC{1.5}\AC{4} & \BSS{psCylinder}\AC{1.5}\AC{4} & \BSS{pstParaboloid}\AC{4}\AC{2} 
\\  \hline 
%-------------------------------------------------------  
\begin{pspicture}(-2,-1.25)(3,4.25)
\pstThreeDCoor[xMin=-2,xMax=1,yMin=-1,yMax=3,zMin=-1,zMax=4]
\pstThreeDBox[linewidth=2pt,drawCoor=false](-1,2,2)(0,0,2)(2,0,0)(0,1,0)
\pstThreeDDot[linecolor=red,drawCoor=true](-1,2,2)
\end{pspicture}
&
\begin{pspicture}(-4,-2.25)(6,4.25)
 \pstThreeDCoor[xMin=-3,xMax=1,yMin=-1,yMax=2,zMin=-1,zMax=4]
 \psBox(-1,1,2){-3}{1}{2}
 \pstThreeDDot[linecolor=red,drawCoor=true](-1,1,2)
 \end{pspicture}
 &
\begin{pspicture}(-4,-2.25)(2,4.25)
 \pstThreeDCoor[xMin=-3,yMax=2]
\pstThreeDSphere(1,-1,2){2}
 \pstThreeDDot[dotstyle=x,linecolor=red,drawCoor=true](1,-1,2)
\end{pspicture} 
 
\\ \hline 
\BSS{pstThreeDBox} 
& 
&  \\
 \Rnode*[fillcolor=green]{A}{(-1,1,2)} \Rnode*[fillcolor=yellow]{B}{(0,0,2)(2,0,0)(0,1,0)} 
 &
\BSS{psBox} \Rnode*[fillcolor=green]{C}{(-1,1,2)} \Rnode*[fillcolor=yellow]{D}{\AC{-3}\AC{1}\AC{2}} 
&\BSS{pstThreeDSphere} 
\Rnode*[fillcolor=green]{E}{(1,-1,2)} \Rnode*[fillcolor=yellow]{F}{\AC{2}} 
\\
 & &  \\
\rnode{AA}{position} \hspace{.5cm} \rnode{BB}{\TFRGB{vecteurs en}{vectors}  X Y Z}   
 
  &
\rnode{CC}{position} \hspace{.5cm} \rnode{DD}{\TFRGB{vecteurs en}{vectors} X Y Z}
&
\rnode{EE}{\TFRGB{centre}{center}} \hspace{.5cm} \rnode{FF}{\TFRGB{rayon}{radius}} 
 \ncline[linecolor=blue]{A}{AA}  \ncline[linecolor=blue]{B}{BB}    
 \ncline[linecolor=blue]{E}{EE}  \ncline[linecolor=blue]{F}{FF} 
 \ncline[linecolor=blue]{C}{CC}  \ncline[linecolor=blue]{D}{DD} \\ \hline 
\end{tabular}


%---------------------------------------------------------
\SbSbSSCT{Portion d ellipse ou de cercle}{Portion of ellipse or circle}

\begin{tabular}{|c|c|c|} \hline 
  \multicolumn{3}{|c|}{ \BS{pstThreeDEllipse}[\RDD{beginAngle}=60](2,-1,2)(-1,1,0)(1,0,-1) } \\ \hline
\begin{pspicture}(-4,-2.25)(2,5)
\pstThreeDCoor[xMax=2,yMax=2,zMax=2]
\pstThreeDEllipse[beginAngle=60](2,-1,2)(-1,1,0)(1,0,-1)
\end{pspicture}
&  
\begin{pspicture}(-4,-2.25)(2,5)
\pstThreeDCoor[xMax=2,yMax=2,zMax=2]
\pstThreeDEllipse[endAngle=300](2,-1,2)(-1,1,0)(1,0,-1)
\end{pspicture}
&  
\begin{pspicture}(-4,-2.25)(2,5)
\pstThreeDCoor[xMax=2,yMax=2,zMax=2]
\pstThreeDEllipse[beginAngle=60,endAngle=300](2,-1,2)(-1,1,0)(1,0,-1)
\end{pspicture}
\\ \hline
 \RDD{beginAngle}=60 & \RDD{endAngle}=300 &  \RDD{beginAngle}=60\\ 
{\blue \dft : beginAngle=0} &{\blue \dft : endAngle=360} &\RDD{endAngle}=300  \\
\hline 
\end{tabular} 
\bigskip

\begin{tabular}{|c|c|c|} \hline 
 \multicolumn{3}{|c|}{ \BS{pstThreeDCircle}[\RDD{endAngle}=300](2,-1,2)(-1,1,0)(1,0,-1) } \\ \hline
\begin{pspicture}(-4,-2.25)(2,5)
\pstThreeDCoor[xMax=2,yMax=2,zMax=2]
\pstThreeDCircle[beginAngle=60](2,-1,2)(-1,1,0)(1,0,-1)
\end{pspicture}
&  
\begin{pspicture}(-4,-2.25)(2,5)
\pstThreeDCoor[xMax=2,yMax=2,zMax=2]
\pstThreeDCircle[endAngle=300](2,-1,2)(-1,1,0)(1,0,-1)
\end{pspicture}
&  
\begin{pspicture}(-4,-2.25)(2,5)
\pstThreeDCoor[xMax=2,yMax=2,zMax=2]
\pstThreeDCircle[beginAngle=60,endAngle=300](2,-1,2)(-1,1,0)(1,0,-1)
\end{pspicture}
\\ \hline
\RDD{beginAngle}=60 & \RDD{endAngle}=300 &  \RDD{beginAngle}=60\\ 
{\blue \dft : beginAngle=0  } &{\blue \dft : endAngle=360} &  \RDD{endAngle}=300 \\
\hline 
\end{tabular}

%---------------------------------------------------
\subsubsection{\emph{increment}}

\begin{tabular}{|c|c|c|c|} \hline 
 \multicolumn{2}{|c|}{ \TFRGB{incrément angulaire}{ angle step} } &  \multicolumn{2}{|c|}{\TFRGB{incrément vertical}{ height  step} } \\ \hline
\begin{pspicture}(-3.5,-2)(3.5,6)
\pstThreeDCoor %[zMax=1]
\psCylinder[increment=45]{1.5}{5}
\end{pspicture}
 & 
\begin{pspicture}(-3.5,-2)(3.5,6)
\pstThreeDCoor 
\psCylinder[increment=20]{1.5}{5}
\end{pspicture}
 & 
\begin{pspicture}(-3.5,-2)(3.5,6)
\pstThreeDCoor
\psCylinder[Hincrement=1]{1.5}{5}
\end{pspicture} 
&
\begin{pspicture}(-3.5,-2)(3.5,6)
\pstThreeDCoor
\psCylinder[Hincrement=.1]{1.5}{5}
\end{pspicture}
 \\ \hline   
\RDD{increment}=45 & \RDD{increment}=20 & \RDD{Hincrement}=1& \RDD{Hincrement}=.1 \\  \hline 
 \multicolumn{2}{|c|}{\blue \dft : increment=.1}  &  \multicolumn{2}{|c|}{\blue \dft : Hincrement=0.5} \\  \hline 
\end{tabular} 
\bigskip

\begin{tabular}{|c|c|c|c|} \hline  
 \multicolumn{2}{|c|}{ \BS{pstThreeDSphere}[\RDD{increment}=3](1,-1,2)\AC{2} } 
 &  \multicolumn{2}{|c|}{ \BS{pstParaboloid}[\RDD{increment}=3](4)\AC{2} } 
 \\ \hline
\begin{pspicture}(-4,-1)(1,4.25)
\pstThreeDSphere[increment=3](1,-1,2){2}
\end{pspicture}
&  
\begin{pspicture}(-4,-1)(1,4.25)
\pstThreeDSphere[increment=20](1,-1,2){2}
\end{pspicture}
&
\begin{pspicture}(-3,-1)(3,5)
\pstParaboloid[increment=3,showInside=false]{4}{2}
\end{pspicture}
&  
\begin{pspicture}(-3,-1)(3,5)
\pstParaboloid[increment=20,showInside=false]{4}{2}
\end{pspicture}
\\ \hline 
increment=3 & increment=20 & increment=3 & increment=20
\\ \hline 
 \multicolumn{4}{|c|}{\blue \dft : increment = 10 }
\\ \hline 
\end{tabular}  

%---------------------------------------------------------------------------
\subsubsection{showInside}


\begin{tabular}{|c|c|c|}
\hline  
\begin{pspicture}(-4,-2.25)(6,4.25)
%\pstThreeDCoor[zMax=6]
 \psBox[showInside=false]{-3}{1}{2}
\end{pspicture} 
&
\begin{pspicture}(-4,-2.25)(6,4.25)
%\pstThreeDCoor[xMax=2,yMax=2,zMin=0,zMax=6,IIIDticks]%
\pstParaboloid[showInside=true]{3}{2}
\end{pspicture}

\\ \hline 
 \BS{psBox}[\RDD{showInside}=false]\AC{-3}\AC{1}\AC{2}  &  \BS{pstParaboloid}[\RDD{showInside}=true]\AC{3}\AC{2}
 \\  \hline 
\begin{pspicture}(-4,-2.25)(6,4.25)
%\pstThreeDCoor[zMax=6]
 \psBox[showInside=true]{-3}{1}{2}
\end{pspicture}
&
\begin{pspicture}(-4,-2.25)(6,4.25)
\pstParaboloid[showInside=false]{3}{2}% Hoehe 5 und Radius 1
\end{pspicture}
\\  \hline 
 \BS{psBox}[\RDD{showInside}=true]\AC{-3}\AC{1}\AC{2}  &  \BS{pstParaboloid}[\RDD{showInside}=false]\AC{3}\AC{2}
 \\  \hline 
\end{tabular} 
%------------------------------------------------------------

\subsubsection{SegmentColor}

\begin{tabular}{|c|c|} \hline  
 \multicolumn{2}{|l|}{\BS{pstParaboloid}[showInside=false, \RDD{SegmentColor}=\AC{[cmyk]\AC{0 0 1 0}}]\AC{4}\AC{5}} \\
 \multicolumn{2}{|l|}{\BS{pstThreeDSphere}[\RDD{SegmentColor}=\AC{[cmyk]\AC{0,1,0,0}}](1,-1,2)\AC{2}}
\\ \hline  

\begin{pspicture}(-4,-1)(4,4.25)
\pstParaboloid[showInside=false, SegmentColor={[cmyk]{0 0 1 0}}]{3}{3}%
\end{pspicture}
&
\begin{pspicture}(-4,-1)(1,4.25)
\pstThreeDSphere[SegmentColor={[cmyk]{0,1,0,0}}](1,-1,2){2}
\end{pspicture}
\\  \hline 
\end{tabular} 




%-----------------------------------------------------------------------
\SbSSCT{Placer des objets en 3D}{How to place objects in 3D picture}
\subsubsection{pstThreeDPut}


\psset{fillcolor=yellow,fillstyle=solid}
\begin{tabular}{|c|c|} \hline  
\begin{pspicture}[shift=*](-4,-2.5)(4,4.25)
\pstThreeDDot[drawCoor=true](2,3,2)
\pstThreeDCoor[linecolor=blue]
\pstThreeDPut(2,3,2){\DFR}
\end{pspicture}
& 
\parbox{8cm}{ 
 \BSS{pstThreeDPu}t(2,3,2)\AC{\BS{DFR}} \\
 
\BS{pstThreeDDot}[drawCoor=true](2,3,2)}
\\ \hline  
\end{tabular} 


\bigskip
\begin{tabular}{|c|c|c|c|c|}
\hline
 \multicolumn{5}{|c|}{  \BS{pstThreeDPut}[\RDD{pOrigin}=lb](2,3,2)\AC{\BS{psframebox}\AC{ texte}} } \\ \hline
\begin{pspicture}(-3,-2.5)(3,4.25)
\pstThreeDDot[linecolor=red,drawCoor=true](2,3,2)
\pstThreeDCoor[linecolor=blue]
\pstThreeDPut[pOrigin=lt](2,3,2){\psframebox{texte}}
\end{pspicture}
&
\begin{pspicture}(-3,-2.5)(3,4.25)
\pstThreeDDot[linecolor=red,drawCoor=true](2,3,2)
\pstThreeDCoor[linecolor=blue]
\pstThreeDPut[pOrigin=lB](2,3,2){\psframebox{texte}}
\end{pspicture}
&
\begin{pspicture}(-3,-2.5)(3,4.25)
\pstThreeDDot[linecolor=red,drawCoor=true](2,3,2)
\pstThreeDCoor[linecolor=blue]
\pstThreeDPut[pOrigin=lb](2,3,2){\psframebox{texte}}

\end{pspicture}
&
\begin{pspicture}(-3,-2.5)(3,4.25)
\pstThreeDDot[linecolor=red,drawCoor=true](2,3,2)
\pstThreeDCoor[linecolor=blue]
\pstThreeDPut[pOrigin=t](2,3,2){\psframebox{texte}}
\end{pspicture}
&
\begin{pspicture}(-3,-2.5)(3,4.25)
\pstThreeDDot[linecolor=red,drawCoor=true](2,3,2)
\pstThreeDCoor[linecolor=blue]
\pstThreeDPut[pOrigin=c](2,3,2){\psframebox{texte}}
\end{pspicture}
\\ \hline
pOrigin=lt & pOrigin=lB & pOrigin=lb &  pOrigin=t & pOrigin=c \\ \hline
\begin{pspicture}(-3,-2.5)(3,4.25)
 \pstThreeDCoor[linecolor=blue]
 \pstThreeDPut[pOrigin=B](2,3,2){\psframebox{texte}}
\pstThreeDDot[linecolor=red,drawCoor=true](2,3,2)
\end{pspicture}
&
\begin{pspicture}(-3,-2.5)(3,4.25)

 \pstThreeDCoor[linecolor=blue]
 \pstThreeDPut[pOrigin=b](2,3,2){\psframebox{texte}}
\pstThreeDDot[linecolor=red,drawCoor=true](2,3,2)
\end{pspicture}
&
\begin{pspicture}(-3,-2.5)(3,4.25)
\pstThreeDDot[linecolor=red,drawCoor=true](2,3,2)
 \pstThreeDCoor[linecolor=blue]
 \pstThreeDPut[pOrigin=rt](2,3,2){\psframebox{texte}}
\end{pspicture}
&
\begin{pspicture}(-3,-2.5)(3,4.25)
\pstThreeDDot[linecolor=red,drawCoor=true](2,3,2)
 \pstThreeDCoor[linecolor=blue]
 \pstThreeDPut[pOrigin=rB](2,3,2){\psframebox{texte}}
\end{pspicture}
&
\begin{pspicture}(-3,-2.5)(3,4.25)
\pstThreeDDot[linecolor=red,drawCoor=true](2,3,2)
 \pstThreeDCoor[linecolor=blue]
 \pstThreeDPut[pOrigin=rb](2,3,2){\psframebox{texte}}

\end{pspicture}
 \\ \hline 
 pOrigin=B & pOrigin=b & pOrigin=rt &  pOrigin=rB & pOrigin=rb \\ \hline
\end{tabular} 



%-------------------------------------------------------------------------
\subsubsection{\BS{pstPlanePut}}

\begin{tabular}{|c|c|c|}
\hline  
\multicolumn{3}{|l|}{   \BSS{pstPlanePut}[\RDD{plane}=xy](0,0,3)\AC{\BS{DFR}}}\\
\multicolumn{3}{|l|}{   \BSS{pstPlanePut}[\RDD{plane}=xy](0,0,-3)\AC{\BS{psframebox}\AC{ texte}} }  
\\ \hline
\begin{pspicture}(-4,-4)(3.5,4.5)
\pstThreeDDot[linecolor=red](0,0,-2)
\pstThreeDDot[linecolor=red](0,0,2)
\pstThreeDCoor[xMin=-4,yMin=-4,zMin=-4]
\pstPlanePut[plane=xy](0,0,-2){\psframebox{texte}}
\pstPlanePut[plane=xy](0,0,2){\DFR}
\end{pspicture}
&
\begin{pspicture}(-4,-4)(3.5,4.5)
\pstThreeDDot[linecolor=red](0,0,-2)
\pstThreeDDot[linecolor=red](0,0,2) 
\pstThreeDCoor[xMin=-4,yMin=-4,zMin=-4]
\pstPlanePut[plane=yz](0,0,-2){\psframebox{texte}}
\pstPlanePut[plane=yz](0,0,2){\DFR }
\end{pspicture}
&
\begin{pspicture}(-4,-4)(3.5,4.5)
\pstThreeDDot[linecolor=red](0,0,-2)
\pstThreeDDot[linecolor=red](0,0,2) 
\pstThreeDCoor[xMin=-4,yMin=-4,zMin=-4]
\pstPlanePut[plane=xz](0,0,-2){\psframebox{texte}}
\pstPlanePut[plane=xz](0,0,2){\DFR}
\end{pspicture}
\\ \hline 
\RDD{plane}=xy & \RDD{plane}=yz &  \RDD{plane}=xz  \\ \hline 
 
\end{tabular}

\bigskip
\begin{tabular}{|c|c|c|} \hline  
 \multicolumn{3}{|l|}{   \BSS{pstPlanePut}[plane=xy,\RDD{planecorr}=normal](0,0,2)\AC{\BS{DFR}}} \\
  \multicolumn{3}{|l|}{   \BSS{pstPlanePut}[plane=xy,\RDD{planecorr}=normal ](0,0,-2)\AC{\BS{psframebox}\AC{texte}} }  \\ \hline
\begin{pspicture}(-4,-4)(3.5,4.5)
 \pstThreeDCoor[xMin=-4,yMin=-4,zMin=-4]
\pstPlanePut[plane=xy,planecorr=normal](0,0,-2){\psframebox{texte}}
\pstPlanePut[plane=xy,planecorr=normal](0,0,2){\DFR}
\pstThreeDDot[linecolor=red](0,0,-2)
\pstThreeDDot[linecolor=red](0,0,2) 
\end{pspicture}
&
\begin{pspicture}(-4,-4)(3.5,4.5)
\pstThreeDCoor[xMin=-4,yMin=-4,zMin=-4]
\pstPlanePut[plane=xy,planecorr=xyrot](0,0,-2){\psframebox{texte}}
\pstPlanePut[plane=xy,planecorr=xyrot](0,0,2){\DFR}
\pstThreeDDot[linecolor=red](0,0,-2)
\pstThreeDDot[linecolor=red](0,0,2) 
\end{pspicture}
&
\begin{pspicture}(-4,-4)(3.5,4.5)
\pstThreeDCoor[xMin=-4,yMin=-4,zMin=-4]
\pstPlanePut[plane=xy,planecorr=off](0,0,-2){\psframebox{texte}}
\pstPlanePut[plane=xy,planecorr=off](0,0,2){\DFR}
\pstThreeDDot[linecolor=red](0,0,-2)
\pstThreeDDot[linecolor=red](0,0,2) 
\end{pspicture}
 \\ \hline 
planecorr=normal &  planecorr=xyrot & planecorr=off
  \\ \hline 
\end{tabular}


\newpage
%\subsection[Créer un graphe en 3D]{Créer un graphe en 3D } 
\SbSSCT{Créer un graphe en 3D}{3D graph}

%\subsection{Graphe en 3D à partir d'équations}
\subsubsection{psplotThreeD}
\psset{unit=.5cm}
\psset{fillstyle=none,fillcolor=yellow,linecolor=blue,dotscale=.25}

\begin{tabular}{|c|} \hline  
\BSS{psplotThreeD}[algebraic]   \Rnode*[fillcolor=green]{A}{(-4,4)(-4,4)} \Rnode*[fillcolor=yellow]{B}{\AC{sin(x)*cos(y)}}\\

\\
 \hspace{1cm} \rnode{AA}{\TFRGB{plages pour x et y}{ranges for x and y}} \hspace{1cm}\rnode{BB}{\TFRGB{fonction en x y}{function y x} }
\\ \hline 
\begin{pspicture}(-6,-4)(6,6)
\psset{Beta=45}
 \psplotThreeD[algebraic](-4,4)(-4,4){sin(x)*cos(y)}
\pstThreeDCoor[xMin=-1,xMax=5,yMin=-1,yMax=5,zMin=-1,zMax=6]
\end{pspicture}

 \\ \hline 
\end{tabular} 
 \ncline[linecolor=blue]{A}{AA}  \ncline[linecolor=blue]{B}{BB}
 
\bigskip 
%\subsubsection{Option plotstyle}
\psset{unit=.25cm}


\begin{tabular}{|c|c|c|} \hline 
 \multicolumn{3}{|c|}{ \BS{psplotThreeD}[algebraic,\RDD{plotstyle}=line](-4,4)(-4,4)\AC{sin(x)*cos(y)} } \\ \hline 

\begin{pspicture}(-6,-4)(6,6)
\psset{Beta=45}
 \psplotThreeD[algebraic,plotstyle=line](-4,4)(-4,4){sin(x)*cos(y)}
\pstThreeDCoor[xMin=-1,xMax=5,yMin=-1,yMax=5,zMin=-1,zMax=6]
\end{pspicture}
&
\begin{pspicture}(-6,-4)(6,6)
\psset{Beta=45}
 \psplotThreeD[algebraic,plotstyle=polygon](-4,4)(-4,4){sin(x)*cos(y)}
\pstThreeDCoor[xMin=-1,xMax=5,yMin=-1,yMax=5,zMin=-1,zMax=6]
\end{pspicture}
&
\begin{pspicture}(-6,-4)(6,6)
\psset{Beta=45}
 \psplotThreeD[algebraic,plotstyle=dots](-4,4)(-4,4){sin(x)*cos(y)}
\pstThreeDCoor[xMin=-1,xMax=5,yMin=-1,yMax=5,zMin=-1,zMax=6]
\end{pspicture}
\\ \hline  
 plotstyle=line & plotstyle=polygon &  plotstyle=dots \\ \hline 
\begin{pspicture}(-6,-4)(6,6)
\psset{Beta=45}
 \psplotThreeD[algebraic,plotstyle=curve](-4,4)(-4,4){sin(x)*cos(y)}
\pstThreeDCoor[xMin=-1,xMax=5,yMin=-1,yMax=5,zMin=-1,zMax=6]
\end{pspicture}
& 
\begin{pspicture}(-6,-4)(6,6)
\psset{Beta=45}
 \psplotThreeD[algebraic,plotstyle=ecurve](-4,4)(-4,4){sin(x)*cos(y)}
\pstThreeDCoor[xMin=-1,xMax=5,yMin=-1,yMax=5,zMin=-1,zMax=6]
\end{pspicture}
&
\begin{pspicture}(-6,-4)(6,6)
\psset{Beta=45}
 \psplotThreeD[algebraic,plotstyle=ccurve](-4,4)(-4,4){sin(x)*cos(y)}
\pstThreeDCoor[xMin=-1,xMax=5,yMin=-1,yMax=5,zMin=-1,zMax=6]
\end{pspicture}

\\ \hline  
plotstyle=curve & plotstyle=ecurve & plotstyle=ccurve  \\ \hline 

\end{tabular} 
\bigskip

\begin{tabular}{|c|c|c|c|} \hline  
 \multicolumn{4}{|c|}{  \BS{psplotThreeD}[algebraic,\RDD{drawStyle}=xLines](-4,4)(-4,4)\AC{sin(x)*cos(y)} } \\ \hline
\begin{pspicture}(-6,-4)(6,6)
\psset{Beta=45}
 \psplotThreeD[algebraic,drawStyle=xLines](-4,4)(-4,4){sin(x)*cos(y)}
\pstThreeDCoor[xMin=-1,xMax=5,yMin=-1,yMax=5,zMin=-1,zMax=6]
\end{pspicture}
& 
\begin{pspicture}(-6,-4)(6,6)
\psset{Beta=45}
 \psplotThreeD[algebraic,drawStyle=yLines](-4,4)(-4,4){sin(x)*cos(y)}
\pstThreeDCoor[xMin=-1,xMax=5,yMin=-1,yMax=5,zMin=-1,zMax=6]
\end{pspicture}
&
\begin{pspicture}(-6,-4)(6,6)
\psset{Beta=45}
 \psplotThreeD[algebraic,drawStyle=xyLines](-4,4)(-4,4){sin(x)*cos(y)}
\pstThreeDCoor[xMin=-1,xMax=5,yMin=-1,yMax=5,zMin=-1,zMax=6]
\end{pspicture}
&
\begin{pspicture}(-6,-4)(6,6)
\psset{Beta=45}
 \psplotThreeD[algebraic,drawStyle=yxLines](-4,4)(-4,4){sin(x)*cos(y)}
\pstThreeDCoor[xMin=-1,xMax=5,yMin=-1,yMax=5,zMin=-1,zMax=6]
\end{pspicture}
\\ \hline  
  drawStyle=xLines & drawStyle=yLines & drawStyle=xyLines & drawStyle=yxLines \\ 
\dft & & & \\ \hline 
\end{tabular} 

\bigskip

\begin{tabular}{|c|c|} \hline 
 \multicolumn{2}{|c|}{  \BS{psplotThreeD}[algebraic,\RDD{showpoints}=false,linewidth=.1pt] } \\ 
 \multicolumn{2}{|c|}{ (-4,4)(-4,4)\AC{sin(x)*cos(y)} } \\ \hline 
\begin{pspicture}(-6,-4)(6,6)
\psset{Beta=45}
 \psplotThreeD[algebraic,showpoints=false,linewidth=.1pt](-4,4)(-4,4){sin(x)*cos(y)}
\pstThreeDCoor[xMin=-1,xMax=5,yMin=-1,yMax=5,zMin=-1,zMax=6]
\end{pspicture}
&  
\begin{pspicture}(-6,-4)(6,6)
\psset{Beta=45}
 \psplotThreeD[algebraic,showpoints=true,linewidth=.1pt](-4,4)(-4,4){sin(x)*cos(y)}
\pstThreeDCoor[xMin=-1,xMax=5,yMin=-1,yMax=5,zMin=-1,zMax=6]
\end{pspicture}
\\ \hline  
showpoints=false  &  showpoints=true \\ \hline 
\dft & \\ \hline 
\end{tabular}

\bigskip
\begin{tabular}{|c|c|c|c|} \hline 
 \multicolumn{4}{|c|}{  \BS{psplotThreeD}[algebraic,\RDD{xPlotpoints}=5,drawStyle=xyLines](-4,4)(-4,4)\AC{sin(x)*cos(y)} } \\ \hline
\begin{pspicture}(-6,-4)(6,6)
\psset{Beta=45}
 \psplotThreeD[algebraic,plotstyle=line,xPlotpoints=5,drawStyle=xyLines](-4,4)(-4,4){sin(x)*cos(y)}
\pstThreeDCoor[xMin=-1,xMax=5,yMin=-1,yMax=5,zMin=-1,zMax=6]
\end{pspicture}
&  
\begin{pspicture}(-6,-4)(6,6)
\psset{Beta=45}
 \psplotThreeD[algebraic,plotstyle=line,yPlotpoints=5,drawStyle=xyLines](-4,4)(-4,4){sin(x)*cos(y)}
\pstThreeDCoor[xMin=-1,xMax=5,yMin=-1,yMax=5,zMin=-1,zMax=6]
\end{pspicture}
&
\begin{pspicture}(-6,-4)(6,6)
\psset{Beta=45}
 \psplotThreeD[algebraic,plotstyle=line,xPlotpoints=5,yPlotpoints=5,drawStyle=xyLines](-4,4)(-4,4){sin(x)*cos(y)}
\pstThreeDCoor[xMin=-1,xMax=5,yMin=-1,yMax=5,zMin=-1,zMax=6]
\end{pspicture}
&
\begin{pspicture}(-6,-4)(6,6)
\psset{Beta=45}
 \psplotThreeD[algebraic,plotstyle=line,xPlotpoints=50,drawStyle=xyLines](-4,4)(-4,4){sin(x)*cos(y)}
\pstThreeDCoor[xMin=-1,xMax=5,yMin=-1,yMax=5,zMin=-1,zMax=6]
\end{pspicture}
\\ \hline  
xPlotpoints=5  &  yPlotpoints=5  & yPlotpoints=5 & xPlotpoints=50 \\ 
  &    &  yPlotpoints=5 &\\ \hline 
 \multicolumn{4}{|c|}{\blue \dft : xPlotpoints=25 \hspace{1cm}  yPlotpoints=25 }
\\ \hline 
\end{tabular}

\bigskip


\begin{tabular}{|c|c|} \hline 
 \multicolumn{2}{|c|}{  \BS{psplotThreeD}[algebraic,\RDD{hiddenLine}=false](-4,4)(-4,4)\AC{sin(x)*cos(y)} } \\ \hline 
\begin{pspicture}(-6,-4)(6,6)
 \psplotThreeD[algebraic,hiddenLine=false,plotstyle=curve](-4,4)(-4,4){sin(x)*cos(y)}
\end{pspicture}
&  
\begin{pspicture}(-6,-4)(6,6)
 \psplotThreeD[algebraic,hiddenLine=true,plotstyle=curve](-4,4)(-4,4){sin(x)*cos(y)}
\end{pspicture}
\begin{pspicture}(-6,-4)(6,6)
 \psplotThreeD[algebraic,hiddenLine=true,plotstyle=curve, drawStyle=yLines ](-4,4)(-4,4){sin(x)*cos(y)}
\end{pspicture}
\\ \hline  
hiddenLine=false  &  hiddenLines=true \\ \hline 
\dft & \\ \hline 
\end{tabular}

%-------------------------------------------------------
\subsubsection{parametricplotThreeD}


\psset{unit=.5cm}

\begin{tabular}{|c|}\hline  
\BSS{parametricplotThreeD}[xPlotpoints=200,plotstyle=curve,algebraic]\\
\Rnode*[fillcolor=green]{A}{(0,50)}  \Rnode*[fillcolor=yellow]{B}{\AC{ t/10*cos(t) | t/10*sin(t) | t/10}}\\
\\
\rnode{AA}{ \TFRGB{plage pour t}{range for t} } \hspace{1cm}\rnode{BB}{\TFRGB{3 fonctions paramétriques}{3 parametric functions}}
\\ \hline  
\begin{pspicture}(-6,-2.25)(6,7)
\pstThreeDCoor[zMax=5]
\parametricplotThreeD[xPlotpoints=200,plotstyle=curve,algebraic](0,50){ t/10*cos(t) | t/10*sin(t) | t/10}
 \end{pspicture}
\\ \hline 
\end{tabular} 
 \ncline[linecolor=blue]{A}{AA}  \ncline[linecolor=blue]{B}{BB}

\newpage
% % %-------------------------------------------------
\SbSSCT{Graphe en 3D à partir d'un fichier de données}{3D graph from a data file}
\subsubsection{fileplotThreeD}

\begin{tabular}{|c|c|} \hline 
\begin{pspicture}[shift=*](-6,-2.25)(6,7)
\pstThreeDCoor[zMax=5]
\fileplotThreeD{data3d.txt}
 \end{pspicture}%
 &  
 \parbox{8cm}{

\BSS{fileplotThreeD}\AC{data3d.txt} \\
 \\
 \% data3d.txt : \TFRGB{fichier de données créé avec}{data file created with} Excel \\
\\

}
 \\  \hline 
\end{tabular} 
% % %-------------------------------------------------------------

\subsubsection{dataplotThreeD}

\begin{tabular}{|c|c|} \hline 
\readdata{\data}{data3d.txt}
\begin{pspicture}[shift=*](-6,-2.25)(6,7)
\pstThreeDCoor[zMax=5]
\dataplotThreeD[plotstyle=line]{\data}
 \end{pspicture}%
 &  
 \parbox{8cm}{

 \BS{readdata}\AC{\BS{data}}\AC{data3d.txt} \\
 \\
 \% data3d.txt : \TFRGB{fichier de données créé avec}{data file created with} Excel \\
\\
\BSS{dataplotThreeD}[plotstyle=line]\AC{\BS{data}} \\
}
 \\  \hline 
\end{tabular} 
%%--------------------------------------------------------
\subsubsection{listplotThreeD}

\begin{tabular}{|c|c|} \hline 
\readdata{\data}{data3d.txt}
\begin{pspicture}[shift=*](-6,-2.25)(6,7)
\pstThreeDCoor[zMax=5]
\listplotThreeD[plotstyle=curve]{\data}
 \end{pspicture}%
 &  
 \parbox{8cm}{

 \BS{readdata}\AC{\BS{data}}\AC{data3d.txt} \\
 \\
 \% data3d.txt : \TFRGB{fichier de données créé avec}{data file created with} Excel \\
\\
\BSS{listplotThreeD}[plotstyle=curve]\AC{\BS{data}} \\
}
\\  \hline  
\end{tabular} 





\newpage
%\section[Créer un dessin en 3D avec pst-solides3d]{Créer un dessin en 3D avec pst-solides3d \cite{pst-sol3d}}
\SSCT{Créer un dessin en 3D avec pst-solides3d}
{3D solid}
Utilisation du module \textbf{pst-solides3d} \label{sol3d}

\emph{Cette partie sera complétée dans uns version ultérieure}

\psset{fillstyle=none,fillcolor=yellow,linecolor=blue,incolor=green!50}


\subsection{Axes}
\psset{unit=.25cm}
\begin{tabular}{|c|c|c|} \hline 
 \begin{pspicture}(-7,-4)(7,10)
 \psframe(-7,-4)(7,10)
 \axesIIID(0,0,0)(2,2,2)
 \end{pspicture}
 &
 \begin{pspicture}(-7,-4)(7,10)
 \psframe(-7,-4)(7,10)
\axesIIID(2,2,2)(2,2,2)
\end{pspicture}
&  
 \begin{pspicture}(-7,-4)(7,10)
 \psframe(-7,-4)(7,10)
\axesIIID(1,1,1)(2,2,2)
\end{pspicture}
\\ \hline
 \BSS{axesIIID}(0,0,0)(2,2,2)  \BSI{axesIIID}{pst-sol3d}
 & 
\BSS{axesIIID}(2,2,2)(2,2,2)
&  
\BSS{axesIIID}(1,1,1)(2,2,2)
\\ \hline 
\end{tabular} 

\bigskip
\begin{tabular}{|c|c|c|} \hline 
 \begin{pspicture}(-7,-4)(7,10)
 \psframe(-7,-4)(7,10)
 \axesIIID[labelsep=0cm](0,0,0)(2,2,2)
 \end{pspicture}
 &
 \begin{pspicture}(-7,-4)(7,10)
 \psframe(-7,-4)(7,10)
\axesIIID[axisnames={a,b,c}](0,0,0)(2,2,2)
\end{pspicture}
&  
 \begin{pspicture}(-7,-4)(7,10)
 \psframe(-7,-4)(7,10)
\axesIIID[showOrigin=false](1,1,1)(2,2,2)
\end{pspicture}
\\ \hline
\RDD{labelsep}=0cm
 & 
\RDD{axisnames}=\AC{a,b,c} \RDI{axisnames}{pst-sol3d}
&  
\RDD{showOrigin}=false \RDI{showorigin}{pst-sol3d}
\\ \hline 
\dft : labelsep=5pt & \dft : axisnames=\AC{x,y,z} & \dft : showOrigin=true\\ \hline 
\end{tabular} 
   
%\subsection{\'Element en 3D}  
\SbSSCT{\'Element en 3D}{3D elements}

\psset{unit=.5cm}


%--------------------point, line, vector,
%\subsubsection{point, line, vector}
\SbSbSSCT{point, ligne, vecteur}{point, line, vector}

\begin{tabular}{|c|c|c|} \hline  
% \multicolumn{3}{|l|}{\BS{psSolid}[object=point,args=1 2 2]} \\
\begin{pspicture}(-4,-4)(4,4)
\psframe(-4,-4)(4,4)
\axesIIID(0,0,0)(1,1,1)
\psSolid[object=point,args=1 2 2,linecolor=red]			%point
\psSolid[object=line,args=0 0 0 1 2 0,linestyle=dotted,linecolor=red]
\psSolid[object=line,args=1 2 2 1 2 0,linestyle=dotted,linecolor=red]
\end{pspicture}
&  
\begin{pspicture}(-4,-4)(4,4)
\psframe(-4,-4)(4,4)
\axesIIID(0,0,0)(1,1,1)

\psSolid[object=line,args=0 -1 0 1 2 2] % line
\psSolid[object=line,args=0 0 0 0 -2 0,linecolor=red]
\psSolid[object=line,args=0 0 0 1 2 0,linestyle=dotted,linecolor=red]
\psSolid[object=line,args= 1 2 0 1 2 2,linestyle=dotted,linecolor=red]
\end{pspicture}
&  
\begin{pspicture}(-4,-4)(4,4)
\psframe(-4,-4)(4,4)
\axesIIID(0,0,0)(1,1,1)
\psSolid[object=vecteur,args=1 2 2,linecolor=red] % vecteur
\psSolid[object=line,args=0 0 0 1 2 0,linestyle=dotted,linecolor=red]
\psSolid[object=line,args=1 2 2 1 2 0,linestyle=dotted,linecolor=red]
\end{pspicture}\\ \hline
[\RDD{object=point},args=1 2 2] \RDI{object=point}{pst-sol3d} & 
[\RDD{object=line},args=0 -1 0 1 2 2] \RDI{object=line}{pst-sol3d}&
[\RDD{object=vecteur},args=1 2 2] \RDI{object=vecteur}{pst-sol3d}
\\ \hline
\end{tabular} 
\bigskip
 %------------------------- plan, grille, cube
 \psset{unit=.2cm}
 
 
% \subsubsection{ plan}
\SbSbSSCT{Plan}{Plane}
 
\begin{tabular}{|c|c|} \hline 
 \multicolumn{2}{|l|}{\BS{psSolid}[\RDD{object=plan},definition=equation,args=\Rnode*[fillcolor=yellow]{A}{\AC{[0 0 1 0]}},base=-2 2 -3 3] \RDI{object=plan}{pst-sol3d} } \\ 
 \multicolumn{2}{|l|}{ } \\ 
 \multicolumn{2}{|c|}{  \rnode{AA}{coeff de l'équation  ax+by+cz+d = 0} } \\  \hline
\begin{pspicture}(-12,-10)(12,10)
\psframe(-12,-10)(12,10)
\psSolid[object=plan,definition=equation,args={[0 0 1 0]},base=-2 2 -3 3]
 \axesIIID[linecolor=red](0,0,0)(1,1,1)
\end{pspicture}
&
\begin{pspicture}(-12,-10)(12,10)
 \psframe(-12,-10)(12,10)
\psSolid[object=plan,definition=equation,args={[0 1 0 0]},base=-2 2 -3 3]
 \axesIIID[linecolor=red](0,0,0)(1,1,1)
\end{pspicture}
\\ \hline
args=\AC{[0 0 1 0]} & args=\AC{[0 1 0 0]}  \\ \hline
\end{tabular} 
\ncline[linecolor=blue]{A}{AA}
\bigskip

\begin{tabular}{|c|c|c|} \hline 
% \multicolumn{3}{|c|}{ \BS{psSolid}[object=cube,a=3,action=draw] } \\  \hline 
\begin{pspicture}(-10,-10)(10,10)
 \psframe(-10,-10)(10,10)
%\axesIIID(1,1,1)(1,1,1)
\psSolid[object=plan,definition=equation,args={[1 0 0 0]},base=-2 2 -3 3]
 \axesIIID[linecolor=red](0,0,0)(1,1,1)
\end{pspicture}
&
\begin{pspicture}(-10,-10)(10,10)
 \psframe(-10,-10)(10,10)
%\axesIIID(1,1,1)(1,1,1)
\psSolid[object=plan,definition=equation,args={[0 0 1 1]},base=-2 2 -3 3]
 \axesIIID[linecolor=red](0,0,0)(1,1,1)
\end{pspicture}
 & 
\begin{pspicture}(-10,-10)(10,10)
 \psframe(-10,-10)(10,10)
%\axesIIID(1,1,1)(1,1,1)
\psSolid[object=plan,definition=equation,args={[1 1 0 0]},base=-2 2 -3 3]
 \axesIIID[linecolor=red](0,0,0)(1,1,1)
\end{pspicture}
\\ \hline
args={[1 0 0 0]} & args={[0 0 1 1]} & args={[1 1 0 0]} \\ \hline
\end{tabular}


%------------------------------------------

%\subsubsection{grille}
\SbSbSSCT{Grille}{Grid}

\begin{tabular}{|c|c|c|} \hline 
\multicolumn{3}{|c|}{\BS{psSolid}[\RDD{object=grille},base=-2 2 -3 3] \RDI{object=grille}{pst-sol3d}} \\ \hline
\begin{pspicture}(-12,-10)(12,10)
 \psframe(-12,-10)(12,10)
\psSolid[object=grille,base=-2 2 -3 3] % grille
 \axesIIID[linecolor=red](0,0,0)(1,1,1)
\end{pspicture}
&
\begin{pspicture}(-12,-10)(12,10)
 \psframe(-12,-10)(12,10)
\psSolid[object=grille,base=-2 2 -3 3,RotX=-90] % grille
 \axesIIID[linecolor=red](0,0,0)(1,1,1)
\end{pspicture}
&
\begin{pspicture}(-12,-10)(12,10)
 \psframe(-12,-10)(12,10)
\psSolid[object=grille,base=-2 2 -3 3,RotY=90] % grille
 \axesIIID[linecolor=red](0,0,0)(1,1,1)
\end{pspicture}

\\ \hline
\dft & RotX=90 & RotY=90  \\ \hline
\end{tabular} 
\bigskip

%------------------------------------------------------------------------------

%
%\begin{pspicture}(-10,-10)(10,10)
% \psframe(-10,-10)(10,10)
%%\axesIIID(1,1,1)(1,1,1)
%\psSolid[object=cube,action=draw] %cube
% \axesIIID[linecolor=red](0,0,0)(1,1,1)
%\end{pspicture}
%
%\begin{pspicture}(-10,-5)(10,15)
% \psframe(-10,-5)(10,15)
%\psSolid[object=cylindre,action=draw](0,0,0)%
% \axesIIID[linecolor=red](0,0,0)(1,1,1)
%\end{pspicture}

\subsubsection{cube}
\begin{tabular}{|c|c|c|} \hline 
 \multicolumn{3}{|c|}{ \BS{psSolid}[object=cube,a=3,action=draw] } \\  \hline 
\begin{pspicture}(-10,-10)(10,10)
 \psframe(-10,-10)(10,10)
%\axesIIID(1,1,1)(1,1,1)
\psSolid[object=cube,a=3,action=draw] %cube
 \axesIIID[linecolor=red](0,0,0)(1,1,1)
\end{pspicture}
&
\begin{pspicture}(-10,-10)(10,10)
 \psframe(-10,-10)(10,10)
%\axesIIID(1,1,1)(1,1,1)
\psSolid[object=cube,a=3,action=draw*] %cube
 \axesIIID[linecolor=red](0,0,0)(1,1,1)
\end{pspicture}
 & 
\begin{pspicture}(-10,-10)(10,10)
 \psframe(-10,-10)(10,10)
%\axesIIID(1,1,1)(1,1,1)
\psSolid[object=cube,a=3,action=draw**] %cube
 \axesIIID[linecolor=red](0,0,0)(1,1,1)
\end{pspicture}
\\ \hline
action=draw & action=draw* & action=draw**\\ \hline
\end{tabular} 
\bigskip


%\subsubsection{cylindre}
\SbSbSSCT{Cylindre}{Cylinder}

\begin{tabular}{|c|c|c|} \hline 
 \multicolumn{3}{|c|}{ \BS{psSolid}[\RDD{object=cylindre},h=3,r=2,action=draw](0,0,0) \RDI{object=cylindre}{pst-sol3d} } \\  \hline 
\begin{pspicture}(-10,-5)(10,15)
 \psframe(-10,-5)(10,15)
\psSolid[object=cylindre,h=3,r=2,action=draw](0,0,0)%
 \axesIIID[linecolor=red](0,0,0)(1,1,1)
\end{pspicture}
&
\begin{pspicture}(-10,-5)(10,15)
 \psframe(-10,-5)(10,15)
\psSolid[object=cylindre,h=3,r=2,action=draw*](0,0,0)%
 \axesIIID[linecolor=red](0,0,0)(1,1,1)
\end{pspicture}
&
\begin{pspicture}(-10,-5)(10,15)
 \psframe(-10,-5)(10,15)
\psSolid[object=cylindre,h=3,r=2,action=draw**](0,0,0)%
 \axesIIID[linecolor=red](0,0,0)(1,1,1)
\end{pspicture}\\ \hline
action=draw & action=draw* & action=draw**\\ \hline
\end{tabular} 


%\subsubsection{cylindrecreux}
\SbSbSSCT{cylindre creux}{Tube}

\begin{tabular}{|c|c|c|} \hline 
 \multicolumn{3}{|c|}{ \BS{psSolid}[\RDD{object=cylindrecreux},h=3,r=2,action=draw](0,0,0) \RDI{object=cylindrecreux}{pst-sol3d} } \\  \hline 
\begin{pspicture}(-10,-5)(10,15)
 \psframe(-10,-5)(10,15)
\psSolid[object=cylindrecreux,h=3,r=2,action=draw](0,0,0)%
 \axesIIID[linecolor=red](0,0,0)(1,1,1)
\end{pspicture}
&
\begin{pspicture}(-10,-5)(10,15)
 \psframe(-10,-5)(10,15)
\psSolid[object=cylindrecreux,h=3,r=2,action=draw*](0,0,0)%
 \axesIIID[linecolor=red](0,0,0)(1,1,1)
\end{pspicture}
&
\begin{pspicture}(-10,-5)(10,15)
 \psframe(-10,-5)(10,15)
\psSolid[object=cylindrecreux,h=3,r=2,action=draw**](0,0,0)%
 \axesIIID[linecolor=red](0,0,0)(1,1,1)
\end{pspicture}\\ \hline
action=draw & action=draw* & action=draw**\\ \hline
\end{tabular} 


%\subsubsection{cone}
\SbSbSSCT{Cône}{Cone}

\begin{tabular}{|c|c|c|} \hline 
 \multicolumn{3}{|c|}{ \BS{psSolid}[\RDD{object=cone},h=3,r=2,action=draw] \RDI{object=cone}{pst-sol3d} } \\  \hline 
\begin{pspicture}(-10,-10)(10,10)
 \psframe(-10,-10)(10,10)
\psSolid[object=cone,h=3,r=2,action=draw]%
 \axesIIID[linecolor=red](0,0,0)(1,1,1)
\end{pspicture}
&
\begin{pspicture}(-10,-10)(10,10)
 \psframe(-10,-10)(10,10)
\psSolid[object=cone,h=3,r=2,action=draw*]%
 \axesIIID[linecolor=red](0,0,0)(1,1,1)
\end{pspicture}
&
\begin{pspicture}(-10,-10)(10,10)
 \psframe(-10,-10)(10,10)
\psSolid[object=cone,h=3,r=2,action=draw**]%
 \axesIIID[linecolor=red](0,0,0)(1,1,1)
\end{pspicture}\\ \hline
action=draw & action=draw* & action=draw**\\ \hline
\end{tabular} 

\subsubsection{conecreux}
\SbSbSSCT{Cône creux}{Empty cone}

\begin{tabular}{|c|c|c|} \hline 
 \multicolumn{3}{|c|}{ \BS{psSolid}[\RDD{object=conecreux},h=4,r=2,action=draw] \RDI{object=conecreux}{pst-sol3d} } \\  \hline 
\begin{pspicture}(-10,-10)(10,10)
 \psframe(-10,-10)(10,10)
\psSolid[object=conecreux,h=4,r=2,action=draw]%
 \axesIIID[linecolor=red](0,0,0)(1,1,1)
\end{pspicture}
&
\begin{pspicture}(-10,-10)(10,10)
 \psframe(-10,-10)(10,10)
\psSolid[object=conecreux,h=4,r=2,action=draw*]%
 \axesIIID[linecolor=red](0,0,0)(1,1,1)
\end{pspicture}
&
\begin{pspicture}(-10,-10)(10,10)
 \psframe(-10,-10)(10,10)
\psSolid[object=conecreux,h=4,r=2,action=draw**]%
 \axesIIID[linecolor=red](0,0,0)(1,1,1)
\end{pspicture}\\ \hline
action=draw & action=draw* & action=draw**\\ \hline
\end{tabular} 
\bigskip

%\subsubsection{tronccone}
\SbSbSSCT{Tronc de cône}{Truncated cone}

\begin{tabular}{|c|c|c|} \hline 
 \multicolumn{3}{|c|}{ \BS{psSolid}[\RDD{object=tronccone},r0=2,r1=1,h=4,action=draw] \RDI{object=troncone}{pst-sol3d} } \\  \hline 
\begin{pspicture}(-10,-5)(10,15)
 \psframe(-10,-5)(10,15)
\psSolid[object=tronccone,r0=2,r1=1,h=4,action=draw]%
 \axesIIID[linecolor=red](0,0,0)(1,1,1)
\end{pspicture}
&
\begin{pspicture}(-10,-5)(10,15)
 \psframe(-10,-5)(10,15)
\psSolid[object=tronccone,r0=2,r1=1,h=4,action=draw*]%
 \axesIIID[linecolor=red](0,0,0)(1,1,1)
\end{pspicture}
&
\begin{pspicture}(-10,-5)(10,15)
 \psframe(-10,-5)(10,15)
\psSolid[object=tronccone,r0=2,r1=1,h=4,action=draw**]%
 \axesIIID[linecolor=red](0,0,0)(1,1,1)
\end{pspicture}\\ \hline
action=draw & action=draw* & action=draw**\\ \hline
\end{tabular} 

%\subsubsection{troncconecreux}
\SbSbSSCT{Tronc de cône creux creux}{Empty truncated cone}

\begin{tabular}{|c|c|c|} \hline 
 \multicolumn{3}{|c|}{ \BS{psSolid}[\RDD{object=troncconecreux},r0=2,r1=1,h=4,action=draw] \RDI{object=tronconecreux}{pst-sol3d} } \\  \hline 
\begin{pspicture}(-10,-5)(10,15)
 \psframe(-10,-5)(10,15)
\psSolid[object=troncconecreux,r0=2,r1=1,h=4,action=draw]%
 \axesIIID[linecolor=red](0,0,0)(1,1,1)
\end{pspicture}
&
\begin{pspicture}(-10,-5)(10,15)
 \psframe(-10,-5)(10,15)
\psSolid[object=troncconecreux,r0=2,r1=1,h=4,action=draw*]%
 \axesIIID[linecolor=red](0,0,0)(1,1,1)
\end{pspicture}
&
\begin{pspicture}(-10,-5)(10,15)
 \psframe(-10,-5)(10,15)
\psSolid[object=troncconecreux,r0=2,r1=1,h=4,action=draw**]%
 \axesIIID[linecolor=red](0,0,0)(1,1,1)
\end{pspicture}\\ \hline
action=draw & action=draw* & action=draw**\\ \hline
\end{tabular} 
\bigskip

\subsubsection{sphere}


\begin{tabular}{|c|c|c|} \hline 
 \multicolumn{3}{|c|}{ \BS{psSolid}[\RDD{object=sphere},r=1,action=draw] \RDI{object=sphere}{pst-sol3d} } \\  \hline 
\begin{pspicture}(-10,-10)(10,10)
 \psframe(-10,-10)(10,10)
\psSolid[object=sphere,r=3,action=draw]%
 \axesIIID[linecolor=red](0,0,0)(1,1,1)
\end{pspicture}
&
\begin{pspicture}(-10,-10)(10,10)
 \psframe(-10,-10)(10,10)
\psSolid[object=sphere,r=3,action=draw*]%
 \axesIIID[linecolor=red](0,0,0)(1,1,1)
\end{pspicture}
&
\begin{pspicture}(-10,-10)(10,10)
 \psframe(-10,-10)(10,10)
\psSolid[object=sphere,r=3,action=draw**]%
 \axesIIID[linecolor=red](0,0,0)(1,1,1)
\end{pspicture}\\ \hline
action=draw & action=draw* & action=draw**\\ \hline
\end{tabular} 
\bigskip

%\subsubsection{calottesphere}
\SbSbSSCT{Calotte sphérique}{Spherical cup}

\begin{tabular}{|c|c|c|} \hline 
 \multicolumn{3}{|c|}{ \BS{psSolid}[\RDD{object=calottesphere},r=3,action=draw] \RDI{object=calottesphere}{pst-sol3d} } \\  \hline 
\begin{pspicture}(-10,-10)(10,10)
 \psframe(-10,-10)(10,10)
\psSolid[object=calottesphere,r=3,action=draw]%
 \axesIIID[linecolor=red](0,0,0)(1,1,1)
\end{pspicture}
&
\begin{pspicture}(-10,-10)(10,10)
 \psframe(-10,-10)(10,10)
\psSolid[object=calottesphere,r=3,action=draw*]%
 \axesIIID[linecolor=red](0,0,0)(1,1,1) \axesIIID[linecolor=red](0,0,0)(1,1,1)
\end{pspicture}
&
\begin{pspicture}(-10,-10)(10,10)
 \psframe(-10,-10)(10,10)
\psSolid[object=calottesphere,r=3,action=draw**]%
 \axesIIID[linecolor=red](0,0,0)(1,1,1)
\end{pspicture}\\ \hline
action=draw & action=draw* & action=draw**\\ \hline
\end{tabular} 

%\subsubsection{calottespherecreuse}
\SbSbSSCT{calotte spherique creuse}{empty spherical cup}

\begin{tabular}{|c|c|c|} \hline 
 \multicolumn{3}{|c|}{ \BS{psSolid}[\RDD{object=calottespherecreuse},r=3,action=draw] \RDI{object=calottespherecreuse}{pst-sol3d} } \\  \hline 
\begin{pspicture}(-10,-10)(10,10)
 \psframe(-10,-10)(10,10)
\psSolid[object=calottespherecreuse,r=3,action=draw]%
\end{pspicture}
&
\begin{pspicture}(-10,-10)(10,10)
 \psframe(-10,-10)(10,10)
\psSolid[object=calottespherecreuse,r=3,action=draw*]%
\end{pspicture}
&
\begin{pspicture}(-10,-10)(10,10)
 \psframe(-10,-10)(10,10)
\psSolid[object=calottespherecreuse,r=3,action=draw**]%
\end{pspicture}\\ \hline
action=draw & action=draw* & action=draw**\\ \hline
\end{tabular} 

%\subsubsection{tore}
\SbSbSSCT{Tore}{Torus}

\begin{tabular}{|c|c|c|} \hline 
 \multicolumn{3}{|c|}{ \BS{psSolid}[r1=2,r0=1, \RDD{object=tore},ngrid=18 36,action=draw] \RDI{object=tore}{pst-sol3d} } \\  \hline 
\begin{pspicture}(-10,-10)(10,10)
 \psframe(-10,-10)(10,10)
\psSolid[r1=2,r0=1,object=tore,ngrid=18 36,action=draw]%
\end{pspicture}
&
\begin{pspicture}(-10,-10)(10,10)
 \psframe(-10,-10)(10,10)
\psSolid[r1=2,r0=1,object=tore,ngrid=18 36,action=draw*]%
\end{pspicture}
&
\begin{pspicture}(-10,-10)(10,10)
 \psframe(-10,-10)(10,10)
\psSolid[r1=2,r0=1,object=tore,ngrid=18 36,action=draw**]%
\end{pspicture}\\ \hline
action=draw & action=draw* & action=draw**\\ \hline
\end{tabular} 
\bigskip

%\subsubsection{anneau}
\SbSbSSCT{Anneau}{Ring}

\begin{tabular}{|c|c|c|} \hline 
 \multicolumn{3}{|c|}{ \BS{psSolid}[\RDD{object=anneau},h=1,R=2,r=1,action=draw] \RDI{object=anneau}{pst-sol3d} } \\  \hline 
\begin{pspicture}(-10,-10)(10,10)
 \psframe(-10,-10)(10,10)
\psSolid[object=anneau,h=1,R=2,r=1,action=draw]%
\end{pspicture}
&
\begin{pspicture}(-10,-10)(10,10)
 \psframe(-10,-10)(10,10)
\psSolid[object=anneau,h=1,R=2,r=1,action=draw*]%
\end{pspicture}
&
\begin{pspicture}(-10,-10)(10,10)
 \psframe(-10,-10)(10,10)
\psSolid[object=anneau,h=1,R=2,r=1,action=draw**]%
\end{pspicture}\\ \hline
action=draw & action=draw* & action=draw**\\ \hline
\end{tabular} 
\bigskip
%------------------------------

%\subsubsection{ conecreux,tronccone, troncconecreux}
%\begin{tabular}{|c|c|c|} \hline 
%\begin{pspicture}(-10,-5)(10,15)
% \psframe(-10,-5)(10,15)
%\axesIIID(1,1,1)(1,1,1)
%\psSolid[object=conecreux,h=3,r=2,RotY=-60,mode=4] % conecreux
%\end{pspicture}
% & 
%\begin{pspicture}(-10,-5)(10,15)
% \psframe(-10,-5)(10,15)
% \axesIIID(1,1,1)(1,1,1)
%\psSolid[object=tronccone,r0=2,r1=1,h=3,mode=4] % tronccone
%\end{pspicture}
%&
%\begin{pspicture}(-10,-5)(10,15)
% \psframe(-10,-5)(10,15)
% \axesIIID(1,1,1)(1,1,1)
%\psSolid[object=tetrahedron,r=3,action=draw**]%
%\end{pspicture}
%\\ \hline 
%\end{tabular} 

\subsubsection{tetrahedron}

\begin{tabular}{|c|c|c|} \hline 
 \multicolumn{3}{|c|}{ \BS{psSolid}[object=tetrahedron,r=1,RotZ=30,action=draw] } \\  \hline 
\begin{pspicture}(-10,-10)(10,10)
 \psframe(-10,-10)(10,10)
%\axesIIID(1,1,1)(1,1,1)
\psSolid[object=tetrahedron,r=3,action=draw,RotZ=30]%
\end{pspicture}
 & 
\begin{pspicture}(-10,-10)(10,10)
 \psframe(-10,-10)(10,10)
% \axesIIID(1,1,1)(1,1,1)
\psSolid[object=tetrahedron,r=3,action=draw*,RotZ=30]%
\end{pspicture}
&
\begin{pspicture}(-10,-10)(10,10)
 \psframe(-10,-10)(10,10)
% \axesIIID(1,1,1)(1,1,1)
\psSolid[object=tetrahedron,r=3,action=draw**,RotZ=30]%
\end{pspicture}\\ \hline 
action=draw & action=draw* & action=draw**\\ \hline

\end{tabular} 


%------------ sphere, calottesphere, calottespherecreuse









\subsubsection{parallelepiped}

\begin{tabular}{|c|c|c|} \hline 
 \multicolumn{3}{|c|}{ \BS{psSolid}[\RDD{object=parallelepiped},a=1,b=2,c=3,action=draw] \RDI{object=parallelepiped}{pst-sol3d} } \\  \hline 
\begin{pspicture}(-10,-10)(10,10)
 \psframe(-10,-10)(10,10)
\psSolid[object=parallelepiped,a=1,b=2,c=3,action=draw]%
\end{pspicture}
&
\begin{pspicture}(-10,-10)(10,10)
 \psframe(-10,-10)(10,10)
\psSolid[object=parallelepiped,a=1,b=2,c=3,action=draw*]%
\end{pspicture}
&
\begin{pspicture}(-10,-10)(10,10)
 \psframe(-10,-10)(10,10)
\psSolid[object=parallelepiped,a=1,b=2,c=3,action=draw**]%
\end{pspicture}\\ \hline
action=draw & action=draw* & action=draw**\\ \hline
\end{tabular} 
\bigskip

\subsubsection{octahedron}
\begin{tabular}{|c|c|c|} \hline 
 \multicolumn{3}{|c|}{ \BS{psSolid}[\RDD{object=octahedron},a=30,action=draw] \RDI{object=octahedron}{pst-sol3d} } \\  \hline 
\begin{pspicture}(-10,-10)(10,10)
 \psframe(-10,-10)(10,10)
\psSolid[object=octahedron,a=3,action=draw]%
\end{pspicture}
&
\begin{pspicture}(-10,-10)(10,10)
 \psframe(-10,-10)(10,10)
\psSolid[object=octahedron,a=3,action=draw*]%
\end{pspicture}
&
\begin{pspicture}(-10,-10)(10,10)
 \psframe(-10,-10)(10,10)
\psSolid[object=octahedron,a=3,action=draw**]%
\end{pspicture}\\ \hline
action=draw & action=draw* & action=draw**\\ \hline
\end{tabular} 
\bigskip

\subsubsection{dodecahedron}

\begin{tabular}{|c|c|c|} \hline 
 \multicolumn{3}{|c|}{ \BS{psSolid}[\RDD{object=dodecahedron},a=2.5,RotZ=90,action=draw] \RDI{object=dodecahedron}{pst-sol3d} } \\  \hline 
\begin{pspicture}(-10,-10)(10,10)
 \psframe(-10,-10)(10,10)
\psSolid[object=dodecahedron,a=2.5,RotZ=90,action=draw]%
\end{pspicture}
&
\begin{pspicture}(-10,-10)(10,10)
 \psframe(-10,-10)(10,10)
\psSolid[object=dodecahedron,a=2.5,RotZ=90,action=draw*]%
\end{pspicture}
&
\begin{pspicture}(-10,-10)(10,10)
 \psframe(-10,-10)(10,10)
\psSolid[object=dodecahedron,a=2.5,RotZ=90,action=draw**]%
\end{pspicture}\\ \hline
action=draw & action=draw* & action=draw**\\ \hline
\end{tabular} 
\bigskip

\subsubsection{icosahedron}
\begin{tabular}{|c|c|c|} \hline 
 \multicolumn{3}{|c|}{ \BS{psSolid}[\RDD{object=icosahedron},a=3,action=draw] \RDI{object=icosahedron}{pst-sol3d} } \\  \hline 
\begin{pspicture}(-10,-10)(10,10)
 \psframe(-10,-10)(10,10)
\psSolid[object=icosahedron,a=3,action=draw]%
\end{pspicture}
&
\begin{pspicture}(-10,-10)(10,10)
 \psframe(-10,-10)(10,10)
\psSolid[object=icosahedron,a=3,action=draw*]%
\end{pspicture}
&
\begin{pspicture}(-10,-10)(10,10)
 \psframe(-10,-10)(10,10)

\psSolid[object=icosahedron,a=3,action=draw**]%
\end{pspicture}\\ \hline
action=draw & action=draw* & action=draw**\\ \hline
\end{tabular} 
\bigskip


%\subsubsection{prisme}
\SbSbSSCT{Prisme}{prism}

\begin{tabular}{|c|c|c|} \hline 
 \multicolumn{3}{|c|}{ \BS{psSolid}[\RDD{object=prisme},action=draw,h=4] \RDI{object=prisme}{pst-sol3d} } \\  \hline 
\begin{pspicture}(-10,-5)(10,10)
 \psframe(-10,-5)(10,10)
\psSolid[object=prisme,action=draw,h=2]%
\end{pspicture}
&
\begin{pspicture}(-10,-5)(10,10)
 \psframe(-10,-5)(10,10)
\psSolid[object=prisme,action=draw*,h=2]%
\end{pspicture}
&
\begin{pspicture}(-10,-5)(10,10)
 \psframe(-10,-5)(10,10)
\psSolid[object=prisme,action=draw**,h=2]%
\end{pspicture}\\ \hline
action=draw & action=draw* & action=draw**\\ \hline
\end{tabular} 

%\subsubsection{prismecreux}
\SbSbSSCT{Prisme creux}{Empty prism}

\begin{tabular}{|c|c|c|} \hline 
 \multicolumn{3}{|c|}{ \BS{psSolid}[\RDD{object=prismecreux},action=draw,h=4] \RDI{object=prismecreux}{pst-sol3d} } \\  \hline 
\begin{pspicture}(-10,-5)(10,10)
 \psframe(-10,-5)(10,10)
\psSolid[object=prismecreux,action=draw,h=2]%
\end{pspicture}
&
\begin{pspicture}(-10,-5)(10,10)
 \psframe(-10,-5)(10,10)
\psSolid[object=prismecreux,action=draw*,h=2]%
\end{pspicture}
&
\begin{pspicture}(-10,-5)(10,10)
 \psframe(-10,-5)(10,10)
\psSolid[object=prismecreux,action=draw**,h=2]%
\end{pspicture}\\ \hline
action=draw & action=draw* & action=draw**\\ \hline
\end{tabular} 
\bigskip

\subsubsection{face,ruban}

\begin{tabular}{|c|c|c|} \hline 
\begin{pspicture}(-10,-10)(10,10)
 \psframe(-10,-10)(10,10)
\psSolid[object=face,fillcolor=yellow,incolor=blue,base=0 0 3 0 1.5 3](0,1,0)
\end{pspicture}
&
\begin{pspicture}(-10,-10)(10,10)
 \psframe(-10,-10)(10,10)
\psSolid[object=ruban,h=3,base=0 0 2 2 4 0 6 2,num=0 1 2 3,show=0 1 2 3,ngrid=3])
\end{pspicture}
&
\begin{pspicture}(-10,-10)(10,10)
 \psframe(-10,-10)(10,10)
 \axesIIID(1,1,1)(4,7,4)
\psSolid[object=ruban,h=3,base=0 0 2 2 4 0 6 2])
\end{pspicture}
\\ \hline
\end{tabular} 

%============================================== A VOIR ========================================

%\newpage
%
%\begin{pspicture}(-10,-10)(10,10)
% \psframe(-10,-10)(10,10)
%\psSurface[ngrid=.25 .25,incolor=white,axesboxed](-4,-4)(4,4){x dup mul y dup mul 3 mul sub x mul 32 div}
%\end{pspicture}
%
%\newpage
%
%\begin{pspicture}(-10,-10)(10,10)
%\psframe(-10,-10)(10,10)
%\psSolid[object=new,action=draw,sommets=
%-3 0 0
%3 0 0
%0 -3 0
%0 3 0
%0 0 -3
%0 0 3 ,
%faces={
%[0 2 1 3]
%[0 4 1 5]}]%
%\end{pspicture}
%
%\begin{pspicture}(-10,-10)(10,10)
%\psframe(-10,-10)(10,10)
%\psSolid[object=new,action=draw,sommets=
%-3 0 0
%3 0 0
%0 -3 0
%0 3 0
%0 0 -3
%0 0 3 ,
%faces={
%[0 2 1 3]
%[0 4 1 5]}
%,show=all,num=all
%]%
%\end{pspicture}
%\newpage
%
%\begin{pspicture}(-10,-10)(10,10)
%\psframe(-10,-10)(10,10)
%\psSolid[object=new,action=draw,sommets=
%2 4 3
%-2 4 3
%-2 -4 3
%2 -4 3
%2 4 0
%-2 4 0
%-2 -4 0
%2 -4 0
%0 4 5
%0 -4 5,
%faces={
%[0 1 2 3]
%[7 6 5 4]
%[0 3 7 4]
%[3 9 2]
%[1 8 0]
%[8 9 3 0]
%[9 8 1 2]
%[6 7 3 2]
%[2 1 5 6]}]%
%\end{pspicture}
%
%\newpage
%\begin{pspicture}(-10,-10)(10,10)
%\psframe(-10,-10)(10,10)
%\defFunction[algebraic]{helice}(t){3*cos(4*t)}{3*sin(4*t)}{t}
%\psSolid[object=courbe,r=0,range=0 6,linecolor=blue,linewidth=0.1,resolution=360,function=helice,fillcolor=red]
%\end{pspicture}
%%
%
%
%\newpage
%\psset{unit=1cm}
%
%
%
%\begin{pspicture}(-2,-2)(3,3)
%\psframe(-2cm,-2cm)(3cm,3cm)
%\psset{viewpoint=100 30 20,Decran=100}
%\psSolid[object=cube,a=2,action=draw*, fillcolor=magenta!20]
%\axesIIID[showOrigin=false](1,1,1)(3,2,2.5)
%\end{pspicture}
%\newpage
\subsection{Mode}

\begin{tabular}{|c|c|c|c|} \hline 
 \multicolumn{4}{|c|}{ \BS{psSolid}[object=cylindre,h=3,r=1.5,\RDD{mode}=1](0,0,0) \RDI{mode}{pst-sol3d} } \\  \hline 
\begin{pspicture}(-7,-5)(7,15)
% \psframe(-7,-5)(7,15)
\psSolid[object=cylindre,h=3,r=1.5,mode=1](0,0,0)%
 \axesIIID[linecolor=red](0,0,0)(1,1,1)
\end{pspicture}
&
\begin{pspicture}(-7,-5)(7,15)
% \psframe(-7,-5)(7,15)
\psSolid[object=cylindre,h=3,r=1.5,mode=2](0,0,0)%
 \axesIIID[linecolor=red](0,0,0)(1,1,1)
\end{pspicture}
&
\begin{pspicture}(-7,-5)(7,15)
% \psframe(-7,-5)(7,15)
\psSolid[object=cylindre,h=3,r=1.5,mode=3](0,0,0)%
 \axesIIID[linecolor=red](0,0,0)(1,1,1)
\end{pspicture}
&
\begin{pspicture}(-7,-5)(7,15)
% \psframe(-7,-5)(7,15)
\psSolid[object=cylindre,h=3,r=1.5,mode=4](0,0,0)%
 \axesIIID[linecolor=red](0,0,0)(1,1,1)
\end{pspicture}\\ \hline
mode=1 & mode=2 & mode=3 & mode=4\\ \hline
\end{tabular}
%-------------------------------------------------

\subsubsection{Options}

\begin{tabular}{|c|c|c|} \hline 
 \multicolumn{3}{|c|}{ \BS{psSolid}[object=cube,a=3,action=draw*,\RDD{trunc}=all,RotZ=30] \RDI{trunc}{pst-sol3d} } \\  \hline 
\begin{pspicture}(-10,-10)(10,10)
\psSolid[object=cube,a=3,action=draw*,trunc=all,RotZ=30] %cube
\end{pspicture}
&
\begin{pspicture}(-10,-10)(10,10)
\psSolid[object=cube,a=3,action=draw*,trunc=0  2 4,RotZ=30] %cube
\end{pspicture}
 & 
\begin{pspicture}(-10,-10)(10,10)
\psSolid[object=cube,a=3,action=draw*,trunccoeff=.5,trunc=all,RotZ=30] %cube
\end{pspicture}
\\ \hline
\RDD{trunc}=all & \RDD{trunc}=0  2 4  & \RDD{trunccoeff}=.5 \RDI{trunccoeff}{pst-sol3d}\\ \hline
\end{tabular} 

\bigskip

\begin{tabular}{|c|c|c|} \hline 
 \multicolumn{3}{|c|}{ \BS{psSolid}[object=cube,a=3,action=draw,\RDD{chanfrein},RotZ=30] \RDI{chanfrein}{pst-sol3d}} \\  \hline 
\begin{pspicture}(-10,-10)(10,10)
\psSolid[object=cube,a=3,action=draw*,chanfrein,RotZ=30] %cube
\end{pspicture}
&
\begin{pspicture}(-10,-10)(10,10)
\psSolid[object=cube,a=3,action=draw*,chanfrein,chanfreincoeff=.2,RotZ=30] %cube
\end{pspicture}
 & 
\begin{pspicture}(-10,-10)(10,10)
\psSolid[object=cube,a=3,action=draw*,chanfreincoeff=.5,chanfrein,RotZ=30] %cube
\end{pspicture}
\\ \hline
\RDD{chanfrein} & chanfrein,\RDD{chanfreincoeff}=.2 \RDI{chanfreincoeff}{pst-sol3d} & chanfrein,\RDD{chanfreincoeff}=.5 \\ \hline
\end{tabular} 
\bigskip

\bigskip

\begin{tabular}{|c|c|c|} \hline 
 \multicolumn{3}{|c|}{ \BS{psSolid}[object=cube,a=3,action=draw**,hollow,affinage=0,RotZ=30] } \\  \hline 
\begin{pspicture}(-10,-10)(10,10)
\psSolid[object=cube,a=3,action=draw**,hollow,affinage=0,RotZ=30] %cube
\end{pspicture}
&
\begin{pspicture}(-10,-10)(10,10)
\psSolid[object=cube,a=3,action=draw**,hollow,affinage=3 4,RotZ=30,] %cube
\end{pspicture}
 & 
\begin{pspicture}(-10,-10)(10,10)
\psSolid[object=cube,a=3,action=draw**,hollow,affinage=all,RotZ=30,] %cube
\end{pspicture}
\\ \hline
\RDD{hollow} \RDI{hollow}{pst-sol3d},\RDD{affinage}=3 \RDI{affinage}{pst-sol3d}& \RDD{hollow},,\RDD{affinage}=3 4  & \RDD{hollow},\RDD{affinage}=all \\ \hline
\end{tabular} 
\bigskip

%========================================================================
\subsection{Positionnement}
\psset{unit=.25cm}

\begin{tabular}{|c|c|c|} \hline 
 \multicolumn{3}{|c|}{ \BS{psSolid}[\RDD{object=parallelepiped},a=1,b=2,c=3,action=draw](1 0 0) \RDI{object=parallelepiped}{pst-sol3d}} \\  \hline 
\begin{pspicture}(-10,-10)(10,10)
 \psframe(-10,-10)(10,10)
 \psSolid[object=parallelepiped,a=1,b=2,c=3,action=draw,linestyle=dotted]% 
\psSolid[object=parallelepiped,a=1,b=2,c=3,action=draw](1 0 0)%
\axesIIID[linecolor=red](0,0,0)(1,1,1)
\end{pspicture}
&
\begin{pspicture}(-10,-10)(10,10)
 \psframe(-10,-10)(10,10)
 \psSolid[object=parallelepiped,a=1,b=2,c=3,action=draw,linestyle=dotted]% 
\psSolid[object=parallelepiped,a=1,b=2,c=3,action=draw](0 1 0)%
\axesIIID[linecolor=red](0,0,0)(1,1,1)
\end{pspicture}
&
\begin{pspicture}(-10,-10)(10,10)
 \psframe(-10,-10)(10,10)
 \psSolid[object=parallelepiped,a=1,b=2,c=3,action=draw,linestyle=dotted]% 
\psSolid[object=parallelepiped,a=1,b=2,c=3,action=draw](0 0 1)%
\axesIIID[linecolor=red](0,0,0)(1,1,1)
\end{pspicture}\\ \hline
(1 0 0) & (0 1 0) & (0 0 1)\\ \hline
\end{tabular} 
\bigskip

\begin{tabular}{|c|c|c|} \hline 
 \multicolumn{3}{|c|}{ \BS{psSolid}[\RDD{object=parallelepiped},a=1,b=2,c=3,action=draw] } \\  \hline 
\begin{pspicture}(-10,-10)(10,10)
 \psframe(-10,-10)(10,10)
 \psSolid[object=parallelepiped,a=1,b=2,c=3,action=draw,linestyle=dotted]% 
\psSolid[object=parallelepiped,a=1,b=2,c=3,action=draw,RotX=30]%
\axesIIID[linecolor=red](0,0,0)(1,1,1)
\end{pspicture}
&
\begin{pspicture}(-10,-10)(10,10)
 \psframe(-10,-10)(10,10)
\psSolid[object=parallelepiped,a=1,b=2,c=3,action=draw,linestyle=dotted]% 
\psSolid[object=parallelepiped,a=1,b=2,c=3,action=draw,RotY=30]%
\axesIIID[linecolor=red](0,0,0)(1,1,1)
\end{pspicture}
&
\begin{pspicture}(-10,-10)(10,10)
 \psframe(-10,-10)(10,10)
 \psSolid[object=parallelepiped,a=1,b=2,c=3,action=draw,linestyle=dotted]% 
\psSolid[object=parallelepiped,a=1,b=2,c=3,action=draw,RotZ=30]%
\axesIIID[linecolor=red](0,0,0)(1,1,1)
\end{pspicture}\\ \hline
RotX=30 & RotY=30 & RotZ=30\\ \hline
\end{tabular}

%\subsection{title}
%
%\psset{unit=.25cm}
%
%\begin{tabular}{|c|c|c|} \hline 
% \multicolumn{3}{|c|}{ \BS{psSolid}[\RDD{object=parallelepiped},a=1,b=2,c=3,action=draw](1 0 0) } \\  \hline 
%\begin{pspicture}(-10,-10)(10,10)
% \psframe(-10,-10)(10,10)
%
%\psSolid[object=parallelepiped,a=1,b=2,c=3,action=draw,show=all]%
%\axesIIID[linecolor=red](0,0,0)(1,1,1)
%\end{pspicture}
%&
%\begin{pspicture}(-10,-10)(10,10)
% \psframe(-10,-10)(10,10)
%
%\psSolid[object=parallelepiped,a=1,b=2,c=3,action=draw,num=all]%
%\axesIIID[linecolor=red](0,0,0)(1,1,1)
%\end{pspicture}
%&
%\begin{pspicture}(-10,-10)(10,10)
% \psframe(-10,-10)(10,10)
%\psSolid[object=parallelepiped,a=1,b=2,c=3,action=draw,show=all,num=all]%
%\axesIIID[linecolor=red](0,0,0)(1,1,1)
%\end{pspicture}\\ \hline
%show=all & num=all & show=all,num=all \\ \hline
%\end{tabular} 
%\bigskip
%
%\begin{tabular}{|c|c|c|} \hline 
% \multicolumn{3}{|c|}{ \BS{psSolid}[\RDD{object=parallelepiped},a=1,b=2,c=3,action=draw](1 0 0) } \\  \hline 
%\begin{pspicture}(-10,-10)(10,10)
% \psframe(-10,-10)(10,10)
%
%\psSolid[object=parallelepiped,a=1,b=2,c=3,action=draw,show=0 1 2 3]%
%\axesIIID[linecolor=red](0,0,0)(1,1,1)
%\end{pspicture}
%&
%\begin{pspicture}(-10,-10)(10,10)
% \psframe(-10,-10)(10,10)
%
%\psSolid[object=parallelepiped,a=1,b=2,c=3,action=draw,num=0 1 2 3]%
%\axesIIID[linecolor=red](0,0,0)(1,1,1)
%\end{pspicture}
%&
%\begin{pspicture}(-10,-10)(10,10)
% \psframe(-10,-10)(10,10)
%\psSolid[object=parallelepiped,a=1,b=2,c=3,action=draw,show=0 1 2 3,num=0 1 2 3]%
%\axesIIID[linecolor=red](0,0,0)(1,1,1)
%\end{pspicture}\\ \hline
%show=0 1 2 3 & num=0 1 2 3 & show=0 1 2 3,num=0 1 2 3 \\ \hline
%\end{tabular} 

%\bigskip
%\psset{unit=.5cm}
%\begin{pspicture}(-3,-2.5)(7,2.5)
%%\psset{viewpoint=50 20 20 rtp2xyz,Decran=40}
%\psSolid[action=draw,object=cube,RotZ=30,show=0 1 2 3,num=0 1 2 3]%
%\end{pspicture}
%
%\begin{pspicture}(-3,-2.5)(7,2.5)
%%\psset{viewpoint=50 20 20 rtp2xyz,Decran=40}
%\psSolid[action=draw,object=parallelepiped,RotZ=30,show=0 1 2 3,num=0 1 2 3]%
%\end{pspicture}


%\subsection{Coloriage numérotation}
\SbSSCT{Coloriage numérotation}coloring and numbering{}

\psset{unit=.15cm}






\begin{tabular}{|c|c|c|}\hline
 \multicolumn{3}{|c|}{ \BS{psSolid}[\RDD{fcol}=0 (green) 1 (red) 4 (cyan) 13 (blue) 40 (black), object=cube,mode=3] \RDI{fcol}{pst-sol3d}} 
 \\  \hline 
\begin{pspicture}(-10,-10)(10,10)  
 \psSolid[fcol=0 (green) 1 (red)  4 (cyan) 13 (blue) 40 (black), object=cube,mode=3]%
 \end{pspicture}
  &  
 \begin{pspicture}(-10,-10)(10,10)
  \psSolid[numfaces=all, object=cube,mode=3]%
  \end{pspicture}  
  & 
  \begin{pspicture}(-10,-10)(10,10)
   \psSolid[numfaces=0 1 2 3, object=cube,mode=3]%
   \end{pspicture} 
   \\ \hline  
\RDD{fcol}=0 (green) 1 (red)  ...  &  \RDD{numfaces}=all  \RDI{numfaces}{pst-sol3d}& \RDD{numfaces}=0 1 2 3
    \\ \hline 
\end{tabular} 
\bigskip

\begin{tabular}{|c|c|c|}\hline
 \multicolumn{3}{|c|}{ \BS{psSolid}[\RDD{fcol}=0 (green) 1 (red) 2 (cyan) 3 (magenta), object=parallelepiped,mode=3] } 
 \\  \hline 
\begin{pspicture}(-10,-10)(10,10)
 \psSolid[fcol=0 (green) 1 (red) 2 (cyan) 3 (magenta), object=parallelepiped,mode=3]%
\end{pspicture} 
&
  \begin{pspicture}(-10,-5)(10,15)
   \psSolid[fcol=0 (green) 1 (red) 2 (cyan) 3 (magenta), object=cylindre,h=4,ngrid=4,mode=3]%
   \end{pspicture}
   & 
\begin{pspicture}(-10,-10)(10,10)
\psSolid[fcol=0 (green) 1 (red) 2 (cyan) 3 (magenta) , object=tore,r1=2,r0=1,ngrid=4,mode=3]%
   \end{pspicture} 
   \\ \hline
\RDD{fcol}= 0 (green)  1 (red) ... & \RDD{numfaces}=all & \RDD{numfaces}=0 1
   \\ \hline

\end{tabular}  
 
 \newpage
 
% \subsection{Dans une prochaine version}
\SbSSCT{Dans une prochaine version}{In a future version}

 
%\subsubsection{Surface d'après une équation}
\SbSbSSCT{Surface d'après une équation}{Equation define surface}

\begin{tabular}{|c|c|c|} \hline 
%\psset{llx=-.5cm,lly=-.5cm,urx=.5cm,ury=0.5cm,fillstyle=none,linewidth=2pt}
\psset{fillcolor=yellow,fillstyle=none,linecolor=blue,unit=.4cm}

% \psset{unit=0.5}
% \psset{lightsrc=30 30 25}
 \psset{viewpoint=50 40 30 rtp2xyz,Decran=50}
\begin{pspicture}(-7,-8)(7,8)
 \psSurface[algebraic,ngrid=.25 .25,hue=0 1](-6,-6)(6,6){sin(x) *cos(y)}
\end{pspicture}
   \\ \hline
 \BSS{psSurface}[algebraic,ngrid=.25 .25,hue=0 1](-6,-6)(6,6)\AC{sin(x) *cos(y)} \BSI{psSurface}{pst-sol3d}
   \\ \hline
\end{tabular}

%\subsubsection{Fusion de 2 solides}
\SbSbSSCT{Fusion de 2 solides}Fusion of two solids{}
  \psset{unit=1cm}
\begin{center} 
%  \begin{pspicture}(-1,1)(10,8)
%  \psframe(-1,-1)(10,10)
 \begin{animateinline}[poster=first, palindrome,autoplay]{5}%
 \multiframe{30}{iAngle=10+30}{\mydessin{\iAngle} }
 \end{animateinline}
%  \end{pspicture}
 
% \begin{pspicture}(-.3\linewidth,-12)(.35\linewidth,30)
%\psframe(-.3\linewidth,-12)(.35\linewidth,30)
%\psset{unit=.15cm} 
%\psset{solidmemory}
%\psSolid[object=cylindrecreux,h=10,r=2,fillcolor=white,mode=4,name=A1,incolor=green!50](0,0,-3)
%\psSolid[object=conecreux,h=15,r=2,RotY=-60,fillcolor=white,incolor=red!50,mode=5,name=B1](4,0,0)%
%\psSolid[object=fusion,action=draw**,base=A1 B1,](0,0,0)
%\composeSolid
% \end{pspicture}

 \psset{unit=1cm}
\bigskip 
 \begin{tabular}{|l|} \hline 
\BSS{psset}\AC{solidmemory}\\
\\
\BS{psSolid}[object=cylindrecreux,h=10,r=2,fillcolor=white,mode=4,name=A1,incolor=green!50](0,0,-3)\\
\BS{psSolid}[object=conecreux,h=15,r=2,RotY=-60,fillcolor=white,incolor=red!50,mode=5,name=B1](4,0,0)\\
\BS{psSolid}[object=fusion,action=draw**,base=A1 B1,](0,0,0)\\
\BSS{composeSolid} \BSI{composeSolid}{pst-sol3d}
    \\ \hline
 \end{tabular}
\end{center} 
 

 \begin{center}

 \end{center} 





\newpage

\appendix

%\section{formules en langage postcript}
\SSCT{formules en langage postcript}
{Formula in postcript}
\label{postcript}

\begin{tabular}{|c|c|c|} \hline  
formule & en PostScript & valeur \\ 
\hline  
\rule[-0.2cm]{0pt}{.7cm} $2 + 3$ & 2 3 add & \makebox[2cm][c]{\psPrintValue {2 3 add}} \\ \hline
\rule[-0.2cm]{0pt}{.7cm}  $2 + 2$ & 2 dup add & \makebox[2cm][c]{\psPrintValue {2 dup add}} \\ \hline
\rule[-0.2cm]{0pt}{.7cm}  $2 - 3$ & 2 -3 add & \makebox[2cm][c]{\psPrintValue {2 -3 add}} \\ \hline 
\rule[-0.2cm]{0pt}{.7cm}$2 * 3$ & 2 3 mul & \makebox[2cm][c]{\psPrintValue {2 3 mul}} \\ \hline 
\rule[-0.2cm]{0pt}{.7cm}$10/2$ & 10 2 div & \makebox[2cm][c]{\psPrintValue {10 2 div}} \\ \hline   
\rule[-0.2cm]{0pt}{.7cm}$3^2$ & 3 2 exp & \makebox[2cm][c]{\psPrintValue {3 2 exp}} \\ \hline  
 \rule[-0.2cm]{0pt}{.7cm}$\sqrt{3}$ & 3 sqrt & \makebox[2cm][c]{\psPrintValue[decimals=2] {3 sqrt}} \\ \hline 
\rule[-0.2cm]{0pt}{.7cm}$\sin({30})$ & 30 sin & \makebox[2cm][c]{\psPrintValue {30 sin}} \\ \hline 
\rule[-0.2cm]{0pt}{.7cm}$\cos({30})$ & 30 cos & \makebox[2cm][c]{\psPrintValue[decimals=2] {30 cos}} \\ \hline 
\rule[-0.2cm]{0pt}{.7cm}$sin^2({30})$ & 30 sin 2 exp & \makebox[2cm][c]{\psPrintValue {30 sin 2 exp}} \\ \hline
\rule[-0.2cm]{0pt}{.7cm}$sin({5}^2)$ & 5 2 exp sin  & \makebox[2cm][c]{\psPrintValue[decimals=2] {5 2 exp sin}} \\ \hline 
\end{tabular}

\newpage

%%%=======================================================================
%\section{Les modules étudiés dans ce document}
\SSCT{Les modules étudiés dans ce document}{Packages studied in this document}

\textbf{Modules chargés automatiquement avec le module pst-all}

\begin{tabular}{|c|c|c|c|}\hline 
name				&  page 				& documentation  \footnotemark[1]\\ \hline 
pst-user			& les bases				& \cite{pst-user} \\
% \cite{pstricks}\\ %& pstricks-doc.pdf 
pstricks-add 		& les additifs		 	& \cite{pstricks-add} \\ %& pstricks-add-doc.pdf
pst-node 			& \pageref{pst-node} 	& \cite{pst-node} \\ %& pst-node-doc.pdf
xcolor				& \pageref{xcolor} 		& \cite{xcolor} \\ %& xcolor.pdf
pst-coil 			& \pageref{pst-coil}	& \cite{pst-coil}\\ %& pst-coil-doc.pdf
pst-grad			& \pageref{pst-grad}	& \cite{pst-grad} \\ %& pst-grad-doc.pdf
pst-fill 			& \pageref{pst-fill}	& \cite{pst-fill}\\ %& pst-fill.pdf
pst-text			& \pageref{pst-text} 	& \cite{pst-text} \\ %& pst-text-doc.pdf
pst-plot 			& \pageref{pst-plot}	& \cite{pst-plot} \\ %& pst-plot-doc.pdf

multido 			& \pageref{multido} 	&\cite{multido} \\ %& multido-doc.pdf
pst-tree			& \pageref{pst-tree}	& \cite{pst-tree} \\ %& pst-tree-doc.pdf
pst-3d 				&		 				& \\ %& pst-3d-doc.pdf
pst-eps 			&  						& \\ %& pst-eps-doc.pdf
\hline 
\end{tabular}  

\bigskip
\textbf{Autres modules}

%
\begin{tabular}{|c|c|c|c|}\hline 
nom 			& voir page 				& documentation\footnotemark[1] \\  \hline
pst-poly		& \pageref{poly} 			& \cite{pst-poly} 	\\ %& pst-poly-doc
pst-bezier		& \pageref{bezier} 			& \cite{pst-bezier} 	\\


pst-fr3d 		& \pageref{pst-fr3d}		& \cite{pst-fr3d}\\ %& pst-fr3d.pdf
pst-slpe 		& \pageref{pst-slpe}		& \cite{pst-slpe}\\  %& pst-slpe.pdf
pst-fun			& \pageref{fun} 			& \cite{pst-fun} 	\\
pst-func			& \pageref{func} 			& \cite{pst-func} 	\\
infix-RPN 		& \pageref{infix-RPN}		& \cite{pst-infix}\\ %& pst-infixplot.pdf
pst-infixplot 	& \pageref{pst-infixplot}  	& \cite{pst-infix}\\ %& pst-infixplot.pdf 
pst-eucl 		& \pageref{pst-eucl} 		& \cite{pst-eucl}\\ %& pst-eucl-doc.pdf
animate 		& \pageref{anim} 			& \cite{anim}	\\ %& animate.pdf
pst-3dplot		& \pageref{3dplot} 			& \cite{pst-3dplot}\\ %& pst-3dplot-doc
pst-solides3d	& \pageref{sol3d} 			& \cite{pst-sol3d}\\ %& pst-solides3d-doc.pdf
\hline 
\end{tabular}
%\bibitem{pst-fun} pst-fun-doc.pdf   \hspace{1cm}	version 0.04 \hspace{1cm}	11 pages 	\hspace{1cm} 	\DGB 
%\bibitem{pst-bezier} pst-bezier-doc.pdf  \hspace{1cm}	version 0.01 \hspace{1cm} 	10 pages 	\hspace{1cm}	\DGB

\bigskip
\textbf{Additifs annuels}

%
\begin{tabular}{|c|c|}\hline 
année						& documentation\footnotemark[1]  \\  \hline 
2005 			& \cite{pst-news05} \\ %pst-news5.pdf	
2008  			& \cite{pst-news08} \\ %pst-news08.pdf \\ 
2010 			& \cite{pst-news10} \\ %pst-news10.pdf 	\\
2013 			& \cite{pst-news13}\\ % pst-news13.pdf 	\\
\hline 
\end{tabular} 
 

\footnotetext[1]{Vous pouvez les trouver pour la distribution Texlive dans le répertoire :  \BS{}texlive\BS{}2011\BS{}tesmf-dist\BS{}doc\BS{}generic}



\newpage

\section{Sources}
\begin{thebibliography}{99}
\bibitem{pst-user} pst-user.pdf  \hspace{1cm}	version 1.51  \hspace{1cm} 	131 pages  \hspace{1cm} 	\DGB
%\bibitem{pstricks} pstricks-doc.pdf  \hspace{1cm}	version 97  \hspace{1cm} 	338 pages  \hspace{1cm} 	\DGB
%\bibitem{pst-add}
\bibitem{pstricks-add}    pstricks-add-doc.pdf  \hspace{1cm}	version 3.61 \hspace{1cm} 	134 pages 	\hspace{1cm}	\DGB
\bibitem{pst-3dplot} pst-3dplot-doc.pdf  \hspace{1cm}	version 1.94 \hspace{1cm} 	69 pages \hspace{1cm}	\DGB
\bibitem{pst-bezier} pst-bezier-doc.pdf  \hspace{1cm}	version 0.01 \hspace{1cm} 	10 pages 	\hspace{1cm}	\DGB
\bibitem{pst-coil} pst-coil-doc.pdf  \hspace{1cm}	version 1.06 \hspace{1cm} 	14 pages 	\hspace{1cm}	\DGB
\bibitem{pst-eucl} pst-eucl-doc.pdf \hspace{1cm}	version 1.51 \hspace{1cm}	52 pages 	\hspace{1cm}	\DGB
\bibitem{pst-fill} pst-fill.pdf     \hspace{1cm}	version 1.00 \hspace{1cm}	37 pages 	\hspace{1cm}  	\DGB
\bibitem{pst-fr3d} pst-fr3d.pdf   \hspace{1cm}	version 1.00 \hspace{1cm}	10 pages 	\hspace{1cm} 	\DGB
\bibitem{pst-fun} pst-fun-doc.pdf   \hspace{1cm}	version 0.04 \hspace{1cm}	11 pages 	\hspace{1cm} 	\DGB
\bibitem{pst-func} pst-func-doc.pdf   \hspace{1cm}	version 0.81 \hspace{1cm}	73 pages 	\hspace{1cm} 	\DGB
\bibitem{pst-grad} pst-grad-doc.pdf \hspace{1cm}	version 1.06 \hspace{1cm}	11 pages 	\hspace{1cm}	\DGB
\bibitem{pst-infix} pst-infixplot.pdf  \hspace{1cm}	version 0.11 \hspace{1cm} 	2 pages 	\hspace{1cm}	\DGB
\bibitem{pst-news05} pst-news05.pdf 	 \hspace{1cm} 	11 pages 	\hspace{1cm}	\DGB
\bibitem{pst-news08} pst-news08.pdf 	\hspace{1cm} 	30 pages 	\hspace{1cm}	\DGB
\bibitem{pst-news10} pst-news10.pdf \hspace{1cm}	28 pages 	\hspace{1cm}	\DGB
\bibitem{pst-news13} pst-news10.pdf \hspace{1cm}	9 pages 	\hspace{1cm}	\DGB
\bibitem{pst-node} pst-node-doc.pdf  \hspace{1cm}	version 1.30 :	53 pages 	\hspace{1cm}  	\DGB
\bibitem{pst-plot} pst-plot-doc.pdf  \hspace{1cm}	version 1.40 :	92 pages 	\hspace{1cm} 	\DGB
\bibitem{pst-poly} pst-poly-doc.pdf  \hspace{1cm}	version 1.61 :	22 pages 	\hspace{1cm} 	\DGB
\bibitem{pst-slpe} pst-slpe.pdf   \hspace{1cm}	version 1.31 \hspace{1cm}	16 pages 	\hspace{1cm}  	\DGB
\bibitem{pst-sol3d} pst-solides3d-doc.pdf   \hspace{1cm}	version v. 4.24 \hspace{1cm}	197 pages 	\hspace{1cm}  	\DGB
\bibitem{pst-text} pst-text-doc.pdf   \hspace{1cm}	version 1.06\hspace{1cm}	11 pages 	\hspace{1cm}  	\DGB
\bibitem{pst-tree} pst-tree-doc.pdf   \hspace{1cm}	version 1.12\hspace{1cm}	24 pages 	\hspace{1cm} 	\DGB
\bibitem{multido} multido-doc.pdf   \hspace{1cm}	version 1.42\hspace{1cm}	4 pages 	\hspace{1cm}  	\DGB
\bibitem{xcolor} xcolor.pdf   		\hspace{1cm}	version 2.11 \hspace{1cm}	65 pages 	\hspace{1cm}  	\DGB
\bibitem{anim} animate.pdf   	\hspace{1cm}	6th December 2012 \hspace{1cm}	23 pages \hspace{1cm}  	\DGB
\end{thebibliography}

\newpage 

\section{Index}

 \printindex 
\end{document}